\documentclass{article}

% Language setting
% Replace `english' with e.g. `spanish' to change the document language
\usepackage[french]{babel}

\usepackage{float}

% Set page size and margins
% Replace `letterpaper' with `a4paper' for UK/EU standard size
\usepackage[a4paper,top=2cm,bottom=2cm,left=3cm,right=3cm,marginparwidth=1.75cm]{geometry}

% Useful packages
\usepackage{siunitx}
\usepackage{amsmath}
\usepackage{pgfplots}
\pgfplotsset{compat=newest}
\usetikzlibrary{plotmarks}
\usetikzlibrary{arrows.meta}
\usepgfplotslibrary{patchplots}
\usepackage{grffile}
\pgfplotsset{plot coordinates/math parser=false}
\newlength\figureheight
\newlength\figurewidth
  
\usepackage{graphicx}
\usepackage{hyperref}
\newcommand{\TF}{\operatorname{TF}}
\newcommand{\sinc}{\operatorname{sinc}}


\includegraphics[width=0.2\textwidth]{n7.png}


\title{Rapport de projet -- Traitement du signal}
\author{Florent Puy, Ewen Le Bihan}

\date{ENSEEIHT, département Sciences du Numérique}

\begin{document}
\maketitle



\tableofcontents

\section{Introduction}

Dans ce projet, nous avons eu à implémenter un modem suivant les règles V21 de l'union internationnale des télécommunications (UIT) en Matlab. Nous utiliserons la méthode de la modulation en fréquence numérique.

\section{Modem en fréquence}

Nous allons tout d'abord réaliser un signal NRZ binaire à partir duquel nous construirons ensuite un signal sinusoïdal modulé fréquence $F_0=\SI{1180}{\hertz}$ pour les bits 0 et de fréquence $F_1=\SI{980}{\hertz}$ pour les bits 1.
Nous comparerons ensuite les densités spéctrales de puissance théorique et pratique du signal NRZ et du signal modulé en fréquence.

\subsection{Génération d'un signal NRZ}

On génère tout d'abord un signal NRZ prennant deux valeurs, 0 ou 1, générées aléatoirement d'une durée $T_s=1/300 s$. On effectue cela sur $N_s$ périodes. Voici les résultats ainsi obtenus.

%%%% This file was created by matlab2tikz.
%
%The latest updates can be retrieved from
%  http://www.mathworks.com/matlabcentral/fileexchange/22022-matlab2tikz-matlab2tikz
%where you can also make suggestions and rate matlab2tikz.
%
\definecolor{mycolor1}{rgb}{0.00000,0.44700,0.74100}%
%
\begin{tikzpicture}

\begin{axis}[%
width=4.521in,
height=3.559in,
at={(0.758in,0.488in)},
scale only axis,
xmin=0,
xmax=4000,
xlabel style={font=\color{white!15!black}},
xlabel={temps [s]},
ymin=-0.1,
ymax=1.1,
ylabel style={font=\color{white!15!black}},
ylabel={$\text{\$m}_\text{i}\text{(t)\$}$},
axis background/.style={fill=white},
title style={font=\bfseries},
title={Signal NRZ aléatoire}
]
\addplot [color=mycolor1, forget plot]
  table[row sep=crcr]{%
1	0\\
160	0\\
161	1\\
480	1\\
481	0\\
800	0\\
801	1\\
1280	1\\
1281	0\\
1440	0\\
1441	1\\
1760	1\\
1761	0\\
1920	0\\
1921	1\\
2080	1\\
2081	0\\
2400	0\\
2401	1\\
3040	1\\
3041	0\\
3200	0\\
3201	1\\
4000	1\\
};
\end{axis}
\end{tikzpicture}%

On calcule ensuite la densité spectrale de puissance de ce signal NRZ en utilisant la fonction $pwelch$ de Matlab utilisant un périodogramme de Welch.

On calcule ensuite la densité spectrale théorique vue en cours d'un signal NRZ:
\[
S_\text{NRZ}(f)=\frac{1}{4} T_s \sinc^2(\pi f T_s)+\frac{1}{4} \delta(f)
\]

On peut déormais comparer la densité spéctrale de puissance théorique à celle calculée:

%%%% This file was created by matlab2tikz.
%
%The latest updates can be retrieved from
%  http://www.mathworks.com/matlabcentral/fileexchange/22022-matlab2tikz-matlab2tikz
%where you can also make suggestions and rate matlab2tikz.
%
\definecolor{mycolor1}{rgb}{0.00000,0.44700,0.74100}%
%
\begin{tikzpicture}

\begin{axis}[%
width=4.521in,
height=3.548in,
at={(0.758in,0.499in)},
scale only axis,
xmin=-25000,
xmax=25000,
xlabel style={font=\color{white!15!black}},
xlabel={fréquence [Hz]},
ymode=log,
ymin=1e-40,
ymax=1,
yminorticks=true,
ylabel style={font=\color{white!15!black}},
ylabel={Densité spectrale de puissance},
axis background/.style={fill=white},
title style={font=\bfseries},
title={Densité spectrale de puissance},
legend style={legend cell align=left, align=left, draw=white!15!black}
]
\addplot [color=mycolor1]
  table[row sep=crcr]{%
-24000	1.63292967995806e-09\\
-23953.0791788856	6.7966459789558e-09\\
-23906.1583577713	2.04807598616413e-08\\
-23859.2375366569	5.52796503081048e-08\\
-23812.3167155425	3.65762502160344e-08\\
-23765.3958944282	9.92843317004177e-09\\
-23718.4750733138	2.94765032381341e-09\\
-23671.5542521994	3.2300533034679e-09\\
-23624.633431085	1.36175432290685e-08\\
-23577.7126099707	4.48747761915906e-08\\
-23530.7917888563	4.6929820779037e-08\\
-23483.8709677419	1.79389740353176e-08\\
-23436.9501466276	5.0155609529083e-09\\
-23390.0293255132	2.59045700537738e-09\\
-23343.1085043988	7.43596877844657e-09\\
-23296.1876832845	2.85948109585483e-08\\
-23249.2668621701	5.64421629755727e-08\\
-23202.3460410557	2.86395558222412e-08\\
-23155.4252199413	8.25396800803398e-09\\
-23108.504398827	1.94329320187427e-09\\
-23061.5835777126	5.08638943875813e-09\\
-23014.6627565982	1.69458838984734e-08\\
-22967.7419354839	4.914913612614e-08\\
-22920.8211143695	4.503518652506e-08\\
-22873.9002932551	1.22306774439153e-08\\
-22826.9794721408	3.99863230483642e-09\\
-22780.0586510264	2.72394801905653e-09\\
-22733.137829912	1.03643576437291e-08\\
-22686.2170087977	3.68030586025632e-08\\
-22639.2961876833	5.32958240320471e-08\\
-22592.3753665689	2.3321315882338e-08\\
-22545.4545454545	5.83998123450221e-09\\
-22498.5337243402	2.32309500826676e-09\\
-22451.6129032258	6.04640707879494e-09\\
-22404.6920821114	2.31953046216101e-08\\
-22357.7712609971	5.38528183115835e-08\\
-22310.8504398827	3.72297236817084e-08\\
-22263.9296187683	9.99812925265712e-09\\
-22217.008797654	2.81356549578239e-09\\
-22170.0879765396	3.45843657466441e-09\\
-22123.1671554252	1.35129848896611e-08\\
-22076.2463343109	4.46819602560227e-08\\
-22029.3255131965	4.95817318415197e-08\\
-21982.4046920821	1.7975867310242e-08\\
-21935.4838709677	4.63857289232878e-09\\
-21888.5630498534	2.5506581825241e-09\\
-21841.642228739	7.70263256632904e-09\\
-21794.7214076246	3.08571530283524e-08\\
-21747.8005865103	5.33072741998019e-08\\
-21700.8797653959	3.20811100386197e-08\\
-21653.9589442815	7.80251761952471e-09\\
-21607.0381231672	2.48381899561016e-09\\
-21560.1173020528	4.25176887380567e-09\\
-21513.1964809384	1.86618704426435e-08\\
-21466.275659824	5.05406859378726e-08\\
-21419.3548387097	4.50032432930166e-08\\
-21372.4340175953	1.28228118670802e-08\\
-21325.5131964809	4.01734106882577e-09\\
-21278.5923753666	2.87260818030064e-09\\
-21184.7507331378	3.66140604361641e-08\\
-21137.8299120235	5.64460386433229e-08\\
-21090.9090909091	2.41604908774575e-08\\
-21043.9882697947	5.42842595687423e-09\\
-20997.0674486804	2.75708938672057e-09\\
-20950.146627566	5.65530334508873e-09\\
-20903.2258064516	2.6523078908794e-08\\
-20856.3049853372	5.2512551747733e-08\\
-20809.3841642229	3.89295935817746e-08\\
-20762.4633431085	1.0114692007242e-08\\
-20715.5425219941	3.35807679341274e-09\\
-20668.6217008798	3.97348137690216e-09\\
-20621.7008797654	1.38695326077387e-08\\
-20574.780058651	4.49852863181559e-08\\
-20527.8592375367	5.25809008067928e-08\\
-20480.9384164223	1.87211842760278e-08\\
-20434.0175953079	4.61449178592977e-09\\
-20387.0967741935	2.38106076700938e-09\\
-20293.2551319648	3.08320737261894e-08\\
-20246.3343108504	5.78016251029166e-08\\
-20199.4134897361	3.23115875615967e-08\\
-20152.4926686217	7.33244227513031e-09\\
-20105.5718475073	3.13798222617853e-09\\
-20058.651026393	4.34580072080288e-09\\
-20011.7302052786	1.92317840035986e-08\\
-19964.8093841642	5.32767164483021e-08\\
-19917.8885630499	4.5466286884904e-08\\
-19870.9677419355	1.49957385709902e-08\\
-19824.0469208211	3.72557595642003e-09\\
-19777.1260997067	2.98631164668493e-09\\
-19683.284457478	3.96103182384384e-08\\
-19636.3636363636	5.80565009802127e-08\\
-19589.4428152493	2.4856632706025e-08\\
-19542.5219941349	6.47623248152758e-09\\
-19495.6011730205	2.52604352584223e-09\\
-19448.6803519062	6.35936796218229e-09\\
-19401.7595307918	2.5660867585257e-08\\
-19354.8387096774	5.75675234171473e-08\\
-19307.917888563	4.07546022079208e-08\\
-19260.9970674487	1.05714047034758e-08\\
-19214.0762463343	3.59292533606328e-09\\
-19167.1554252199	3.86342636524016e-09\\
-19120.2346041056	1.47576623210414e-08\\
-19073.3137829912	4.78479473805968e-08\\
-19026.3929618768	5.51642986361981e-08\\
-18979.4721407625	1.96390723081506e-08\\
-18932.5513196481	5.15077964557016e-09\\
-18885.6304985337	2.71134938618296e-09\\
-18838.7096774194	8.74087396109374e-09\\
-18791.788856305	3.31771249494441e-08\\
-18744.8680351906	6.13864820937776e-08\\
-18697.9472140762	3.34970399501938e-08\\
-18651.0263929619	8.37072089261064e-09\\
-18604.1055718475	2.87922041076996e-09\\
-18557.1847507331	5.43577748908519e-09\\
-18510.2639296188	1.85661246399783e-08\\
-18463.3431085044	5.82329957814198e-08\\
-18416.42228739	4.90240654430953e-08\\
-18369.5014662757	1.54300743529304e-08\\
-18322.5806451613	3.92992905541118e-09\\
-18275.6598240469	2.96251171265624e-09\\
-18228.7390029326	1.18683618630837e-08\\
-18181.8181818182	4.37292248464953e-08\\
-18134.8973607038	6.03262610286499e-08\\
-18087.9765395894	2.68915781436197e-08\\
-18041.0557184751	7.08178641222401e-09\\
-17994.1348973607	2.68312175901679e-09\\
-17947.2140762463	7.13643194676664e-09\\
-17900.293255132	2.72343762349005e-08\\
-17853.3724340176	6.12812707912036e-08\\
-17806.4516129032	4.4547319307051e-08\\
-17759.5307917889	1.2143119359815e-08\\
-17712.6099706745	3.04866750553495e-09\\
-17665.6891495601	4.08187825655986e-09\\
-17618.7683284457	1.61131559971182e-08\\
-17571.8475073314	5.13912623789616e-08\\
-17524.926686217	6.12153829866175e-08\\
-17478.0058651026	1.95814833371346e-08\\
-17431.0850439883	5.76248727810637e-09\\
-17384.1642228739	3.07933928062263e-09\\
-17337.2434017595	9.01308904767838e-09\\
-17290.3225806452	3.62316319035331e-08\\
-17243.4017595308	6.65923415622805e-08\\
-17196.4809384164	3.60964470944908e-08\\
-17149.5601173021	9.55381024357878e-09\\
-17102.6392961877	2.98305069884662e-09\\
-17055.7184750733	5.71751781462428e-09\\
-17008.7976539589	2.19120690358788e-08\\
-16961.8768328446	6.17628233135847e-08\\
-16914.9560117302	5.37192353350696e-08\\
-16868.0351906158	1.66315504551796e-08\\
-16821.1143695015	4.37698929346378e-09\\
-16774.1935483871	4.10945496359856e-09\\
-16727.2727272727	1.21656068750467e-08\\
-16680.3519061584	4.70922815040199e-08\\
-16633.431085044	6.6709709644672e-08\\
-16586.5102639296	2.98295311518424e-08\\
-16539.5894428152	7.42911763430195e-09\\
-16492.6686217009	2.97384225352324e-09\\
-16445.7478005865	7.68452853010408e-09\\
-16398.8269794721	2.96411382418797e-08\\
-16351.9061583578	6.94504696371247e-08\\
-16304.9853372434	4.75187901595118e-08\\
-16211.1436950147	3.61912196127011e-09\\
-16164.2228739003	4.55837007944116e-09\\
-16117.3020527859	1.84482742237976e-08\\
-16070.3812316716	5.6156739930779e-08\\
-16023.4604105572	6.59832259500854e-08\\
-15976.5395894428	2.39041041783412e-08\\
-15929.6187683284	5.42815660825681e-09\\
-15882.6979472141	3.95503769024861e-09\\
-15835.7771260997	9.30777685952316e-09\\
-15788.8563049853	4.12012458932618e-08\\
-15741.935483871	7.392992538558e-08\\
-15695.0146627566	3.95508853514749e-08\\
-15601.1730205279	3.08977805024238e-09\\
-15554.2521994135	6.04237141339169e-09\\
-15507.3313782991	2.46421578349954e-08\\
-15460.4105571848	6.94656322159664e-08\\
-15413.4897360704	5.95974695238638e-08\\
-15366.568914956	1.8445806984027e-08\\
-15319.6480938416	5.31487976280728e-09\\
-15272.7272727273	4.53600759315367e-09\\
-15225.8064516129	1.37455728625863e-08\\
-15178.8856304985	5.31300638112411e-08\\
-15131.9648093842	7.1873091808668e-08\\
-15085.0439882698	3.64240419489675e-08\\
-15038.1231671554	7.8060731745481e-09\\
-14991.2023460411	3.8376488358571e-09\\
-14944.2815249267	7.6165580743415e-09\\
-14897.3607038123	3.40596751226429e-08\\
-14850.4398826979	7.98501423111432e-08\\
-14803.5190615836	5.19436601900925e-08\\
-14709.6774193548	4.1180578766633e-09\\
-14662.7565982405	5.81767455504117e-09\\
-14615.8357771261	1.86769691798484e-08\\
-14568.9149560117	6.58065886694649e-08\\
-14521.9941348974	7.41189268903372e-08\\
-14475.073313783	2.74713136450293e-08\\
-14428.1524926686	6.29231510552849e-09\\
-14381.2316715543	3.70899477158701e-09\\
-14334.3108504399	1.17371427237958e-08\\
-14287.3900293255	4.85033227484484e-08\\
-14240.4692082111	8.08519531394101e-08\\
-14193.5483870968	4.69464613441758e-08\\
-14146.6275659824	1.17770131258618e-08\\
-14099.706744868	3.92786018675943e-09\\
-14052.7859237537	7.2106903284573e-09\\
-14005.8651026393	2.81017053957105e-08\\
-13958.9442815249	7.77989806975849e-08\\
-13912.0234604106	7.03187062477802e-08\\
-13865.1026392962	2.13566902986816e-08\\
-13818.1818181818	5.49777860911888e-09\\
-13771.2609970674	4.51915248755835e-09\\
-13724.3401759531	1.6159609974527e-08\\
-13677.4193548387	6.04548953109271e-08\\
-13630.4985337243	8.77201048826482e-08\\
-13583.57771261	3.79110767805788e-08\\
-13536.6568914956	9.93727091180737e-09\\
-13489.7360703812	3.85101567695233e-09\\
-13442.8152492669	9.76414496353388e-09\\
-13395.8944281525	3.91982583646509e-08\\
-13348.9736070381	8.98962678095017e-08\\
-13302.0527859238	6.22697731756055e-08\\
-13255.1319648094	1.76194892029727e-08\\
-13208.211143695	4.66131398751507e-09\\
-13161.2903225806	6.91462402234334e-09\\
-13114.3695014663	2.12846157471882e-08\\
-13067.4486803519	7.8689429189867e-08\\
-13020.5278592375	8.61567155568366e-08\\
-12973.6070381232	2.98156991781079e-08\\
-12926.6862170088	9.00375302105324e-09\\
-12879.7653958944	3.47464099338566e-09\\
-12832.8445747801	1.4866585108076e-08\\
-12785.9237536657	5.18670517623186e-08\\
-12739.0029325513	1.02559513270975e-07\\
-12692.082111437	5.2136384514168e-08\\
-12645.1612903226	1.36265351317618e-08\\
-12598.2404692082	4.78304431366226e-09\\
-12551.3196480938	9.35883419258259e-09\\
-12504.3988269795	3.36820543048697e-08\\
-12457.4780058651	8.84493457660992e-08\\
-12410.5571847507	8.48512159952241e-08\\
-12363.6363636364	2.58662440843274e-08\\
-12316.715542522	6.17284104748679e-09\\
-12269.7947214076	5.69531920574659e-09\\
-12222.8739002933	1.93246142530775e-08\\
-12175.9530791789	7.153995371067e-08\\
-12129.0322580645	1.08499616607086e-07\\
-12082.1114369501	4.03376440377801e-08\\
-12035.1906158358	1.34695121618723e-08\\
-11988.2697947214	3.2659689282723e-09\\
-11941.348973607	1.37189058855145e-08\\
-11894.4281524927	4.16014142566862e-08\\
-11847.5073313783	1.12678383440928e-07\\
-11800.5865102639	7.48208869015089e-08\\
-11753.6656891496	2.04119942761367e-08\\
-11706.7448680352	6.1060362975308e-09\\
-11659.8240469208	6.76813166287461e-09\\
-11612.9032258065	2.87161566319852e-08\\
-11565.9824046921	9.50849540337011e-08\\
-11519.0615835777	9.97835167854541e-08\\
-11472.1407624633	3.82964248331909e-08\\
-11425.219941349	1.07786207078719e-08\\
-11378.2991202346	5.630540752072e-09\\
-11331.3782991202	1.62968036777413e-08\\
-11284.4574780059	6.30159336639295e-08\\
-11237.5366568915	1.24850847225958e-07\\
-11190.6158357771	6.35857755418504e-08\\
-11143.6950146628	1.8437259158053e-08\\
-11096.7741935484	4.37655254485961e-09\\
-11049.853372434	1.15852930273047e-08\\
-11002.9325513196	3.88657252650569e-08\\
-10956.0117302053	1.13200692976089e-07\\
-10909.0909090909	1.04111940335329e-07\\
-10862.1700879765	2.84074266710156e-08\\
-10815.2492668622	9.35561728132487e-09\\
-10768.3284457478	6.45188633876638e-09\\
-10721.4076246334	2.47405341318603e-08\\
-10674.4868035191	8.83463245795813e-08\\
-10627.5659824047	1.28410817130616e-07\\
-10580.6451612903	5.64290614222548e-08\\
-10533.724340176	1.42208422085054e-08\\
-10486.8035190616	5.7157086311442e-09\\
-10439.8826979472	1.50344994670452e-08\\
-10392.9618768328	5.8070171498379e-08\\
-10346.0410557185	1.35390001916383e-07\\
-10299.1202346041	9.39609730563783e-08\\
-10252.1994134897	2.53819457408481e-08\\
-10205.2785923754	7.19890368806271e-09\\
-10158.357771261	8.96598196488779e-09\\
-10111.4369501466	3.52942332809654e-08\\
-10064.5161290323	1.1734145487644e-07\\
-10017.5953079179	1.30709453330319e-07\\
-9970.67448680352	4.76179497906183e-08\\
-9923.75366568915	1.23784948898146e-08\\
-9876.83284457478	6.89145093221316e-09\\
-9829.91202346041	2.10110368226544e-08\\
-9782.99120234604	8.47286437581117e-08\\
-9736.07038123167	1.46984434416603e-07\\
-9689.1495601173	8.88371767768556e-08\\
-9642.22873900293	2.17556914107965e-08\\
-9595.30791788856	6.99531160674242e-09\\
-9548.38709677419	1.2124291761384e-08\\
-9501.46627565983	5.36110742451109e-08\\
-9454.54545454546	1.45957787163644e-07\\
-9407.62463343109	1.30498139368501e-07\\
-9360.70381231671	3.73906301549354e-08\\
-9313.78299120235	1.18104772007711e-08\\
-9266.86217008798	8.56862580009428e-09\\
-9173.02052785924	1.10882193349516e-07\\
-9126.09970674487	1.71656105930741e-07\\
-9079.1788856305	7.3850378592307e-08\\
-9032.25806451613	1.67096193144771e-08\\
-8985.33724340176	8.58751784775143e-09\\
-8938.41642228739	1.7839440526982e-08\\
-8891.49560117302	8.43105481678183e-08\\
-8844.57478005865	1.67786793821608e-07\\
-8797.65395894428	1.2493126568205e-07\\
-8750.73313782991	3.26742654755232e-08\\
-8703.81231671554	1.09491938461992e-08\\
-8656.89149560118	1.31562988357027e-08\\
-8609.97067448681	4.6340823174733e-08\\
-8563.04985337243	1.51227426545887e-07\\
-8516.12903225806	1.7755141436175e-07\\
-8469.2082111437	6.35911852136742e-08\\
-8422.28739002933	1.58066975015546e-08\\
-8375.36656891496	8.27684598124778e-09\\
-8281.52492668622	1.09231982021092e-07\\
-8234.60410557185	2.05779852504288e-07\\
-8187.68328445748	1.15622244159888e-07\\
-8140.76246334311	2.64388946811332e-08\\
-8093.84164222874	1.14509889408833e-08\\
-8046.92082111437	1.61069578958849e-08\\
-8000	7.18831210604878e-08\\
-7953.07917888563	2.00374849201041e-07\\
-7906.15835777126	1.71816891198543e-07\\
-7859.23753665689	5.70653525909103e-08\\
-7812.31671554252	1.43149049781667e-08\\
-7765.39589442815	1.16714758906804e-08\\
-7671.55425219941	1.57678655745871e-07\\
-7624.63343108504	2.32307335557649e-07\\
-7577.71260997068	1.00040677870297e-07\\
-7530.79178885631	2.63051861115643e-08\\
-7483.87096774194	1.04111060434204e-08\\
-7436.95014662757	2.6595439046279e-08\\
-7390.0293255132	1.08325349458173e-07\\
-7343.10850439883	2.444767722202e-07\\
-7296.18768328446	1.74024659419371e-07\\
-7249.26686217009	4.54994642803126e-08\\
-7202.34604105572	1.56480642387427e-08\\
-7155.42521994135	1.71447883525274e-08\\
-7108.50439882698	6.61803121691524e-08\\
-7061.58357771261	2.16201876894052e-07\\
-7014.66275659824	2.50644781374838e-07\\
-6967.74193548387	8.98621012918124e-08\\
-6920.8211143695	2.38154113903945e-08\\
-6873.90029325513	1.27668089724597e-08\\
-6826.97947214076	4.17262600248256e-08\\
-6780.05865102639	1.59906727979708e-07\\
-6733.13782991202	2.9770366212917e-07\\
-6686.21700879766	1.6346058002735e-07\\
-6639.29618768329	4.12471540135661e-08\\
-6592.37536656891	1.43935200382511e-08\\
-6545.45454545455	2.77025342492599e-08\\
-6498.53372434018	9.56958473858157e-08\\
-6451.61290322581	3.0233308030269e-07\\
-6404.69208211144	2.56113229014012e-07\\
-6357.77126099707	8.13089205823553e-08\\
-6310.8504398827	2.09619376616614e-08\\
-6263.92961876833	1.6139952808651e-08\\
-6217.00879765396	6.552899802552e-08\\
-6170.08797653959	2.43891720161373e-07\\
-6123.16715542522	3.38703581653104e-07\\
-6076.24633431085	1.52070795930442e-07\\
-6029.32551319648	4.05242171980202e-08\\
-5982.40469208211	1.56385150022363e-08\\
-5935.48387096774	4.23812364205566e-08\\
-5888.56304985337	1.63724533030608e-07\\
-5841.64222873901	3.71180888772674e-07\\
-5794.72140762463	2.71722829745159e-07\\
-5747.80058651026	7.48841669425628e-08\\
-5700.87976539589	1.90771195263387e-08\\
-5653.95894428153	2.61552386320575e-08\\
-5607.03812316716	1.04697123252699e-07\\
-5560.11730205279	3.37275275092751e-07\\
-5513.19648093842	4.04697730602102e-07\\
-5466.27565982405	1.3060141296732e-07\\
-5419.35483870968	3.8959587399495e-08\\
-5372.43401759531	2.13199548693056e-08\\
-5325.51319648094	6.35734129435394e-08\\
-5278.59237536657	2.58766460313171e-07\\
-5231.6715542522	4.79488596551454e-07\\
-5184.75073313783	2.62057530997236e-07\\
-5137.82991202346	7.02693308660746e-08\\
-5090.90909090909	2.23536666542292e-08\\
-5043.98826979472	4.396507517625e-08\\
-4997.06744868035	1.71051028774799e-07\\
-4950.14662756598	4.87089143723891e-07\\
-4903.22580645161	4.27153983931986e-07\\
-4856.30498533724	1.33770207983787e-07\\
-4809.38416422287	3.57696049184982e-08\\
-4762.46334310851	3.45755940334036e-08\\
-4715.54252199414	1.04357642345076e-07\\
-4668.62170087976	4.09147616299963e-07\\
-4621.7008797654	5.84868894954649e-07\\
-4574.78005865103	2.64170525704045e-07\\
-4527.85923753666	6.68557781738678e-08\\
-4480.93841642229	2.74419346757138e-08\\
-4434.01759530792	7.2800534420113e-08\\
-4387.09677419355	2.85541547249204e-07\\
-4340.17595307918	6.76227343203146e-07\\
-4293.25513196481	4.67207975816167e-07\\
-4199.41348973607	3.68624653521296e-08\\
-4152.4926686217	4.80329624949477e-08\\
-4105.57184750733	1.98264575237626e-07\\
-4058.65102639296	6.12320639614469e-07\\
-4011.73020527859	7.26697319859692e-07\\
-3964.80938416422	2.66806198970501e-07\\
-3917.88856304986	6.17196297653782e-08\\
-3870.96774193548	4.65233114751478e-08\\
-3824.04692082111	1.12613361434595e-07\\
-3777.12609970674	5.07276404180957e-07\\
-3730.20527859238	9.21376187305949e-07\\
-3683.28445747801	4.98763998879378e-07\\
-3589.44281524927	4.08397329850772e-08\\
-3542.5219941349	8.29621427430603e-08\\
-3495.60117302053	3.45915037119533e-07\\
-3448.68035190616	9.90623145681188e-07\\
-3401.75953079179	8.60360036623149e-07\\
-3354.83870967742	2.70883872125703e-07\\
-3307.91788856305	8.00525187388106e-08\\
-3260.99706744868	7.14308576769107e-08\\
-3214.07624633431	2.22729099787069e-07\\
-3167.15542521994	8.78167591194005e-07\\
-3120.23460410557	1.20439657229042e-06\\
-3073.3137829912	6.20231798879656e-07\\
-3026.39296187683	1.36318662024007e-07\\
-2979.47214076246	6.9817811925861e-08\\
-2932.55131964809	1.44440025488377e-07\\
-2885.63049853372	6.61712147795903e-07\\
-2838.70967741936	1.57971692261491e-06\\
-2791.78885630499	1.04351703418893e-06\\
-2697.94721407625	8.75500122244917e-08\\
-2651.02639296188	1.30895637839849e-07\\
-2604.10557184751	4.34678996484143e-07\\
-2557.18475073314	1.56649530909088e-06\\
-2510.26392961877	1.79504709402612e-06\\
-2463.3431085044	6.80605310638206e-07\\
-2416.42228739003	1.61114348198402e-07\\
-2369.50146627566	1.00384188670056e-07\\
-2322.58064516129	3.33240833111429e-07\\
-2275.65982404692	1.42309304586673e-06\\
-2228.73900293255	2.42388519006517e-06\\
-2181.81818181818	1.43708694081246e-06\\
-2134.89736070381	3.7321294956581e-07\\
-2087.97653958944	1.30834116981878e-07\\
-2041.05571847507	2.57420186722671e-07\\
-1994.13489736071	1.04671207345637e-06\\
-1947.21407624634	2.98085454347638e-06\\
-1900.29325513196	2.75701681020195e-06\\
-1853.37243401759	8.66490590493051e-07\\
-1806.45161290323	2.33846368022205e-07\\
-1759.53079178886	2.09663610697347e-07\\
-1665.68914956012	3.0891824145614e-06\\
-1618.76832844575	4.61178194249124e-06\\
-1571.84750733138	2.05808116159071e-06\\
-1524.92668621701	5.69005036337926e-07\\
-1478.00586510264	2.39923624491242e-07\\
-1431.08504398827	6.66005378224616e-07\\
-1384.1642228739	2.82910421976991e-06\\
-1337.24340175953	6.74578276554042e-06\\
-1290.32258064516	4.84034363204172e-06\\
-1243.40175953079	1.44862730448142e-06\\
-1196.48093841642	4.14748468724272e-07\\
-1149.56011730205	7.02710132915608e-07\\
-1102.63929618768	2.36523257250344e-06\\
-1055.71847507331	9.22804422469072e-06\\
-1008.79765395894	1.05941632128154e-05\\
-961.876832844573	3.88218066098419e-06\\
-914.956011730206	1.28523964515546e-06\\
-868.035190615836	5.88427548963181e-07\\
-821.114369501465	2.86482473605217e-06\\
-774.193548387098	1.10986869042613e-05\\
-727.272727272728	2.3396542373272e-05\\
-680.351906158357	1.28048966837301e-05\\
-633.431085043987	3.75513216665831e-06\\
-586.51026392962	1.55656131140366e-06\\
-539.58944281525	4.03947609222608e-06\\
-492.668621700879	1.75555355341337e-05\\
-445.747800586509	5.21406116962008e-05\\
-398.826979472142	5.54651214436025e-05\\
-351.906158357771	2.02646641849441e-05\\
-304.985337243401	6.32451119377884e-06\\
-258.064516129034	9.92479796124793e-06\\
-211.143695014664	5.68988762309767e-05\\
-164.222873900293	0.000307621250998231\\
-117.302052785923	0.000616957306531322\\
-70.3812316715557	0.000576453183126342\\
-23.4604105571852	0.00213340607671187\\
23.4604105571852	0.00512914876701211\\
70.3812316715557	0.00213340607671187\\
117.302052785923	0.000576453183126342\\
164.222873900293	0.000616957306531322\\
211.143695014664	0.000307621250998231\\
258.064516129034	5.68988762309767e-05\\
304.985337243401	9.92479796124793e-06\\
351.906158357771	6.32451119377884e-06\\
398.826979472142	2.02646641849441e-05\\
445.747800586509	5.54651214436025e-05\\
492.668621700879	5.21406116962008e-05\\
539.58944281525	1.75555355341337e-05\\
586.51026392962	4.03947609222608e-06\\
633.431085043987	1.55656131140366e-06\\
680.351906158357	3.75513216665831e-06\\
727.272727272728	1.28048966837301e-05\\
774.193548387098	2.3396542373272e-05\\
821.114369501465	1.10986869042613e-05\\
868.035190615836	2.86482473605217e-06\\
914.956011730206	5.88427548963181e-07\\
961.876832844573	1.28523964515546e-06\\
1008.79765395894	3.88218066098419e-06\\
1055.71847507331	1.05941632128154e-05\\
1102.63929618768	9.22804422469072e-06\\
1149.56011730205	2.36523257250344e-06\\
1196.48093841642	7.02710132915608e-07\\
1243.40175953079	4.14748468724272e-07\\
1290.32258064516	1.44862730448142e-06\\
1337.24340175953	4.84034363204172e-06\\
1384.1642228739	6.74578276554042e-06\\
1431.08504398827	2.82910421976991e-06\\
1478.00586510264	6.66005378224616e-07\\
1524.92668621701	2.39923624491242e-07\\
1571.84750733138	5.69005036337926e-07\\
1618.76832844575	2.05808116159071e-06\\
1665.68914956012	4.61178194249124e-06\\
1712.60997067449	3.0891824145614e-06\\
1806.45161290323	2.09663610697347e-07\\
1853.37243401759	2.33846368022205e-07\\
1900.29325513196	8.66490590493051e-07\\
1947.21407624634	2.75701681020195e-06\\
1994.13489736071	2.98085454347638e-06\\
2041.05571847507	1.04671207345637e-06\\
2087.97653958944	2.57420186722671e-07\\
2134.89736070381	1.30834116981878e-07\\
2181.81818181818	3.7321294956581e-07\\
2228.73900293255	1.43708694081246e-06\\
2275.65982404692	2.42388519006517e-06\\
2322.58064516129	1.42309304586673e-06\\
2369.50146627566	3.33240833111429e-07\\
2416.42228739003	1.00384188670056e-07\\
2463.3431085044	1.61114348198402e-07\\
2510.26392961877	6.80605310638206e-07\\
2557.18475073314	1.79504709402612e-06\\
2604.10557184751	1.56649530909088e-06\\
2651.02639296188	4.34678996484143e-07\\
2697.94721407625	1.30895637839849e-07\\
2744.86803519062	8.75500122244917e-08\\
2838.70967741936	1.04351703418893e-06\\
2885.63049853372	1.57971692261491e-06\\
2932.55131964809	6.61712147795903e-07\\
2979.47214076246	1.44440025488377e-07\\
3026.39296187683	6.9817811925861e-08\\
3073.3137829912	1.36318662024007e-07\\
3120.23460410557	6.20231798879656e-07\\
3167.15542521994	1.20439657229042e-06\\
3214.07624633431	8.78167591194005e-07\\
3260.99706744868	2.22729099787069e-07\\
3307.91788856305	7.14308576769107e-08\\
3354.83870967742	8.00525187388106e-08\\
3401.75953079179	2.70883872125703e-07\\
3448.68035190616	8.60360036623149e-07\\
3495.60117302053	9.90623145681188e-07\\
3542.5219941349	3.45915037119533e-07\\
3589.44281524927	8.29621427430603e-08\\
3636.36363636364	4.08397329850772e-08\\
3730.20527859238	4.98763998879378e-07\\
3777.12609970674	9.21376187305949e-07\\
3824.04692082111	5.07276404180957e-07\\
3870.96774193548	1.12613361434595e-07\\
3917.88856304986	4.65233114751478e-08\\
3964.80938416422	6.17196297653782e-08\\
4011.73020527859	2.66806198970501e-07\\
4058.65102639296	7.26697319859692e-07\\
4105.57184750733	6.12320639614469e-07\\
4152.4926686217	1.98264575237626e-07\\
4199.41348973607	4.80329624949477e-08\\
4246.33431085044	3.68624653521296e-08\\
4340.17595307918	4.67207975816167e-07\\
4387.09677419355	6.76227343203146e-07\\
4434.01759530792	2.85541547249204e-07\\
4480.93841642229	7.2800534420113e-08\\
4527.85923753666	2.74419346757138e-08\\
4574.78005865103	6.68557781738678e-08\\
4621.7008797654	2.64170525704045e-07\\
4668.62170087976	5.84868894954649e-07\\
4715.54252199414	4.09147616299963e-07\\
4762.46334310851	1.04357642345076e-07\\
4809.38416422287	3.45755940334036e-08\\
4856.30498533724	3.57696049184982e-08\\
4903.22580645161	1.33770207983787e-07\\
4950.14662756598	4.27153983931986e-07\\
4997.06744868035	4.87089143723891e-07\\
5043.98826979472	1.71051028774799e-07\\
5090.90909090909	4.396507517625e-08\\
5137.82991202346	2.23536666542292e-08\\
5184.75073313783	7.02693308660746e-08\\
5231.6715542522	2.62057530997236e-07\\
5278.59237536657	4.79488596551454e-07\\
5325.51319648094	2.58766460313171e-07\\
5372.43401759531	6.35734129435394e-08\\
5419.35483870968	2.13199548693056e-08\\
5466.27565982405	3.8959587399495e-08\\
5513.19648093842	1.3060141296732e-07\\
5560.11730205279	4.04697730602102e-07\\
5607.03812316716	3.37275275092751e-07\\
5653.95894428153	1.04697123252699e-07\\
5700.87976539589	2.61552386320575e-08\\
5747.80058651026	1.90771195263387e-08\\
5794.72140762463	7.48841669425628e-08\\
5841.64222873901	2.71722829745159e-07\\
5888.56304985337	3.71180888772674e-07\\
5935.48387096774	1.63724533030608e-07\\
5982.40469208211	4.23812364205566e-08\\
6029.32551319648	1.56385150022363e-08\\
6076.24633431085	4.05242171980202e-08\\
6123.16715542522	1.52070795930442e-07\\
6170.08797653959	3.38703581653104e-07\\
6217.00879765396	2.43891720161373e-07\\
6263.92961876833	6.552899802552e-08\\
6310.8504398827	1.6139952808651e-08\\
6357.77126099707	2.09619376616614e-08\\
6404.69208211144	8.13089205823553e-08\\
6451.61290322581	2.56113229014012e-07\\
6498.53372434018	3.0233308030269e-07\\
6545.45454545455	9.56958473858157e-08\\
6592.37536656891	2.77025342492599e-08\\
6639.29618768329	1.43935200382511e-08\\
6686.21700879766	4.12471540135661e-08\\
6733.13782991202	1.6346058002735e-07\\
6780.05865102639	2.9770366212917e-07\\
6826.97947214076	1.59906727979708e-07\\
6873.90029325513	4.17262600248256e-08\\
6920.8211143695	1.27668089724597e-08\\
6967.74193548387	2.38154113903945e-08\\
7014.66275659824	8.98621012918124e-08\\
7061.58357771261	2.50644781374838e-07\\
7108.50439882698	2.16201876894052e-07\\
7155.42521994135	6.61803121691524e-08\\
7202.34604105572	1.71447883525274e-08\\
7249.26686217009	1.56480642387427e-08\\
7296.18768328446	4.54994642803126e-08\\
7343.10850439883	1.74024659419371e-07\\
7390.0293255132	2.444767722202e-07\\
7436.95014662757	1.08325349458173e-07\\
7483.87096774194	2.6595439046279e-08\\
7530.79178885631	1.04111060434204e-08\\
7577.71260997068	2.63051861115643e-08\\
7624.63343108504	1.00040677870297e-07\\
7671.55425219941	2.32307335557649e-07\\
7718.47507331378	1.57678655745871e-07\\
7812.31671554252	1.16714758906804e-08\\
7859.23753665689	1.43149049781667e-08\\
7906.15835777126	5.70653525909103e-08\\
7953.07917888563	1.71816891198543e-07\\
8000	2.00374849201041e-07\\
8046.92082111437	7.18831210604878e-08\\
8093.84164222874	1.61069578958849e-08\\
8140.76246334311	1.14509889408833e-08\\
8187.68328445748	2.64388946811332e-08\\
8234.60410557185	1.15622244159888e-07\\
8281.52492668622	2.05779852504288e-07\\
8328.44574780059	1.09231982021092e-07\\
8422.28739002933	8.27684598124778e-09\\
8469.2082111437	1.58066975015546e-08\\
8516.12903225806	6.35911852136742e-08\\
8563.04985337243	1.7755141436175e-07\\
8609.97067448681	1.51227426545887e-07\\
8656.89149560118	4.6340823174733e-08\\
8703.81231671554	1.31562988357027e-08\\
8750.73313782991	1.09491938461992e-08\\
8797.65395894428	3.26742654755232e-08\\
8844.57478005865	1.2493126568205e-07\\
8891.49560117302	1.67786793821608e-07\\
8938.41642228739	8.43105481678183e-08\\
8985.33724340176	1.7839440526982e-08\\
9032.25806451613	8.58751784775143e-09\\
9079.1788856305	1.67096193144771e-08\\
9126.09970674487	7.3850378592307e-08\\
9173.02052785924	1.71656105930741e-07\\
9219.94134897361	1.10882193349516e-07\\
9313.78299120235	8.56862580009428e-09\\
9360.70381231671	1.18104772007711e-08\\
9407.62463343109	3.73906301549354e-08\\
9454.54545454546	1.30498139368501e-07\\
9501.46627565983	1.45957787163644e-07\\
9548.38709677419	5.36110742451109e-08\\
9595.30791788856	1.2124291761384e-08\\
9642.22873900293	6.99531160674242e-09\\
9689.1495601173	2.17556914107965e-08\\
9736.07038123167	8.88371767768556e-08\\
9782.99120234604	1.46984434416603e-07\\
9829.91202346041	8.47286437581117e-08\\
9876.83284457478	2.10110368226544e-08\\
9923.75366568915	6.89145093221316e-09\\
9970.67448680352	1.23784948898146e-08\\
10017.5953079179	4.76179497906183e-08\\
10064.5161290323	1.30709453330319e-07\\
10111.4369501466	1.1734145487644e-07\\
10158.357771261	3.52942332809654e-08\\
10205.2785923754	8.96598196488779e-09\\
10252.1994134897	7.19890368806271e-09\\
10299.1202346041	2.53819457408481e-08\\
10346.0410557185	9.39609730563783e-08\\
10392.9618768328	1.35390001916383e-07\\
10439.8826979472	5.8070171498379e-08\\
10486.8035190616	1.50344994670452e-08\\
10533.724340176	5.7157086311442e-09\\
10580.6451612903	1.42208422085054e-08\\
10627.5659824047	5.64290614222548e-08\\
10674.4868035191	1.28410817130616e-07\\
10721.4076246334	8.83463245795813e-08\\
10768.3284457478	2.47405341318603e-08\\
10815.2492668622	6.45188633876638e-09\\
10862.1700879765	9.35561728132487e-09\\
10909.0909090909	2.84074266710156e-08\\
10956.0117302053	1.04111940335329e-07\\
11002.9325513196	1.13200692976089e-07\\
11049.853372434	3.88657252650569e-08\\
11096.7741935484	1.15852930273047e-08\\
11143.6950146628	4.37655254485961e-09\\
11190.6158357771	1.8437259158053e-08\\
11237.5366568915	6.35857755418504e-08\\
11284.4574780059	1.24850847225958e-07\\
11331.3782991202	6.30159336639295e-08\\
11378.2991202346	1.62968036777413e-08\\
11425.219941349	5.630540752072e-09\\
11472.1407624633	1.07786207078719e-08\\
11519.0615835777	3.82964248331909e-08\\
11565.9824046921	9.97835167854541e-08\\
11612.9032258065	9.50849540337011e-08\\
11659.8240469208	2.87161566319852e-08\\
11706.7448680352	6.76813166287461e-09\\
11753.6656891496	6.1060362975308e-09\\
11800.5865102639	2.04119942761367e-08\\
11847.5073313783	7.48208869015089e-08\\
11894.4281524927	1.12678383440928e-07\\
11941.348973607	4.16014142566862e-08\\
11988.2697947214	1.37189058855145e-08\\
12035.1906158358	3.2659689282723e-09\\
12082.1114369501	1.34695121618723e-08\\
12129.0322580645	4.03376440377801e-08\\
12175.9530791789	1.08499616607086e-07\\
12222.8739002933	7.153995371067e-08\\
12269.7947214076	1.93246142530775e-08\\
12316.715542522	5.69531920574659e-09\\
12363.6363636364	6.17284104748679e-09\\
12410.5571847507	2.58662440843274e-08\\
12457.4780058651	8.48512159952241e-08\\
12504.3988269795	8.84493457660992e-08\\
12551.3196480938	3.36820543048697e-08\\
12598.2404692082	9.35883419258259e-09\\
12645.1612903226	4.78304431366226e-09\\
12692.082111437	1.36265351317618e-08\\
12739.0029325513	5.2136384514168e-08\\
12785.9237536657	1.02559513270975e-07\\
12832.8445747801	5.18670517623186e-08\\
12879.7653958944	1.4866585108076e-08\\
12926.6862170088	3.47464099338566e-09\\
12973.6070381232	9.00375302105324e-09\\
13020.5278592375	2.98156991781079e-08\\
13067.4486803519	8.61567155568366e-08\\
13114.3695014663	7.8689429189867e-08\\
13161.2903225806	2.12846157471882e-08\\
13208.211143695	6.91462402234334e-09\\
13255.1319648094	4.66131398751507e-09\\
13302.0527859238	1.76194892029727e-08\\
13348.9736070381	6.22697731756055e-08\\
13395.8944281525	8.98962678095017e-08\\
13442.8152492669	3.91982583646509e-08\\
13489.7360703812	9.76414496353388e-09\\
13536.6568914956	3.85101567695233e-09\\
13583.57771261	9.93727091180737e-09\\
13630.4985337243	3.79110767805788e-08\\
13677.4193548387	8.77201048826482e-08\\
13724.3401759531	6.04548953109271e-08\\
13771.2609970674	1.6159609974527e-08\\
13818.1818181818	4.51915248755835e-09\\
13865.1026392962	5.49777860911888e-09\\
13912.0234604106	2.13566902986816e-08\\
13958.9442815249	7.03187062477802e-08\\
14005.8651026393	7.77989806975849e-08\\
14052.7859237537	2.81017053957105e-08\\
14099.706744868	7.2106903284573e-09\\
14146.6275659824	3.92786018675943e-09\\
14193.5483870968	1.17770131258618e-08\\
14240.4692082111	4.69464613441758e-08\\
14287.3900293255	8.08519531394101e-08\\
14334.3108504399	4.85033227484484e-08\\
14381.2316715543	1.17371427237958e-08\\
14428.1524926686	3.70899477158701e-09\\
14475.073313783	6.29231510552849e-09\\
14521.9941348974	2.74713136450293e-08\\
14568.9149560117	7.41189268903372e-08\\
14615.8357771261	6.58065886694649e-08\\
14662.7565982405	1.86769691798484e-08\\
14709.6774193548	5.81767455504117e-09\\
14756.5982404692	4.1180578766633e-09\\
14850.4398826979	5.19436601900925e-08\\
14897.3607038123	7.98501423111432e-08\\
14944.2815249267	3.40596751226429e-08\\
14991.2023460411	7.6165580743415e-09\\
15038.1231671554	3.8376488358571e-09\\
15085.0439882698	7.8060731745481e-09\\
15131.9648093842	3.64240419489675e-08\\
15178.8856304985	7.1873091808668e-08\\
15225.8064516129	5.31300638112411e-08\\
15272.7272727273	1.37455728625863e-08\\
15319.6480938416	4.53600759315367e-09\\
15366.568914956	5.31487976280728e-09\\
15413.4897360704	1.8445806984027e-08\\
15460.4105571848	5.95974695238638e-08\\
15507.3313782991	6.94656322159664e-08\\
15554.2521994135	2.46421578349954e-08\\
15601.1730205279	6.04237141339169e-09\\
15648.0938416422	3.08977805024238e-09\\
15741.935483871	3.95508853514749e-08\\
15788.8563049853	7.392992538558e-08\\
15835.7771260997	4.12012458932618e-08\\
15882.6979472141	9.30777685952316e-09\\
15929.6187683284	3.95503769024861e-09\\
15976.5395894428	5.42815660825681e-09\\
16023.4604105572	2.39041041783412e-08\\
16070.3812316716	6.59832259500854e-08\\
16117.3020527859	5.6156739930779e-08\\
16164.2228739003	1.84482742237976e-08\\
16211.1436950147	4.55837007944116e-09\\
16258.064516129	3.61912196127011e-09\\
16351.9061583578	4.75187901595118e-08\\
16398.8269794721	6.94504696371247e-08\\
16445.7478005865	2.96411382418797e-08\\
16492.6686217009	7.68452853010408e-09\\
16539.5894428152	2.97384225352324e-09\\
16586.5102639296	7.42911763430195e-09\\
16633.431085044	2.98295311518424e-08\\
16680.3519061584	6.6709709644672e-08\\
16727.2727272727	4.70922815040199e-08\\
16774.1935483871	1.21656068750467e-08\\
16821.1143695015	4.10945496359856e-09\\
16868.0351906158	4.37698929346378e-09\\
16914.9560117302	1.66315504551796e-08\\
16961.8768328446	5.37192353350696e-08\\
17008.7976539589	6.17628233135847e-08\\
17055.7184750733	2.19120690358788e-08\\
17102.6392961877	5.71751781462428e-09\\
17149.5601173021	2.98305069884662e-09\\
17196.4809384164	9.55381024357878e-09\\
17243.4017595308	3.60964470944908e-08\\
17290.3225806452	6.65923415622805e-08\\
17337.2434017595	3.62316319035331e-08\\
17384.1642228739	9.01308904767838e-09\\
17431.0850439883	3.07933928062263e-09\\
17478.0058651026	5.76248727810637e-09\\
17524.926686217	1.95814833371346e-08\\
17571.8475073314	6.12153829866175e-08\\
17618.7683284457	5.13912623789616e-08\\
17665.6891495601	1.61131559971182e-08\\
17712.6099706745	4.08187825655986e-09\\
17759.5307917889	3.04866750553495e-09\\
17806.4516129032	1.2143119359815e-08\\
17853.3724340176	4.4547319307051e-08\\
17900.293255132	6.12812707912036e-08\\
17947.2140762463	2.72343762349005e-08\\
17994.1348973607	7.13643194676664e-09\\
18041.0557184751	2.68312175901679e-09\\
18087.9765395894	7.08178641222401e-09\\
18134.8973607038	2.68915781436197e-08\\
18181.8181818182	6.03262610286499e-08\\
18228.7390029326	4.37292248464953e-08\\
18275.6598240469	1.18683618630837e-08\\
18322.5806451613	2.96251171265624e-09\\
18369.5014662757	3.92992905541118e-09\\
18416.42228739	1.54300743529304e-08\\
18463.3431085044	4.90240654430953e-08\\
18510.2639296188	5.82329957814198e-08\\
18557.1847507331	1.85661246399783e-08\\
18604.1055718475	5.43577748908519e-09\\
18651.0263929619	2.87922041076996e-09\\
18697.9472140762	8.37072089261064e-09\\
18744.8680351906	3.34970399501938e-08\\
18791.788856305	6.13864820937776e-08\\
18838.7096774194	3.31771249494441e-08\\
18885.6304985337	8.74087396109374e-09\\
18932.5513196481	2.71134938618296e-09\\
18979.4721407625	5.15077964557016e-09\\
19026.3929618768	1.96390723081506e-08\\
19073.3137829912	5.51642986361981e-08\\
19120.2346041056	4.78479473805968e-08\\
19167.1554252199	1.47576623210414e-08\\
19214.0762463343	3.86342636524016e-09\\
19260.9970674487	3.59292533606328e-09\\
19307.917888563	1.05714047034758e-08\\
19354.8387096774	4.07546022079208e-08\\
19401.7595307918	5.75675234171473e-08\\
19448.6803519062	2.5660867585257e-08\\
19495.6011730205	6.35936796218229e-09\\
19542.5219941349	2.52604352584223e-09\\
19589.4428152493	6.47623248152758e-09\\
19636.3636363636	2.4856632706025e-08\\
19683.284457478	5.80565009802127e-08\\
19730.2052785924	3.96103182384384e-08\\
19824.0469208211	2.98631164668493e-09\\
19870.9677419355	3.72557595642003e-09\\
19917.8885630499	1.49957385709902e-08\\
19964.8093841642	4.5466286884904e-08\\
20011.7302052786	5.32767164483021e-08\\
20058.651026393	1.92317840035986e-08\\
20105.5718475073	4.34580072080288e-09\\
20152.4926686217	3.13798222617853e-09\\
20199.4134897361	7.33244227513031e-09\\
20246.3343108504	3.23115875615967e-08\\
20293.2551319648	5.78016251029166e-08\\
20340.1759530792	3.08320737261894e-08\\
20434.0175953079	2.38106076700938e-09\\
20480.9384164223	4.61449178592977e-09\\
20527.8592375367	1.87211842760278e-08\\
20574.780058651	5.25809008067928e-08\\
20621.7008797654	4.49852863181559e-08\\
20668.6217008798	1.38695326077387e-08\\
20715.5425219941	3.97348137690216e-09\\
20762.4633431085	3.35807679341274e-09\\
20809.3841642229	1.0114692007242e-08\\
20856.3049853372	3.89295935817746e-08\\
20903.2258064516	5.2512551747733e-08\\
20950.146627566	2.6523078908794e-08\\
20997.0674486804	5.65530334508873e-09\\
21043.9882697947	2.75708938672057e-09\\
21090.9090909091	5.42842595687423e-09\\
21137.8299120235	2.41604908774575e-08\\
21184.7507331378	5.64460386433229e-08\\
21231.6715542522	3.66140604361641e-08\\
21325.5131964809	2.87260818030064e-09\\
21372.4340175953	4.01734106882577e-09\\
21419.3548387097	1.28228118670802e-08\\
21466.275659824	4.50032432930166e-08\\
21513.1964809384	5.05406859378726e-08\\
21560.1173020528	1.86618704426435e-08\\
21607.0381231672	4.25176887380567e-09\\
21653.9589442815	2.48381899561016e-09\\
21700.8797653959	7.80251761952471e-09\\
21747.8005865103	3.20811100386197e-08\\
21794.7214076246	5.33072741998019e-08\\
21841.642228739	3.08571530283524e-08\\
21888.5630498534	7.70263256632904e-09\\
21935.4838709677	2.5506581825241e-09\\
21982.4046920821	4.63857289232878e-09\\
22029.3255131965	1.7975867310242e-08\\
22076.2463343109	4.95817318415197e-08\\
22123.1671554252	4.46819602560227e-08\\
22170.0879765396	1.35129848896611e-08\\
22217.008797654	3.45843657466441e-09\\
22263.9296187683	2.81356549578239e-09\\
22310.8504398827	9.99812925265712e-09\\
22357.7712609971	3.72297236817084e-08\\
22404.6920821114	5.38528183115835e-08\\
22451.6129032258	2.31953046216101e-08\\
22498.5337243402	6.04640707879494e-09\\
22545.4545454545	2.32309500826676e-09\\
22592.3753665689	5.83998123450221e-09\\
22639.2961876833	2.3321315882338e-08\\
22686.2170087977	5.32958240320471e-08\\
22733.137829912	3.68030586025632e-08\\
22780.0586510264	1.03643576437291e-08\\
22826.9794721408	2.72394801905653e-09\\
22873.9002932551	3.99863230483642e-09\\
22920.8211143695	1.22306774439153e-08\\
22967.7419354839	4.503518652506e-08\\
23014.6627565982	4.914913612614e-08\\
23061.5835777126	1.69458838984734e-08\\
23108.504398827	5.08638943875813e-09\\
23155.4252199413	1.94329320187427e-09\\
23202.3460410557	8.25396800803398e-09\\
23249.2668621701	2.86395558222412e-08\\
23296.1876832845	5.64421629755727e-08\\
23343.1085043988	2.85948109585483e-08\\
23390.0293255132	7.43596877844657e-09\\
23436.9501466276	2.59045700537738e-09\\
23483.8709677419	5.0155609529083e-09\\
23530.7917888563	1.79389740353176e-08\\
23577.7126099707	4.6929820779037e-08\\
23624.633431085	4.48747761915906e-08\\
23671.5542521994	1.36175432290685e-08\\
23718.4750733138	3.2300533034679e-09\\
23765.3958944282	2.94765032381341e-09\\
23812.3167155425	9.92843317004177e-09\\
23859.2375366569	3.65762502160344e-08\\
23906.1583577713	5.52796503081048e-08\\
23953.0791788856	2.04807598616413e-08\\
24000	6.7966459789558e-09\\
};
\addlegendentry{Pratique}

\addplot [color=red]
  table[row sep=crcr]{%
-24000	1.26631196965301e-36\\
-23953.0791788856	2.94843294621119e-09\\
-23906.1583577713	9.20430684265351e-09\\
-23859.2375366569	1.32244628263452e-08\\
-23812.3167155425	1.14207989811639e-08\\
-23765.3958944282	5.38323193650717e-09\\
-23718.4750733138	4.99335130363089e-10\\
-23671.5542521994	1.16819700264575e-09\\
-23624.633431085	6.85998003453254e-09\\
-23577.7126099707	1.255038148357e-08\\
-23530.7917888563	1.31764034101032e-08\\
-23483.8709677419	8.16225373893329e-09\\
-23436.9501466276	1.96999445480361e-09\\
-23390.0293255132	1.50878548329001e-10\\
-23343.1085043988	4.3912848758425e-09\\
-23296.1876832845	1.0973980823713e-08\\
-23249.2668621701	1.40578011024185e-08\\
-23202.3460410557	1.08826063181257e-08\\
-23155.4252199413	4.2618721256523e-09\\
-23108.504398827	1.12567711644758e-10\\
-23061.5835777126	2.1903706296943e-09\\
-23014.6627565982	8.71430081514548e-09\\
-22967.7419354839	1.39138212944516e-08\\
-22920.8211143695	1.31554323921928e-08\\
-22873.9002932551	7.09467046867304e-09\\
-22826.9794721408	1.1334500672449e-09\\
-22780.0586510264	6.29355507531859e-10\\
-22733.137829912	6.1050974487918e-09\\
-22686.2170087977	1.27412008425165e-08\\
-22639.2961876833	1.46408960085715e-08\\
-22592.3753665689	1.00935731499371e-08\\
-22545.4545454545	3.13909147118045e-09\\
-22498.5337243402	3.53921904939615e-12\\
-22451.6129032258	3.5505082131424e-09\\
-22404.6920821114	1.06924885881024e-08\\
-22357.7712609971	1.51016737666314e-08\\
-22310.8504398827	1.28408117773952e-08\\
-22263.9296187683	5.90390606041906e-09\\
-22217.008797654	4.83278248060095e-10\\
-22170.0879765396	1.46799548364856e-09\\
-22123.1671554252	8.06106129461077e-09\\
-22076.2463343109	1.44438023698404e-08\\
-22029.3255131965	1.49362370380492e-08\\
-21982.4046920821	9.07739000075142e-09\\
-21935.4838709677	2.08211584986131e-09\\
-21888.5630498534	2.26426456790727e-10\\
-21841.642228739	5.24434084630197e-09\\
-21794.7214076246	1.27381730534595e-08\\
-21747.8005865103	1.60583318226674e-08\\
-21700.8797653959	1.22303588679944e-08\\
-21653.9589442815	4.64672567009898e-09\\
-21607.0381231672	8.82582836955969e-11\\
-21560.1173020528	2.68958625721717e-09\\
-21513.1964809384	1.0218496921197e-08\\
-21466.275659824	1.60166114862768e-08\\
-21419.3548387097	1.49153252430842e-08\\
-21372.4340175953	7.87127480657984e-09\\
-21325.5131964809	1.16463880756861e-09\\
-21278.5923753666	8.29434876339403e-10\\
-21231.6715542522	7.25483990276086e-09\\
-21184.7507331378	1.47871751053481e-08\\
-21137.8299120235	1.67327625036502e-08\\
-21090.9090909091	1.13353279077496e-08\\
-21043.9882697947	3.39076808187827e-09\\
-20997.0674486804	1.62501053267456e-11\\
-20950.146627566	4.30539747317656e-09\\
-20903.2258064516	1.25252455180207e-08\\
-20856.3049853372	1.73937068099602e-08\\
-20809.3841642229	1.45608837462681e-08\\
-20762.4633431085	6.52624412416381e-09\\
-20715.5425219941	4.64976215478828e-10\\
-20668.6217008798	1.85253810259587e-09\\
-20621.7008797654	9.55162314954285e-09\\
-20574.780058651	1.67691963336907e-08\\
-20527.8592375367	1.70810751984579e-08\\
-20480.9384164223	1.01816649117118e-08\\
-20434.0175953079	2.21344803691576e-09\\
-20387.0967741935	3.31787369708964e-10\\
-20340.1759530792	6.31367734625314e-09\\
-20293.2551319648	1.49184955227334e-08\\
-20246.3343108504	1.85109451249227e-08\\
-20199.4134897361	1.3868941476874e-08\\
-20152.4926686217	5.10699400944169e-09\\
-20105.5718475073	6.39280456193345e-11\\
-20058.651026393	3.32529820655912e-09\\
-20011.7302052786	1.20907617712596e-08\\
-19964.8093841642	1.86098927922294e-08\\
-19917.8885630499	1.7069134654615e-08\\
-19870.9677419355	8.81094478723086e-09\\
-19824.0469208211	1.20047648353443e-09\\
-19777.1260997067	1.09358041374903e-09\\
-19730.2052785924	8.69845048952297e-09\\
-19683.284457478	1.73260551953019e-08\\
-19636.3636363636	1.93087029308977e-08\\
-19589.4428152493	1.28507306549207e-08\\
-19542.5219941349	3.69137613258697e-09\\
-19495.6011730205	4.23947470129056e-11\\
-19448.6803519062	5.26483820535347e-09\\
-19401.7595307918	1.48150247257085e-08\\
-19354.8387096774	2.02333065351424e-08\\
-19307.917888563	1.66751573579267e-08\\
-19260.9970674487	7.28095341514147e-09\\
-19214.0762463343	4.44023175579727e-10\\
-19167.1554252199	2.35199826980033e-09\\
-19120.2346041056	1.14286742805807e-08\\
-19073.3137829912	1.96680709240235e-08\\
-19026.3929618768	1.97347813139601e-08\\
-18979.4721407625	1.15343449190619e-08\\
-18885.6304985337	4.78805070896611e-10\\
-18838.7096774194	7.67403677455527e-09\\
-18791.788856305	1.76545870726849e-08\\
-18744.8680351906	2.15646203932857e-08\\
-18697.9472140762	1.58922321113116e-08\\
-18651.0263929619	5.66588926955118e-09\\
-18604.1055718475	4.05584053743759e-11\\
-18557.1847507331	4.1466021163547e-09\\
-18510.2639296188	1.4457803771731e-08\\
-18463.3431085044	2.18591810683732e-08\\
-18416.42228739	1.97474195142822e-08\\
-18369.5014662757	9.96611427772032e-09\\
-18322.5806451613	1.24224845904081e-09\\
-18275.6598240469	1.44641995133746e-09\\
-18228.7390029326	1.05400114438394e-08\\
-18181.8181818182	2.05282187311429e-08\\
-18134.8973607038	2.25331956538859e-08\\
-18087.9765395894	1.4730711249046e-08\\
-18041.0557184751	4.05617305845558e-09\\
-17994.1348973607	8.84223138747357e-11\\
-17947.2140762463	6.50380340246039e-09\\
-17900.293255132	1.77238063584767e-08\\
-17853.3724340176	2.38110865541511e-08\\
-17806.4516129032	1.93183639707894e-08\\
-17759.5307917889	8.21166538979011e-09\\
-17712.6099706745	4.1990976618359e-10\\
-17665.6891495601	3.01058284640455e-09\\
-17618.7683284457	1.3832950713401e-08\\
-17571.8475073314	2.33449601308756e-08\\
-17524.926686217	2.30758495135726e-08\\
-17478.0058651026	1.32204283445573e-08\\
-17431.0850439883	2.55786963649986e-09\\
-17384.1642228739	6.85180768683946e-10\\
-17337.2434017595	9.43471829020149e-09\\
-17290.3225806452	2.11496576900453e-08\\
-17243.4017595308	2.54355130176038e-08\\
-17196.4809384164	1.84359451551755e-08\\
-17149.5601173021	6.35665121761528e-09\\
-17102.6392961877	1.98407531714791e-11\\
-17055.7184750733	5.22599047484223e-09\\
-17008.7976539589	1.75053477763237e-08\\
-16961.8768328446	2.60062822483292e-08\\
-16914.9560117302	2.31403303283351e-08\\
-16868.0351906158	1.14127120137596e-08\\
-16821.1143695015	1.29174398062791e-09\\
-16774.1935483871	1.92483586787159e-09\\
-16727.2727272727	1.29332859474327e-08\\
-16680.3519061584	2.46443402363855e-08\\
-16633.431085044	2.66475360700914e-08\\
-16586.5102639296	1.7108125071599e-08\\
-16539.5894428152	4.50715310333232e-09\\
-16492.6686217009	1.64344320310107e-10\\
-16445.7478005865	8.13405887827855e-09\\
-16398.8269794721	2.14916791767361e-08\\
-16351.9061583578	2.84087214914391e-08\\
-16304.9853372434	2.2688954938905e-08\\
-16258.064516129	9.38218418860919e-09\\
-16211.1436950147	3.91992325798081e-10\\
-16164.2228739003	3.89496637686729e-09\\
-16117.3020527859	1.69747116042045e-08\\
-16070.3812316716	2.81043562746152e-08\\
-16023.4604105572	2.73692425401226e-08\\
-15976.5395894428	1.53654950039361e-08\\
-15929.6187683284	2.78978123823597e-09\\
-15882.6979472141	9.78155331374201e-10\\
-15835.7771260997	1.1760457977683e-08\\
-15788.8563049853	2.57085612148816e-08\\
-15741.935483871	3.04469761158933e-08\\
-15695.0146627566	2.17023852228689e-08\\
-15648.0938416422	7.22835604053321e-09\\
-15601.1730205279	4.71078873603341e-12\\
-15554.2521994135	6.67387855185638e-09\\
-15507.3313782991	2.15142862506303e-08\\
-15460.4105571848	3.14157588473143e-08\\
-15413.4897360704	2.75337080341024e-08\\
-15366.568914956	1.32642275720533e-08\\
-15319.6480938416	1.35152708337457e-09\\
-15272.7272727273	2.58584445984407e-09\\
-15225.8064516129	1.61128830946149e-08\\
-15178.8856304985	3.00554058482492e-08\\
-15131.9648093842	3.20174519259081e-08\\
-15085.0439882698	2.01835425967106e-08\\
-15038.1231671554	5.07706708334125e-09\\
-14991.2023460411	2.86187395380787e-10\\
-14944.2815249267	1.03278506209546e-08\\
-14897.3607038123	2.64866235683659e-08\\
-14850.4398826979	3.44566987958346e-08\\
-14803.5190615836	2.70892871866185e-08\\
-14756.5982404692	1.08890216966466e-08\\
-14709.6774193548	3.59454821828681e-10\\
-14662.7565982405	5.10920541171401e-09\\
-14615.8357771261	2.11791241206796e-08\\
-14568.9149560117	3.44155870463105e-08\\
-14521.9941348974	3.30222532590576e-08\\
-14475.073313783	1.81616924549671e-08\\
-14428.1524926686	3.08185523550224e-09\\
-14381.2316715543	1.40113023601688e-09\\
-14334.3108504399	1.490855453272e-08\\
-14287.3900293255	3.18053414137659e-08\\
-14240.4692082111	3.71004851666891e-08\\
-14193.5483870968	2.60042409284694e-08\\
-14146.6275659824	8.35612719061392e-09\\
-14099.706744868	3.60486056272737e-13\\
-14052.7859237537	8.66473671352555e-09\\
-14005.8651026393	2.6925481990283e-08\\
-13958.9442815249	3.86581191950147e-08\\
-13912.0234604106	3.33736339314251e-08\\
-13865.1026392962	1.56967799820872e-08\\
-13818.1818181818	1.42540146204612e-09\\
-13771.2609970674	3.52097605968305e-09\\
-13724.3401759531	2.04507363571638e-08\\
-13677.4193548387	3.73632002373573e-08\\
-13630.4985337243	3.92191627790634e-08\\
-13583.57771261	2.42717699981792e-08\\
-13536.6568914956	5.81660281311295e-09\\
-13489.7360703812	4.80580456965759e-10\\
-13442.8152492669	1.33613436545961e-08\\
-13395.8944281525	3.32955676675219e-08\\
-13348.9736070381	4.26398428464821e-08\\
-13302.0527859238	3.29989568318478e-08\\
-13255.1319648094	1.28841713297172e-08\\
-13208.211143695	3.21267974083335e-10\\
-13161.2903225806	6.82246263937947e-09\\
-13114.3695014663	2.69702319474117e-08\\
-13067.4486803519	4.30326947383211e-08\\
-13020.5278592375	4.06876285466308e-08\\
-12973.6070381232	2.19155935686072e-08\\
-12926.6862170088	3.46005634314609e-09\\
-12879.7653958944	2.02620121269168e-09\\
-12832.8445747801	1.92984026128656e-08\\
-12785.9237536657	4.0209500706313e-08\\
-12739.0029325513	4.62081123007468e-08\\
-12692.082111437	3.18462678628312e-08\\
-12645.1612903226	9.86006080095412e-09\\
-12598.2404692082	1.62533751602375e-11\\
-12551.3196480938	1.14866991111604e-08\\
-12504.3988269795	3.44623872226517e-08\\
-12457.4780058651	4.86670784691333e-08\\
-12410.5571847507	4.13883434904203e-08\\
-12363.6363636364	1.89963186330952e-08\\
-12316.715542522	1.51925143295691e-09\\
-12269.7947214076	4.88401874838794e-09\\
-12222.8739002933	2.65643512500772e-08\\
-12175.9530791789	4.75635021315855e-08\\
-12129.0322580645	4.92039880439729e-08\\
-12082.1114369501	2.98906820904132e-08\\
-12035.1906158358	6.80796322960748e-09\\
-11988.2697947214	7.9383723537795e-10\\
-11941.348973607	1.76990385259504e-08\\
-11894.4281524927	4.29009428777117e-08\\
-11847.5073313783	5.41017850601774e-08\\
-11800.5865102639	4.12167604091959e-08\\
-11753.6656891496	1.56190099729607e-08\\
-11706.7448680352	2.76166518129845e-10\\
-11659.8240469208	9.32373425221981e-09\\
-11612.9032258065	3.5235997885552e-08\\
-11565.9824046921	5.52298490911124e-08\\
-11519.0615835777	5.1465961154726e-08\\
-11472.1407624633	2.71419275685875e-08\\
-11425.219941349	3.96695071779788e-09\\
-11378.2991202346	2.97940630475905e-09\\
-11331.3782991202	2.56501582014599e-08\\
-11284.4574780059	5.22374317145674e-08\\
-11237.5366568915	5.9156195932314e-08\\
-11190.6158357771	4.00876947645693e-08\\
-11143.6950146628	1.19426260817229e-08\\
-11096.7741935484	7.03941659286391e-11\\
-11049.853372434	1.56420779646231e-08\\
-11002.9325513196	4.53787882947894e-08\\
-10956.0117302053	6.3057125076491e-08\\
-10909.0909090909	5.28342871030355e-08\\
-10862.1700879765	2.36536295575937e-08\\
-10815.2492668622	1.64266439144788e-09\\
-10768.3284457478	6.94828889460671e-09\\
-10721.4076246334	3.55389065891409e-08\\
-10674.4868035191	6.24011234321204e-08\\
-10627.5659824047	6.36356958232267e-08\\
-10580.6451612903	3.79430477440727e-08\\
-10533.724340176	8.19239173405136e-09\\
-10486.8035190616	1.31124055246936e-09\\
-10439.8826979472	2.41698692921186e-08\\
-10392.9618768328	5.7048339234264e-08\\
-10346.0410557185	7.08706571557452e-08\\
-10299.1202346041	5.31561017090129e-08\\
-10252.1994134897	1.9535304730583e-08\\
-10205.2785923754	2.22721883605852e-10\\
-10158.357771261	1.31381489186315e-08\\
-10111.4369501466	4.75777154005631e-08\\
-10064.5161290323	7.32995369759685e-08\\
-10017.5953079179	6.73339598903793e-08\\
-9970.67448680352	3.47615976476198e-08\\
-9923.75366568915	4.67683602643126e-09\\
-9876.83284457478	4.49563950704745e-09\\
-9829.91202346041	3.52835243535808e-08\\
-9782.99120234604	7.02936640422158e-08\\
-9736.07038123167	7.84729767506899e-08\\
-9689.1495601173	5.22916668486766e-08\\
-9642.22873900293	1.49686789655161e-08\\
-9595.30791788856	1.99105197745068e-10\\
-9548.38709677419	2.20681417110534e-08\\
-9501.46627565983	6.20017644480956e-08\\
-9454.54545454546	8.48195248875459e-08\\
-9407.62463343109	7.00354409845519e-08\\
-9360.70381231671	3.05721041137414e-08\\
-9313.78299120235	1.81232279334637e-09\\
-9266.86217008798	1.02352418255839e-08\\
-9219.94134897361	4.94233488273995e-08\\
-9173.02052785924	8.51627099322532e-08\\
-9126.09970674487	8.56440665971707e-08\\
-9079.1788856305	5.01224129420864e-08\\
-9032.25806451613	1.02310426635353e-08\\
-8985.33724340176	2.20173696786503e-09\\
-8938.41642228739	3.43686567840311e-08\\
-8891.49560117302	7.90837358538688e-08\\
-8844.57478005865	9.68289179120287e-08\\
-8797.65395894428	7.15180963156464e-08\\
-8750.73313782991	2.54719000942624e-08\\
-8703.81231671554	1.59789017423426e-10\\
-8656.89149560118	1.92981659174138e-08\\
-8609.97067448681	6.71223569857598e-08\\
-8563.04985337243	1.01711159478848e-07\\
-8516.12903225806	9.21410532873322e-08\\
-8469.2082111437	4.65620370905772e-08\\
-8422.28739002933	5.73029561687681e-09\\
-8375.36656891496	7.05027848601463e-09\\
-8328.44574780059	5.08279343091034e-08\\
-8281.52492668622	9.915711322473e-08\\
-8234.60410557185	1.09178701999831e-07\\
-8187.68328445748	7.15565591570492e-08\\
-8140.76246334311	1.96548341196358e-08\\
-8093.84164222874	4.81766970056708e-10\\
-8046.92082111437	3.26746481522892e-08\\
-8000	8.90518215609265e-08\\
-7953.07917888563	1.20016034114531e-07\\
-7906.15835777126	9.76968552571976e-08\\
-7859.23753665689	4.15730396261852e-08\\
-7812.31671554252	2.05986737664702e-09\\
-7765.39589442815	1.58362804980635e-08\\
-7718.47507331378	7.24662566046066e-08\\
-7671.55425219941	1.2265297657208e-07\\
-7624.63343108504	1.21705684156255e-07\\
-7577.71260997068	6.99263334112556e-08\\
-7530.79178885631	1.34555960434486e-08\\
-7483.87096774194	3.83397241093388e-09\\
-7436.95014662757	5.16827122206275e-08\\
-7390.0293255132	1.16094249944747e-07\\
-7343.10850439883	1.40196698415707e-07\\
-7296.18768328446	1.02017073731775e-07\\
-7249.26686217009	3.51932884357787e-08\\
-7202.34604105572	8.84004804483417e-11\\
-7155.42521994135	3.00573904343058e-08\\
-7108.50439882698	1.00656654750734e-07\\
-7061.58357771261	1.50158939668907e-07\\
-7014.66275659824	1.34235588918774e-07\\
-6967.74193548387	6.64104187979461e-08\\
-6920.8211143695	7.42366372834707e-09\\
-6873.90029325513	1.17176997981171e-08\\
-6826.97947214076	7.81431421444177e-08\\
-6780.05865102639	1.49463018242459e-07\\
-6733.13782991202	1.62447859198343e-07\\
-6686.21700879766	1.04774079261669e-07\\
-6639.29618768329	2.75817127597968e-08\\
-6592.37536656891	1.11237119001193e-09\\
-6545.45454545455	5.18441260234193e-08\\
-6498.53372434018	1.37327156090676e-07\\
-6451.61290322581	1.82516029394662e-07\\
-6404.69208211144	1.46586514037284e-07\\
-6357.77126099707	6.08116648238379e-08\\
-6310.8504398827	2.4528321989262e-09\\
-6263.92961876833	2.63434415265087e-08\\
-6217.00879765396	1.14679480196104e-07\\
-6170.08797653959	1.90904917520172e-07\\
-6123.16715542522	1.87093075392426e-07\\
-6076.24633431085	1.05596749692858e-07\\
-6029.32551319648	1.91019136282015e-08\\
-5982.40469208211	7.12755296178205e-09\\
-5935.48387096774	8.4358635696611e-08\\
-5888.56304985337	1.85312677501665e-07\\
-5841.64222873901	2.20983874565184e-07\\
-5794.72140762463	1.58572662511318e-07\\
-5747.80058651026	5.29786237383229e-08\\
-5700.87976539589	1.9845244546515e-11\\
-5653.95894428153	5.11078166359495e-08\\
-5607.03812316716	1.65254913334039e-07\\
-5560.11730205279	2.43057766744156e-07\\
-5513.19648093842	2.14675426070204e-07\\
-5466.27565982405	1.04053635326746e-07\\
-5419.35483870968	1.04844947438441e-08\\
-5372.43401759531	2.13366633531731e-08\\
-5325.51319648094	1.32530385388584e-07\\
-5278.59237536657	2.48996955101481e-07\\
-5231.6715542522	2.67536086021841e-07\\
-5184.75073313783	1.7000825906952e-07\\
-5137.82991202346	4.28638159751534e-08\\
-5090.90909090909	2.6492948793743e-09\\
-5043.98826979472	9.15037719720006e-08\\
-4997.06744868035	2.36185265996162e-07\\
-4950.14662756598	3.10116138417367e-07\\
-4903.22580645161	2.46119962278863e-07\\
-4856.30498533724	9.96266647152229e-08\\
-4809.38416422287	3.16251468115288e-09\\
-4762.46334310851	4.91566264838672e-08\\
-4715.54252199414	2.04495642344742e-07\\
-4668.62170087976	3.35554711821561e-07\\
-4621.7008797654	3.25420990386005e-07\\
-4574.78005865103	1.80711552506629e-07\\
-4527.85923753666	3.06630155110207e-08\\
-4480.93841642229	1.48806843873657e-08\\
-4434.01759530792	1.56974824384879e-07\\
-4387.09677419355	3.38180001199365e-07\\
-4340.17595307918	3.99155635973961e-07\\
-4293.25513196481	2.83045557471193e-07\\
-4246.33431085044	9.16731636359196e-08\\
-4199.41348973607	1.62550102743481e-11\\
-4152.4926686217	1.00367331998179e-07\\
-4105.57184750733	3.14627701904653e-07\\
-4058.65102639296	4.57539811020391e-07\\
-4011.73020527859	4.00308986401697e-07\\
-3964.80938416422	1.90508801534068e-07\\
-3917.88856304986	1.71703625145706e-08\\
-3870.96774193548	4.54484435873974e-08\\
-3824.04692082111	2.65013625162872e-07\\
-3777.12609970674	4.90809144346716e-07\\
-3730.20527859238	5.2299751595627e-07\\
-3683.28445747801	3.2841705344907e-07\\
-3636.36363636364	7.93824118448216e-08\\
-3589.44281524927	7.17944440942902e-09\\
-3542.5219941349	1.94133867823101e-07\\
-3495.60117302053	4.90536842629885e-07\\
-3448.68035190616	6.38811448915435e-07\\
-3401.75953079179	5.02875515370924e-07\\
-3354.83870967742	1.99238140808355e-07\\
-3307.91788856305	4.76357041172684e-09\\
-3260.99706744868	1.12726717102963e-07\\
-3214.07624633431	4.51223259052829e-07\\
-3167.15542521994	7.3338167381664e-07\\
-3120.23460410557	7.07115126919184e-07\\
-3073.3137829912	3.88066429072602e-07\\
-3026.39296187683	6.17813756141332e-08\\
-2979.47214076246	3.89517853391433e-08\\
-2932.55131964809	3.72230763671006e-07\\
-2885.63049853372	7.91334987146907e-07\\
-2838.70967741936	9.29896146639201e-07\\
-2791.78885630499	6.55394182394294e-07\\
-2744.86803519062	2.06753225904335e-07\\
-2697.94721407625	4.82361104702807e-10\\
-2651.02639296188	2.60315443678972e-07\\
-2604.10557184751	7.9773106365113e-07\\
-2557.18475073314	1.1555174620202e-06\\
-2510.26392961877	1.00889958515531e-06\\
-2463.3431085044	4.74798625732497e-07\\
-2416.42228739003	3.81113855783707e-08\\
-2369.50146627566	1.33427593616724e-07\\
-2322.58064516129	7.40027124576974e-07\\
-2275.65982404692	1.3639674297254e-06\\
-2228.73900293255	1.45523988441598e-06\\
-2181.81818181818	9.1175935365922e-07\\
-2134.89736070381	2.12926560923834e-07\\
-2087.97653958944	2.7487186195759e-08\\
-2041.05571847507	6.11106688031465e-07\\
-1994.13489736071	1.53101283605391e-06\\
-1947.21407624634	2.00246264798814e-06\\
-1900.29325513196	1.58385747506734e-06\\
-1853.37243401759	6.22048777086201e-07\\
-1806.45161290323	1.06130891101384e-08\\
-1759.53079178886	4.15066031411637e-07\\
-1712.60997067449	1.62783145664103e-06\\
-1665.68914956012	2.66561509024641e-06\\
-1618.76832844575	2.60066590691021e-06\\
-1571.84750733138	1.4363753733379e-06\\
-1524.92668621701	2.17652420253759e-07\\
-1478.00586510264	1.81295515958513e-07\\
-1431.08504398827	1.61944072741733e-06\\
-1384.1642228739	3.47992465334824e-06\\
-1337.24340175953	4.17414237440862e-06\\
-1290.32258064516	3.001956758107e-06\\
-1243.40175953079	9.47321563795512e-07\\
-1196.48093841642	7.20550319795157e-09\\
-1149.56011730205	1.46060101651337e-06\\
-1102.63929618768	4.53574320857232e-06\\
-1055.71847507331	6.79370893211522e-06\\
-1008.79765395894	6.16114925760834e-06\\
-961.876832844573	2.9921245759359e-06\\
-914.956011730206	2.20849286322483e-07\\
-868.035190615836	1.08844578015071e-06\\
-821.114369501465	6.09349793285367e-06\\
-774.193548387098	1.18818026050399e-05\\
-727.272727272728	1.35684751446976e-05\\
-680.351906158357	9.12667644840757e-06\\
-633.431085043987	2.22794843953179e-06\\
-586.51026392962	4.37908764301979e-07\\
-539.58944281525	9.12407332062155e-06\\
-492.668621700879	2.5462144718253e-05\\
-445.747800586509	3.81699582764901e-05\\
-398.826979472142	3.53187908643942e-05\\
-351.906158357771	1.64135158052956e-05\\
-304.985337243401	2.22461737867166e-07\\
-258.064516129034	2.0626513867124e-05\\
-211.143695014664	0.000109614171797184\\
-164.222873900293	0.00027556507604712\\
-117.302052785923	0.000490008604388786\\
-70.3812316715557	0.000692955513443509\\
-23.4604105571852	0.000816701813950125\\
23.4604105571852	0.000816701813950125\\
70.3812316715557	0.000692955513443509\\
117.302052785923	0.000490008604388786\\
164.222873900293	0.00027556507604712\\
211.143695014664	0.000109614171797184\\
258.064516129034	2.0626513867124e-05\\
304.985337243401	2.22461737867166e-07\\
351.906158357771	1.64135158052956e-05\\
398.826979472142	3.53187908643942e-05\\
445.747800586509	3.81699582764901e-05\\
492.668621700879	2.5462144718253e-05\\
539.58944281525	9.12407332062155e-06\\
586.51026392962	4.37908764301979e-07\\
633.431085043987	2.22794843953179e-06\\
680.351906158357	9.12667644840757e-06\\
727.272727272728	1.35684751446976e-05\\
774.193548387098	1.18818026050399e-05\\
821.114369501465	6.09349793285367e-06\\
868.035190615836	1.08844578015071e-06\\
914.956011730206	2.20849286322483e-07\\
961.876832844573	2.9921245759359e-06\\
1008.79765395894	6.16114925760834e-06\\
1055.71847507331	6.79370893211522e-06\\
1102.63929618768	4.53574320857232e-06\\
1149.56011730205	1.46060101651337e-06\\
1196.48093841642	7.20550319795157e-09\\
1243.40175953079	9.47321563795512e-07\\
1290.32258064516	3.001956758107e-06\\
1337.24340175953	4.17414237440862e-06\\
1384.1642228739	3.47992465334824e-06\\
1431.08504398827	1.61944072741733e-06\\
1478.00586510264	1.81295515958513e-07\\
1524.92668621701	2.17652420253759e-07\\
1571.84750733138	1.4363753733379e-06\\
1618.76832844575	2.60066590691021e-06\\
1665.68914956012	2.66561509024641e-06\\
1712.60997067449	1.62783145664103e-06\\
1759.53079178886	4.15066031411637e-07\\
1806.45161290323	1.06130891101384e-08\\
1853.37243401759	6.22048777086201e-07\\
1900.29325513196	1.58385747506734e-06\\
1947.21407624634	2.00246264798814e-06\\
1994.13489736071	1.53101283605391e-06\\
2041.05571847507	6.11106688031465e-07\\
2087.97653958944	2.7487186195759e-08\\
2134.89736070381	2.12926560923834e-07\\
2181.81818181818	9.1175935365922e-07\\
2228.73900293255	1.45523988441598e-06\\
2275.65982404692	1.3639674297254e-06\\
2322.58064516129	7.40027124576974e-07\\
2369.50146627566	1.33427593616724e-07\\
2416.42228739003	3.81113855783707e-08\\
2463.3431085044	4.74798625732497e-07\\
2510.26392961877	1.00889958515531e-06\\
2557.18475073314	1.1555174620202e-06\\
2604.10557184751	7.9773106365113e-07\\
2651.02639296188	2.60315443678972e-07\\
2697.94721407625	4.82361104702807e-10\\
2744.86803519062	2.06753225904335e-07\\
2791.78885630499	6.55394182394294e-07\\
2838.70967741936	9.29896146639201e-07\\
2885.63049853372	7.91334987146907e-07\\
2932.55131964809	3.72230763671006e-07\\
2979.47214076246	3.89517853391433e-08\\
3026.39296187683	6.17813756141332e-08\\
3073.3137829912	3.88066429072602e-07\\
3120.23460410557	7.07115126919184e-07\\
3167.15542521994	7.3338167381664e-07\\
3214.07624633431	4.51223259052829e-07\\
3260.99706744868	1.12726717102963e-07\\
3307.91788856305	4.76357041172684e-09\\
3354.83870967742	1.99238140808355e-07\\
3401.75953079179	5.02875515370924e-07\\
3448.68035190616	6.38811448915435e-07\\
3495.60117302053	4.90536842629885e-07\\
3542.5219941349	1.94133867823101e-07\\
3589.44281524927	7.17944440942902e-09\\
3636.36363636364	7.93824118448216e-08\\
3683.28445747801	3.2841705344907e-07\\
3730.20527859238	5.2299751595627e-07\\
3777.12609970674	4.90809144346716e-07\\
3824.04692082111	2.65013625162872e-07\\
3870.96774193548	4.54484435873974e-08\\
3917.88856304986	1.71703625145706e-08\\
3964.80938416422	1.90508801534068e-07\\
4011.73020527859	4.00308986401697e-07\\
4058.65102639296	4.57539811020391e-07\\
4105.57184750733	3.14627701904653e-07\\
4152.4926686217	1.00367331998179e-07\\
4199.41348973607	1.62550102743481e-11\\
4246.33431085044	9.16731636359196e-08\\
4293.25513196481	2.83045557471193e-07\\
4340.17595307918	3.99155635973961e-07\\
4387.09677419355	3.38180001199365e-07\\
4434.01759530792	1.56974824384879e-07\\
4480.93841642229	1.48806843873657e-08\\
4527.85923753666	3.06630155110207e-08\\
4574.78005865103	1.80711552506629e-07\\
4621.7008797654	3.25420990386005e-07\\
4668.62170087976	3.35554711821561e-07\\
4715.54252199414	2.04495642344742e-07\\
4762.46334310851	4.91566264838672e-08\\
4809.38416422287	3.16251468115288e-09\\
4856.30498533724	9.96266647152229e-08\\
4903.22580645161	2.46119962278863e-07\\
4950.14662756598	3.10116138417367e-07\\
4997.06744868035	2.36185265996162e-07\\
5043.98826979472	9.15037719720006e-08\\
5090.90909090909	2.6492948793743e-09\\
5137.82991202346	4.28638159751534e-08\\
5184.75073313783	1.7000825906952e-07\\
5231.6715542522	2.67536086021841e-07\\
5278.59237536657	2.48996955101481e-07\\
5325.51319648094	1.32530385388584e-07\\
5372.43401759531	2.13366633531731e-08\\
5419.35483870968	1.04844947438441e-08\\
5466.27565982405	1.04053635326746e-07\\
5513.19648093842	2.14675426070204e-07\\
5560.11730205279	2.43057766744156e-07\\
5607.03812316716	1.65254913334039e-07\\
5653.95894428153	5.11078166359495e-08\\
5700.87976539589	1.9845244546515e-11\\
5747.80058651026	5.29786237383229e-08\\
5794.72140762463	1.58572662511318e-07\\
5841.64222873901	2.20983874565184e-07\\
5888.56304985337	1.85312677501665e-07\\
5935.48387096774	8.4358635696611e-08\\
5982.40469208211	7.12755296178205e-09\\
6029.32551319648	1.91019136282015e-08\\
6076.24633431085	1.05596749692858e-07\\
6123.16715542522	1.87093075392426e-07\\
6170.08797653959	1.90904917520172e-07\\
6217.00879765396	1.14679480196104e-07\\
6263.92961876833	2.63434415265087e-08\\
6310.8504398827	2.4528321989262e-09\\
6357.77126099707	6.08116648238379e-08\\
6404.69208211144	1.46586514037284e-07\\
6451.61290322581	1.82516029394662e-07\\
6498.53372434018	1.37327156090676e-07\\
6545.45454545455	5.18441260234193e-08\\
6592.37536656891	1.11237119001193e-09\\
6639.29618768329	2.75817127597968e-08\\
6686.21700879766	1.04774079261669e-07\\
6733.13782991202	1.62447859198343e-07\\
6780.05865102639	1.49463018242459e-07\\
6826.97947214076	7.81431421444177e-08\\
6873.90029325513	1.17176997981171e-08\\
6920.8211143695	7.42366372834707e-09\\
6967.74193548387	6.64104187979461e-08\\
7014.66275659824	1.34235588918774e-07\\
7061.58357771261	1.50158939668907e-07\\
7108.50439882698	1.00656654750734e-07\\
7155.42521994135	3.00573904343058e-08\\
7202.34604105572	8.84004804483417e-11\\
7249.26686217009	3.51932884357787e-08\\
7296.18768328446	1.02017073731775e-07\\
7343.10850439883	1.40196698415707e-07\\
7390.0293255132	1.16094249944747e-07\\
7436.95014662757	5.16827122206275e-08\\
7483.87096774194	3.83397241093388e-09\\
7530.79178885631	1.34555960434486e-08\\
7577.71260997068	6.99263334112556e-08\\
7624.63343108504	1.21705684156255e-07\\
7671.55425219941	1.2265297657208e-07\\
7718.47507331378	7.24662566046066e-08\\
7765.39589442815	1.58362804980635e-08\\
7812.31671554252	2.05986737664702e-09\\
7859.23753665689	4.15730396261852e-08\\
7906.15835777126	9.76968552571976e-08\\
7953.07917888563	1.20016034114531e-07\\
8000	8.90518215609265e-08\\
8046.92082111437	3.26746481522892e-08\\
8093.84164222874	4.81766970056708e-10\\
8140.76246334311	1.96548341196358e-08\\
8187.68328445748	7.15565591570492e-08\\
8234.60410557185	1.09178701999831e-07\\
8281.52492668622	9.915711322473e-08\\
8328.44574780059	5.08279343091034e-08\\
8375.36656891496	7.05027848601463e-09\\
8422.28739002933	5.73029561687681e-09\\
8469.2082111437	4.65620370905772e-08\\
8516.12903225806	9.21410532873322e-08\\
8563.04985337243	1.01711159478848e-07\\
8609.97067448681	6.71223569857598e-08\\
8656.89149560118	1.92981659174138e-08\\
8703.81231671554	1.59789017423426e-10\\
8750.73313782991	2.54719000942624e-08\\
8797.65395894428	7.15180963156464e-08\\
8844.57478005865	9.68289179120287e-08\\
8891.49560117302	7.90837358538688e-08\\
8938.41642228739	3.43686567840311e-08\\
8985.33724340176	2.20173696786503e-09\\
9032.25806451613	1.02310426635353e-08\\
9079.1788856305	5.01224129420864e-08\\
9126.09970674487	8.56440665971707e-08\\
9173.02052785924	8.51627099322532e-08\\
9219.94134897361	4.94233488273995e-08\\
9266.86217008798	1.02352418255839e-08\\
9313.78299120235	1.81232279334637e-09\\
9360.70381231671	3.05721041137414e-08\\
9407.62463343109	7.00354409845519e-08\\
9454.54545454546	8.48195248875459e-08\\
9501.46627565983	6.20017644480956e-08\\
9548.38709677419	2.20681417110534e-08\\
9595.30791788856	1.99105197745068e-10\\
9642.22873900293	1.49686789655161e-08\\
9689.1495601173	5.22916668486766e-08\\
9736.07038123167	7.84729767506899e-08\\
9782.99120234604	7.02936640422158e-08\\
9829.91202346041	3.52835243535808e-08\\
9876.83284457478	4.49563950704745e-09\\
9923.75366568915	4.67683602643126e-09\\
9970.67448680352	3.47615976476198e-08\\
10017.5953079179	6.73339598903793e-08\\
10064.5161290323	7.32995369759685e-08\\
10111.4369501466	4.75777154005631e-08\\
10158.357771261	1.31381489186315e-08\\
10205.2785923754	2.22721883605852e-10\\
10252.1994134897	1.9535304730583e-08\\
10299.1202346041	5.31561017090129e-08\\
10346.0410557185	7.08706571557452e-08\\
10392.9618768328	5.7048339234264e-08\\
10439.8826979472	2.41698692921186e-08\\
10486.8035190616	1.31124055246936e-09\\
10533.724340176	8.19239173405136e-09\\
10580.6451612903	3.79430477440727e-08\\
10627.5659824047	6.36356958232267e-08\\
10674.4868035191	6.24011234321204e-08\\
10721.4076246334	3.55389065891409e-08\\
10768.3284457478	6.94828889460671e-09\\
10815.2492668622	1.64266439144788e-09\\
10862.1700879765	2.36536295575937e-08\\
10909.0909090909	5.28342871030355e-08\\
10956.0117302053	6.3057125076491e-08\\
11002.9325513196	4.53787882947894e-08\\
11049.853372434	1.56420779646231e-08\\
11096.7741935484	7.03941659286391e-11\\
11143.6950146628	1.19426260817229e-08\\
11190.6158357771	4.00876947645693e-08\\
11237.5366568915	5.9156195932314e-08\\
11284.4574780059	5.22374317145674e-08\\
11331.3782991202	2.56501582014599e-08\\
11378.2991202346	2.97940630475905e-09\\
11425.219941349	3.96695071779788e-09\\
11472.1407624633	2.71419275685875e-08\\
11519.0615835777	5.1465961154726e-08\\
11565.9824046921	5.52298490911124e-08\\
11612.9032258065	3.5235997885552e-08\\
11659.8240469208	9.32373425221981e-09\\
11706.7448680352	2.76166518129845e-10\\
11753.6656891496	1.56190099729607e-08\\
11800.5865102639	4.12167604091959e-08\\
11847.5073313783	5.41017850601774e-08\\
11894.4281524927	4.29009428777117e-08\\
11941.348973607	1.76990385259504e-08\\
11988.2697947214	7.9383723537795e-10\\
12035.1906158358	6.80796322960748e-09\\
12082.1114369501	2.98906820904132e-08\\
12129.0322580645	4.92039880439729e-08\\
12175.9530791789	4.75635021315855e-08\\
12222.8739002933	2.65643512500772e-08\\
12269.7947214076	4.88401874838794e-09\\
12316.715542522	1.51925143295691e-09\\
12363.6363636364	1.89963186330952e-08\\
12410.5571847507	4.13883434904203e-08\\
12457.4780058651	4.86670784691333e-08\\
12504.3988269795	3.44623872226517e-08\\
12551.3196480938	1.14866991111604e-08\\
12598.2404692082	1.62533751602375e-11\\
12645.1612903226	9.86006080095412e-09\\
12692.082111437	3.18462678628312e-08\\
12739.0029325513	4.62081123007468e-08\\
12785.9237536657	4.0209500706313e-08\\
12832.8445747801	1.92984026128656e-08\\
12879.7653958944	2.02620121269168e-09\\
12926.6862170088	3.46005634314609e-09\\
12973.6070381232	2.19155935686072e-08\\
13020.5278592375	4.06876285466308e-08\\
13067.4486803519	4.30326947383211e-08\\
13114.3695014663	2.69702319474117e-08\\
13161.2903225806	6.82246263937947e-09\\
13208.211143695	3.21267974083335e-10\\
13255.1319648094	1.28841713297172e-08\\
13302.0527859238	3.29989568318478e-08\\
13348.9736070381	4.26398428464821e-08\\
13395.8944281525	3.32955676675219e-08\\
13442.8152492669	1.33613436545961e-08\\
13489.7360703812	4.80580456965759e-10\\
13536.6568914956	5.81660281311295e-09\\
13583.57771261	2.42717699981792e-08\\
13630.4985337243	3.92191627790634e-08\\
13677.4193548387	3.73632002373573e-08\\
13724.3401759531	2.04507363571638e-08\\
13771.2609970674	3.52097605968305e-09\\
13818.1818181818	1.42540146204612e-09\\
13865.1026392962	1.56967799820872e-08\\
13912.0234604106	3.33736339314251e-08\\
13958.9442815249	3.86581191950147e-08\\
14005.8651026393	2.6925481990283e-08\\
14052.7859237537	8.66473671352555e-09\\
14099.706744868	3.60486056272737e-13\\
14146.6275659824	8.35612719061392e-09\\
14193.5483870968	2.60042409284694e-08\\
14240.4692082111	3.71004851666891e-08\\
14287.3900293255	3.18053414137659e-08\\
14334.3108504399	1.490855453272e-08\\
14381.2316715543	1.40113023601688e-09\\
14428.1524926686	3.08185523550224e-09\\
14475.073313783	1.81616924549671e-08\\
14521.9941348974	3.30222532590576e-08\\
14568.9149560117	3.44155870463105e-08\\
14615.8357771261	2.11791241206796e-08\\
14662.7565982405	5.10920541171401e-09\\
14709.6774193548	3.59454821828681e-10\\
14756.5982404692	1.08890216966466e-08\\
14803.5190615836	2.70892871866185e-08\\
14850.4398826979	3.44566987958346e-08\\
14897.3607038123	2.64866235683659e-08\\
14944.2815249267	1.03278506209546e-08\\
14991.2023460411	2.86187395380787e-10\\
15038.1231671554	5.07706708334125e-09\\
15085.0439882698	2.01835425967106e-08\\
15131.9648093842	3.20174519259081e-08\\
15178.8856304985	3.00554058482492e-08\\
15225.8064516129	1.61128830946149e-08\\
15272.7272727273	2.58584445984407e-09\\
15319.6480938416	1.35152708337457e-09\\
15366.568914956	1.32642275720533e-08\\
15413.4897360704	2.75337080341024e-08\\
15460.4105571848	3.14157588473143e-08\\
15507.3313782991	2.15142862506303e-08\\
15554.2521994135	6.67387855185638e-09\\
15601.1730205279	4.71078873603341e-12\\
15648.0938416422	7.22835604053321e-09\\
15695.0146627566	2.17023852228689e-08\\
15741.935483871	3.04469761158933e-08\\
15788.8563049853	2.57085612148816e-08\\
15835.7771260997	1.1760457977683e-08\\
15882.6979472141	9.78155331374201e-10\\
15929.6187683284	2.78978123823597e-09\\
15976.5395894428	1.53654950039361e-08\\
16023.4604105572	2.73692425401226e-08\\
16070.3812316716	2.81043562746152e-08\\
16117.3020527859	1.69747116042045e-08\\
16164.2228739003	3.89496637686729e-09\\
16211.1436950147	3.91992325798081e-10\\
16258.064516129	9.38218418860919e-09\\
16304.9853372434	2.2688954938905e-08\\
16351.9061583578	2.84087214914391e-08\\
16398.8269794721	2.14916791767361e-08\\
16445.7478005865	8.13405887827855e-09\\
16492.6686217009	1.64344320310107e-10\\
16539.5894428152	4.50715310333232e-09\\
16586.5102639296	1.7108125071599e-08\\
16633.431085044	2.66475360700914e-08\\
16680.3519061584	2.46443402363855e-08\\
16727.2727272727	1.29332859474327e-08\\
16774.1935483871	1.92483586787159e-09\\
16821.1143695015	1.29174398062791e-09\\
16868.0351906158	1.14127120137596e-08\\
16914.9560117302	2.31403303283351e-08\\
16961.8768328446	2.60062822483292e-08\\
17008.7976539589	1.75053477763237e-08\\
17055.7184750733	5.22599047484223e-09\\
17102.6392961877	1.98407531714791e-11\\
17149.5601173021	6.35665121761528e-09\\
17196.4809384164	1.84359451551755e-08\\
17243.4017595308	2.54355130176038e-08\\
17290.3225806452	2.11496576900453e-08\\
17337.2434017595	9.43471829020149e-09\\
17384.1642228739	6.85180768683946e-10\\
17431.0850439883	2.55786963649986e-09\\
17478.0058651026	1.32204283445573e-08\\
17524.926686217	2.30758495135726e-08\\
17571.8475073314	2.33449601308756e-08\\
17618.7683284457	1.3832950713401e-08\\
17665.6891495601	3.01058284640455e-09\\
17712.6099706745	4.1990976618359e-10\\
17759.5307917889	8.21166538979011e-09\\
17806.4516129032	1.93183639707894e-08\\
17853.3724340176	2.38110865541511e-08\\
17900.293255132	1.77238063584767e-08\\
17947.2140762463	6.50380340246039e-09\\
17994.1348973607	8.84223138747357e-11\\
18041.0557184751	4.05617305845558e-09\\
18087.9765395894	1.4730711249046e-08\\
18134.8973607038	2.25331956538859e-08\\
18181.8181818182	2.05282187311429e-08\\
18228.7390029326	1.05400114438394e-08\\
18275.6598240469	1.44641995133746e-09\\
18322.5806451613	1.24224845904081e-09\\
18369.5014662757	9.96611427772032e-09\\
18416.42228739	1.97474195142822e-08\\
18463.3431085044	2.18591810683732e-08\\
18510.2639296188	1.4457803771731e-08\\
18557.1847507331	4.1466021163547e-09\\
18604.1055718475	4.05584053743759e-11\\
18651.0263929619	5.66588926955118e-09\\
18697.9472140762	1.58922321113116e-08\\
18744.8680351906	2.15646203932857e-08\\
18791.788856305	1.76545870726849e-08\\
18838.7096774194	7.67403677455527e-09\\
18885.6304985337	4.78805070896611e-10\\
18979.4721407625	1.15343449190619e-08\\
19026.3929618768	1.97347813139601e-08\\
19073.3137829912	1.96680709240235e-08\\
19120.2346041056	1.14286742805807e-08\\
19167.1554252199	2.35199826980033e-09\\
19214.0762463343	4.44023175579727e-10\\
19260.9970674487	7.28095341514147e-09\\
19307.917888563	1.66751573579267e-08\\
19354.8387096774	2.02333065351424e-08\\
19401.7595307918	1.48150247257085e-08\\
19448.6803519062	5.26483820535347e-09\\
19495.6011730205	4.23947470129056e-11\\
19542.5219941349	3.69137613258697e-09\\
19589.4428152493	1.28507306549207e-08\\
19636.3636363636	1.93087029308977e-08\\
19683.284457478	1.73260551953019e-08\\
19730.2052785924	8.69845048952297e-09\\
19777.1260997067	1.09358041374903e-09\\
19824.0469208211	1.20047648353443e-09\\
19870.9677419355	8.81094478723086e-09\\
19917.8885630499	1.7069134654615e-08\\
19964.8093841642	1.86098927922294e-08\\
20011.7302052786	1.20907617712596e-08\\
20058.651026393	3.32529820655912e-09\\
20105.5718475073	6.39280456193345e-11\\
20152.4926686217	5.10699400944169e-09\\
20199.4134897361	1.3868941476874e-08\\
20246.3343108504	1.85109451249227e-08\\
20293.2551319648	1.49184955227334e-08\\
20340.1759530792	6.31367734625314e-09\\
20387.0967741935	3.31787369708964e-10\\
20434.0175953079	2.21344803691576e-09\\
20480.9384164223	1.01816649117118e-08\\
20527.8592375367	1.70810751984579e-08\\
20574.780058651	1.67691963336907e-08\\
20621.7008797654	9.55162314954285e-09\\
20668.6217008798	1.85253810259587e-09\\
20715.5425219941	4.64976215478828e-10\\
20762.4633431085	6.52624412416381e-09\\
20809.3841642229	1.45608837462681e-08\\
20856.3049853372	1.73937068099602e-08\\
20903.2258064516	1.25252455180207e-08\\
20950.146627566	4.30539747317656e-09\\
20997.0674486804	1.62501053267456e-11\\
21043.9882697947	3.39076808187827e-09\\
21090.9090909091	1.13353279077496e-08\\
21137.8299120235	1.67327625036502e-08\\
21184.7507331378	1.47871751053481e-08\\
21231.6715542522	7.25483990276086e-09\\
21278.5923753666	8.29434876339403e-10\\
21325.5131964809	1.16463880756861e-09\\
21372.4340175953	7.87127480657984e-09\\
21419.3548387097	1.49153252430842e-08\\
21466.275659824	1.60166114862768e-08\\
21513.1964809384	1.0218496921197e-08\\
21560.1173020528	2.68958625721717e-09\\
21607.0381231672	8.82582836955969e-11\\
21653.9589442815	4.64672567009898e-09\\
21700.8797653959	1.22303588679944e-08\\
21747.8005865103	1.60583318226674e-08\\
21794.7214076246	1.27381730534595e-08\\
21841.642228739	5.24434084630197e-09\\
21888.5630498534	2.26426456790727e-10\\
21935.4838709677	2.08211584986131e-09\\
21982.4046920821	9.07739000075142e-09\\
22029.3255131965	1.49362370380492e-08\\
22076.2463343109	1.44438023698404e-08\\
22123.1671554252	8.06106129461077e-09\\
22170.0879765396	1.46799548364856e-09\\
22217.008797654	4.83278248060095e-10\\
22263.9296187683	5.90390606041906e-09\\
22310.8504398827	1.28408117773952e-08\\
22357.7712609971	1.51016737666314e-08\\
22404.6920821114	1.06924885881024e-08\\
22451.6129032258	3.5505082131424e-09\\
22498.5337243402	3.53921904939615e-12\\
22545.4545454545	3.13909147118045e-09\\
22592.3753665689	1.00935731499371e-08\\
22639.2961876833	1.46408960085715e-08\\
22686.2170087977	1.27412008425165e-08\\
22733.137829912	6.1050974487918e-09\\
22780.0586510264	6.29355507531859e-10\\
22826.9794721408	1.1334500672449e-09\\
22873.9002932551	7.09467046867304e-09\\
22920.8211143695	1.31554323921928e-08\\
22967.7419354839	1.39138212944516e-08\\
23014.6627565982	8.71430081514548e-09\\
23061.5835777126	2.1903706296943e-09\\
23108.504398827	1.12567711644758e-10\\
23155.4252199413	4.2618721256523e-09\\
23202.3460410557	1.08826063181257e-08\\
23249.2668621701	1.40578011024185e-08\\
23296.1876832845	1.0973980823713e-08\\
23343.1085043988	4.3912848758425e-09\\
23390.0293255132	1.50878548329001e-10\\
23436.9501466276	1.96999445480361e-09\\
23483.8709677419	8.16225373893329e-09\\
23530.7917888563	1.31764034101032e-08\\
23577.7126099707	1.255038148357e-08\\
23624.633431085	6.85998003453254e-09\\
23671.5542521994	1.16819700264575e-09\\
23718.4750733138	4.99335130363089e-10\\
23765.3958944282	5.38323193650717e-09\\
23812.3167155425	1.14207989811639e-08\\
23859.2375366569	1.32244628263452e-08\\
23906.1583577713	9.20430684265351e-09\\
23953.0791788856	2.94843294621119e-09\\
24000	1.26631196965301e-36\\
};
\addlegendentry{Théorique}

\end{axis}
\end{tikzpicture}%



\subsection{Génération d'un signal modulé en fréquence}

Pour construire notre signal modulé en fréquence nous allons nous baser sur le signal NRZ et sur une simple sinusoïdale. Lorsque le signal NRZ vaut 0, la fréquence de la sinusoïdale sera $F_0=\SI{1180}{Hz}$ et lorsqu'il vaut 1 la sinusoïdale sera de fréquence $F_1=\SI{980}{Hz}$. Au final, notre signal modulé suit la formule suivante:
\[
\operatorname{modulé}(t)=\operatorname{NRZ}(t) \cos(2\pi F_1 t + \phi_1) + (1-\operatorname{NRZ}(t)) \cos(2 \pi F_0 t + \phi_1)
\]
$\phi_1$ et $\phi_2$ étant des déphasages tirés aléatoirement dans $[0, 2\pi]$



On obtient ainsi le signal suivant:

%%%% This file was created by matlab2tikz.
%
%The latest updates can be retrieved from
%  http://www.mathworks.com/matlabcentral/fileexchange/22022-matlab2tikz-matlab2tikz
%where you can also make suggestions and rate matlab2tikz.
%
\definecolor{mycolor1}{rgb}{0.00000,0.44700,0.74100}%
%
\begin{tikzpicture}

\begin{axis}[%
width=4.521in,
height=3.559in,
at={(0.758in,0.488in)},
scale only axis,
unbounded coords=jump,
xmin=0,
xmax=0.09,
xlabel style={font=\color{white!15!black}},
xlabel={Temps [s]},
ymin=-0.1,
ymax=1.1,
ylabel style={font=\color{white!15!black}},
ylabel={Amplitude},
axis background/.style={fill=white},
title style={font=\bfseries},
title={NRZ modulé en fréquence}
]
\addplot [color=mycolor1, forget plot]
  table[row sep=crcr]{%
0	0.395553951938314\\
2.08333333333588e-05	0.929135973038245\\
4.16666666667176e-05	0.918442742421095\\
7.85542774606185e-05	-0.22\\
nan	nan\\
0.000149909614509891	-0.22\\
0.000187500000000007	0.929135973038245\\
0.000208333333333366	0.918442742421095\\
0.000245220944127267	-0.22\\
nan	nan\\
0.000316576281176539	-0.22\\
0.000354166666666655	0.929135973038245\\
0.000375000000000014	0.918442742421095\\
0.000411887610794026	-0.22\\
nan	nan\\
0.000483242947843188	-0.22\\
0.000520833333333304	0.929135973038245\\
0.000541666666666663	0.918442742421097\\
0.000578554277460674	-0.22\\
nan	nan\\
0.000649909614509836	-0.22\\
0.000687499999999952	0.929135973038245\\
0.000708333333333311	0.918442742421097\\
0.000745220944127323	-0.22\\
nan	nan\\
0.000816576281176484	-0.22\\
0.000854166666666711	0.929135973038245\\
0.000874999999999959	0.918442742421094\\
0.000911887610793971	-0.22\\
nan	nan\\
0.000983242947843133	-0.22\\
0.00102083333333336	0.929135973038245\\
0.00104166666666672	0.918442742421094\\
0.00107855427746062	-0.22\\
nan	nan\\
0.00114990961450989	-0.22\\
0.00118750000000001	0.929135973038244\\
0.00120833333333337	0.918442742421094\\
0.00124522094412727	-0.22\\
nan	nan\\
0.00131657628117654	-0.22\\
0.00135416666666666	0.929135973038247\\
0.00137500000000002	0.918442742421097\\
0.00141188761079403	-0.22\\
nan	nan\\
0.00148324294784319	-0.22\\
0.0015208333333333	0.929135973038244\\
0.00154166666666666	0.918442742421094\\
0.00157855427746068	-0.22\\
nan	nan\\
0.00164990961450984	-0.22\\
0.00168749999999995	0.929135973038242\\
0.00170833333333331	0.9184427424211\\
0.00174522094412732	-0.22\\
nan	nan\\
0.00181657628117649	-0.22\\
0.00185416666666671	0.929135973038239\\
0.00187499999999996	0.918442742421098\\
0.00191188761079397	-0.22\\
nan	nan\\
0.00198324294784313	-0.22\\
0.00202083333333336	0.929135973038241\\
0.00204166666666672	0.9184427424211\\
0.00207855427746062	-0.22\\
nan	nan\\
0.00214990961450989	-0.22\\
0.00218750000000001	0.929135973038239\\
0.00220833333333337	0.918442742421098\\
0.00224522094412727	-0.22\\
nan	nan\\
0.00231657628117654	-0.22\\
0.00235416666666666	0.929135973038241\\
0.00237500000000002	0.918442742421101\\
0.00241188761079403	-0.22\\
nan	nan\\
0.00248324294784319	-0.22\\
0.00252083333333331	0.929135973038238\\
0.00254166666666666	0.918442742421104\\
0.00257855427746068	-0.22\\
nan	nan\\
0.00264990961450984	-0.22\\
0.00268749999999995	0.929135973038241\\
0.00270833333333331	0.918442742421095\\
0.00274522094412732	-0.22\\
nan	nan\\
0.00281657628117649	-0.22\\
0.00285416666666671	0.929135973038244\\
0.00287499999999996	0.918442742421098\\
0.00291188761079397	-0.22\\
nan	nan\\
0.00298324294784313	-0.22\\
0.00302083333333336	0.929135973038241\\
0.00304166666666672	0.918442742421101\\
0.00307855427746062	-0.22\\
nan	nan\\
0.00314990961450989	-0.22\\
0.00318750000000001	0.929135973038238\\
0.00320833333333337	0.918442742421098\\
0.00324522094412727	-0.22\\
nan	nan\\
0.00340431386650408	-0.22\\
0.00352083333333331	0.946001075789602\\
0.00354166666666667	0.99766655428657\\
0.00356250000000002	0.981342705830844\\
0.00360416666666663	0.753734350346538\\
0.00368749999999995	-0.192266725443062\\
nan	nan\\
0.00390431386650414	-0.22\\
0.00402083333333336	0.9460010757896\\
0.00404166666666672	0.99766655428657\\
0.00406249999999997	0.981342705830843\\
0.00410416666666669	0.753734350346539\\
0.00418750000000001	-0.192266725443069\\
nan	nan\\
0.00440431386650408	-0.22\\
0.00452083333333331	0.946001075789602\\
0.00454166666666667	0.99766655428657\\
0.00456250000000002	0.981342705830844\\
0.00460416666666663	0.753734350346543\\
0.00468749999999996	-0.192266725443069\\
nan	nan\\
0.00490431386650414	-0.22\\
0.00502083333333336	0.946001075789602\\
0.00504166666666672	0.99766655428657\\
0.00506249999999997	0.981342705830843\\
0.00510416666666669	0.753734350346534\\
0.00518750000000001	-0.192266725443061\\
nan	nan\\
0.00540431386650408	-0.22\\
0.00552083333333331	0.946001075789604\\
0.00554166666666667	0.997666554286569\\
0.00556250000000003	0.981342705830844\\
0.00560416666666663	0.753734350346539\\
0.00568749999999996	-0.192266725443068\\
nan	nan\\
0.00590431386650414	-0.22\\
0.00602083333333336	0.946001075789602\\
0.00604166666666661	0.997666554286569\\
0.00606249999999997	0.981342705830843\\
0.00610416666666669	0.753734350346534\\
0.00618750000000001	-0.192266725443061\\
nan	nan\\
0.00640431386650409	-0.22\\
0.00652083333333331	0.946001075789604\\
0.00654166666666667	0.997666554286569\\
0.00656250000000003	0.981342705830844\\
0.00660416666666663	0.753734350346539\\
0.00668749999999996	-0.192266725443068\\
nan	nan\\
0.00690431386650414	-0.22\\
0.00702083333333337	0.946001075789601\\
0.00704166666666661	0.997666554286569\\
0.00706249999999997	0.981342705830846\\
0.00710416666666669	0.753734350346535\\
0.00718750000000001	-0.19226672544306\\
nan	nan\\
0.00740431386650409	-0.22\\
0.00752083333333331	0.946001075789604\\
0.00754166666666667	0.997666554286569\\
0.00756250000000003	0.981342705830845\\
0.00760416666666663	0.75373435034654\\
0.00768749999999996	-0.192266725443067\\
nan	nan\\
0.00790431386650414	-0.22\\
0.00802083333333337	0.946001075789601\\
0.00804166666666661	0.997666554286568\\
0.00806249999999997	0.981342705830843\\
0.00810416666666669	0.753734350346544\\
0.00818750000000001	-0.19226672544306\\
nan	nan\\
0.00840431386650409	-0.22\\
0.00852083333333331	0.946001075789599\\
0.00854166666666667	0.997666554286569\\
0.00856250000000003	0.981342705830847\\
0.00860416666666663	0.75373435034654\\
0.00868749999999996	-0.192266725443053\\
nan	nan\\
0.00890431386650414	-0.22\\
0.00902083333333337	0.946001075789601\\
0.00904166666666661	0.997666554286569\\
0.00906249999999997	0.981342705830846\\
0.00910416666666669	0.753734350346545\\
0.00918750000000002	-0.192266725443059\\
nan	nan\\
0.00940431386650409	-0.22\\
0.00952083333333331	0.946001075789599\\
0.00954166666666667	0.99766655428657\\
0.00956250000000003	0.981342705830845\\
0.00960416666666664	0.75373435034654\\
0.00968749999999996	-0.192266725443052\\
nan	nan\\
0.00990431386650414	-0.22\\
0.00997916666666665	0.657179221451548\\
0.01	0.395553951938348\\
0.0100208333333334	0.929135973038248\\
0.0100416666666666	0.918442742421105\\
0.0100785542774606	-0.22\\
nan	nan\\
0.0101499096145099	-0.22\\
0.0101875	0.929135973038261\\
0.0102083333333334	0.918442742421091\\
0.0102452209441273	-0.22\\
nan	nan\\
0.0103165762811765	-0.22\\
0.0103541666666667	0.929135973038232\\
0.010375	0.918442742421122\\
0.010411887610794	-0.22\\
nan	nan\\
0.0104832429478432	-0.22\\
0.0105208333333333	0.929135973038245\\
0.0105416666666667	0.918442742421108\\
0.0105785542774607	-0.22\\
nan	nan\\
0.0106499096145098	-0.22\\
0.0106875	0.929135973038237\\
0.0107083333333333	0.918442742421094\\
0.0107452209441273	-0.22\\
nan	nan\\
0.0108165762811765	-0.22\\
0.0108541666666667	0.92913597303825\\
0.010875	0.918442742421103\\
0.010911887610794	-0.22\\
nan	nan\\
0.0109832429478431	-0.22\\
0.0110208333333334	0.929135973038221\\
0.0110416666666666	0.918442742421089\\
0.0110785542774606	-0.22\\
nan	nan\\
0.0111499096145098	-0.22\\
0.0111875	0.929135973038234\\
0.0112083333333334	0.91844274242112\\
0.0112452209441273	-0.22\\
nan	nan\\
0.0113165762811765	-0.22\\
0.0113541666666667	0.929135973038247\\
0.011375	0.918442742421106\\
0.0114118876107939	-0.22\\
nan	nan\\
0.0114832429478432	-0.22\\
0.0115208333333333	0.929135973038239\\
0.0115416666666667	0.918442742421092\\
0.0115785542774607	-0.22\\
nan	nan\\
0.0116499096145098	-0.22\\
0.0116875	0.929135973038252\\
0.0117083333333333	0.9184427424211\\
0.0117452209441273	-0.22\\
nan	nan\\
0.0118165762811765	-0.22\\
0.0118541666666667	0.929135973038223\\
0.011875	0.918442742421086\\
0.011911887610794	-0.22\\
nan	nan\\
0.0119832429478431	-0.22\\
0.0120208333333334	0.929135973038236\\
0.0120416666666666	0.918442742421117\\
0.0120785542774606	-0.22\\
nan	nan\\
0.0121499096145098	-0.22\\
0.0121875	0.929135973038249\\
0.0122083333333334	0.918442742421104\\
0.0122452209441273	-0.22\\
nan	nan\\
0.0123165762811766	-0.22\\
0.0123541666666667	0.929135973038241\\
0.012375	0.918442742421112\\
0.0124118876107939	-0.22\\
nan	nan\\
0.0124832429478432	-0.22\\
0.0125208333333333	0.929135973038254\\
0.0125416666666667	0.918442742421098\\
0.0125785542774607	-0.22\\
nan	nan\\
0.0126499096145098	-0.22\\
0.0126875	0.929135973038225\\
0.0127083333333333	0.918442742421084\\
0.0127452209441273	-0.22\\
nan	nan\\
0.0128165762811765	-0.22\\
0.0128541666666667	0.929135973038238\\
0.012875	0.918442742421115\\
0.012911887610794	-0.22\\
nan	nan\\
0.0129832429478431	-0.22\\
0.0130208333333334	0.92913597303823\\
0.0130416666666666	0.918442742421101\\
0.0130785542774606	-0.22\\
nan	nan\\
0.0131499096145098	-0.22\\
0.0131875	0.929135973038243\\
0.0132083333333334	0.91844274242111\\
0.0132452209441273	-0.22\\
nan	nan\\
0.0133165762811766	-0.22\\
0.0133541666666667	0.929135973038256\\
0.013375	0.918442742421096\\
0.0134118876107939	-0.22\\
nan	nan\\
0.0134832429478432	-0.22\\
0.0135208333333333	0.929135973038206\\
0.0135416666666667	0.918442742421149\\
0.0135785542774607	-0.22\\
nan	nan\\
0.0136499096145098	-0.22\\
0.0136875	0.92913597303824\\
0.0137083333333333	0.918442742421113\\
0.0137452209441273	-0.22\\
nan	nan\\
0.0138165762811765	-0.22\\
0.0138541666666666	0.92913597303819\\
0.013875	0.918442742421121\\
0.013911887610794	-0.22\\
nan	nan\\
0.0139832429478431	-0.22\\
0.0140208333333334	0.929135973038224\\
0.0140416666666666	0.91844274242113\\
0.0140785542774606	-0.22\\
nan	nan\\
0.0141499096145098	-0.22\\
0.0141875	0.929135973038216\\
0.0142083333333334	0.918442742421138\\
0.0142452209441273	-0.22\\
nan	nan\\
0.0143165762811766	-0.22\\
0.0143541666666667	0.929135973038208\\
0.014375	0.918442742421147\\
0.0144118876107939	-0.22\\
nan	nan\\
0.0144832429478432	-0.22\\
0.0145208333333333	0.929135973038242\\
0.0145416666666667	0.91844274242111\\
0.0145785542774607	-0.22\\
nan	nan\\
0.0146499096145098	-0.22\\
0.0146875	0.929135973038193\\
0.0147083333333333	0.918442742421119\\
0.0147452209441273	-0.22\\
nan	nan\\
0.0148165762811765	-0.22\\
0.0148541666666666	0.929135973038227\\
0.014875	0.918442742421128\\
0.014911887610794	-0.22\\
nan	nan\\
0.0149832429478431	-0.22\\
0.0150208333333334	0.929135973038219\\
0.0150416666666666	0.918442742421136\\
0.0150785542774606	-0.22\\
nan	nan\\
0.0151499096145098	-0.22\\
0.0151875	0.929135973038211\\
0.0152083333333334	0.918442742421145\\
0.0152452209441273	-0.22\\
nan	nan\\
0.0153165762811766	-0.22\\
0.0153541666666667	0.929135973038245\\
0.015375	0.918442742421108\\
0.0154118876107939	-0.22\\
nan	nan\\
0.0154832429478432	-0.22\\
0.0155208333333333	0.929135973038195\\
0.0155416666666667	0.918442742421162\\
0.0155785542774607	-0.22\\
nan	nan\\
0.0156499096145098	-0.22\\
0.0156875	0.929135973038229\\
0.0157083333333333	0.91844274242117\\
0.0157452209441273	-0.22\\
nan	nan\\
0.0158165762811765	-0.22\\
0.0158541666666666	0.929135973038221\\
0.015875	0.918442742421134\\
0.015911887610794	-0.22\\
nan	nan\\
0.0159832429478431	-0.22\\
0.0160208333333334	0.929135973038213\\
0.0160416666666666	0.918442742421142\\
0.0160785542774606	-0.22\\
nan	nan\\
0.0161499096145098	-0.22\\
0.0161875	0.929135973038247\\
0.0162083333333334	0.918442742421151\\
0.0162452209441273	-0.22\\
nan	nan\\
0.0163165762811766	-0.22\\
0.0163541666666667	0.929135973038197\\
0.016375	0.918442742421114\\
0.0164118876107939	-0.22\\
nan	nan\\
0.0164832429478432	-0.22\\
0.0165208333333333	0.929135973038189\\
0.0165416666666667	0.918442742421123\\
0.0165785542774607	-0.22\\
nan	nan\\
0.01665295537323	-0.22\\
0.0166666666666667	0.0682747863999734\\
0.0166875	-0.192266725443077\\
nan	nan\\
0.016904313866504	-0.22\\
0.0170208333333334	0.946001075789595\\
0.0170416666666666	0.997666554286569\\
0.0170625	0.98134270583085\\
0.0171041666666667	0.753734350346529\\
0.0171875	-0.192266725443028\\
nan	nan\\
0.0174043138665041	-0.22\\
0.0175208333333333	0.946001075789597\\
0.0175416666666667	0.99766655428657\\
0.0175625	0.981342705830848\\
0.0176041666666666	0.753734350346524\\
0.0176875	-0.192266725443034\\
nan	nan\\
0.017904313866504	-0.22\\
0.0180208333333334	0.9460010757896\\
0.0180416666666666	0.99766655428657\\
0.0180625	0.981342705830847\\
0.0181041666666667	0.753734350346538\\
0.0181875	-0.192266725443041\\
nan	nan\\
0.0184043138665041	-0.22\\
0.0185208333333333	0.946001075789602\\
0.0185416666666667	0.997666554286569\\
0.0185625	0.981342705830846\\
0.0186041666666666	0.753734350346553\\
0.0186875	-0.192266725443048\\
nan	nan\\
0.018904313866504	-0.22\\
0.0190208333333334	0.946001075789604\\
0.0190416666666666	0.997666554286569\\
0.0190625	0.981342705830844\\
0.0191041666666667	0.753734350346548\\
0.0191875	-0.192266725443055\\
nan	nan\\
0.0194043138665041	-0.22\\
0.0195208333333333	0.946001075789588\\
0.0195416666666667	0.99766655428657\\
0.0195625	0.981342705830843\\
0.0196041666666666	0.753734350346543\\
0.0196875	-0.192266725443061\\
nan	nan\\
0.019904313866504	-0.22\\
0.0200208333333334	0.946001075789581\\
0.0200416666666666	0.99766655428657\\
0.0200625	0.981342705830842\\
0.0201041666666667	0.753734350346558\\
0.0201875	-0.192266725443068\\
nan	nan\\
0.0204043138665041	-0.22\\
0.0205208333333333	0.946001075789592\\
0.0205416666666667	0.997666554286569\\
0.0205625	0.981342705830846\\
0.0206041666666666	0.753734350346572\\
0.0206875	-0.192266725443047\\
nan	nan\\
0.020904313866504	-0.22\\
0.0210208333333334	0.946001075789585\\
0.0210416666666666	0.997666554286567\\
0.0210625	0.98134270583085\\
0.0211041666666667	0.753734350346549\\
0.0211875	-0.192266725443026\\
nan	nan\\
0.0214043138665041	-0.22\\
0.0215208333333333	0.946001075789597\\
0.0215416666666667	0.997666554286569\\
0.0215625	0.981342705830854\\
0.0216041666666666	0.753734350346563\\
0.0216875	-0.192266725443005\\
nan	nan\\
0.021904313866504	-0.22\\
0.0220208333333334	0.94600107578959\\
0.0220416666666666	0.997666554286564\\
0.0220625	0.981342705830858\\
0.0221041666666667	0.753734350346577\\
0.0221875	-0.192266725443039\\
nan	nan\\
0.0224043138665041	-0.22\\
0.0225208333333333	0.946001075789583\\
0.0225416666666667	0.997666554286566\\
0.0225625	0.981342705830851\\
0.0226041666666666	0.753734350346591\\
0.0226875	-0.192266725443018\\
nan	nan\\
0.022904313866504	-0.22\\
0.0230208333333334	0.946001075789594\\
0.0230416666666666	0.997666554286565\\
0.0230625	0.981342705830856\\
0.0231041666666667	0.753734350346568\\
0.0231875	-0.192266725443053\\
nan	nan\\
0.0234043138665041	-0.22\\
0.0235208333333333	0.946001075789587\\
0.0235416666666667	0.997666554286567\\
0.0235625	0.981342705830849\\
0.0236041666666666	0.753734350346582\\
0.0236875	-0.192266725443031\\
nan	nan\\
0.023904313866504	-0.22\\
0.0240208333333334	0.94600107578958\\
0.0240416666666666	0.997666554286566\\
0.0240625	0.981342705830853\\
0.0241041666666667	0.753734350346559\\
0.0241875	-0.19226672544301\\
nan	nan\\
0.0244043138665041	-0.22\\
0.0245208333333333	0.946001075789592\\
0.0245416666666667	0.997666554286568\\
0.0245625	0.981342705830857\\
0.0246041666666666	0.753734350346573\\
0.0246875	-0.192266725442989\\
nan	nan\\
0.024904313866504	-0.22\\
0.0250208333333334	0.946001075789585\\
0.0250416666666666	0.997666554286567\\
0.0250625	0.981342705830861\\
0.0251041666666667	0.75373435034655\\
0.0251875	-0.192266725443024\\
nan	nan\\
0.0254043138665041	-0.22\\
0.0255208333333333	0.946001075789596\\
0.0255416666666667	0.997666554286569\\
0.0255625	0.981342705830854\\
0.0256041666666667	0.753734350346564\\
0.0256875	-0.192266725443003\\
nan	nan\\
0.025904313866504	-0.22\\
0.0260208333333334	0.946001075789589\\
0.0260416666666666	0.997666554286564\\
0.0260625	0.981342705830859\\
0.0261041666666667	0.753734350346578\\
0.0261875	-0.192266725443037\\
nan	nan\\
0.0264043138665041	-0.22\\
0.0265208333333333	0.946001075789582\\
0.0265416666666667	0.997666554286566\\
0.0265625	0.981342705830852\\
0.0266041666666667	0.753734350346592\\
0.0266458333333334	0.324163484379289\\
0.0266666666666666	0.395553951938143\\
0.0266875	0.929135973038207\\
0.0267083333333333	0.918442742421104\\
0.0267452209441273	-0.22\\
nan	nan\\
0.0268165762811765	-0.22\\
0.0268541666666666	0.929135973038199\\
0.026875	0.918442742421157\\
0.026911887610794	-0.22\\
nan	nan\\
0.0269832429478432	-0.22\\
0.0270208333333334	0.929135973038191\\
0.0270416666666666	0.918442742421121\\
0.0270785542774606	-0.22\\
nan	nan\\
0.0271499096145098	-0.22\\
0.0271875	0.929135973038267\\
0.0272083333333333	0.918442742421129\\
0.0272452209441273	-0.22\\
nan	nan\\
0.0273165762811766	-0.22\\
0.0273541666666667	0.929135973038259\\
0.027375	0.918442742421048\\
0.0274118876107939	-0.22\\
nan	nan\\
0.0274832429478432	-0.22\\
0.0275208333333333	0.929135973038251\\
0.0275416666666667	0.918442742421056\\
0.0275785542774607	-0.22\\
nan	nan\\
0.0276499096145099	-0.22\\
0.0276875	0.929135973038243\\
0.0277083333333333	0.918442742421155\\
0.0277452209441273	-0.22\\
nan	nan\\
0.0278165762811765	-0.22\\
0.0278541666666666	0.929135973038235\\
0.027875	0.918442742421073\\
0.027911887610794	-0.22\\
nan	nan\\
0.0279832429478432	-0.22\\
0.0280208333333334	0.929135973038227\\
0.0280416666666666	0.918442742421082\\
0.0280785542774606	-0.22\\
nan	nan\\
0.0281499096145098	-0.22\\
0.0281875	0.929135973038303\\
0.0282083333333333	0.91844274242109\\
0.0282452209441273	-0.22\\
nan	nan\\
0.0283165762811766	-0.22\\
0.0283541666666667	0.929135973038211\\
0.028375	0.918442742421099\\
0.0284118876107939	-0.22\\
nan	nan\\
0.0284832429478432	-0.22\\
0.0285208333333333	0.929135973038203\\
0.0285416666666667	0.918442742421107\\
0.0285785542774607	-0.22\\
nan	nan\\
0.0286499096145099	-0.22\\
0.0286875	0.929135973038195\\
0.0287083333333333	0.918442742421116\\
0.0287452209441273	-0.22\\
nan	nan\\
0.0288165762811765	-0.22\\
0.0288541666666666	0.929135973038271\\
0.028875	0.918442742421124\\
0.028911887610794	-0.22\\
nan	nan\\
0.0289832429478432	-0.22\\
0.0290208333333334	0.929135973038264\\
0.0290416666666666	0.918442742421043\\
0.0290785542774606	-0.22\\
nan	nan\\
0.0291499096145098	-0.22\\
0.0291875	0.929135973038256\\
0.0292083333333333	0.918442742421142\\
0.0292452209441273	-0.22\\
nan	nan\\
0.0293165762811765	-0.22\\
0.0293541666666667	0.929135973038248\\
0.029375	0.91844274242115\\
0.0294118876107939	-0.22\\
nan	nan\\
0.0294832429478432	-0.22\\
0.0295208333333333	0.92913597303824\\
0.0295416666666667	0.918442742421159\\
0.0295785542774607	-0.22\\
nan	nan\\
0.0296499096145099	-0.22\\
0.0296875	0.929135973038232\\
0.0297083333333333	0.918442742421077\\
0.0297452209441273	-0.22\\
nan	nan\\
0.0298165762811765	-0.22\\
0.0298541666666666	0.929135973038224\\
0.029875	0.918442742421086\\
0.029911887610794	-0.22\\
nan	nan\\
0.0299817671434905	-0.22\\
0.0300208333333334	0.946001075789589\\
0.0300416666666666	0.997666554286564\\
0.0300625	0.981342705830859\\
0.0301041666666667	0.75373435034658\\
0.0301875	-0.192266725443035\\
nan	nan\\
0.0304043138665041	-0.22\\
0.0305208333333333	0.946001075789582\\
0.0305416666666667	0.997666554286566\\
0.0305625	0.981342705830852\\
0.0306041666666667	0.753734350346594\\
0.0306875	-0.192266725443014\\
nan	nan\\
0.0309043138665041	-0.22\\
0.0310208333333334	0.946001075789575\\
0.0310416666666666	0.997666554286565\\
0.0310625	0.981342705830856\\
0.0311041666666667	0.753734350346571\\
0.0311875	-0.192266725442993\\
nan	nan\\
0.0314043138665041	-0.22\\
0.0315208333333333	0.946001075789586\\
0.0315416666666667	0.997666554286563\\
0.0315625	0.98134270583085\\
0.0316041666666667	0.753734350346622\\
0.0316875	-0.192266725443028\\
nan	nan\\
0.0319043138665041	-0.22\\
0.0320208333333334	0.946001075789579\\
0.0320416666666666	0.997666554286562\\
0.0320625	0.981342705830854\\
0.0321041666666667	0.753734350346599\\
0.0321875	-0.192266725443006\\
nan	nan\\
0.0324043138665041	-0.22\\
0.0325208333333333	0.94600107578959\\
0.0325416666666667	0.997666554286564\\
0.0325625	0.981342705830847\\
0.0326041666666667	0.753734350346613\\
0.0326875	-0.192266725443041\\
nan	nan\\
0.0329043138665041	-0.22\\
0.0330208333333334	0.946001075789583\\
0.0330416666666666	0.997666554286563\\
0.0330625	0.981342705830851\\
0.0331041666666667	0.75373435034659\\
0.0331875	-0.19226672544302\\
nan	nan\\
0.0334043138665041	-0.22\\
0.0335208333333333	0.946001075789576\\
0.0335416666666667	0.997666554286565\\
0.0335625000000001	0.981342705830844\\
0.0336041666666667	0.753734350346604\\
0.0336875	-0.192266725443055\\
nan	nan\\
0.0339043138665041	-0.22\\
0.0340208333333334	0.946001075789588\\
0.0340416666666666	0.997666554286564\\
0.0340625	0.981342705830859\\
0.0341041666666667	0.753734350346581\\
0.0341875	-0.192266725442978\\
nan	nan\\
0.0344043138665041	-0.22\\
0.0345208333333333	0.946001075789581\\
0.0345416666666667	0.997666554286566\\
0.0345625000000001	0.981342705830863\\
0.0346041666666667	0.753734350346595\\
0.0346875	-0.192266725442957\\
nan	nan\\
0.0349043138665041	-0.22\\
0.0350208333333334	0.946001075789592\\
0.0350416666666666	0.997666554286565\\
0.0350625	0.981342705830857\\
0.0351041666666667	0.753734350346572\\
0.0351875	-0.192266725442991\\
nan	nan\\
0.0354043138665041	-0.22\\
0.0355208333333333	0.946001075789585\\
0.0355416666666667	0.997666554286567\\
0.0355625000000001	0.981342705830861\\
0.0356041666666667	0.753734350346586\\
0.0356875	-0.19226672544297\\
nan	nan\\
0.0359043138665041	-0.22\\
0.0360208333333334	0.946001075789597\\
0.0360416666666666	0.997666554286566\\
0.0360625	0.981342705830865\\
0.0361041666666667	0.753734350346563\\
0.0361875	-0.192266725443005\\
nan	nan\\
0.0364043138665041	-0.22\\
0.0365208333333333	0.946001075789571\\
0.0365416666666667	0.997666554286568\\
0.0365625000000001	0.981342705830858\\
0.0366041666666667	0.753734350346577\\
0.0366458333333334	0.324163484379374\\
0.0366666666666666	0.395553951938416\\
0.0366875	0.929135973038233\\
0.0367083333333333	0.918442742421166\\
0.0367452209441274	-0.22\\
nan	nan\\
0.0368165762811765	-0.22\\
0.0368541666666666	0.929135973038225\\
0.036875	0.918442742421084\\
0.036911887610794	-0.22\\
nan	nan\\
0.0369832429478432	-0.22\\
0.0370208333333333	0.929135973038217\\
0.0370416666666666	0.918442742421093\\
0.0370785542774607	-0.22\\
nan	nan\\
0.0371499096145098	-0.22\\
0.0371875	0.929135973038209\\
0.0372083333333333	0.918442742421191\\
0.0372452209441273	-0.22\\
nan	nan\\
0.0373165762811765	-0.22\\
0.0373541666666667	0.929135973038285\\
0.037375	0.91844274242111\\
0.0374118876107939	-0.22\\
nan	nan\\
0.0374832429478432	-0.22\\
0.0375208333333333	0.929135973038193\\
0.0375416666666667	0.918442742421118\\
0.0375785542774606	-0.22\\
nan	nan\\
0.0376499096145099	-0.22\\
0.0376875	0.929135973038269\\
0.0377083333333333	0.918442742421127\\
0.0377452209441274	-0.22\\
nan	nan\\
0.0378165762811765	-0.22\\
0.0378541666666666	0.929135973038261\\
0.037875	0.918442742421135\\
0.037911887610794	-0.22\\
nan	nan\\
0.0379832429478432	-0.22\\
0.0380208333333333	0.929135973038169\\
0.0380416666666666	0.918442742421054\\
0.0380785542774607	-0.22\\
nan	nan\\
0.0381499096145098	-0.22\\
0.0381875	0.929135973038245\\
0.0382083333333333	0.918442742421152\\
0.0382452209441273	-0.22\\
nan	nan\\
0.0383165762811765	-0.22\\
0.0383541666666667	0.929135973038237\\
0.038375	0.918442742421071\\
0.0384118876107939	-0.22\\
nan	nan\\
0.0384832429478432	-0.22\\
0.0385208333333333	0.929135973038229\\
0.0385416666666667	0.918442742421079\\
0.0385785542774606	-0.22\\
nan	nan\\
0.0386499096145099	-0.22\\
0.0386875	0.929135973038221\\
0.0387083333333333	0.918442742421178\\
0.0387452209441274	-0.22\\
nan	nan\\
0.0388165762811765	-0.22\\
0.0388541666666666	0.929135973038298\\
0.038875	0.918442742421097\\
0.038911887610794	-0.22\\
nan	nan\\
0.0389832429478432	-0.22\\
0.0390208333333333	0.929135973038205\\
0.0390416666666666	0.918442742421195\\
0.0390785542774607	-0.22\\
nan	nan\\
0.0391499096145098	-0.22\\
0.0391875	0.929135973038282\\
0.0392083333333333	0.918442742421114\\
0.0392452209441273	-0.22\\
nan	nan\\
0.0393165762811765	-0.22\\
0.0393541666666667	0.92913597303819\\
0.039375	0.918442742421122\\
0.0394118876107939	-0.22\\
nan	nan\\
0.0394832429478432	-0.22\\
0.0395208333333333	0.929135973038182\\
0.0395416666666667	0.918442742421131\\
0.0395785542774606	-0.22\\
nan	nan\\
0.0396499096145099	-0.22\\
0.0396875	0.929135973038258\\
0.0397083333333333	0.918442742421139\\
0.0397452209441274	-0.22\\
nan	nan\\
0.0398165762811765	-0.22\\
0.0398541666666666	0.929135973038166\\
0.039875	0.918442742421058\\
0.039911887610794	-0.22\\
nan	nan\\
0.0399817671434906	-0.22\\
0.0400208333333333	0.946001075789578\\
0.0400416666666666	0.997666554286562\\
0.0400625	0.981342705830854\\
0.0401041666666667	0.753734350346601\\
0.0401875	-0.192266725443003\\
nan	nan\\
0.0404043138665041	-0.22\\
0.0405208333333333	0.946001075789571\\
0.0405416666666667	0.997666554286568\\
0.0405625000000001	0.981342705830837\\
0.0406041666666667	0.753734350346578\\
0.0406875	-0.192266725443037\\
nan	nan\\
0.0409043138665041	-0.22\\
0.0410208333333333	0.946001075789564\\
0.0410416666666666	0.997666554286566\\
0.0410625	0.981342705830841\\
0.0411041666666667	0.753734350346592\\
0.0411875	-0.192266725443016\\
nan	nan\\
0.0414043138665041	-0.22\\
0.0415208333333333	0.946001075789594\\
0.0415416666666667	0.997666554286565\\
0.0415624999999999	0.981342705830867\\
0.0416041666666667	0.753734350346607\\
0.0416875	-0.192266725442995\\
nan	nan\\
0.0419043138665041	-0.22\\
0.0420208333333333	0.946001075789587\\
0.0420416666666666	0.997666554286563\\
0.0420625	0.981342705830871\\
0.0421041666666667	0.753734350346621\\
0.0421875	-0.192266725442974\\
nan	nan\\
0.0424043138665041	-0.22\\
0.0425208333333333	0.94600107578958\\
0.0425416666666667	0.997666554286562\\
0.0425624999999999	0.981342705830875\\
0.0426041666666667	0.75373435034656\\
0.0426875	-0.192266725442953\\
nan	nan\\
0.0429043138665041	-0.22\\
0.0430208333333333	0.946001075789573\\
0.0430416666666666	0.997666554286568\\
0.0430625	0.981342705830857\\
0.0431041666666667	0.753734350346574\\
0.0431875	-0.192266725442932\\
nan	nan\\
0.0433240196055433	-0.22\\
0.0433541666666667	0.92913597303825\\
0.0433750000000001	0.918442742421147\\
0.043411887610794	-0.22\\
nan	nan\\
0.0434832429478432	-0.22\\
0.0435208333333333	0.929135973038242\\
0.0435416666666667	0.918442742421065\\
0.0435785542774606	-0.22\\
nan	nan\\
0.0436499096145099	-0.22\\
0.0436875	0.929135973038235\\
0.0437083333333333	0.918442742421074\\
0.0437452209441274	-0.22\\
nan	nan\\
0.0438165762811765	-0.22\\
0.0438541666666666	0.929135973038227\\
0.043875	0.918442742421173\\
0.043911887610794	-0.22\\
nan	nan\\
0.0439832429478432	-0.22\\
0.0440208333333333	0.929135973038303\\
0.0440416666666666	0.918442742421091\\
0.0440785542774607	-0.22\\
nan	nan\\
0.0441499096145098	-0.22\\
0.0441875	0.929135973038211\\
0.0442083333333333	0.9184427424211\\
0.0442452209441273	-0.22\\
nan	nan\\
0.0443165762811765	-0.22\\
0.0443541666666667	0.929135973038203\\
0.0443750000000001	0.918442742421108\\
0.044411887610794	-0.22\\
nan	nan\\
0.0444832429478432	-0.22\\
0.0445208333333333	0.929135973038195\\
0.0445416666666667	0.918442742421117\\
0.0445785542774606	-0.22\\
nan	nan\\
0.0446499096145099	-0.22\\
0.0446875	0.929135973038187\\
0.0447083333333333	0.918442742421125\\
0.0447452209441274	-0.22\\
nan	nan\\
0.0448165762811765	-0.22\\
0.0448541666666666	0.929135973038263\\
0.044875	0.918442742421134\\
0.044911887610794	-0.22\\
nan	nan\\
0.0449832429478432	-0.22\\
0.0450208333333333	0.929135973038171\\
0.0450416666666666	0.918442742421052\\
0.0450785542774607	-0.22\\
nan	nan\\
0.0451499096145098	-0.22\\
0.0451875	0.929135973038247\\
0.0452083333333333	0.918442742421151\\
0.0452452209441273	-0.22\\
nan	nan\\
0.0453165762811765	-0.22\\
0.0453541666666667	0.929135973038239\\
0.0453750000000001	0.918442742421159\\
0.045411887610794	-0.22\\
nan	nan\\
0.0454832429478432	-0.22\\
0.0455208333333333	0.929135973038231\\
0.0455416666666667	0.918442742421078\\
0.0455785542774606	-0.22\\
nan	nan\\
0.0456499096145099	-0.22\\
0.0456875	0.929135973038223\\
0.0457083333333334	0.918442742421176\\
0.0457452209441274	-0.22\\
nan	nan\\
0.0458165762811765	-0.22\\
0.0458541666666666	0.929135973038299\\
0.045875	0.918442742421095\\
0.045911887610794	-0.22\\
nan	nan\\
0.0459832429478432	-0.22\\
0.0460208333333333	0.929135973038207\\
0.0460416666666666	0.918442742421104\\
0.0460785542774607	-0.22\\
nan	nan\\
0.0461499096145098	-0.22\\
0.0461875	0.929135973038199\\
0.0462083333333333	0.918442742421112\\
0.0462452209441273	-0.22\\
nan	nan\\
0.0463165762811765	-0.22\\
0.0463541666666667	0.929135973038275\\
0.0463750000000001	0.918442742421121\\
0.046411887610794	-0.22\\
nan	nan\\
0.0464832429478432	-0.22\\
0.0465208333333333	0.929135973038183\\
0.0465416666666667	0.918442742421039\\
0.0465785542774606	-0.22\\
nan	nan\\
0.0466499096145099	-0.22\\
0.0466875	0.929135973038259\\
0.0467083333333334	0.918442742421138\\
0.0467452209441274	-0.22\\
nan	nan\\
0.0468165762811765	-0.22\\
0.0468541666666666	0.929135973038251\\
0.046875	0.918442742421146\\
0.046911887610794	-0.22\\
nan	nan\\
0.0469832429478432	-0.22\\
0.0470208333333333	0.929135973038243\\
0.0470416666666666	0.918442742421155\\
0.0470785542774607	-0.22\\
nan	nan\\
0.0471499096145098	-0.22\\
0.0471875	0.929135973038235\\
0.0472083333333333	0.918442742421163\\
0.0472452209441273	-0.22\\
nan	nan\\
0.0473165762811765	-0.22\\
0.0473541666666667	0.929135973038143\\
0.0473750000000001	0.918442742421082\\
0.047411887610794	-0.22\\
nan	nan\\
0.0474832429478432	-0.22\\
0.0475208333333333	0.929135973038219\\
0.0475416666666667	0.91844274242118\\
0.0475785542774606	-0.22\\
nan	nan\\
0.0476499096145099	-0.22\\
0.0476875	0.929135973038211\\
0.0477083333333334	0.918442742421099\\
0.0477452209441274	-0.22\\
nan	nan\\
0.0478165762811765	-0.22\\
0.0478541666666666	0.929135973038203\\
0.047875	0.918442742421107\\
0.047911887610794	-0.22\\
nan	nan\\
0.0479832429478432	-0.22\\
0.0480208333333333	0.929135973038195\\
0.0480416666666666	0.918442742421116\\
0.0480785542774607	-0.22\\
nan	nan\\
0.0481499096145098	-0.22\\
0.0481875	0.929135973038271\\
0.0482083333333333	0.918442742421124\\
0.0482452209441273	-0.22\\
nan	nan\\
0.0483165762811765	-0.22\\
0.0483541666666667	0.929135973038179\\
0.0483750000000001	0.918442742421043\\
0.048411887610794	-0.22\\
nan	nan\\
0.0484832429478432	-0.22\\
0.0485208333333333	0.929135973038256\\
0.0485416666666667	0.918442742421142\\
0.0485785542774606	-0.22\\
nan	nan\\
0.0486499096145099	-0.22\\
0.0486875	0.929135973038248\\
0.0487083333333334	0.91844274242115\\
0.0487452209441274	-0.22\\
nan	nan\\
0.0488165762811765	-0.22\\
0.0488541666666666	0.92913597303824\\
0.048875	0.918442742421069\\
0.048911887610794	-0.22\\
nan	nan\\
0.0489832429478432	-0.22\\
0.0490208333333333	0.929135973038232\\
0.0490416666666667	0.918442742421167\\
0.0490785542774607	-0.22\\
nan	nan\\
0.0491499096145098	-0.22\\
0.0491875000000001	0.929135973038308\\
0.0492083333333333	0.918442742421086\\
0.0492452209441273	-0.22\\
nan	nan\\
0.0493165762811765	-0.22\\
0.0493541666666667	0.929135973038216\\
0.0493750000000001	0.918442742421094\\
0.049411887610794	-0.22\\
nan	nan\\
0.0494832429478432	-0.22\\
0.0495208333333333	0.929135973038208\\
0.0495416666666667	0.918442742421103\\
0.0495785542774606	-0.22\\
nan	nan\\
0.0496499096145099	-0.22\\
0.0496875	0.9291359730382\\
0.0497083333333334	0.918442742421111\\
0.0497452209441274	-0.22\\
nan	nan\\
0.0498165762811765	-0.22\\
0.0498541666666666	0.929135973038192\\
0.049875	0.91844274242121\\
0.049911887610794	-0.22\\
nan	nan\\
0.0499817671434906	-0.22\\
0.0500208333333333	0.946001075789585\\
0.0500416666666667	0.997666554286563\\
0.0500625	0.981342705830872\\
0.0501041666666666	0.753734350346623\\
0.0501875000000001	-0.19226672544297\\
nan	nan\\
0.0504043138665041	-0.22\\
0.0505208333333333	0.946001075789578\\
0.0505416666666667	0.997666554286562\\
0.0505625	0.981342705830876\\
0.0506041666666667	0.753734350346563\\
0.0506875	-0.192266725442949\\
nan	nan\\
0.0509043138665041	-0.22\\
0.0510208333333333	0.946001075789571\\
0.0510416666666667	0.997666554286568\\
0.0510625	0.981342705830858\\
0.0511041666666666	0.753734350346577\\
0.0511875000000001	-0.192266725442928\\
nan	nan\\
0.0514043138665041	-0.22\\
0.0515208333333333	0.946001075789601\\
0.0515416666666667	0.997666554286566\\
0.0515625	0.981342705830862\\
0.0516041666666667	0.753734350346591\\
0.0516875	-0.192266725442906\\
nan	nan\\
0.0519043138665041	-0.22\\
0.0520208333333333	0.946001075789557\\
0.0520416666666667	0.997666554286565\\
0.0520625	0.981342705830867\\
0.0521041666666666	0.753734350346605\\
0.0521875000000001	-0.192266725442997\\
nan	nan\\
0.0524043138665041	-0.22\\
0.0525208333333333	0.946001075789587\\
0.0525416666666667	0.997666554286564\\
0.0525625	0.981342705830871\\
0.0526041666666667	0.75373435034662\\
0.0526875	-0.192266725442976\\
nan	nan\\
0.0529043138665041	-0.22\\
0.0530208333333333	0.94600107578958\\
0.0530416666666667	0.997666554286562\\
0.0530625	0.981342705830853\\
0.0531041666666666	0.753734350346634\\
0.0531875000000001	-0.192266725442955\\
nan	nan\\
0.0534043138665041	-0.22\\
0.0535208333333334	0.94600107578961\\
0.0535416666666667	0.997666554286568\\
0.0535625	0.981342705830857\\
0.0536041666666667	0.753734350346648\\
0.0536875	-0.192266725442933\\
nan	nan\\
0.0539043138665041	-0.22\\
0.0540208333333333	0.946001075789566\\
0.0540416666666667	0.997666554286567\\
0.0540625	0.981342705830861\\
0.0541041666666666	0.753734350346662\\
0.0541875000000001	-0.192266725443024\\
nan	nan\\
0.0544043138665041	-0.22\\
0.0545208333333334	0.946001075789559\\
0.0545416666666667	0.997666554286565\\
0.0545625	0.981342705830865\\
0.0546041666666667	0.753734350346601\\
0.0546875	-0.192266725443003\\
nan	nan\\
0.0549043138665041	-0.22\\
0.0550208333333333	0.946001075789552\\
0.0550416666666667	0.997666554286564\\
0.0550625	0.981342705830848\\
0.0551041666666666	0.753734350346616\\
0.0551875000000001	-0.192266725442982\\
nan	nan\\
0.0554043138665041	-0.22\\
0.0555208333333334	0.946001075789582\\
0.0555416666666667	0.99766655428657\\
0.0555625	0.981342705830852\\
0.0556041666666667	0.75373435034663\\
0.0556875	-0.19226672544296\\
nan	nan\\
0.0559043138665041	-0.22\\
0.0560208333333333	0.946001075789575\\
0.0560416666666667	0.997666554286569\\
0.0560625	0.981342705830856\\
0.0561041666666666	0.753734350346644\\
0.0561875000000001	-0.192266725443051\\
nan	nan\\
0.0564043138665041	-0.22\\
0.0565208333333334	0.946001075789568\\
0.0565416666666667	0.99766655428656\\
0.0565625	0.98134270583086\\
0.0566041666666667	0.753734350346583\\
0.0566875	-0.19226672544303\\
nan	nan\\
0.0569043138665041	-0.22\\
0.0570208333333333	0.946001075789561\\
0.0570416666666667	0.997666554286566\\
0.0570625	0.981342705830864\\
0.0571041666666666	0.753734350346598\\
0.0571875000000001	-0.192266725443008\\
nan	nan\\
0.0574043138665041	-0.22\\
0.0575208333333334	0.946001075789591\\
0.0575416666666667	0.997666554286564\\
0.0575625	0.981342705830847\\
0.0576041666666667	0.753734350346612\\
0.0576875	-0.192266725442987\\
nan	nan\\
0.0579043138665041	-0.22\\
0.0580208333333333	0.946001075789547\\
0.0580416666666667	0.997666554286563\\
0.0580625	0.981342705830851\\
0.0581041666666666	0.753734350346626\\
0.0581875000000001	-0.192266725443078\\
nan	nan\\
0.0584043138665041	-0.22\\
0.0585208333333334	0.94600107578954\\
0.0585416666666667	0.997666554286561\\
0.0585625	0.981342705830855\\
0.0586041666666667	0.753734350346565\\
0.0586875	-0.192266725443057\\
nan	nan\\
0.0589043138665041	-0.22\\
0.0590208333333333	0.94600107578957\\
0.0590416666666667	0.99766655428656\\
0.0590625	0.981342705830859\\
0.0591041666666666	0.75373435034658\\
0.0591874999999999	-0.192266725442924\\
nan	nan\\
0.0594043138665041	-0.22\\
0.0595208333333334	0.946001075789563\\
0.0595416666666667	0.997666554286566\\
0.0595625	0.981342705830863\\
0.0596041666666667	0.753734350346594\\
0.0596875	-0.192266725442903\\
nan	nan\\
0.0599043138665041	-0.22\\
0.0600208333333333	0.946001075789556\\
0.0600416666666667	0.997666554286565\\
0.0600625	0.981342705830867\\
0.0601041666666666	0.753734350346608\\
0.0601874999999999	-0.192266725442993\\
nan	nan\\
0.0604043138665041	-0.22\\
0.0605208333333334	0.946001075789549\\
0.0605416666666667	0.997666554286556\\
0.0605625	0.981342705830871\\
0.0606041666666667	0.753734350346622\\
0.0606875	-0.192266725442972\\
nan	nan\\
0.0609043138665041	-0.22\\
0.0610208333333333	0.946001075789542\\
0.0610416666666667	0.997666554286562\\
0.0610625	0.981342705830854\\
0.0611041666666666	0.753734350346562\\
0.0611875	-0.192266725442951\\
nan	nan\\
0.0614043138665041	-0.22\\
0.0615208333333334	0.946001075789572\\
0.0615416666666667	0.99766655428656\\
0.0615625	0.981342705830858\\
0.0616041666666667	0.753734350346576\\
0.0616875	-0.19226672544293\\
nan	nan\\
0.0619043138665041	-0.22\\
0.0620208333333333	0.946001075789565\\
0.0620416666666667	0.997666554286559\\
0.0620625	0.981342705830862\\
0.0621041666666666	0.75373435034659\\
0.0621875	-0.19226672544302\\
nan	nan\\
0.0624043138665041	-0.22\\
0.0625208333333334	0.946001075789558\\
0.0625416666666667	0.997666554286557\\
0.0625625	0.981342705830866\\
0.0626041666666667	0.753734350346604\\
0.0626875	-0.192266725442887\\
nan	nan\\
0.0629043138665041	-0.22\\
0.0630208333333333	0.946001075789551\\
0.0630416666666667	0.997666554286564\\
0.0630625	0.98134270583087\\
0.0631041666666666	0.753734350346543\\
0.0631875	-0.192266725442866\\
nan	nan\\
0.0633240196055433	-0.22\\
0.0633541666666667	0.929135973038303\\
0.063375	0.918442742421181\\
0.063411887610794	-0.22\\
nan	nan\\
0.0634832429478431	-0.22\\
0.0635208333333334	0.929135973038211\\
0.0635416666666667	0.9184427424211\\
0.0635785542774606	-0.22\\
nan	nan\\
0.0636499096145099	-0.22\\
0.0636875	0.929135973038119\\
0.0637083333333334	0.918442742421018\\
0.0637452209441273	-0.22\\
nan	nan\\
0.0638165762811765	-0.22\\
0.0638541666666667	0.929135973038195\\
0.063875	0.918442742421117\\
0.063911887610794	-0.22\\
nan	nan\\
0.0639832429478432	-0.22\\
0.0640208333333333	0.929135973038271\\
0.0640416666666667	0.918442742421215\\
0.0640785542774607	-0.22\\
nan	nan\\
0.0641499096145098	-0.22\\
0.0641875	0.929135973038179\\
0.0642083333333333	0.918442742420954\\
0.0642452209441273	-0.22\\
nan	nan\\
0.0643165762811765	-0.22\\
0.0643541666666667	0.929135973038087\\
0.064375	0.918442742421052\\
0.064411887610794	-0.22\\
nan	nan\\
0.0644832429478431	-0.22\\
0.0645208333333334	0.929135973038163\\
0.0645416666666667	0.918442742421151\\
0.0645785542774606	-0.22\\
nan	nan\\
0.0646499096145099	-0.22\\
0.0646875	0.929135973038239\\
0.0647083333333334	0.918442742421069\\
0.0647452209441273	-0.22\\
nan	nan\\
0.0648165762811765	-0.22\\
0.0648541666666667	0.929135973038147\\
0.064875	0.918442742421168\\
0.064911887610794	-0.22\\
nan	nan\\
0.0649832429478432	-0.22\\
0.0650208333333333	0.929135973038223\\
0.0650416666666667	0.918442742421086\\
0.0650785542774607	-0.22\\
nan	nan\\
0.0651499096145098	-0.22\\
0.0651875	0.929135973038131\\
0.0652083333333333	0.918442742421185\\
0.0652452209441273	-0.22\\
nan	nan\\
0.0653165762811765	-0.22\\
0.0653541666666667	0.929135973038039\\
0.065375	0.918442742421104\\
0.065411887610794	-0.22\\
nan	nan\\
0.0654832429478431	-0.22\\
0.0655208333333334	0.929135973038283\\
0.0655416666666667	0.918442742421202\\
0.0655785542774606	-0.22\\
nan	nan\\
0.0656499096145099	-0.22\\
0.0656875	0.929135973038191\\
0.0657083333333334	0.9184427424213\\
0.0657452209441273	-0.22\\
nan	nan\\
0.0658165762811765	-0.22\\
0.0658541666666667	0.929135973038099\\
0.065875	0.918442742421039\\
0.065911887610794	-0.22\\
nan	nan\\
0.0659832429478432	-0.22\\
0.0660208333333333	0.929135973038175\\
0.0660416666666667	0.918442742421138\\
0.0660785542774607	-0.22\\
nan	nan\\
0.0661499096145098	-0.22\\
0.0661875	0.929135973038251\\
0.0662083333333333	0.918442742421236\\
0.0662452209441273	-0.22\\
nan	nan\\
0.0663165762811765	-0.22\\
0.0663541666666667	0.929135973038159\\
0.066375	0.918442742421335\\
0.066411887610794	-0.22\\
nan	nan\\
0.0664832429478431	-0.22\\
0.0665208333333334	0.929135973038235\\
0.0665416666666667	0.918442742421073\\
0.0665785542774606	-0.22\\
nan	nan\\
0.0666529553732299	-0.22\\
0.0666666666666667	0.0682747864000546\\
0.0666875	-0.192266725442941\\
nan	nan\\
0.0669043138665041	-0.22\\
0.0670208333333333	0.946001075789569\\
0.0670416666666667	0.99766655428656\\
0.0670625	0.98134270583086\\
0.0671041666666666	0.753734350346657\\
0.0671875	-0.19226672544292\\
nan	nan\\
0.0674043138665041	-0.22\\
0.0675208333333334	0.946001075789599\\
0.0675416666666667	0.997666554286566\\
0.0675625	0.981342705830864\\
0.0676041666666667	0.753734350346671\\
0.0676875	-0.192266725442899\\
nan	nan\\
0.0679043138665041	-0.22\\
0.0680208333333333	0.946001075789592\\
0.0680416666666667	0.997666554286564\\
0.0680625	0.981342705830868\\
0.0681041666666666	0.753734350346685\\
0.0681875	-0.192266725442989\\
nan	nan\\
0.0684043138665041	-0.22\\
0.0685208333333334	0.946001075789585\\
0.0685416666666667	0.997666554286563\\
0.0685625	0.981342705830872\\
0.0686041666666667	0.753734350346625\\
0.0686875	-0.192266725442968\\
nan	nan\\
0.0689043138665041	-0.22\\
0.0690208333333333	0.946001075789578\\
0.0690416666666667	0.997666554286562\\
0.0690625	0.981342705830854\\
0.0691041666666666	0.753734350346639\\
0.0691875	-0.192266725442947\\
nan	nan\\
0.0694043138665041	-0.22\\
0.0695208333333334	0.946001075789571\\
0.0695416666666667	0.997666554286568\\
0.0695625	0.981342705830859\\
0.0696041666666667	0.753734350346653\\
0.0696875	-0.192266725442926\\
nan	nan\\
0.0699043138665041	-0.22\\
0.0700208333333333	0.946001075789601\\
0.0700416666666667	0.997666554286566\\
0.0700625	0.981342705830863\\
0.0701041666666666	0.753734350346667\\
0.0701875	-0.192266725443016\\
nan	nan\\
0.0704043138665041	-0.22\\
0.0705208333333334	0.946001075789594\\
0.0705416666666667	0.997666554286565\\
0.0705625	0.981342705830867\\
0.0706041666666667	0.753734350346607\\
0.0706875	-0.192266725442995\\
nan	nan\\
0.0709043138665041	-0.22\\
0.0710208333333333	0.946001075789587\\
0.0710416666666667	0.997666554286563\\
0.0710625	0.981342705830871\\
0.0711041666666666	0.753734350346621\\
0.0711875	-0.192266725442974\\
nan	nan\\
0.0714043138665041	-0.22\\
0.0715208333333334	0.94600107578958\\
0.0715416666666667	0.99766655428657\\
0.0715625	0.981342705830853\\
0.0716041666666667	0.753734350346635\\
0.0716875	-0.192266725442953\\
nan	nan\\
0.0719043138665041	-0.22\\
0.0720208333333333	0.946001075789536\\
0.0720416666666667	0.997666554286568\\
0.0720625	0.981342705830857\\
0.0721041666666666	0.753734350346649\\
0.0721875	-0.192266725442932\\
nan	nan\\
0.0724043138665041	-0.22\\
0.0725208333333334	0.946001075789529\\
0.0725416666666667	0.997666554286567\\
0.0725625	0.981342705830862\\
0.0726041666666667	0.753734350346589\\
0.0726875	-0.192266725443022\\
nan	nan\\
0.0729043138665041	-0.22\\
0.0730208333333333	0.946001075789559\\
0.0730416666666667	0.997666554286565\\
0.0730625	0.981342705830866\\
0.0731041666666666	0.753734350346603\\
0.0731875	-0.192266725443001\\
nan	nan\\
0.0734043138665041	-0.22\\
0.0735208333333334	0.946001075789552\\
0.0735416666666667	0.997666554286564\\
0.0735625	0.981342705830848\\
0.0736041666666667	0.753734350346617\\
0.0736875	-0.19226672544298\\
nan	nan\\
0.0739043138665041	-0.22\\
0.0740208333333333	0.946001075789545\\
0.0740416666666667	0.99766655428657\\
0.0740625	0.981342705830852\\
0.0741041666666666	0.753734350346631\\
0.0741875	-0.192266725442958\\
nan	nan\\
0.0744043138665041	-0.22\\
0.0745208333333334	0.946001075789538\\
0.0745416666666666	0.997666554286553\\
0.0745625	0.981342705830856\\
0.0746041666666667	0.753734350346645\\
0.0746875	-0.192266725443049\\
nan	nan\\
0.0749043138665041	-0.22\\
0.0750208333333333	0.946001075789531\\
0.0750416666666667	0.997666554286559\\
0.0750625	0.98134270583086\\
0.0751041666666666	0.753734350346585\\
0.0751875	-0.192266725442805\\
nan	nan\\
0.0754043138665041	-0.22\\
0.0755208333333334	0.946001075789561\\
0.0755416666666666	0.997666554286558\\
0.0755625	0.981342705830843\\
0.0756041666666667	0.753734350346599\\
0.0756875	-0.192266725442895\\
nan	nan\\
0.0759043138665041	-0.22\\
0.0760208333333333	0.946001075789554\\
0.0760416666666667	0.997666554286556\\
0.0760625	0.981342705830847\\
0.0761041666666666	0.753734350346613\\
0.0761875	-0.192266725442874\\
nan	nan\\
0.0764043138665041	-0.22\\
0.0765208333333334	0.946001075789547\\
0.0765416666666666	0.997666554286555\\
0.0765625	0.981342705830873\\
0.0766041666666667	0.753734350346627\\
0.0766875	-0.192266725442853\\
nan	nan\\
0.0769043138665041	-0.22\\
0.0770208333333333	0.94600107578954\\
0.0770416666666667	0.997666554286561\\
0.0770625	0.981342705830877\\
0.0771041666666666	0.753734350346567\\
0.0771875	-0.192266725442832\\
nan	nan\\
0.0774043138665041	-0.22\\
0.0775208333333334	0.946001075789569\\
0.0775416666666666	0.99766655428656\\
0.0775625	0.981342705830881\\
0.0776041666666667	0.753734350346581\\
0.0776875	-0.192266725442922\\
nan	nan\\
0.0779043138665041	-0.22\\
0.0780208333333333	0.946001075789562\\
0.0780416666666667	0.997666554286558\\
0.0780625	0.981342705830885\\
0.0781041666666666	0.753734350346595\\
0.0781875	-0.192266725442901\\
nan	nan\\
0.0784043138665041	-0.22\\
0.0785208333333334	0.946001075789556\\
0.0785416666666666	0.997666554286557\\
0.0785625	0.98134270583089\\
0.0786041666666667	0.753734350346609\\
0.0786875	-0.19226672544288\\
nan	nan\\
0.0789043138665041	-0.22\\
0.0790208333333333	0.946001075789549\\
0.0790416666666667	0.997666554286555\\
0.0790625	0.981342705830872\\
0.0791041666666666	0.753734350346549\\
0.0791875	-0.192266725442858\\
nan	nan\\
0.0794043138665042	-0.22\\
0.0795208333333334	0.946001075789578\\
0.0795416666666666	0.997666554286562\\
0.0795625	0.981342705830876\\
0.0796041666666667	0.753734350346563\\
0.0796875	-0.192266725442837\\
nan	nan\\
0.0799043138665041	-0.22\\
0.0800208333333333	0.946001075789571\\
0.0800416666666667	0.99766655428656\\
0.0800625	0.98134270583088\\
0.0801041666666666	0.753734350346577\\
0.0801875	-0.192266725442928\\
nan	nan\\
0.0804043138665042	-0.22\\
0.0805208333333334	0.946001075789564\\
0.0805416666666666	0.997666554286559\\
0.0805625	0.981342705830884\\
0.0806041666666667	0.753734350346591\\
0.0806875	-0.192266725442906\\
nan	nan\\
0.0809043138665041	-0.22\\
0.0810208333333333	0.946001075789557\\
0.0810416666666667	0.997666554286565\\
0.0810625	0.981342705830867\\
0.0811041666666666	0.75373435034668\\
0.0811875	-0.192266725442885\\
nan	nan\\
0.081404313866504	-0.22\\
0.0815208333333334	0.946001075789587\\
0.0815416666666666	0.997666554286564\\
0.0815625	0.981342705830871\\
0.0816041666666667	0.753734350346545\\
0.0816875	-0.192266725443087\\
nan	nan\\
0.0819043138665041	-0.22\\
0.0820208333333333	0.94600107578958\\
0.0820416666666667	0.997666554286562\\
0.0820625	0.981342705830831\\
0.0821041666666666	0.753734350346559\\
0.0821875	-0.192266725443066\\
nan	nan\\
0.0824043138665042	-0.22\\
0.0825208333333334	0.946001075789573\\
0.0825416666666666	0.997666554286576\\
0.0825625	0.981342705830835\\
0.0826041666666667	0.753734350346573\\
0.0826875	-0.192266725443045\\
nan	nan\\
0.0829043138665041	-0.22\\
0.0830208333333333	0.94600107578964\\
0.0830416666666667	0.997666554286575\\
0.0830625	0.981342705830839\\
0.0831041666666666	0.753734350346587\\
0.0831875	-0.192266725443024\\
};
\end{axis}
\end{tikzpicture}%

On calcule ensuite la densité spectrale de puissance de ce signal NRZ en utilisant la fonction \verb|pwelch| de Matlab utilisant un périodogramme de Welch. 
On obtient ceci:

%%%% This file was created by matlab2tikz.
%
%The latest updates can be retrieved from
%  http://www.mathworks.com/matlabcentral/fileexchange/22022-matlab2tikz-matlab2tikz
%where you can also make suggestions and rate matlab2tikz.
%
\definecolor{mycolor1}{rgb}{0.00000,0.44700,0.74100}%
\definecolor{mycolor2}{rgb}{0.85000,0.32500,0.09800}%
%
\begin{tikzpicture}

\begin{axis}[%
width=4.521in,
height=3.548in,
at={(0.758in,0.499in)},
scale only axis,
xmin=-25000,
xmax=25000,
xlabel style={font=\color{white!15!black}},
xlabel={Fréquence [Hz]},
ymode=log,
ymin=1e-10,
ymax=1254817.63975307,
yminorticks=true,
ylabel style={font=\color{white!15!black}},
ylabel={Amplitude},
axis background/.style={fill=white},
title style={font=\bfseries},
title={Densité spectrale de puissance},
legend style={legend cell align=left, align=left, draw=white!15!black}
]
\addplot [color=mycolor1]
  table[row sep=crcr]{%
-24000	2.93245261130722\\
-23987.9969992498	2.89077613629063\\
-23975.9939984996	0.337796928205646\\
-23963.9909977494	0.252653981571482\\
-23951.9879969993	1.53699011475019\\
-23939.9849962491	1.98665469309985\\
-23927.9819954989	0.394623283010367\\
-23915.9789947487	0.870337217110672\\
-23903.9759939985	3.00092352595294\\
-23891.9729932483	0.160460131251335\\
-23879.9699924981	2.25955922471988\\
-23867.9669917479	0.832737079710491\\
-23855.9639909978	3.08142848265614\\
-23843.9609902476	3.10122960312069\\
-23831.9579894974	0.843761970569595\\
-23819.9549887472	2.25126663272548\\
-23807.951987997	0.14428884542775\\
-23795.9489872468	2.97071149792039\\
-23783.9459864966	0.879844311478809\\
-23771.9429857464	0.40969038670103\\
-23759.9399849963	2.04328236692584\\
-23747.9369842461	1.54251014600905\\
-23735.9339834959	0.251638597891975\\
-23723.9309827457	0.334129463129934\\
-23711.9279819955	2.90730361210714\\
-23699.9249812453	2.93362346983434\\
-23687.9219804951	2.87659758604057\\
-23675.9189797449	0.341872951072737\\
-23663.9159789948	0.254021202081828\\
-23651.9129782446	1.53274486098469\\
-23639.9099774944	1.93238531873766\\
-23627.9069767442	0.380143920252613\\
-23615.903975994	0.862469512157413\\
-23603.9009752438	3.03376473991942\\
-23591.8979744936	0.178617862139056\\
-23579.8949737434	2.27039786273808\\
-23567.8919729932	0.822436560860676\\
-23555.8889722431	3.06486551928798\\
-23543.8859714929	3.12433730621235\\
-23531.8829707427	0.855548672428561\\
-23519.8799699925	2.24548809355008\\
-23507.8769692423	0.130048188515074\\
-23495.8739684921	2.94301868688397\\
-23483.8709677419	0.891022239754996\\
-23471.8679669918	0.425398353582142\\
-23459.8649662416	2.10246844356079\\
-23447.8619654914	1.54932479334973\\
-23435.8589647412	0.250972923397418\\
-23423.855963991	0.330858458587292\\
-23411.8529632408	2.92623292368317\\
-23399.8499624906	2.93713984994984\\
-23387.8469617404	2.86472286748413\\
-23375.8439609902	0.346371699782668\\
-23363.8409602401	0.255744374404866\\
-23351.8379594899	1.52975899991699\\
-23339.8349587397	1.8802864234572\\
-23327.8319579895	0.366202712931229\\
-23315.8289572393	0.856220767638056\\
-23303.8259564891	3.06935514495098\\
-23291.8229557389	0.198835750047045\\
-23279.8199549888	2.28382723593382\\
-23267.8169542386	0.812825658769073\\
-23255.8139534884	3.05148741696331\\
-23243.8109527382	3.15083567093718\\
-23231.807951988	0.868137588296084\\
-23219.8049512378	2.24220403474953\\
-23207.8019504876	0.117699817448339\\
-23195.7989497374	2.91774426678561\\
-23183.7959489872	0.90391369194543\\
-23171.7929482371	0.441804259965013\\
-23159.7899474869	2.16442688506746\\
-23147.7869467367	1.55745850841423\\
-23135.7839459865	0.250656814945863\\
-23123.7809452363	0.327973801136335\\
-23111.7779444861	2.94762516256161\\
-23099.7749437359	2.94301319199335\\
-23087.7719429857	2.85511439434091\\
-23075.7689422356	0.351309522629181\\
-23063.7659414854	0.257829627324626\\
-23051.7629407352	1.52802147662469\\
-23039.759939985	1.83018103634207\\
-23027.7569392348	0.352753239423854\\
-23015.7539384846	0.851581486757717\\
-23003.7509377344	3.10782576528664\\
-22991.7479369842	0.221206093624287\\
-22979.7449362341	2.29990514756874\\
-22967.7419354839	0.803872044077969\\
-22955.7389347337	3.04125564701331\\
-22943.7359339835	3.18082512216222\\
-22931.7329332333	0.881572391845358\\
-22919.7299324831	2.24140716751569\\
-22907.7269317329	0.107222702020436\\
-22895.7239309827	2.89479573803742\\
-22883.7209302326	0.918572983476628\\
-22871.7179294824	0.458969603980683\\
-22859.7149287322	2.22938725277313\\
-22847.711927982	1.56694055296454\\
-22835.7089272318	0.250692134705682\\
-22823.7059264816	0.325467295617561\\
-22811.7029257314	2.97155002775214\\
-22799.6999249812	2.9512626526222\\
-22787.6969242311	2.84774183209058\\
-22775.6939234809	0.356705081745803\\
-22763.6909227307	0.260285153662721\\
-22751.6879219805	1.52752547576998\\
-22739.6849212303	1.78190172508886\\
-22727.6819204801	0.339751908742888\\
-22715.6789197299	0.848553168649402\\
-22703.6759189797	3.14931973178476\\
-22691.6729182296	0.245840749125871\\
-22679.6699174794	2.31870308236285\\
-22667.6669167292	0.795545574760647\\
-22655.663915979	3.03414624327087\\
-22643.6609152288	3.21442330231724\\
-22631.6579144786	0.895900364485954\\
-22619.6549137284	2.24310248670085\\
-22607.6519129782	0.0986132276773663\\
-22595.6489122281	2.87408817598231\\
-22583.6459114779	0.935066591091585\\
-22571.6429107277	0.476960830573465\\
-22559.6399099775	2.29759663770237\\
-22547.6369092273	1.5778052364942\\
-22535.6339084771	0.251082795483984\\
-22523.6309077269	0.323332620345191\\
-22511.6279069767	2.99808631779213\\
-22499.6249062266	2.96191524010352\\
-22487.6219054764	2.84258189199473\\
-22475.6189047262	0.362579512589777\\
-22463.615903976	0.263121285251713\\
-22451.6129032258	1.52826836375683\\
-22439.6099024756	1.73528939132622\\
-22427.6069017254	0.327157662611023\\
-22415.6039009752	0.847148522318217\\
-22403.6009002251	3.19399348991841\\
-22391.5978994749	0.272872399820162\\
-22379.5948987247	2.34030693395138\\
-22367.5918979745	0.787818073458179\\
-22355.5888972243	3.03014975368597\\
-22343.5858964741	3.25176613610279\\
-22331.5828957239	0.911172767609649\\
-22319.5798949737	2.24730740866239\\
-22307.5768942236	0.0918855680595813\\
-22295.5738934734	2.85554353404008\\
-22283.5708927232	0.953473854743925\\
-22271.567891973	0.495849922209744\\
-22259.5648912228	2.36932181112546\\
-22247.5618904726	1.59009219784202\\
-22235.5588897224	0.25183483148984\\
-22223.5558889722	0.321565302703761\\
-22211.5528882221	3.02732249496194\\
-22199.5498874719	2.97500600599022\\
-22187.5468867217	2.83961817230211\\
-22175.5438859715	0.368956612048867\\
-22163.5408852213	0.266350593821455\\
-22151.5378844711	1.53025166210949\\
-22139.5348837209	1.6901921566127\\
-22127.5318829707	0.31493170972127\\
-22115.5288822206	0.847391837789556\\
-22103.5258814704	3.24201814704913\\
-22091.5228807202	0.302456149584776\\
-22079.51987997	2.36481790495073\\
-22067.5168792198	0.780663118087681\\
-22055.5138784696	3.02927135030606\\
-22043.5108777194	3.29300910618743\\
-22031.5078769692	0.927445255208966\\
-22019.5048762191	2.25405205169897\\
-22007.5018754689	0.0870723406057673\\
-21995.4988747187	2.83908999025313\\
-21983.4958739685	0.973887859007472\\
-21971.4928732183	0.515715064693339\\
-21959.4898724681	2.44485162955433\\
-21947.4868717179	1.60384673534202\\
-21935.4838709677	0.252956496442217\\
-21923.4808702176	0.320162715214279\\
-21911.4778694674	3.05935732877506\\
-21899.4748687172	2.99057829578593\\
-21887.471867967	2.83884104510435\\
-21875.4688672168	0.375863058017318\\
-21863.4658664666	0.269988019536878\\
-21851.4628657164	1.53348105240496\\
-21839.4598649662	1.64646432699618\\
-21827.4568642161	0.303037289421408\\
-21815.4538634659	0.849319525185231\\
-21803.4508627157	3.29358098059591\\
-21791.4478619655	0.334771474148968\\
-21779.4448612153	2.39235359757888\\
-21767.4418604651	0.774055842596389\\
-21755.4388597149	3.03153110304759\\
-21743.4358589647	3.33832876159324\\
-21731.4328582146	0.94477833316084\\
-21719.4298574644	2.26337966505291\\
-21707.4268567142	0.0842255655337081\\
-21695.423855964	2.82466132898844\\
-21683.4208552138	0.996416511598691\\
-21671.4178544636	0.536641399707359\\
-21659.4148537134	2.524499733883\\
-21647.4118529632	1.61912019038794\\
-21635.4088522131	0.254458391167531\\
-21623.4058514629	0.319124092591839\\
-21611.4028507127	3.09430062702897\\
-21599.3998499625	3.00868406125447\\
-21587.3968492123	2.84024758691216\\
-21575.3938484621	0.383328664168927\\
-21563.3908477119	0.274051029749709\\
-21551.3878469617	1.53796641185047\\
-21539.3848462116	1.6039654240035\\
-21527.3818454614	0.291439462880619\\
-21515.3788447112	0.852980835876522\\
-21503.375843961	3.34888713104217\\
-21491.3728432108	0.370024572727591\\
-21479.3698424606	2.42304931532396\\
-21467.3668417104	0.767972745637406\\
-21455.3638409602	3.03696442544864\\
-21443.3608402101	3.38792448660613\\
-21431.3578394599	0.963237871767052\\
-21419.3548387097	2.27534721748129\\
-21407.3518379595	0.0834179564915747\\
-21395.3488372093	2.81219634727796\\
-21383.3458364591	1.02118384144275\\
-21371.3428357089	0.558721877208388\\
-21359.3398349587	2.60860758985221\\
-21347.3368342086	1.63597039030113\\
-21335.3338334584	0.256353623830698\\
-21323.3308327082	0.318450570627264\\
-21311.327831958	3.13227406394839\\
-21299.3248312078	3.02938423918046\\
-21287.3218304576	2.84384155291145\\
-21275.3188297074	0.391386674453737\\
-21263.3158289572	0.278559810965881\\
-21251.3128282071	1.54372188095474\\
-21239.3098274569	1.56255927295202\\
-21227.3068267067	0.280104930694768\\
-21215.3038259565	0.858438785145219\\
-21203.3008252063	3.40816150887026\\
-21191.2978244561	0.408451172983491\\
-21179.2948237059	2.45705960221935\\
-21167.2918229557	0.762391504202411\\
-21155.2888222056	3.04562270512531\\
-21143.2858214554	3.44202056246928\\
-21131.2828207052	0.982895679754912\\
-21119.279819955	2.29002615870067\\
-21107.2768192048	0.084744580100556\\
-21095.2738184546	2.80163827718837\\
-21083.2708177044	1.04833154438325\\
-21071.2678169542	0.582058223901524\\
-21059.2648162041	2.69754792442408\\
-21047.2618154539	1.65446215816371\\
-21035.2588147037	0.258658006462785\\
-21023.2558139535	0.318145248328034\\
-21011.2528132033	3.17341211606033\\
-20999.2498124531	3.05274920031633\\
-20987.2468117029	2.8496333933778\\
-20975.2438109527	0.400074102447876\\
-20963.2408102026	0.283537497742674\\
-20951.2378094524	1.55076596324986\\
-20939.2348087022	1.52211313981007\\
-20927.231807952	0.269001876018504\\
-20915.2288072018	0.865771300912716\\
-20903.2258064516	3.47165094873928\\
-20891.2228057014	0.450319854839245\\
-20879.2198049512	2.4945600512902\\
-20867.2168042011	0.757290789632899\\
-20855.2138034509	3.05757413543375\\
-20843.2108027007	3.50086856181878\\
-20831.2078019505	1.00383014906014\\
-20819.2048012003	2.30750337135593\\
-20807.2018004501	0.0883249320957411\\
-20795.1987996999	2.79293421422777\\
-20783.1957989497	1.07802081053724\\
-20771.1927981996	0.606762045931557\\
-20759.1897974494	2.79172862259404\\
-20747.1867966992	1.67466789753795\\
-20735.183795949	0.261390292482912\\
-20723.1807951988	0.318213275127336\\
-20711.1777944486	3.21786311963121\\
-20699.1747936984	3.07885927526928\\
-20687.1717929482	2.8576403131137\\
-20675.1687921981	0.409432121916176\\
-20663.1657914479	0.289010443369654\\
-20651.1627906977	1.55912165934814\\
-20639.1597899475	1.48249690978292\\
-20627.1567891973	0.258099833910652\\
-20615.1537884471	0.875072629152447\\
-20603.1507876969	3.53962664895553\\
-20591.1477869467	0.495935973087574\\
-20579.1447861965	2.5357494209647\\
-20567.1417854464	0.752650083401107\\
-20555.1387846962	3.07290476960184\\
-20543.135783946	3.56475012342522\\
-20531.1327831958	1.026126980915\\
-20519.1297824456	2.32788233598872\\
-20507.1267816954	0.0943054900974129\\
-20495.1237809452	2.78603454227265\\
-20483.1207801951	1.1104344752866\\
-20471.1177794449	0.632956088527443\\
-20459.1147786947	2.8915971605371\\
-20447.1117779445	1.69666826262449\\
-20435.1087771943	0.264572461077699\\
-20423.1057764441	0.318661965741071\\
-20411.1027756939	3.26579046333322\\
-20399.0997749437	3.10780536409922\\
-20387.0967741935	2.86788637404379\\
-20375.0937734434	0.419506516093412\\
-20363.0907726932	0.295008537963749\\
-20351.087771943	1.56881663647505\\
-20339.0847711928	1.44358230214392\\
-20327.0817704426	0.24736958768143\\
-20315.0787696924	0.886455034176004\\
-20303.0757689422	3.6123869433376\\
-20291.0727681921	0.545646276459063\\
-20279.0697674419	2.58085210587486\\
-20267.0667666917	0.748449489669195\\
-20255.0637659415	3.09171982473972\\
-20243.0607651913	3.63398016417145\\
-20231.0577644411	1.04988000593579\\
-20219.0547636909	2.35128453715096\\
-20207.0517629407	0.102862817871057\\
-20195.0487621905	2.78089234420007\\
-20183.0457614404	1.14577954501821\\
-20171.0427606902	0.660775678532786\\
-20159.03975994	2.99764566544042\\
-20147.0367591898	1.72055292531827\\
-20135.0337584396	0.26823005534972\\
-20123.0307576894	0.319500945634534\\
-20111.0277569392	3.31737393436929\\
-20099.0247561891	3.13968963752821\\
-20087.0217554389	2.88040264175482\\
-20075.0187546887	0.430348194728744\\
-20063.0157539385	0.301565580866704\\
-20051.0127531883	1.57988343654722\\
-20039.0097524381	1.40524211730885\\
-20027.0067516879	0.236783094533201\\
-20015.0037509377	0.900050841135244\\
-20003.0007501875	3.69026045785642\\
-19990.9977494374	0.599844342011941\\
-19978.9947486872	2.63012101814283\\
-19966.991747937	0.744669541674435\\
-19954.9887471868	3.11414526913762\\
-19942.9857464366	3.70891059758193\\
-19930.9827456864	1.075192112965\\
-19918.9797449362	2.37785114434818\\
-19906.976744186	0.114207313866424\\
-19894.9737434359	2.77746278640288\\
-19882.9707426857	1.18429015925098\\
-19870.9677419355	0.690370381566154\\
-19858.9647411853	3.11041670890767\\
-19846.9617404351	1.74642145348663\\
-19834.9587396849	0.272392582977159\\
-19822.9557389347	0.320742331354262\\
-19810.9527381845	3.37281123717681\\
-19798.9497374344	3.17462633970935\\
-19786.9467366842	2.89522737750446\\
-19774.943735934	0.442013789222864\\
-19762.9407351838	0.30871971589534\\
-19750.9377344336	1.59235972541214\\
-19738.9347336834	1.36734951381196\\
-19726.9317329332	0.226313443773468\\
-19714.928732183	0.916014878494747\\
-19702.9257314329	3.77360971742176\\
-19690.9227306827	0.658976969876397\\
-19678.9197299325	2.6838409475752\\
-19666.9167291823	0.741290998363332\\
-19654.9137284321	3.14032973539789\\
-19642.9107276819	3.78993464039873\\
-19630.9077269317	1.1021763038014\\
-19618.9047261815	2.40774501100118\\
-19606.9017254314	0.12858771869931\\
-19594.8987246812	2.77570246497625\\
-19582.895723931	1.22623106472266\\
-19570.8927231808	0.721905911220483\\
-19558.8897224306	3.23050996167224\\
-19546.8867216804	1.77438431716776\\
-19534.8837209302	0.277093989784586\\
-19522.88072018	0.322400950335164\\
-19510.8777194299	3.43231970843754\\
-19498.8747186797	3.21274270391862\\
-19486.8717179295	2.91240627755092\\
-19474.8687171793	0.454566338949626\\
-19462.8657164291	0.316513939445428\\
-19450.8627156789	1.60628858733854\\
-19438.8597149287	1.3297773154529\\
-19426.8567141785	0.215934852390033\\
-19414.8537134284	0.934527392083682\\
-19402.8507126782	3.8628352781238\\
-19390.847711928	0.723551715007864\\
-19378.8447111778	2.74233248241033\\
-19366.8417104276	0.738294627765743\\
-19354.8387096774	3.17044681036325\\
-19342.8357089272	3.87749180729751\\
-19330.832708177	1.13095689440678\\
-19318.8297074269	2.44115304251094\\
-19306.8267066767	0.146296523250586\\
-19294.8237059265	2.775568698467\\
-19282.8207051763	1.27190169303104\\
-19270.8177044261	0.755566335747538\\
-19258.8147036759	3.35858986330469\\
-19246.8117029257	1.80456404252673\\
-19234.8087021755	0.282373219131097\\
-19222.8057014254	0.324494606456034\\
-19210.8027006752	3.49613825507882\\
-19198.799699925	3.2541799941403\\
-19186.7966991748	2.93199276170989\\
-19174.7936984246	0.468076083758969\\
-19162.7906976744	0.324996693888832\\
-19150.7876969242	1.62171886930168\\
-19138.784696174	1.29239735089394\\
-19126.7816954239	0.205622704596033\\
-19114.7786946737	0.955797518446937\\
-19102.7756939235	3.95838047607897\\
-19090.7726931733	0.794145772440985\\
-19078.7696924231	2.80595659126195\\
-19066.7666916729	0.735660972692131\\
-19054.7636909227	3.20469776637005\\
-19042.7606901725	3.97207371312968\\
-19030.7576894224	1.16167088669828\\
-19018.7546886722	2.47828899747129\\
-19006.751687922	0.167676452559105\\
-18994.7486871718	2.77701875067976\\
-18982.7456864216	1.32164095364564\\
-18970.7426856714	0.791556636496227\\
-18958.7396849212	3.4953944900999\\
-18946.736684171	1.83709653730523\\
-18934.7336834209	0.288274872658785\\
-18922.7306826707	0.327044398778369\\
-18910.7276819205	3.56452954676861\\
-18898.7246811703	3.29909468817787\\
-18886.7216804201	2.95404831458768\\
-18874.7186796699	0.482621380969839\\
-18862.7156789197	0.33422256111154\\
-18850.7126781695	1.63870558069412\\
-18838.7096774194	1.25507983186245\\
-18826.7066766692	0.195353644222132\\
-18814.703675919	0.980067425644669\\
-18802.7006751688	4.06073690224595\\
-18790.6976744186	0.871416481049296\\
-18778.6946736684	2.87511998654533\\
-18766.6916729182	0.733370093940196\\
-18754.688672168	3.24331481160339\\
-18742.6856714179	4.07423082816403\\
-18730.6826706677	1.19446953979532\\
-18718.6796699175	2.51939679967665\\
-18706.6766691673	0.193128242419996\\
-18694.6736684171	2.78000896460039\\
-18682.6706676669	1.37583287932981\\
-18670.6676669167	0.830105684434148\\
-18658.6646661665	3.64174584483648\\
-18646.6616654164	1.87213261637083\\
-18634.6586646662	0.294849991604903\\
-18622.655663916	0.330075102592631\\
-18610.6526631658	3.63778249945562\\
-18598.6496624156	3.34765982025145\\
-18586.6466616654	2.97864288199267\\
-18574.6436609152	0.498289768955769\\
-18562.640660165	0.34425307444445\\
-18550.6376594149	1.65731035521319\\
-18538.6346586647	1.21769278018761\\
-18526.6316579145	0.185105731752627\\
-18514.6286571643	1.00761725542145\\
-18502.6256564141	4.17045073432626\\
-18490.6226556639	0.95611377251038\\
-18478.6196549137	2.95028141877267\\
-18466.6166541635	0.731401285412841\\
-18454.6136534134	3.28656495702003\\
-18442.6106526632	4.18458036231621\\
-18430.607651913	1.22952017499338\\
-18418.6046511628	2.56475445708164\\
-18406.6016504126	0.22311997817133\\
-18394.5986496624	2.78449378430372\\
-18382.5956489122	1.43491329181687\\
-18370.592648162	0.871469714657782\\
-18358.5896474119	3.79856183698395\\
-18346.5866466617	1.90983976135725\\
-18334.5836459115	0.302156982276855\\
-18322.5806451613	0.333615624265347\\
-18310.5776444111	3.71621509310959\\
-18298.5746436609	3.40006650401525\\
-18286.5716429107	3.00585532709311\\
-18274.5686421605	0.515179203803809\\
-18262.5656414104	0.355157671084619\\
-18250.5626406602	1.6776019831216\\
-18238.55963991	1.1801015198164\\
-18226.5566391598	0.174858681461028\\
-18214.5536384096	1.03877103221554\\
-18202.5506376594	4.2881300854646\\
-18190.5476369092	1.04909496845018\\
-18178.544636159	3.03195908222619\\
-18166.5416354089	0.729732754571867\\
-18154.5386346587	3.33475461928573\\
-18142.5356339085	4.30381549402339\\
-18130.5326331583	1.26700825599285\\
-18118.5296324081	2.61467870607284\\
-18106.5266316579	0.258198331448931\\
-18094.5236309077	2.79042463881846\\
-18082.5206301575	1.49937769535782\\
-18070.5176294074	0.91593639719924\\
-18058.5146286572	3.96687028347379\\
-18046.511627907	1.95040415591237\\
-18034.5086271568	0.310262714514593\\
-18022.5056264066	0.337699543601463\\
-18010.5026256564	3.80017757409102\\
-17998.4996249062	3.45652566053732\\
-17986.496624156	3.03577395022086\\
-17974.4936234059	0.533399501356621\\
-17962.4906226557	0.367014812224264\\
-17950.4876219055	1.69965702376065\\
-17938.4846211553	1.14216825779236\\
-17926.4816204051	0.164594199029792\\
-17914.4786196549	1.07390374548904\\
-17902.4756189047	4.41445356266593\\
-17890.4726181545	1.15134242627058\\
-17878.4696174044	3.12073935596759\\
-17866.4666166542	0.728341260402118\\
-17854.463615904	3.38823510873593\\
-17842.4606151538	4.43271620678633\\
-17830.4576144036	1.30713979370436\\
-17818.4546136534	2.66953052746018\\
-17806.4516129032	0.299002114966017\\
-17794.448612153	2.79774865604249\\
-17782.4456114029	1.56979065401681\\
-17770.4426106527	0.963829624207807\\
-17758.4396099025	4.14782533117069\\
-17746.4366091523	1.99403304656538\\
-17734.4336084021	0.319243829215505\\
-17722.4306076519	0.342365761227889\\
-17710.4276069017	3.89005610104703\\
-17698.4246061515	3.5172699800935\\
-17686.4216054014	3.06849707792309\\
-17674.4186046512	0.553074023950918\\
-17662.415603901	0.379913304282826\\
-17650.4126031508	1.7235605102726\\
-17638.4096024006	1.103751788016\\
-17626.4066016504	0.154296446193408\\
-17614.4036009002	1.11344986281221\\
-17602.40060015	4.55018027228014\\
-17590.3975993999	1.2639846570428\\
-17578.3945986497	3.21728715620454\\
-17566.3915978995	0.727201699761482\\
-17554.3885971493	3.44740918780186\\
-17542.3855963991	4.57216205887797\\
-17530.3825956489	1.35014413620623\\
-17518.3795948987	2.72972171682586\\
-17506.3765941485	0.346278683148299\\
-17494.3735933984	2.80640716931347\\
-17482.3705926482	1.64679697244824\\
-17470.367591898	1.01551516169958\\
-17458.3645911478	4.34272679741532\\
-17446.3615903976	2.04095749007616\\
-17434.3585896474	0.329188299391074\\
-17422.3555888972	0.347659272273716\\
-17410.352588147	3.9862769048681\\
-17398.3495873968	3.58255615132788\\
-17386.3465866467	3.10413372623156\\
-17374.3435858965	0.574341659898433\\
-17362.3405851463	0.393953862481109\\
-17350.3375843961	1.74940676087718\\
-17338.3345836459	1.06470736652485\\
-17326.3315828957	0.143952667267176\\
-17314.3285821455	1.1579135972243\\
-17302.3255813954	4.6961615623874\\
-17290.3225806452	1.38832169620003\\
-17278.319579895	3.32235824304272\\
-17266.3165791448	0.726286631235113\\
-17254.3135783946	3.5127389319852\\
-17242.3105776444	4.72314728644694\\
-17230.3075768942	1.39627721652269\\
-17218.304576144	2.79572273714989\\
-17206.3015753938	0.400903844976777\\
-17194.2985746437	2.81633397183839\\
-17170.2925731433	1.07140734950238\\
-17158.2895723931	4.55304303976174\\
-17146.2865716429	2.09143556219913\\
-17134.2835708927	0.340197300452029\\
-17122.2805701425	0.353632093005603\\
-17110.2775693924	4.08931104546132\\
-17098.2745686422	3.65266739742845\\
-17086.271567892	3.1428043450838\\
-17074.2685671418	0.597359154716488\\
-17062.2655663916	0.409250968046063\\
-17050.2625656414	1.77730031491013\\
-17038.2595648912	1.02488682508657\\
-17026.256564141	0.133554023216886\\
-17014.2535633908	1.20788133681646\\
-17002.2505626407	4.85335486075228\\
-16990.2475618905	1.52585571012185\\
-16978.2445611403	3.43681390906028\\
-16966.2415603901	0.725565723118654\\
-16954.2385596399	3.58475518524921\\
-16942.2355588897	4.88679874006375\\
-16930.2325581395	1.44582534784908\\
-16918.2295573893	2.86807213973052\\
-16906.2265566392	0.463906134582099\\
-16894.223555889	2.82745326460418\\
-16882.2205551388	1.82365310781483\\
-16870.2175543886	1.13197707820017\\
-16858.2145536384	4.78043811660771\\
-16846.2115528882	2.14575611965378\\
-16834.208552138	0.352387459667311\\
-16822.2055513878	0.360344373863411\\
-16810.2025506377	4.19967986343538\\
-16798.1995498875	3.72791636592719\\
-16786.1965491373	3.18464165076022\\
-16774.1935483871	0.622303867191709\\
-16762.1905476369	0.425935083\\
-16750.1875468867	1.80735701503445\\
-16738.1845461365	0.984139016559549\\
-16726.1815453863	0.123096693422037\\
-16714.1785446362	1.26403675381\\
-16702.175543886	5.02284005419032\\
-16690.1725431358	1.67832808282091\\
-16678.1695423856	3.56163858847724\\
-16666.1665416354	0.725005109906327\\
-16654.1635408852	3.66406898024526\\
-16642.160540135	5.06439727883304\\
-16630.1575393848	1.4991096756099\\
-16618.1545386347	2.94738791585562\\
-16606.1515378845	0.536496518478439\\
-16594.1485371343	2.83967723372668\\
-16582.1455363841	1.92532832777296\\
-16570.1425356339	1.19776132936454\\
-16558.1395348837	5.02680418914111\\
-16546.1365341335	2.20424322909943\\
-16534.1335333833	0.365893572864264\\
-16522.1305326332	0.36786574073041\\
-16510.127531883	4.31796124476257\\
-16498.1245311328	3.80864842700066\\
-16486.1215303826	3.22979155357747\\
-16474.1185296324	0.649377040780701\\
-16462.1155288822	0.4441553028725\\
-16450.112528132	1.83970526295401\\
-16438.1095273818	0.942310720931241\\
-16426.1065266317	0.112583324507562\\
-16414.1035258815	1.3271792532743\\
-16402.1005251313	5.20583896375493\\
-16390.0975243811	1.84776456697944\\
-16378.0945236309	3.69796106435066\\
-16366.0915228807	0.724566638093823\\
-16354.0885221305	3.7513853926285\\
-16342.0855213803	5.25740340641825\\
-16330.0825206302	1.55649142179634\\
-16318.07951988	3.03438123889291\\
-16306.0765191298	0.620104924211657\\
-16294.0735183796	2.85290317882687\\
-16282.0705176294	2.03729069904243\\
-16270.0675168792	1.26937464027202\\
-16258.064516129	5.29430036129656\\
-16246.0615153788	2.26726140501211\\
-16234.0585146287	0.380871900428471\\
-16222.0555138785	0.376276917537784\\
-16210.0525131283	4.44479683761791\\
-16198.0495123781	3.89524544571692\\
-16186.0465116279	3.27841418978998\\
-16174.0435108777	0.678807703433737\\
-16162.0405101275	0.464082548700711\\
-16150.0375093773	1.874487482752\\
-16138.0345086272	0.899248190201014\\
-16126.031507877	0.102024931560881\\
-16114.0285071268	1.39824660829945\\
-16102.0255063766	5.403738615551\\
-16090.0225056264	2.03653052948703\\
-16078.0195048762	3.8470801368937\\
-16066.016504126	0.724206977952939\\
-16054.0135033758	3.84752043281363\\
-16042.0105026257	5.46748814351946\\
-16030.0075018755	1.61837808954209\\
-16018.0045011253	3.12987318553927\\
-16006.0015003751	0.716425378775076\\
-15993.9984996249	2.86701009904426\\
-15981.9954988747	2.16085161238712\\
-15969.9924981245	1.3475229505027\\
-15957.9894973743	5.58539947046212\\
-15945.9864966242	2.33522183405326\\
-15933.983495874	0.397504185341772\\
-15921.9804951238	0.385671697582169\\
-15909.9774943736	4.58090039108285\\
-15897.9744936234	3.98813010552263\\
-15885.9714928732	3.33068506562354\\
-15873.968492123	0.710857337980498\\
-15861.9654913728	0.485913427404359\\
-15849.9624906227	1.91186183416966\\
-15837.9594898725	0.854799578385036\\
-15825.9564891223	0.0914433921845732\\
-15813.9534883721	1.47834287565246\\
-15801.9504876219	5.61811918898473\\
-15789.9474868717	2.24739890912728\\
-15777.9444861215	4.0104958561232\\
-15765.9414853713	0.723876573122558\\
-15753.9384846212	3.95342175158643\\
-15741.935483871	5.69657040206045\\
-15729.9324831208	1.6852308376003\\
-15717.9294823706	3.23481519487994\\
-15705.9264816204	0.827472081725154\\
-15693.9234808702	2.88185462397889\\
-15681.92048012	2.29753912400257\\
-15669.9174793698	1.43302041494205\\
-15657.9144786197	5.90294475699814\\
-15645.9114778695	2.40858981040784\\
-15633.9084771193	0.416002577134535\\
-15621.9054763691	0.396159349794397\\
-15609.9024756189	4.72706741750824\\
-15597.8994748687	4.08777087490176\\
-15585.8964741185	3.3867963231968\\
-15573.8934733683	0.745825502459487\\
-15561.8904726182	0.509874925481994\\
-15549.887471868	1.95200423105381\\
-15537.8844711178	0.80881859840047\\
-15525.8814703676	0.0808747215956818\\
-15513.8784696174	1.56877301466476\\
-15501.8754688672	5.85078776668188\\
-15489.872468117	2.48363428651386\\
-15477.8694673668	4.18994773824854\\
-15465.8664666167	0.723518393369857\\
-15453.8634658665	4.07019416841062\\
-15441.8604651163	5.94686248605972\\
-15429.8574643661	1.7575732886387\\
-15417.8544636159	3.35031424915575\\
-15405.8514628657	0.955649456982197\\
-15393.8484621155	2.8972661547969\\
-15381.8454613653	2.44914141360818\\
-15369.8424606152	1.52680993365274\\
-15357.839459865	6.25021888458185\\
-15345.8364591148	2.4878936671092\\
-15333.8334583646	0.436615699382783\\
-15321.8304576144	0.407867571083815\\
-15309.8274568642	4.88418642501778\\
-15297.824456114	4.19468772853521\\
-15285.8214553638	3.44695813657863\\
-15273.8184546137	0.784056628992388\\
-15261.8154538635	0.536230149013388\\
-15249.8124531133	1.99511073461993\\
-15237.8094523631	0.76116988060405\\
-15225.8064516129	0.0703733828613696\\
-15213.8034508627	1.67108607329736\\
-15201.8004501125	6.10381932756549\\
-15189.7974493623	2.749097513825\\
-15177.7944486122	4.38746181016747\\
-15165.791447862	0.723066448290983\\
-15153.7884471118	4.19913134241968\\
-15141.7854463616	6.22092581829699\\
-15129.7824456114	1.83600210500209\\
-15117.7794448612	3.47766406437864\\
-15105.776444111	1.10384019834461\\
-15093.7734433608	2.91304105489712\\
-15081.7704426107	2.61776071751754\\
-15069.7674418605	1.62998837087504\\
-15057.7644411103	6.63102850654245\\
-15045.7614403601	2.57373556918752\\
-15033.7584396099	0.459636170861809\\
-15021.7554388597	0.420946128903554\\
-15009.7524381095	5.05325201842424\\
-14997.7494373593	4.30945875406472\\
-14985.7464366092	3.51140024753655\\
-14973.743435859	0.825948293633354\\
-14961.7404351088	0.565285386520868\\
-14949.7374343586	2.04140041387096\\
-14937.7344336084	0.711736696789746\\
-14925.7314328582	0.0600179772116611\\
-14913.728432108	1.78712940211996\\
-14901.7254313578	6.37960684587277\\
-14889.7224306077	3.0483767693291\\
-14877.7194298575	4.60540889114397\\
-14865.7164291073	0.722444010040282\\
-14853.7134283571	4.34175532759873\\
-14841.7104276069	6.52173962924542\\
-14829.7074268567	1.92119975735525\\
-14817.7044261065	3.61838398849386\\
-14805.7014253563	1.27551664686054\\
-14793.6984246062	2.92893569941829\\
-14781.695423856	2.80588070605893\\
-14769.6924231058	1.74383773200086\\
-14757.6894223556	7.04980853818686\\
-14745.6864216054	2.66680464063965\\
-14733.6834208552	0.485409987726249\\
-14721.680420105	0.435571382512894\\
-14709.6774193548	5.23538023304272\\
-14697.6744186047	4.43272780543316\\
-14685.6714178545	3.58037364702878\\
-14673.6684171043	0.871961335325911\\
-14661.6654163541	0.597398857290335\\
-14649.6624156039	2.09111879356082\\
-14637.6594148537	0.660431981467325\\
-14625.6564141035	0.0499187865355322\\
-14613.6534133533	1.91911717193586\\
-14601.6504126032	6.68092292248847\\
-14589.647411853	3.38695283443171\\
-14577.6444111028	4.84657728946613\\
-14565.6414103526	0.721561481459283\\
-14553.6384096024	4.49986632926117\\
-14541.6354088522	6.85278620279552\\
-14529.632408102	2.01395003287055\\
-14517.6294073518	3.77426786639376\\
-14505.6264066017	1.47488265746769\\
-14493.6234058515	2.94465816160668\\
-14481.6204051013	3.01645125290151\\
-14469.6174043511	1.86986396798403\\
-14457.6144036009	7.51175160096191\\
-14445.6114028507	2.76789304110181\\
-14433.6084021005	0.514348306633325\\
-14421.6054013503	0.451951931429755\\
-14409.6024006002	5.43182654938885\\
-14397.59939985	4.56521339496864\\
-14385.5963990998	3.6541524083675\\
-14373.5933983496	0.922632315568945\\
-14361.5903975994	0.632991624482647\\
-14349.5873968492	2.14454205206576\\
-14337.584396099	0.607213967239312\\
-14325.5813953488	0.0402278190610738\\
-14313.5783945987	2.06971759385929\\
-14301.5753938485	7.01099613931119\\
-14289.5723930983	3.77140905222519\\
-14277.5693923481	5.11426414150924\\
-14265.5663915979	0.720313831220128\\
-14253.5633908477	4.67560577280093\\
-14241.5603900975	7.21815744465963\\
-14229.5573893473	2.11515699102698\\
-14217.5543885972	3.94744590630141\\
-14205.551387847	1.70705563810083\\
-14193.5483870968	2.95985827994908\\
-14181.5453863466	3.25299589374377\\
-14169.5423855964	2.00984562440444\\
-14157.5393848462	8.02296987854119\\
-14145.536384096	2.87791580248455\\
-14133.5333833458	0.546942350828341\\
-14121.5303825957	0.470335721405171\\
-14109.5273818455	5.64400714082511\\
-14097.5243810953	4.70771905948427\\
-14085.5213803451	3.733035672365\\
-14073.5183795949	0.978588964613066\\
-14061.5153788447	0.672561312918731\\
-14049.5123780945	2.20198218785555\\
-14037.5093773443	0.552108307625342\\
-14025.5063765942	0.031152263871371\\
-14013.503375844	2.24216480388986\\
-14001.5003750938	7.37360636701385\\
-13989.4973743436	4.2097001322755\\
-13977.4943735934	5.41239106977465\\
-13965.4913728432	0.718577500024581\\
-13953.488372093	4.87153690300681\\
-13941.4853713428	7.6226891542992\\
-13929.4823705927	2.22586829235236\\
-13917.4793698425	4.14046366346698\\
-13905.4763690923	1.97830198843691\\
-13893.4733683421	2.97411582267073\\
-13881.4703675919	3.51974915209079\\
-13869.4673668417	2.16589530709245\\
-13857.4643660915	8.59069907378147\\
-13845.4613653413	2.99793550243108\\
-13833.4583645911	0.583782413506157\\
-13821.455363841	0.491019051641469\\
-13809.4523630908	5.87352403973737\\
-13797.4493623406	4.8611454866919\\
-13785.4463615904	3.81734977905233\\
-13773.4433608402	1.04056946983559\\
-13761.44036009	0.716699494378227\\
-13749.4373593398	2.26379345645391\\
-13737.4343585896	0.495239374134625\\
-13725.4313578395	0.022972625602839\\
-13713.4283570893	2.44040357583785\\
-13701.4253563391	7.77320469713244\\
-13689.4223555889	4.71149917436627\\
-13677.4193548387	5.74565186086304\\
-13665.4163540885	0.71620665980202\\
-13653.4133533383	5.09074869400916\\
-13641.4103525881	8.07213162789325\\
-13629.407351838	2.34730412099819\\
-13617.4043510878	4.35638378371988\\
-13605.4013503376	2.29634419217367\\
-13593.3983495874	2.98692644807562\\
-13581.3953488372	3.82183356068795\\
-13569.392348087	2.34053798990569\\
-13557.3893473368	9.22355755938164\\
-13545.3863465866	3.12919322263422\\
-13533.3833458365	0.625582287744656\\
-13521.3803450863	0.51435808479786\\
-13509.3773443361	6.1221950774753\\
-13497.3743435859	5.02650475305943\\
-13485.3713428357	3.90745052933352\\
-13473.3683420855	1.10944675137984\\
-13461.3653413353	0.766113916823531\\
-13449.3623405851	2.33038049836731\\
-13437.359339835	0.436874613766323\\
-13425.3563390848	0.0160673422557117\\
-13413.3533383346	2.66927815799728\\
-13401.3503375844	8.21506567623998\\
-13389.3473368342	5.28864970187505\\
-13377.344336084	6.11970272415649\\
-13365.3413353338	0.713028682430017\\
-13353.3383345836	5.33699107650189\\
-13341.3353338335	8.5733683916917\\
-13329.3323330833	2.48089332897744\\
-13317.3293323331	4.59891833282712\\
-13305.3263315829	2.67076503398581\\
-13293.3233308327	2.99768517533637\\
-13281.3203300825	4.16548998959254\\
-13269.3173293323	2.53681167611453\\
-13257.3143285821	9.93187861813464\\
-13245.311327832	3.2731477595435\\
-13233.3083270818	0.673210955223951\\
-13221.3053263316	0.540783684627421\\
-13209.3023255814	6.39208967454943\\
-13197.2993248312	5.20493710713731\\
-13185.296324081	4.00372554560805\\
-13173.2933233308	1.18625927465993\\
-13161.2903225806	0.821657200411344\\
-13149.2873218305	2.40220875446074\\
-13137.2843210803	0.377487643594294\\
-13125.2813203301	0.0109464705367402\\
-13113.2783195799	2.93478104757407\\
-13101.2753188297	8.70548236065099\\
-13089.2723180795	5.95576018048629\\
-13077.2693173293	6.54140978327585\\
-13065.2663165791	0.708838644068309\\
-13053.263315829	5.61485270847944\\
-13041.2603150788	9.13469938979204\\
-13029.2573143286	2.62831899596026\\
-13017.2543135784	4.87260269666214\\
-13005.2513128282	3.11354490510415\\
-12993.248312078	3.00566715335129\\
-12981.2453113278	4.55838035734409\\
-12969.2423105776	2.75839805043651\\
-12957.2393098275	10.7281405148079\\
-12945.2363090773	3.43152579738281\\
-12933.2333083271	0.727734090903073\\
-12921.2303075769	0.570820725644466\\
-12909.2273068267	6.68557184148504\\
-12897.2243060765	5.39773083360063\\
-12885.2213053263	4.10659667393028\\
-12873.2183045761	1.27225052006909\\
-12861.215303826	0.884364267480336\\
-12849.2123030758	2.47981802802665\\
-12837.2093023256	0.317848462668484\\
-12825.2063015754	0.00829818716607979\\
-12813.2033008252	3.24438411472487\\
-12801.200300075	9.25201878047832\\
-12789.1972993248	6.73099410882311\\
-12777.1942985746	7.01917438497811\\
-12765.1912978245	0.703392657996505\\
-12753.1882970743	5.929997232941\\
-12741.1852963241	9.76621150608567\\
-12729.1822955739	2.79157639778293\\
-12717.1792948237	5.18302667767846\\
-12705.1762940735	3.63978363025611\\
-12693.1732933233	3.01000571006751\\
-12681.1702925731	5.00998981419989\\
-12669.167291823	3.00979383882218\\
-12657.1642910728	11.6275290551505\\
-12645.1612903226	3.6063868250355\\
-12633.1582895724	0.790468996738265\\
-12621.1552888222	0.605113482847921\\
-12609.152288072	7.00535212792447\\
-12597.1492873218	5.60634586801406\\
-12585.1462865716	4.21652233658715\\
-12573.1432858215	1.36892004968128\\
-12561.1402850713	0.95550171789125\\
-12549.1372843211	2.56384045300655\\
-12537.1342835709	0.2591532999421\\
-12525.1312828207	0.00905361268163909\\
-12513.1282820705	3.60748426904432\\
-12501.1252813203	9.86384031725303\\
-12489.1222805701	7.63713190541408\\
-12477.11927982	7.56336553890934\\
-12465.1162790698	0.69639979511695\\
-12453.1132783196	6.28948098034324\\
-12441.1102775694	10.4802689264567\\
-12429.1072768192	2.97304752068495\\
-12417.104276069	5.53714533196242\\
-12405.1012753188	4.26868139190442\\
-12393.0982745686	3.00966807950791\\
-12381.0952738185	5.53216740777129\\
-12369.0922730683	3.29653812464909\\
-12357.0892723181	12.6486818472051\\
-12345.0862715679	3.80020815107782\\
-12333.0832708177	0.863058142512662\\
-12321.0802700675	0.644459392587098\\
-12309.0772693173	7.35455075544326\\
-12297.0742685671	5.83244199936263\\
-12285.071267817	4.33399968685032\\
-12273.0682670668	1.47809030139939\\
-12261.0652663166	1.03663376623378\\
-12249.0622655664	2.65502474763927\\
-12237.0592648162	0.203213043158643\\
-12225.056264066	0.0144781599566918\\
-12213.0532633158	4.03601058559293\\
-12201.0502625656	10.5521512363236\\
-12189.0472618155	8.70301564870925\\
-12177.0442610653	8.18690185188349\\
-12165.0412603151	0.687512325073661\\
-12153.0382595649	6.70218567470684\\
-12141.0352588147	11.2921702280327\\
-12129.0322580645	3.17559792118897\\
-12117.0292573143	5.94370257557944\\
-12105.0262565641	5.02488851633062\\
-12093.023255814	3.0034300333604\\
-12081.0202550638	6.13986227459951\\
-12069.0172543136	3.62551764579351\\
-12057.0142535634	13.814685208943\\
-12045.0112528132	4.01599772195875\\
-12033.008252063	0.947568845452151\\
-12021.0052513128	0.689854498231162\\
-12009.0022505626	7.73677483879039\\
-11996.9992498125	6.07791271792038\\
-11984.9962490623	4.45956633744998\\
-11972.9932483121	1.60199500274712\\
-11960.9902475619	1.12971148540392\\
-11948.9872468117	2.75426961007431\\
-11936.9842460615	0.152729332906147\\
-11924.9812453113	0.0263018117384255\\
-11912.9782445611	4.54526235348326\\
-11900.975243811	11.3307817368783\\
-11888.9722430608	9.96554108423102\\
-11876.9692423106	8.9060451598027\\
-11864.9662415604	0.676314006596958\\
-11852.9632408102	7.1794160032852\\
-11840.96024006	12.2210409231355\\
-11828.9572393098	3.40270418990488\\
-11816.9542385596	6.41381676624935\\
-11804.9512378095	5.94038832983916\\
-11792.9482370593	2.98985223540954\\
-11780.9452363091	6.85214016645221\\
-11768.9422355589	4.00538240670633\\
-11756.9392348087	15.1544276843479\\
-11744.9362340585	4.25744602068843\\
-11732.9332333083	1.04663022717722\\
-11720.9302325581	0.742555453870472\\
-11708.927231808	8.15621351271775\\
-11696.9242310578	6.34492604898587\\
-11684.9212303076	4.59380131041775\\
-11672.9182295574	1.74339774810104\\
-11660.9152288072	1.2371953852836\\
-11648.912228057	2.86267068608478\\
-11636.9092273068	0.111703684233571\\
-11624.9062265566	0.0469074006181004\\
-11612.9032258065	5.15508267373656\\
-11600.9002250563	12.2169869523593\\
-11588.8972243061	11.4724445641577\\
-11576.8942235559	9.74149880356196\\
-11564.8912228057	0.662306208264595\\
-11552.8882220555	7.73573745597702\\
-11540.8852213053	13.2910640184864\\
-11528.8822205551	3.65862398116905\\
-11516.879219805	6.96180290746126\\
-11504.8762190548	7.05716315575\\
-11492.8732183046	2.96726414552185\\
-11480.8702175544	7.69360870715812\\
-11468.8672168042	4.44711993379273\\
-11456.864216054	16.7044652690978\\
-11444.8612153038	4.52913386839908\\
-11432.8582145536	1.16362421789612\\
-11420.8552138035	0.804165382226572\\
-11408.8522130533	8.61775604079786\\
-11396.8492123031	6.63597408771547\\
-11384.8462115529	4.73732467082627\\
-11372.8432108027	1.90575335280288\\
-11360.8402100525	1.36222652792498\\
-11348.8372093023	2.98158812521018\\
-11336.8342085521	0.0860516391515543\\
-11324.831207802	0.0796067603261735\\
-11312.8282070518	5.89152819925784\\
-11300.8252063016	13.2325516710512\\
-11288.8222055514	13.2862651743762\\
-11276.8192048012	10.7199521068494\\
-11264.816204051	0.644891815825951\\
-11252.8132033008	8.39017071725602\\
-11240.8102025506	14.5332046343629\\
-11228.8072018005	3.94862626141677\\
-11216.8042010503	7.60634700209779\\
-11204.8012003001	8.43103179034572\\
-11192.7981995499	2.93376695959668\\
-11180.7951987997	8.69644967401912\\
-11168.7921980495	4.96486183775285\\
-11156.7891972993	18.511634343443\\
-11144.7861965491	4.83682174644358\\
-11132.783195799	1.30295636330265\\
-11120.7801950488	0.876754732505848\\
-11108.7771942986	9.12713975815642\\
-11096.7741935484	6.95393344141011\\
-11084.7711927982	4.89079502784522\\
-11072.768192048	2.09343098082249\\
-11060.7651912978	1.50886971488601\\
-11048.7621905476	3.1127460917687\\
-11036.7591897975	0.0845384877165072\\
-11024.7561890473	0.129052476231427\\
-11012.7531882971	6.78928667973493\\
-11000.7501875469	14.4053445580094\\
-10988.7471867967	15.4900782489517\\
-10976.7441860465	11.8762912360118\\
-10964.7411852963	0.623357318987673\\
-10952.7381845461	9.16792580333784\\
-10940.735183796	15.9876712641712\\
-10928.7321830458	4.27930837672672\\
-10916.7291822956	8.37221529822094\\
-10904.7261815454	10.137274641161\\
-10892.7231807952	2.88727775974249\\
-10880.720180045	9.90337133619218\\
-10868.7171792948	5.57703743115809\\
-10856.7141785446	20.6367793026915\\
-10844.7111777945	5.18786054601216\\
-10832.7081770443	1.47044686354106\\
-10820.7051762941	0.963034542806348\\
-10808.7021755439	9.69113725420593\\
-10796.6991747937	7.30213944413258\\
-10784.6961740435	5.05490366544579\\
-10772.6931732933	2.31202829344515\\
-10760.6901725431	1.68246604573315\\
-10748.6871717929	3.25838309740241\\
-10736.6841710428	0.120229315039123\\
-10724.6811702926	0.201863226444501\\
-10712.6781695424	7.8952458526125\\
-10700.6751687922	15.7715474753889\\
-10688.672168042	18.1959570033971\\
-10676.6691672918	13.2568269799939\\
-10664.6661665416	0.596854452607379\\
-10652.6631657915	10.1029714599025\\
-10640.6601650413	17.7074998277444\\
-10628.6571642911	4.65904095797317\\
-10616.6541635409	9.29279457961604\\
-10604.6511627907	12.2790478516458\\
-10592.6481620405	2.82565745466678\\
-10580.6451612903	11.3719871631561\\
-10568.6421605401	6.30805749813235\\
-10556.6391597899	23.1601795284139\\
-10544.6361590398	5.59178749759313\\
-10532.6331582896	1.67390687228021\\
-10520.6301575394	1.06660996274169\\
-10508.6271567892	10.317795966885\\
-10496.624156039	7.68447789069201\\
-10484.6211552888	5.23036342785146\\
-10472.6181545386	2.56882277030477\\
-10460.6151537884	1.89015550985819\\
-10448.6121530383	3.42148531897376\\
-10436.6091522881	0.212780047810677\\
-10424.6061515379	0.307593379969194\\
-10412.6031507877	9.27387750082935\\
-10400.6001500375	17.3789230249586\\
-10388.5971492873	21.5577377348723\\
-10376.5941485371	14.9241121828621\\
-10364.5911477869	0.564384887926606\\
-10352.5881470368	11.2419292417462\\
-10340.5851462866	19.7638917057415\\
-10328.5821455364	5.09860556584151\\
-10316.5791447862	10.4139565328011\\
-10304.576144036	15.0002574755748\\
-10292.5731432858	2.74700469102109\\
-10280.5701425356	13.1814623237829\\
-10268.5671417854	7.19082833554948\\
-10256.5641410353	26.1896325526302\\
-10244.5611402851	6.06121201644021\\
-10232.5581395349	1.92400753187111\\
-10220.5551387847	1.19235985958183\\
-10208.5521380345	11.0167490952072\\
-10196.5491372843	8.10549923593695\\
-10184.5461365341	5.41788953024198\\
-10172.5431357839	2.87343506341704\\
-10160.5401350338	2.14167116343142\\
-10148.5371342836	3.60615935213191\\
-10136.5341335334	0.392138490543219\\
-10124.5311327832	0.460271989309994\\
-10112.528132033	11.0155603428263\\
-10100.5251312828	19.2917217905856\\
-10088.5221305326	25.7907590256709\\
-10076.5191297824	16.964312171042\\
-10064.5161290323	0.524795995533192\\
-10052.5131282821	12.6501270749644\\
-10040.5101275319	22.2543684662373\\
-10028.5071267817	5.61213079211737\\
-10016.5041260315	11.8000920810526\\
-10004.5011252813	18.5057669730928\\
-9992.49812453114	2.65027647975809\\
-9980.49512378095	15.4428694006981\\
-9968.49212303076	8.27060308188116\\
-9956.48912228057	29.8728076400726\\
-9944.48612153038	6.61316901379724\\
-9932.4831207802	2.23562489894563\\
-9920.48012003001	1.34702025871412\\
-9908.47711927982	11.7996257790858\\
-9896.47411852963	8.57056185167392\\
-9884.47111777945	5.61816804936317\\
-9872.46811702926	3.23882965270338\\
-9860.46511627907	2.45058049762819\\
-9848.46211552888	3.8182467556474\\
-9836.4591147787	0.704677613347027\\
-9824.45611402851	0.680911561520054\\
-9812.45311327832	13.2498061616397\\
-9800.45011252813	21.5982582474245\\
-9788.44711177795	31.203251527018\\
-9776.44411102776	19.4988025105329\\
-9764.44111027757	0.476805414846802\\
-9752.43810952738	14.4212843480583\\
-9740.4351087772	25.3155933912502\\
-9728.43210802701	6.21850475779677\\
-9716.42910727682	13.5438171187554\\
-9704.42610652663	23.0940479968427\\
-9692.42310577644	2.53655651269392\\
-9680.42010502626	18.3158088568301\\
-9668.41710427607	9.6110566691565\\
-9656.41410352588	34.4166881638873\\
-9644.4111027757	7.27125149286106\\
-9632.40810202551	2.62998390286609\\
-9620.40510127532	1.54010825657324\\
-9608.40210052513	12.6806034160122\\
-9596.39909977495	9.08601321246373\\
-9584.39609902476	5.83180583561867\\
-9572.39309827457	3.68286996610972\\
-9560.39009752438	2.83628902382305\\
-9548.3870967742	4.06637285585873\\
-9536.38409602401	1.22366387436733\\
-9524.38109527382	1.00172478010146\\
-9512.37809452363	16.1669516312161\\
-9500.37509377344	24.4229814762273\\
-9488.37209302326	38.247877765742\\
-9476.36909227307	22.703014730311\\
-9464.36609152288	0.419093395409326\\
-9452.36309077269	16.6935243003375\\
-9440.36009002251	29.1442056515819\\
-9428.35708927232	6.94357493221771\\
-9416.35408852213	15.7821196378409\\
-9404.35108777195	29.2117188672116\\
-9392.34808702176	2.41163017823292\\
-9380.34508627157	22.0360058391203\\
-9368.34208552138	11.3041955299191\\
-9356.33908477119	40.1192204111278\\
-9344.33608402101	8.06909257614698\\
-9332.33308327082	3.13819454864333\\
-9320.33008252063	1.78543637092468\\
-9308.32708177044	13.677170915697\\
-9296.32408102026	9.65942108587862\\
-9284.32108027007	6.05925297939305\\
-9272.31807951988	4.23082002462319\\
-9260.31507876969	3.32739540262402\\
-9248.31207801951	4.36380706887022\\
-9236.30907726932	2.06774959624283\\
-9224.30607651913	1.47346766786875\\
-9212.30307576894	20.0550540079733\\
-9200.30007501876	27.9464210904662\\
-9188.29707426857	47.6095416005692\\
-9176.29407351838	26.838217945205\\
-9164.29107276819	0.350549653407224\\
-9152.28807201801	19.676869701407\\
-9140.28507126782	34.0319885126066\\
-9128.28207051763	7.82370250061856\\
-9116.27906976744	18.724290620904\\
-9104.27606901725	37.5482172597433\\
-9092.27306826707	2.29126819118466\\
-9080.27006751688	26.9629704001974\\
-9068.26706676669	13.4871655992304\\
-9056.26406601651	47.4228918149261\\
-9044.26106526632	9.05628889321871\\
-9032.25806451613	3.80732046366561\\
-9020.25506376594	2.10369633584456\\
-9008.25206301576	14.8112196924436\\
-8996.24906226557	10.2998713476577\\
-8984.24606151538	6.30068651914661\\
-8972.24306076519	4.9195344721075\\
-8960.24006001501	3.96755526424877\\
-8948.23705926482	4.7319103630701\\
-8936.23405851463	3.43498862993601\\
-8924.23105776444	2.17876934175199\\
-8912.22805701425	25.3653730844459\\
-8900.22505626407	32.4395796717305\\
-8888.22205551388	60.3615892463681\\
-8876.21905476369	32.3074932752214\\
-8864.2160540135	0.270873394820217\\
-8852.21305326332	23.7025939237955\\
-8840.21005251313	40.4279376315195\\
-8828.20705176294	8.91175931768331\\
-8816.20405101276	22.7024881686582\\
-8804.20105026257	49.2077277856975\\
-8792.19804951238	2.21233596565326\\
-8780.19504876219	33.6661921503752\\
-8768.19204801201	16.3720982505147\\
-8756.18904726182	57.0097140789033\\
-8744.18604651163	10.3089725022209\\
-8732.18304576144	4.71129168349648\\
-8720.18004501125	2.52708144027304\\
-8708.17704426107	16.1106747982632\\
-8696.17404351088	11.0183556030825\\
-8684.17104276069	6.55584542625844\\
-8672.1680420105	5.80482024233565\\
-8660.16504126032	4.82625802690968\\
-8636.15903975994	5.66849540670326\\
-8624.15603900975	3.2575512350301\\
-8612.15303825957	32.8345871572342\\
-8600.15003750938	38.3263022554407\\
-8588.14703675919	78.2582473798779\\
-8576.14403600901	39.7595230613699\\
-8564.14103525882	0.182001174205365\\
-8552.13803450863	29.3165710685871\\
-8540.13503375844	49.0537516092493\\
-8528.13203300825	10.2877817435924\\
-8516.12903225806	28.2683511416087\\
-8504.12603150788	66.0386347586357\\
-8492.12303075769	2.25704968844555\\
-8480.1200300075	43.0898193674297\\
-8468.11702925732	20.3020885732209\\
-8456.11402850713	69.979141922505\\
-8444.11102775694	11.9497866104598\\
-8432.10802700675	5.9716576883716\\
-8420.10502625657	3.10803834947061\\
-8408.10202550638	17.6120860917267\\
-8396.09902475619	11.828281372926\\
-8384.096024006	6.82382628189279\\
-8372.09302325582	6.97513630342797\\
-8360.09002250563	6.01985943678356\\
-8348.08702175544	5.85354238367733\\
-8336.08402100525	9.39079092304157\\
-8324.08102025506	4.9584691305036\\
-8312.07801950488	43.7271455262028\\
-8300.07501875469	46.3034648965835\\
-8288.0720180045	104.31692084451\\
-8276.06901725431	50.2933618797993\\
-8264.06601650413	0.0915641728607937\\
-8252.06301575394	37.4673901792955\\
-8240.06001500375	61.132889436927\\
-8228.05701425357	12.0791192536396\\
-8216.05401350338	36.3900291875932\\
-8204.05101275319	91.3072871775631\\
-8192.048012003	2.60878459013396\\
-8180.04501125282	56.8889268296372\\
-8168.04201050263	25.863406789248\\
-8156.03900975244	88.2044335410151\\
-8144.03600900225	14.188327198119\\
-8132.03300825206	7.79983922040432\\
-8120.03000750187	3.93702432922546\\
-8108.02700675169	19.3650968011518\\
-8096.0240060015	12.7461518464611\\
-8084.02100525131	7.10294087960241\\
-8072.01800450113	8.57894994631742\\
-8060.01500375094	7.75574746305156\\
-8048.01200300075	6.79656319247069\\
-8036.00900225056	15.7991759220366\\
-8024.00600150038	7.75092094459824\\
-8012.00300075019	60.3529944237652\\
-8000	57.5905491618052\\
-7987.99699924981	144.068998096081\\
-7975.99399849963	65.894537468142\\
-7963.99099774944	0.0216567802662538\\
-7951.98799699925	49.9162022808783\\
-7939.98499624906	78.8821044761983\\
-7927.98199549887	14.5016076637538\\
-7915.97899474869	48.8876679369547\\
-7903.9759939985	131.191678539495\\
-7891.97299324831	3.68985096078387\\
-7879.96999249813	78.1752096653655\\
-7867.96699174794	34.1298249201248\\
-7855.96399099775	115.108490866786\\
-7843.96099024756	17.4103332341565\\
-7831.95798949738	10.5907921933888\\
-7819.95498874719	5.18178511652411\\
-7807.951987997	21.4410478605861\\
-7795.94898724681	13.7924837257285\\
-7783.94598649663	7.39111856012421\\
-7771.94298574644	10.884415970017\\
-7759.93998499625	10.4339720697103\\
-7747.93698424606	8.29571946277411\\
-7735.93398349587	27.3769218186231\\
-7723.93098274569	12.5925356560003\\
-7711.9279819955	87.2824739164312\\
-7699.92498124531	74.4991671679466\\
-7687.92198049512	208.516964309664\\
-7675.91897974494	90.4581744436843\\
-7663.91597899475	0.0328224293021715\\
-7651.91297824456	70.2225415119331\\
-7639.90997749438	106.677134731817\\
-7627.90697674419	17.9524195972579\\
-7615.903975994	69.4996287096762\\
-7603.90097524381	198.451764781738\\
-7591.89797449363	6.53544848689516\\
-7579.89497374344	113.354420909654\\
-7567.89197299325	47.2540623662791\\
-7555.88897224306	157.543662406582\\
-7543.88597149287	22.3971346248773\\
-7531.88297074269	15.1546749918648\\
-7519.8799699925	7.18444690366655\\
-7507.87696924231	23.9521408346512\\
-7495.87396849212	14.9930648555116\\
-7483.87096774194	7.68896490580707\\
-7471.86796699175	14.4253870467553\\
-7459.86496624156	14.9104660754109\\
-7447.86196549137	10.9531565282398\\
-7435.85896474119	49.8216808973087\\
-7423.855963991	21.6499529536902\\
-7411.85296324081	134.569271980964\\
-7399.84996249063	101.903430849355\\
-7387.84696174044	322.061111989046\\
-7375.84396099025	132.514276772867\\
-7363.84096024006	0.297441738273464\\
-7351.83795948987	106.445423555861\\
-7339.83495873968	154.206510226298\\
-7327.8319579895	23.2485723107867\\
-7315.82895723931	106.871053461318\\
-7303.82595648912	322.767797969255\\
-7291.82295573893	13.9501435568715\\
-7279.81995498875	177.335405817937\\
-7267.81695423856	70.1168312708647\\
-7255.81395348837	231.046376578061\\
-7243.81095273819	30.9504793225378\\
-7231.807951988	23.3805814114456\\
-7219.80495123781	10.7401511306422\\
-7207.80195048762	27.1030541762658\\
-7195.79894973744	16.3807055137734\\
-7183.79594898725	8.01538849454986\\
-7171.79294823706	20.4187015480953\\
-7159.78994748687	23.2972276519095\\
-7147.78694673668	16.380584967287\\
-7135.7839459865	98.1988035492384\\
-7123.78094523631	40.5907020377518\\
-7111.77794448612	227.86710267593\\
-7099.77494373593	151.737880944277\\
-7087.77194298575	547.762032909865\\
-7075.76894223556	213.878276931909\\
-7063.76594148537	1.36472183312161\\
-7051.76294073519	180.010543674522\\
-7039.759939985	246.719958311213\\
-7027.75693923481	32.3616005721091\\
-7015.75393848462	184.674343396568\\
-7003.75093773444	586.012103557058\\
-6991.74793698425	34.8251126629626\\
-6979.74493623406	311.525671695138\\
-6967.74193548387	115.942252473006\\
-6955.73893473368	377.828936535114\\
-6943.7359339835	48.0813180803669\\
-6931.73293323331	40.5821835150989\\
-6919.72993248312	18.0869678609387\\
-6907.72693173293	31.3703001243038\\
-6895.72393098275	17.9977212426819\\
-6883.72093023256	8.49175002988929\\
-6871.71792948237	32.2581745386869\\
-6859.71492873218	42.0321454279244\\
-6847.711927982	29.9186055719936\\
-6835.70892723181	221.890348982356\\
-6823.70592648162	87.8858550054902\\
-6811.70292573144	449.148500093267\\
-6799.69992498125	260.691840931838\\
-6787.69692423106	1091.51647749567\\
-6775.69392348087	405.335293887113\\
-6763.69092273068	5.39976915483158\\
-6751.6879219805	363.936116724089\\
-6739.68492123031	468.965478763273\\
-6727.68192048012	51.3997003378818\\
-6715.67891972993	387.352371837069\\
-6703.67591897974	1283.74716156373\\
-6691.67291822956	105.79713469105\\
-6679.66991747937	668.27118324171\\
-6667.66691672918	232.399617520824\\
-6655.663915979	751.38903236234\\
-6643.66091522881	92.6259642163298\\
-6631.65791447862	87.0492014882281\\
-6619.65491372843	37.8422565650849\\
-6607.65191297825	38.4105689748921\\
-6595.64891222806	19.8995238135775\\
-6583.64591147787	9.94206964271097\\
-6571.64291072768	63.2318547456983\\
-6559.63990997749	99.1150296418652\\
-6547.63690922731	77.000467591446\\
-6535.63390847712	655.389483207184\\
-6523.63090772693	251.661186231798\\
-6511.62790697674	1185.8962035196\\
-6499.62490622656	596.869212211002\\
-6487.62190547637	2969.69398580231\\
-6475.61890472618	1057.8994667953\\
-6463.615903976	25.3683764137763\\
-6451.61290322581	1048.599924978\\
-6439.60990247562	1272.67283838944\\
-6427.60690172543	110.412095015824\\
-6415.60390097525	1208.12895325619\\
-6403.60090022506	4189.14731736652\\
-6391.59789947487	479.886585093787\\
-6379.59489872468	2211.10619586468\\
-6367.59189797449	719.368659842024\\
-6355.58889722431	2344.03546410861\\
-6343.58589647412	293.54955723265\\
-6331.58289572393	306.712997646573\\
-6319.57989497374	133.095061955608\\
-6307.57689422356	60.4649747971562\\
-6295.57389347337	22.1599325585346\\
-6283.57089272318	22.6476824161627\\
-6271.56789197299	222.540342088662\\
-6259.56489122281	462.508360632168\\
-6247.56189047262	434.166024104446\\
-6235.55888972243	4036.53888926657\\
-6223.55588897225	1569.81740206505\\
-6211.55288822206	7158.44455127447\\
-6199.54988747187	3226.69454055639\\
-6187.54688672168	20513.3744100669\\
-6175.54388597149	7432.74551728567\\
-6163.54088522131	312.237629773571\\
-6151.53788447112	9514.45289364669\\
-6139.53488372093	11732.5242482912\\
-6127.53188297074	796.893726228774\\
-6115.52888222055	15996.5116292317\\
-6103.52588147037	64999.2238904828\\
-6091.52288072018	11821.4873418129\\
-6079.51987996999	48571.0873417871\\
-6067.51687921981	17978.7406016359\\
-6055.51387846962	77032.9851270426\\
-6043.51087771943	14585.0588117928\\
-6031.50787696924	26019.6575621363\\
-6019.50487621906	22876.5342829717\\
-6007.50187546887	16674.2439913723\\
-5995.49887471868	415989.961795309\\
-5983.49587396849	17176.5330496123\\
-5971.49287321831	22184.6116983552\\
-5959.48987246812	26113.7888288079\\
-5947.48687171793	15144.2404333163\\
-5935.48387096774	78443.7069433976\\
-5923.48087021755	18305.8160113776\\
-5911.47786946737	50157.3814193085\\
-5899.47486871718	12652.7372086062\\
-5887.47186796699	68356.1739508387\\
-5875.4688672168	16512.6073475319\\
-5863.46586646662	753.597870039406\\
-5851.46286571643	12029.7606974502\\
-5839.45986496624	9760.65488712021\\
-5827.45686421606	275.012592269145\\
-5815.45386346587	7851.48577758974\\
-5803.45086271568	22447.0608533422\\
-5791.44786196549	3741.91621184304\\
-5779.44486121531	7803.61499321949\\
-5767.44186046512	1649.40836802786\\
-5755.43885971493	4298.17889033651\\
-5743.43585896474	530.334665564024\\
-5731.43285821455	464.858963065025\\
-5719.42985746437	162.998155731759\\
-5707.42685671418	81.1541716822097\\
-5695.42385596399	28.1967485576963\\
-5683.4208552138	59.9658760808478\\
-5671.41785446362	82.6850241234607\\
-5659.41485371343	326.993975071962\\
-5647.41185296324	385.199061033971\\
-5635.40885221306	2646.81351770095\\
-5623.40585146287	797.996690971374\\
-5611.40285071268	2547.77266456231\\
-5599.39984996249	629.711669467886\\
-5587.3968492123	5063.80453723383\\
-5575.39384846212	1329.53676368348\\
-5563.39084771193	100.380704360802\\
-5551.38784696174	1385.19501432\\
-5539.38484621155	1129.49069981747\\
-5527.38184546136	13.1515772019442\\
-5515.37884471118	1250.20445116463\\
-5503.37584396099	3748.89001061307\\
-5491.3728432108	868.179112959118\\
-5479.36984246062	1473.49532942079\\
-5467.36684171043	276.170428037186\\
-5455.36384096024	734.91384937573\\
-5443.36084021005	120.458286015263\\
-5431.35783945987	96.0306718929393\\
-5419.35483870968	35.6944732237895\\
-5407.35183795949	67.3727150368459\\
-5395.3488372093	32.311154931515\\
-5383.34583645912	34.2429531166745\\
-5371.34283570893	13.1003366141904\\
-5359.33983495874	104.790005731287\\
-5347.33683420855	153.972264545076\\
-5335.33383345836	958.461204135293\\
-5323.33083270818	285.05332725338\\
-5311.32783195799	853.825544350504\\
-5299.3248312078	167.157889638012\\
-5287.32183045761	1783.8361042533\\
-5275.31882970743	445.496171170098\\
-5263.31582895724	48.3472350983942\\
-5251.31282820705	549.745281367917\\
-5239.30982745687	415.155179805867\\
-5227.30682670668	0.571012411150904\\
-5215.30382595649	551.751307932153\\
-5203.3008252063	1595.88721490466\\
-5191.29782445612	480.789582830957\\
-5179.29482370593	647.20516481097\\
-5167.29182295574	97.3849344692588\\
-5155.28882220555	251.477380241268\\
-5143.28582145536	60.7653361958045\\
-5131.28282070518	37.5779951178741\\
-5119.27981995499	15.3488699979314\\
-5107.2768192048	75.1285913665575\\
-5095.27381845461	37.5145022447793\\
-5083.27081770443	29.6572854764467\\
-5071.26781695424	3.40034296230673\\
-5059.26481620405	59.1769661352503\\
-5047.26181545386	102.002938579236\\
-5035.25881470368	553.714306511756\\
-5023.25581395349	159.288669263072\\
-5011.2528132033	439.374785267883\\
-4999.24981245312	64.2437401153623\\
-4987.24681170293	937.440093599289\\
-4975.24381095274	218.061307895442\\
-4963.24081020255	32.8792737256408\\
-4951.23780945236	317.42168243976\\
-4939.23480870217	222.548158694877\\
-4927.23180795199	1.78030453779239\\
-4915.2288072018	346.610111878938\\
-4903.22580645161	932.886872965243\\
-4891.22280570143	359.377605351047\\
-4879.21980495124	390.028150754825\\
-4867.21680420105	43.1601372361597\\
-4855.21380345086	105.748276069437\\
-4831.20780195049	18.3303851945902\\
-4819.2048012003	9.32586998383322\\
-4807.20180045012	89.6609800468902\\
-4795.19879969993	44.2519232532845\\
-4783.19579894974	29.4124054754103\\
-4771.19279819955	2.50940196351206\\
-4759.18979744936	44.5642531816593\\
-4747.18679669917	84.7992150241992\\
-4735.18379594899	395.623938162547\\
-4723.1807951988	109.228273951667\\
-4711.17779444861	275.219949236071\\
-4699.17479369842	28.4417359085184\\
-4687.17179294824	596.256670084421\\
-4675.16879219805	127.484753474155\\
-4663.16579144786	26.242482119837\\
-4651.16279069768	222.810612930144\\
-4639.15978994749	148.726880455346\\
-4627.1567891973	6.06540176539309\\
-4615.15378844711	265.443995367329\\
-4603.15078769693	644.063527745001\\
-4591.14778694674	314.94749134601\\
-4579.14478619655	282.716728601797\\
-4567.14178544636	20.3171060886286\\
-4555.13878469617	46.2000336374256\\
-4543.13578394599	39.1706445839603\\
-4531.1327831958	9.78134597656525\\
-4519.12978244561	7.62369095087756\\
-4507.12678169542	111.118591824137\\
-4495.12378094524	53.2321165275032\\
-4483.12078019505	31.2102362921418\\
-4471.11777944486	4.7273604343885\\
-4459.11477869467	40.7757841429322\\
-4447.11177794449	80.5849269010613\\
-4435.1087771943	321.055083257789\\
-4423.10577644411	84.5310454168912\\
-4411.10277569393	193.153580986571\\
-4399.09977494374	13.4986548811859\\
-4387.09677419355	425.571969111472\\
-4375.09377344336	82.523591769912\\
-4363.09077269317	22.9807823982761\\
-4351.08777194299	177.983846092569\\
-4339.0847711928	118.134635324519\\
-4327.08177044261	12.4316848604376\\
-4315.07876969242	233.790993807486\\
-4303.07576894224	495.178108186849\\
-4291.07276819205	306.826303062272\\
-4279.06976744186	235.209089378895\\
-4267.06676669167	8.93266880114429\\
-4255.06376594149	18.6980953505266\\
-4243.0607651913	40.3488076858754\\
-4231.05776444111	5.44098749496943\\
-4219.05476369093	8.22211568605637\\
-4207.05176294074	142.601549755779\\
-4195.04876219055	65.6455638963502\\
-4183.04576144036	34.7283180399026\\
-4171.04276069017	9.42224221279152\\
-4159.03975993999	43.3249292837658\\
-4147.0367591898	84.1608399150897\\
-4135.03375843961	285.3861036471\\
-4123.03075768942	71.1189562246339\\
-4111.02775693923	146.067877931145\\
-4099.02475618905	7.1953335267733\\
-4087.02175543886	329.29619246326\\
-4075.01875468867	57.0446910735069\\
-4063.01575393849	21.4116448275435\\
-4051.0127531883	157.279151978094\\
-4039.00975243811	109.901034612675\\
-4027.00675168792	21.7034823980208\\
-4015.00375093774	231.004439686822\\
-4003.00075018755	412.793205240698\\
-3990.99774943736	325.101191319061\\
-3978.99474868717	221.101800153823\\
-3966.99174793698	2.92289263221811\\
-3954.9887471868	7.18107236652279\\
-3942.98574643661	46.050027883676\\
-3930.98274568642	3.38862303466865\\
-3918.97974493623	11.0567205295599\\
-3906.97674418605	190.548907097736\\
-3894.97374343586	83.6317722571398\\
-3882.97074268567	40.4147384775241\\
-3870.96774193548	17.6320988882801\\
-3858.9647411853	52.1288346840283\\
-3846.96174043511	95.0921441365677\\
-3834.95873968492	273.725121119921\\
-3822.95573893474	63.8994298813724\\
-3810.95273818455	116.587258105001\\
-3798.94973743436	5.50727764700814\\
-3786.94673668417	271.757231369426\\
-3774.94373593398	41.4387449872398\\
-3762.9407351838	20.9083727628189\\
-3750.93773443361	152.250011345562\\
-3738.93473368342	119.11009655356\\
-3726.93173293323	36.0977970704221\\
-3714.92873218305	254.017253664364\\
-3702.92573143286	369.10646607038\\
-3690.92273068267	372.543456140291\\
-3678.91972993248	234.85531942574\\
-3666.9167291823	0.298780843822358\\
-3654.91372843211	7.43765909417665\\
-3642.91072768192	57.2507112484103\\
-3630.90772693173	3.34072286807268\\
-3618.90472618155	17.3691228482934\\
-3606.90172543136	268.209211031706\\
-3594.89872468117	111.388382643921\\
-3582.89572393098	49.6074333099087\\
-3570.8927231808	32.3236847325038\\
-3558.88972243061	70.2748462021107\\
-3546.88672168042	116.124618619055\\
-3534.88372093023	282.843594361355\\
-3522.88072018004	60.9203077975623\\
-3510.87771942986	97.2806455119367\\
-3498.87471867967	7.28896230207912\\
-3486.87171792948	237.992796285156\\
-3474.8687171793	31.8725927567343\\
-3462.86571642911	21.2687649889213\\
-3450.86271567892	162.113985449041\\
-3438.85971492873	150.050577203921\\
-3426.85671417855	60.6027960145881\\
-3414.85371342836	312.367904688549\\
-3402.85071267817	354.326203814128\\
-3390.84771192798	465.052411079322\\
-3378.84471117779	285.126768813529\\
-3366.84171042761	1.12668023414777\\
-3354.83870967742	22.4526065496014\\
-3342.83570892723	77.6449841303765\\
-3330.83270817704	6.444174112309\\
-3318.82970742686	30.7918677737095\\
-3306.82670667667	405.791006405988\\
-3294.82370592648	158.190297886927\\
-3282.82070517629	65.6158990410385\\
-3270.81770442611	60.978269197122\\
-3258.81470367592	106.701740350909\\
-3246.81170292573	155.299529450962\\
-3234.80870217555	318.738462006716\\
-3222.80570142536	61.873400744921\\
-3210.80270067517	85.203199618673\\
-3198.79969992498	13.2228721227622\\
-3186.79669917479	222.568609937941\\
-3174.79369842461	27.6905513076417\\
-3162.79069767442	22.5857771031194\\
-3150.78769692423	193.868590793149\\
-3138.78469617404	221.205781153048\\
-3126.78169542386	107.828950604568\\
-3114.77869467367	438.650002279071\\
-3102.77569392348	370.713281965847\\
-3090.77269317329	647.555461594823\\
-3078.76969242311	405.914394399683\\
-3066.76669167292	8.2324548974403\\
-3054.76369092273	67.4530047084267\\
-3042.76069017254	117.516046144624\\
-3030.75768942236	17.2051283659121\\
-3018.75468867217	62.0546411265472\\
-3006.75168792198	684.65292607533\\
-2994.74868717179	248.352463601993\\
-2982.74568642161	97.7318125503331\\
-2970.74268567142	125.353576475808\\
-2958.73968492123	187.859205837514\\
-2946.73668417104	235.56629287134\\
-2934.73368342085	405.620444800856\\
-2922.73068267067	68.3204419841412\\
-2910.72768192048	81.5640173554586\\
-2898.72468117029	27.1266160630387\\
-2886.72168042011	228.082528688792\\
-2874.71867966992	33.1324025500583\\
-2862.71567891973	25.4495971643784\\
-2850.71267816954	273.987705377942\\
-2838.70967741936	394.826556812027\\
-2826.70667666917	217.374881368585\\
-2814.70367591898	740.470320883476\\
-2802.70067516879	441.168059324333\\
-2790.69767441861	1063.46961039503\\
-2778.69467366842	716.652313946452\\
-2766.69167291823	33.5372796395775\\
-2754.68867216804	200.951339790939\\
-2742.68567141785	211.078296954458\\
-2730.68267066767	53.7755221234057\\
-2718.67966991748	151.178451400976\\
-2706.67666916729	1390.78827736398\\
-2694.6736684171	465.437853116604\\
-2682.67066766692	180.186111719513\\
-2670.66766691673	309.051787975187\\
-2658.66466616654	415.407356857186\\
-2646.66166541636	441.39573666105\\
-2634.65866466617	637.049037498022\\
-2622.65566391598	87.2588551678566\\
-2610.65266316579	101.282809772533\\
-2598.6496624156	65.0407223352798\\
-2586.64666166542	276.9965676163\\
-2574.64366091523	74.4988515758087\\
-2562.64066016504	32.3276433347356\\
-2550.63765941485	519.408338539251\\
-2538.63465866467	953.920785882657\\
-2526.63165791448	570.761605681868\\
-2514.62865716429	1717.81479508401\\
-2502.6256564141	675.0271276999\\
-2490.62265566392	2384.34914629863\\
-2478.61965491373	1821.34973250257\\
-2466.61665416354	144.550377487954\\
-2454.61365341335	750.176356600957\\
-2442.61065266317	535.268671137278\\
-2430.60765191298	228.830079389373\\
-2418.60465116279	539.838966351621\\
-2406.6016504126	4184.48970273351\\
-2394.59864966242	1289.76899291061\\
-2382.59564891223	528.71650452921\\
-2370.59264816204	1170.62925541546\\
-2358.58964741185	1478.2786081453\\
-2346.58664666166	1327.36316387643\\
-2334.58364591148	1636.29600852365\\
-2322.58064516129	162.777573030401\\
-2310.5776444111	293.690000404737\\
-2298.57464366092	249.823026061336\\
-2286.57164291073	541.5296834029\\
-2274.56864216054	431.984797246397\\
-2262.56564141035	62.9595442368674\\
-2250.56264066017	2045.75562263922\\
-2238.55963990998	4862.15031539566\\
-2226.55663915979	3189.35264223319\\
-2214.5536384096	9200.13536970897\\
-2202.55063765942	2249.56777411428\\
-2190.54763690923	13305.9841181874\\
-2178.54463615904	12693.8588329447\\
-2166.54163540885	1530.99408760964\\
-2154.53863465866	7887.99200611938\\
-2142.53563390848	4411.56552991538\\
-2130.53263315829	3358.43540256001\\
-2118.5296324081	7766.74451511817\\
-2106.52663165791	57175.3775135916\\
-2094.52363090773	18315.9096161847\\
-2082.52063015754	10340.7410642349\\
-2070.51762940735	33449.9142367282\\
-2058.51462865717	51973.7558222462\\
-2046.51162790698	52915.3409810299\\
-2034.50862715679	86473.5903500981\\
-2022.5056264066	8612.52774421358\\
-2010.50262565642	148615.148575999\\
-1998.49962490623	1254817.63975307\\
-1986.49662415604	285525.975146389\\
-1974.49362340585	127912.334251226\\
-1962.49062265566	1743.76739141453\\
-1950.48762190547	61562.6681171969\\
-1938.48462115529	96010.71267751\\
-1926.4816204051	39158.4484940062\\
-1914.47861965491	68377.768570207\\
-1902.47561890473	5570.45789373118\\
-1890.47261815454	41067.8834419847\\
-1878.46961740435	33194.4211808259\\
-1866.46661665416	3588.37338143758\\
-1854.46361590398	12485.0185332248\\
-1842.46061515379	3699.8177396908\\
-1830.4576144036	3199.97243889319\\
-1818.45461365341	4770.73566849177\\
-1806.45161290323	20548.4937009295\\
-1794.44861215304	3960.6305102977\\
-1782.44561140285	1991.82007377456\\
-1770.44261065266	4497.83312289733\\
-1758.43960990247	4215.99429036524\\
-1746.43660915229	1998.9207356601\\
-1734.4336084021	1575.05660234622\\
-1722.43060765191	20.7393714093196\\
-1710.42760690173	794.682757362885\\
-1698.42460615154	171.334657930866\\
-1686.42160540135	69.0562870671193\\
-1674.41860465116	757.831004318536\\
-1662.41560390098	27.2274236299987\\
-1650.41260315079	976.645829965037\\
-1638.4096024006	2496.67079364454\\
-1626.40660165041	1402.47045890148\\
-1614.40360090023	3049.02495273082\\
-1602.40060015004	110.890396339234\\
-1590.39759939985	2348.42678262497\\
-1578.39459864966	2750.11420771217\\
-1566.39159789947	395.301041192142\\
-1554.38859714929	1480.39381767807\\
-1542.3855963991	361.965535928666\\
-1518.37959489872	813.123908535361\\
-1506.37659414854	3099.68168022417\\
-1494.37359339835	529.051443797479\\
-1482.37059264816	428.367363817297\\
-1470.36759189798	964.336966315178\\
-1458.36459114779	914.594319204027\\
-1446.3615903976	345.546769658615\\
-1434.35858964741	298.588267863719\\
-1422.35558889723	5.8035422870369\\
-1410.35258814704	320.587110871699\\
-1398.34958739685	24.8124899205931\\
-1386.34658664666	17.0829582879134\\
-1374.34358589647	300.316393418084\\
-1362.34058514629	23.9423452171935\\
-1350.3375843961	300.288682118698\\
-1338.33458364591	779.500879536047\\
-1326.33158289572	444.570462197048\\
-1314.32858214554	980.079494054344\\
-1302.32558139535	5.44429087564687\\
-1290.32258064516	659.33082023822\\
-1278.31957989497	1009.22359774256\\
-1266.31657914479	162.159015451522\\
-1254.3135783946	598.015636600673\\
-1242.31057764441	110.096445003573\\
-1230.30757689423	270.764939769726\\
-1218.30457614404	385.22388141961\\
-1206.30157539385	1214.27974319503\\
-1194.29857464366	164.114877637743\\
-1182.29557389347	235.678807257349\\
-1170.29257314328	461.279459619642\\
-1158.2895723931	427.093200783087\\
-1146.28657164291	119.025807460778\\
-1134.28357089272	138.290496306005\\
-1122.28057014253	11.2298694450726\\
-1110.27756939235	228.28550278468\\
-1098.27456864216	4.23974001639406\\
-1086.27156789197	17.0874979812031\\
-1074.26856714179	199.784577728801\\
-1062.2655663916	27.4211057817752\\
-1050.26256564141	167.653372276009\\
-1038.25956489122	392.674827383647\\
-1026.25656414104	220.402435045799\\
-1014.25356339085	497.045749547717\\
-1002.25056264066	0.878012786623213\\
-990.247561890472	271.17148143796\\
-978.244561140284	554.194336018002\\
-966.241560390099	93.8366683110406\\
-954.238559639911	339.6969235168\\
-942.235558889723	46.113254911095\\
-930.232558139534	188.371962972728\\
-918.229557389346	255.591526812424\\
-906.226556639162	655.081583900136\\
-894.223555888973	64.1485072569037\\
-882.220555138785	186.254412274187\\
-870.217554388597	296.815859874222\\
-858.214553638409	267.186449417016\\
-846.21155288822	52.2511906348925\\
-834.208552138036	102.424022403243\\
-822.205551387848	16.7308816837745\\
-810.202550637659	199.650295376869\\
-798.199549887471	1.78985380864282\\
-786.196549137283	18.4831359179368\\
-774.193548387098	165.538645939171\\
-762.19054763691	32.7547166983199\\
-750.187546886722	126.939461931147\\
-738.184546136534	246.181194053298\\
-726.181545386346	133.370702241357\\
-714.178544636161	314.424938569226\\
-702.175543885973	5.80472516604021\\
-690.172543135785	130.613820873598\\
-678.169542385596	370.910496857288\\
-666.166541635408	63.3943025629697\\
-654.16354088522	226.610585995708\\
-642.160540135035	22.796843128938\\
-630.157539384847	155.838975155157\\
-618.154538634659	202.803976406557\\
-606.151537884471	417.997253429013\\
-594.148537134282	26.5123076237698\\
-582.145536384098	178.882755502976\\
-570.14253563391	225.171390344203\\
-558.139534883721	196.95539413113\\
-546.136534133533	26.1226478222066\\
-534.133533383345	98.6555178220733\\
-522.130532633157	22.1598662352976\\
-510.127531882972	194.828107445533\\
-498.124531132784	6.14694339012918\\
-486.121530382596	18.5269619252877\\
-474.118529632407	156.316007970626\\
-462.115528882219	40.0858761886185\\
-450.112528132035	116.201571200597\\
-438.109527381846	176.398099305682\\
-426.106526631658	90.362131282304\\
-414.10352588147	228.680700953497\\
-402.100525131282	10.8179885039515\\
-390.097524381094	67.0341946430206\\
-378.094523630909	280.515901077299\\
-366.091522880721	46.5636422886566\\
-354.088522130533	166.256278881579\\
-342.085521380344	13.6099356822436\\
-330.082520630156	144.599870445025\\
-318.079519879971	181.390765679172\\
-306.076519129783	297.712098164646\\
-294.073518379595	11.156529450847\\
-282.070517629407	194.286464054714\\
-270.067516879219	191.366747937678\\
-258.064516129034	162.564267526847\\
-246.061515378846	15.0455462973026\\
-234.058514628658	109.044402712267\\
-222.055513878469	28.2202371549694\\
-210.052513128281	205.330473957219\\
-198.049512378093	16.1946231966873\\
-186.046511627908	17.2173009398858\\
-174.04351087772	162.675693136553\\
-162.040510127532	50.4398134930174\\
-150.037509377344	121.171409874319\\
-138.034508627155	139.694801858971\\
-126.031507876971	65.8342414429898\\
-114.028507126783	184.525271035159\\
-102.025506376594	14.6504424810177\\
-90.0225056264062	34.3068346314708\\
-78.019504876218	231.563475592144\\
-66.0165041260298	35.8505219023024\\
-54.0135033758452	130.533205530017\\
-42.010502625657	11.5098583056741\\
-30.0075018754687	146.966845733747\\
-18.0045011252805	178.374334628846\\
-6.00150037509229	231.047790791167\\
6.00150037509229	6.89717904989903\\
18.0045011252805	231.047790791167\\
30.0075018754687	178.374334628846\\
42.010502625657	146.966845733747\\
54.0135033758452	11.5098583056741\\
66.0165041260298	130.533205530017\\
78.019504876218	35.8505219023024\\
90.0225056264062	231.563475592144\\
102.025506376594	34.3068346314708\\
114.028507126783	14.6504424810177\\
126.031507876971	184.525271035159\\
138.034508627155	65.8342414429898\\
150.037509377344	139.694801858971\\
162.040510127532	121.171409874319\\
174.04351087772	50.4398134930174\\
186.046511627908	162.675693136553\\
198.049512378093	17.2173009398858\\
210.052513128281	16.1946231966873\\
222.055513878469	205.330473957219\\
234.058514628658	28.2202371549694\\
246.061515378846	109.044402712267\\
258.064516129034	15.0455462973026\\
270.067516879219	162.564267526847\\
282.070517629407	191.366747937678\\
294.073518379595	194.286464054714\\
306.076519129783	11.156529450847\\
318.079519879971	297.712098164646\\
330.082520630156	181.390765679172\\
342.085521380344	144.599870445025\\
354.088522130533	13.6099356822436\\
366.091522880721	166.256278881579\\
378.094523630909	46.5636422886566\\
390.097524381094	280.515901077299\\
402.100525131282	67.0341946430206\\
414.10352588147	10.8179885039515\\
426.106526631658	228.680700953497\\
438.109527381846	90.362131282304\\
450.112528132035	176.398099305682\\
462.115528882219	116.201571200597\\
474.118529632407	40.0858761886185\\
486.121530382596	156.316007970626\\
498.124531132784	18.5269619252877\\
510.127531882972	6.14694339012918\\
522.130532633157	194.828107445533\\
534.133533383345	22.1598662352976\\
546.136534133533	98.6555178220733\\
558.139534883721	26.1226478222066\\
570.14253563391	196.95539413113\\
582.145536384098	225.171390344203\\
594.148537134282	178.882755502976\\
606.151537884471	26.5123076237698\\
618.154538634659	417.997253429013\\
630.157539384847	202.803976406557\\
642.160540135035	155.838975155157\\
654.16354088522	22.796843128938\\
666.166541635408	226.610585995708\\
678.169542385596	63.3943025629697\\
690.172543135785	370.910496857288\\
702.175543885973	130.613820873598\\
714.178544636161	5.80472516604021\\
726.181545386346	314.424938569226\\
738.184546136534	133.370702241357\\
750.187546886722	246.181194053298\\
762.19054763691	126.939461931147\\
774.193548387098	32.7547166983199\\
786.196549137283	165.538645939171\\
798.199549887471	18.4831359179368\\
810.202550637659	1.78985380864282\\
822.205551387848	199.650295376869\\
834.208552138036	16.7308816837745\\
846.21155288822	102.424022403243\\
858.214553638409	52.2511906348925\\
870.217554388597	267.186449417016\\
882.220555138785	296.815859874222\\
894.223555888973	186.254412274187\\
906.226556639162	64.1485072569037\\
918.229557389346	655.081583900136\\
930.232558139534	255.591526812424\\
942.235558889723	188.371962972728\\
954.238559639911	46.113254911095\\
966.241560390099	339.6969235168\\
978.244561140284	93.8366683110406\\
990.247561890472	554.194336018002\\
1002.25056264066	271.17148143796\\
1014.25356339085	0.878012786623213\\
1026.25656414104	497.045749547717\\
1038.25956489122	220.402435045799\\
1050.26256564141	392.674827383647\\
1062.2655663916	167.653372276009\\
1074.26856714179	27.4211057817752\\
1086.27156789197	199.784577728801\\
1098.27456864216	17.0874979812031\\
1110.27756939235	4.23974001639406\\
1122.28057014253	228.28550278468\\
1134.28357089272	11.2298694450726\\
1146.28657164291	138.290496306005\\
1158.2895723931	119.025807460778\\
1170.29257314328	427.093200783087\\
1182.29557389347	461.279459619642\\
1194.29857464366	235.678807257349\\
1206.30157539385	164.114877637743\\
1218.30457614404	1214.27974319503\\
1230.30757689423	385.22388141961\\
1242.31057764441	270.764939769726\\
1254.3135783946	110.096445003573\\
1266.31657914479	598.015636600673\\
1278.31957989497	162.159015451522\\
1290.32258064516	1009.22359774256\\
1302.32558139535	659.33082023822\\
1314.32858214554	5.44429087564687\\
1326.33158289572	980.079494054344\\
1338.33458364591	444.570462197048\\
1350.3375843961	779.500879536047\\
1362.34058514629	300.288682118698\\
1374.34358589647	23.9423452171935\\
1386.34658664666	300.316393418084\\
1398.34958739685	17.0829582879134\\
1410.35258814704	24.8124899205931\\
1422.35558889723	320.587110871699\\
1434.35858964741	5.8035422870369\\
1446.3615903976	298.588267863719\\
1458.36459114779	345.546769658615\\
1470.36759189798	914.594319204027\\
1482.37059264816	964.336966315178\\
1494.37359339835	428.367363817297\\
1506.37659414854	529.051443797479\\
1518.37959489872	3099.68168022417\\
1530.38259564891	813.123908535361\\
1554.38859714929	361.965535928666\\
1566.39159789947	1480.39381767807\\
1578.39459864966	395.301041192142\\
1590.39759939985	2750.11420771217\\
1602.40060015004	2348.42678262497\\
1614.40360090023	110.890396339234\\
1626.40660165041	3049.02495273082\\
1638.4096024006	1402.47045890148\\
1650.41260315079	2496.67079364454\\
1662.41560390098	976.645829965037\\
1674.41860465116	27.2274236299987\\
1686.42160540135	757.831004318536\\
1698.42460615154	69.0562870671193\\
1710.42760690173	171.334657930866\\
1722.43060765191	794.682757362885\\
1734.4336084021	20.7393714093196\\
1746.43660915229	1575.05660234622\\
1758.43960990247	1998.9207356601\\
1770.44261065266	4215.99429036524\\
1782.44561140285	4497.83312289733\\
1794.44861215304	1991.82007377456\\
1806.45161290323	3960.6305102977\\
1818.45461365341	20548.4937009295\\
1830.4576144036	4770.73566849177\\
1842.46061515379	3199.97243889319\\
1854.46361590398	3699.8177396908\\
1866.46661665416	12485.0185332248\\
1878.46961740435	3588.37338143758\\
1890.47261815454	33194.4211808259\\
1902.47561890473	41067.8834419847\\
1914.47861965491	5570.45789373118\\
1926.4816204051	68377.768570207\\
1938.48462115529	39158.4484940062\\
1950.48762190547	96010.71267751\\
1962.49062265566	61562.6681171969\\
1974.49362340585	1743.76739141453\\
1986.49662415604	127912.334251226\\
1998.49962490623	285525.975146389\\
2010.50262565642	1254817.63975307\\
2022.5056264066	148615.148575999\\
2034.50862715679	8612.52774421358\\
2046.51162790698	86473.5903500981\\
2058.51462865717	52915.3409810299\\
2070.51762940735	51973.7558222462\\
2082.52063015754	33449.9142367282\\
2094.52363090773	10340.7410642349\\
2106.52663165791	18315.9096161847\\
2118.5296324081	57175.3775135916\\
2130.53263315829	7766.74451511817\\
2142.53563390848	3358.43540256001\\
2154.53863465866	4411.56552991538\\
2166.54163540885	7887.99200611938\\
2178.54463615904	1530.99408760964\\
2190.54763690923	12693.8588329447\\
2202.55063765942	13305.9841181874\\
2214.5536384096	2249.56777411428\\
2226.55663915979	9200.13536970897\\
2238.55963990998	3189.35264223319\\
2250.56264066017	4862.15031539566\\
2262.56564141035	2045.75562263922\\
2274.56864216054	62.9595442368674\\
2286.57164291073	431.984797246397\\
2298.57464366092	541.5296834029\\
2310.5776444111	249.823026061336\\
2322.58064516129	293.690000404737\\
2334.58364591148	162.777573030401\\
2346.58664666166	1636.29600852365\\
2358.58964741185	1327.36316387643\\
2370.59264816204	1478.2786081453\\
2382.59564891223	1170.62925541546\\
2394.59864966242	528.71650452921\\
2406.6016504126	1289.76899291061\\
2418.60465116279	4184.48970273351\\
2430.60765191298	539.838966351621\\
2442.61065266317	228.830079389373\\
2454.61365341335	535.268671137278\\
2466.61665416354	750.176356600957\\
2478.61965491373	144.550377487954\\
2490.62265566392	1821.34973250257\\
2502.6256564141	2384.34914629863\\
2514.62865716429	675.0271276999\\
2526.63165791448	1717.81479508401\\
2538.63465866467	570.761605681868\\
2550.63765941485	953.920785882657\\
2562.64066016504	519.408338539251\\
2574.64366091523	32.3276433347356\\
2586.64666166542	74.4988515758087\\
2598.6496624156	276.9965676163\\
2610.65266316579	65.0407223352798\\
2622.65566391598	101.282809772533\\
2634.65866466617	87.2588551678566\\
2646.66166541636	637.049037498022\\
2658.66466616654	441.39573666105\\
2670.66766691673	415.407356857186\\
2682.67066766692	309.051787975187\\
2694.6736684171	180.186111719513\\
2706.67666916729	465.437853116604\\
2718.67966991748	1390.78827736398\\
2730.68267066767	151.178451400976\\
2742.68567141785	53.7755221234057\\
2754.68867216804	211.078296954458\\
2766.69167291823	200.951339790939\\
2778.69467366842	33.5372796395775\\
2790.69767441861	716.652313946452\\
2802.70067516879	1063.46961039503\\
2814.70367591898	441.168059324333\\
2826.70667666917	740.470320883476\\
2838.70967741936	217.374881368585\\
2850.71267816954	394.826556812027\\
2862.71567891973	273.987705377942\\
2874.71867966992	25.4495971643784\\
2886.72168042011	33.1324025500583\\
2898.72468117029	228.082528688792\\
2910.72768192048	27.1266160630387\\
2922.73068267067	81.5640173554586\\
2934.73368342085	68.3204419841412\\
2946.73668417104	405.620444800856\\
2958.73968492123	235.56629287134\\
2970.74268567142	187.859205837514\\
2982.74568642161	125.353576475808\\
2994.74868717179	97.7318125503331\\
3006.75168792198	248.352463601993\\
3018.75468867217	684.65292607533\\
3030.75768942236	62.0546411265472\\
3042.76069017254	17.2051283659121\\
3054.76369092273	117.516046144624\\
3066.76669167292	67.4530047084267\\
3078.76969242311	8.2324548974403\\
3090.77269317329	405.914394399683\\
3102.77569392348	647.555461594823\\
3114.77869467367	370.713281965847\\
3126.78169542386	438.650002279071\\
3138.78469617404	107.828950604568\\
3150.78769692423	221.205781153048\\
3162.79069767442	193.868590793149\\
3174.79369842461	22.5857771031194\\
3186.79669917479	27.6905513076417\\
3198.79969992498	222.568609937941\\
3210.80270067517	13.2228721227622\\
3222.80570142536	85.203199618673\\
3234.80870217555	61.873400744921\\
3246.81170292573	318.738462006716\\
3258.81470367592	155.299529450962\\
3270.81770442611	106.701740350909\\
3282.82070517629	60.978269197122\\
3294.82370592648	65.6158990410385\\
3306.82670667667	158.190297886927\\
3318.82970742686	405.791006405988\\
3330.83270817704	30.7918677737095\\
3342.83570892723	6.444174112309\\
3354.83870967742	77.6449841303765\\
3366.84171042761	22.4526065496014\\
3378.84471117779	1.12668023414777\\
3390.84771192798	285.126768813529\\
3402.85071267817	465.052411079322\\
3414.85371342836	354.326203814128\\
3426.85671417855	312.367904688549\\
3438.85971492873	60.6027960145881\\
3450.86271567892	150.050577203921\\
3462.86571642911	162.113985449041\\
3474.8687171793	21.2687649889213\\
3486.87171792948	31.8725927567343\\
3498.87471867967	237.992796285156\\
3510.87771942986	7.28896230207912\\
3522.88072018004	97.2806455119367\\
3534.88372093023	60.9203077975623\\
3546.88672168042	282.843594361355\\
3558.88972243061	116.124618619055\\
3570.8927231808	70.2748462021107\\
3582.89572393098	32.3236847325038\\
3594.89872468117	49.6074333099087\\
3606.90172543136	111.388382643921\\
3618.90472618155	268.209211031706\\
3630.90772693173	17.3691228482934\\
3642.91072768192	3.34072286807268\\
3654.91372843211	57.2507112484103\\
3666.9167291823	7.43765909417665\\
3678.91972993248	0.298780843822358\\
3690.92273068267	234.85531942574\\
3702.92573143286	372.543456140291\\
3714.92873218305	369.10646607038\\
3726.93173293323	254.017253664364\\
3738.93473368342	36.0977970704221\\
3750.93773443361	119.11009655356\\
3762.9407351838	152.250011345562\\
3774.94373593398	20.9083727628189\\
3786.94673668417	41.4387449872398\\
3798.94973743436	271.757231369426\\
3810.95273818455	5.50727764700814\\
3822.95573893474	116.587258105001\\
3834.95873968492	63.8994298813724\\
3846.96174043511	273.725121119921\\
3858.9647411853	95.0921441365677\\
3870.96774193548	52.1288346840283\\
3882.97074268567	17.6320988882801\\
3894.97374343586	40.4147384775241\\
3906.97674418605	83.6317722571398\\
3918.97974493623	190.548907097736\\
3930.98274568642	11.0567205295599\\
3942.98574643661	3.38862303466865\\
3954.9887471868	46.050027883676\\
3966.99174793698	7.18107236652279\\
3978.99474868717	2.92289263221811\\
3990.99774943736	221.101800153823\\
4003.00075018755	325.101191319061\\
4015.00375093774	412.793205240698\\
4027.00675168792	231.004439686822\\
4039.00975243811	21.7034823980208\\
4051.0127531883	109.901034612675\\
4063.01575393849	157.279151978094\\
4075.01875468867	21.4116448275435\\
4087.02175543886	57.0446910735069\\
4099.02475618905	329.29619246326\\
4111.02775693923	7.1953335267733\\
4123.03075768942	146.067877931145\\
4135.03375843961	71.1189562246339\\
4147.0367591898	285.3861036471\\
4159.03975993999	84.1608399150897\\
4171.04276069017	43.3249292837658\\
4183.04576144036	9.42224221279152\\
4195.04876219055	34.7283180399026\\
4207.05176294074	65.6455638963502\\
4219.05476369093	142.601549755779\\
4231.05776444111	8.22211568605637\\
4243.0607651913	5.44098749496943\\
4255.06376594149	40.3488076858754\\
4267.06676669167	18.6980953505266\\
4279.06976744186	8.93266880114429\\
4291.07276819205	235.209089378895\\
4303.07576894224	306.826303062272\\
4315.07876969242	495.178108186849\\
4327.08177044261	233.790993807486\\
4339.0847711928	12.4316848604376\\
4351.08777194299	118.134635324519\\
4363.09077269317	177.983846092569\\
4375.09377344336	22.9807823982761\\
4387.09677419355	82.523591769912\\
4399.09977494374	425.571969111472\\
4411.10277569393	13.4986548811859\\
4423.10577644411	193.153580986571\\
4435.1087771943	84.5310454168912\\
4447.11177794449	321.055083257789\\
4459.11477869467	80.5849269010613\\
4471.11777944486	40.7757841429322\\
4483.12078019505	4.7273604343885\\
4495.12378094524	31.2102362921418\\
4507.12678169542	53.2321165275032\\
4519.12978244561	111.118591824137\\
4531.1327831958	7.62369095087756\\
4543.13578394599	9.78134597656525\\
4555.13878469617	39.1706445839603\\
4567.14178544636	46.2000336374256\\
4579.14478619655	20.3171060886286\\
4591.14778694674	282.716728601797\\
4603.15078769693	314.94749134601\\
4615.15378844711	644.063527745001\\
4627.1567891973	265.443995367329\\
4639.15978994749	6.06540176539309\\
4651.16279069768	148.726880455346\\
4663.16579144786	222.810612930144\\
4675.16879219805	26.242482119837\\
4687.17179294824	127.484753474155\\
4699.17479369842	596.256670084421\\
4711.17779444861	28.4417359085184\\
4723.1807951988	275.219949236071\\
4735.18379594899	109.228273951667\\
4747.18679669917	395.623938162547\\
4759.18979744936	84.7992150241992\\
4771.19279819955	44.5642531816593\\
4783.19579894974	2.50940196351206\\
4795.19879969993	29.4124054754103\\
4807.20180045012	44.2519232532845\\
4819.2048012003	89.6609800468902\\
4831.20780195049	9.32586998383322\\
4843.21080270068	18.3303851945902\\
4867.21680420105	105.748276069437\\
4879.21980495124	43.1601372361597\\
4891.22280570143	390.028150754825\\
4903.22580645161	359.377605351047\\
4915.2288072018	932.886872965243\\
4927.23180795199	346.610111878938\\
4939.23480870217	1.78030453779239\\
4951.23780945236	222.548158694877\\
4963.24081020255	317.42168243976\\
4975.24381095274	32.8792737256408\\
4987.24681170293	218.061307895442\\
4999.24981245312	937.440093599289\\
5011.2528132033	64.2437401153623\\
5023.25581395349	439.374785267883\\
5035.25881470368	159.288669263072\\
5047.26181545386	553.714306511756\\
5059.26481620405	102.002938579236\\
5071.26781695424	59.1769661352503\\
5083.27081770443	3.40034296230673\\
5095.27381845461	29.6572854764467\\
5107.2768192048	37.5145022447793\\
5119.27981995499	75.1285913665575\\
5131.28282070518	15.3488699979314\\
5143.28582145536	37.5779951178741\\
5155.28882220555	60.7653361958045\\
5167.29182295574	251.477380241268\\
5179.29482370593	97.3849344692588\\
5191.29782445612	647.20516481097\\
5203.3008252063	480.789582830957\\
5215.30382595649	1595.88721490466\\
5227.30682670668	551.751307932153\\
5239.30982745687	0.571012411150904\\
5251.31282820705	415.155179805867\\
5263.31582895724	549.745281367917\\
5275.31882970743	48.3472350983942\\
5287.32183045761	445.496171170098\\
5299.3248312078	1783.8361042533\\
5311.32783195799	167.157889638012\\
5323.33083270818	853.825544350504\\
5335.33383345836	285.05332725338\\
5347.33683420855	958.461204135293\\
5359.33983495874	153.972264545076\\
5371.34283570893	104.790005731287\\
5383.34583645912	13.1003366141904\\
5395.3488372093	34.2429531166745\\
5407.35183795949	32.311154931515\\
5419.35483870968	67.3727150368459\\
5431.35783945987	35.6944732237895\\
5443.36084021005	96.0306718929393\\
5455.36384096024	120.458286015263\\
5467.36684171043	734.91384937573\\
5479.36984246062	276.170428037186\\
5491.3728432108	1473.49532942079\\
5503.37584396099	868.179112959118\\
5515.37884471118	3748.89001061307\\
5527.38184546136	1250.20445116463\\
5539.38484621155	13.1515772019442\\
5551.38784696174	1129.49069981747\\
5563.39084771193	1385.19501432\\
5575.39384846212	100.380704360802\\
5587.3968492123	1329.53676368348\\
5599.39984996249	5063.80453723383\\
5611.40285071268	629.711669467886\\
5623.40585146287	2547.77266456231\\
5635.40885221306	797.996690971374\\
5647.41185296324	2646.81351770095\\
5659.41485371343	385.199061033971\\
5671.41785446362	326.993975071962\\
5683.4208552138	82.6850241234607\\
5695.42385596399	59.9658760808478\\
5707.42685671418	28.1967485576963\\
5719.42985746437	81.1541716822097\\
5731.43285821455	162.998155731759\\
5743.43585896474	464.858963065025\\
5755.43885971493	530.334665564024\\
5767.44186046512	4298.17889033651\\
5779.44486121531	1649.40836802786\\
5791.44786196549	7803.61499321949\\
5803.45086271568	3741.91621184304\\
5815.45386346587	22447.0608533422\\
5827.45686421606	7851.48577758974\\
5839.45986496624	275.012592269145\\
5851.46286571643	9760.65488712021\\
5863.46586646662	12029.7606974502\\
5875.4688672168	753.597870039406\\
5887.47186796699	16512.6073475319\\
5899.47486871718	68356.1739508387\\
5911.47786946737	12652.7372086062\\
5923.48087021755	50157.3814193085\\
5935.48387096774	18305.8160113776\\
5947.48687171793	78443.7069433976\\
5959.48987246812	15144.2404333163\\
5971.49287321831	26113.7888288079\\
5983.49587396849	22184.6116983552\\
5995.49887471868	17176.5330496123\\
6007.50187546887	415989.961795309\\
6019.50487621906	16674.2439913723\\
6031.50787696924	22876.5342829717\\
6043.51087771943	26019.6575621363\\
6055.51387846962	14585.0588117928\\
6067.51687921981	77032.9851270426\\
6079.51987996999	17978.7406016359\\
6091.52288072018	48571.0873417871\\
6103.52588147037	11821.4873418129\\
6115.52888222055	64999.2238904828\\
6127.53188297074	15996.5116292317\\
6139.53488372093	796.893726228774\\
6151.53788447112	11732.5242482912\\
6163.54088522131	9514.45289364669\\
6175.54388597149	312.237629773571\\
6187.54688672168	7432.74551728567\\
6199.54988747187	20513.3744100669\\
6211.55288822206	3226.69454055639\\
6223.55588897225	7158.44455127447\\
6235.55888972243	1569.81740206505\\
6247.56189047262	4036.53888926657\\
6259.56489122281	434.166024104446\\
6271.56789197299	462.508360632168\\
6283.57089272318	222.540342088662\\
6295.57389347337	22.6476824161627\\
6307.57689422356	22.1599325585346\\
6319.57989497374	60.4649747971562\\
6331.58289572393	133.095061955608\\
6343.58589647412	306.712997646573\\
6355.58889722431	293.54955723265\\
6367.59189797449	2344.03546410861\\
6379.59489872468	719.368659842024\\
6391.59789947487	2211.10619586468\\
6403.60090022506	479.886585093787\\
6415.60390097525	4189.14731736652\\
6427.60690172543	1208.12895325619\\
6439.60990247562	110.412095015824\\
6451.61290322581	1272.67283838944\\
6463.615903976	1048.599924978\\
6475.61890472618	25.3683764137763\\
6487.62190547637	1057.8994667953\\
6499.62490622656	2969.69398580231\\
6511.62790697674	596.869212211002\\
6523.63090772693	1185.8962035196\\
6535.63390847712	251.661186231798\\
6547.63690922731	655.389483207184\\
6559.63990997749	77.000467591446\\
6571.64291072768	99.1150296418652\\
6583.64591147787	63.2318547456983\\
6595.64891222806	9.94206964271097\\
6607.65191297825	19.8995238135775\\
6619.65491372843	38.4105689748921\\
6631.65791447862	37.8422565650849\\
6643.66091522881	87.0492014882281\\
6655.663915979	92.6259642163298\\
6667.66691672918	751.38903236234\\
6679.66991747937	232.399617520824\\
6691.67291822956	668.27118324171\\
6703.67591897974	105.79713469105\\
6715.67891972993	1283.74716156373\\
6727.68192048012	387.352371837069\\
6739.68492123031	51.3997003378818\\
6751.6879219805	468.965478763273\\
6763.69092273068	363.936116724089\\
6775.69392348087	5.39976915483158\\
6787.69692423106	405.335293887113\\
6799.69992498125	1091.51647749567\\
6811.70292573144	260.691840931838\\
6823.70592648162	449.148500093267\\
6835.70892723181	87.8858550054902\\
6847.711927982	221.890348982356\\
6859.71492873218	29.9186055719936\\
6871.71792948237	42.0321454279244\\
6883.72093023256	32.2581745386869\\
6895.72393098275	8.49175002988929\\
6907.72693173293	17.9977212426819\\
6919.72993248312	31.3703001243038\\
6931.73293323331	18.0869678609387\\
6943.7359339835	40.5821835150989\\
6955.73893473368	48.0813180803669\\
6967.74193548387	377.828936535114\\
6979.74493623406	115.942252473006\\
6991.74793698425	311.525671695138\\
7003.75093773444	34.8251126629626\\
7015.75393848462	586.012103557058\\
7027.75693923481	184.674343396568\\
7039.759939985	32.3616005721091\\
7051.76294073519	246.719958311213\\
7063.76594148537	180.010543674522\\
7075.76894223556	1.36472183312161\\
7087.77194298575	213.878276931909\\
7099.77494373593	547.762032909865\\
7111.77794448612	151.737880944277\\
7123.78094523631	227.86710267593\\
7135.7839459865	40.5907020377518\\
7147.78694673668	98.1988035492384\\
7159.78994748687	16.380584967287\\
7171.79294823706	23.2972276519095\\
7183.79594898725	20.4187015480953\\
7195.79894973744	8.01538849454986\\
7207.80195048762	16.3807055137734\\
7219.80495123781	27.1030541762658\\
7231.807951988	10.7401511306422\\
7243.81095273819	23.3805814114456\\
7255.81395348837	30.9504793225378\\
7267.81695423856	231.046376578061\\
7279.81995498875	70.1168312708647\\
7291.82295573893	177.335405817937\\
7303.82595648912	13.9501435568715\\
7315.82895723931	322.767797969255\\
7327.8319579895	106.871053461318\\
7339.83495873968	23.2485723107867\\
7351.83795948987	154.206510226298\\
7363.84096024006	106.445423555861\\
7375.84396099025	0.297441738273464\\
7387.84696174044	132.514276772867\\
7399.84996249063	322.061111989046\\
7411.85296324081	101.903430849355\\
7423.855963991	134.569271980964\\
7435.85896474119	21.6499529536902\\
7447.86196549137	49.8216808973087\\
7459.86496624156	10.9531565282398\\
7471.86796699175	14.9104660754109\\
7483.87096774194	14.4253870467553\\
7495.87396849212	7.68896490580707\\
7507.87696924231	14.9930648555116\\
7519.8799699925	23.9521408346512\\
7531.88297074269	7.18444690366655\\
7543.88597149287	15.1546749918648\\
7555.88897224306	22.3971346248773\\
7567.89197299325	157.543662406582\\
7579.89497374344	47.2540623662791\\
7591.89797449363	113.354420909654\\
7603.90097524381	6.53544848689516\\
7615.903975994	198.451764781738\\
7627.90697674419	69.4996287096762\\
7639.90997749438	17.9524195972579\\
7651.91297824456	106.677134731817\\
7663.91597899475	70.2225415119331\\
7675.91897974494	0.0328224293021715\\
7687.92198049512	90.4581744436843\\
7699.92498124531	208.516964309664\\
7711.9279819955	74.4991671679466\\
7723.93098274569	87.2824739164312\\
7735.93398349587	12.5925356560003\\
7747.93698424606	27.3769218186231\\
7759.93998499625	8.29571946277411\\
7771.94298574644	10.4339720697103\\
7783.94598649663	10.884415970017\\
7795.94898724681	7.39111856012421\\
7807.951987997	13.7924837257285\\
7819.95498874719	21.4410478605861\\
7831.95798949738	5.18178511652411\\
7843.96099024756	10.5907921933888\\
7855.96399099775	17.4103332341565\\
7867.96699174794	115.108490866786\\
7879.96999249813	34.1298249201248\\
7891.97299324831	78.1752096653655\\
7903.9759939985	3.68985096078387\\
7915.97899474869	131.191678539495\\
7927.98199549887	48.8876679369547\\
7939.98499624906	14.5016076637538\\
7951.98799699925	78.8821044761983\\
7963.99099774944	49.9162022808783\\
7975.99399849963	0.0216567802662538\\
7987.99699924981	65.894537468142\\
8000	144.068998096081\\
8012.00300075019	57.5905491618052\\
8024.00600150038	60.3529944237652\\
8036.00900225056	7.75092094459824\\
8048.01200300075	15.7991759220366\\
8060.01500375094	6.79656319247069\\
8072.01800450113	7.75574746305156\\
8084.02100525131	8.57894994631742\\
8096.0240060015	7.10294087960241\\
8108.02700675169	12.7461518464611\\
8120.03000750187	19.3650968011518\\
8132.03300825206	3.93702432922546\\
8144.03600900225	7.79983922040432\\
8156.03900975244	14.188327198119\\
8168.04201050263	88.2044335410151\\
8180.04501125282	25.863406789248\\
8192.048012003	56.8889268296372\\
8204.05101275319	2.60878459013396\\
8216.05401350338	91.3072871775631\\
8228.05701425357	36.3900291875932\\
8240.06001500375	12.0791192536396\\
8252.06301575394	61.132889436927\\
8264.06601650413	37.4673901792955\\
8276.06901725431	0.0915641728607937\\
8288.0720180045	50.2933618797993\\
8300.07501875469	104.31692084451\\
8312.07801950488	46.3034648965835\\
8324.08102025506	43.7271455262028\\
8336.08402100525	4.9584691305036\\
8348.08702175544	9.39079092304157\\
8360.09002250563	5.85354238367733\\
8372.09302325582	6.01985943678356\\
8384.096024006	6.97513630342797\\
8396.09902475619	6.82382628189279\\
8408.10202550638	11.828281372926\\
8420.10502625657	17.6120860917267\\
8432.10802700675	3.10803834947061\\
8444.11102775694	5.9716576883716\\
8456.11402850713	11.9497866104598\\
8468.11702925732	69.979141922505\\
8480.1200300075	20.3020885732209\\
8492.12303075769	43.0898193674297\\
8504.12603150788	2.25704968844555\\
8516.12903225806	66.0386347586357\\
8528.13203300825	28.2683511416087\\
8540.13503375844	10.2877817435924\\
8552.13803450863	49.0537516092493\\
8564.14103525882	29.3165710685871\\
8576.14403600901	0.182001174205365\\
8588.14703675919	39.7595230613699\\
8600.15003750938	78.2582473798779\\
8612.15303825957	38.3263022554407\\
8624.15603900975	32.8345871572342\\
8636.15903975994	3.2575512350301\\
8648.16204051013	5.66849540670326\\
8672.1680420105	4.82625802690968\\
8684.17104276069	5.80482024233565\\
8696.17404351088	6.55584542625844\\
8708.17704426107	11.0183556030825\\
8720.18004501125	16.1106747982632\\
8732.18304576144	2.52708144027304\\
8744.18604651163	4.71129168349648\\
8756.18904726182	10.3089725022209\\
8768.19204801201	57.0097140789033\\
8780.19504876219	16.3720982505147\\
8792.19804951238	33.6661921503752\\
8804.20105026257	2.21233596565326\\
8816.20405101276	49.2077277856975\\
8828.20705176294	22.7024881686582\\
8840.21005251313	8.91175931768331\\
8852.21305326332	40.4279376315195\\
8864.2160540135	23.7025939237955\\
8876.21905476369	0.270873394820217\\
8888.22205551388	32.3074932752214\\
8900.22505626407	60.3615892463681\\
8912.22805701425	32.4395796717305\\
8924.23105776444	25.3653730844459\\
8936.23405851463	2.17876934175199\\
8948.23705926482	3.43498862993601\\
8960.24006001501	4.7319103630701\\
8972.24306076519	3.96755526424877\\
8984.24606151538	4.9195344721075\\
8996.24906226557	6.30068651914661\\
9008.25206301576	10.2998713476577\\
9020.25506376594	14.8112196924436\\
9032.25806451613	2.10369633584456\\
9044.26106526632	3.80732046366561\\
9056.26406601651	9.05628889321871\\
9068.26706676669	47.4228918149261\\
9080.27006751688	13.4871655992304\\
9092.27306826707	26.9629704001974\\
9104.27606901725	2.29126819118466\\
9116.27906976744	37.5482172597433\\
9128.28207051763	18.724290620904\\
9140.28507126782	7.82370250061856\\
9152.28807201801	34.0319885126066\\
9164.29107276819	19.676869701407\\
9176.29407351838	0.350549653407224\\
9188.29707426857	26.838217945205\\
9200.30007501876	47.6095416005692\\
9212.30307576894	27.9464210904662\\
9224.30607651913	20.0550540079733\\
9236.30907726932	1.47346766786875\\
9248.31207801951	2.06774959624283\\
9260.31507876969	4.36380706887022\\
9272.31807951988	3.32739540262402\\
9284.32108027007	4.23082002462319\\
9296.32408102026	6.05925297939305\\
9308.32708177044	9.65942108587862\\
9320.33008252063	13.677170915697\\
9332.33308327082	1.78543637092468\\
9344.33608402101	3.13819454864333\\
9356.33908477119	8.06909257614698\\
9368.34208552138	40.1192204111278\\
9380.34508627157	11.3041955299191\\
9392.34808702176	22.0360058391203\\
9404.35108777195	2.41163017823292\\
9416.35408852213	29.2117188672116\\
9428.35708927232	15.7821196378409\\
9440.36009002251	6.94357493221771\\
9452.36309077269	29.1442056515819\\
9464.36609152288	16.6935243003375\\
9476.36909227307	0.419093395409326\\
9488.37209302326	22.703014730311\\
9500.37509377344	38.247877765742\\
9512.37809452363	24.4229814762273\\
9524.38109527382	16.1669516312161\\
9536.38409602401	1.00172478010146\\
9548.3870967742	1.22366387436733\\
9560.39009752438	4.06637285585873\\
9572.39309827457	2.83628902382305\\
9584.39609902476	3.68286996610972\\
9596.39909977495	5.83180583561867\\
9608.40210052513	9.08601321246373\\
9620.40510127532	12.6806034160122\\
9632.40810202551	1.54010825657324\\
9644.4111027757	2.62998390286609\\
9656.41410352588	7.27125149286106\\
9668.41710427607	34.4166881638873\\
9680.42010502626	9.6110566691565\\
9692.42310577644	18.3158088568301\\
9704.42610652663	2.53655651269392\\
9716.42910727682	23.0940479968427\\
9728.43210802701	13.5438171187554\\
9740.4351087772	6.21850475779677\\
9752.43810952738	25.3155933912502\\
9764.44111027757	14.4212843480583\\
9776.44411102776	0.476805414846802\\
9788.44711177795	19.4988025105329\\
9800.45011252813	31.203251527018\\
9812.45311327832	21.5982582474245\\
9824.45611402851	13.2498061616397\\
9836.4591147787	0.680911561520054\\
9848.46211552888	0.704677613347027\\
9860.46511627907	3.8182467556474\\
9872.46811702926	2.45058049762819\\
9884.47111777945	3.23882965270338\\
9896.47411852963	5.61816804936317\\
9908.47711927982	8.57056185167392\\
9920.48012003001	11.7996257790858\\
9932.4831207802	1.34702025871412\\
9944.48612153038	2.23562489894563\\
9956.48912228057	6.61316901379724\\
9968.49212303076	29.8728076400726\\
9980.49512378095	8.27060308188116\\
9992.49812453114	15.4428694006981\\
10004.5011252813	2.65027647975809\\
10016.5041260315	18.5057669730928\\
10028.5071267817	11.8000920810526\\
10040.5101275319	5.61213079211737\\
10052.5131282821	22.2543684662373\\
10064.5161290323	12.6501270749644\\
10076.5191297824	0.524795995533192\\
10088.5221305326	16.964312171042\\
10100.5251312828	25.7907590256709\\
10112.528132033	19.2917217905856\\
10124.5311327832	11.0155603428263\\
10136.5341335334	0.460271989309994\\
10148.5371342836	0.392138490543219\\
10160.5401350338	3.60615935213191\\
10172.5431357839	2.14167116343142\\
10184.5461365341	2.87343506341704\\
10196.5491372843	5.41788953024198\\
10208.5521380345	8.10549923593695\\
10220.5551387847	11.0167490952072\\
10232.5581395349	1.19235985958183\\
10244.5611402851	1.92400753187111\\
10256.5641410353	6.06121201644021\\
10268.5671417854	26.1896325526302\\
10280.5701425356	7.19082833554948\\
10292.5731432858	13.1814623237829\\
10304.576144036	2.74700469102109\\
10316.5791447862	15.0002574755748\\
10328.5821455364	10.4139565328011\\
10340.5851462866	5.09860556584151\\
10352.5881470368	19.7638917057415\\
10364.5911477869	11.2419292417462\\
10376.5941485371	0.564384887926606\\
10388.5971492873	14.9241121828621\\
10400.6001500375	21.5577377348723\\
10412.6031507877	17.3789230249586\\
10424.6061515379	9.27387750082935\\
10436.6091522881	0.307593379969194\\
10448.6121530383	0.212780047810677\\
10460.6151537884	3.42148531897376\\
10472.6181545386	1.89015550985819\\
10484.6211552888	2.56882277030477\\
10496.624156039	5.23036342785146\\
10508.6271567892	7.68447789069201\\
10520.6301575394	10.317795966885\\
10532.6331582896	1.06660996274169\\
10544.6361590398	1.67390687228021\\
10556.6391597899	5.59178749759313\\
10568.6421605401	23.1601795284139\\
10580.6451612903	6.30805749813235\\
10592.6481620405	11.3719871631561\\
10604.6511627907	2.82565745466678\\
10616.6541635409	12.2790478516458\\
10628.6571642911	9.29279457961604\\
10640.6601650413	4.65904095797317\\
10652.6631657915	17.7074998277444\\
10664.6661665416	10.1029714599025\\
10676.6691672918	0.596854452607379\\
10688.672168042	13.2568269799939\\
10700.6751687922	18.1959570033971\\
10712.6781695424	15.7715474753889\\
10724.6811702926	7.8952458526125\\
10736.6841710428	0.201863226444501\\
10748.6871717929	0.120229315039123\\
10760.6901725431	3.25838309740241\\
10772.6931732933	1.68246604573315\\
10784.6961740435	2.31202829344515\\
10796.6991747937	5.05490366544579\\
10808.7021755439	7.30213944413258\\
10820.7051762941	9.69113725420593\\
10832.7081770443	0.963034542806348\\
10844.7111777945	1.47044686354106\\
10856.7141785446	5.18786054601216\\
10868.7171792948	20.6367793026915\\
10880.720180045	5.57703743115809\\
10892.7231807952	9.90337133619218\\
10904.7261815454	2.88727775974249\\
10916.7291822956	10.137274641161\\
10928.7321830458	8.37221529822094\\
10940.735183796	4.27930837672672\\
10952.7381845461	15.9876712641712\\
10964.7411852963	9.16792580333784\\
10976.7441860465	0.623357318987673\\
10988.7471867967	11.8762912360118\\
11000.7501875469	15.4900782489517\\
11012.7531882971	14.4053445580094\\
11024.7561890473	6.78928667973493\\
11036.7591897975	0.129052476231427\\
11048.7621905476	0.0845384877165072\\
11060.7651912978	3.1127460917687\\
11072.768192048	1.50886971488601\\
11084.7711927982	2.09343098082249\\
11096.7741935484	4.89079502784522\\
11108.7771942986	6.95393344141011\\
11120.7801950488	9.12713975815642\\
11132.783195799	0.876754732505848\\
11144.7861965491	1.30295636330265\\
11156.7891972993	4.83682174644358\\
11168.7921980495	18.511634343443\\
11180.7951987997	4.96486183775285\\
11192.7981995499	8.69644967401912\\
11204.8012003001	2.93376695959668\\
11216.8042010503	8.43103179034572\\
11228.8072018005	7.60634700209779\\
11240.8102025506	3.94862626141677\\
11252.8132033008	14.5332046343629\\
11264.816204051	8.39017071725602\\
11276.8192048012	0.644891815825951\\
11288.8222055514	10.7199521068494\\
11300.8252063016	13.2862651743762\\
11312.8282070518	13.2325516710512\\
11324.831207802	5.89152819925784\\
11336.8342085521	0.0796067603261735\\
11348.8372093023	0.0860516391515543\\
11360.8402100525	2.98158812521018\\
11372.8432108027	1.36222652792498\\
11384.8462115529	1.90575335280288\\
11396.8492123031	4.73732467082627\\
11408.8522130533	6.63597408771547\\
11420.8552138035	8.61775604079786\\
11432.8582145536	0.804165382226572\\
11444.8612153038	1.16362421789612\\
11456.864216054	4.52913386839908\\
11468.8672168042	16.7044652690978\\
11480.8702175544	4.44711993379273\\
11492.8732183046	7.69360870715812\\
11504.8762190548	2.96726414552185\\
11516.879219805	7.05716315575\\
11528.8822205551	6.96180290746126\\
11540.8852213053	3.65862398116905\\
11552.8882220555	13.2910640184864\\
11564.8912228057	7.73573745597702\\
11576.8942235559	0.662306208264595\\
11588.8972243061	9.74149880356196\\
11600.9002250563	11.4724445641577\\
11612.9032258065	12.2169869523593\\
11624.9062265566	5.15508267373656\\
11636.9092273068	0.0469074006181004\\
11648.912228057	0.111703684233571\\
11660.9152288072	2.86267068608478\\
11672.9182295574	1.2371953852836\\
11684.9212303076	1.74339774810104\\
11696.9242310578	4.59380131041775\\
11708.927231808	6.34492604898587\\
11720.9302325581	8.15621351271775\\
11732.9332333083	0.742555453870472\\
11744.9362340585	1.04663022717722\\
11756.9392348087	4.25744602068843\\
11768.9422355589	15.1544276843479\\
11780.9452363091	4.00538240670633\\
11792.9482370593	6.85214016645221\\
11804.9512378095	2.98985223540954\\
11816.9542385596	5.94038832983916\\
11828.9572393098	6.41381676624935\\
11840.96024006	3.40270418990488\\
11852.9632408102	12.2210409231355\\
11864.9662415604	7.1794160032852\\
11876.9692423106	0.676314006596958\\
11888.9722430608	8.9060451598027\\
11900.975243811	9.96554108423102\\
11912.9782445611	11.3307817368783\\
11924.9812453113	4.54526235348326\\
11936.9842460615	0.0263018117384255\\
11948.9872468117	0.152729332906147\\
11960.9902475619	2.75426961007431\\
11972.9932483121	1.12971148540392\\
11984.9962490623	1.60199500274712\\
11996.9992498125	4.45956633744998\\
12009.0022505626	6.07791271792038\\
12021.0052513128	7.73677483879039\\
12033.008252063	0.689854498231162\\
12045.0112528132	0.947568845452151\\
12057.0142535634	4.01599772195875\\
12069.0172543136	13.814685208943\\
12081.0202550638	3.62551764579351\\
12093.023255814	6.13986227459951\\
12105.0262565641	3.0034300333604\\
12117.0292573143	5.02488851633062\\
12129.0322580645	5.94370257557944\\
12141.0352588147	3.17559792118897\\
12153.0382595649	11.2921702280327\\
12165.0412603151	6.70218567470684\\
12177.0442610653	0.687512325073661\\
12189.0472618155	8.18690185188349\\
12201.0502625656	8.70301564870925\\
12213.0532633158	10.5521512363236\\
12225.056264066	4.03601058559293\\
12237.0592648162	0.0144781599566918\\
12249.0622655664	0.203213043158643\\
12261.0652663166	2.65502474763927\\
12273.0682670668	1.03663376623378\\
12285.071267817	1.47809030139939\\
12297.0742685671	4.33399968685032\\
12309.0772693173	5.83244199936263\\
12321.0802700675	7.35455075544326\\
12333.0832708177	0.644459392587098\\
12345.0862715679	0.863058142512662\\
12357.0892723181	3.80020815107782\\
12369.0922730683	12.6486818472051\\
12381.0952738185	3.29653812464909\\
12393.0982745686	5.53216740777129\\
12405.1012753188	3.00966807950791\\
12417.104276069	4.26868139190442\\
12429.1072768192	5.53714533196242\\
12441.1102775694	2.97304752068495\\
12453.1132783196	10.4802689264567\\
12465.1162790698	6.28948098034324\\
12477.11927982	0.69639979511695\\
12489.1222805701	7.56336553890934\\
12501.1252813203	7.63713190541408\\
12513.1282820705	9.86384031725303\\
12525.1312828207	3.60748426904432\\
12537.1342835709	0.00905361268163909\\
12549.1372843211	0.2591532999421\\
12561.1402850713	2.56384045300655\\
12573.1432858215	0.95550171789125\\
12585.1462865716	1.36892004968128\\
12597.1492873218	4.21652233658715\\
12609.152288072	5.60634586801406\\
12621.1552888222	7.00535212792447\\
12633.1582895724	0.605113482847921\\
12645.1612903226	0.790468996738265\\
12657.1642910728	3.6063868250355\\
12669.167291823	11.6275290551505\\
12681.1702925731	3.00979383882218\\
12693.1732933233	5.00998981419989\\
12705.1762940735	3.01000571006751\\
12717.1792948237	3.63978363025611\\
12729.1822955739	5.18302667767846\\
12741.1852963241	2.79157639778293\\
12753.1882970743	9.76621150608567\\
12765.1912978245	5.929997232941\\
12777.1942985746	0.703392657996505\\
12789.1972993248	7.01917438497811\\
12801.200300075	6.73099410882311\\
12813.2033008252	9.25201878047832\\
12825.2063015754	3.24438411472487\\
12837.2093023256	0.00829818716607979\\
12849.2123030758	0.317848462668484\\
12861.215303826	2.47981802802665\\
12873.2183045761	0.884364267480336\\
12885.2213053263	1.27225052006909\\
12897.2243060765	4.10659667393028\\
12909.2273068267	5.39773083360063\\
12921.2303075769	6.68557184148504\\
12933.2333083271	0.570820725644466\\
12945.2363090773	0.727734090903073\\
12957.2393098275	3.43152579738281\\
12969.2423105776	10.7281405148079\\
12981.2453113278	2.75839805043651\\
12993.248312078	4.55838035734409\\
13005.2513128282	3.00566715335129\\
13017.2543135784	3.11354490510415\\
13029.2573143286	4.87260269666214\\
13041.2603150788	2.62831899596026\\
13053.263315829	9.13469938979204\\
13065.2663165791	5.61485270847944\\
13077.2693173293	0.708838644068309\\
13089.2723180795	6.54140978327585\\
13101.2753188297	5.95576018048629\\
13113.2783195799	8.70548236065099\\
13125.2813203301	2.93478104757407\\
13137.2843210803	0.0109464705367402\\
13149.2873218305	0.377487643594294\\
13161.2903225806	2.40220875446074\\
13173.2933233308	0.821657200411344\\
13185.296324081	1.18625927465993\\
13197.2993248312	4.00372554560805\\
13209.3023255814	5.20493710713731\\
13221.3053263316	6.39208967454943\\
13233.3083270818	0.540783684627421\\
13245.311327832	0.673210955223951\\
13257.3143285821	3.2731477595435\\
13269.3173293323	9.93187861813464\\
13281.3203300825	2.53681167611453\\
13293.3233308327	4.16548998959254\\
13305.3263315829	2.99768517533637\\
13317.3293323331	2.67076503398581\\
13329.3323330833	4.59891833282712\\
13341.3353338335	2.48089332897744\\
13353.3383345836	8.5733683916917\\
13365.3413353338	5.33699107650189\\
13377.344336084	0.713028682430017\\
13389.3473368342	6.11970272415649\\
13401.3503375844	5.28864970187505\\
13413.3533383346	8.21506567623998\\
13425.3563390848	2.66927815799728\\
13437.359339835	0.0160673422557117\\
13449.3623405851	0.436874613766323\\
13461.3653413353	2.33038049836731\\
13473.3683420855	0.766113916823531\\
13485.3713428357	1.10944675137984\\
13497.3743435859	3.90745052933352\\
13509.3773443361	5.02650475305943\\
13521.3803450863	6.1221950774753\\
13533.3833458365	0.51435808479786\\
13545.3863465866	0.625582287744656\\
13557.3893473368	3.12919322263422\\
13569.392348087	9.22355755938164\\
13581.3953488372	2.34053798990569\\
13593.3983495874	3.82183356068795\\
13605.4013503376	2.98692644807562\\
13617.4043510878	2.29634419217367\\
13629.407351838	4.35638378371988\\
13641.4103525881	2.34730412099819\\
13653.4133533383	8.07213162789325\\
13665.4163540885	5.09074869400916\\
13677.4193548387	0.71620665980202\\
13689.4223555889	5.74565186086304\\
13701.4253563391	4.71149917436627\\
13713.4283570893	7.77320469713244\\
13725.4313578395	2.44040357583785\\
13737.4343585896	0.022972625602839\\
13749.4373593398	0.495239374134625\\
13761.44036009	2.26379345645391\\
13773.4433608402	0.716699494378227\\
13785.4463615904	1.04056946983559\\
13797.4493623406	3.81734977905233\\
13809.4523630908	4.8611454866919\\
13821.455363841	5.87352403973737\\
13833.4583645911	0.491019051641469\\
13845.4613653413	0.583782413506157\\
13857.4643660915	2.99793550243108\\
13869.4673668417	8.59069907378147\\
13881.4703675919	2.16589530709245\\
13893.4733683421	3.51974915209079\\
13905.4763690923	2.97411582267073\\
13917.4793698425	1.97830198843691\\
13929.4823705927	4.14046366346698\\
13941.4853713428	2.22586829235236\\
13953.488372093	7.6226891542992\\
13965.4913728432	4.87153690300681\\
13977.4943735934	0.718577500024581\\
13989.4973743436	5.41239106977465\\
14001.5003750938	4.2097001322755\\
14013.503375844	7.37360636701385\\
14025.5063765942	2.24216480388986\\
14037.5093773443	0.031152263871371\\
14049.5123780945	0.552108307625342\\
14061.5153788447	2.20198218785555\\
14073.5183795949	0.672561312918731\\
14085.5213803451	0.978588964613066\\
14097.5243810953	3.733035672365\\
14109.5273818455	4.70771905948427\\
14121.5303825957	5.64400714082511\\
14133.5333833458	0.470335721405171\\
14145.536384096	0.546942350828341\\
14157.5393848462	2.87791580248455\\
14169.5423855964	8.02296987854119\\
14181.5453863466	2.00984562440444\\
14193.5483870968	3.25299589374377\\
14205.551387847	2.95985827994908\\
14217.5543885972	1.70705563810083\\
14229.5573893473	3.94744590630141\\
14241.5603900975	2.11515699102698\\
14253.5633908477	7.21815744465963\\
14265.5663915979	4.67560577280093\\
14277.5693923481	0.720313831220128\\
14289.5723930983	5.11426414150924\\
14301.5753938485	3.77140905222519\\
14313.5783945987	7.01099613931119\\
14325.5813953488	2.06971759385929\\
14337.584396099	0.0402278190610738\\
14349.5873968492	0.607213967239312\\
14361.5903975994	2.14454205206576\\
14373.5933983496	0.632991624482647\\
14385.5963990998	0.922632315568945\\
14397.59939985	3.6541524083675\\
14409.6024006002	4.56521339496864\\
14421.6054013503	5.43182654938885\\
14433.6084021005	0.451951931429755\\
14445.6114028507	0.514348306633325\\
14457.6144036009	2.76789304110181\\
14469.6174043511	7.51175160096191\\
14481.6204051013	1.86986396798403\\
14493.6234058515	3.01645125290151\\
14505.6264066017	2.94465816160668\\
14517.6294073518	1.47488265746769\\
14529.632408102	3.77426786639376\\
14541.6354088522	2.01395003287055\\
14553.6384096024	6.85278620279552\\
14565.6414103526	4.49986632926117\\
14577.6444111028	0.721561481459283\\
14589.647411853	4.84657728946613\\
14601.6504126032	3.38695283443171\\
14613.6534133533	6.68092292248847\\
14625.6564141035	1.91911717193586\\
14637.6594148537	0.0499187865355322\\
14649.6624156039	0.660431981467325\\
14661.6654163541	2.09111879356082\\
14673.6684171043	0.597398857290335\\
14685.6714178545	0.871961335325911\\
14697.6744186047	3.58037364702878\\
14709.6774193548	4.43272780543316\\
14721.680420105	5.23538023304272\\
14733.6834208552	0.435571382512894\\
14745.6864216054	0.485409987726249\\
14757.6894223556	2.66680464063965\\
14769.6924231058	7.04980853818686\\
14781.695423856	1.74383773200086\\
14793.6984246062	2.80588070605893\\
14805.7014253563	2.92893569941829\\
14817.7044261065	1.27551664686054\\
14829.7074268567	3.61838398849386\\
14841.7104276069	1.92119975735525\\
14853.7134283571	6.52173962924542\\
14865.7164291073	4.34175532759873\\
14877.7194298575	0.722444010040282\\
14889.7224306077	4.60540889114397\\
14901.7254313578	3.0483767693291\\
14913.728432108	6.37960684587277\\
14925.7314328582	1.78712940211996\\
14937.7344336084	0.0600179772116611\\
14949.7374343586	0.711736696789746\\
14961.7404351088	2.04140041387096\\
14973.743435859	0.565285386520868\\
14985.7464366092	0.825948293633354\\
14997.7494373593	3.51140024753655\\
15009.7524381095	4.30945875406472\\
15021.7554388597	5.05325201842424\\
15033.7584396099	0.420946128903554\\
15045.7614403601	0.459636170861809\\
15057.7644411103	2.57373556918752\\
15069.7674418605	6.63102850654245\\
15081.7704426107	1.62998837087504\\
15093.7734433608	2.61776071751754\\
15105.776444111	2.91304105489712\\
15117.7794448612	1.10384019834461\\
15129.7824456114	3.47766406437864\\
15141.7854463616	1.83600210500209\\
15153.7884471118	6.22092581829699\\
15165.791447862	4.19913134241968\\
15177.7944486122	0.723066448290983\\
15189.7974493623	4.38746181016747\\
15201.8004501125	2.749097513825\\
15213.8034508627	6.10381932756549\\
15225.8064516129	1.67108607329736\\
15237.8094523631	0.0703733828613696\\
15249.8124531133	0.76116988060405\\
15261.8154538635	1.99511073461993\\
15273.8184546137	0.536230149013388\\
15285.8214553638	0.784056628992388\\
15297.824456114	3.44695813657863\\
15309.8274568642	4.19468772853521\\
15321.8304576144	4.88418642501778\\
15333.8334583646	0.407867571083815\\
15345.8364591148	0.436615699382783\\
15357.839459865	2.4878936671092\\
15369.8424606152	6.25021888458185\\
15381.8454613653	1.52680993365274\\
15393.8484621155	2.44914141360818\\
15405.8514628657	2.8972661547969\\
15417.8544636159	0.955649456982197\\
15429.8574643661	3.35031424915575\\
15441.8604651163	1.7575732886387\\
15453.8634658665	5.94686248605972\\
15465.8664666167	4.07019416841062\\
15477.8694673668	0.723518393369857\\
15489.872468117	4.18994773824854\\
15501.8754688672	2.48363428651386\\
15513.8784696174	5.85078776668188\\
15525.8814703676	1.56877301466476\\
15537.8844711178	0.0808747215956818\\
15549.887471868	0.80881859840047\\
15561.8904726182	1.95200423105381\\
15573.8934733683	0.509874925481994\\
15585.8964741185	0.745825502459487\\
15597.8994748687	3.3867963231968\\
15609.9024756189	4.08777087490176\\
15621.9054763691	4.72706741750824\\
15633.9084771193	0.396159349794397\\
15645.9114778695	0.416002577134535\\
15657.9144786197	2.40858981040784\\
15669.9174793698	5.90294475699814\\
15681.92048012	1.43302041494205\\
15693.9234808702	2.29753912400257\\
15705.9264816204	2.88185462397889\\
15717.9294823706	0.827472081725154\\
15729.9324831208	3.23481519487994\\
15741.935483871	1.6852308376003\\
15753.9384846212	5.69657040206045\\
15765.9414853713	3.95342175158643\\
15777.9444861215	0.723876573122558\\
15789.9474868717	4.0104958561232\\
15801.9504876219	2.24739890912728\\
15813.9534883721	5.61811918898473\\
15825.9564891223	1.47834287565246\\
15837.9594898725	0.0914433921845732\\
15849.9624906227	0.854799578385036\\
15861.9654913728	1.91186183416966\\
15873.968492123	0.485913427404359\\
15885.9714928732	0.710857337980498\\
15897.9744936234	3.33068506562354\\
15909.9774943736	3.98813010552263\\
15921.9804951238	4.58090039108285\\
15933.983495874	0.385671697582169\\
15945.9864966242	0.397504185341772\\
15957.9894973743	2.33522183405326\\
15969.9924981245	5.58539947046212\\
15981.9954988747	1.3475229505027\\
15993.9984996249	2.16085161238712\\
16006.0015003751	2.86701009904426\\
16018.0045011253	0.716425378775076\\
16030.0075018755	3.12987318553927\\
16042.0105026257	1.61837808954209\\
16054.0135033758	5.46748814351946\\
16066.016504126	3.84752043281363\\
16078.0195048762	0.724206977952939\\
16090.0225056264	3.8470801368937\\
16102.0255063766	2.03653052948703\\
16114.0285071268	5.403738615551\\
16126.031507877	1.39824660829945\\
16138.0345086272	0.102024931560881\\
16150.0375093773	0.899248190201014\\
16162.0405101275	1.874487482752\\
16174.0435108777	0.464082548700711\\
16186.0465116279	0.678807703433737\\
16198.0495123781	3.27841418978998\\
16210.0525131283	3.89524544571692\\
16222.0555138785	4.44479683761791\\
16234.0585146287	0.376276917537784\\
16246.0615153788	0.380871900428471\\
16258.064516129	2.26726140501211\\
16270.0675168792	5.29430036129656\\
16282.0705176294	1.26937464027202\\
16294.0735183796	2.03729069904243\\
16306.0765191298	2.85290317882687\\
16318.07951988	0.620104924211657\\
16330.0825206302	3.03438123889291\\
16342.0855213803	1.55649142179634\\
16354.0885221305	5.25740340641825\\
16366.0915228807	3.7513853926285\\
16378.0945236309	0.724566638093823\\
16390.0975243811	3.69796106435066\\
16402.1005251313	1.84776456697944\\
16414.1035258815	5.20583896375493\\
16426.1065266317	1.3271792532743\\
16438.1095273818	0.112583324507562\\
16450.112528132	0.942310720931241\\
16462.1155288822	1.83970526295401\\
16474.1185296324	0.4441553028725\\
16486.1215303826	0.649377040780701\\
16498.1245311328	3.22979155357747\\
16510.127531883	3.80864842700066\\
16522.1305326332	4.31796124476257\\
16534.1335333833	0.36786574073041\\
16546.1365341335	0.365893572864264\\
16558.1395348837	2.20424322909943\\
16570.1425356339	5.02680418914111\\
16582.1455363841	1.19776132936454\\
16594.1485371343	1.92532832777296\\
16606.1515378845	2.83967723372668\\
16618.1545386347	0.536496518478439\\
16630.1575393848	2.94738791585562\\
16642.160540135	1.4991096756099\\
16654.1635408852	5.06439727883304\\
16666.1665416354	3.66406898024526\\
16678.1695423856	0.725005109906327\\
16690.1725431358	3.56163858847724\\
16702.175543886	1.67832808282091\\
16714.1785446362	5.02284005419032\\
16726.1815453863	1.26403675381\\
16738.1845461365	0.123096693422037\\
16750.1875468867	0.984139016559549\\
16762.1905476369	1.80735701503445\\
16774.1935483871	0.425935083\\
16786.1965491373	0.622303867191709\\
16798.1995498875	3.18464165076022\\
16810.2025506377	3.72791636592719\\
16822.2055513878	4.19967986343538\\
16834.208552138	0.360344373863411\\
16846.2115528882	0.352387459667311\\
16858.2145536384	2.14575611965378\\
16870.2175543886	4.78043811660771\\
16882.2205551388	1.13197707820017\\
16894.223555889	1.82365310781483\\
16906.2265566392	2.82745326460418\\
16918.2295573893	0.463906134582099\\
16930.2325581395	2.86807213973052\\
16942.2355588897	1.44582534784908\\
16954.2385596399	4.88679874006375\\
16966.2415603901	3.58475518524921\\
16978.2445611403	0.725565723118654\\
16990.2475618905	3.43681390906028\\
17002.2505626407	1.52585571012185\\
17014.2535633908	4.85335486075228\\
17026.256564141	1.20788133681646\\
17038.2595648912	0.133554023216886\\
17050.2625656414	1.02488682508657\\
17062.2655663916	1.77730031491013\\
17074.2685671418	0.409250968046063\\
17086.271567892	0.597359154716488\\
17098.2745686422	3.1428043450838\\
17110.2775693924	3.65266739742845\\
17122.2805701425	4.08931104546132\\
17134.2835708927	0.353632093005603\\
17146.2865716429	0.340197300452029\\
17158.2895723931	2.09143556219913\\
17170.2925731433	4.55304303976174\\
17182.2955738935	1.07140734950238\\
17206.3015753938	2.81633397183839\\
17218.304576144	0.400903844976777\\
17230.3075768942	2.79572273714989\\
17242.3105776444	1.39627721652269\\
17254.3135783946	4.72314728644694\\
17266.3165791448	3.5127389319852\\
17278.319579895	0.726286631235113\\
17290.3225806452	3.32235824304272\\
17302.3255813954	1.38832169620003\\
17314.3285821455	4.6961615623874\\
17326.3315828957	1.1579135972243\\
17338.3345836459	0.143952667267176\\
17350.3375843961	1.06470736652485\\
17362.3405851463	1.74940676087718\\
17374.3435858965	0.393953862481109\\
17386.3465866467	0.574341659898433\\
17398.3495873968	3.10413372623156\\
17410.352588147	3.58255615132788\\
17422.3555888972	3.9862769048681\\
17434.3585896474	0.347659272273716\\
17446.3615903976	0.329188299391074\\
17458.3645911478	2.04095749007616\\
17470.367591898	4.34272679741532\\
17482.3705926482	1.01551516169958\\
17494.3735933984	1.64679697244824\\
17506.3765941485	2.80640716931347\\
17518.3795948987	0.346278683148299\\
17530.3825956489	2.72972171682586\\
17542.3855963991	1.35014413620623\\
17554.3885971493	4.57216205887797\\
17566.3915978995	3.44740918780186\\
17578.3945986497	0.727201699761482\\
17590.3975993999	3.21728715620454\\
17602.40060015	1.2639846570428\\
17614.4036009002	4.55018027228014\\
17626.4066016504	1.11344986281221\\
17638.4096024006	0.154296446193408\\
17650.4126031508	1.103751788016\\
17662.415603901	1.7235605102726\\
17674.4186046512	0.379913304282826\\
17686.4216054014	0.553074023950918\\
17698.4246061515	3.06849707792309\\
17710.4276069017	3.5172699800935\\
17722.4306076519	3.89005610104703\\
17734.4336084021	0.342365761227889\\
17746.4366091523	0.319243829215505\\
17758.4396099025	1.99403304656538\\
17770.4426106527	4.14782533117069\\
17782.4456114029	0.963829624207807\\
17794.448612153	1.56979065401681\\
17806.4516129032	2.79774865604249\\
17818.4546136534	0.299002114966017\\
17830.4576144036	2.66953052746018\\
17842.4606151538	1.30713979370436\\
17854.463615904	4.43271620678633\\
17866.4666166542	3.38823510873593\\
17878.4696174044	0.728341260402118\\
17890.4726181545	3.12073935596759\\
17902.4756189047	1.15134242627058\\
17914.4786196549	4.41445356266593\\
17926.4816204051	1.07390374548904\\
17938.4846211553	0.164594199029792\\
17950.4876219055	1.14216825779236\\
17962.4906226557	1.69965702376065\\
17974.4936234059	0.367014812224264\\
17986.496624156	0.533399501356621\\
17998.4996249062	3.03577395022086\\
18010.5026256564	3.45652566053732\\
18022.5056264066	3.80017757409102\\
18034.5086271568	0.337699543601463\\
18046.511627907	0.310262714514593\\
18058.5146286572	1.95040415591237\\
18070.5176294074	3.96687028347379\\
18082.5206301575	0.91593639719924\\
18094.5236309077	1.49937769535782\\
18106.5266316579	2.79042463881846\\
18118.5296324081	0.258198331448931\\
18130.5326331583	2.61467870607284\\
18142.5356339085	1.26700825599285\\
18154.5386346587	4.30381549402339\\
18166.5416354089	3.33475461928573\\
18178.544636159	0.729732754571867\\
18190.5476369092	3.03195908222619\\
18202.5506376594	1.04909496845018\\
18214.5536384096	4.2881300854646\\
18226.5566391598	1.03877103221554\\
18238.55963991	0.174858681461028\\
18250.5626406602	1.1801015198164\\
18262.5656414104	1.6776019831216\\
18274.5686421605	0.355157671084619\\
18286.5716429107	0.515179203803809\\
18298.5746436609	3.00585532709311\\
18310.5776444111	3.40006650401525\\
18322.5806451613	3.71621509310959\\
18334.5836459115	0.333615624265347\\
18346.5866466617	0.302156982276855\\
18358.5896474119	1.90983976135725\\
18370.592648162	3.79856183698395\\
18382.5956489122	0.871469714657782\\
18394.5986496624	1.43491329181687\\
18406.6016504126	2.78449378430372\\
18418.6046511628	0.22311997817133\\
18430.607651913	2.56475445708164\\
18442.6106526632	1.22952017499338\\
18454.6136534134	4.18458036231621\\
18466.6166541635	3.28656495702003\\
18478.6196549137	0.731401285412841\\
18490.6226556639	2.95028141877267\\
18502.6256564141	0.95611377251038\\
18514.6286571643	4.17045073432626\\
18526.6316579145	1.00761725542145\\
18538.6346586647	0.185105731752627\\
18550.6376594149	1.21769278018761\\
18562.640660165	1.65731035521319\\
18574.6436609152	0.34425307444445\\
18586.6466616654	0.498289768955769\\
18598.6496624156	2.97864288199267\\
18610.6526631658	3.34765982025145\\
18622.655663916	3.63778249945562\\
18634.6586646662	0.330075102592631\\
18646.6616654164	0.294849991604903\\
18658.6646661665	1.87213261637083\\
18670.6676669167	3.64174584483648\\
18682.6706676669	0.830105684434148\\
18694.6736684171	1.37583287932981\\
18706.6766691673	2.78000896460039\\
18718.6796699175	0.193128242419996\\
18730.6826706677	2.51939679967665\\
18742.6856714179	1.19446953979532\\
18754.688672168	4.07423082816403\\
18766.6916729182	3.24331481160339\\
18778.6946736684	0.733370093940196\\
18790.6976744186	2.87511998654533\\
18802.7006751688	0.871416481049296\\
18814.703675919	4.06073690224595\\
18826.7066766692	0.980067425644669\\
18838.7096774194	0.195353644222132\\
18850.7126781695	1.25507983186245\\
18862.7156789197	1.63870558069412\\
18874.7186796699	0.33422256111154\\
18886.7216804201	0.482621380969839\\
18898.7246811703	2.95404831458768\\
18910.7276819205	3.29909468817787\\
18922.7306826707	3.56452954676861\\
18934.7336834209	0.327044398778369\\
18946.736684171	0.288274872658785\\
18958.7396849212	1.83709653730523\\
18970.7426856714	3.4953944900999\\
18982.7456864216	0.791556636496227\\
18994.7486871718	1.32164095364564\\
19006.751687922	2.77701875067976\\
19018.7546886722	0.167676452559105\\
19030.7576894224	2.47828899747129\\
19042.7606901725	1.16167088669828\\
19054.7636909227	3.97207371312968\\
19066.7666916729	3.20469776637005\\
19078.7696924231	0.735660972692131\\
19090.7726931733	2.80595659126195\\
19102.7756939235	0.794145772440985\\
19114.7786946737	3.95838047607897\\
19126.7816954239	0.955797518446937\\
19138.784696174	0.205622704596033\\
19150.7876969242	1.29239735089394\\
19162.7906976744	1.62171886930168\\
19174.7936984246	0.324996693888832\\
19186.7966991748	0.468076083758969\\
19198.799699925	2.93199276170989\\
19210.8027006752	3.2541799941403\\
19222.8057014254	3.49613825507882\\
19234.8087021755	0.324494606456034\\
19246.8117029257	0.282373219131097\\
19258.8147036759	1.80456404252673\\
19270.8177044261	3.35858986330469\\
19282.8207051763	0.755566335747538\\
19294.8237059265	1.27190169303104\\
19306.8267066767	2.775568698467\\
19318.8297074269	0.146296523250586\\
19330.832708177	2.44115304251094\\
19342.8357089272	1.13095689440678\\
19354.8387096774	3.87749180729751\\
19366.8417104276	3.17044681036325\\
19378.8447111778	0.738294627765743\\
19390.847711928	2.74233248241033\\
19402.8507126782	0.723551715007864\\
19414.8537134284	3.8628352781238\\
19426.8567141785	0.934527392083682\\
19438.8597149287	0.215934852390033\\
19450.8627156789	1.3297773154529\\
19462.8657164291	1.60628858733854\\
19474.8687171793	0.316513939445428\\
19486.8717179295	0.454566338949626\\
19498.8747186797	2.91240627755092\\
19510.8777194299	3.21274270391862\\
19522.88072018	3.43231970843754\\
19534.8837209302	0.322400950335164\\
19546.8867216804	0.277093989784586\\
19558.8897224306	1.77438431716776\\
19570.8927231808	3.23050996167224\\
19582.895723931	0.721905911220483\\
19594.8987246812	1.22623106472266\\
19606.9017254314	2.77570246497625\\
19618.9047261815	0.12858771869931\\
19630.9077269317	2.40774501100118\\
19642.9107276819	1.1021763038014\\
19654.9137284321	3.78993464039873\\
19666.9167291823	3.14032973539789\\
19678.9197299325	0.741290998363332\\
19690.9227306827	2.6838409475752\\
19702.9257314329	0.658976969876397\\
19714.928732183	3.77360971742176\\
19726.9317329332	0.916014878494747\\
19738.9347336834	0.226313443773468\\
19750.9377344336	1.36734951381196\\
19762.9407351838	1.59235972541214\\
19774.943735934	0.30871971589534\\
19786.9467366842	0.442013789222864\\
19798.9497374344	2.89522737750446\\
19810.9527381845	3.17462633970935\\
19822.9557389347	3.37281123717681\\
19834.9587396849	0.320742331354262\\
19846.9617404351	0.272392582977159\\
19858.9647411853	1.74642145348663\\
19870.9677419355	3.11041670890767\\
19882.9707426857	0.690370381566154\\
19894.9737434359	1.18429015925098\\
19906.976744186	2.77746278640288\\
19918.9797449362	0.114207313866424\\
19930.9827456864	2.37785114434818\\
19942.9857464366	1.075192112965\\
19954.9887471868	3.70891059758193\\
19966.991747937	3.11414526913762\\
19978.9947486872	0.744669541674435\\
19990.9977494374	2.63012101814283\\
20003.0007501875	0.599844342011941\\
20015.0037509377	3.69026045785642\\
20027.0067516879	0.900050841135244\\
20039.0097524381	0.236783094533201\\
20051.0127531883	1.40524211730885\\
20063.0157539385	1.57988343654722\\
20075.0187546887	0.301565580866704\\
20087.0217554389	0.430348194728744\\
20099.0247561891	2.88040264175482\\
20111.0277569392	3.13968963752821\\
20123.0307576894	3.31737393436929\\
20135.0337584396	0.319500945634534\\
20147.0367591898	0.26823005534972\\
20159.03975994	1.72055292531827\\
20171.0427606902	2.99764566544042\\
20183.0457614404	0.660775678532786\\
20195.0487621905	1.14577954501821\\
20207.0517629407	2.78089234420007\\
20219.0547636909	0.102862817871057\\
20231.0577644411	2.35128453715096\\
20243.0607651913	1.04988000593579\\
20255.0637659415	3.63398016417145\\
20267.0667666917	3.09171982473972\\
20279.0697674419	0.748449489669195\\
20291.0727681921	2.58085210587486\\
20303.0757689422	0.545646276459063\\
20315.0787696924	3.6123869433376\\
20327.0817704426	0.886455034176004\\
20339.0847711928	0.24736958768143\\
20351.087771943	1.44358230214392\\
20363.0907726932	1.56881663647505\\
20375.0937734434	0.295008537963749\\
20387.0967741935	0.419506516093412\\
20399.0997749437	2.86788637404379\\
20411.1027756939	3.10780536409922\\
20423.1057764441	3.26579046333322\\
20435.1087771943	0.318661965741071\\
20447.1117779445	0.264572461077699\\
20459.1147786947	1.69666826262449\\
20471.1177794449	2.8915971605371\\
20483.1207801951	0.632956088527443\\
20495.1237809452	1.1104344752866\\
20507.1267816954	2.78603454227265\\
20519.1297824456	0.0943054900974129\\
20531.1327831958	2.32788233598872\\
20543.135783946	1.026126980915\\
20555.1387846962	3.56475012342522\\
20567.1417854464	3.07290476960184\\
20579.1447861965	0.752650083401107\\
20591.1477869467	2.5357494209647\\
20603.1507876969	0.495935973087574\\
20615.1537884471	3.53962664895553\\
20627.1567891973	0.875072629152447\\
20639.1597899475	0.258099833910652\\
20651.1627906977	1.48249690978292\\
20663.1657914479	1.55912165934814\\
20675.1687921981	0.289010443369654\\
20687.1717929482	0.409432121916176\\
20699.1747936984	2.8576403131137\\
20711.1777944486	3.07885927526928\\
20723.1807951988	3.21786311963121\\
20735.183795949	0.318213275127336\\
20747.1867966992	0.261390292482912\\
20759.1897974494	1.67466789753795\\
20771.1927981996	2.79172862259404\\
20783.1957989497	0.606762045931557\\
20795.1987996999	1.07802081053724\\
20807.2018004501	2.79293421422777\\
20819.2048012003	0.0883249320957411\\
20831.2078019505	2.30750337135593\\
20843.2108027007	1.00383014906014\\
20855.2138034509	3.50086856181878\\
20867.2168042011	3.05757413543375\\
20879.2198049512	0.757290789632899\\
20891.2228057014	2.4945600512902\\
20903.2258064516	0.450319854839245\\
20915.2288072018	3.47165094873928\\
20927.231807952	0.865771300912716\\
20939.2348087022	0.269001876018504\\
20951.2378094524	1.52211313981007\\
20963.2408102026	1.55076596324986\\
20975.2438109527	0.283537497742674\\
20987.2468117029	0.400074102447876\\
20999.2498124531	2.8496333933778\\
21011.2528132033	3.05274920031633\\
21023.2558139535	3.17341211606033\\
21035.2588147037	0.318145248328034\\
21047.2618154539	0.258658006462785\\
21059.2648162041	1.65446215816371\\
21071.2678169542	2.69754792442408\\
21083.2708177044	0.582058223901524\\
21095.2738184546	1.04833154438325\\
21107.2768192048	2.80163827718837\\
21119.279819955	0.084744580100556\\
21131.2828207052	2.29002615870067\\
21143.2858214554	0.982895679754912\\
21155.2888222056	3.44202056246928\\
21167.2918229557	3.04562270512531\\
21179.2948237059	0.762391504202411\\
21191.2978244561	2.45705960221935\\
21203.3008252063	0.408451172983491\\
21215.3038259565	3.40816150887026\\
21227.3068267067	0.858438785145219\\
21239.3098274569	0.280104930694768\\
21251.3128282071	1.56255927295202\\
21263.3158289572	1.54372188095474\\
21275.3188297074	0.278559810965881\\
21287.3218304576	0.391386674453737\\
21299.3248312078	2.84384155291145\\
21311.327831958	3.02938423918046\\
21323.3308327082	3.13227406394839\\
21335.3338334584	0.318450570627264\\
21347.3368342086	0.256353623830698\\
21359.3398349587	1.63597039030113\\
21371.3428357089	2.60860758985221\\
21383.3458364591	0.558721877208388\\
21395.3488372093	1.02118384144275\\
21407.3518379595	2.81219634727796\\
21419.3548387097	0.0834179564915747\\
21431.3578394599	2.27534721748129\\
21443.3608402101	0.963237871767052\\
21455.3638409602	3.38792448660613\\
21467.3668417104	3.03696442544864\\
21479.3698424606	0.767972745637406\\
21491.3728432108	2.42304931532396\\
21503.375843961	0.370024572727591\\
21515.3788447112	3.34888713104217\\
21527.3818454614	0.852980835876522\\
21539.3848462116	0.291439462880619\\
21551.3878469617	1.6039654240035\\
21563.3908477119	1.53796641185047\\
21575.3938484621	0.274051029749709\\
21587.3968492123	0.383328664168927\\
21599.3998499625	2.84024758691216\\
21611.4028507127	3.00868406125447\\
21623.4058514629	3.09430062702897\\
21635.4088522131	0.319124092591839\\
21647.4118529632	0.254458391167531\\
21659.4148537134	1.61912019038794\\
21671.4178544636	2.524499733883\\
21683.4208552138	0.536641399707359\\
21695.423855964	0.996416511598691\\
21707.4268567142	2.82466132898844\\
21719.4298574644	0.0842255655337081\\
21731.4328582146	2.26337966505291\\
21743.4358589647	0.94477833316084\\
21755.4388597149	3.33832876159324\\
21767.4418604651	3.03153110304759\\
21779.4448612153	0.774055842596389\\
21791.4478619655	2.39235359757888\\
21803.4508627157	0.334771474148968\\
21815.4538634659	3.29358098059591\\
21827.4568642161	0.849319525185231\\
21839.4598649662	0.303037289421408\\
21851.4628657164	1.64646432699618\\
21863.4658664666	1.53348105240496\\
21875.4688672168	0.269988019536878\\
21887.471867967	0.375863058017318\\
21899.4748687172	2.83884104510435\\
21911.4778694674	2.99057829578593\\
21923.4808702176	3.05935732877506\\
21935.4838709677	0.320162715214279\\
21947.4868717179	0.252956496442217\\
21959.4898724681	1.60384673534202\\
21971.4928732183	2.44485162955433\\
21983.4958739685	0.515715064693339\\
21995.4988747187	0.973887859007472\\
22007.5018754689	2.83908999025313\\
22019.5048762191	0.0870723406057673\\
22031.5078769692	2.25405205169897\\
22043.5108777194	0.927445255208966\\
22055.5138784696	3.29300910618743\\
22067.5168792198	3.02927135030606\\
22079.51987997	0.780663118087681\\
22091.5228807202	2.36481790495073\\
22103.5258814704	0.302456149584776\\
22115.5288822206	3.24201814704913\\
22127.5318829707	0.847391837789556\\
22139.5348837209	0.31493170972127\\
22151.5378844711	1.6901921566127\\
22163.5408852213	1.53025166210949\\
22175.5438859715	0.266350593821455\\
22187.5468867217	0.368956612048867\\
22199.5498874719	2.83961817230211\\
22211.5528882221	2.97500600599022\\
22223.5558889722	3.02732249496194\\
22235.5588897224	0.321565302703761\\
22247.5618904726	0.25183483148984\\
22259.5648912228	1.59009219784202\\
22271.567891973	2.36932181112546\\
22283.5708927232	0.495849922209744\\
22295.5738934734	0.953473854743925\\
22307.5768942236	2.85554353404008\\
22319.5798949737	0.0918855680595813\\
22331.5828957239	2.24730740866239\\
22343.5858964741	0.911172767609649\\
22355.5888972243	3.25176613610279\\
22367.5918979745	3.03014975368597\\
22379.5948987247	0.787818073458179\\
22391.5978994749	2.34030693395138\\
22403.6009002251	0.272872399820162\\
22415.6039009752	3.19399348991841\\
22427.6069017254	0.847148522318217\\
22439.6099024756	0.327157662611023\\
22451.6129032258	1.73528939132622\\
22463.615903976	1.52826836375683\\
22475.6189047262	0.263121285251713\\
22487.6219054764	0.362579512589777\\
22499.6249062266	2.84258189199473\\
22511.6279069767	2.96191524010352\\
22523.6309077269	2.99808631779213\\
22535.6339084771	0.323332620345191\\
22547.6369092273	0.251082795483984\\
22559.6399099775	1.5778052364942\\
22571.6429107277	2.29759663770237\\
22583.6459114779	0.476960830573465\\
22595.6489122281	0.935066591091585\\
22607.6519129782	2.87408817598231\\
22619.6549137284	0.0986132276773663\\
22631.6579144786	2.24310248670085\\
22643.6609152288	0.895900364485954\\
22655.663915979	3.21442330231724\\
22667.6669167292	3.03414624327087\\
22679.6699174794	0.795545574760647\\
22691.6729182296	2.31870308236285\\
22703.6759189797	0.245840749125871\\
22715.6789197299	3.14931973178476\\
22727.6819204801	0.848553168649402\\
22739.6849212303	0.339751908742888\\
22751.6879219805	1.78190172508886\\
22763.6909227307	1.52752547576998\\
22775.6939234809	0.260285153662721\\
22787.6969242311	0.356705081745803\\
22799.6999249812	2.84774183209058\\
22811.7029257314	2.9512626526222\\
22823.7059264816	2.97155002775214\\
22835.7089272318	0.325467295617561\\
22847.711927982	0.250692134705682\\
22859.7149287322	1.56694055296454\\
22871.7179294824	2.22938725277313\\
22883.7209302326	0.458969603980683\\
22895.7239309827	0.918572983476628\\
22907.7269317329	2.89479573803742\\
22919.7299324831	0.107222702020436\\
22931.7329332333	2.24140716751569\\
22943.7359339835	0.881572391845358\\
22955.7389347337	3.18082512216222\\
22967.7419354839	3.04125564701331\\
22979.7449362341	0.803872044077969\\
22991.7479369842	2.29990514756874\\
23003.7509377344	0.221206093624287\\
23015.7539384846	3.10782576528664\\
23027.7569392348	0.851581486757717\\
23039.759939985	0.352753239423854\\
23051.7629407352	1.83018103634207\\
23063.7659414854	1.52802147662469\\
23075.7689422356	0.257829627324626\\
23087.7719429857	0.351309522629181\\
23099.7749437359	2.85511439434091\\
23111.7779444861	2.94301319199335\\
23123.7809452363	2.94762516256161\\
23135.7839459865	0.327973801136335\\
23147.7869467367	0.250656814945863\\
23159.7899474869	1.55745850841423\\
23171.7929482371	2.16442688506746\\
23183.7959489872	0.441804259965013\\
23195.7989497374	0.90391369194543\\
23207.8019504876	2.91774426678561\\
23219.8049512378	0.117699817448339\\
23231.807951988	2.24220403474953\\
23243.8109527382	0.868137588296084\\
23255.8139534884	3.15083567093718\\
23267.8169542386	3.05148741696331\\
23279.8199549888	0.812825658769073\\
23291.8229557389	2.28382723593382\\
23303.8259564891	0.198835750047045\\
23315.8289572393	3.06935514495098\\
23327.8319579895	0.856220767638056\\
23339.8349587397	0.366202712931229\\
23351.8379594899	1.8802864234572\\
23363.8409602401	1.52975899991699\\
23375.8439609902	0.255744374404866\\
23387.8469617404	0.346371699782668\\
23399.8499624906	2.86472286748413\\
23411.8529632408	2.93713984994984\\
23423.855963991	2.92623292368317\\
23435.8589647412	0.330858458587292\\
23447.8619654914	0.250972923397418\\
23459.8649662416	1.54932479334973\\
23471.8679669918	2.10246844356079\\
23483.8709677419	0.425398353582142\\
23495.8739684921	0.891022239754996\\
23507.8769692423	2.94301868688397\\
23519.8799699925	0.130048188515074\\
23531.8829707427	2.24548809355008\\
23543.8859714929	0.855548672428561\\
23555.8889722431	3.12433730621235\\
23567.8919729932	3.06486551928798\\
23579.8949737434	0.822436560860676\\
23591.8979744936	2.27039786273808\\
23603.9009752438	0.178617862139056\\
23615.903975994	3.03376473991942\\
23627.9069767442	0.862469512157413\\
23639.9099774944	0.380143920252613\\
23651.9129782446	1.93238531873766\\
23663.9159789948	1.53274486098469\\
23675.9189797449	0.254021202081828\\
23687.9219804951	0.341872951072737\\
23699.9249812453	2.87659758604057\\
23711.9279819955	2.93362346983434\\
23723.9309827457	2.90730361210714\\
23735.9339834959	0.334129463129934\\
23747.9369842461	0.251638597891975\\
23759.9399849963	1.54251014600905\\
23771.9429857464	2.04328236692584\\
23783.9459864966	0.40969038670103\\
23795.9489872468	0.879844311478809\\
23807.951987997	2.97071149792039\\
23819.9549887472	0.14428884542775\\
23831.9579894974	2.25126663272548\\
23843.9609902476	0.843761970569595\\
23855.9639909978	3.10122960312069\\
23867.9669917479	3.08142848265614\\
23879.9699924981	0.832737079710491\\
23891.9729932483	2.25955922471988\\
23903.9759939985	0.160460131251335\\
23915.9789947487	3.00092352595294\\
23927.9819954989	0.870337217110672\\
23939.9849962491	0.394623283010367\\
23951.9879969993	1.98665469309985\\
23963.9909977494	1.53699011475019\\
23975.9939984996	0.252653981571482\\
23987.9969992498	0.337796928205646\\
24000	2.89077613629063\\
};
\addlegendentry{Théorique}

\addplot [color=mycolor2]
  table[row sep=crcr]{%
-24000	1.02603713315632e-08\\
-23953.0791788856	5.34393847255656e-09\\
-23906.1583577713	5.42413927183039e-09\\
-23859.2375366569	7.06774407895119e-09\\
-23812.3167155425	5.64973238778839e-09\\
-23765.3958944282	5.05746122659395e-09\\
-23718.4750733138	8.98536656231893e-09\\
-23671.5542521994	7.51908440929027e-09\\
-23624.633431085	4.74083020421153e-09\\
-23577.7126099707	6.29998560168605e-09\\
-23530.7917888563	6.60837217254921e-09\\
-23483.8709677419	5.55926636051907e-09\\
-23436.9501466276	5.72810055214935e-09\\
-23390.0293255132	1.06515815446735e-08\\
-23343.1085043988	4.20392279312558e-09\\
-23296.1876832845	5.99102583098236e-09\\
-23249.2668621701	6.40325848302216e-09\\
-23202.3460410557	6.3417593009874e-09\\
-23155.4252199413	4.80809648060119e-09\\
-23108.504398827	1.01276976184229e-08\\
-23061.5835777126	6.06421456600552e-09\\
-23014.6627565982	5.02480237057881e-09\\
-22967.7419354839	6.89312843150773e-09\\
-22920.8211143695	6.05776913483648e-09\\
-22873.9002932551	5.0958738899701e-09\\
-22826.9794721408	8.45633150931201e-09\\
-22780.0586510264	8.21234650072718e-09\\
-22733.137829912	4.8690987647466e-09\\
-22686.2170087977	6.00326646783531e-09\\
-22639.2961876833	6.84046426932177e-09\\
-22592.3753665689	5.76473477928261e-09\\
-22545.4545454545	5.89794938544493e-09\\
-22498.5337243402	1.01587715494472e-08\\
-22451.6129032258	5.47687574771774e-09\\
-22404.6920821114	4.95811004243854e-09\\
-22357.7712609971	7.30929459174194e-09\\
-22310.8504398827	5.95385653238616e-09\\
-22263.9296187683	5.27777615739597e-09\\
-22217.008797654	9.45308701637496e-09\\
-22170.0879765396	7.36708150156432e-09\\
-22123.1671554252	4.19904953100462e-09\\
-22076.2463343109	7.16299642360804e-09\\
-22029.3255131965	6.2652016493601e-09\\
-21982.4046920821	5.93921942829406e-09\\
-21935.4838709677	6.65084162185285e-09\\
-21888.5630498534	1.03817025551138e-08\\
-21841.642228739	4.17547573945203e-09\\
-21794.7214076246	5.66631490157603e-09\\
-21747.8005865103	7.33156300299514e-09\\
-21700.8797653959	6.1253922599098e-09\\
-21653.9589442815	5.76766937967642e-09\\
-21607.0381231672	1.01374113790018e-08\\
-21560.1173020528	6.09338151824363e-09\\
-21513.1964809384	5.22968410751253e-09\\
-21466.275659824	6.55663995266352e-09\\
-21419.3548387097	7.05928626997468e-09\\
-21372.4340175953	5.312023953531e-09\\
-21325.5131964809	9.08864510372774e-09\\
-21278.5923753666	8.3927326533293e-09\\
-21231.6715542522	4.38861331152365e-09\\
-21184.7507331378	6.48938712799215e-09\\
-21137.8299120235	7.34535366459983e-09\\
-21090.9090909091	5.61050062894851e-09\\
-21043.9882697947	7.21158755747846e-09\\
-20997.0674486804	1.02478900824027e-08\\
-20950.146627566	5.16473754216914e-09\\
-20903.2258064516	5.39245820103201e-09\\
-20856.3049853372	7.35433742424774e-09\\
-20809.3841642229	6.57611067176519e-09\\
-20762.4633431085	6.20316276135553e-09\\
-20715.5425219941	9.46404796303222e-09\\
-20668.6217008798	7.81275897819787e-09\\
-20621.7008797654	4.17427421772841e-09\\
-20574.780058651	7.31738091003574e-09\\
-20527.8592375367	6.92432958599583e-09\\
-20480.9384164223	6.12350410724847e-09\\
-20434.0175953079	8.29237347580613e-09\\
-20387.0967741935	1.00347534616679e-08\\
-20340.1759530792	4.04119336604137e-09\\
-20293.2551319648	6.49254240605003e-09\\
-20246.3343108504	7.83829170368848e-09\\
-20199.4134897361	6.39923402437033e-09\\
-20152.4926686217	6.5355028140615e-09\\
-20105.5718475073	1.13600808867735e-08\\
-20058.651026393	5.38111297235994e-09\\
-20011.7302052786	5.76009655876556e-09\\
-19917.8885630499	8.05751524749803e-09\\
-19870.9677419355	6.16103856524798e-09\\
-19824.0469208211	9.30667100277566e-09\\
-19777.1260997067	9.34003373794343e-09\\
-19730.2052785924	4.16381658805751e-09\\
-19683.284457478	6.67165259273348e-09\\
-19636.3636363636	8.49843901656963e-09\\
-19589.4428152493	6.10557076235259e-09\\
-19542.5219941349	8.10332102863004e-09\\
-19495.6011730205	1.11462946802491e-08\\
-19448.6803519062	4.53921628927836e-09\\
-19401.7595307918	6.26181933601159e-09\\
-19354.8387096774	8.0023760125389e-09\\
-19307.917888563	7.13731911980736e-09\\
-19260.9970674487	7.47798387876937e-09\\
-19214.0762463343	1.00226736518046e-08\\
-19167.1554252199	8.1092263490313e-09\\
-19120.2346041056	4.79203561769186e-09\\
-19073.3137829912	7.35863916731919e-09\\
-19026.3929618768	8.184624590396e-09\\
-18979.4721407625	7.23586589970701e-09\\
-18932.5513196481	8.9424613602976e-09\\
-18885.6304985337	1.08575596667479e-08\\
-18838.7096774194	3.77439668496412e-09\\
-18791.788856305	7.88236179316005e-09\\
-18697.9472140762	7.87988675077825e-09\\
-18651.0263929619	7.48578650664353e-09\\
-18604.1055718475	1.26299000267282e-08\\
-18557.1847507331	5.18014238619997e-09\\
-18510.2639296188	5.94896768734084e-09\\
-18463.3431085044	8.61314250193111e-09\\
-18369.5014662757	7.69238769636575e-09\\
-18322.5806451613	1.06098991530897e-08\\
-18275.6598240469	9.54405326767685e-09\\
-18228.7390029326	4.26540136934464e-09\\
-18181.8181818182	7.84515715733362e-09\\
-18134.8973607038	9.19993340939328e-09\\
-18087.9765395894	7.92453273553115e-09\\
-18041.0557184751	9.08391947836229e-09\\
-17994.1348973607	1.22850308078516e-08\\
-17947.2140762463	4.28291575286981e-09\\
-17900.293255132	6.82393296991384e-09\\
-17853.3724340176	9.837130681176e-09\\
-17806.4516129032	7.50381204518002e-09\\
-17759.5307917889	9.22219969912846e-09\\
-17712.6099706745	1.22788153457188e-08\\
-17665.6891495601	7.47405512703119e-09\\
-17618.7683284457	5.20947442802222e-09\\
-17571.8475073314	9.10404741606395e-09\\
-17524.926686217	9.07667112276403e-09\\
-17478.0058651026	9.45621548026606e-09\\
-17431.0850439883	9.49301728577223e-09\\
-17384.1642228739	1.26686029518294e-08\\
-17337.2434017595	3.62653836333819e-09\\
-17290.3225806452	8.63791786349936e-09\\
-17243.4017595308	9.33461396119493e-09\\
-17196.4809384164	9.68967061830808e-09\\
-17149.5601173021	9.09690165260772e-09\\
-17102.6392961877	1.36020161407537e-08\\
-17055.7184750733	6.04650512139218e-09\\
-17008.7976539589	6.47542298194798e-09\\
-16961.8768328446	1.02594059331852e-08\\
-16914.9560117302	9.41698279185228e-09\\
-16868.0351906158	9.38462079846234e-09\\
-16821.1143695015	1.27393529421788e-08\\
-16774.1935483871	9.87191406989546e-09\\
-16727.2727272727	5.10078749970215e-09\\
-16680.3519061584	9.42256788487313e-09\\
-16633.431085044	1.06095544072376e-08\\
-16586.5102639296	9.70449248800454e-09\\
-16539.5894428152	1.1218052798535e-08\\
-16492.6686217009	1.32237156688076e-08\\
-16445.7478005865	5.1115927902975e-09\\
-16398.8269794721	8.42181055801457e-09\\
-16351.9061583578	1.10751518071119e-08\\
-16304.9853372434	1.0016079827131e-08\\
-16258.064516129	1.10780093295583e-08\\
-16211.1436950147	1.40422953979675e-08\\
-16164.2228739003	8.15683097146531e-09\\
-16117.3020527859	5.75053049556208e-09\\
-16070.3812316716	1.177785823999e-08\\
-16023.4604105572	1.07694236590659e-08\\
-15976.5395894428	1.17167204623824e-08\\
-15929.6187683284	1.21122211549134e-08\\
-15882.6979472141	1.35412413059574e-08\\
-15835.7771260997	4.36894984671782e-09\\
-15788.8563049853	9.83018098886207e-09\\
-15741.935483871	1.26137140667231e-08\\
-15695.0146627566	1.10584351326549e-08\\
-15648.0938416422	1.26037229925183e-08\\
-15601.1730205279	1.49051189026328e-08\\
-15554.2521994135	6.3435175461757e-09\\
-15507.3313782991	8.48223278088209e-09\\
-15460.4105571848	1.18486350529457e-08\\
-15413.4897360704	1.29309599630819e-08\\
-15366.568914956	1.1649670680257e-08\\
-15319.6480938416	1.62296653970125e-08\\
-15272.7272727273	1.03623733908574e-08\\
-15225.8064516129	5.00335031990951e-09\\
-15178.8856304985	1.28150776462371e-08\\
-15131.9648093842	1.36430395180245e-08\\
-15085.0439882698	1.1292658318779e-08\\
-15038.1231671554	1.5711651190732e-08\\
-14991.2023460411	1.5013777889799e-08\\
-14944.2815249267	5.22011727344737e-09\\
-14897.3607038123	1.09152127193123e-08\\
-14850.4398826979	1.39708447758081e-08\\
-14803.5190615836	1.31876279689514e-08\\
-14756.5982404692	1.51002516864968e-08\\
-14709.6774193548	1.58292306512751e-08\\
-14662.7565982405	1.00324581736508e-08\\
-14615.8357771261	7.24004684624492e-09\\
-14568.9149560117	1.46994951628982e-08\\
-14521.9941348974	1.46546967366002e-08\\
-14475.073313783	1.42121205432596e-08\\
-14428.1524926686	1.69413700131026e-08\\
-14381.2316715543	1.4755618868879e-08\\
-14334.3108504399	4.66789794407063e-09\\
-14287.3900293255	1.41882332031365e-08\\
-14240.4692082111	1.58761752786376e-08\\
-14193.5483870968	1.49977226623331e-08\\
-14146.6275659824	1.62439111876617e-08\\
-14099.706744868	1.89785948921268e-08\\
-14052.7859237537	5.98168358045544e-09\\
-14005.8651026393	1.20856600801594e-08\\
-13958.9442815249	1.56239875267249e-08\\
-13912.0234604106	1.72467557395218e-08\\
-13865.1026392962	1.70472611820232e-08\\
-13818.1818181818	1.85548519833941e-08\\
-13771.2609970674	1.32801162104692e-08\\
-13724.3401759531	6.65455338773458e-09\\
-13677.4193548387	1.70473204274414e-08\\
-13630.4985337243	1.85651027444051e-08\\
-13583.57771261	1.67808966687737e-08\\
-13536.6568914956	2.03242869409035e-08\\
-13489.7360703812	1.80823238882197e-08\\
-13442.8152492669	5.41834616929053e-09\\
-13395.8944281525	1.58668628005814e-08\\
-13348.9736070381	1.88268518204064e-08\\
-13302.0527859238	1.87188699482414e-08\\
-13255.1319648094	2.19997717647761e-08\\
-13208.211143695	1.95658285266713e-08\\
-13161.2903225806	1.11565254742218e-08\\
-13114.3695014663	1.07927590912689e-08\\
-13067.4486803519	1.9655396496382e-08\\
-13020.5278592375	2.1495295864443e-08\\
-12973.6070381232	2.074253538068e-08\\
-12926.6862170088	2.26503517380257e-08\\
-12879.7653958944	1.76143690319771e-08\\
-12832.8445747801	5.01479846346175e-09\\
-12785.9237536657	2.21488607430761e-08\\
-12739.0029325513	2.19161632922405e-08\\
-12692.082111437	2.23483448885506e-08\\
-12598.2404692082	2.42152099666679e-08\\
-12551.3196480938	6.60012406535348e-09\\
-12504.3988269795	1.64130495857123e-08\\
-12457.4780058651	2.56971885125041e-08\\
-12410.5571847507	2.3860392024481e-08\\
-12363.6363636364	2.54239777586999e-08\\
-12316.715542522	2.57179326157166e-08\\
-12269.7947214076	1.48736159455046e-08\\
-12222.8739002933	1.0597983516623e-08\\
-12175.9530791789	2.57831588177429e-08\\
-12129.0322580645	2.7144524326165e-08\\
-12082.1114369501	2.67322031219518e-08\\
-12035.1906158358	2.89583168220711e-08\\
-11988.2697947214	2.31131320402121e-08\\
-11941.348973607	6.96472953077274e-09\\
-11894.4281524927	2.41178067942771e-08\\
-11847.5073313783	3.16232000473092e-08\\
-11800.5865102639	2.71023301124417e-08\\
-11753.6656891496	3.2311238472463e-08\\
-11706.7448680352	2.88153049102502e-08\\
-11659.8240469208	1.19878460020813e-08\\
-11612.9032258065	1.61905997708234e-08\\
-11565.9824046921	3.47042889564821e-08\\
-11519.0615835777	2.98095127670374e-08\\
-11472.1407624633	3.75123024382737e-08\\
-11378.2991202346	2.4441745509206e-08\\
-11331.3782991202	7.9076304626689e-09\\
-11284.4574780059	3.485340358973e-08\\
-11237.5366568915	3.59558033244945e-08\\
-11190.6158357771	3.61808351086405e-08\\
-11143.6950146628	4.03425609957779e-08\\
-11096.7741935484	3.02826788541389e-08\\
-11049.853372434	9.90029865733861e-09\\
-11002.9325513196	2.88963005537674e-08\\
-10956.0117302053	4.24306486366156e-08\\
-10909.0909090909	4.00429098876157e-08\\
-10862.1700879765	4.47378980757308e-08\\
-10815.2492668622	3.76907577860877e-08\\
-10768.3284457478	1.8932248554338e-08\\
-10721.4076246334	1.72387353917198e-08\\
-10674.4868035191	4.70671482106103e-08\\
-10627.5659824047	4.62185480805263e-08\\
-10580.6451612903	4.93633497427878e-08\\
-10533.724340176	4.55152194891002e-08\\
-10486.8035190616	3.19903008572003e-08\\
-10439.8826979472	1.10745517272777e-08\\
-10392.9618768328	4.6282768305519e-08\\
-10346.0410557185	5.43191833995695e-08\\
-10299.1202346041	5.27905340439602e-08\\
-10252.1994134897	5.68892779083095e-08\\
-10205.2785923754	4.16173005085735e-08\\
-10158.357771261	1.60436233362698e-08\\
-10111.4369501466	3.5333437236351e-08\\
-10064.5161290323	6.27947552676032e-08\\
-10017.5953079179	6.2342063473542e-08\\
-9970.67448680352	6.49505593773051e-08\\
-9923.75366568915	5.44437385046072e-08\\
-9876.83284457478	3.14881848096605e-08\\
-9829.91202346041	1.88463044868344e-08\\
-9782.99120234604	7.32050340215601e-08\\
-9736.07038123167	6.88416478784331e-08\\
-9689.1495601173	7.79439520731478e-08\\
-9642.22873900293	7.05847059854173e-08\\
-9595.30791788856	4.55993006307002e-08\\
-9548.38709677419	1.41484831960854e-08\\
-9501.46627565983	6.90979807464925e-08\\
-9454.54545454546	8.34385459876946e-08\\
-9407.62463343109	8.96027703897166e-08\\
-9360.70381231671	8.83313136683399e-08\\
-9313.78299120235	6.79445330131112e-08\\
-9266.86217008798	2.83855483199704e-08\\
-9219.94134897361	4.20905921914118e-08\\
-9173.02052785924	1.07911362910523e-07\\
-9126.09970674487	1.07920179915439e-07\\
-9079.1788856305	1.01120191924377e-07\\
-9032.25806451613	1.00128760183149e-07\\
-8985.33724340176	4.55672413814161e-08\\
-8938.41642228739	2.22896673780151e-08\\
-8891.49560117302	1.22166506042206e-07\\
-8844.57478005865	1.27936532785726e-07\\
-8750.73313782991	1.32951365869494e-07\\
-8703.81231671554	7.46272662529645e-08\\
-8656.89149560118	2.64929611986854e-08\\
-8609.97067448681	9.32015089559614e-08\\
-8563.04985337243	1.76156057319242e-07\\
-8516.12903225806	1.61993317358226e-07\\
-8469.2082111437	1.69405433543537e-07\\
-8422.28739002933	1.27069793346546e-07\\
-8375.36656891496	4.58932639994873e-08\\
-8328.44574780059	7.0777830080727e-08\\
-8281.52492668622	2.10372043908419e-07\\
-8187.68328445748	2.23713027748611e-07\\
-8140.76246334311	1.92387825278476e-07\\
-8093.84164222874	8.3706389115546e-08\\
-8046.92082111437	4.83876443411985e-08\\
-8000	2.3395701984614e-07\\
-7953.07917888563	2.76149152601123e-07\\
-7906.15835777126	3.11095186798455e-07\\
-7859.23753665689	2.96337275633038e-07\\
-7812.31671554252	1.54912639512525e-07\\
-7765.39589442815	6.38146716065391e-08\\
-7718.47507331378	2.05099686123048e-07\\
-7671.55425219941	4.21457537983705e-07\\
-7624.63343108504	3.9634114356172e-07\\
-7577.71260997068	4.57536295777958e-07\\
-7530.79178885631	3.0426709666391e-07\\
-7483.87096774194	1.07821964485611e-07\\
-7436.95014662757	1.56781071038957e-07\\
-7390.0293255132	6.01424046928942e-07\\
-7343.10850439883	5.8596413506373e-07\\
-7296.18768328446	6.76108430046057e-07\\
-7249.26686217009	6.1710888873066e-07\\
-7202.34604105572	2.12271642269584e-07\\
-7155.42521994135	1.51235336373787e-07\\
-7108.50439882698	7.59083054603577e-07\\
-7061.58357771261	1.03361918052752e-06\\
-7014.66275659824	1.06290454192735e-06\\
-6967.74193548387	1.21910494668306e-06\\
-6920.8211143695	5.62600594927029e-07\\
-6873.90029325513	2.23923744297885e-07\\
-6826.97947214076	7.84086337363442e-07\\
-6780.05865102639	2.1562757619958e-06\\
-6733.13782991202	1.97453180773176e-06\\
-6686.21700879766	2.60659990474006e-06\\
-6639.29618768329	1.80060224479824e-06\\
-6592.37536656891	4.56376707492178e-07\\
-6545.45454545455	1.15742440325113e-06\\
-6498.53372434018	4.48724240620297e-06\\
-6451.61290322581	5.8358966607996e-06\\
-6404.69208211144	7.4087768341554e-06\\
-6357.77126099707	7.37200895749211e-06\\
-6310.8504398827	2.67205564010241e-06\\
-6263.92961876833	3.04239206132891e-06\\
-6217.00879765396	1.86644044475933e-05\\
-6170.08797653959	4.15296108442013e-05\\
-6123.16715542522	7.858288155943e-05\\
-6076.24633431085	0.000178435857235696\\
-6029.32551319648	0.000288351996574989\\
-5982.40469208211	0.000477202636990849\\
-5935.48387096774	0.00029159271267292\\
-5888.56304985337	0.000184250448189294\\
-5841.64222873901	8.11724173205951e-05\\
-5794.72140762463	4.42043899681536e-05\\
-5747.80058651026	2.00423636992604e-05\\
-5700.87976539589	3.19537296015236e-06\\
-5653.95894428153	2.94165191960185e-06\\
-5607.03812316716	8.41436088075064e-06\\
-5560.11730205279	8.41906544547786e-06\\
-5513.19648093842	6.6967411944617e-06\\
-5466.27565982405	5.46692985311638e-06\\
-5419.35483870968	1.31000593117918e-06\\
-5372.43401759531	5.53324560814903e-07\\
-5325.51319648094	2.27228529047155e-06\\
-5278.59237536657	3.30363537777456e-06\\
-5231.6715542522	2.35753209917986e-06\\
-5184.75073313783	2.99344229813634e-06\\
-5137.82991202346	1.03957225832537e-06\\
-5090.90909090909	2.84258665291669e-07\\
-5043.98826979472	8.32017026315036e-07\\
-4997.06744868035	1.71119950157069e-06\\
-4950.14662756598	1.3920819373164e-06\\
-4903.22580645161	1.55354127417014e-06\\
-4856.30498533724	1.22783209424341e-06\\
-4809.38416422287	2.40690166298545e-07\\
-4762.46334310851	3.76155787733959e-07\\
-4715.54252199414	9.87855810041806e-07\\
-4668.62170087976	1.00309467109088e-06\\
-4621.7008797654	8.83271023403924e-07\\
-4574.78005865103	1.20089341702439e-06\\
-4527.85923753666	2.71984061507331e-07\\
-4480.93841642229	2.87966732691493e-07\\
-4434.01759530792	6.07188840196977e-07\\
-4387.09677419355	7.72629845431228e-07\\
-4340.17595307918	6.25114082939412e-07\\
-4293.25513196481	9.58192846378021e-07\\
-4246.33431085044	5.15515513324233e-07\\
-4199.41348973607	2.36023296917446e-07\\
-4152.4926686217	4.34996766931919e-07\\
-4105.57184750733	6.27151993781503e-07\\
-4058.65102639296	5.30414505449077e-07\\
-4011.73020527859	6.84101014043453e-07\\
-3964.80938416422	7.91304654612891e-07\\
-3917.88856304986	2.01404062550897e-07\\
-3870.96774193548	4.66705168642708e-07\\
-3824.04692082111	5.38810761137301e-07\\
-3777.12609970674	4.36751035047882e-07\\
-3730.20527859238	5.29664324394367e-07\\
-3683.28445747801	9.51801115468675e-07\\
-3636.36363636364	3.58719539815892e-07\\
-3589.44281524927	4.94260578264859e-07\\
-3542.5219941349	6.30737882693339e-07\\
-3495.60117302053	4.26722609590499e-07\\
-3448.68035190616	3.86854254909757e-07\\
-3401.75953079179	9.72340286009337e-07\\
-3354.83870967742	7.79702963789538e-07\\
-3307.91788856305	5.45381429405032e-07\\
-3260.99706744868	8.33923590194376e-07\\
-3214.07624633431	5.79280844957501e-07\\
-3167.15542521994	3.31096815752775e-07\\
-3120.23460410557	8.38929178820121e-07\\
-3073.3137829912	1.46756379825546e-06\\
-3026.39296187683	8.18367268667463e-07\\
-2979.47214076246	1.31521811038625e-06\\
-2932.55131964809	1.08414573121172e-06\\
-2885.63049853372	3.06068652922502e-07\\
-2838.70967741936	7.10809199537747e-07\\
-2791.78885630499	2.45777206356903e-06\\
-2744.86803519062	1.83579120335909e-06\\
-2697.94721407625	2.39115734070584e-06\\
-2651.02639296188	2.47855170534526e-06\\
-2604.10557184751	8.13041043068779e-07\\
-2557.18475073314	7.65348692682439e-07\\
-2510.26392961877	3.80773091034487e-06\\
-2463.3431085044	5.98199934776512e-06\\
-2416.42228739003	6.18303009699151e-06\\
-2369.50146627566	8.21687260879674e-06\\
-2322.58064516129	4.42384200609482e-06\\
-2275.65982404692	1.30452412105667e-06\\
-2228.73900293255	1.26120771862e-05\\
-2181.81818181818	3.66036400468008e-05\\
-2134.89736070381	5.54384591356083e-05\\
-2087.97653958944	0.000145709373751773\\
-2041.05571847507	0.000266517255855899\\
-1994.13489736071	0.0014705339666445\\
-1947.21407624634	0.00101234343634565\\
-1900.29325513196	0.000176203545808736\\
-1853.37243401759	0.00011624502702585\\
-1806.45161290323	5.0592904943238e-05\\
-1759.53079178886	2.9869024620017e-05\\
-1665.68914956012	2.00563316144162e-06\\
-1618.76832844575	6.81761469431199e-06\\
-1571.84750733138	9.08474164899482e-06\\
-1524.92668621701	7.27683045345931e-06\\
-1478.00586510264	6.38441653573042e-06\\
-1431.08504398827	3.01085129009052e-06\\
-1384.1642228739	8.18135895659929e-07\\
-1337.24340175953	1.66692498044665e-06\\
-1290.32258064516	3.12433256016748e-06\\
-1243.40175953079	2.95026991183351e-06\\
-1196.48093841642	2.93942842976404e-06\\
-1149.56011730205	2.07096882087833e-06\\
-1102.63929618768	7.67436675737542e-07\\
-1055.71847507331	7.58009221366795e-07\\
-1008.79765395894	1.43823331711502e-06\\
-961.876832844573	1.67267464607066e-06\\
-914.956011730206	1.80541803620733e-06\\
-868.035190615836	1.43166366232957e-06\\
-821.114369501465	9.006838005002e-07\\
-774.193548387098	5.28005986103776e-07\\
-727.272727272728	8.26127240719219e-07\\
-680.351906158357	1.1069996230526e-06\\
-586.51026392962	1.21302633312049e-06\\
-539.58944281525	8.96096171323359e-07\\
-492.668621700879	5.9133452161008e-07\\
-445.747800586509	5.91856488572415e-07\\
-398.826979472142	7.99252489392156e-07\\
-351.906158357771	8.56192583648022e-07\\
-304.985337243401	9.60996360515391e-07\\
-258.064516129034	9.31186058559911e-07\\
-211.143695014664	6.86458753683663e-07\\
-164.222873900293	5.25311268338609e-07\\
-117.302052785923	7.11186037987794e-07\\
-70.3812316715557	6.50465790348401e-07\\
-23.4604105571852	8.06007850932632e-07\\
23.4604105571852	9.45668731638636e-07\\
70.3812316715557	8.06007850932632e-07\\
117.302052785923	6.50465790348401e-07\\
164.222873900293	7.11186037987794e-07\\
211.143695014664	5.25311268338609e-07\\
258.064516129034	6.86458753683663e-07\\
304.985337243401	9.31186058559911e-07\\
351.906158357771	9.60996360515391e-07\\
398.826979472142	8.56192583648022e-07\\
445.747800586509	7.99252489392156e-07\\
492.668621700879	5.91856488572415e-07\\
539.58944281525	5.9133452161008e-07\\
586.51026392962	8.96096171323359e-07\\
633.431085043987	1.21302633312049e-06\\
727.272727272728	1.1069996230526e-06\\
774.193548387098	8.26127240719219e-07\\
821.114369501465	5.28005986103776e-07\\
868.035190615836	9.006838005002e-07\\
914.956011730206	1.43166366232957e-06\\
961.876832844573	1.80541803620733e-06\\
1008.79765395894	1.67267464607066e-06\\
1055.71847507331	1.43823331711502e-06\\
1102.63929618768	7.58009221366795e-07\\
1149.56011730205	7.67436675737542e-07\\
1196.48093841642	2.07096882087833e-06\\
1243.40175953079	2.93942842976404e-06\\
1290.32258064516	2.95026991183351e-06\\
1337.24340175953	3.12433256016748e-06\\
1384.1642228739	1.66692498044665e-06\\
1431.08504398827	8.18135895659929e-07\\
1478.00586510264	3.01085129009052e-06\\
1524.92668621701	6.38441653573042e-06\\
1571.84750733138	7.27683045345931e-06\\
1618.76832844575	9.08474164899482e-06\\
1665.68914956012	6.81761469431199e-06\\
1712.60997067449	2.00563316144162e-06\\
1806.45161290323	2.9869024620017e-05\\
1853.37243401759	5.0592904943238e-05\\
1900.29325513196	0.00011624502702585\\
1947.21407624634	0.000176203545808736\\
1994.13489736071	0.00101234343634565\\
2041.05571847507	0.0014705339666445\\
2087.97653958944	0.000266517255855899\\
2134.89736070381	0.000145709373751773\\
2181.81818181818	5.54384591356083e-05\\
2228.73900293255	3.66036400468008e-05\\
2275.65982404692	1.26120771862e-05\\
2322.58064516129	1.30452412105667e-06\\
2369.50146627566	4.42384200609482e-06\\
2416.42228739003	8.21687260879674e-06\\
2463.3431085044	6.18303009699151e-06\\
2510.26392961877	5.98199934776512e-06\\
2557.18475073314	3.80773091034487e-06\\
2604.10557184751	7.65348692682439e-07\\
2651.02639296188	8.13041043068779e-07\\
2697.94721407625	2.47855170534526e-06\\
2744.86803519062	2.39115734070584e-06\\
2791.78885630499	1.83579120335909e-06\\
2838.70967741936	2.45777206356903e-06\\
2885.63049853372	7.10809199537747e-07\\
2932.55131964809	3.06068652922502e-07\\
2979.47214076246	1.08414573121172e-06\\
3026.39296187683	1.31521811038625e-06\\
3073.3137829912	8.18367268667463e-07\\
3120.23460410557	1.46756379825546e-06\\
3167.15542521994	8.38929178820121e-07\\
3214.07624633431	3.31096815752775e-07\\
3260.99706744868	5.79280844957501e-07\\
3307.91788856305	8.33923590194376e-07\\
3354.83870967742	5.45381429405032e-07\\
3401.75953079179	7.79702963789538e-07\\
3448.68035190616	9.72340286009337e-07\\
3495.60117302053	3.86854254909757e-07\\
3542.5219941349	4.26722609590499e-07\\
3589.44281524927	6.30737882693339e-07\\
3636.36363636364	4.94260578264859e-07\\
3683.28445747801	3.58719539815892e-07\\
3730.20527859238	9.51801115468675e-07\\
3777.12609970674	5.29664324394367e-07\\
3824.04692082111	4.36751035047882e-07\\
3870.96774193548	5.38810761137301e-07\\
3917.88856304986	4.66705168642708e-07\\
3964.80938416422	2.01404062550897e-07\\
4011.73020527859	7.91304654612891e-07\\
4058.65102639296	6.84101014043453e-07\\
4105.57184750733	5.30414505449077e-07\\
4152.4926686217	6.27151993781503e-07\\
4199.41348973607	4.34996766931919e-07\\
4246.33431085044	2.36023296917446e-07\\
4293.25513196481	5.15515513324233e-07\\
4340.17595307918	9.58192846378021e-07\\
4387.09677419355	6.25114082939412e-07\\
4434.01759530792	7.72629845431228e-07\\
4480.93841642229	6.07188840196977e-07\\
4527.85923753666	2.87966732691493e-07\\
4574.78005865103	2.71984061507331e-07\\
4621.7008797654	1.20089341702439e-06\\
4668.62170087976	8.83271023403924e-07\\
4715.54252199414	1.00309467109088e-06\\
4762.46334310851	9.87855810041806e-07\\
4809.38416422287	3.76155787733959e-07\\
4856.30498533724	2.40690166298545e-07\\
4903.22580645161	1.22783209424341e-06\\
4950.14662756598	1.55354127417014e-06\\
4997.06744868035	1.3920819373164e-06\\
5043.98826979472	1.71119950157069e-06\\
5090.90909090909	8.32017026315036e-07\\
5137.82991202346	2.84258665291669e-07\\
5184.75073313783	1.03957225832537e-06\\
5231.6715542522	2.99344229813634e-06\\
5278.59237536657	2.35753209917986e-06\\
5325.51319648094	3.30363537777456e-06\\
5372.43401759531	2.27228529047155e-06\\
5419.35483870968	5.53324560814903e-07\\
5466.27565982405	1.31000593117918e-06\\
5513.19648093842	5.46692985311638e-06\\
5560.11730205279	6.6967411944617e-06\\
5607.03812316716	8.41906544547786e-06\\
5653.95894428153	8.41436088075064e-06\\
5700.87976539589	2.94165191960185e-06\\
5747.80058651026	3.19537296015236e-06\\
5794.72140762463	2.00423636992604e-05\\
5841.64222873901	4.42043899681536e-05\\
5888.56304985337	8.11724173205951e-05\\
5935.48387096774	0.000184250448189294\\
5982.40469208211	0.00029159271267292\\
6029.32551319648	0.000477202636990849\\
6076.24633431085	0.000288351996574989\\
6123.16715542522	0.000178435857235696\\
6170.08797653959	7.858288155943e-05\\
6217.00879765396	4.15296108442013e-05\\
6263.92961876833	1.86644044475933e-05\\
6310.8504398827	3.04239206132891e-06\\
6357.77126099707	2.67205564010241e-06\\
6404.69208211144	7.37200895749211e-06\\
6451.61290322581	7.4087768341554e-06\\
6498.53372434018	5.8358966607996e-06\\
6545.45454545455	4.48724240620297e-06\\
6592.37536656891	1.15742440325113e-06\\
6639.29618768329	4.56376707492178e-07\\
6686.21700879766	1.80060224479824e-06\\
6733.13782991202	2.60659990474006e-06\\
6780.05865102639	1.97453180773176e-06\\
6826.97947214076	2.1562757619958e-06\\
6873.90029325513	7.84086337363442e-07\\
6920.8211143695	2.23923744297885e-07\\
6967.74193548387	5.62600594927029e-07\\
7014.66275659824	1.21910494668306e-06\\
7061.58357771261	1.06290454192735e-06\\
7108.50439882698	1.03361918052752e-06\\
7155.42521994135	7.59083054603577e-07\\
7202.34604105572	1.51235336373787e-07\\
7249.26686217009	2.12271642269584e-07\\
7296.18768328446	6.1710888873066e-07\\
7343.10850439883	6.76108430046057e-07\\
7390.0293255132	5.8596413506373e-07\\
7436.95014662757	6.01424046928942e-07\\
7483.87096774194	1.56781071038957e-07\\
7530.79178885631	1.07821964485611e-07\\
7577.71260997068	3.0426709666391e-07\\
7624.63343108504	4.57536295777958e-07\\
7671.55425219941	3.9634114356172e-07\\
7718.47507331378	4.21457537983705e-07\\
7765.39589442815	2.05099686123048e-07\\
7812.31671554252	6.38146716065391e-08\\
7859.23753665689	1.54912639512525e-07\\
7906.15835777126	2.96337275633038e-07\\
7953.07917888563	3.11095186798455e-07\\
8000	2.76149152601123e-07\\
8046.92082111437	2.3395701984614e-07\\
8093.84164222874	4.83876443411985e-08\\
8140.76246334311	8.3706389115546e-08\\
8187.68328445748	1.92387825278476e-07\\
8234.60410557185	2.23713027748611e-07\\
8328.44574780059	2.10372043908419e-07\\
8375.36656891496	7.0777830080727e-08\\
8422.28739002933	4.58932639994873e-08\\
8469.2082111437	1.27069793346546e-07\\
8516.12903225806	1.69405433543537e-07\\
8563.04985337243	1.61993317358226e-07\\
8609.97067448681	1.76156057319242e-07\\
8656.89149560118	9.32015089559614e-08\\
8703.81231671554	2.64929611986854e-08\\
8750.73313782991	7.46272662529645e-08\\
8797.65395894428	1.32951365869494e-07\\
8891.49560117302	1.27936532785726e-07\\
8938.41642228739	1.22166506042206e-07\\
8985.33724340176	2.22896673780151e-08\\
9032.25806451613	4.55672413814161e-08\\
9079.1788856305	1.00128760183149e-07\\
9126.09970674487	1.01120191924377e-07\\
9173.02052785924	1.07920179915439e-07\\
9219.94134897361	1.07911362910523e-07\\
9266.86217008798	4.20905921914118e-08\\
9313.78299120235	2.83855483199704e-08\\
9360.70381231671	6.79445330131112e-08\\
9407.62463343109	8.83313136683399e-08\\
9454.54545454546	8.96027703897166e-08\\
9501.46627565983	8.34385459876946e-08\\
9548.38709677419	6.90979807464925e-08\\
9595.30791788856	1.41484831960854e-08\\
9642.22873900293	4.55993006307002e-08\\
9689.1495601173	7.05847059854173e-08\\
9736.07038123167	7.79439520731478e-08\\
9782.99120234604	6.88416478784331e-08\\
9829.91202346041	7.32050340215601e-08\\
9876.83284457478	1.88463044868344e-08\\
9923.75366568915	3.14881848096605e-08\\
9970.67448680352	5.44437385046072e-08\\
10017.5953079179	6.49505593773051e-08\\
10064.5161290323	6.2342063473542e-08\\
10111.4369501466	6.27947552676032e-08\\
10158.357771261	3.5333437236351e-08\\
10205.2785923754	1.60436233362698e-08\\
10252.1994134897	4.16173005085735e-08\\
10299.1202346041	5.68892779083095e-08\\
10346.0410557185	5.27905340439602e-08\\
10392.9618768328	5.43191833995695e-08\\
10439.8826979472	4.6282768305519e-08\\
10486.8035190616	1.10745517272777e-08\\
10533.724340176	3.19903008572003e-08\\
10580.6451612903	4.55152194891002e-08\\
10627.5659824047	4.93633497427878e-08\\
10674.4868035191	4.62185480805263e-08\\
10721.4076246334	4.70671482106103e-08\\
10768.3284457478	1.72387353917198e-08\\
10815.2492668622	1.8932248554338e-08\\
10862.1700879765	3.76907577860877e-08\\
10909.0909090909	4.47378980757308e-08\\
10956.0117302053	4.00429098876157e-08\\
11002.9325513196	4.24306486366156e-08\\
11049.853372434	2.88963005537674e-08\\
11096.7741935484	9.90029865733861e-09\\
11143.6950146628	3.02826788541389e-08\\
11190.6158357771	4.03425609957779e-08\\
11237.5366568915	3.61808351086405e-08\\
11284.4574780059	3.59558033244945e-08\\
11331.3782991202	3.485340358973e-08\\
11378.2991202346	7.9076304626689e-09\\
11425.219941349	2.4441745509206e-08\\
11519.0615835777	3.75123024382737e-08\\
11565.9824046921	2.98095127670374e-08\\
11612.9032258065	3.47042889564821e-08\\
11659.8240469208	1.61905997708234e-08\\
11706.7448680352	1.19878460020813e-08\\
11753.6656891496	2.88153049102502e-08\\
11800.5865102639	3.2311238472463e-08\\
11847.5073313783	2.71023301124417e-08\\
11894.4281524927	3.16232000473092e-08\\
11941.348973607	2.41178067942771e-08\\
11988.2697947214	6.96472953077274e-09\\
12035.1906158358	2.31131320402121e-08\\
12082.1114369501	2.89583168220711e-08\\
12129.0322580645	2.67322031219518e-08\\
12175.9530791789	2.7144524326165e-08\\
12222.8739002933	2.57831588177429e-08\\
12269.7947214076	1.0597983516623e-08\\
12316.715542522	1.48736159455046e-08\\
12363.6363636364	2.57179326157166e-08\\
12410.5571847507	2.54239777586999e-08\\
12457.4780058651	2.3860392024481e-08\\
12504.3988269795	2.56971885125041e-08\\
12551.3196480938	1.64130495857123e-08\\
12598.2404692082	6.60012406535348e-09\\
12645.1612903226	2.42152099666679e-08\\
12785.9237536657	2.19161632922405e-08\\
12832.8445747801	2.21488607430761e-08\\
12879.7653958944	5.01479846346175e-09\\
12926.6862170088	1.76143690319771e-08\\
12973.6070381232	2.26503517380257e-08\\
13020.5278592375	2.074253538068e-08\\
13067.4486803519	2.1495295864443e-08\\
13114.3695014663	1.9655396496382e-08\\
13161.2903225806	1.07927590912689e-08\\
13208.211143695	1.11565254742218e-08\\
13255.1319648094	1.95658285266713e-08\\
13302.0527859238	2.19997717647761e-08\\
13348.9736070381	1.87188699482414e-08\\
13395.8944281525	1.88268518204064e-08\\
13442.8152492669	1.58668628005814e-08\\
13489.7360703812	5.41834616929053e-09\\
13536.6568914956	1.80823238882197e-08\\
13583.57771261	2.03242869409035e-08\\
13630.4985337243	1.67808966687737e-08\\
13677.4193548387	1.85651027444051e-08\\
13724.3401759531	1.70473204274414e-08\\
13771.2609970674	6.65455338773458e-09\\
13818.1818181818	1.32801162104692e-08\\
13865.1026392962	1.85548519833941e-08\\
13912.0234604106	1.70472611820232e-08\\
13958.9442815249	1.72467557395218e-08\\
14005.8651026393	1.56239875267249e-08\\
14052.7859237537	1.20856600801594e-08\\
14099.706744868	5.98168358045544e-09\\
14146.6275659824	1.89785948921268e-08\\
14193.5483870968	1.62439111876617e-08\\
14240.4692082111	1.49977226623331e-08\\
14287.3900293255	1.58761752786376e-08\\
14334.3108504399	1.41882332031365e-08\\
14381.2316715543	4.66789794407063e-09\\
14428.1524926686	1.4755618868879e-08\\
14475.073313783	1.69413700131026e-08\\
14521.9941348974	1.42121205432596e-08\\
14568.9149560117	1.46546967366002e-08\\
14615.8357771261	1.46994951628982e-08\\
14662.7565982405	7.24004684624492e-09\\
14709.6774193548	1.00324581736508e-08\\
14756.5982404692	1.58292306512751e-08\\
14803.5190615836	1.51002516864968e-08\\
14850.4398826979	1.31876279689514e-08\\
14897.3607038123	1.39708447758081e-08\\
14944.2815249267	1.09152127193123e-08\\
14991.2023460411	5.22011727344737e-09\\
15038.1231671554	1.5013777889799e-08\\
15085.0439882698	1.5711651190732e-08\\
15131.9648093842	1.1292658318779e-08\\
15178.8856304985	1.36430395180245e-08\\
15225.8064516129	1.28150776462371e-08\\
15272.7272727273	5.00335031990951e-09\\
15319.6480938416	1.03623733908574e-08\\
15366.568914956	1.62296653970125e-08\\
15413.4897360704	1.1649670680257e-08\\
15460.4105571848	1.29309599630819e-08\\
15507.3313782991	1.18486350529457e-08\\
15554.2521994135	8.48223278088209e-09\\
15601.1730205279	6.3435175461757e-09\\
15648.0938416422	1.49051189026328e-08\\
15695.0146627566	1.26037229925183e-08\\
15741.935483871	1.10584351326549e-08\\
15788.8563049853	1.26137140667231e-08\\
15835.7771260997	9.83018098886207e-09\\
15882.6979472141	4.36894984671782e-09\\
15929.6187683284	1.35412413059574e-08\\
15976.5395894428	1.21122211549134e-08\\
16023.4604105572	1.17167204623824e-08\\
16070.3812316716	1.07694236590659e-08\\
16117.3020527859	1.177785823999e-08\\
16164.2228739003	5.75053049556208e-09\\
16211.1436950147	8.15683097146531e-09\\
16258.064516129	1.40422953979675e-08\\
16304.9853372434	1.10780093295583e-08\\
16351.9061583578	1.0016079827131e-08\\
16398.8269794721	1.10751518071119e-08\\
16445.7478005865	8.42181055801457e-09\\
16492.6686217009	5.1115927902975e-09\\
16539.5894428152	1.32237156688076e-08\\
16586.5102639296	1.1218052798535e-08\\
16633.431085044	9.70449248800454e-09\\
16680.3519061584	1.06095544072376e-08\\
16727.2727272727	9.42256788487313e-09\\
16774.1935483871	5.10078749970215e-09\\
16821.1143695015	9.87191406989546e-09\\
16868.0351906158	1.27393529421788e-08\\
16914.9560117302	9.38462079846234e-09\\
16961.8768328446	9.41698279185228e-09\\
17008.7976539589	1.02594059331852e-08\\
17055.7184750733	6.47542298194798e-09\\
17102.6392961877	6.04650512139218e-09\\
17149.5601173021	1.36020161407537e-08\\
17196.4809384164	9.09690165260772e-09\\
17243.4017595308	9.68967061830808e-09\\
17290.3225806452	9.33461396119493e-09\\
17337.2434017595	8.63791786349936e-09\\
17384.1642228739	3.62653836333819e-09\\
17431.0850439883	1.26686029518294e-08\\
17478.0058651026	9.49301728577223e-09\\
17524.926686217	9.45621548026606e-09\\
17571.8475073314	9.07667112276403e-09\\
17618.7683284457	9.10404741606395e-09\\
17665.6891495601	5.20947442802222e-09\\
17712.6099706745	7.47405512703119e-09\\
17759.5307917889	1.22788153457188e-08\\
17806.4516129032	9.22219969912846e-09\\
17853.3724340176	7.50381204518002e-09\\
17900.293255132	9.837130681176e-09\\
17947.2140762463	6.82393296991384e-09\\
17994.1348973607	4.28291575286981e-09\\
18041.0557184751	1.22850308078516e-08\\
18087.9765395894	9.08391947836229e-09\\
18134.8973607038	7.92453273553115e-09\\
18181.8181818182	9.19993340939328e-09\\
18228.7390029326	7.84515715733362e-09\\
18275.6598240469	4.26540136934464e-09\\
18322.5806451613	9.54405326767685e-09\\
18369.5014662757	1.06098991530897e-08\\
18416.42228739	7.69238769636575e-09\\
18510.2639296188	8.61314250193111e-09\\
18557.1847507331	5.94896768734084e-09\\
18604.1055718475	5.18014238619997e-09\\
18651.0263929619	1.26299000267282e-08\\
18697.9472140762	7.48578650664353e-09\\
18744.8680351906	7.87988675077825e-09\\
18838.7096774194	7.88236179316005e-09\\
18885.6304985337	3.77439668496412e-09\\
18932.5513196481	1.08575596667479e-08\\
18979.4721407625	8.9424613602976e-09\\
19026.3929618768	7.23586589970701e-09\\
19073.3137829912	8.184624590396e-09\\
19120.2346041056	7.35863916731919e-09\\
19167.1554252199	4.79203561769186e-09\\
19214.0762463343	8.1092263490313e-09\\
19260.9970674487	1.00226736518046e-08\\
19307.917888563	7.47798387876937e-09\\
19354.8387096774	7.13731911980736e-09\\
19401.7595307918	8.0023760125389e-09\\
19448.6803519062	6.26181933601159e-09\\
19495.6011730205	4.53921628927836e-09\\
19542.5219941349	1.11462946802491e-08\\
19589.4428152493	8.10332102863004e-09\\
19636.3636363636	6.10557076235259e-09\\
19683.284457478	8.49843901656963e-09\\
19730.2052785924	6.67165259273348e-09\\
19777.1260997067	4.16381658805751e-09\\
19824.0469208211	9.34003373794343e-09\\
19870.9677419355	9.30667100277566e-09\\
19917.8885630499	6.16103856524798e-09\\
19964.8093841642	8.05751524749803e-09\\
20058.651026393	5.76009655876556e-09\\
20105.5718475073	5.38111297235994e-09\\
20152.4926686217	1.13600808867735e-08\\
20199.4134897361	6.5355028140615e-09\\
20246.3343108504	6.39923402437033e-09\\
20293.2551319648	7.83829170368848e-09\\
20340.1759530792	6.49254240605003e-09\\
20387.0967741935	4.04119336604137e-09\\
20434.0175953079	1.00347534616679e-08\\
20480.9384164223	8.29237347580613e-09\\
20527.8592375367	6.12350410724847e-09\\
20574.780058651	6.92432958599583e-09\\
20621.7008797654	7.31738091003574e-09\\
20668.6217008798	4.17427421772841e-09\\
20715.5425219941	7.81275897819787e-09\\
20762.4633431085	9.46404796303222e-09\\
20809.3841642229	6.20316276135553e-09\\
20856.3049853372	6.57611067176519e-09\\
20903.2258064516	7.35433742424774e-09\\
20950.146627566	5.39245820103201e-09\\
20997.0674486804	5.16473754216914e-09\\
21043.9882697947	1.02478900824027e-08\\
21090.9090909091	7.21158755747846e-09\\
21137.8299120235	5.61050062894851e-09\\
21184.7507331378	7.34535366459983e-09\\
21231.6715542522	6.48938712799215e-09\\
21278.5923753666	4.38861331152365e-09\\
21325.5131964809	8.3927326533293e-09\\
21372.4340175953	9.08864510372774e-09\\
21419.3548387097	5.312023953531e-09\\
21466.275659824	7.05928626997468e-09\\
21513.1964809384	6.55663995266352e-09\\
21560.1173020528	5.22968410751253e-09\\
21607.0381231672	6.09338151824363e-09\\
21653.9589442815	1.01374113790018e-08\\
21700.8797653959	5.76766937967642e-09\\
21747.8005865103	6.1253922599098e-09\\
21794.7214076246	7.33156300299514e-09\\
21841.642228739	5.66631490157603e-09\\
21888.5630498534	4.17547573945203e-09\\
21935.4838709677	1.03817025551138e-08\\
21982.4046920821	6.65084162185285e-09\\
22029.3255131965	5.93921942829406e-09\\
22076.2463343109	6.2652016493601e-09\\
22123.1671554252	7.16299642360804e-09\\
22170.0879765396	4.19904953100462e-09\\
22217.008797654	7.36708150156432e-09\\
22263.9296187683	9.45308701637496e-09\\
22310.8504398827	5.27777615739597e-09\\
22357.7712609971	5.95385653238616e-09\\
22404.6920821114	7.30929459174194e-09\\
22451.6129032258	4.95811004243854e-09\\
22498.5337243402	5.47687574771774e-09\\
22545.4545454545	1.01587715494472e-08\\
22592.3753665689	5.89794938544493e-09\\
22639.2961876833	5.76473477928261e-09\\
22686.2170087977	6.84046426932177e-09\\
22733.137829912	6.00326646783531e-09\\
22780.0586510264	4.8690987647466e-09\\
22826.9794721408	8.21234650072718e-09\\
22873.9002932551	8.45633150931201e-09\\
22920.8211143695	5.0958738899701e-09\\
22967.7419354839	6.05776913483648e-09\\
23014.6627565982	6.89312843150773e-09\\
23061.5835777126	5.02480237057881e-09\\
23108.504398827	6.06421456600552e-09\\
23155.4252199413	1.01276976184229e-08\\
23202.3460410557	4.80809648060119e-09\\
23249.2668621701	6.3417593009874e-09\\
23296.1876832845	6.40325848302216e-09\\
23343.1085043988	5.99102583098236e-09\\
23390.0293255132	4.20392279312558e-09\\
23436.9501466276	1.06515815446735e-08\\
23483.8709677419	5.72810055214935e-09\\
23530.7917888563	5.55926636051907e-09\\
23577.7126099707	6.60837217254921e-09\\
23624.633431085	6.29998560168605e-09\\
23671.5542521994	4.74083020421153e-09\\
23718.4750733138	7.51908440929027e-09\\
23765.3958944282	8.98536656231893e-09\\
23812.3167155425	5.05746122659395e-09\\
23859.2375366569	5.64973238778839e-09\\
23906.1583577713	7.06774407895119e-09\\
23953.0791788856	5.42413927183039e-09\\
24000	5.34393847255656e-09\\
};
\addlegendentry{Expérimentale}

\end{axis}
\end{tikzpicture}%


\section{Canal de transmission à bruit additif, blanc et Gaussien}

Dans cette section, nous allons tenter de simuler un bruit blanc Gaussion que nous additionnerons à notre signal modulé en fréquence afin de modéliser le signal reçu par le modem.
Le bruit simulé sera généré aléatoirement grâce au module \verb|rand| de Matlab et sera de puissance $\sigma^2$ avec:

\[
\sigma=\sqrt{\frac{S_\text{module}}{10^{\operatorname{SNR}/10}}} \huge
\]

avec $S_\text{module}$ représentant la densité spectrale de puissance du signal modulé en fréquence et $\operatorname{SNR}$ le rapport signal sur bruit (signal to noise ratio) que nous fixerons à 10 par la suite.

%%%% This file was created by matlab2tikz.
%
%The latest updates can be retrieved from
%  http://www.mathworks.com/matlabcentral/fileexchange/22022-matlab2tikz-matlab2tikz
%where you can also make suggestions and rate matlab2tikz.
%
\definecolor{mycolor1}{rgb}{0.00000,0.44700,0.74100}%
%
\begin{tikzpicture}

\begin{axis}[%
width=4.521in,
height=3.559in,
at={(0.758in,0.488in)},
scale only axis,
xmin=0,
xmax=4000,
xlabel style={font=\color{white!15!black}},
xlabel={Temps [s]},
ymin=-1,
ymax=1,
ylabel style={font=\color{white!15!black}},
ylabel={Amplitude},
axis background/.style={fill=white},
title style={font=\bfseries},
title={Signal bruité (SNR = 5000 dB)}
]
\addplot [color=mycolor1, forget plot]
  table[row sep=crcr]{%
1	0.981603928409186\\
2	0.997963789571259\\
3	0.997923473341416\\
4	0.981483642261537\\
5	0.948914462592711\\
6	0.900751164501344\\
7	0.837785246276326\\
8	0.761051467133711\\
9	0.671810842356081\\
10	0.57152992022111\\
11	0.461856681276004\\
12	0.344593456019084\\
13	0.221667306048403\\
14	0.0950983554284903\\
16	-0.160622310922918\\
17	-0.285571605349105\\
18	-0.40582791889301\\
19	-0.519415002620462\\
20	-0.624466207319529\\
21	-0.719255159336626\\
22	-0.802224131183721\\
23	-0.872009640667329\\
24	-0.927464857865743\\
25	-0.967678451717802\\
26	-0.991989566509346\\
27	-0.999998682138084\\
28	-0.991574179683084\\
29	-0.966854504383264\\
30	-0.926245890476366\\
31	-0.87041568529412\\
32	-0.800281382314552\\
33	-0.716995543406938\\
34	-0.62192685804348\\
35	-0.516637650749544\\
36	-0.402858206426572\\
37	-0.282458335472711\\
38	-0.157416645991816\\
41	0.224832337547468\\
42	0.347639745867127\\
43	0.464734167843289\\
44	0.574191315929966\\
45	0.674212410788641\\
46	0.763153741779206\\
47	0.839553679110395\\
48	0.902156693736288\\
49	0.949933990267709\\
50	0.98210041382481\\
51	0.998127352992469\\
52	0.997751426825289\\
53	0.980978813158799\\
54	0.94808514708302\\
55	0.89961099125685\\
56	0.836352952493598\\
57	0.759350590612485\\
58	0.669869334688428\\
59	0.56937968744478\\
60	0.459533059542082\\
61	0.342134630889632\\
62	0.219113684970125\\
63	0.0924919036897336\\
65	-0.163205759226003\\
66	-0.288079597586147\\
67	-0.408219239622667\\
68	-0.521650353739915\\
69	-0.626508853873474\\
70	-0.721071533207578\\
71	-0.803784382737831\\
72	-0.873288129293087\\
73	-0.928440573341049\\
74	-0.968335359486446\\
75	-0.992316871185267\\
76	-0.999991004914591\\
77	-0.991231646725282\\
78	-0.966182744753496\\
79	-0.925255943630873\\
80	-0.869123819662491\\
81	-0.798708827951486\\
82	-0.715168143104165\\
83	-0.619874642638933\\
84	-0.51439434561189\\
85	-0.400460677224601\\
86	-0.279945982337267\\
87	-0.154830756029696\\
90	0.227382530744308\\
91	0.35009325656074\\
92	0.467050675923019\\
93	0.576332752736562\\
94	0.676143584721103\\
95	0.764842916620637\\
96	0.840973095570462\\
97	0.903283025636028\\
98	0.95074872787427\\
99	0.982590168034221\\
100	0.998284075360061\\
101	0.997572541832596\\
102	0.980467260536898\\
103	0.94724933350426\\
104	0.898464652179882\\
105	0.834914926441343\\
106	0.757644509587863\\
107	0.667923235811941\\
108	0.567225552204036\\
109	0.457206288219822\\
110	0.339673460807717\\
111	0.216558562111004\\
112	0.089884818021801\\
114	-0.165788088934733\\
115	-0.290585615357486\\
116	-0.410607762463314\\
117	-0.5238821295261\\
118	-0.628547206405983\\
119	-0.722882964935252\\
120	-0.805339125244245\\
121	-0.874560632499652\\
122	-0.929409925388882\\
123	-0.968985630393036\\
124	-0.992637374632977\\
125	-0.99997647386499\\
126	-0.990882319976208\\
127	-0.965504363015043\\
128	-0.924259655184414\\
129	-0.867825997152977\\
130	-0.797130799327078\\
131	-0.713335841118351\\
132	-0.617818178682683\\
133	-0.51214751487305\\
134	-0.398060403309955\\
135	-0.277431710483597\\
136	-0.152243804874615\\
138	0.103539357279715\\
139	0.229931165486505\\
140	0.352544367754035\\
141	0.469363982889718\\
142	0.578470239423041\\
143	0.678070124441092\\
144	0.766526849314687\\
145	0.842386748095123\\
146	0.904403166534848\\
147	0.951556949155474\\
148	0.98307318768093\\
149	0.998433955600831\\
150	0.997386819589337\\
151	0.979948987902389\\
152	0.946407027583973\\
153	0.897312155126201\\
154	0.833471177975753\\
155	0.755933235752309\\
156	0.665972559064357\\
157	0.56506752926316\\
158	0.454876383257215\\
159	0.337209962641737\\
160	0.21400195498336\\
161	0.0872771162930803\\
163	-0.168369282350795\\
164	-0.293089641487768\\
165	-0.412993471044501\\
166	-0.526110314683137\\
167	-0.630581250946307\\
168	-0.724689442104136\\
169	-0.806888348045959\\
170	-0.875827141565424\\
171	-0.930372907364927\\
172	-0.969629259980138\\
173	-0.992951074654684\\
174	-0.999955089088417\\
175	-0.990526201830562\\
176	-0.964819363818151\\
177	-0.923257031965932\\
178	-0.8665222266618\\
179	-0.795547307257493\\
180	-0.711498650009162\\
181	-0.615757480270076\\
182	-0.509897173932586\\
183	-0.395657401134031\\
184	-0.274915537144352\\
185	-0.149655810257627\\
187	0.106142922577874\\
188	0.232478224305851\\
189	0.354993062648191\\
190	0.471674072888163\\
191	0.580603761339262\\
192	0.679992016744109\\
193	0.768205528319413\\
194	0.843794626995532\\
195	0.905517108756158\\
196	0.952358648572044\\
197	0.983549469453465\\
198	0.998576992686594\\
199	0.99719426136835\\
200	0.979423998807306\\
201	0.945558235096541\\
202	0.896153507995678\\
203	0.832021716992131\\
204	0.754216780835122\\
205	0.664017317815251\\
206	0.562905633412356\\
207	0.452543360622258\\
208	0.334744153276915\\
209	0.211443881110881\\
210	0.0846688163769613\\
212	-0.170949321782246\\
213	-0.295591658813919\\
214	-0.415376349014878\\
215	-0.528334893938791\\
216	-0.632610973553255\\
217	-0.726490952332369\\
218	-0.808432040525531\\
219	-0.877087647810185\\
220	-0.931329512669436\\
221	-0.970266243837159\\
222	-0.993257969101251\\
223	-0.999926850731754\\
224	-0.990163294729427\\
225	-0.964127751857177\\
226	-0.922248080847112\\
227	-0.865212517124291\\
228	-0.793958362595731\\
229	-0.709656582367188\\
230	-0.613692561524658\\
231	-0.507643338214166\\
232	-0.393251687166412\\
233	-0.272397479564461\\
234	-0.147066789917062\\
236	0.108745760384863\\
237	0.235023689745958\\
238	0.357439324459392\\
239	0.473980930085418\\
240	0.582733303861914\\
241	0.681909248458396\\
242	0.769878942129708\\
243	0.845196722621495\\
244	0.906624844664293\\
245	0.953153820629268\\
246	0.984019010088559\\
247	0.998713185637826\\
248	0.996994868489764\\
249	0.978892296850063\\
250	0.944702961858184\\
251	0.894988718729564\\
252	0.830566553424887\\
253	0.752495156600162\\
254	0.662057525465571\\
255	0.560739879469565\\
256	0.450207236306142\\
257	0.332276049612574\\
258	0.208884358024989\\
259	0.0820599361504719\\
261	-0.173528189546232\\
262	-0.29809165018878\\
263	-0.417756380041737\\
264	-0.530555852046291\\
265	-0.634636360315653\\
266	-0.728287483273107\\
267	-0.80997019210281\\
268	-0.878342142594192\\
269	-0.932279734745407\\
270	-0.970896577598069\\
271	-0.993558055869016\\
272	-0.999891758988269\\
273	-0.989793601160272\\
274	-0.963429531872862\\
275	-0.921232808742843\\
276	-0.863896877517163\\
277	-0.792363976232537\\
278	-0.707809650818945\\
279	-0.611623436599075\\
280	-0.505386023164647\\
281	-0.390843277895783\\
282	-0.269877555002495\\
283	-0.144476761596707\\
285	0.111347852860035\\
286	0.237567544359536\\
287	0.359883136421104\\
288	0.476284538669915\\
289	0.584858852396337\\
290	0.68382180644312\\
291	0.771547079275479\\
292	0.8465930253642\\
293	0.907726366666793\\
294	0.953942459876544\\
295	0.984481806367057\\
296	0.998842533520474\\
297	0.996788642320098\\
298	0.978353885674096\\
299	0.943841213731957\\
300	0.893817795310497\\
301	0.829105697247087\\
302	0.750768374848121\\
303	0.660093195448098\\
304	0.558570282278652\\
305	0.447868026320066\\
306	0.329805668565768\\
307	0.206323403269835\\
308	0.0794504934942779\\
309	-0.0487280770685175\\
310	-0.176105867967635\\
311	-0.300589598476108\\
312	-0.420133547813748\\
313	-0.532773173783426\\
314	-0.636657397351428\\
315	-0.730079022612699\\
316	-0.811502792234478\\
317	-0.879590617319991\\
318	-0.933223567080859\\
319	-0.971520256942767\\
320	-0.993851332901158\\
321	-0.999849814098525\\
322	-0.98941712365604\\
323	-0.962724708650057\\
324	-0.920211222612579\\
325	-0.862575316857146\\
326	-0.790764159095033\\
327	-0.705957868022779\\
328	-0.609550119674623\\
329	-0.503125244256353\\
330	-0.388432189829018\\
331	-0.267355780730213\\
332	-0.141885743048533\\
334	0.113949182169563\\
335	0.240109770712479\\
336	0.362324481784526\\
337	0.478584882854193\\
338	0.586980392373334\\
339	0.685729677589279\\
340	0.773209928324377\\
341	0.847983525653035\\
342	0.908821667214852\\
343	0.954724560909199\\
344	0.984937855118005\\
345	0.998965035448236\\
346	0.996575584272705\\
347	0.97780876897059\\
348	0.942972996623666\\
349	0.892640745764766\\
350	0.827639158471811\\
351	0.7490364474138\\
352	0.65812434122563\\
353	0.556396856709398\\
354	0.445525746696603\\
355	0.327333027067652\\
356	0.203761034396848\\
357	0.0768405062931379\\
358	-0.05134279100821\\
359	-0.178682339379066\\
360	-0.303085486556483\\
361	-0.422507836036857\\
362	-0.534986843952538\\
363	-0.63867407080852\\
364	-0.731865558072968\\
365	-0.813029830417236\\
366	-0.880833063429691\\
367	-0.9341610032061\\
368	-0.972137277596175\\
369	-0.994137798187239\\
370	-0.999801016350375\\
371	-0.989033864797875\\
372	-0.96201328702\\
373	-0.91918332945761\\
374	-0.861247844202808\\
375	-0.789158922148999\\
376	-0.704101246670234\\
377	-0.607472624962156\\
378	-0.50086101698389\\
379	-0.386018439491636\\
380	-0.264832174031653\\
381	-0.13929375203179\\
383	0.116549730484621\\
384	0.242650351379325\\
385	0.364763343816776\\
386	0.480881946870795\\
387	0.589097909252359\\
388	0.687632848821522\\
389	0.774867477878615\\
390	0.849368213957405\\
391	0.909910738801045\\
392	0.955500118366672\\
393	0.985387153215015\\
394	0.99908069058165\\
395	0.996355695807779\\
396	0.977256950474839\\
397	0.94209831648368\\
398	0.891457578159134\\
399	0.826166947150796\\
400	0.747299386167015\\
401	0.656150976292793\\
402	0.554219617657964\\
403	0.443180413489699\\
404	0.32485814206575\\
405	0.201197268968372\\
406	0.0742299924359031\\
407	-0.0539571530503054\\
408	-0.181257586121774\\
409	-0.305579297322765\\
410	-0.424879228439295\\
411	-0.537196847381892\\
412	-0.640686366864884\\
413	-0.733647077408023\\
414	-0.814551296184163\\
415	-0.882069472408602\\
416	-0.935092036696915\\
417	-0.972747635329597\\
418	-0.994417449764114\\
419	-0.999745366077605\\
420	-0.9886438272124\\
421	-0.961295271858944\\
422	-0.918149136322882\\
423	-0.859914468651823\\
424	-0.787548276395682\\
425	-0.702239799486506\\
426	-0.605390966700725\\
427	-0.498593356865968\\
428	-0.38360204342689\\
429	-0.262306752202676\\
430	-0.136700806311183\\
432	0.119149479980479\\
433	0.245189268948252\\
434	0.36719970580134\\
435	0.483175714976824\\
436	0.591211388520605\\
437	0.689531307094967\\
438	0.77651971657815\\
439	0.850747080787642\\
440	0.910993573960695\\
441	0.956269126933421\\
442	0.985829697578538\\
443	0.999189498128089\\
444	0.996128978432353\\
445	0.976698433969887\\
446	0.941217179307841\\
447	0.890268300603111\\
448	0.82468907337352\\
449	0.745557203014414\\
450	0.654173114174228\\
451	0.552038580047792\\
452	0.440832042773764\\
453	0.322381030522592\\
454	0.198632124556752\\
455	0.0716189698146081\\
456	-0.0565711452759388\\
457	-0.183831590545196\\
458	-0.308071013682365\\
459	-0.427247708766572\\
460	-0.539403168924309\\
461	-0.642694271728942\\
462	-0.735423568408805\\
463	-0.816067179108359\\
464	-0.883299835782054\\
465	-0.936016661171379\\
466	-0.973351325959356\\
467	-0.994690285715478\\
468	-0.999682863662201\\
469	-0.988247013572618\\
470	-0.960570668087257\\
471	-0.917108650296541\\
472	-0.858575199343704\\
473	-0.785932232874529\\
474	-0.700373539229531\\
475	-0.603305159157117\\
476	-0.496322279444939\\
477	-0.381183018197135\\
478	-0.259779532552784\\
479	-0.134106923658692\\
481	0.121748412839224\\
482	0.247726506016988\\
483	0.369633551040351\\
484	0.485466171450298\\
485	0.593320815692095\\
486	0.691425039397927\\
487	0.778166633098408\\
488	0.852120116692276\\
489	0.91207016527278\\
490	0.95703158133847\\
491	0.986265485175863\\
492	0.999291457341315\\
493	0.995895433700298\\
494	0.976133223282886\\
495	0.940329591134287\\
496	0.889072921248044\\
497	0.823205547269936\\
498	0.743809909895845\\
499	0.652190768426863\\
500	0.549853758826885\\
501	0.438480650644578\\
502	0.319901709416172\\
503	0.196065618742523\\
504	0.0690074563249254\\
505	-0.0591847497694289\\
506	-0.18640433500741\\
507	-0.310560618558156\\
508	-0.429613260786027\\
509	-0.54160579345762\\
510	-0.644697771638675\\
511	-0.737195018898547\\
512	-0.817577468798845\\
513	-0.884524145117211\\
514	-0.936934870292134\\
515	-0.973948345348163\\
516	-0.994956304170955\\
517	-0.999613509532537\\
518	-0.987843426598829\\
519	-0.95983948067169\\
520	-0.916061878510391\\
521	-0.857230045457527\\
522	-0.784310802662276\\
523	-0.698502478691353\\
524	-0.601215216627224\\
525	-0.494047800286808\\
526	-0.378761380381093\\
527	-0.25725053240285\\
528	-0.131512121852211\\
530	0.12434651124795\\
531	0.250262045195996\\
532	0.372064862852767\\
533	0.487753300592885\\
534	0.595426176309047\\
535	0.693314032751005\\
536	0.77980821615165\\
537	0.853487312261223\\
538	0.91314050535766\\
539	0.957787476356316\\
540	0.986694513020211\\
541	0.999386567522834\\
542	0.99565506321278\\
543	0.975561322287831\\
544	0.939435558047535\\
545	0.887871448286205\\
546	0.821716379007739\\
547	0.742057518787533\\
548	0.65020395263673\\
549	0.54766516897007\\
550	0.436126253217935\\
551	0.317420195738578\\
552	0.19349776911622\\
553	0.0663954698657108\\
554	-0.0617979486169133\\
555	-0.188975801875586\\
556	-0.313048094886199\\
557	-0.431975868284098\\
558	-0.543804705885577\\
559	-0.646696852861623\\
560	-0.738961416735947\\
561	-0.819082154905118\\
562	-0.885742392022621\\
563	-0.937846657766386\\
564	-0.974538689404199\\
565	-0.995215503307008\\
566	-0.999537304163823\\
567	-0.987433069056351\\
568	-0.959101714624012\\
569	-0.91500882813898\\
570	-0.855879016211929\\
571	-0.782683996871128\\
572	-0.696626630694936\\
573	-0.59912115343559\\
574	-0.491769934979857\\
575	-0.376337146577498\\
576	-0.254719769086933\\
577	-0.128916418675999\\
579	0.126943757399658\\
580	0.252795869107103\\
581	0.374493624573915\\
582	0.490037086729444\\
583	0.597527455941417\\
584	0.695198274207542\\
585	0.781444454486973\\
586	0.854848658123501\\
587	0.914204586879805\\
588	0.958536806806023\\
589	0.98711677817073\\
590	0.999474828020539\\
591	0.995407868616894\\
592	0.974982734905097\\
593	0.938535086174397\\
594	0.886663889953525\\
595	0.820221578793735\\
596	0.740300041700266\\
597	0.648212680421238\\
598	0.545472825477191\\
599	0.433768866630544\\
600	0.31493650649918\\
601	0.190928593278386\\
602	0.0637830283394578\\
603	-0.0644107239086225\\
604	-0.191545973524171\\
605	-0.315533425617559\\
606	-0.434335515068142\\
607	-0.545999891137399\\
608	-0.648691501696703\\
609	-0.740722749814722\\
610	-0.820581227114417\\
611	-0.886954568149122\\
612	-0.938752017344996\\
613	-0.975122354081577\\
614	-0.995467881347849\\
615	-0.999454248078109\\
616	-0.987015943758706\\
617	-0.95835737500056\\
618	-0.913949506399149\\
619	-0.854522120867841\\
620	-0.781051826651947\\
621	-0.694746008097809\\
622	-0.597022983934949\\
623	-0.489488699137382\\
624	-0.373910333401\\
625	-0.252187259950006\\
626	-0.126319831921592\\
628	0.129540133493265\\
629	0.255327960383966\\
630	0.37691981955777\\
631	0.492317514206206\\
632	0.599624640187812\\
633	0.697077750853168\\
634	0.783075336889397\\
635	0.856204144949061\\
636	0.915262402546432\\
637	0.959279567551675\\
638	0.987532277732953\\
639	0.99955623823007\\
640	0.995153851607029\\
641	0.998723106531997\\
642	0.979060572355593\\
643	0.936085626590284\\
644	0.870821545682702\\
645	0.784822332683234\\
646	0.68013571487927\\
647	0.559254385359509\\
648	0.425056649498401\\
649	0.280737889625016\\
650	0.129734479775834\\
652	-0.177870566152251\\
653	-0.327147818387857\\
654	-0.468635353819536\\
655	-0.598964212697865\\
656	-0.715031134523542\\
657	-0.814072449785272\\
658	-0.893729885694938\\
659	-0.95210671902214\\
660	-0.987812938970364\\
661	-0.999998344725554\\
662	-0.988372789587174\\
663	-0.953213089652309\\
664	-0.895356432547032\\
665	-0.816180443151552\\
666	-0.717570380973484\\
667	-0.601874250233323\\
668	-0.471846891536188\\
669	-0.330584386359988\\
670	-0.181450336258422\\
672	0.126125399308421\\
673	0.277243398839346\\
674	0.42175995571597\\
675	0.556233986157167\\
676	0.677463528983481\\
677	0.782561987512963\\
678	0.869026862317241\\
679	0.934799338236189\\
680	0.97831330682493\\
681	0.998532656964926\\
682	0.994975945703118\\
683	0.967727861892399\\
684	0.917437209665422\\
685	0.845301459756683\\
686	0.753038236530756\\
687	0.642844419623088\\
688	0.517343834042549\\
689	0.379524774270521\\
690	0.232668849982019\\
691	0.0802728476351149\\
692	-0.0740345315116429\\
693	-0.22657907440589\\
694	-0.373728542889694\\
695	-0.511979160924056\\
696	-0.63803904300039\\
697	-0.748906577227899\\
698	-0.84194189671507\\
699	-0.914929737440616\\
700	-0.966132185907554\\
701	-0.994330060598713\\
702	-0.998851941911653\\
703	-0.979590159329291\\
704	-0.937003355162233\\
705	-0.872105563817513\\
706	-0.786442066624659\\
707	-0.682052597142047\\
708	-0.561422773059348\\
709	-0.427424911155413\\
710	-0.283249634559979\\
711	-0.132329900829518\\
713	0.175293712478378\\
714	0.324672766104868\\
715	0.466321036397858\\
716	0.596865736350082\\
717	0.71319846607048\\
718	0.812549226896408\\
719	0.89255237783118\\
720	0.951302963823309\\
721	0.987402074651527\\
722	0.999990154376974\\
723	0.98876746822998\\
724	0.954001239592799\\
725	0.896519287146475\\
726	0.817690313679122\\
727	0.719391315901248\\
728	0.60396289127948\\
729	0.474153506068433\\
730	0.333054051578529\\
731	0.184024246963872\\
733	-0.123527882669805\\
734	-0.274727083835842\\
735	-0.419384758324668\\
736	-0.554056462203789\\
737	-0.675535527491775\\
738	-0.780929416126583\\
739	-0.867728594194432\\
740	-0.933866286439297\\
741	-0.977767688251333\\
742	-0.998387463337622\\
743	-0.995234634228382\\
744	-0.968384272937328\\
745	-0.918475713426233\\
746	-0.846697328420305\\
747	-0.754758233067605\\
748	-0.644847589195706\\
749	-0.519582479179462\\
750	-0.381945590591386\\
751	-0.235214195428853\\
752	-0.0828821149866599\\
753	0.0714234715210296\\
754	0.224028393726257\\
755	0.371298975779609\\
756	0.509728557807193\\
757	0.636020992987596\\
758	0.747169132108866\\
759	0.840526426796259\\
760	0.913869946477007\\
761	0.965453308581345\\
762	0.994048261658463\\
763	0.998973931272303\\
764	0.980113032301688\\
765	0.937914661617469\\
766	0.873383604638093\\
767	0.78805641038025\\
768	0.68396480469255\\
769	0.56358731283035\\
770	0.429790243289517\\
771	0.285759438133482\\
772	0.13492441490871\\
774	-0.17271565736155\\
775	-0.322195488551188\\
776	-0.464003522864004\\
777	-0.594763169151065\\
778	-0.711360909435371\\
779	-0.811020434886359\\
780	-0.891368752513245\\
781	-0.950492688500162\\
782	-0.986984442789435\\
783	-0.999975110208197\\
784	-0.989155369970831\\
785	-0.954782850915308\\
786	-0.897675997103306\\
787	-0.819194579848954\\
788	-0.721207320201756\\
789	-0.606047392831897\\
790	-0.476456870805123\\
791	-0.335521434081329\\
792	-0.186596896387528\\
793	-0.0332292999369201\\
794	0.120929519385299\\
795	0.272208885883629\\
796	0.417006686517198\\
797	0.551875140809443\\
798	0.6736028959549\\
799	0.779291492337961\\
800	0.866424378757074\\
801	0.932926834027967\\
802	0.977215368166981\\
803	0.998235426874999\\
804	0.995486501526557\\
805	0.9690340467846\\
806	0.919507922057164\\
807	0.848087393915648\\
808	0.756473056576397\\
809	0.646846339055173\\
810	0.521817563156219\\
811	0.384363789100007\\
812	0.237757928744031\\
813	0.0854908142737258\\
814	-0.0688119220017143\\
815	-0.22147617758128\\
816	-0.368866863828316\\
817	-0.507474461067432\\
818	-0.633998583758512\\
819	-0.745426565976686\\
820	-0.839105196003402\\
821	-0.912803891950716\\
822	-0.964767814145944\\
823	-0.993759649622916\\
824	-0.999089073777213\\
825	-0.980629187688919\\
826	-0.938819539711403\\
827	-0.874655659385553\\
828	-0.789665352885095\\
829	-0.685872324424963\\
830	-0.565747989837291\\
831	-0.432152629689881\\
832	-0.288267283144251\\
833	-0.137518004230969\\
835	0.17013641847052\\
836	0.319716002705263\\
837	0.461682829101846\\
838	0.592656525512211\\
839	0.709518477211986\\
840	0.809486084232958\\
841	0.89017901785337\\
842	0.949675898606529\\
843	0.986560046246268\\
844	0.999953212321998\\
845	0.989536492151728\\
846	0.955557918263366\\
847	0.89882655448946\\
848	0.82069323135056\\
849	0.723018381428119\\
850	0.608127740603322\\
851	0.478756969960159\\
852	0.337986516958154\\
853	0.189168266897468\\
854	0.0358457311749589\\
855	-0.118330327262811\\
856	-0.269688822242188\\
857	-0.414625756592613\\
858	-0.549690036924858\\
859	-0.671665647619648\\
860	-0.777648227372083\\
861	-0.865114224944136\\
862	-0.931980987440511\\
863	-0.976656350358098\\
864	-0.998076548617973\\
865	-0.995731545871877\\
866	-0.969677178981328\\
867	-0.920533828484167\\
868	-0.849471646714392\\
869	-0.758182695303731\\
870	-0.648840655502227\\
871	-0.524049070653746\\
872	-0.386779353222664\\
873	-0.240300032492996\\
874	-0.0880989276161017\\
875	0.0661999008534622\\
876	0.218922443463271\\
877	0.366432223704578\\
878	0.505216886154813\\
879	0.631971829173835\\
880	0.743678890773936\\
881	0.837678214077641\\
882	0.911731581168624\\
883	0.964075707300253\\
884	0.993464226469769\\
885	0.999197368637397\\
886	0.981138621953505\\
887	0.939717983240826\\
888	0.875921719340568\\
889	0.79126888311248\\
890	0.687775143265299\\
891	0.56790478927087\\
892	0.434512054164315\\
893	0.290773152402835\\
894	0.140110651020223\\
896	-0.167556013484045\\
897	-0.317234325561913\\
898	-0.459358971017082\\
899	-0.590545819872204\\
900	-0.707671182028662\\
901	-0.807946185452693\\
902	-0.888983182006541\\
903	-0.948852599740803\\
904	-0.9861288879315\\
905	-0.999924460868442\\
906	-0.989910832160604\\
907	-0.956326436324161\\
908	-0.899970951418709\\
909	-0.822186257913017\\
910	-0.724824487167098\\
911	-0.610203920334698\\
912	-0.48105378776836\\
913	-0.340449283313319\\
914	-0.191738340868596\\
915	-0.0384619167307392\\
916	0.115730324117976\\
917	0.267166910183732\\
918	0.41224198486907\\
919	0.547501165526228\\
920	0.66972379576282\\
921	0.77599963249304\\
922	0.86379814173506\\
923	0.931028753159353\\
924	0.976090638656387\\
925	0.997910829656121\\
926	0.99596976558496\\
927	0.970313665119193\\
928	0.921553425675484\\
929	0.850850077329596\\
930	0.759887137531678\\
931	0.650830524868525\\
932	0.526276986377525\\
933	0.389192266402461\\
934	0.242840489252558\\
935	0.0907064371385786\\
936	-0.0635874259783122\\
937	-0.216367208875454\\
938	-0.363995072095349\\
939	-0.502955848542115\\
940	-0.629940743124735\\
941	-0.741926118479569\\
942	-0.836245490799229\\
943	-0.910653021480357\\
944	-0.963376992787744\\
945	-0.99316199422401\\
946	-0.999298815110706\\
947	-0.981641331603441\\
948	-0.940609986048457\\
949	-0.877181775826557\\
950	-0.792866990071161\\
951	-0.689673248171857\\
952	-0.570057696348158\\
953	-0.436868500542005\\
954	-0.293277028734792\\
955	-0.142702337506762\\
957	0.164974460086796\\
958	0.314750474130051\\
959	0.457031964537236\\
960	0.588431066696558\\
961	0.705819036546018\\
962	0.806400749099339\\
963	0.887781253168669\\
964	0.948022797545491\\
965	0.985690970799169\\
966	0.999888856044436\\
967	0.990278387431317\\
968	0.957088399830809\\
969	0.901109180048024\\
970	0.823673649302691\\
971	0.726625625040469\\
972	0.612275917797433\\
973	0.483347308487737\\
974	0.342909716267059\\
975	0.194307100687183\\
976	0.0410778386726633\\
977	-0.113129527770525\\
978	-0.26464316699321\\
979	-0.409855387685639\\
980	-0.545308541616123\\
981	-0.667777353694419\\
982	-0.774345718999029\\
983	-0.862476138150214\\
984	-0.930070137711482\\
985	-0.975518236938115\\
986	-0.997738271125399\\
987	-0.996201159032353\\
988	-0.970943500835801\\
989	-0.922566706642556\\
990	-0.852222676313886\\
991	-0.761586371578233\\
992	-0.65281593351483\\
993	-0.528501295058049\\
994	-0.391602512102054\\
995	-0.24537928161044\\
996	-0.0933133249691309\\
997	0.0609745152819414\\
998	0.21381049133106\\
999	0.361555425704864\\
1000	0.500691363726673\\
1001	0.627905339532845\\
1002	0.740168261106646\\
1003	0.834807035987524\\
1004	0.909568220278743\\
1005	0.962671675396905\\
1006	0.992852954957016\\
1007	0.999393412501831\\
1008	0.982137313194016\\
1009	0.941495542020675\\
1010	0.87843582020696\\
1011	0.794459662808549\\
1012	0.691566626134772\\
1013	0.57220669631397\\
1014	0.439221952671687\\
1015	0.295778894978866\\
1016	0.145293045927247\\
1018	-0.162391775973447\\
1019	-0.312264465434055\\
1020	-0.454701825612119\\
1021	-0.586312280480797\\
1022	-0.703962053458781\\
1023	-0.804849785766237\\
1024	-0.886573239577046\\
1025	-0.947186497708117\\
1026	-0.985246297851518\\
1027	-0.999846398094178\\
1028	-0.990639155444569\\
1029	-0.957843803560991\\
1030	-0.902241232576216\\
1031	-0.82515539532551\\
1032	-0.728421782703208\\
1033	-0.614343718789314\\
1034	-0.485637516398583\\
1035	-0.345367798955976\\
1036	-0.196874528746321\\
1037	-0.0436934790718624\\
1038	0.110527956045644\\
1039	0.2621176099683\\
1040	0.407465981399127\\
1041	0.543112180222124\\
1042	0.665826334754456\\
1043	0.772686498225994\\
1044	0.861148223250893\\
1045	0.929105147667087\\
1046	0.974939149127295\\
1047	0.997558874208153\\
1048	0.996425724628807\\
1049	0.971566681814238\\
1050	0.923573664440482\\
1051	0.853589434259447\\
1052	0.763280385796861\\
1053	0.654796867833738\\
1054	0.53072198144946\\
1055	0.394010073801837\\
1056	0.247916392166189\\
1057	0.0959195732402804\\
1058	-0.0583611866732099\\
1059	-0.211252308353323\\
1060	-0.359113301254183\\
1061	-0.498423447228561\\
1062	-0.625865632347995\\
1063	-0.738405330703245\\
1064	-0.833362859502358\\
1065	-0.908477184997992\\
1066	-0.96195975996261\\
1067	-0.992537110787453\\
1068	-0.999481160162304\\
1069	-0.982626563325539\\
1070	-0.942374645087966\\
1071	-0.879683843886596\\
1072	-0.796046890407979\\
1073	-0.693455264177828\\
1074	-0.574351774439492\\
1075	-0.441572394423019\\
1076	-0.2982787339879\\
1077	-0.147882758526066\\
1079	0.159807978845038\\
1080	0.309776316511943\\
1081	0.452368570211092\\
1082	0.584189475745916\\
1083	0.702100245494876\\
1084	0.803293306082196\\
1085	0.885359149511714\\
1086	0.946343705960771\\
1087	0.984794872136263\\
1088	0.999797087308707\\
1089	0.990993133728352\\
1090	0.958592642336498\\
1091	0.903367101243475\\
1092	0.826631485825601\\
1093	0.730212947844393\\
1094	0.61640730913814\\
1095	0.487924395804384\\
1096	0.347823514532593\\
1097	0.19944060744956\\
1098	0.0463088200008315\\
1099	-0.107925626774886\\
1100	-0.259590256418505\\
1101	-0.405073782386353\\
1102	-0.540912096397733\\
1103	-0.663870752315688\\
1104	-0.771021981546255\\
1105	-0.859814406137502\\
1106	-0.928133789640469\\
1107	-0.974353379192962\\
1108	-0.997372640134017\\
1109	-0.996643460834548\\
1110	-0.972183203783516\\
1111	-0.924574292168472\\
1112	-0.854950341798485\\
1113	-0.764969168577409\\
1114	-0.656773314248312\\
1115	-0.532939030332273\\
1116	-0.396414935000848\\
1117	-0.250451803530723\\
1118	-0.098525164090006\\
1119	0.0557474580637063\\
1120	0.208692677475938\\
1121	0.51493363569034\\
1122	0.620368088176292\\
1123	0.715607634994285\\
1124	0.799087143743236\\
1125	0.869434742250178\\
1126	0.925494363425059\\
1127	0.966344743649188\\
1128	0.991314562490516\\
1129	0.99999347494213\\
1130	0.992238854883453\\
1131	0.968178138946769\\
1132	0.928206732269246\\
1133	0.872981510548925\\
1134	0.803410025183894\\
1135	0.720635588900677\\
1136	0.626018486965222\\
1137	0.521113622743542\\
1138	0.407644964981955\\
1139	0.287477216723346\\
1140	0.16258517145161\\
1143	-0.219727345716819\\
1144	-0.342725593261093\\
1145	-0.460091611875214\\
1146	-0.569896650692044\\
1147	-0.67033621326209\\
1148	-0.759759711998413\\
1149	-0.836697593337703\\
1150	-0.899885487856409\\
1151	-0.948284988459363\\
1152	-0.981100715190223\\
1153	-0.997793386218291\\
1154	-0.998088680201818\\
1155	-0.981981744384939\\
1156	-0.94973727434899\\
1157	-0.901885164098076\\
1158	-0.839211797972439\\
1159	-0.762747127493185\\
1160	-0.673747745506716\\
1161	-0.573676235794665\\
1162	-0.464177137495881\\
1163	-0.347049919343135\\
1164	-0.22421940784443\\
1165	-0.0977041553737763\\
1167	0.158037761732885\\
1168	0.283061655836718\\
1169	0.403433816663892\\
1170	0.517176091489091\\
1171	0.622419280743998\\
1172	0.717433855771105\\
1173	0.800658381274843\\
1174	0.87072517538536\\
1175	0.926482785650023\\
1176	0.967014911589558\\
1177	0.991655462847575\\
1178	0.999999505482265\\
1179	0.991909916503118\\
1180	0.967519637300484\\
1181	0.927229488936518\\
1182	0.871701585193932\\
1183	0.801848451639216\\
1184	0.718818029502927\\
1185	0.623974810830987\\
1186	0.518877414910548\\
1187	0.405252974483574\\
1188	0.284968752670011\\
1189	0.160001457037652\\
1192	-0.222280603375566\\
1193	-0.345183852489754\\
1194	-0.462414474528032\\
1195	-0.572045943680678\\
1196	-0.672276615879582\\
1197	-0.761459336367352\\
1198	-0.838128508443788\\
1199	-0.901024178555872\\
1200	-0.949112741919635\\
1201	-0.981603928409186\\
1202	-0.997963789571259\\
1203	-0.997923473341416\\
1204	-0.981483642261537\\
1205	-0.948914462592711\\
1206	-0.900751164501344\\
1207	-0.837785246276326\\
1208	-0.761051467133711\\
1209	-0.671810842356081\\
1210	-0.57152992022111\\
1211	-0.461856681276004\\
1212	-0.344593456019084\\
1213	-0.221667306048403\\
1214	-0.0950983554284903\\
1216	0.160622310922918\\
1217	0.285571605349105\\
1218	0.40582791889301\\
1219	0.519415002620462\\
1220	0.624466207319529\\
1221	0.719255159336626\\
1222	0.802224131183721\\
1223	0.872009640667329\\
1224	0.927464857865743\\
1225	0.967678451717802\\
1226	0.991989566509346\\
1227	0.999998682138084\\
1228	0.991574179683084\\
1229	0.966854504383264\\
1230	0.926245890476366\\
1231	0.87041568529412\\
1232	0.800281382314552\\
1233	0.716995543406938\\
1234	0.62192685804348\\
1235	0.516637650749544\\
1236	0.402858206426572\\
1237	0.282458335472711\\
1238	0.157416645991816\\
1241	-0.224832337547468\\
1242	-0.347639745867127\\
1243	-0.464734167843289\\
1244	-0.574191315929966\\
1245	-0.674212410788641\\
1246	-0.763153741779206\\
1247	-0.839553679110395\\
1248	-0.902156693736288\\
1249	-0.949933990267709\\
1250	-0.98210041382481\\
1251	-0.998127352992469\\
1252	-0.997751426825289\\
1253	-0.980978813158799\\
1254	-0.94808514708302\\
1255	-0.89961099125685\\
1256	-0.836352952493598\\
1257	-0.759350590612485\\
1258	-0.669869334688428\\
1259	-0.56937968744478\\
1260	-0.459533059542082\\
1261	-0.342134630889632\\
1262	-0.219113684970125\\
1263	-0.0924919036897336\\
1265	0.163205759226003\\
1266	0.288079597586147\\
1267	0.408219239622667\\
1268	0.521650353739915\\
1269	0.626508853873474\\
1270	0.721071533207578\\
1271	0.803784382737831\\
1272	0.873288129293087\\
1273	0.928440573341049\\
1274	0.968335359486446\\
1275	0.992316871185267\\
1276	0.999991004914591\\
1277	0.991231646725282\\
1278	0.966182744753496\\
1279	0.925255943630873\\
1280	0.869123819662491\\
1281	-0.0541528667249622\\
1282	0.100114272070641\\
1283	0.251997589903112\\
1284	0.397880593044647\\
1285	0.534289661263301\\
1286	0.657976758193399\\
1287	0.765996770332094\\
1288	0.855777633137222\\
1289	0.925181574456929\\
1290	0.97255601702318\\
1291	0.996772927969232\\
1292	0.997255678421425\\
1293	0.973992773608188\\
1294	0.927538126563832\\
1295	0.858997868902861\\
1296	0.770004012721529\\
1297	0.662675590756862\\
1298	0.539568200097165\\
1299	0.403613150859201\\
1300	0.258047668762174\\
1301	0.106337813544997\\
1302	-0.0479040513828295\\
1303	-0.201005272928796\\
1304	-0.349320357859597\\
1305	-0.489317775641666\\
1306	-0.617664047717881\\
1307	-0.731303120964185\\
1308	-0.827529135372515\\
1309	-0.904050853284389\\
1310	-0.959046216050865\\
1311	-0.99120572907259\\
1312	-0.999763642174912\\
1313	-0.984516182886637\\
1314	-0.945826408465109\\
1315	-0.88461556114089\\
1316	-0.80234113241886\\
1317	-0.700962158746734\\
1318	-0.58289257489696\\
1319	-0.450943735767396\\
1320	-0.308257475207029\\
1321	-0.158231295809856\\
1323	0.149462014473556\\
1324	0.299802659854322\\
1325	0.443004704009127\\
1326	0.575658363020466\\
1327	0.694605020573363\\
1328	0.797012437767989\\
1329	0.880442191562452\\
1330	0.942907736062807\\
1331	0.982921704170622\\
1332	0.999531323278006\\
1333	0.992341101740294\\
1334	0.961522245924243\\
1335	0.907808583614951\\
1336	0.832479090842753\\
1337	0.737327438182547\\
1338	0.624619281665218\\
1339	0.497038315240388\\
1340	0.357622369338969\\
1341	0.20969107709152\\
1342	0.0567668305425286\\
1343	-0.0975090910255858\\
1344	-0.249463223597559\\
1345	-0.395477387275605\\
1346	-0.532074838763947\\
1347	-0.656003056102691\\
1348	-0.764311184471808\\
1349	-0.854420298992409\\
1350	-0.924184811511168\\
1351	-0.97194355919919\\
1352	-0.99655935850069\\
1353	-0.997446082611077\\
1354	-0.974582617741817\\
1355	-0.928513365862273\\
1356	-0.860335281943208\\
1357	-0.77167175436125\\
1358	-0.664633950399548\\
1359	-0.541770547240503\\
1360	-0.406007045414754\\
1361	-0.260576109706108\\
1362	-0.10894059616885\\
1363	0.0452889019366012\\
1364	0.198440025984837\\
1365	0.346866094512734\\
1366	0.487032934358012\\
1367	0.615603032853869\\
1368	0.729515007345071\\
1369	0.826056499772676\\
1370	0.902928760629948\\
1371	0.958301384491278\\
1372	0.990855893791831\\
1373	0.999817133100805\\
1374	0.98497172634643\\
1375	0.946673157513942\\
1376	0.885833353836915\\
1377	0.803900971899566\\
1378	0.702826903673667\\
1379	0.585017823828366\\
1380	0.453278884400333\\
1381	0.310746921316422\\
1382	0.160815763193568\\
1384	-0.146872918054214\\
1385	-0.297304065895332\\
1386	-0.440656106532515\\
1387	-0.573515684481663\\
1388	-0.692719280291385\\
1389	-0.795428537105636\\
1390	-0.879197844777991\\
1391	-0.942032572299013\\
1392	-0.98243656195109\\
1393	-0.999447754325047\\
1394	-0.992661095914173\\
1395	-0.962238183843056\\
1396	-0.908903418081536\\
1397	-0.833926752752632\\
1398	-0.739093457257695\\
1399	-0.626661607225287\\
1400	-0.499308317470422\\
1401	-0.360065997215315\\
1402	-0.212250145389589\\
1403	-0.0593804052869018\\
1404	0.0949032416642694\\
1405	0.246927147498809\\
1406	0.393071470948598\\
1407	0.529856369482786\\
1408	0.654024857840795\\
1409	0.762620360107576\\
1410	0.853057108745816\\
1411	0.923181714305883\\
1412	0.97132443978262\\
1413	0.996338958726028\\
1414	0.997629650416457\\
1415	0.975165782194836\\
1416	0.929482241233927\\
1417	0.861666798341503\\
1418	0.773334207049174\\
1419	0.666587754715238\\
1420	0.543969181148896\\
1421	0.408398157243482\\
1422	0.263102764690302\\
1423	0.111542632126657\\
1424	-0.0426734420852881\\
1425	-0.195873418954761\\
1426	-0.34440945378492\\
1427	-0.484744755005522\\
1428	-0.613537798715697\\
1429	-0.727721893711987\\
1430	-0.824578202473731\\
1431	-0.901800479403391\\
1432	-0.95754998484108\\
1433	-0.990499267295036\\
1434	-0.999863771392029\\
1435	-0.985420518920819\\
1436	-0.947513418171184\\
1437	-0.887045075130118\\
1438	-0.805455301533129\\
1439	-0.70468683150375\\
1440	-0.587139063113227\\
1441	-0.681909248458396\\
1442	-0.769878942129708\\
1443	-0.845196722621495\\
1444	-0.906624844664293\\
1445	-0.953153820629268\\
1446	-0.984019010088559\\
1447	-0.998713185637826\\
1448	-0.996994868489764\\
1449	-0.978892296850063\\
1450	-0.944702961858184\\
1451	-0.894988718729564\\
1452	-0.830566553424887\\
1453	-0.752495156600162\\
1454	-0.662057525465571\\
1455	-0.560739879469565\\
1456	-0.450207236306142\\
1457	-0.332276049612574\\
1458	-0.208884358025443\\
1459	-0.0820599361504719\\
1461	0.173528189546232\\
1462	0.29809165018878\\
1463	0.417756380041737\\
1464	0.530555852046291\\
1465	0.634636360315653\\
1466	0.728287483273107\\
1467	0.80997019210281\\
1468	0.878342142594192\\
1469	0.932279734745407\\
1470	0.970896577598069\\
1471	0.993558055869016\\
1472	0.999891758988269\\
1473	0.989793601160272\\
1474	0.963429531872862\\
1475	0.921232808742843\\
1476	0.863896877517163\\
1477	0.792363976232537\\
1478	0.707809650818945\\
1479	0.611623436599075\\
1480	0.505386023164647\\
1481	0.390843277895783\\
1482	0.269877555002495\\
1483	0.144476761596707\\
1485	-0.111347852860035\\
1486	-0.237567544359536\\
1487	-0.359883136421104\\
1488	-0.476284538669915\\
1489	-0.584858852396337\\
1490	-0.68382180644312\\
1491	-0.771547079275479\\
1492	-0.8465930253642\\
1493	-0.907726366666793\\
1494	-0.953942459876544\\
1495	-0.984481806367057\\
1496	-0.998842533520474\\
1497	-0.996788642320098\\
1498	-0.978353885674096\\
1499	-0.943841213731957\\
1500	-0.893817795310497\\
1501	-0.829105697247087\\
1502	-0.750768374848121\\
1503	-0.660093195448098\\
1504	-0.558570282278652\\
1505	-0.447868026320066\\
1506	-0.329805668565768\\
1507	-0.206323403269835\\
1508	-0.0794504934942779\\
1509	0.0487280770685175\\
1510	0.176105867967635\\
1511	0.300589598476108\\
1512	0.420133547813748\\
1513	0.532773173783426\\
1514	0.636657397351428\\
1515	0.730079022612699\\
1516	0.811502792234478\\
1517	0.879590617319991\\
1518	0.933223567080859\\
1519	0.971520256942767\\
1520	0.993851332901158\\
1521	0.999849814098525\\
1522	0.98941712365604\\
1523	0.962724708650057\\
1524	0.920211222612579\\
1525	0.862575316857146\\
1526	0.790764159095033\\
1527	0.705957868022779\\
1528	0.609550119674623\\
1529	0.503125244256353\\
1530	0.388432189829018\\
1531	0.267355780730213\\
1532	0.141885743048533\\
1534	-0.113949182169563\\
1535	-0.240109770712479\\
1536	-0.362324481784526\\
1537	-0.478584882854193\\
1538	-0.586980392373334\\
1539	-0.685729677589279\\
1540	-0.773209928324377\\
1541	-0.847983525653035\\
1542	-0.908821667214852\\
1543	-0.954724560909199\\
1544	-0.984937855118005\\
1545	-0.998965035448236\\
1546	-0.996575584272705\\
1547	-0.97780876897059\\
1548	-0.942972996623666\\
1549	-0.892640745764766\\
1550	-0.827639158471811\\
1551	-0.7490364474138\\
1552	-0.65812434122563\\
1553	-0.556396856709398\\
1554	-0.445525746696603\\
1555	-0.327333027067652\\
1556	-0.203761034396848\\
1557	-0.0768405062931379\\
1558	0.05134279100821\\
1559	0.178682339379066\\
1560	0.303085486556483\\
1561	0.422507836036857\\
1562	0.534986843952538\\
1563	0.63867407080852\\
1564	0.731865558072968\\
1565	0.813029830417236\\
1566	0.880833063429691\\
1567	0.9341610032061\\
1568	0.972137277596175\\
1569	0.994137798187239\\
1570	0.999801016350375\\
1571	0.989033864797875\\
1572	0.96201328702\\
1573	0.91918332945761\\
1574	0.861247844202808\\
1575	0.789158922148999\\
1576	0.704101246670234\\
1577	0.607472624962156\\
1578	0.50086101698389\\
1579	0.386018439491636\\
1580	0.264832174031653\\
1581	0.13929375203179\\
1583	-0.116549730484621\\
1584	-0.242650351379325\\
1585	-0.364763343816776\\
1586	-0.480881946870795\\
1587	-0.589097909252359\\
1588	-0.687632848821522\\
1589	-0.774867477878615\\
1590	-0.849368213957405\\
1591	-0.909910738801045\\
1592	-0.955500118366672\\
1593	-0.985387153215015\\
1594	-0.99908069058165\\
1595	-0.996355695807779\\
1596	-0.977256950474839\\
1597	-0.94209831648368\\
1598	-0.891457578159134\\
1599	-0.826166947150796\\
1600	-0.747299386167015\\
1601	-0.656150976292793\\
1602	-0.554219617657964\\
1603	-0.443180413489699\\
1604	-0.32485814206575\\
1605	-0.201197268968372\\
1606	-0.0742299924359031\\
1607	0.0539571530503054\\
1608	0.181257586121774\\
1609	0.305579297322765\\
1610	0.424879228439295\\
1611	0.537196847381892\\
1612	0.640686366864884\\
1613	0.733647077408023\\
1614	0.814551296184163\\
1615	0.882069472408602\\
1616	0.935092036696915\\
1617	0.972747635329597\\
1618	0.994417449764114\\
1619	0.999745366077605\\
1620	0.9886438272124\\
1621	0.961295271858944\\
1622	0.918149136322882\\
1623	0.859914468651823\\
1624	0.787548276395682\\
1625	0.702239799486506\\
1626	0.605390966700725\\
1627	0.498593356865968\\
1628	0.38360204342689\\
1629	0.262306752202676\\
1630	0.136700806311183\\
1632	-0.119149479980479\\
1633	-0.245189268947797\\
1634	-0.36719970580134\\
1635	-0.483175714976824\\
1636	-0.591211388520605\\
1637	-0.689531307094967\\
1638	-0.77651971657815\\
1639	-0.850747080787642\\
1640	-0.910993573960695\\
1641	-0.956269126933421\\
1642	-0.985829697578538\\
1643	-0.999189498128089\\
1644	-0.996128978432353\\
1645	-0.976698433969887\\
1646	-0.941217179307841\\
1647	-0.890268300603111\\
1648	-0.82468907337352\\
1649	-0.745557203014414\\
1650	-0.654173114174228\\
1651	-0.552038580047792\\
1652	-0.440832042773764\\
1653	-0.322381030522592\\
1654	-0.198632124556752\\
1655	-0.0716189698146081\\
1656	0.0565711452759388\\
1657	0.183831590545196\\
1658	0.308071013682365\\
1659	0.427247708766572\\
1660	0.539403168924309\\
1661	0.642694271728942\\
1662	0.735423568408805\\
1663	0.816067179108359\\
1664	0.883299835782054\\
1665	0.936016661171379\\
1666	0.973351325959356\\
1667	0.994690285715478\\
1668	0.999682863662201\\
1669	0.988247013572618\\
1670	0.960570668087257\\
1671	0.917108650296541\\
1672	0.858575199343704\\
1673	0.785932232874529\\
1674	0.700373539229531\\
1675	0.603305159157117\\
1676	0.496322279444939\\
1677	0.381183018197135\\
1678	0.259779532552784\\
1679	0.134106923658692\\
1681	-0.121748412839224\\
1682	-0.247726506016988\\
1683	-0.369633551040806\\
1684	-0.485466171450298\\
1685	-0.593320815692095\\
1686	-0.691425039397927\\
1687	-0.778166633098408\\
1688	-0.852120116692276\\
1689	-0.91207016527278\\
1690	-0.95703158133847\\
1691	-0.986265485175863\\
1692	-0.999291457341315\\
1693	-0.995895433700298\\
1694	-0.976133223282886\\
1695	-0.940329591134287\\
1696	-0.889072921248044\\
1697	-0.823205547269936\\
1698	-0.743809909895845\\
1699	-0.652190768426863\\
1700	-0.549853758826885\\
1701	-0.438480650644578\\
1702	-0.319901709416172\\
1703	-0.196065618742523\\
1704	-0.0690074563249254\\
1705	0.0591847497694289\\
1706	0.18640433500741\\
1707	0.310560618558156\\
1708	0.429613260786027\\
1709	0.54160579345762\\
1710	0.644697771638675\\
1711	0.737195018898547\\
1712	0.817577468798845\\
1713	0.884524145117211\\
1714	0.936934870292134\\
1715	0.973948345348163\\
1716	0.994956304170955\\
1717	0.999613509532537\\
1718	0.987843426598829\\
1719	0.95983948067169\\
1720	0.916061878510391\\
1721	0.857230045457527\\
1722	0.784310802662276\\
1723	0.698502478691353\\
1724	0.601215216627224\\
1725	0.494047800286353\\
1726	0.378761380381093\\
1727	0.25725053240285\\
1728	0.131512121852211\\
1730	-0.12434651124795\\
1731	-0.250262045195996\\
1732	-0.372064862852767\\
1733	-0.487753300592885\\
1734	-0.595426176309047\\
1735	-0.693314032751005\\
1736	-0.77980821615165\\
1737	-0.853487312261223\\
1738	-0.91314050535766\\
1739	-0.957787476356316\\
1740	-0.986694513020211\\
1741	-0.999386567522834\\
1742	-0.99565506321278\\
1743	-0.975561322287831\\
1744	-0.939435558047535\\
1745	-0.887871448286205\\
1746	-0.821716379007739\\
1747	-0.742057518787533\\
1748	-0.65020395263673\\
1749	-0.54766516897007\\
1750	-0.436126253217935\\
1751	-0.317420195738578\\
1752	-0.193497769116675\\
1753	-0.0663954698657108\\
1754	0.0617979486169133\\
1755	0.188975801875586\\
1756	0.313048094886199\\
1757	0.431975868284098\\
1758	0.543804705885577\\
1759	0.646696852861623\\
1760	0.738961416735947\\
1761	-0.966395146252125\\
1762	-0.915341345268644\\
1763	-0.842492351231613\\
1764	-0.749582771561109\\
1765	-0.638824876286435\\
1766	-0.512855921686423\\
1767	-0.374675354577448\\
1768	-0.227573392485738\\
1769	-0.0750526802758031\\
1770	0.0792551113277113\\
1771	0.231675759450809\\
1772	0.378579975994398\\
1773	0.516469824613068\\
1774	0.642062010079826\\
1775	0.752366056826759\\
1776	0.844755515163797\\
1777	0.917030499670091\\
1778	0.967470070666877\\
1779	0.994873211514005\\
1780	0.998587426016002\\
1781	0.978524275010841\\
1782	0.935161482192598\\
1783	0.86953155903484\\
1784	0.783197219657723\\
1785	0.678214171042782\\
1786	0.557082164592885\\
1787	0.42268547455069\\
1788	0.278224220544416\\
1789	0.12713816953692\\
1791	-0.18044620072078\\
1792	-0.329620628436714\\
1793	-0.470946459267452\\
1794	-0.601058583812119\\
1795	-0.716858902232616\\
1796	-0.815590093112405\\
1797	-0.894901268034573\\
1798	-0.952903948588755\\
1799	-0.988217032930152\\
1800	-0.999999681197096\\
1801	-0.98797133674816\\
1802	-0.952418406496236\\
1803	-0.89418744127488\\
1804	-0.814664978614474\\
1805	-0.715744527898096\\
1806	-0.599781484008417\\
1807	-0.46953704301859\\
1808	-0.328112455353221\\
1809	-0.178875181913099\\
1811	0.128722051497334\\
1812	0.279757813647848\\
1813	0.424132262411604\\
1814	0.558407697745224\\
1815	0.679386887215969\\
1816	0.784189195306681\\
1817	0.870319174227006\\
1818	0.935725983022621\\
1819	0.978852220149292\\
1820	0.998671006761015\\
1821	0.994710437724279\\
1822	0.967064818149538\\
1823	0.916392417892439\\
1824	0.843899797491758\\
1825	0.75131307875381\\
1826	0.640836844066598\\
1827	0.515101643088656\\
1828	0.377101356729327\\
1829	0.230121909848549\\
1830	0.0776630301020305\\
1831	-0.0766450840778816\\
1832	-0.229128202137872\\
1833	-0.376155548506176\\
1834	-0.514226254992991\\
1835	-0.640052719964842\\
1836	-0.750638889424863\\
1837	-0.843351596058255\\
1838	-0.915983257578773\\
1839	-0.96680444147205\\
1840	-0.994605044512355\\
1841	-0.998723106531997\\
1842	-0.979060572355593\\
1843	-0.936085626590284\\
1844	-0.870821545682702\\
1845	-0.784822332683234\\
1846	-0.68013571487927\\
1847	-0.559254385359509\\
1848	-0.425056649498401\\
1849	-0.280737889625016\\
1850	-0.129734479776289\\
1852	0.177870566152251\\
1853	0.327147818387857\\
1854	0.468635353819536\\
1855	0.598964212697865\\
1856	0.715031134523542\\
1857	0.814072449785272\\
1858	0.893729885694938\\
1859	0.95210671902214\\
1860	0.987812938970364\\
1861	0.999998344725554\\
1862	0.988372789587174\\
1863	0.953213089652309\\
1864	0.895356432547032\\
1865	0.816180443151552\\
1866	0.717570380973484\\
1867	0.601874250233323\\
1868	0.471846891536188\\
1869	0.330584386359988\\
1870	0.181450336258422\\
1872	-0.126125399307966\\
1873	-0.277243398839346\\
1874	-0.42175995571597\\
1875	-0.556233986157167\\
1876	-0.677463528983481\\
1877	-0.782561987512963\\
1878	-0.869026862317241\\
1879	-0.934799338236189\\
1880	-0.97831330682493\\
1881	-0.998532656964926\\
1882	-0.994975945703118\\
1883	-0.967727861892399\\
1884	-0.917437209665422\\
1885	-0.845301459756683\\
1886	-0.753038236530756\\
1887	-0.642844419623088\\
1888	-0.517343834042549\\
1889	-0.379524774270521\\
1890	-0.232668849982019\\
1891	-0.0802728476351149\\
1892	0.0740345315116429\\
1893	0.22657907440589\\
1894	0.373728542889694\\
1895	0.511979160924056\\
1896	0.63803904300039\\
1897	0.748906577227899\\
1898	0.84194189671507\\
1899	0.914929737440616\\
1900	0.966132185907554\\
1901	0.994330060598713\\
1902	0.998851941911653\\
1903	0.979590159329291\\
1904	0.937003355162233\\
1905	0.872105563817513\\
1906	0.786442066624659\\
1907	0.682052597142047\\
1908	0.561422773059348\\
1909	0.427424911154958\\
1910	0.283249634559979\\
1911	0.132329900829518\\
1913	-0.175293712478378\\
1914	-0.324672766104868\\
1915	-0.466321036397858\\
1916	-0.596865736350082\\
1917	-0.71319846607048\\
1918	-0.812549226896408\\
1919	-0.89255237783118\\
1920	-0.951302963823309\\
1921	-0.987402074651527\\
1922	-0.999990154376974\\
1923	-0.98876746822998\\
1924	-0.954001239592799\\
1925	-0.896519287146475\\
1926	-0.817690313679122\\
1927	-0.719391315901248\\
1928	-0.60396289127948\\
1929	-0.474153506068433\\
1930	-0.333054051578529\\
1931	-0.184024246963872\\
1933	0.123527882669805\\
1934	0.274727083835842\\
1935	0.419384758324668\\
1936	0.554056462203789\\
1937	0.675535527491775\\
1938	0.780929416126583\\
1939	0.867728594194432\\
1940	0.933866286439297\\
1941	0.977767688251333\\
1942	0.998387463337622\\
1943	0.995234634228382\\
1944	0.968384272937328\\
1945	0.918475713426233\\
1946	0.846697328420305\\
1947	0.75475823306806\\
1948	0.644847589195706\\
1949	0.519582479179462\\
1950	0.381945590591386\\
1951	0.235214195428853\\
1952	0.0828821149866599\\
1953	-0.0714234715210296\\
1954	-0.224028393726257\\
1955	-0.371298975779609\\
1956	-0.509728557806739\\
1957	-0.636020992987596\\
1958	-0.747169132108866\\
1959	-0.840526426796259\\
1960	-0.913869946477007\\
1961	-0.965453308581345\\
1962	-0.994048261658463\\
1963	-0.998973931272303\\
1964	-0.980113032301688\\
1965	-0.937914661617469\\
1966	-0.873383604638093\\
1967	-0.78805641038025\\
1968	-0.68396480469255\\
1969	-0.56358731283035\\
1970	-0.429790243289517\\
1971	-0.285759438133482\\
1972	-0.13492441490871\\
1974	0.17271565736155\\
1975	0.322195488551188\\
1976	0.464003522864004\\
1977	0.594763169151065\\
1978	0.711360909435371\\
1979	0.811020434886359\\
1980	0.891368752513245\\
1981	0.950492688500162\\
1982	0.986984442789435\\
1983	0.999975110208197\\
1984	0.989155369970831\\
1985	0.954782850915308\\
1986	0.897675997103306\\
1987	0.819194579848954\\
1988	0.721207320201756\\
1989	0.606047392831897\\
1990	0.476456870805123\\
1991	0.335521434081329\\
1992	0.186596896387528\\
1993	0.0332292999369201\\
1994	-0.120929519385299\\
1995	-0.272208885883629\\
1996	-0.417006686517198\\
1997	-0.551875140809898\\
1998	-0.6736028959549\\
1999	-0.779291492337961\\
2000	-0.866424378757074\\
2001	-0.932926834027967\\
2002	-0.977215368166981\\
2003	-0.998235426874999\\
2004	-0.995486501526557\\
2005	-0.9690340467846\\
2006	-0.919507922057164\\
2007	-0.848087393915648\\
2008	-0.756473056576397\\
2009	-0.646846339055173\\
2010	-0.521817563156219\\
2011	-0.384363789100007\\
2012	-0.237757928744031\\
2013	-0.0854908142737258\\
2014	0.0688119220017143\\
2015	0.22147617758128\\
2016	0.368866863828316\\
2017	0.507474461067432\\
2018	0.633998583758512\\
2019	0.745426565976686\\
2020	0.839105196003402\\
2021	0.912803891950716\\
2022	0.964767814145944\\
2023	0.993759649622916\\
2024	0.999089073777213\\
2025	0.980629187688919\\
2026	0.938819539711403\\
2027	0.874655659385553\\
2028	0.789665352885095\\
2029	0.685872324424963\\
2030	0.565747989837291\\
2031	0.432152629689881\\
2032	0.288267283144251\\
2033	0.137518004230969\\
2035	-0.17013641847052\\
2036	-0.319716002705263\\
2037	-0.461682829101846\\
2038	-0.592656525512211\\
2039	-0.709518477211986\\
2040	-0.809486084232958\\
2041	-0.89017901785337\\
2042	-0.949675898606529\\
2043	-0.986560046246268\\
2044	-0.999953212321998\\
2045	-0.989536492151728\\
2046	-0.955557918263366\\
2047	-0.89882655448946\\
2048	-0.82069323135056\\
2049	-0.723018381428119\\
2050	-0.608127740603322\\
2051	-0.478756969960159\\
2052	-0.337986516958154\\
2053	-0.189168266897468\\
2054	-0.0358457311749589\\
2055	0.118330327262811\\
2056	0.269688822242188\\
2057	0.414625756592613\\
2058	0.549690036924858\\
2059	0.671665647619648\\
2060	0.777648227372083\\
2061	0.865114224944136\\
2062	0.931980987440511\\
2063	0.976656350358098\\
2064	0.998076548617973\\
2065	0.995731545871877\\
2066	0.969677178981328\\
2067	0.920533828484167\\
2068	0.849471646714392\\
2069	0.758182695303731\\
2070	0.648840655502227\\
2071	0.524049070653746\\
2072	0.386779353222664\\
2073	0.240300032492996\\
2074	0.0880989276161017\\
2075	-0.0661999008534622\\
2076	-0.218922443463271\\
2077	-0.366432223704578\\
2078	-0.505216886154813\\
2079	-0.631971829173835\\
2080	-0.743678890773936\\
2081	-0.837678214077641\\
2082	-0.911731581168624\\
2083	-0.964075707300253\\
2084	-0.993464226469769\\
2085	-0.999197368637397\\
2086	-0.981138621953505\\
2087	-0.939717983240826\\
2088	-0.875921719340568\\
2089	-0.79126888311248\\
2090	-0.687775143265299\\
2091	-0.567904789270415\\
2092	-0.434512054164315\\
2093	-0.290773152402835\\
2094	-0.140110651020223\\
2096	0.167556013484045\\
2097	0.317234325562367\\
2098	0.459358971017082\\
2099	0.590545819872204\\
2100	0.707671182028662\\
2101	0.807946185452693\\
2102	0.888983182006541\\
2103	0.948852599740803\\
2104	0.9861288879315\\
2105	0.999924460868442\\
2106	0.989910832160604\\
2107	0.956326436324161\\
2108	0.899970951418709\\
2109	0.822186257913017\\
2110	0.724824487167098\\
2111	0.610203920334698\\
2112	0.48105378776836\\
2113	0.340449283313319\\
2114	0.191738340868596\\
2115	0.0384619167307392\\
2116	-0.115730324117976\\
2117	-0.267166910183732\\
2118	-0.41224198486907\\
2119	-0.547501165526228\\
2120	-0.66972379576282\\
2121	-0.77599963249304\\
2122	-0.86379814173506\\
2123	-0.931028753159353\\
2124	-0.976090638656387\\
2125	-0.997910829656121\\
2126	-0.99596976558496\\
2127	-0.970313665119193\\
2128	-0.921553425675484\\
2129	-0.850850077329596\\
2130	-0.759887137531678\\
2131	-0.650830524868525\\
2132	-0.526276986377525\\
2133	-0.389192266402461\\
2134	-0.242840489252558\\
2135	-0.0907064371385786\\
2136	0.0635874259783122\\
2137	0.216367208875454\\
2138	0.363995072095349\\
2139	0.502955848542115\\
2140	0.629940743124735\\
2141	0.741926118479569\\
2142	0.836245490799229\\
2143	0.910653021480357\\
2144	0.963376992787744\\
2145	0.99316199422401\\
2146	0.999298815110706\\
2147	0.981641331603441\\
2148	0.940609986048457\\
2149	0.877181775826557\\
2150	0.792866990071161\\
2151	0.689673248171857\\
2152	0.570057696348158\\
2153	0.436868500542005\\
2154	0.293277028734792\\
2155	0.142702337506762\\
2157	-0.164974460086796\\
2158	-0.314750474130051\\
2159	-0.457031964537236\\
2160	-0.588431066696558\\
2161	-0.705819036546018\\
2162	-0.806400749099339\\
2163	-0.887781253168669\\
2164	-0.948022797545491\\
2165	-0.985690970799169\\
2166	-0.999888856044436\\
2167	-0.990278387431317\\
2168	-0.957088399830809\\
2169	-0.901109180048024\\
2170	-0.823673649302691\\
2171	-0.726625625040469\\
2172	-0.612275917797433\\
2173	-0.483347308487737\\
2174	-0.342909716267059\\
2175	-0.194307100687183\\
2176	-0.0410778386731181\\
2177	0.113129527770525\\
2178	0.26464316699321\\
2179	0.409855387685639\\
2180	0.545308541616123\\
2181	0.667777353694419\\
2182	0.774345718999029\\
2183	0.862476138150214\\
2184	0.930070137711482\\
2185	0.975518236938115\\
2186	0.997738271125399\\
2187	0.996201159032353\\
2188	0.970943500835801\\
2189	0.922566706642556\\
2190	0.852222676313886\\
2191	0.761586371578233\\
2192	0.65281593351483\\
2193	0.528501295058049\\
2194	0.391602512102054\\
2195	0.24537928161044\\
2196	0.0933133249686762\\
2197	-0.0609745152819414\\
2198	-0.21381049133106\\
2199	-0.361555425704864\\
2200	-0.500691363726673\\
2201	-0.627905339532845\\
2202	-0.740168261106646\\
2203	-0.834807035987524\\
2204	-0.909568220278743\\
2205	-0.962671675396905\\
2206	-0.992852954957016\\
2207	-0.999393412501831\\
2208	-0.982137313194016\\
2209	-0.941495542020675\\
2210	-0.87843582020696\\
2211	-0.794459662808549\\
2212	-0.691566626134772\\
2213	-0.57220669631397\\
2214	-0.439221952671687\\
2215	-0.295778894978866\\
2216	-0.145293045927247\\
2218	0.162391775973447\\
2219	0.312264465434055\\
2220	0.454701825612119\\
2221	0.586312280480797\\
2222	0.703962053458781\\
2223	0.804849785766237\\
2224	0.886573239577046\\
2225	0.947186497708117\\
2226	0.985246297851518\\
2227	0.999846398094178\\
2228	0.990639155444569\\
2229	0.957843803560536\\
2230	0.902241232576216\\
2231	0.82515539532551\\
2232	0.728421782703208\\
2233	0.614343718789314\\
2234	0.485637516398583\\
2235	0.345367798955976\\
2236	0.196874528746321\\
2237	0.0436934790718624\\
2238	-0.110527956045644\\
2239	-0.2621176099683\\
2240	-0.407465981399127\\
2241	-0.543112180222124\\
2242	-0.665826334754456\\
2243	-0.772686498225994\\
2244	-0.861148223250893\\
2245	-0.929105147667087\\
2246	-0.974939149127295\\
2247	-0.997558874208153\\
2248	-0.996425724628807\\
2249	-0.971566681814238\\
2250	-0.923573664440482\\
2251	-0.853589434259447\\
2252	-0.763280385796861\\
2253	-0.654796867833738\\
2254	-0.53072198144946\\
2255	-0.394010073801837\\
2256	-0.247916392166189\\
2257	-0.0959195732402804\\
2258	0.0583611866732099\\
2259	0.211252308353323\\
2260	0.359113301254183\\
2261	0.498423447228561\\
2262	0.625865632347541\\
2263	0.738405330703245\\
2264	0.833362859502358\\
2265	0.908477184997992\\
2266	0.96195975996261\\
2267	0.992537110787453\\
2268	0.999481160162304\\
2269	0.982626563325539\\
2270	0.942374645087966\\
2271	0.879683843886596\\
2272	0.796046890407979\\
2273	0.693455264177828\\
2274	0.574351774439492\\
2275	0.441572394423019\\
2276	0.2982787339879\\
2277	0.147882758526066\\
2279	-0.159807978845038\\
2280	-0.309776316511943\\
2281	-0.452368570211092\\
2282	-0.584189475745916\\
2283	-0.702100245494876\\
2284	-0.803293306082196\\
2285	-0.885359149511714\\
2286	-0.946343705960771\\
2287	-0.984794872136263\\
2288	-0.999797087308707\\
2289	-0.990993133728352\\
2290	-0.958592642336498\\
2291	-0.903367101243475\\
2292	-0.826631485825601\\
2293	-0.730212947844393\\
2294	-0.61640730913814\\
2295	-0.487924395804384\\
2296	-0.347823514532593\\
2297	-0.19944060744956\\
2298	-0.0463088200008315\\
2299	0.107925626774886\\
2300	0.259590256418505\\
2301	0.405073782386353\\
2302	0.540912096397733\\
2303	0.663870752315688\\
2304	0.771021981546255\\
2305	0.859814406137502\\
2306	0.928133789640469\\
2307	0.974353379192507\\
2308	0.997372640134017\\
2309	0.996643460834548\\
2310	0.972183203783516\\
2311	0.924574292168472\\
2312	0.854950341798485\\
2313	0.764969168577409\\
2314	0.656773314248312\\
2315	0.532939030332273\\
2316	0.396414935000848\\
2317	0.250451803530723\\
2318	0.098525164090006\\
2319	-0.0557474580637063\\
2320	-0.208692677475938\\
2321	-0.356668715480737\\
2322	-0.496152114591496\\
2323	-0.62382163555003\\
2324	-0.736637339352455\\
2325	-0.831912971241309\\
2326	-0.907379923116423\\
2327	-0.96124125136339\\
2328	-0.992214463879009\\
2329	-0.999562057490948\\
2330	-0.983109078644702\\
2331	-0.943247289224473\\
2332	-0.880925838311214\\
2333	-0.797628661991439\\
2334	-0.695339149356187\\
2335	-0.576492916021834\\
2336	-0.443919809687031\\
2337	-0.300776528627466\\
2338	-0.150471457552612\\
2340	0.15722308641034\\
2341	0.307286044417651\\
2342	0.450032214326711\\
2343	0.58206266704201\\
2344	0.700233625414057\\
2345	0.801731320716044\\
2346	0.884138991294094\\
2347	0.945494428079201\\
2348	0.98433669674705\\
2349	0.9997409240259\\
2350	0.991340319856135\\
2351	0.959334911025508\\
2352	0.904486778334103\\
2353	0.828101910685746\\
2354	0.73199910818812\\
2355	0.618466674700358\\
2356	0.490207931031364\\
2357	0.350276846165798\\
2358	0.202005319209093\\
2359	0.0489238435343395\\
2360	-0.105322557793897\\
2361	-0.257061123665608\\
2362	-0.40267880704323\\
2363	-0.538708305222826\\
2364	-0.661910619780883\\
2365	-0.769352180368514\\
2366	-0.858474695952737\\
2367	-0.92715607028822\\
2368	-0.973760931149172\\
2369	-0.997179570179469\\
2370	-0.996854366158004\\
2371	-0.972793062518122\\
2372	-0.925568582967571\\
2373	-0.856305389603676\\
2374	-0.766652708344736\\
2375	-0.658745259212083\\
2376	-0.535152426510194\\
2377	-0.398817079216315\\
2378	-0.25298549832678\\
2379	-0.101130079659015\\
2380	0.0531333473668383\\
2381	0.206131616242146\\
2382	0.354221685140146\\
2383	0.493877381384209\\
2384	0.62177336314835\\
2385	0.734864299171477\\
2386	0.830457381141514\\
2387	0.906276442154194\\
2388	0.960516154524612\\
2389	0.991885016444485\\
2390	0.999636103932971\\
2391	0.983584855843674\\
2392	0.944113468450269\\
2393	0.882161794969306\\
2394	0.799204966717753\\
2395	0.697218268757752\\
2396	0.578630106386299\\
2397	0.446264182374307\\
2398	0.303272261778602\\
2399	0.153059125264463\\
2401	-0.154637116386766\\
2402	-0.304793666219211\\
2403	-0.447692773971994\\
2404	-0.579931868946005\\
2405	-0.698362206010188\\
2406	-0.800163840372988\\
2407	-0.882912773286989\\
2408	-0.94463866988508\\
2409	-0.983871774824365\\
2410	-0.999677908630929\\
2411	-0.991680711448225\\
2412	-0.960070604540761\\
2413	-0.905600256173329\\
2414	-0.829566659828288\\
2415	-0.73378025149168\\
2416	-0.620521801361065\\
2417	-0.492488106427572\\
2418	-0.352727777040855\\
2419	-0.20456864644666\\
2420	-0.0515385317494292\\
2421	0.102718766943781\\
2422	0.254530229045031\\
2423	0.400281071784775\\
2424	0.536500821801383\\
2425	0.659945950585097\\
2426	0.767677106136944\\
2427	0.8571291018784\\
2428	0.926171996312405\\
2429	0.973161809057274\\
2430	0.996979665667823\\
2431	0.997058439153079\\
2432	0.973396253837564\\
2433	0.926556530022935\\
2434	0.857654568387261\\
2435	0.768330993560539\\
2436	0.660712689209959\\
2437	0.537362154813763\\
2438	0.401216489984108\\
2439	0.255517459188468\\
2440	0.103734302093471\\
2441	-0.050518872500561\\
2442	-0.203569142205652\\
2443	-0.351772227003494\\
2444	-0.491599263195894\\
2445	-0.61972082918146\\
2446	-0.733086222313432\\
2447	-0.828996099180131\\
2448	-0.905166749675118\\
2449	-0.959784474415301\\
2450	-0.991548770740792\\
2451	-0.999703298981331\\
2452	-0.984053891662825\\
2453	-0.944973176827716\\
2454	-0.883391705388931\\
2455	-0.800775793782122\\
2456	-0.699092609503168\\
2457	-0.58076333088502\\
2458	-0.448605496416803\\
2459	-0.305765916335531\\
2460	-0.15564574392647\\
2462	0.152050086497184\\
2463	0.302299198999208\\
2464	0.445350265180878\\
2465	0.577797096061659\\
2466	0.696486000110326\\
2467	0.79859087579689\\
2468	0.881680503894586\\
2469	0.943776437242832\\
2470	0.983400109554623\\
2471	0.999608041555348\\
2472	0.992014306171768\\
2473	0.960799717839109\\
2474	0.906707527130038\\
2475	0.831025723214225\\
2476	0.735556365547382\\
2477	0.622572675034917\\
2478	0.494764906365617\\
2479	0.355176290359395\\
2480	0.207130571594007\\
2481	0.0541528667249622\\
2482	-0.100114272070641\\
2483	-0.251997589903112\\
2484	-0.397880593044647\\
2485	-0.534289661263301\\
2486	-0.657976758193399\\
2487	-0.765996770332094\\
2488	-0.855777633137222\\
2489	-0.925181574456929\\
2490	-0.97255601702318\\
2491	-0.996772927969232\\
2492	-0.997255678421425\\
2493	-0.973992773608188\\
2494	-0.927538126563832\\
2495	-0.858997868902861\\
2496	-0.770004012721529\\
2497	-0.662675590756862\\
2498	-0.539568200097165\\
2499	-0.403613150859201\\
2500	-0.258047668762174\\
2501	-0.106337813544997\\
2502	0.0479040513832842\\
2503	0.201005272928796\\
2504	0.349320357859597\\
2505	0.489317775641666\\
2506	0.617664047717881\\
2507	0.731303120964185\\
2508	0.827529135372515\\
2509	0.904050853284389\\
2510	0.959046216050865\\
2511	0.99120572907259\\
2512	0.999763642174912\\
2513	0.984516182886637\\
2514	0.945826408465109\\
2515	0.88461556114089\\
2516	0.80234113241886\\
2517	0.700962158746734\\
2518	0.58289257489696\\
2519	0.450943735767396\\
2520	0.308257475207029\\
2521	0.158231295809856\\
2523	-0.149462014473556\\
2524	-0.299802659854322\\
2525	-0.443004704009127\\
2526	-0.575658363020466\\
2527	-0.694605020573363\\
2528	-0.797012437767989\\
2529	-0.880442191562452\\
2530	-0.942907736062807\\
2531	-0.982921704170622\\
2532	-0.999531323278006\\
2533	-0.992341101740294\\
2534	-0.961522245924243\\
2535	-0.907808583614951\\
2536	-0.832479090842753\\
2537	-0.737327438182547\\
2538	-0.624619281665218\\
2539	-0.497038315240843\\
2540	-0.357622369338969\\
2541	-0.20969107709152\\
2542	-0.0567668305425286\\
2543	0.0975090910255858\\
2544	0.249463223597104\\
2545	0.395477387275605\\
2546	0.532074838763947\\
2547	0.656003056102691\\
2548	0.764311184471808\\
2549	0.854420298991954\\
2550	0.924184811511168\\
2551	0.97194355919919\\
2552	0.99655935850069\\
2553	0.997446082611077\\
2554	0.974582617741817\\
2555	0.928513365862273\\
2556	0.860335281943208\\
2557	0.77167175436125\\
2558	0.664633950399548\\
2559	0.541770547240503\\
2560	0.406007045414754\\
2561	0.0872771162930803\\
2563	-0.168369282350341\\
2564	-0.293089641487768\\
2565	-0.412993471044501\\
2566	-0.526110314683137\\
2567	-0.630581250946307\\
2568	-0.724689442103681\\
2569	-0.806888348045959\\
2570	-0.875827141565424\\
2571	-0.930372907364927\\
2572	-0.969629259980138\\
2573	-0.992951074654684\\
2574	-0.999955089088417\\
2575	-0.990526201830562\\
2576	-0.964819363818151\\
2577	-0.923257031965932\\
2578	-0.8665222266618\\
2579	-0.795547307257493\\
2580	-0.711498650009162\\
2581	-0.615757480270076\\
2582	-0.509897173932586\\
2583	-0.395657401134031\\
2584	-0.274915537144352\\
2585	-0.149655810258082\\
2587	0.106142922577874\\
2588	0.232478224305851\\
2589	0.354993062648191\\
2590	0.471674072888163\\
2591	0.580603761339262\\
2592	0.679992016744109\\
2593	0.768205528319413\\
2594	0.843794626995532\\
2595	0.905517108756158\\
2596	0.952358648572044\\
2597	0.983549469453465\\
2598	0.998576992686594\\
2599	0.99719426136835\\
2600	0.979423998807306\\
2601	0.945558235096541\\
2602	0.896153507995678\\
2603	0.832021716992131\\
2604	0.754216780835122\\
2605	0.664017317815251\\
2606	0.562905633412356\\
2607	0.452543360622258\\
2608	0.334744153276915\\
2609	0.211443881110881\\
2610	0.0846688163769613\\
2612	-0.170949321782246\\
2613	-0.295591658813919\\
2614	-0.415376349014878\\
2615	-0.528334893938791\\
2616	-0.632610973553255\\
2617	-0.726490952332369\\
2618	-0.808432040525531\\
2619	-0.877087647810185\\
2620	-0.931329512669436\\
2621	-0.970266243837159\\
2622	-0.993257969101251\\
2623	-0.999926850731754\\
2624	-0.990163294729427\\
2625	-0.964127751857177\\
2626	-0.922248080847112\\
2627	-0.865212517124291\\
2628	-0.793958362595731\\
2629	-0.709656582367643\\
2630	-0.613692561524658\\
2631	-0.507643338214166\\
2632	-0.393251687166412\\
2633	-0.272397479564461\\
2634	-0.147066789917062\\
2636	0.108745760384409\\
2637	0.235023689745958\\
2638	0.357439324459392\\
2639	0.473980930084963\\
2640	0.582733303861914\\
2641	0.681909248458396\\
2642	0.769878942129708\\
2643	0.845196722621495\\
2644	0.906624844664293\\
2645	0.953153820629268\\
2646	0.984019010088559\\
2647	0.998713185637826\\
2648	0.996994868489764\\
2649	0.978892296850063\\
2650	0.944702961858184\\
2651	0.894988718729564\\
2652	0.830566553424887\\
2653	0.752495156600617\\
2654	0.662057525465571\\
2655	0.560739879469565\\
2656	0.450207236306142\\
2657	0.332276049612574\\
2658	0.208884358025443\\
2659	0.0820599361504719\\
2661	-0.173528189546232\\
2662	-0.29809165018878\\
2663	-0.417756380041737\\
2664	-0.530555852046291\\
2665	-0.634636360315653\\
2666	-0.728287483273107\\
2667	-0.80997019210281\\
2668	-0.878342142594192\\
2669	-0.932279734745407\\
2670	-0.970896577598069\\
2671	-0.993558055869016\\
2672	-0.999891758988269\\
2673	-0.989793601160272\\
2674	-0.963429531872862\\
2675	-0.921232808743298\\
2676	-0.863896877517163\\
2677	-0.792363976232537\\
2678	-0.707809650818945\\
2679	-0.611623436599075\\
2680	-0.505386023164647\\
2681	-0.390843277895783\\
2682	-0.269877555002495\\
2683	-0.144476761596707\\
2685	0.111347852860035\\
2686	0.237567544359536\\
2687	0.359883136421104\\
2688	0.476284538669915\\
2689	0.584858852396337\\
2690	0.68382180644312\\
2691	0.771547079275479\\
2692	0.8465930253642\\
2693	0.907726366666793\\
2694	0.953942459876544\\
2695	0.984481806367057\\
2696	0.998842533520474\\
2697	0.996788642320098\\
2698	0.978353885674096\\
2699	0.943841213731957\\
2700	0.893817795310497\\
2701	0.829105697247087\\
2702	0.750768374848121\\
2703	0.660093195448098\\
2704	0.558570282278652\\
2705	0.447868026320066\\
2706	0.329805668565768\\
2707	0.206323403269835\\
2708	0.0794504934942779\\
2709	-0.0487280770685175\\
2710	-0.176105867967635\\
2711	-0.300589598476108\\
2712	-0.420133547813748\\
2713	-0.532773173782971\\
2714	-0.636657397351428\\
2715	-0.730079022612699\\
2716	-0.811502792234478\\
2717	-0.879590617319991\\
2718	-0.933223567080859\\
2719	-0.971520256942767\\
2720	-0.993851332901158\\
2721	0.63073335917079\\
2722	0.503838039785478\\
2723	0.36494583267995\\
2724	0.217363900240343\\
2725	0.0646063159056212\\
2726	-0.0896896094373005\\
2727	-0.241849935462142\\
2728	-0.388251572599529\\
2729	-0.525408551422515\\
2730	-0.650055027066173\\
2731	-0.759223042260601\\
2732	-0.850313197362993\\
2733	-0.921156544661244\\
2734	-0.970066233192028\\
2735	-0.995877674346502\\
2736	-0.99797627189082\\
2737	-0.976312056119241\\
2738	-0.931400873679195\\
2739	-0.864312104749843\\
2740	-0.776643200028957\\
2741	-0.670481643833682\\
2742	-0.548355249011365\\
2743	-0.413171967185917\\
2744	-0.268150647523271\\
2745	-0.116744392710189\\
2746	0.037441662873789\\
2747	0.190736195011596\\
2748	0.339489107550889\\
2749	0.480158444846438\\
2750	0.609394729264295\\
2751	0.724120715596655\\
2752	0.821604663347443\\
2753	0.899525382206775\\
2754	0.956027501914377\\
2755	0.989765650479967\\
2756	0.999936488838102\\
2757	0.986297839153394\\
2758	0.949174451317958\\
2759	0.889450270330599\\
2760	0.80854738868311\\
2761	0.708392184915283\\
2762	0.591369454615233\\
2763	0.460265626049477\\
2764	0.31820241252035\\
2765	0.168562481248955\\
2767	-0.13909965990797\\
2768	-0.289796126456622\\
2769	-0.433592257349119\\
2770	-0.567064122868487\\
2771	-0.687033624384185\\
2772	-0.790644168058861\\
2773	-0.875428683204973\\
2774	-0.939368365706287\\
2775	-0.98094074776327\\
2776	-0.999155949374199\\
2777	-0.993580248367834\\
2778	-0.96434640776215\\
2779	-0.912150514545374\\
2780	-0.838235405146406\\
2781	-0.744361072258016\\
2782	-0.632762757649289\\
2783	-0.506097728824443\\
2784	-0.367382006822481\\
2785	-0.219918551743831\\
2786	-0.0672186159617922\\
2787	0.0870818623056948\\
2788	0.239308834322856\\
2789	0.385837623612588\\
2790	0.523179233127848\\
2791	0.648063421762345\\
2792	0.757516572062286\\
2793	0.848932495032386\\
2794	0.920134486102597\\
2795	0.96942715464138\\
2796	0.995636792903497\\
2797	0.99813932318375\\
2798	0.976875157733502\\
2799	0.932350617602424\\
2800	0.86562587662911\\
2801	0.778289717641201\\
2802	0.672421701947769\\
2803	0.550542652903005\\
2804	0.415554632580552\\
2805	0.270671840774412\\
2806	0.119344081684176\\
2807	-0.0348253793713411\\
2808	-0.188165613308229\\
2809	-0.337025435768737\\
2810	-0.47786034547471\\
2811	-0.607316922347763\\
2812	-0.722312675796729\\
2813	-0.820109441900058\\
2814	-0.898378581830457\\
2815	-0.955256429072961\\
2816	-0.989388665189381\\
2817	-0.999962567495004\\
2818	-0.986726360799366\\
2819	-0.949995212423346\\
2820	-0.890643727752831\\
2821	-0.810085125006935\\
2822	-0.710237585100913\\
2823	-0.593478577837686\\
2824	-0.462588251975831\\
2825	-0.320683237103367\\
2826	-0.17114243357446\\
2828	0.136506643426856\\
2829	0.287289484632311\\
2830	0.431231675830531\\
2831	0.564905809463198\\
2832	0.685128970693768\\
2833	0.789038525791966\\
2834	0.874160284310619\\
2835	0.938467412031059\\
2836	0.980428691915222\\
2837	0.99904498391561\\
2838	0.993873015496774\\
2839	0.965035936400454\\
2840	0.913220386326429\\
2841	0.839660145353264\\
2842	0.746106756400877\\
2843	0.634787819242774\\
2844	0.508353949126104\\
2845	0.369815662970268\\
2846	0.222471695949935\\
2847	0.0698304553093294\\
2848	-0.0844735183245575\\
2849	-0.236766092987637\\
2850	-0.38342103013747\\
2851	-0.520946329021172\\
2852	-0.646067374704671\\
2853	-0.755804909930703\\
2854	-0.847545974213972\\
2855	-0.919106121044933\\
2856	-0.968781431745811\\
2857	-0.995389087477633\\
2858	-0.998295533341661\\
2859	-0.977431563955179\\
2860	-0.933293971298554\\
2861	-0.866933715605228\\
2862	-0.779930900942873\\
2863	-0.674357151359345\\
2864	-0.552726283437096\\
2865	-0.417934449810218\\
2866	-0.273191178871002\\
2867	-0.121942952687277\\
2869	0.185593741939101\\
2870	0.334559454051941\\
2871	0.475558970901602\\
2872	0.605234952949104\\
2873	0.720499685346113\\
2874	0.818608599513936\\
2875	0.897225624067687\\
2876	0.954478809010652\\
2877	0.989004898739495\\
2878	0.999981792520884\\
2879	0.987148119532776\\
2880	0.950809462367943\\
2881	0.121748412839224\\
2882	0.247726506016988\\
2883	0.369633551040351\\
2884	0.485466171450298\\
2885	0.593320815692095\\
2886	0.691425039397927\\
2887	0.778166633098408\\
2888	0.852120116692276\\
2889	0.91207016527278\\
2890	0.95703158133847\\
2891	0.986265485175863\\
2892	0.999291457341315\\
2893	0.995895433700298\\
2894	0.976133223282886\\
2895	0.940329591134287\\
2896	0.889072921248044\\
2897	0.823205547269936\\
2898	0.743809909895845\\
2899	0.652190768426863\\
2900	0.549853758826885\\
2901	0.438480650644578\\
2902	0.319901709416172\\
2903	0.196065618742523\\
2904	0.0690074563249254\\
2905	-0.0591847497694289\\
2906	-0.18640433500741\\
2907	-0.310560618558156\\
2908	-0.429613260786027\\
2909	-0.54160579345762\\
2910	-0.644697771638675\\
2911	-0.737195018898547\\
2912	-0.817577468798845\\
2913	-0.884524145117211\\
2914	-0.936934870292134\\
2915	-0.973948345348163\\
2916	-0.994956304170955\\
2917	-0.999613509532537\\
2918	-0.987843426598829\\
2919	-0.95983948067169\\
2920	-0.916061878510391\\
2921	-0.857230045457527\\
2922	-0.784310802662276\\
2923	-0.698502478691353\\
2924	-0.601215216627224\\
2925	-0.494047800286808\\
2926	-0.378761380381093\\
2927	-0.25725053240285\\
2928	-0.131512121852211\\
2930	0.12434651124795\\
2931	0.250262045195996\\
2932	0.372064862852767\\
2933	0.487753300592885\\
2934	0.595426176309047\\
2935	0.693314032751005\\
2936	0.77980821615165\\
2937	0.853487312261223\\
2938	0.91314050535766\\
2939	0.957787476356316\\
2940	0.986694513020211\\
2941	0.999386567522834\\
2942	0.99565506321278\\
2943	0.975561322287831\\
2944	0.939435558047535\\
2945	0.887871448286205\\
2946	0.821716379007739\\
2947	0.742057518787533\\
2948	0.65020395263673\\
2949	0.54766516897007\\
2950	0.436126253217935\\
2951	0.317420195738578\\
2952	0.193497769116675\\
2953	0.0663954698657108\\
2954	-0.0617979486169133\\
2955	-0.188975801875586\\
2956	-0.313048094886199\\
2957	-0.431975868284098\\
2958	-0.543804705885577\\
2959	-0.646696852861623\\
2960	-0.738961416735947\\
2961	-0.819082154905118\\
2962	-0.885742392022621\\
2963	-0.937846657766386\\
2964	-0.974538689404199\\
2965	-0.995215503307008\\
2966	-0.999537304163823\\
2967	-0.987433069056351\\
2968	-0.959101714624012\\
2969	-0.91500882813898\\
2970	-0.855879016211929\\
2971	-0.782683996871128\\
2972	-0.696626630694936\\
2973	-0.599121153436045\\
2974	-0.491769934979857\\
2975	-0.376337146577498\\
2976	-0.254719769086933\\
2977	-0.128916418676454\\
2979	0.126943757399658\\
2980	0.252795869107103\\
2981	0.374493624573915\\
2982	0.490037086729444\\
2983	0.597527455941417\\
2984	0.695198274207542\\
2985	0.781444454486973\\
2986	0.854848658123501\\
2987	0.914204586879805\\
2988	0.958536806806023\\
2989	0.98711677817073\\
2990	0.999474828020539\\
2991	0.995407868616894\\
2992	0.974982734905097\\
2993	0.938535086174397\\
2994	0.886663889953525\\
2995	0.820221578793735\\
2996	0.740300041700266\\
2997	0.648212680421693\\
2998	0.545472825477191\\
2999	0.433768866630544\\
3000	0.31493650649918\\
3001	0.190928593278386\\
3002	0.0637830283399126\\
3003	-0.0644107239086225\\
3004	-0.191545973524171\\
3005	-0.315533425617559\\
3006	-0.434335515068142\\
3007	-0.545999891137399\\
3008	-0.648691501696703\\
3009	-0.740722749814722\\
3010	-0.820581227113962\\
3011	-0.886954568149122\\
3012	-0.938752017344996\\
3013	-0.975122354081122\\
3014	-0.995467881347849\\
3015	-0.999454248078109\\
3016	-0.987015943758706\\
3017	-0.95835737500056\\
3018	-0.913949506399149\\
3019	-0.854522120867841\\
3020	-0.781051826651947\\
3021	-0.694746008097809\\
3022	-0.597022983934949\\
3023	-0.489488699137382\\
3024	-0.373910333401\\
3025	-0.252187259950006\\
3026	-0.126319831921592\\
3028	0.12954013349281\\
3029	0.255327960383511\\
3030	0.37691981955777\\
3031	0.492317514206206\\
3032	0.599624640187812\\
3033	0.697077750853168\\
3034	0.783075336889397\\
3035	0.856204144949061\\
3036	0.915262402546432\\
3037	0.959279567551675\\
3038	0.987532277732953\\
3039	0.99955623823007\\
3040	0.995153851607029\\
3041	0.974397465099628\\
3042	0.937628181686705\\
3043	0.885450254525495\\
3044	0.818721156872471\\
3045	0.738537490679391\\
3046	0.646216965429176\\
3047	0.543276743374918\\
3048	0.431408507040032\\
3049	0.312450658719627\\
3050	0.188358108837292\\
3051	0.0611701496518435\\
3052	-0.0670230577361508\\
3053	-0.194114832338528\\
3054	-0.318016593718312\\
3055	-0.436692184965068\\
3056	-0.548191334166404\\
3057	-0.6506817044733\\
3058	-0.742479006062695\\
3059	-0.822074675151271\\
3060	-0.888160665188479\\
3061	-0.939650942822027\\
3062	-0.975699335378977\\
3063	-0.99571343656271\\
3064	-0.999364341845194\\
3065	-0.986592053563982\\
3066	-0.957606466903144\\
3067	-0.912883920551849\\
3068	-0.853159368724846\\
3069	-0.779414303190606\\
3070	-0.692860623789329\\
3071	-0.594920722505321\\
3072	-0.487204108394053\\
3073	-0.371480957484891\\
3074	-0.249653022349776\\
3075	-0.123722379385072\\
3077	0.132135621732687\\
3078	0.257858301671149\\
3079	0.379343431175585\\
3080	0.49459456739396\\
3081	0.601717714674123\\
3082	0.698952449805347\\
3083	0.784700852180777\\
3084	0.857553763447413\\
3085	0.916313945106594\\
3086	0.960015753502375\\
3087	0.987941008860162\\
3088	0.999630797592545\\
3089	0.994893013923502\\
3090	0.973805516883203\\
3091	0.936714850800854\\
3092	0.884230550320353\\
3093	0.817215123528058\\
3094	0.736769877804818\\
3095	0.644216821337977\\
3096	0.541076937714479\\
3097	0.429045190624038\\
3098	0.309962669438846\\
3099	0.185786333410761\\
3100	0.0585568517103638\\
3101	-0.0696349321951857\\
3102	-0.196682360711293\\
3103	-0.320497582168628\\
3104	-0.439045861822706\\
3105	-0.550379019953652\\
3106	-0.652667447549902\\
3107	-0.744230173442702\\
3108	-0.823562488780681\\
3109	-0.889360674873842\\
3110	-0.940543428036563\\
3111	-0.976269629342369\\
3112	-0.995952167269479\\
3113	-0.999267586080805\\
3114	-0.986161401377558\\
3115	-0.95684899547814\\
3116	-0.911812077900322\\
3117	-0.851790769122999\\
3118	-0.777771437711181\\
3119	-0.690970490691598\\
3120	-0.592814383556288\\
3121	-0.484916178408184\\
3122	-0.369049035479748\\
3123	-0.247117073655772\\
3124	-0.121124078869343\\
3126	0.134730204330026\\
3127	0.260386875626409\\
3128	0.381764442815893\\
3129	0.496868230685777\\
3130	0.603806665054435\\
3131	0.700822358216101\\
3132	0.786320989220712\\
3133	0.858897504367633\\
3134	0.917359207353911\\
3135	0.960745359613156\\
3136	0.988342968749748\\
3137	0.999698505597735\\
3138	0.994625357355289\\
3139	0.97320689431308\\
3140	0.935795099775987\\
3141	0.883004785698176\\
3142	0.815703489082807\\
3143	0.734997215192379\\
3144	0.642212261856912\\
3145	0.538873423573023\\
3146	0.426678933580661\\
3147	0.307472555708046\\
3148	0.183213284625253\\
3149	0.0559431524270622\\
3150	-0.0722463293841429\\
3151	-0.199248541045563\\
3152	-0.322976373964593\\
3153	-0.441396529509348\\
3154	-0.55256293350476\\
3155	-0.654648717317286\\
3156	-0.745976239952142\\
3157	-0.825044657804938\\
3158	-0.890554588981104\\
3159	-0.94142946687225\\
3160	-0.9768332320632\\
3161	-0.996184071831976\\
3162	-0.999163981448419\\
3163	-0.985723990151655\\
3164	-0.956084965916943\\
3165	-0.910733985791012\\
3166	-0.850416331442375\\
3167	-0.776123241473215\\
3168	-0.689075621759457\\
3169	-0.590703981523347\\
3170	-0.482624924861284\\
3171	-0.366614584054332\\
3172	-0.244579431248894\\
3173	-0.118524948182767\\
3175	0.137323863501024\\
3176	0.262913664919324\\
3177	0.384182837885419\\
3178	0.499138488498829\\
3179	0.605891477011482\\
3180	0.70268746326883\\
3181	0.787935736904728\\
3182	0.860235358501086\\
3183	0.918398182123383\\
3184	0.961468380882252\\
3185	0.988738154647763\\
3186	0.999759361781344\\
3187	0.994350883735933\\
3188	0.972601601491533\\
3189	0.934868934916267\\
3190	0.881772969059966\\
3191	0.814186263897227\\
3192	0.733219514991106\\
3193	0.640203300725261\\
3194	0.536666216054073\\
3195	0.424309752127101\\
3196	0.304980334595257\\
3197	0.180638980116782\\
3198	0.0533290697158009\\
3199	-0.0748572314050762\\
3200	-0.20181335575262\\
3201	0.932926834027967\\
3202	0.977215368166981\\
3203	0.998235426874999\\
3204	0.995486501526557\\
3205	0.9690340467846\\
3206	0.919507922057164\\
3207	0.848087393915648\\
3208	0.756473056576397\\
3209	0.646846339055173\\
3210	0.521817563156219\\
3211	0.384363789100462\\
3212	0.237757928744031\\
3213	0.0854908142737258\\
3214	-0.0688119220017143\\
3215	-0.22147617758128\\
3216	-0.368866863828316\\
3217	-0.507474461067432\\
3218	-0.633998583758512\\
3219	-0.745426565976686\\
3220	-0.839105196003402\\
3221	-0.912803891950716\\
3222	-0.964767814145944\\
3223	-0.993759649622916\\
3224	-0.999089073777213\\
3225	-0.980629187688919\\
3226	-0.938819539711403\\
3227	-0.874655659385553\\
3228	-0.789665352885095\\
3229	-0.685872324424963\\
3230	-0.565747989837291\\
3231	-0.432152629689881\\
3232	-0.288267283144251\\
3233	-0.137518004230969\\
3235	0.17013641847052\\
3236	0.319716002705263\\
3237	0.461682829101846\\
3238	0.592656525512211\\
3239	0.709518477211986\\
3240	0.809486084232958\\
3241	0.89017901785337\\
3242	0.949675898606529\\
3243	0.986560046246268\\
3244	0.999953212321998\\
3245	0.989536492151728\\
3246	0.955557918263366\\
3247	0.89882655448946\\
3248	0.82069323135056\\
3249	0.723018381428119\\
3250	0.608127740603322\\
3251	0.478756969960159\\
3252	0.337986516958154\\
3253	0.189168266897013\\
3254	0.0358457311749589\\
3255	-0.118330327262811\\
3256	-0.269688822242188\\
3257	-0.414625756592613\\
3258	-0.549690036924858\\
3259	-0.671665647619648\\
3260	-0.777648227372083\\
3261	-0.865114224944136\\
3262	-0.931980987440511\\
3263	-0.976656350358098\\
3264	-0.998076548617973\\
3265	-0.995731545871877\\
3266	-0.969677178981328\\
3267	-0.920533828484167\\
3268	-0.849471646714392\\
3269	-0.758182695303731\\
3270	-0.648840655502227\\
3271	-0.524049070653746\\
3272	-0.386779353222664\\
3273	-0.240300032492996\\
3274	-0.0880989276161017\\
3275	0.0661999008534622\\
3276	0.218922443463271\\
3277	0.366432223704578\\
3278	0.505216886154813\\
3279	0.631971829173835\\
3280	0.743678890773936\\
3281	0.837678214077641\\
3282	0.911731581168624\\
3283	0.964075707300253\\
3284	0.993464226469769\\
3285	0.999197368637397\\
3286	0.981138621953505\\
3287	0.939717983240826\\
3288	0.875921719340568\\
3289	0.79126888311248\\
3290	0.687775143265299\\
3291	0.56790478927087\\
3292	0.434512054164315\\
3293	0.290773152402835\\
3294	0.140110651020223\\
3296	-0.167556013484045\\
3297	-0.317234325561913\\
3298	-0.459358971017082\\
3299	-0.59054581987175\\
3300	-0.707671182028662\\
3301	-0.807946185452693\\
3302	-0.888983182006541\\
3303	-0.948852599740803\\
3304	-0.9861288879315\\
3305	-0.999924460868442\\
3306	-0.989910832160604\\
3307	-0.956326436324161\\
3308	-0.899970951418709\\
3309	-0.822186257913017\\
3310	-0.724824487167098\\
3311	-0.610203920335152\\
3312	-0.48105378776836\\
3313	-0.340449283313319\\
3314	-0.191738340868596\\
3315	-0.0384619167307392\\
3316	0.115730324117976\\
3317	0.267166910183732\\
3318	0.41224198486907\\
3319	0.547501165526228\\
3320	0.66972379576282\\
3321	0.775999632492585\\
3322	0.86379814173506\\
3323	0.931028753159353\\
3324	0.976090638656387\\
3325	0.997910829656121\\
3326	0.99596976558496\\
3327	0.970313665119193\\
3328	0.921553425675484\\
3329	0.850850077330051\\
3330	0.759887137531678\\
3331	0.650830524868525\\
3332	0.526276986377525\\
3333	0.389192266402461\\
3334	0.242840489252558\\
3335	0.0907064371385786\\
3336	-0.0635874259783122\\
3337	-0.216367208875454\\
3338	-0.363995072095804\\
3339	-0.502955848542115\\
3340	-0.629940743124735\\
3341	-0.741926118479569\\
3342	-0.836245490799229\\
3343	-0.910653021480357\\
3344	-0.963376992787744\\
3345	-0.99316199422401\\
3346	-0.999298815110706\\
3347	-0.981641331603441\\
3348	-0.940609986048457\\
3349	-0.877181775826557\\
3350	-0.792866990071161\\
3351	-0.689673248171857\\
3352	-0.570057696348158\\
3353	-0.436868500542005\\
3354	-0.293277028734792\\
3355	-0.142702337506762\\
3357	0.164974460086796\\
3358	0.314750474130051\\
3359	0.457031964537236\\
3360	0.588431066696558\\
3361	-0.906359271198653\\
3362	-0.844860388938741\\
3363	-0.769477375416045\\
3364	-0.681449047916885\\
3365	-0.582222032262507\\
3366	-0.473426989480686\\
3367	-0.356851818110499\\
3368	-0.234412272520331\\
3369	-0.108120480092566\\
3371	0.147688910493798\\
3372	0.273002629285202\\
3373	0.393829921219094\\
3374	0.508185154108105\\
3375	0.614189055253064\\
3376	0.710099594731219\\
3377	0.794340613289478\\
3378	0.865527724377444\\
3379	0.922491064657152\\
3380	0.964294519116265\\
3381	0.990251104843992\\
3382	0.99993426066203\\
3383	0.993184857075448\\
3384	0.970113811351894\\
3385	0.9311002647446\\
3386	0.87678535182431\\
3387	0.808061664301931\\
3388	0.726058582500173\\
3389	0.632123715522539\\
3390	0.527800755132375\\
3391	0.414804107277632\\
3392	0.294990718160989\\
3393	0.170329557854984\\
3396	-0.212058606871778\\
3397	-0.335336789032226\\
3398	-0.453104167209403\\
3399	-0.563425394737806\\
3400	-0.664487492308581\\
3401	-0.754629641817701\\
3402	-0.83237047965622\\
3403	-0.896432440900753\\
3404	-0.945762754359293\\
3405	-0.97955074343281\\
3406	-0.997241148480498\\
3407	-0.998543251759656\\
3408	-0.983435654973619\\
3409	-0.952166630925603\\
3410	-0.905250043491378\\
3411	-0.843456902967318\\
3412	-0.76780269556275\\
3413	-0.679530695267204\\
3414	-0.580091532332972\\
3415	-0.471119354139773\\
3416	-0.354404970205451\\
3417	-0.231866422668645\\
3418	-0.105517465873163\\
3420	0.150277686017489\\
3421	0.275520235739805\\
3422	0.396234985166302\\
3423	0.510438151594371\\
3424	0.616252961328883\\
3425	0.711940491907626\\
3426	0.795928248925065\\
3427	0.866836007840902\\
3428	0.923498496088996\\
3429	0.964984542751154\\
3430	0.990612381082428\\
3431	0.999960852420372\\
3432	0.992876327354224\\
3433	0.969475230417174\\
3434	0.930142126806004\\
3435	0.875523402575254\\
3436	0.806516642161569\\
3437	0.724255877805717\\
3438	0.630092953274016\\
3439	0.525575308144653\\
3440	0.412420547744659\\
3441	0.292488216641232\\
3442	0.167749239557452\\
3445	-0.214616329949877\\
3446	-0.337802044659384\\
3447	-0.455436442262908\\
3448	-0.565586361449732\\
3449	-0.666441638129982\\
3450	-0.756344853020892\\
3451	-0.833818569077721\\
3452	-0.897589611160583\\
3453	-0.946609988939144\\
3454	-0.980074119183882\\
3455	-0.997432064434634\\
3456	-0.998398570473\\
3457	-0.982957754085874\\
3458	-0.951363364087683\\
3459	-0.904134611301743\\
3460	-0.842047636036568\\
3461	-0.766122753275795\\
3462	-0.677607685190196\\
3463	-0.577957056520972\\
3464	-0.468808489799358\\
3465	-0.351955693248328\\
3466	-0.229318983630037\\
3467	-0.102913728449039\\
3469	0.152865431554801\\
3470	0.278035953809194\\
3471	0.398637333363695\\
3472	0.512687650594216\\
3473	0.618312643675381\\
3474	0.713776509523086\\
3475	0.797510429357771\\
3476	0.868138350107074\\
3477	0.924499597965223\\
3478	0.965667952490094\\
3479	0.990966867774659\\
3480	0.999980590558607\\
3481	0.992560992569906\\
3482	0.968830004808297\\
3483	0.929177613777483\\
3484	0.874255452586567\\
3485	0.804966092246559\\
3486	0.72244820914284\\
3487	0.628057872438603\\
3488	0.523346258922629\\
3489	0.410034161527619\\
3490	0.289983710440538\\
3491	0.165167771525375\\
3494	-0.217172582071726\\
3495	-0.34026498502908\\
3496	-0.457765595805995\\
3497	-0.56774345169606\\
3498	-0.66839121623525\\
3499	-0.758054880321197\\
3500	-0.835260943600133\\
3501	-0.898740629442273\\
3502	-0.947450735559869\\
3503	-0.980590777616726\\
3504	-0.997616144101357\\
3505	-0.998247046274173\\
3506	-0.982473116115671\\
3507	-0.950553576712537\\
3508	-0.903012982274959\\
3509	-0.840632597805325\\
3510	-0.764437560068927\\
3511	-0.67568003086626\\
3512	-0.575818619456186\\
3513	-0.466494412297834\\
3514	-0.349504004026585\\
3515	-0.226769972865441\\
3516	-0.100309285665844\\
3518	0.155452129369678\\
3519	0.28054976625117\\
3520	0.401036949344871\\
3521	0.356668715480737\\
3522	0.496152114591496\\
3523	0.623821635549575\\
3524	0.736637339352001\\
3525	0.831912971240854\\
3526	0.907379923116423\\
3527	0.96124125136339\\
3528	0.992214463879009\\
3529	0.999562057490948\\
3530	0.983109078644702\\
3531	0.943247289224473\\
3532	0.880925838311668\\
3533	0.797628661991894\\
3534	0.695339149356187\\
3535	0.576492916021834\\
3536	0.443919809687031\\
3537	0.300776528627466\\
3538	0.150471457552612\\
3540	-0.15722308641034\\
3541	-0.307286044417651\\
3542	-0.450032214326711\\
3543	-0.58206266704201\\
3544	-0.700233625414057\\
3545	-0.801731320716044\\
3546	-0.884138991294094\\
3547	-0.945494428079201\\
3548	-0.98433669674705\\
3549	-0.9997409240259\\
3550	-0.991340319856135\\
3551	-0.959334911025508\\
3552	-0.904486778334103\\
3553	-0.828101910685746\\
3554	-0.73199910818812\\
3555	-0.618466674700358\\
3556	-0.490207931031364\\
3557	-0.350276846165798\\
3558	-0.202005319209093\\
3559	-0.0489238435343395\\
3560	0.105322557793897\\
3561	0.257061123665608\\
3562	0.40267880704323\\
3563	0.538708305222826\\
3564	0.661910619780883\\
3565	0.769352180368514\\
3566	0.858474695952737\\
3567	0.92715607028822\\
3568	0.973760931149172\\
3569	0.997179570179469\\
3570	0.996854366158004\\
3571	0.972793062518122\\
3572	0.925568582967571\\
3573	0.856305389603676\\
3574	0.766652708344736\\
3575	0.658745259212537\\
3576	0.535152426510649\\
3577	0.398817079216769\\
3578	0.25298549832678\\
3579	0.101130079659015\\
3580	-0.0531333473668383\\
3581	-0.206131616242146\\
3582	-0.354221685139692\\
3583	-0.493877381383754\\
3584	-0.62177336314835\\
3585	-0.734864299171477\\
3586	-0.830457381141514\\
3587	-0.906276442154194\\
3588	-0.960516154524612\\
3589	-0.991885016444485\\
3590	-0.999636103932971\\
3591	-0.983584855844128\\
3592	-0.944113468450269\\
3593	-0.882161794969306\\
3594	-0.799204966717753\\
3595	-0.697218268757752\\
3596	-0.578630106386754\\
3597	-0.446264182374307\\
3598	-0.303272261778602\\
3599	-0.153059125264463\\
3601	0.154637116386766\\
3602	0.304793666219211\\
3603	0.447692773971539\\
3604	0.579931868946005\\
3605	0.698362206010188\\
3606	0.800163840372988\\
3607	0.882912773286989\\
3608	0.94463866988508\\
3609	0.983871774823911\\
3610	0.999677908630929\\
3611	0.991680711448225\\
3612	0.960070604540761\\
3613	0.905600256173329\\
3614	0.829566659828743\\
3615	0.73378025149168\\
3616	0.620521801361065\\
3617	0.492488106427572\\
3618	0.352727777040855\\
3619	0.20456864644666\\
3620	0.0515385317494292\\
3621	-0.102718766943781\\
3622	-0.254530229045031\\
3623	-0.400281071784775\\
3624	-0.536500821801383\\
3625	-0.659945950585097\\
3626	-0.767677106136944\\
3627	-0.8571291018784\\
3628	-0.926171996312405\\
3629	-0.973161809057274\\
3630	-0.996979665667823\\
3631	-0.997058439153079\\
3632	-0.973396253837564\\
3633	-0.926556530022935\\
3634	-0.857654568387261\\
3635	-0.768330993560539\\
3636	-0.660712689209959\\
3637	-0.537362154813763\\
3638	-0.401216489984563\\
3639	-0.255517459188468\\
3640	-0.103734302093471\\
3641	0.050518872500561\\
3642	0.203569142205652\\
3643	0.351772227003494\\
3644	0.491599263195894\\
3645	0.61972082918146\\
3646	0.733086222312977\\
3647	0.828996099180131\\
3648	0.905166749675118\\
3649	0.959784474415301\\
3650	0.991548770740792\\
3651	0.999703298981331\\
3652	0.984053891662825\\
3653	0.944973176827716\\
3654	0.883391705388931\\
3655	0.800775793782122\\
3656	0.699092609503623\\
3657	0.58076333088502\\
3658	0.448605496416803\\
3659	0.305765916335531\\
3660	0.15564574392647\\
3662	-0.152050086497184\\
3663	-0.302299198999208\\
3664	-0.445350265180878\\
3665	-0.577797096061659\\
3666	-0.696486000110326\\
3667	-0.79859087579689\\
3668	-0.881680503894131\\
3669	-0.943776437242832\\
3670	-0.983400109554623\\
3671	-0.999608041555348\\
3672	-0.992014306171768\\
3673	-0.960799717839109\\
3674	-0.906707527130038\\
3675	-0.831025723214225\\
3676	-0.735556365547382\\
3677	-0.622572675034917\\
3678	-0.494764906366072\\
3679	-0.355176290359395\\
3680	-0.207130571594007\\
3681	-0.0541528667249622\\
3682	0.100114272070641\\
3683	0.251997589903112\\
3684	0.397880593044647\\
3685	0.534289661263301\\
3686	0.657976758193399\\
3687	0.765996770332094\\
3688	0.855777633137222\\
3689	0.925181574456929\\
3690	0.97255601702318\\
3691	0.996772927969232\\
3692	0.997255678421425\\
3693	0.973992773608643\\
3694	0.927538126563832\\
3695	0.858997868902861\\
3696	0.770004012721529\\
3697	0.662675590756862\\
3698	0.539568200097165\\
3699	0.403613150859201\\
3700	0.258047668762174\\
3701	0.106337813544997\\
3702	-0.0479040513832842\\
3703	-0.201005272928796\\
3704	-0.349320357859597\\
3705	-0.489317775641666\\
3706	-0.617664047717881\\
3707	-0.731303120964185\\
3708	-0.827529135372515\\
3709	-0.904050853284389\\
3710	-0.959046216050865\\
3711	-0.99120572907259\\
3712	-0.999763642174912\\
3713	-0.984516182886637\\
3714	-0.945826408465109\\
3715	-0.88461556114089\\
3716	-0.80234113241886\\
3717	-0.700962158746734\\
3718	-0.58289257489696\\
3719	-0.450943735767396\\
3720	-0.308257475207029\\
3721	-0.158231295809856\\
3723	0.149462014473556\\
3724	0.299802659854322\\
3725	0.443004704009127\\
3726	0.575658363020466\\
3727	0.694605020573363\\
3728	0.797012437767989\\
3729	0.880442191562452\\
3730	0.942907736062807\\
3731	0.982921704170622\\
3732	0.999531323278006\\
3733	0.992341101740294\\
3734	0.961522245924243\\
3735	0.907808583614951\\
3736	0.832479090842753\\
3737	0.737327438182547\\
3738	0.624619281665218\\
3739	0.497038315240843\\
3740	0.357622369338969\\
3741	0.20969107709152\\
3742	0.0567668305425286\\
3743	-0.0975090910255858\\
3744	-0.249463223597559\\
3745	-0.395477387275605\\
3746	-0.532074838763947\\
3747	-0.656003056102691\\
3748	-0.764311184471353\\
3749	-0.854420298991954\\
3750	-0.924184811511168\\
3751	-0.97194355919919\\
3752	-0.99655935850069\\
3753	-0.997446082611077\\
3754	-0.974582617741817\\
3755	-0.928513365862273\\
3756	-0.860335281943208\\
3757	-0.77167175436125\\
3758	-0.664633950399548\\
3759	-0.541770547240503\\
3760	-0.406007045414754\\
3761	-0.260576109706108\\
3762	-0.108940596169305\\
3763	0.0452889019366012\\
3764	0.198440025984382\\
3765	0.346866094512734\\
3766	0.487032934358012\\
3767	0.615603032853869\\
3768	0.729515007345071\\
3769	0.826056499772676\\
3770	0.902928760629948\\
3771	0.958301384491278\\
3772	0.990855893791377\\
3773	0.999817133100805\\
3774	0.98497172634643\\
3775	0.946673157513942\\
3776	0.885833353836915\\
3777	0.803900971899566\\
3778	0.702826903673667\\
3779	0.585017823828366\\
3780	0.453278884400333\\
3781	0.310746921316422\\
3782	0.160815763194023\\
3784	-0.146872918053759\\
3785	-0.297304065895332\\
3786	-0.440656106532515\\
3787	-0.573515684481663\\
3788	-0.692719280291385\\
3789	-0.795428537105636\\
3790	-0.879197844777991\\
3791	-0.942032572299013\\
3792	-0.98243656195109\\
3793	-0.999447754325047\\
3794	-0.992661095914173\\
3795	-0.962238183843056\\
3796	-0.908903418081536\\
3797	-0.833926752752632\\
3798	-0.739093457257695\\
3799	-0.626661607225287\\
3800	-0.499308317470422\\
3801	-0.360065997215315\\
3802	-0.212250145389589\\
3803	-0.0593804052869018\\
3804	0.0949032416642694\\
3805	0.246927147498809\\
3806	0.393071470948598\\
3807	0.529856369482786\\
3808	0.654024857840795\\
3809	0.762620360107576\\
3810	0.853057108745816\\
3811	0.923181714305883\\
3812	0.97132443978262\\
3813	0.996338958726028\\
3814	0.997629650416457\\
3815	0.975165782194836\\
3816	0.929482241233927\\
3817	0.861666798341503\\
3818	0.773334207049174\\
3819	0.666587754715238\\
3820	0.54396918114935\\
3821	0.408398157243482\\
3822	0.263102764690302\\
3823	0.111542632126202\\
3824	-0.0426734420852881\\
3825	-0.195873418954307\\
3826	-0.34440945378492\\
3827	-0.484744755005067\\
3828	-0.613537798715242\\
3829	-0.727721893711987\\
3830	-0.824578202473731\\
3831	-0.901800479403391\\
3832	-0.957549984840625\\
3833	-0.990499267295036\\
3834	-0.999863771392029\\
3835	-0.985420518920819\\
3836	-0.947513418171184\\
3837	-0.887045075130118\\
3838	-0.805455301533129\\
3839	-0.70468683150375\\
3840	-0.587139063113227\\
3841	-0.455610926310328\\
3842	-0.313234237601137\\
3843	-0.16339912836429\\
3845	0.144282814983853\\
3846	0.294803434248024\\
3847	0.438304488848644\\
3848	0.571369075130406\\
3849	0.690828792188768\\
3850	0.793839184664648\\
3851	0.877947472070446\\
3852	0.941150951949112\\
3853	0.981944686221595\\
3854	0.999357335269451\\
3855	0.992974286499702\\
3856	0.962947526689732\\
3857	0.909992023025552\\
3858	0.835368699021728\\
3859	0.740854410669272\\
3860	0.628699637716636\\
3861	0.501574897496994\\
3862	0.362507157239634\\
3863	0.214807758949519\\
3864	0.0619935730446741\\
3865	-0.0922967418464395\\
3866	-0.244389378989126\\
3867	-0.390662860553675\\
3868	-0.527634268625661\\
3869	-0.652042176965551\\
3870	-0.76092430882909\\
3871	-0.851688071741592\\
3872	-0.9221722897164\\
3873	-0.970698663017629\\
3874	-0.996111730155462\\
3875	-0.997806380580641\\
3876	-0.97574226297138\\
3877	-0.930444746038575\\
3878	-0.862992408972332\\
3879	-0.774991359391151\\
3880	-0.668536990313441\\
3881	-0.546164086753834\\
3882	-0.410786469957202\\
3883	-0.265627616397978\\
3884	-0.114143903583226\\
3885	0.0400576897550309\\
3886	0.193305469430015\\
3887	0.341950452513174\\
3888	0.482453253266158\\
3889	0.611468359457831\\
3890	0.725923792354479\\
3891	0.82309425360836\\
3892	0.900666017337244\\
3893	0.956792022250283\\
3894	0.990135852027834\\
3895	0.999903556728441\\
3896	0.985862557532982\\
3897	0.948347184677004\\
3898	0.888250716715902\\
3899	0.807004110666185\\
3900	0.706541929489049\\
3901	0.58925627821236\\
3902	0.45793984551392\\
3903	0.315719407013603\\
3904	0.165981373615523\\
3906	-0.141691723015356\\
3907	-0.292300782050461\\
3908	-0.435949867074669\\
3909	-0.56921854967959\\
3910	-0.688933569223536\\
3911	-0.792244391338954\\
3912	-0.87669108200862\\
3913	-0.940262881055787\\
3914	-0.98144608035318\\
3915	-0.99926006673013\\
3916	-0.993280671350931\\
3917	-0.963650269601203\\
3918	-0.911074390986414\\
3919	-0.836804919767474\\
3920	-0.742610286347372\\
3921	-0.63073335917079\\
3922	-0.503838039785478\\
3923	-0.36494583267995\\
3924	-0.217363900240343\\
3925	-0.0646063159056212\\
3926	0.0896896094373005\\
3927	0.241849935462142\\
3928	0.388251572599529\\
3929	0.525408551422515\\
3930	0.650055027066173\\
3931	0.759223042260601\\
3932	0.850313197362993\\
3933	0.921156544661244\\
3934	0.970066233192028\\
3935	0.995877674346502\\
3936	0.99797627189082\\
3937	0.976312056119241\\
3938	0.931400873679195\\
3939	0.864312104749843\\
3940	0.776643200028957\\
3941	0.670481643833682\\
3942	0.548355249011365\\
3943	0.413171967185917\\
3944	0.268150647523271\\
3945	0.116744392710189\\
3946	-0.037441662873789\\
3947	-0.190736195011141\\
3948	-0.339489107550889\\
3949	-0.480158444846438\\
3950	-0.609394729264295\\
3951	-0.724120715596655\\
3952	-0.821604663347443\\
3953	-0.899525382206775\\
3954	-0.956027501914377\\
3955	-0.989765650479967\\
3956	-0.999936488838102\\
3957	-0.986297839153394\\
3958	-0.949174451317958\\
3959	-0.889450270330599\\
3960	-0.80854738868311\\
3961	-0.708392184915283\\
3962	-0.591369454615233\\
3963	-0.460265626049477\\
3964	-0.31820241252035\\
3965	-0.16856248124941\\
3967	0.13909965990797\\
3968	0.289796126456622\\
3969	0.433592257349119\\
3970	0.567064122868487\\
3971	0.687033624384185\\
3972	0.790644168058861\\
3973	0.875428683204973\\
3974	0.939368365706287\\
3975	0.98094074776327\\
3976	0.999155949374199\\
3977	0.993580248367834\\
3978	0.96434640776215\\
3979	0.912150514544919\\
3980	0.838235405146406\\
3981	0.744361072258016\\
3982	0.632762757649743\\
3983	0.506097728824443\\
3984	0.367382006822936\\
3985	0.219918551743831\\
3986	0.0672186159617922\\
3987	-0.0870818623056948\\
3988	-0.239308834322856\\
3989	-0.385837623612588\\
3990	-0.523179233127848\\
3991	-0.648063421762345\\
3992	-0.757516572062286\\
3993	-0.848932495032386\\
3994	-0.920134486102597\\
3995	-0.96942715464138\\
3996	-0.995636792903497\\
3997	-0.99813932318375\\
3998	-0.976875157733502\\
3999	-0.932350617602424\\
4000	-0.86562587662911\\
};
\end{axis}
\end{tikzpicture}%

\section{Démodulation par filtrage}

On souhaite désormais reconstituer le signal de départ. Pour cela, nous allons procéder à un filtrage passe bas d'une part et passe haut d'autre part, avec une fréquence de coupure $F_c=\frac{F_0+F_1}{2}$. Nous ferons ensuite passer chacun des signaux filtrés par un détecteur d'énergie qui permettra de reproduire de signal binaire initial de manière fidèle.

\paragraph{Filtre passe haut}
Pour le passe haut, nous allons utiliser un filtre de réponse impulsionnelle suivante:
\[
h_\text{haut}(t) = \frac{2 F_c}{F_e} \sinc(2 F_c t)
\]
\\
et donc \[H_\text{haut}(f) = \TF(h_\text{haut}(t))
\]

\paragraph{Filtre passe bas}

\[
H_\text{bas}(f) = 1-H_\text{haut}(f)
\]

Les réponses des filtres sont les suivantes:

%%%% This file was created by matlab2tikz.
%
%The latest updates can be retrieved from
%  http://www.mathworks.com/matlabcentral/fileexchange/22022-matlab2tikz-matlab2tikz
%where you can also make suggestions and rate matlab2tikz.
%
\definecolor{mycolor1}{rgb}{0.00000,0.44700,0.74100}%
\definecolor{mycolor2}{rgb}{0.85000,0.32500,0.09800}%
%
\begin{tikzpicture}

\begin{axis}[%
width=4.521in,
height=3.548in,
at={(0.758in,0.499in)},
scale only axis,
xmin=-0.0025,
xmax=0.0025,
ymin=-0.2,
ymax=1,
axis background/.style={fill=white}
]
\addplot [color=mycolor1, forget plot]
  table[row sep=crcr]{%
-0.00208333333333333	-0.00318309886183787\\
-0.00206249999999997	-0.00318317492734366\\
-0.00204166666666672	-0.00311909156929424\\
-0.00202083333333336	-0.00299081106446764\\
-0.002	-0.00279956172779527\\
-0.00197916666666664	-0.00254783933014335\\
-0.00195833333333328	-0.00223938409131019\\
-0.00193750000000004	-0.00187913323343325\\
-0.00191666666666668	-0.00147314954171185\\
-0.00189583333333332	-0.00102852684019938\\
-0.00185416666666671	-5.61774650087443e-05\\
-0.00177083333333339	0.00195666421959917\\
-0.00175000000000003	0.00241545663842524\\
-0.00172916666666667	0.00283652879109197\\
-0.00170833333333331	0.00321058443246536\\
-0.00168749999999995	0.00352902537568234\\
-0.00166666666666671	0.00378413364320329\\
-0.00164583333333335	0.00396924098084706\\
-0.00162499999999999	0.00407888230089704\\
-0.00160416666666663	0.00410892985197497\\
-0.00158333333333338	0.00405670520761037\\
-0.00156250000000002	0.00392106651788271\\
-0.00154166666666666	0.00370246887363823\\
-0.0015208333333333	0.00340299608403061\\
-0.00149999999999995	0.0030263626579945\\
-0.0014791666666667	0.00257788530046443\\
-0.00145833333333334	0.00206442377572313\\
-0.00143749999999998	0.00149429154376235\\
-0.00139583333333337	0.000223797745730003\\
-0.00129166666666669	-0.00314668345950841\\
-0.00127083333333333	-0.00374732295256697\\
-0.00124999999999997	-0.00429196845667068\\
-0.00122916666666661	-0.00476798876485895\\
-0.00120833333333337	-0.00516364914617817\\
-0.00118750000000001	-0.00546835616486996\\
-0.00116666666666665	-0.00567288882739969\\
-0.00114583333333329	-0.00576961169154977\\
-0.00112500000000004	-0.00575266580402922\\
-0.00110416666666668	-0.00561813364122188\\
-0.00108333333333333	-0.00536417460743954\\
-0.00106249999999997	-0.0049911280907049\\
-0.00104166666666672	-0.00450158158078551\\
-0.00102083333333336	-0.00390040190991692\\
-0.001	-0.00319472827441547\\
-0.000979166666666642	-0.00239392632536617\\
-0.000958333333333283	-0.00150950326827604\\
-0.000916666666666677	0.000454246434654526\\
-0.000812499999999994	0.00567989546351644\\
-0.000791666666666635	0.00661879347518046\\
-0.000770833333333387	0.00747281077017503\\
-0.000750000000000028	0.00822102909442923\\
-0.00072916666666667	0.00884328453400929\\
-0.000708333333333311	0.00932050634240345\\
-0.000687499999999952	0.00963504613361477\\
-0.000666666666666704	0.00977099196810116\\
-0.000645833333333345	0.00971446202866277\\
-0.000624999999999987	0.00945387284344712\\
-0.000604166666666628	0.00898017735982271\\
-0.00058333333333338	0.00828706860093087\\
-0.000562500000000021	0.00737114513983217\\
-0.000541666666666663	0.00623203519655136\\
-0.000520833333333304	0.00487247679202218\\
-0.000499999999999945	0.00329835206993767\\
-0.000479166666666697	0.00151867461190636\\
-0.000458333333333338	-0.000454470688504793\\
-0.00043749999999998	-0.00260603487341049\\
-0.000395833333333373	-0.00737036774129618\\
-0.000354166666666655	-0.0126015609750331\\
-0.000229166666666614	-0.0289336924620972\\
-0.000187500000000007	-0.0338042628617465\\
-0.000166666666666648	-0.0360019245108092\\
-0.00014583333333329	-0.0380065351526726\\
-0.000125000000000042	-0.0397946281998789\\
-0.000104166666666683	-0.0413451468692532\\
-8.33333333333242e-05	-0.0426397415309975\\
-6.24999999999654e-05	-0.043663029550196\\
-4.16666666667176e-05	-0.0444028136048185\\
-2.08333333333588e-05	-0.0448502551012079\\
0	0.955\\
2.08333333333588e-05	-0.0448502551012079\\
4.16666666667176e-05	-0.0444028136048185\\
6.24999999999654e-05	-0.043663029550196\\
8.33333333333242e-05	-0.0426397415309975\\
0.000104166666666683	-0.0413451468692532\\
0.000125000000000042	-0.0397946281998789\\
0.00014583333333329	-0.0380065351526726\\
0.000166666666666648	-0.0360019245108092\\
0.000187500000000007	-0.0338042628617465\\
0.000229166666666614	-0.0289336924620972\\
0.000270833333333331	-0.0236175509297717\\
0.000375000000000014	-0.0099398164461062\\
0.000416666666666621	-0.0049181582154173\\
0.00043749999999998	-0.00260603487341049\\
0.000458333333333338	-0.000454470688504793\\
0.000479166666666697	0.00151867461190636\\
0.000499999999999945	0.00329835206993767\\
0.000520833333333304	0.00487247679202218\\
0.000541666666666663	0.00623203519655136\\
0.000562500000000021	0.00737114513983217\\
0.00058333333333338	0.00828706860093087\\
0.000604166666666628	0.00898017735982271\\
0.000624999999999987	0.00945387284344712\\
0.000645833333333345	0.00971446202866277\\
0.000666666666666704	0.00977099196810116\\
0.000687499999999952	0.00963504613361477\\
0.000708333333333311	0.00932050634240345\\
0.00072916666666667	0.00884328453400929\\
0.000750000000000028	0.00822102909442923\\
0.000770833333333387	0.00747281077017503\\
0.000791666666666635	0.00661879347518046\\
0.000833333333333353	0.00467744641894319\\
0.000937500000000036	-0.000554984574026918\\
0.000958333333333283	-0.00150950326827604\\
0.000979166666666642	-0.00239392632536617\\
0.001	-0.00319472827441547\\
0.00102083333333336	-0.00390040190991692\\
0.00104166666666672	-0.00450158158078551\\
0.00106249999999997	-0.0049911280907049\\
0.00108333333333333	-0.00536417460743954\\
0.00110416666666668	-0.00561813364122188\\
0.00112500000000004	-0.00575266580402922\\
0.00114583333333329	-0.00576961169154977\\
0.00116666666666665	-0.00567288882739969\\
0.00118750000000001	-0.00546835616486996\\
0.00120833333333337	-0.00516364914617817\\
0.00122916666666661	-0.00476798876485895\\
0.00124999999999997	-0.00429196845667068\\
0.00127083333333333	-0.00374732295256697\\
0.00129166666666669	-0.00314668345950841\\
0.0013333333333333	-0.00183090131978936\\
0.00141666666666662	0.000877137131158112\\
0.00143749999999998	0.00149429154376235\\
0.00145833333333334	0.00206442377572313\\
0.0014791666666667	0.00257788530046443\\
0.00149999999999995	0.0030263626579945\\
0.0015208333333333	0.00340299608403061\\
0.00154166666666666	0.00370246887363823\\
0.00156250000000002	0.00392106651788271\\
0.00158333333333338	0.00405670520761037\\
0.00160416666666663	0.00410892985197497\\
0.00162499999999999	0.00407888230089704\\
0.00164583333333335	0.00396924098084706\\
0.00166666666666671	0.00378413364320329\\
0.00168749999999995	0.00352902537568234\\
0.00170833333333331	0.00321058443246536\\
0.00172916666666667	0.00283652879109197\\
0.00175000000000003	0.00241545663842524\\
0.00179166666666664	0.00146995464960564\\
0.00189583333333332	-0.00102852684019938\\
0.00191666666666668	-0.00147314954171185\\
0.00193750000000004	-0.00187913323343325\\
0.00195833333333328	-0.00223938409131019\\
0.00197916666666664	-0.00254783933014335\\
0.002	-0.00279956172779527\\
0.00202083333333336	-0.00299081106446764\\
0.00204166666666672	-0.00311909156929424\\
0.00206249999999997	-0.00318317492734366\\
0.00208333333333333	-0.00318309886183787\\
};
\addplot [color=mycolor2, forget plot]
  table[row sep=crcr]{%
-0.00208333333333333	0.00318309886183791\\
-0.0020625	0.00318317492734361\\
-0.00204166666666666	0.00311909156929429\\
-0.00202083333333333	0.00299081106446766\\
-0.002	0.00279956172779523\\
-0.00197916666666666	0.00254783933014338\\
-0.00195833333333333	0.00223938409131023\\
-0.0019375	0.00187913323343322\\
-0.00191666666666666	0.00147314954171188\\
-0.00189583333333333	0.00102852684019933\\
-0.00185416666666666	5.61774650087374e-05\\
-0.00177083333333333	-0.00195666421959916\\
-0.00175	-0.00241545663842527\\
-0.00172916666666666	-0.00283652879109197\\
-0.00170833333333333	-0.00321058443246536\\
-0.0016875	-0.00352902537568231\\
-0.00166666666666666	-0.00378413364320329\\
-0.00164583333333333	-0.00396924098084705\\
-0.001625	-0.00407888230089703\\
-0.00160416666666666	-0.00410892985197494\\
-0.00158333333333333	-0.0040567052076104\\
-0.0015625	-0.00392106651788268\\
-0.00154166666666666	-0.00370246887363825\\
-0.00152083333333333	-0.00340299608403059\\
-0.0015	-0.0030263626579945\\
-0.00147916666666666	-0.00257788530046447\\
-0.00145833333333333	-0.0020644237757231\\
-0.0014375	-0.00149429154376232\\
-0.00139583333333333	-0.000223797745729983\\
-0.00129166666666667	0.00314668345950841\\
-0.00127083333333333	0.00374732295256701\\
-0.00125	0.0042919684566707\\
-0.00122916666666667	0.00476798876485895\\
-0.00120833333333333	0.00516364914617822\\
-0.0011875	0.00546835616486993\\
-0.00116666666666667	0.00567288882739972\\
-0.00114583333333333	0.00576961169154982\\
-0.001125	0.00575266580402919\\
-0.00110416666666667	0.00561813364122191\\
-0.00108333333333333	0.00536417460743957\\
-0.0010625	0.00499112809070489\\
-0.00104166666666667	0.00450158158078553\\
-0.00102083333333333	0.0039004019099169\\
-0.001	0.00319472827441542\\
-0.00097916666666667	0.00239392632536613\\
-0.000958333333333332	0.00150950326827601\\
-0.00091666666666667	-0.000454246434654471\\
-0.000812500000000001	-0.00567989546351649\\
-0.00079166666666667	-0.00661879347518047\\
-0.000770833333333332	-0.00747281077017498\\
-0.000750000000000001	-0.00822102909442928\\
-0.00072916666666667	-0.00884328453400932\\
-0.000708333333333332	-0.00932050634240345\\
-0.000687500000000001	-0.00963504613361472\\
-0.00066666666666667	-0.00977099196810118\\
-0.000645833333333332	-0.00971446202866279\\
-0.000625000000000001	-0.00945387284344715\\
-0.00060416666666667	-0.00898017735982275\\
-0.000583333333333332	-0.00828706860093084\\
-0.0005625	-0.0073711451398322\\
-0.000541666666666669	-0.0062320351965514\\
-0.000520833333333331	-0.00487247679202217\\
-0.0005	-0.0032983520699377\\
-0.000479166666666669	-0.0015186746119064\\
-0.000458333333333331	0.000454470688504793\\
-0.0004375	0.0026060348734105\\
-0.000395833333333331	0.0073703677412962\\
-0.000354166666666669	0.0126015609750331\\
-0.000229166666666669	0.0289336924620972\\
-0.0001875	0.0338042628617465\\
-0.000166666666666669	0.0360019245108092\\
-0.000145833333333331	0.0380065351526726\\
-0.000125	0.0397946281998788\\
-0.000104166666666669	0.0413451468692531\\
-8.33333333333311e-05	0.0426397415309974\\
-6.25000000000001e-05	0.0436630295501959\\
-4.1666666666669e-05	0.0444028136048185\\
-2.0833333333331e-05	0.0448502551012078\\
0	0.045\\
2.0833333333331e-05	0.0448502551012078\\
4.1666666666669e-05	0.0444028136048185\\
6.25000000000001e-05	0.0436630295501959\\
8.33333333333311e-05	0.0426397415309974\\
0.000104166666666669	0.0413451468692531\\
0.000125	0.0397946281998788\\
0.000145833333333331	0.0380065351526726\\
0.000166666666666669	0.0360019245108092\\
0.0001875	0.0338042628617465\\
0.000229166666666669	0.0289336924620972\\
0.000270833333333331	0.0236175509297717\\
0.000375	0.00993981644610623\\
0.000416666666666669	0.00491815821541732\\
0.0004375	0.0026060348734105\\
0.000458333333333331	0.000454470688504793\\
0.000479166666666669	-0.0015186746119064\\
0.0005	-0.0032983520699377\\
0.000520833333333331	-0.00487247679202217\\
0.000541666666666669	-0.0062320351965514\\
0.0005625	-0.0073711451398322\\
0.000583333333333332	-0.00828706860093084\\
0.00060416666666667	-0.00898017735982275\\
0.000625000000000001	-0.00945387284344715\\
0.000645833333333332	-0.00971446202866279\\
0.00066666666666667	-0.00977099196810118\\
0.000687500000000001	-0.00963504613361472\\
0.000708333333333332	-0.00932050634240345\\
0.00072916666666667	-0.00884328453400932\\
0.000750000000000001	-0.00822102909442928\\
0.000770833333333332	-0.00747281077017498\\
0.00079166666666667	-0.00661879347518047\\
0.000833333333333332	-0.0046774464189432\\
0.000937500000000001	0.000554984574026959\\
0.000958333333333332	0.00150950326827601\\
0.00097916666666667	0.00239392632536613\\
0.001	0.00319472827441542\\
0.00102083333333333	0.0039004019099169\\
0.00104166666666667	0.00450158158078553\\
0.0010625	0.00499112809070489\\
0.00108333333333333	0.00536417460743957\\
0.00110416666666667	0.00561813364122191\\
0.001125	0.00575266580402919\\
0.00114583333333333	0.00576961169154982\\
0.00116666666666667	0.00567288882739972\\
0.0011875	0.00546835616486993\\
0.00120833333333333	0.00516364914617822\\
0.00122916666666667	0.00476798876485895\\
0.00125	0.0042919684566707\\
0.00127083333333333	0.00374732295256701\\
0.00129166666666667	0.00314668345950841\\
0.00133333333333333	0.00183090131978936\\
0.00141666666666667	-0.000877137131158105\\
0.0014375	-0.00149429154376232\\
0.00145833333333333	-0.0020644237757231\\
0.00147916666666666	-0.00257788530046447\\
0.0015	-0.0030263626579945\\
0.00152083333333333	-0.00340299608403059\\
0.00154166666666666	-0.00370246887363825\\
0.0015625	-0.00392106651788268\\
0.00158333333333333	-0.0040567052076104\\
0.00160416666666666	-0.00410892985197494\\
0.001625	-0.00407888230089703\\
0.00164583333333333	-0.00396924098084705\\
0.00166666666666666	-0.00378413364320329\\
0.0016875	-0.00352902537568231\\
0.00170833333333333	-0.00321058443246536\\
0.00172916666666666	-0.00283652879109197\\
0.00175	-0.00241545663842527\\
0.00179166666666666	-0.00146995464960569\\
0.00189583333333333	0.00102852684019933\\
0.00191666666666666	0.00147314954171188\\
0.0019375	0.00187913323343322\\
0.00195833333333333	0.00223938409131023\\
0.00197916666666666	0.00254783933014338\\
0.002	0.00279956172779523\\
0.00202083333333333	0.00299081106446766\\
0.00204166666666666	0.00311909156929429\\
0.0020625	0.00318317492734361\\
0.00208333333333333	0.00318309886183791\\
};
\end{axis}
\end{tikzpicture}%

\paragraph{}

%%%% This file was created by matlab2tikz.
%
%The latest updates can be retrieved from
%  http://www.mathworks.com/matlabcentral/fileexchange/22022-matlab2tikz-matlab2tikz
%where you can also make suggestions and rate matlab2tikz.
%
\definecolor{mycolor1}{rgb}{0.00000,0.44700,0.74100}%
\definecolor{mycolor2}{rgb}{0.85000,0.32500,0.09800}%
%
\begin{tikzpicture}

\begin{axis}[%
width=4.521in,
height=3.548in,
at={(0.758in,0.499in)},
scale only axis,
unbounded coords=jump,
xmin=-25000,
xmax=25000,
ymode=log,
ymin=1e-08,
ymax=100,
yminorticks=true,
axis background/.style={fill=white}
]
\addplot [color=mycolor1, forget plot]
  table[row sep=crcr]{%
-24000	1.00284134051483\\
-9707.23965327799	1.00784668917116\\
-9443.53558784031	1.00634391250713\\
-9086.07007691369	1.00219220929594\\
-8664.1435722134	1.00853005589454\\
-8294.95788060066	1.0023857812845\\
-7878.89146624344	1.011376305327\\
-7509.70577463069	1.00240781529685\\
-7081.91917958735	1.01506004269969\\
-6730.31375900378	1.00217791442368\\
-6267.36662190209	1.01974901170664\\
-5968.50201440605	1.00096663918285\\
-5821.99975582957	1.01714225808576\\
-5611.03650347943	1.04099858046966\\
-5487.97460627518	1.03356722618652\\
-5177.38981809303	1.0013161485488\\
-5077.76828226102	1.01832174822303\\
-4802.34403613722	1.08121372328696\\
-4720.30277133439	1.07315522551496\\
-4626.54132584544	1.04486686922427\\
-4468.31888658283	0.996003132367466\\
-4409.71798315224	0.997418650516003\\
-4356.97717006471	1.0152401761057\\
-4304.23635697717	1.05120091934154\\
-4245.63545354658	1.11272091443234\\
-4169.4542790868	1.22121006145484\\
-4040.53229153949	1.43112405022432\\
-3981.9313881089	1.50941537509066\\
-3929.19057502137	1.55563075140529\\
-3882.30985227689	1.57034780457899\\
-3835.42912953241	1.55547668029589\\
-3794.408497131	1.515929511999\\
-3753.38786472958	1.45093497498028\\
-3718.22732267122	1.37546094305573\\
-3683.06678061287	1.28298625797261\\
-3647.90623855451	1.17540319674613\\
-3618.60578683921	1.07600980623875\\
-3589.30533512392	0.969613088671864\\
-3560.00488340862	0.858179393470407\\
-3536.56452203638	0.766835138670634\\
-3513.12416066415	0.674776579192027\\
-3489.68379929191	0.583176467998959\\
-3472.10352826273	0.51548717129855\\
-3454.52325723355	0.449213427234115\\
-3436.94298620437	0.384840462855972\\
-3419.36271517519	0.322840402176546\\
-3407.64253448907	0.283050703579549\\
-3395.92235380295	0.244647760125015\\
-3384.20217311684	0.20775640496713\\
-3372.48199243072	0.172496231612077\\
-3360.7618117446	0.138981138678295\\
-3354.90172140154	0.122911928192072\\
-3349.04163105848	0.107318904642378\\
-3343.18154071542	0.0922145195238283\\
-3337.32145037236	0.0776107944725708\\
-3331.4613600293	0.0635193108314036\\
-3325.60126968624	0.049951199770272\\
-3319.74117934318	0.0369171329732436\\
-3313.88108900012	0.0244273139028829\\
-3308.02099865706	0.0124914696441364\\
-3302.160908314	0.00111884333889737\\
-3296.30081797094	-0.00968181278675526\\
-3290.44072762788	-0.019902243794301\\
-3284.58063728482	-0.0295346978027968\\
-3278.72054694176	-0.0385719311511417\\
-3272.8604565987	-0.0470072129547623\\
-3267.00036625564	-0.0548343290500605\\
-3261.14027591259	-0.0620475853296203\\
-3255.28018556953	-0.068641810471287\\
-3249.42009522647	-0.0746123580559206\\
-3243.56000488341	-0.0799551080805467\\
-3237.69991454035	-0.0846664678712921\\
-3231.83982419729	-0.0887433723951299\\
-3220.11964351117	-0.0949841914586418\\
-3208.39946282505	-0.0986635725869213\\
-3196.67928213893	-0.0997758869663738\\
-3184.95910145281	-0.0983237780917271\\
-3173.2389207667	-0.0943180352165872\\
-3161.51874008058	-0.0877774331808311\\
-3155.65864973752	-0.083564563990025\\
-3149.79855939446	-0.0787285401467724\\
-3143.9384690514	-0.0732738638425516\\
-3138.07837870834	-0.0672054949415363\\
-3132.21828836528	-0.0605288424918328\\
-3126.35819802222	-0.0532497558352813\\
-3120.49810767916	-0.0453745153341191\\
-3114.6380173361	-0.036909822733815\\
-3108.77792699304	-0.0278627911743993\\
-3102.91783664998	-0.0182409348719961\\
-3097.05774630692	-0.00805215848653547\\
-3091.19765596386	0.00269525380486109\\
-3085.3375656208	0.0139926495122191\\
-3079.47747527774	0.0258310192938315\\
-3073.61738493469	0.0382010086571355\\
-3067.75729459163	0.051092929879601\\
-3061.89720424857	0.0644967741307062\\
-3056.03711390551	0.0784022237753743\\
-3050.17702356245	0.0927986648406104\\
-3044.31693321939	0.107675199625698\\
-3038.45684287633	0.123020659436337\\
-3032.59675253327	0.138823617423142\\
-3020.87657184715	0.171755107382671\\
-3009.15639116103	0.206373584339024\\
-2997.43621047491	0.242579665254818\\
-2985.71602978879	0.280271087478488\\
-2973.99584910267	0.319343125452807\\
-2956.41557807349	0.380305805263977\\
-2938.83530704432	0.443767362927457\\
-2921.25503601514	0.509355604286703\\
-2897.8146746429	0.599467263808462\\
-2874.37431327066	0.691807769167712\\
-2850.93395189842	0.785502690256061\\
-2821.63350018313	0.903228572744357\\
-2792.33304846783	1.02013891276253\\
-2757.17250640948	1.15723788805243\\
-2722.01196435112	1.2885943760245\\
-2686.85142229276	1.41205386257115\\
-2645.83078989134	1.54359360095767\\
-2604.81015748993	1.65924664745704\\
-2557.92943474545	1.76928374103485\\
-2511.04871200098	1.85336415918948\\
-2464.1679892565	1.90968229154065\\
-2411.42717616897	1.93838352068749\\
-2358.68636308143	1.92960648867852\\
-2311.80564033696	1.89039774211996\\
-2264.92491759248	1.8223753875837\\
-2218.044194848	1.72697563229736\\
-2171.16347210353	1.6064104532494\\
-2130.14283970211	1.48264976897469\\
-2089.1222073007	1.34453613945666\\
-2053.96166524234	1.21706267650456\\
-2018.80112318398	1.08353685194144\\
-1983.64058112563	0.946409026532242\\
-1954.34012941033	0.831279261727773\\
-1925.03967769503	0.71706815215248\\
-1901.59931632279	0.62745598313109\\
-1878.15895495056	0.54031517120234\\
-1854.71859357832	0.456479003249757\\
-1837.13832254914	0.396262514792276\\
-1819.55805151996	0.338699265508007\\
-1801.97778049078	0.284112704628335\\
-1784.3975094616	0.232812600360447\\
-1772.67732877548	0.200584104372002\\
-1760.95714808936	0.17002973578995\\
-1749.23696740325	0.141228143238201\\
-1737.51678671713	0.114253773031927\\
-1731.65669637407	0.101473846214038\\
-1725.79660603101	0.089176604251686\\
-1719.93651568795	0.0773700555255055\\
-1714.07642534489	0.066061898860856\\
-1708.21633500183	0.0552595166417749\\
-1702.35624465877	0.0449699681951868\\
-1696.49615431571	0.0351999834814129\\
-1690.63606397265	0.0259559570753332\\
-1684.77597362959	0.0172439424659039\\
-1678.91588328653	0.00906964666595916\\
-1673.05579294348	0.00143842514990056\\
-1667.19570260042	-0.0056447228804401\\
-1661.33561225736	-0.0121751588931026\\
-1655.4755219143	-0.0181486090815703\\
-1649.61543157124	-0.0235611680734111\\
-1643.75534122818	-0.0284093022881394\\
-1637.89525088512	-0.0326898529396692\\
-1632.03516054206	-0.03640003867962\\
-1626.175070199	-0.0395374578802563\\
-1620.31497985594	-0.0421000905512338\\
-1614.45488951288	-0.0440862998926634\\
-1608.59479916982	-0.0454948334797297\\
-1602.73470882676	-0.0463248240782844\\
-1596.8746184837	-0.0465757900940885\\
-1591.01452814064	-0.0462476356531751\\
-1585.15443779758	-0.0453406503121169\\
-1579.29434745452	-0.0438555084071351\\
-1573.43425711146	-0.0417932680339995\\
-1567.57416676841	-0.0391553696685281\\
-1561.71407642535	-0.0359436344278396\\
-1555.85398608229	-0.0321602619750341\\
-1549.99389573923	-0.0278078280714245\\
-1544.13380539617	-0.022889281782841\\
-1538.27371505311	-0.0174079423414557\\
-1532.41362471005	-0.0113674956705275\\
-1526.55353436699	-0.00477199057791364\\
-1520.69344402393	0.00237416537717297\\
-1514.83335368087	0.0100662103354489\\
-1508.97326333781	0.018299032901186\\
-1503.11317299475	0.0270671772101309\\
-1497.25308265169	0.036364848295125\\
-1491.39299230863	0.0461859177382354\\
-1485.53290196557	0.0565239296041517\\
-1479.67281162251	0.0673721066447664\\
-1473.81272127945	0.0787233567635679\\
-1467.95263093639	0.0905702797341611\\
-1462.09254059333	0.10290517415985\\
-1450.37235990721	0.129006609313262\\
-1438.6521792211	0.156960306392116\\
-1426.93199853498	0.186694573222286\\
-1415.21181784886	0.218133640395567\\
-1403.49163716274	0.251197931015164\\
-1385.91136613356	0.30365897324134\\
-1368.33109510438	0.359295171602081\\
-1350.7508240752	0.417798500253555\\
-1333.17055304603	0.478849030566786\\
-1309.73019167379	0.563641378644391\\
-1286.28983030155	0.651575047520894\\
-1262.84946892931	0.741835985968858\\
-1233.54901721402	0.856688017009298\\
-1204.24856549872	0.97230985832773\\
-1169.08802344036	1.1098882741476\\
-1133.927481382	1.24376503761306\\
-1098.76693932365	1.37152395594126\\
-1057.74630692223	1.50989246662556\\
-1016.72567452081	1.63377494418881\\
-969.84495177634	1.75417409627901\\
-922.964229031862	1.84885086771245\\
-876.083506287388	1.91531165235299\\
-829.202783542914	1.95185225640328\\
-782.322060798437	1.95754641903256\\
-735.441338053963	1.93223937442644\\
-688.560615309485	1.87655128693148\\
-641.679892565011	1.79189033012841\\
-594.799169820533	1.68047018598328\\
-553.778537419119	1.56338990081338\\
-512.757905017701	1.43067408083479\\
-471.737272616287	1.28532719492052\\
-436.576730557928	1.15328171802749\\
-401.416188499574	1.01685311452385\\
-372.115736784275	0.90164800901762\\
-342.815285068977	0.786705924625891\\
-313.514833353682	0.673604005400197\\
-290.074471981443	0.58551619082351\\
-266.634110609204	0.500433436681232\\
-243.193749236969	0.419150848011362\\
-225.61347820779	0.361148050051534\\
-208.033207178611	0.306031303990695\\
-190.452936149431	0.254101786166131\\
-178.732755463312	0.221394895211982\\
-167.012574777196	0.190312965254502\\
-155.292394091077	0.160933171753497\\
-143.572213404957	0.133328829423154\\
-131.852032718838	0.107569144067225\\
-125.991942375778	0.0954014925732736\\
-120.131852032719	0.0837189760648634\\
-114.271761689659	0.0725289864276536\\
-108.411671346599	0.0618386167720763\\
-102.551581003539	0.0516546552306847\\
-96.6914906604798	0.0419835790077327\\
-90.83140031742	0.0328315486793118\\
-84.9713099743603	0.0242044027596892\\
-79.1112196313043	0.0161076525315924\\
-73.2511292882446	0.00854647715867775\\
-67.3910389451848	0.00152571907700035\\
-61.5309486021251	-0.00495012032126534\\
-55.6708582590654	-0.0108768847103271\\
-49.8107679160057	-0.0162507665513866\\
-43.950677572946	-0.021068310507595\\
-38.0905872298863	-0.0253264165483282\\
-32.2304968868266	-0.0290223427374395\\
-26.3704065437669	-0.032153707704045\\
-20.5103162007072	-0.0347184927898288\\
-14.6502258576475	-0.0367150438710176\\
-8.79013551458775	-0.0381420728534124\\
-2.93004517152804	-0.0389986588368639\\
2.93004517152804	-0.039284248946786\\
8.79013551458775	-0.0389986588368639\\
14.6502258576475	-0.0381420728534124\\
20.5103162007072	-0.0367150438710176\\
26.3704065437669	-0.0347184927898288\\
32.2304968868266	-0.032153707704045\\
38.0905872298863	-0.0290223427374395\\
43.950677572946	-0.0253264165483282\\
49.8107679160057	-0.021068310507595\\
55.6708582590654	-0.0162507665513866\\
61.5309486021251	-0.0108768847103271\\
67.3910389451848	-0.00495012032126534\\
73.2511292882446	0.00152571907700035\\
79.1112196313043	0.00854647715867775\\
84.9713099743603	0.0161076525315924\\
90.83140031742	0.0242044027596892\\
96.6914906604798	0.0328315486793118\\
102.551581003539	0.0419835790077327\\
108.411671346599	0.0516546552306847\\
114.271761689659	0.0618386167720763\\
120.131852032719	0.0725289864276536\\
125.991942375778	0.0837189760648634\\
131.852032718838	0.0954014925732736\\
137.712123061898	0.107569144067225\\
149.432303748017	0.133328829423154\\
161.152484434137	0.160933171753497\\
172.872665120252	0.190312965254502\\
184.592845806372	0.221394895211982\\
196.313026492491	0.254101786166131\\
213.89329752167	0.306031303990695\\
231.473568550849	0.361148050051534\\
249.053839580029	0.419150848011362\\
272.494200952264	0.500433436681232\\
295.934562324503	0.58551619082351\\
319.374923696741	0.673604005400197\\
348.675375412036	0.786705924625891\\
377.975827127335	0.90164800901762\\
407.276278842633	1.01685311452385\\
442.436820900988	1.15328171802749\\
477.597362959346	1.28532719492052\\
518.617995360761	1.43067408083479\\
559.638627762179	1.56338990081338\\
600.659260163593	1.68047018598328\\
647.539982908071	1.79189033012841\\
694.420705652545	1.87655128693148\\
741.301428397022	1.93223937442644\\
788.182151141496	1.95754641903256\\
835.062873885974	1.95185225640328\\
881.943596630448	1.91531165235299\\
928.824319374922	1.84885086771245\\
975.7050421194	1.75417409627901\\
1022.58576486387	1.63377494418881\\
1063.60639726529	1.50989246662556\\
1104.62702966671	1.37152395594126\\
1145.64766206812	1.22181408965479\\
1180.80820412648	1.08714779382094\\
1215.96874618484	0.949199155478187\\
1245.26919790014	0.833605562337403\\
1274.56964961543	0.719097343665334\\
1298.01001098767	0.62934150291742\\
1321.45037235991	0.54211722323692\\
1344.89073373214	0.458236000732144\\
1362.47100476132	0.397998349569571\\
1380.0512757905	0.340415724088584\\
1397.63154681968	0.285804339612201\\
1415.21181784886	0.234467747490499\\
1426.93199853498	0.202205894311817\\
1438.6521792211	0.171609511729068\\
1450.37235990721	0.142756307220438\\
1462.09254059333	0.115720044667061\\
1467.95263093639	0.10290517415985\\
1473.81272127945	0.0905702797341611\\
1479.67281162251	0.0787233567635679\\
1485.53290196557	0.0673721066447664\\
1491.39299230863	0.0565239296041517\\
1497.25308265169	0.0461859177382354\\
1503.11317299475	0.036364848295125\\
1508.97326333781	0.0270671772101309\\
1514.83335368087	0.018299032901186\\
1520.69344402393	0.0100662103354489\\
1526.55353436699	0.00237416537717297\\
1532.41362471005	-0.00477199057791364\\
1538.27371505311	-0.0113674956705275\\
1544.13380539617	-0.0174079423414557\\
1549.99389573923	-0.022889281782841\\
1555.85398608229	-0.0278078280714245\\
1561.71407642535	-0.0321602619750341\\
1567.57416676841	-0.0359436344278396\\
1573.43425711146	-0.0391553696685281\\
1579.29434745452	-0.0417932680339995\\
1585.15443779758	-0.0438555084071351\\
1591.01452814064	-0.0453406503121169\\
1596.8746184837	-0.0462476356531751\\
1602.73470882676	-0.0465757900940885\\
1608.59479916982	-0.0463248240782844\\
1614.45488951288	-0.0454948334797297\\
1620.31497985594	-0.0440862998926634\\
1626.175070199	-0.0421000905512338\\
1632.03516054206	-0.0395374578802563\\
1637.89525088512	-0.03640003867962\\
1643.75534122818	-0.0326898529396692\\
1649.61543157124	-0.0284093022881394\\
1655.4755219143	-0.0235611680734111\\
1661.33561225736	-0.0181486090815703\\
1667.19570260042	-0.0121751588931026\\
1673.05579294348	-0.0056447228804401\\
1678.91588328653	0.00143842514990056\\
1684.77597362959	0.00906964666595916\\
1690.63606397265	0.0172439424659039\\
1696.49615431571	0.0259559570753332\\
1702.35624465877	0.0351999834814129\\
1708.21633500183	0.0449699681951868\\
1714.07642534489	0.0552595166417749\\
1719.93651568795	0.066061898860856\\
1725.79660603101	0.0773700555255055\\
1731.65669637407	0.089176604251686\\
1737.51678671713	0.101473846214038\\
1743.37687706019	0.114253773031927\\
1755.09705774631	0.141228143238201\\
1766.81723843242	0.17002973578995\\
1778.53741911854	0.200584104372002\\
1790.25759980466	0.232812600360447\\
1801.97778049078	0.26663265139941\\
1819.55805151996	0.320157400015862\\
1837.13832254914	0.376763733173268\\
1854.71859357832	0.43613249110897\\
1878.15895495056	0.519014062109301\\
1901.59931632279	0.605406347429255\\
1925.03967769503	0.694479336996092\\
1954.34012941033	0.808311718652108\\
1983.64058112563	0.923386843517411\\
2012.94103284092	1.03810107619541\\
2048.10157489928	1.17310981081574\\
2083.26211695764	1.30285987535807\\
2124.28274935905	1.44452721005051\\
2165.30338176047	1.57267556300461\\
2206.32401416188	1.68456770292539\\
2253.20473690636	1.78970764655892\\
2300.08545965084	1.86820488191576\\
2346.96618239531	1.91834094544438\\
2399.70699548285	1.93963704193348\\
2452.44780857038	1.92338067716478\\
2499.32853131486	1.87781407428582\\
2546.20925405933	1.8039664386101\\
2593.08997680381	1.70341094936495\\
2639.97069954828	1.57835899148366\\
2680.9913319497	1.4511168886958\\
2722.01196435112	1.30977379476686\\
2757.17250640948	1.17959075230625\\
2792.33304846783	1.04328853805076\\
2821.63350018313	0.926725599194983\\
2850.93395189842	0.809037455065404\\
2874.37431327066	0.715138286231451\\
2897.8146746429	0.622377867674489\\
2921.25503601514	0.531626214573054\\
2938.83530704432	0.465412288772645\\
2956.41557807349	0.401200303594941\\
2973.99584910267	0.339363576894535\\
2985.71602978879	0.299641159725105\\
2997.43621047491	0.261246201189658\\
3009.15639116103	0.224284503124075\\
3020.87657184715	0.188859611535734\\
3032.59675253327	0.155072401510128\\
3038.45684287633	0.138823617423142\\
3044.31693321939	0.123020659436337\\
3050.17702356245	0.107675199625698\\
3056.03711390551	0.0927986648406104\\
3061.89720424857	0.0784022237753743\\
3067.75729459163	0.0644967741307062\\
3073.61738493469	0.051092929879601\\
3079.47747527774	0.0382010086571355\\
3085.3375656208	0.0258310192938315\\
3091.19765596386	0.0139926495122191\\
3097.05774630692	0.00269525380486109\\
3102.91783664998	-0.00805215848653547\\
3108.77792699304	-0.0182409348719961\\
3114.6380173361	-0.0278627911743993\\
3120.49810767916	-0.036909822733815\\
3126.35819802222	-0.0453745153341191\\
3132.21828836528	-0.0532497558352813\\
3138.07837870834	-0.0605288424918328\\
3143.9384690514	-0.0672054949415363\\
3149.79855939446	-0.0732738638425516\\
3155.65864973752	-0.0787285401467724\\
3161.51874008058	-0.083564563990025\\
3167.37883042364	-0.0877774331808311\\
3179.09901110976	-0.0943180352165872\\
3190.81919179587	-0.0983237780917271\\
3202.53937248199	-0.0997758869663738\\
3214.25955316811	-0.0986635725869213\\
3225.97973385423	-0.0949841914586418\\
3237.69991454035	-0.0887433723951299\\
3243.56000488341	-0.0846664678712921\\
3249.42009522647	-0.0799551080805467\\
3255.28018556953	-0.0746123580559206\\
3261.14027591259	-0.068641810471287\\
3267.00036625564	-0.0620475853296203\\
3272.8604565987	-0.0548343290500605\\
3278.72054694176	-0.0470072129547623\\
3284.58063728482	-0.0385719311511417\\
3290.44072762788	-0.0295346978027968\\
3296.30081797094	-0.019902243794301\\
3302.160908314	-0.00968181278675526\\
3308.02099865706	0.00111884333889737\\
3313.88108900012	0.0124914696441364\\
3319.74117934318	0.0244273139028829\\
3325.60126968624	0.0369171329732436\\
3331.4613600293	0.049951199770272\\
3337.32145037236	0.0635193108314036\\
3343.18154071542	0.0776107944725708\\
3349.04163105848	0.0922145195238283\\
3354.90172140154	0.107318904642378\\
3360.7618117446	0.122911928192072\\
3366.62190208766	0.138981138678295\\
3378.34208277378	0.172496231612077\\
3390.06226345989	0.20775640496713\\
3401.78244414601	0.244647760125015\\
3413.50262483213	0.283050703579549\\
3425.22280551825	0.322840402176546\\
3442.80307654743	0.384840462855972\\
3460.38334757661	0.449213427234115\\
3477.96361860579	0.51548717129855\\
3501.40397997803	0.605968936689599\\
3524.84434135026	0.697794319707401\\
3548.2847027225	0.789783546201827\\
3577.5851544378	0.903223892767596\\
3606.88560615309	1.01289063691337\\
3636.18605786839	1.11671087935326\\
3671.34659992675	1.23093964877382\\
3706.5071419851	1.33124592697357\\
3741.66768404346	1.41541503731414\\
3782.68831644488	1.4911806242377\\
3823.7089488463	1.5416340264825\\
3864.72958124771	1.56683176098259\\
3911.61030399219	1.56643960107551\\
3958.49102673666	1.53854026129323\\
4011.2318398242	1.4808441855122\\
4069.83274325479	1.3946357643897\\
4157.73409840068	1.2494405716888\\
4257.3556342327	1.10560367653108\\
4315.95653766329	1.04628045551816\\
4368.69735075082	1.01239104258024\\
4421.43816383836	0.996473557808924\\
4480.03906726895	0.996743029888078\\
4556.22024172873	1.01643403533274\\
4749.60322304969	1.07729493343014\\
4831.64448785252	1.08019849875199\\
4925.40593334147	1.06374336103988\\
5212.55036015138	0.999660690850909\\
5318.03198632646	1.00562490882278\\
5640.33695519473	1.04025027433413\\
5780.97912342815	1.02442129157638\\
5997.80246612135	1.00012251061262\\
6138.44463435478	1.0058972902668\\
6437.30924185081	1.02581498795457\\
6618.97204248565	1.01181235050572\\
6824.07520449274	0.999955278799365\\
7011.59809547064	1.00971356554002\\
7240.14161884996	1.01839349847716\\
7480.40532291539	1.00430903784712\\
7679.64839457942	1.00099296556962\\
8183.61616408253	1.00837668572661\\
8476.6206812355	1.00075409056754\\
9015.74899279697	1.00527477136276\\
9302.89341960689	1.00124068173093\\
9783.42082773776	1.00553615605415\\
10129.1661579783	1.00170145932106\\
10574.5330240508	1.00496281748533\\
10955.4388963497	1.00211668208415\\
11377.3654010499	1.00421894135086\\
11775.851544378	1.00233799705817\\
12186.0578683921	1.00350960004946\\
12596.2641924063	1.00252579591822\\
13006.4705164205	1.00270184819896\\
13416.6768404346	1.00268272959236\\
13826.8831644488	1.00205310528284\\
14242.949578806	1.0029115097727\\
14659.0159931632	1.00133875633935\\
15086.8025882066	1.00332738854949\\
15532.1694542791	1.00034476341141\\
16516.6646319131	1.00064416803259\\
17249.1759247955	1.00009481476089\\
19399.8290806983	1.00165341641159\\
19991.6982053473	1.00305185813636\\
24000	1.00283983599301\\
};
\addplot [color=mycolor2, forget plot]
  table[row sep=crcr]{%
-24000	-0.00284134051541254\\
-23976.5596386278	-0.00281733185141145\\
-23953.1192772555	-0.00274611709168084\\
-23929.6789158833	-0.00263010254264651\\
-23906.238554511	-0.00247320833745444\\
-23882.7981931388	-0.00228073607346066\\
-23859.3578317666	-0.00205918980113674\\
-23841.7775607374	-0.00187842412257593\\
-23824.1972897082	-0.00168894126064206\\
-23806.617018679	-0.00149434793701714\\
-23789.0367476499	-0.0012983483943787\\
-23771.4564766207	-0.00110467394012934\\
-23759.7362959346	-0.000978674225344446\\
-23748.0161152484	-0.000856416311212569\\
-23736.2959345623	-0.000738935829724721\\
-23724.5757538762	-0.000627228047288775\\
-23712.8555731901	-0.000522239440905372\\
-23701.135392504	-0.000424859686585813\\
-23689.4152118178	-0.000335914127787074\\
-23683.5551214748	-0.000294842768824362\\
-23677.6950311317	-0.000256156787581085\\
-23671.8349407887	-0.000219938245436648\\
-23665.9748504456	-0.000186263983670552\\
-23660.1147601026	-0.00015520546063389\\
-23654.2546697595	-0.000126828600321508\\
-23648.3945794164	-0.000101193652633783\\
-23642.5344890734	-7.83550656602712e-05\\
-23636.6743987303	-5.83613702124949e-05\\
-23630.8143083873	-4.12550768870749e-05\\
-23624.9542180442	-2.7072585857694e-05\\
-23619.0941277011	-1.58441095928798e-05\\
-23613.2340373581	-7.59360866299011e-06\\
-23607.373947015	-2.33874077388739e-06\\
-23601.513856672	-9.0823133571775e-08\\
-23595.6537663289	-8.54808233626689e-07\\
-23589.7936759858	-4.62927309620972e-06\\
-23583.9335856428	-1.14064220109323e-05\\
-23578.0734952997	-2.1172102755075e-05\\
-23572.2134049567	-3.39058362646826e-05\\
-23566.3533146136	-4.95808596941853e-05\\
-23560.4932242705	-6.81641827723735e-05\\
-23554.6331339275	-8.96166573390238e-05\\
-23548.7730435844	-0.000113893059913721\\
-23542.9129532414	-0.000140942187121089\\
-23537.0528628983	-0.000170706963778393\\
-23531.1927725552	-0.000203124563407944\\
-23525.3326822122	-0.000238126540927679\\
-23519.4725918691	-0.000275638977228373\\
-23513.6125015261	-0.000315582635360732\\
-23501.8923208399	-0.000402421095557273\\
-23490.1721401538	-0.000497908301027104\\
-23478.4519594677	-0.000601237361657966\\
-23466.7317787816	-0.000711534937569924\\
-23455.0115980955	-0.000827868621857204\\
-23443.2914174093	-0.000949254823742937\\
-23425.7111463802	-0.00113855262994018\\
-23408.130875351	-0.00133330200090315\\
-23390.5506043218	-0.00152979955553598\\
-23372.9703332926	-0.00172430801620584\\
-23355.3900622635	-0.00191312729735072\\
-23337.8097912343	-0.00209266490066311\\
-23314.369429862	-0.00231173858207379\\
-23290.9290684898	-0.00250086139780925\\
-23267.4887071176	-0.002653639369865\\
-23244.0483457453	-0.00276490487818962\\
-23220.6079843731	-0.00283089160029641\\
-23197.1676230009	-0.00284936218507815\\
-23173.7272616286	-0.00281968433563411\\
-23150.2869002564	-0.00274285271512935\\
-23126.8465388841	-0.0026214559182546\\
-23103.4061775119	-0.00245958960357174\\
-23079.9658161397	-0.00226271870021161\\
-23056.5254547674	-0.00203749332194539\\
-23038.9451837382	-0.00185451617663497\\
-23021.3649127091	-0.00166334005929498\\
-23003.7846416799	-0.00146760582281903\\
-22986.2043706507	-0.00127104229099939\\
-22968.6240996215	-0.00107739534806444\\
-22956.9039189354	-0.000951745131603513\\
-22945.1837382493	-0.000830100429676494\\
-22933.4635575632	-0.000713493527961728\\
-22921.7433768771	-0.000602914317891804\\
-22910.0231961909	-0.00049930191015217\\
-22898.3030155048	-0.000403536675847783\\
-22886.5828348187	-0.000316432782750981\\
-22880.7227444756	-0.000276363452316818\\
-22874.8626541326	-0.000238731289862733\\
-22869.0025637895	-0.000203616472647206\\
-22863.1424734465	-0.000171093858826122\\
-22857.2823831034	-0.000141232828365827\\
-22851.4222927603	-0.000114097135518203\\
-22845.5622024173	-8.97447731970047e-05\\
-22839.7021120742	-6.82278495082941e-05\\
-22833.8420217312	-4.95924767370388e-05\\
-22827.9819313881	-3.38786729934728e-05\\
-22822.121841045	-2.11202767477588e-05\\
-22816.261750702	-1.13448744295206e-05\\
-22810.4016603589	-4.57374124570659e-06\\
-22804.5415700159	-8.21795347178264e-07\\
-22798.6814796728	-9.75654389382338e-08\\
-22792.8213893298	-2.40317190797747e-06\\
-22786.9612989867	-7.73432150630204e-06\\
-22781.1012086436	-1.60803156061958e-05\\
-22775.2411183006	-2.74240720093062e-05\\
-22769.3810279575	-4.17421602660924e-05\\
-22763.5209376145	-5.90048504323481e-05\\
-22757.6608472714	-7.91761751625779e-05\\
-22751.8007569283	-0.000102214005008151\\
-22745.9406665853	-0.000128070136766963\\
-22740.0805762422	-0.000156690394694511\\
-22734.2204858992	-0.000188014744367536\\
-22728.3603955561	-0.000221977418955338\\
-22722.500305213	-0.000258507057649185\\
-22716.64021487	-0.000297526855936778\\
-22710.7801245269	-0.00033895472741485\\
-22699.0599438408	-0.000428680983837285\\
-22687.3397631547	-0.000526930340305221\\
-22675.6195824686	-0.000632874964801485\\
-22663.8994017824	-0.000745621665077606\\
-22652.1792210963	-0.000864219428830506\\
-22640.4590404102	-0.000987667455107538\\
-22622.878769381	-0.00117964492654475\\
-22605.2984983518	-0.00137653821552113\\
-22587.7182273227	-0.00157460737881394\\
-22570.1379562935	-0.00177008848521255\\
-22552.5576852643	-0.00195926521608474\\
-22534.9774142351	-0.00213853959297041\\
-22511.5370528629	-0.00235630774269005\\
-22488.0966914907	-0.00254307525048042\\
-22464.6563301184	-0.00269252401529647\\
-22441.2159687462	-0.00279959325577728\\
-22417.7756073739	-0.00286065131759589\\
-22394.3352460017	-0.0028736193259035\\
-22370.8948846295	-0.00283804247544993\\
-22347.4545232572	-0.00275510651585241\\
-22324.014161885	-0.00262759883751439\\
-22300.5738005128	-0.00245981543344755\\
-22277.1334391405	-0.00225741684033257\\
-22253.6930777683	-0.00202723788587094\\
-22236.1128067391	-0.0018410241368157\\
-22218.5325357099	-0.0016470921792747\\
-22200.9522646807	-0.00144913758640776\\
-22183.3719936516	-0.00125093471025438\\
-22165.7917226224	-0.00105626491351087\\
-22154.0715419363	-0.000930287966284631\\
-22142.3513612502	-0.000808606329948592\\
-22130.631180564	-0.000692254594489036\\
-22118.9109998779	-0.000582222661009546\\
-22107.1908191918	-0.000479447346117513\\
-22095.4706385057	-0.000384804432001888\\
-22089.6105481626	-0.00034078777638871\\
-22083.7504578196	-0.000299101229714281\\
-22077.8903674765	-0.000259833881528768\\
-22072.0302771334	-0.000223069718647253\\
-22066.1701867904	-0.000188887446328079\\
-22060.3100964473	-0.000157360320569037\\
-22054.4500061043	-0.000128555991917456\\
-22048.5899157612	-0.000102536361090831\\
-22042.7298254181	-7.93574467543084e-05\\
-22036.8697350751	-5.90692656961301e-05\\
-22031.009644732	-4.17157256969613e-05\\
-22025.149554389	-2.73345312889737e-05\\
-22019.2894640459	-1.59571026288878e-05\\
-22013.4293737028	-7.60850765080387e-06\\
-22007.5692833598	-2.30740764265196e-06\\
-22001.7091930167	-6.60163722716593e-08\\
-21995.8491026737	-8.90072845246499e-07\\
-21989.9890123306	-4.77882775783598e-06\\
-21984.1289219875	-1.1725043677168e-05\\
-21978.2688316445	-2.17150089512198e-05\\
-21972.4087413014	-3.47285653200621e-05\\
-21966.5486509584	-5.07391491756901e-05\\
-21960.6885606153	-6.97138463841963e-05\\
-21954.8284702722	-9.16134605594766e-05\\
-21948.9683799292	-0.000116392594641658\\
-21943.1082895861	-0.000143999745620105\\
-21937.2481992431	-0.000174377412194723\\
-21931.3881089	-0.000207462215150681\\
-21925.5280185569	-0.000243185030209084\\
-21919.6679282139	-0.000281471133043174\\
-21913.8078378708	-0.000322240356207789\\
-21902.0876571847	-0.000410881300080587\\
-21890.3674764986	-0.000508364338302483\\
-21878.6472958125	-0.000613870769674534\\
-21866.9271151264	-0.000726513609956042\\
-21855.2069344402	-0.000845345067356821\\
-21843.4867537541	-0.000969364536003562\\
-21825.9064827249	-0.00116282540728234\\
-21808.3262116958	-0.00136193264305908\\
-21790.7459406666	-0.00156290913142642\\
-21773.1656696374	-0.00176193957099966\\
-21755.5853986082	-0.00195524304155938\\
-21738.0051275791	-0.00213914497301313\\
-21714.5647662068	-0.00236371040728484\\
-21691.1244048346	-0.00255777979373768\\
-21667.6840434623	-0.00271478547671585\\
-21644.2436820901	-0.00282940627672859\\
-21620.8033207179	-0.00289774855537648\\
-21597.3629593456	-0.00291747924063\\
-21573.9225979734	-0.00288790627088031\\
-21550.4822366011	-0.00281000368688873\\
-21527.0418752289	-0.00268638046715559\\
-21503.6015138567	-0.00252119409978788\\
-21480.1611524844	-0.00232001175123109\\
-21456.7207911122	-0.0020896236641936\\
-21439.140520083	-0.00190231942998053\\
-21421.5602490538	-0.00170651903154536\\
-21403.9799780247	-0.00150595524508733\\
-21386.3997069955	-0.00130445474806615\\
-21368.8194359663	-0.00110586537972965\\
-21357.0992552802	-0.000976968564233929\\
-21345.3790745941	-0.000852152069549209\\
-21333.6588939079	-0.000732478899955249\\
-21321.9387132218	-0.00061896914622417\\
-21310.2185325357	-0.000512591326913035\\
-21298.4983518496	-0.000414254161034963\\
-21286.7781711635	-0.000324798841846169\\
-21280.9180808204	-0.000283644509092519\\
-21275.0579904774	-0.000244991877513645\\
-21269.1979001343	-0.000208924019328482\\
-21263.3378097912	-0.000175518559451495\\
-21257.4777194482	-0.000144847509567032\\
-21251.6176291051	-0.000116977114014739\\
-21245.7575387621	-9.19677078558858e-05\\
-21239.897448419	-6.98735873809977e-05\\
-21234.0373580759	-5.07428933736981e-05\\
-21228.1772677329	-3.46175073520181e-05\\
-21222.3171773898	-2.15329610357304e-05\\
-21216.4570870468	-1.1518359214173e-05\\
-21210.5969967037	-4.5963161948821e-06\\
-21204.7369063606	-7.82905968610599e-07\\
-21198.8768160176	-8.76261961876233e-08\\
-21193.0167256745	-2.51337610274945e-06\\
-21187.1566353315	-8.05644832862293e-06\\
-21181.2965449884	-1.6706534756168e-05\\
-21175.4364546453	-2.84467463083629e-05\\
-21169.5763643023	-4.32536466786409e-05\\
-21163.7162739592	-6.1097299924636e-05\\
-21157.8561836162	-8.19413318335964e-05\\
-21151.9960932731	-0.000105743004929365\\
-21146.13600293	-0.000132453306964239\\
-21140.275912587	-0.000162017052727458\\
-21134.4158222439	-0.000194372998928431\\
-21128.5557319009	-0.000229453971963535\\
-21122.6956415578	-0.000267187008242536\\
-21116.8355512147	-0.000307493506831298\\
-21110.9754608717	-0.000350289394060698\\
-21099.2552801856	-0.00044298674486833\\
-21087.5350994994	-0.000544503992459431\\
-21075.8149188133	-0.000653990920535655\\
-21064.0947381272	-0.000770529313278586\\
-21052.3745574411	-0.000893140683809489\\
-21040.654376755	-0.00102079451760563\\
-21023.0741057258	-0.001219371205951\\
-21005.4938346966	-0.0014231102353403\\
-20987.9135636674	-0.00162815076132841\\
-20970.3332926383	-0.001830603462985\\
-20952.7530216091	-0.00202662453111203\\
-20935.1727505799	-0.0022124888636091\\
-20911.7323892077	-0.00243843387551629\\
-20888.2920278354	-0.00263242710895824\\
-20864.8516664632	-0.00278789881637128\\
-20841.411305091	-0.00289957289827388\\
-20817.9709437187	-0.00296364698937434\\
-20794.5305823465	-0.00297792299951192\\
-20771.0902209742	-0.00294188363257645\\
-20747.649859602	-0.00285671222316417\\
-20724.2094982298	-0.00272525514233901\\
-20700.7691368575	-0.00255192796067947\\
-20677.3287754853	-0.00234256845891991\\
-20653.8884141131	-0.00210424137417143\\
-20636.3081430839	-0.00191130174259514\\
-20618.7278720547	-0.00171026079202968\\
-20601.1476010255	-0.00150495390425756\\
-20583.5673299963	-0.00129930204105443\\
-20565.9870589672	-0.00109723721700748\\
-20554.266878281	-0.000966435241031535\\
-20542.5466975949	-0.000840064734401158\\
-20530.8265169088	-0.000719204212635943\\
-20519.1063362227	-0.000604886375442926\\
-20507.3861555366	-0.000498089330591291\\
-20495.6659748504	-0.000399728272984745\\
-20489.8058845074	-0.000353978104927964\\
-20483.9457941643	-0.000310647689780881\\
-20478.0857038213	-0.000269830358505824\\
-20472.2256134782	-0.000231614157814808\\
-20466.3655231351	-0.00019608166175015\\
-20460.5054327921	-0.000163309794756426\\
-20454.645342449	-0.000133369666660804\\
-20448.785252106	-0.000106326419876124\\
-20442.9251617629	-8.22390891847683e-05\\
-20437.0650714198	-6.11604743852655e-05\\
-20431.2049810768	-4.31370260855577e-05\\
-20425.3448907337	-2.82087448914954e-05\\
-20419.4848003907	-1.64090941985982e-05\\
-20413.6247100476	-7.76492678911273e-06\\
-20407.7646197046	-2.29642538632653e-06\\
-20401.9045293615	-1.70573014601601e-08\\
-20396.0444390184	-9.33543273517639e-07\\
-20390.1843486754	-5.04584057474679e-06\\
-20384.3242583323	-1.23471404209158e-05\\
-20378.4641679893	-2.28238796987978e-05\\
-20372.6040776462	-3.64557669931793e-05\\
-20366.7439873031	-5.32158228612711e-05\\
-20360.8838969601	-7.30704342765375e-05\\
-20355.023806617	-9.59794231308124e-05\\
-20349.163716274	-0.000121896128662621\\
-20343.3036259309	-0.000150767503627476\\
-20337.4435355878	-0.000182534224030026\\
-20331.5834452448	-0.000217130812177925\\
-20325.7233549017	-0.000254485772812664\\
-20319.8632645587	-0.000294521742025921\\
-20314.0031742156	-0.000337155648650876\\
-20302.2829935295	-0.000429857506229384\\
-20290.5628128434	-0.000531819517057388\\
-20278.8426321572	-0.000642190833221445\\
-20267.1224514711	-0.000760048725137236\\
-20255.402270785	-0.000884406333322104\\
-20243.6820900989	-0.00101422096927117\\
-20226.1018190697	-0.00121678043811988\\
-20208.5215480405	-0.00142532986583139\\
-20190.9412770114	-0.00163592296076743\\
-20173.3610059822	-0.00184456966670262\\
-20155.780734953	-0.00204731205105907\\
-20138.2004639238	-0.00224029967761424\\
-20114.7601025516	-0.00247613493676936\\
-20091.3197411793	-0.00268016006111007\\
-20067.8793798071	-0.00284546356189473\\
-20044.4390184349	-0.00296643163928709\\
-20020.9986570626	-0.00303894019690302\\
-19997.5582956904	-0.00306049685036157\\
-19974.1179343182	-0.00303032804499619\\
-19950.6775729459	-0.00294940824394311\\
-19927.2372115737	-0.0028204300975864\\
-19903.7968502014	-0.00264771649768994\\
-19880.3564888292	-0.00243707738301771\\
-19856.916127457	-0.00219561603434716\\
-19839.3358564278	-0.00199916885448461\\
-19821.7555853986	-0.00179370166096843\\
-19804.1753143694	-0.00158313571932502\\
-19786.5950433403	-0.00137149493905905\\
-19769.0147723111	-0.00116282942486231\\
-19757.294591625	-0.00102735134846444\\
-19745.5744109388	-0.000896132343363573\\
-19733.8542302527	-0.000770294123527442\\
-19722.1340495666	-0.000650913919879023\\
-19710.4138688805	-0.000539015319244037\\
-19698.6936881944	-0.000435559548498835\\
-19686.9735075082	-0.00034143727816274\\
-19681.1134171652	-0.000298133362220343\\
-19675.2533268221	-0.000257461015274727\\
-19669.3932364791	-0.000219508446775727\\
-19663.533146136	-0.000184358150436336\\
-19657.6730557929	-0.000152086726849342\\
-19651.8129654499	-0.000122764718492462\\
-19645.9528751068	-9.64564574642837e-05\\
-19640.0927847638	-7.32199263042771e-05\\
-19634.2326944207	-5.31066321693841e-05\\
-19628.3726040776	-3.61614946600077e-05\\
-19622.5125137346	-2.24227475303034e-05\\
-19616.6524233915	-1.19218545007031e-05\\
-19610.7923330485	-4.6834393550504e-06\\
-19604.9322427054	-7.25230481056351e-07\\
-19599.0721523623	-5.8019972881849e-08\\
-19593.2120620193	-2.68563739496502e-06\\
-19587.3519716762	-8.60493826722153e-06\\
-19581.4918813332	-1.78058073018194e-05\\
-19575.6317909901	-3.02711763954224e-05\\
-19569.7717006471	-4.59770573428382e-05\\
-19563.911610304	-6.48925892133559e-05\\
-19558.0515199609	-8.69801002897976e-05\\
-19552.1914296179	-0.000112195184460174\\
-19546.3313392748	-0.000140486791887828\\
-19540.4712489318	-0.000171797333797657\\
-19534.6111585887	-0.000206062801140336\\
-19528.7510682456	-0.000243212896917138\\
-19522.8909779026	-0.000283171181858921\\
-19517.0308875595	-0.000325855233192107\\
-19511.1707972165	-0.000371176816142205\\
-19499.4506165303	-0.000469351693087314\\
-19487.7304358442	-0.000576880978879248\\
-19476.0102551581	-0.000692869827713654\\
-19464.290074472	-0.000816350876680697\\
-19452.5698937859	-0.000946292361867014\\
-19440.8497130997	-0.00108160678768538\\
-19423.2694420706	-0.00129216150219575\\
-19405.6891710414	-0.00150827166609827\\
-19388.1089000122	-0.00172585220064794\\
-19370.528628983	-0.00194078384828432\\
-19352.9483579539	-0.00214899153281789\\
-19335.3680869247	-0.00234652200028608\\
-19311.9277255524	-0.00258683256160068\\
-19288.4873641802	-0.0027933871035165\\
-19265.047002808	-0.00295918251593026\\
-19241.6066414357	-0.00307857974996451\\
-19218.1662800635	-0.00314749735773166\\
-19194.7259186912	-0.0031635527994453\\
-19171.285557319	-0.00312614663227855\\
-19147.8451959468	-0.00303648661415781\\
-19124.4048345745	-0.00289755077932135\\
-19100.9644732023	-0.0027139906008825\\
-19077.5241118301	-0.00249197738222417\\
-19054.0837504578	-0.00223899694348069\\
-19036.5034794286	-0.00203405071237181\\
-19018.9232083995	-0.00182038724457605\\
-19001.3429373703	-0.00160208759142087\\
-18983.7626663411	-0.00138332807492845\\
-18966.1823953119	-0.00116830094451714\\
-18954.4622146258	-0.00102906680304969\\
-18942.7420339397	-0.000894520485818125\\
-18931.0218532536	-0.000765814873823112\\
-18919.3016725674	-0.000644054721529025\\
-18907.5814918813	-0.000530287249178608\\
-18895.8613111952	-0.00042549321124583\\
-18890.0012208522	-0.000376747467148975\\
-18884.1411305091	-0.000330578517176891\\
-18878.281040166	-0.000287086632907327\\
-18872.420949823	-0.000246366475928381\\
-18866.5608594799	-0.000208506894207686\\
-18860.7007691369	-0.000173590730640519\\
-18854.8406787938	-0.000141694644190633\\
-18848.9805884507	-0.000112888944024474\\
-18843.1204981077	-8.72374369815938e-05\\
-18837.2604077646	-6.47972887199579e-05\\
-18831.4003174216	-4.56188988320399e-05\\
-18825.5402270785	-2.97457902099721e-05\\
-18819.6801367354	-1.72145128995256e-05\\
-18813.8200463924	-8.054562652008e-06\\
-18807.9599560493	-2.28831436094514e-06\\
-18802.0998657063	6.90294731979261e-08\\
-18796.2397753632	-9.90524872404033e-07\\
-18790.3796850201	-5.46774119205874e-06\\
-18784.5195946771	-1.33561474888933e-05\\
-18778.659504334	-2.46420454223196e-05\\
-18772.799413991	-3.93045350541097e-05\\
-18766.9393236479	-5.73155548495061e-05\\
-18761.0792333048	-7.86399368576844e-05\\
-18755.2191429618	-0.000103235476965918\\
-18749.3590526187	-0.000131053020082639\\
-18743.4989622757	-0.000162036560084215\\
-18737.6388719326	-0.00019612335432281\\
-18731.7787815895	-0.000233244052445928\\
-18725.9186912465	-0.000273322839292687\\
-18720.0586009034	-0.000316277591548625\\
-18714.1985105604	-0.000362020047838695\\
-18702.4783298743	-0.000461485448615252\\
-18690.7581491881	-0.000570897465252642\\
-18679.037968502	-0.000689349321815146\\
-18667.3177878159	-0.000815856617485192\\
-18655.5976071298	-0.000949365567354957\\
-18643.8774264437	-0.00108876184037811\\
-18626.2971554145	-0.00130633382351794\\
-18608.7168843853	-0.00153042270630186\\
-18591.1366133561	-0.00175679955082239\\
-18573.5563423269	-0.00198118421492547\\
-18555.9760712978	-0.00219932674509548\\
-18538.3958002686	-0.00240708834446833\\
-18514.9554388963	-0.00266116777054495\\
-18491.5150775241	-0.00288121078107255\\
-18468.0747161519	-0.00305975541951707\\
-18444.6343547796	-0.00319072521218933\\
-18421.1939934074	-0.00326963812941574\\
-18397.7536320352	-0.00329376217527827\\
-18374.3132706629	-0.00326221221261659\\
-18350.8729092907	-0.00317598460315991\\
-18327.4325479184	-0.00303792833347716\\
-18303.9921865462	-0.00285265343988584\\
-18280.551825174	-0.0026263796659168\\
-18257.1114638017	-0.00236673031190815\\
-18239.5311927726	-0.0021553327840013\\
-18221.9509217434	-0.00193411008510074\\
-18204.3706507142	-0.0017072890074432\\
-18186.790379685	-0.00147921149032229\\
-18169.2101086558	-0.00125425218065833\\
-18157.4899279697	-0.00110815159609184\\
-18145.7697472836	-0.000966613188023733\\
-18134.0495665975	-0.00083085174866964\\
-18122.3293859114	-0.000702034797422572\\
-18110.6092052252	-0.000581272632318132\\
-18098.8890245391	-0.000469608852506454\\
-18087.168843853	-0.000368011432396361\\
-18081.3087535099	-0.000321267364520118\\
-18075.4486631669	-0.000277364423876758\\
-18069.5885728238	-0.000236398757207158\\
-18063.7284824808	-0.000198460357582423\\
-18057.8683921377	-0.000163632869696458\\
-18052.0083017947	-0.000131993408469772\\
-18046.1482114516	-0.000103612391332356\\
-18040.2881211085	-7.85533845819969e-05\\
-18034.4280307655	-5.68729641296078e-05\\
-18028.5679404224	-3.86205909431193e-05\\
-18022.7078500794	-2.38385014690801e-05\\
-18016.8477597363	-1.256161326216e-05\\
-18010.9876693932	-4.81744604422046e-06\\
-18005.1275790502	-6.26058363293904e-07\\
nan	nan\\
-17993.4073983641	-2.94428023271915e-06\\
-17987.547308021	-9.45635203853158e-06\\
-17981.6872176779	-1.95261122754273e-05\\
-17975.8271273349	-3.31359178562001e-05\\
-17969.9670369918	-5.0260617890078e-05\\
-17964.1069466488	-7.08676017363276e-05\\
-17958.2468563057	-9.49168628787064e-05\\
-17952.3867659626	-0.000122361078494211\\
-17946.5266756196	-0.00015314570456439\\
-17940.6665852765	-0.000187209086336214\\
-17934.8064949335	-0.000224482583897122\\
-17928.9464045904	-0.000264890712639448\\
-17923.0863142473	-0.00030835129829279\\
-17917.2262239043	-0.000354775646231425\\
-17911.3661335612	-0.000404068724687542\\
-17899.6459528751	-0.000510850454125667\\
-17887.925772189	-0.000627817291522315\\
-17876.2055915029	-0.000754002566785259\\
-17864.4854108167	-0.000888360196592567\\
-17852.7652301306	-0.00102977342971059\\
-17841.0450494445	-0.00117706420270445\\
-17823.4647784153	-0.00140631971993997\\
-17805.8845073862	-0.00164171214863288\\
-17788.304236357	-0.00187880394986281\\
-17770.7239653278	-0.00211311584269286\\
-17753.1436942986	-0.00234021200491103\\
-17735.5634232694	-0.00255578463012177\\
-17712.1230618972	-0.00281824737597132\\
-17688.682700525	-0.00304409270797898\\
-17665.2423391527	-0.00322565414337754\\
-17641.8019777805	-0.00335674158194467\\
-17618.3616164083	-0.00343285499842081\\
-17594.921255036	-0.0034513416214765\\
-17571.4808936638	-0.00341149112122121\\
-17548.0405322915	-0.00331456540429649\\
-17524.6001709193	-0.0031637618160903\\
-17501.1598095471	-0.00296411079482232\\
-17477.7194481748	-0.00272231124081693\\
-17454.2790868026	-0.00244650897752015\\
-17436.6988157734	-0.00222291492083514\\
-17419.1185447442	-0.00198968700024413\\
-17401.5382737151	-0.00175128573374146\\
-17383.9580026859	-0.00151228062114851\\
-17366.3777316567	-0.00127726338193065\\
-17354.6575509706	-0.00112504134473395\\
-17342.9373702845	-0.000977913898139755\\
-17331.2171895983	-0.00083714722528202\\
-17319.4970089122	-0.000703955630570854\\
-17307.7768282261	-0.000579491176827423\\
-17296.05664754	-0.00046483383308878\\
-17290.1965571969	-0.000411498312991594\\
-17284.3364668539	-0.000360982217294307\\
-17278.4763765108	-0.000313396253729843\\
-17272.6162861677	-0.000268845013054363\\
-17266.7561958247	-0.000227426742778135\\
-17260.8961054816	-0.0001892331341979\\
-17255.0360151386	-0.000154349123092558\\
-17249.1759247955	-0.000122852704579042\\
-17243.3158344524	-9.48147625138472e-05\\
-17237.4557441094	-7.02989137991503e-05\\
-17231.5956537663	-4.93613679520115e-05\\
-17225.7355634233	-3.20508022419686e-05\\
-17219.8754730802	-1.84082526626683e-05\\
-17214.0153827372	-8.467020996523e-06\\
-17208.1552923941	-2.25259817053585e-06\\
-17202.295202051	2.17395917606153e-07\\
-17196.435111708	-1.06674403792564e-06\\
-17190.5750213649	-6.10678190334502e-06\\
-17184.7149310219	-1.48965300430669e-05\\
-17178.8548406788	-2.74218560787596e-05\\
-17172.9947503357	-4.36607064672648e-05\\
-17167.1346599927	-6.35831468621092e-05\\
-17161.2745696496	-8.71514191847529e-05\\
-17155.4144793066	-0.0001143200153026\\
-17149.5543889635	-0.000145035767163624\\
-17143.6942986204	-0.000179237953221132\\
-17137.8342082774	-0.000216858420919931\\
-17131.9741179343	-0.000257821725022549\\
-17126.1140275913	-0.000302045281462762\\
-17120.2539372482	-0.000349439536447095\\
-17114.3938469051	-0.00039990815044365\\
-17102.673666219	-0.000509650373644626\\
-17090.9534855329	-0.000630373411576671\\
-17079.2333048468	-0.000761084283144257\\
-17067.5131241607	-0.000900703820765691\\
-17055.7929434745	-0.00104807566845909\\
-17044.0727627884	-0.00120197594788073\\
-17026.4924917592	-0.00144225164530541\\
-17008.9122207301	-0.00168981450016808\\
-16991.3319497009	-0.00194000642106733\\
-16973.7516786717	-0.00218810788894791\\
-16956.1714076425	-0.00242942768344873\\
-16938.5911366134	-0.00265939229299421\\
-16915.1507752411	-0.00294083786981587\\
-16891.7104138689	-0.0031848439498138\\
-16868.2700524966	-0.00338312646340058\\
-16844.8296911244	-0.00352892142605913\\
-16821.3893297522	-0.00361721888623165\\
-16797.9489683799	-0.00364493834982327\\
-16774.5086070077	-0.00361103955292185\\
-16751.0682456355	-0.00351656461909876\\
-16727.6278842632	-0.00336460994172886\\
-16704.187522891	-0.00316022850187203\\
-16680.7471615187	-0.00291026568369465\\
-16657.3068001465	-0.00262313390868933\\
-16639.7265291173	-0.00238918878177739\\
-16622.1462580881	-0.00214423714665986\\
-16604.565987059	-0.00189296510718047\\
-16586.9857160298	-0.00164019175297522\\
-16569.4054450006	-0.00139077773104547\\
-16557.6852643145	-0.00122874806979584\\
-16545.9650836284	-0.00107174565352221\\
-16534.2449029423	-0.000921124013042951\\
-16522.5247222561	-0.00077818507271979\\
-16510.80454157	-0.00064416803343282\\
-16499.0843608839	-0.000520238766584936\\
-16487.3641801978	-0.000407479809935973\\
-16481.5040898547	-0.000355602396256526\\
-16475.6439995117	-0.000306881050907357\\
-16469.7839091686	-0.000261423630404602\\
-16463.9238188255	-0.000219331177446408\\
-16458.0637284825	-0.000180697700851392\\
-16452.2036381394	-0.000145609970177582\\
-16446.3435477964	-0.000114147325468672\\
-16440.4834574533	-8.63815025671135e-05\\
-16434.6233671102	-6.2376474350329e-05\\
-16428.7632767672	-4.21883082618168e-05\\
-16422.9031864241	-2.58650404379933e-05\\
-16417.0430960811	-1.344656672264e-05\\
-16411.183005738	-4.96455080859517e-06\\
-16405.3229153949	-4.42349719875634e-07\\
-16399.4628250519	1.05043195003138e-07\\
-16393.6027347088	-3.32896238068003e-06\\
-16387.7426443658	-1.07425321233418e-05\\
-16381.8825540227	-2.21254032706356e-05\\
-16376.0224636796	-3.74588986587459e-05\\
-16370.1623733366	-5.6715958578948e-05\\
-16364.3022829935	-7.98611904064394e-05\\
-16358.4421926505	-0.000106850935909349\\
-16352.5821023074	-0.000137633356111667\\
-16346.7220119643	-0.00017214853354908\\
-16340.8619216213	-0.000210328591719765\\
-16335.0018312782	-0.000252097831477033\\
-16329.1417409352	-0.000297372884124726\\
-16323.2816505921	-0.000346062880858526\\
-16317.4215602491	-0.000398069638266874\\
-16311.561469906	-0.00045328785944793\\
-16299.8412892199	-0.00057290325122892\\
-16288.1211085338	-0.000703932999889502\\
-16276.4009278476	-0.000845302577438621\\
-16264.6807471615	-0.000995847901314616\\
-16252.9605664754	-0.00115432502617893\\
-16241.2403857893	-0.00131942052976436\\
-16223.6601147601	-0.0015764590329577\\
-16206.0798437309	-0.00184047531004613\\
-16188.4995727017	-0.00210650727740089\\
-16170.9193016726	-0.00236954064236121\\
-16153.3390306434	-0.00262460426284351\\
-16135.7587596142	-0.00286686495555605\\
-16112.318398242	-0.00316205242376054\\
-16088.8780368697	-0.00341634258813389\\
-16065.4376754975	-0.00362109170310465\\
-16041.9973141253	-0.00376930052668983\\
-16018.556952753	-0.00385585748052977\\
-15995.1165913808	-0.00387771890197957\\
-15971.6762300085	-0.00383402005160874\\
-15948.2358686363	-0.00372611285472477\\
-15924.7955072641	-0.00355752881660454\\
-15901.3551458918	-0.00333386807227135\\
-15877.9147845196	-0.00306261803341563\\
-15854.4744231474	-0.00275290748649402\\
-15836.8941521182	-0.002501641152423\\
-15819.313881089	-0.0022394075006136\\
-15801.7336100598	-0.00197122912742261\\
-15784.1533390306	-0.00170225724090561\\
-15766.5730680015	-0.00143767392686483\\
-15754.8528873153	-0.00126625449902387\\
-15743.1327066292	-0.00110054047939344\\
-15731.4125259431	-0.000941965042732009\\
-15719.692345257	-0.000791903830388467\\
-15707.9721645709	-0.000651663183166302\\
-15696.2519838848	-0.000522468936991754\\
-15690.3918935417	-0.000462372664789287\\
-15684.5318031986	-0.000405455877350871\\
-15678.6717128556	-0.000351844588788032\\
-15672.8116225125	-0.000301657943000936\\
-15666.9515321695	-0.000255007954396796\\
-15661.0914418264	-0.000211999263428327\\
-15655.2313514833	-0.00017272890747811\\
-15649.3712611403	-0.000137286107612723\\
-15643.5111707972	-0.000105752071654079\\
-15637.6510804542	-7.81998140127464e-05\\
-15631.7909901111	-5.46939926726159e-05\\
-15625.930899768	-3.52907636955655e-05\\
-15620.070809425	-2.00376535697323e-05\\
-15614.2107190819	-8.97344969084332e-06\\
-15608.3506287389	-2.1281092255186e-06\\
-15602.4905383958	4.773134301401e-07\\
-15596.6304480527	-1.16927957373461e-06\\
-15590.7703577097	-7.07099466777598e-06\\
-15584.9102673666	-1.72219309810085e-05\\
-15579.0501770236	-3.16071835082891e-05\\
-15573.1900866805	-5.0202865334378e-05\\
-15567.3299963374	-7.2976148895099e-05\\
-15561.4699059944	-9.98853261971183e-05\\
-15555.6098156513	-0.000130879887902337\\
-15549.7497253083	-0.000165900621115956\\
-15543.8896349652	-0.000204879725700828\\
-15538.0295446221	-0.000247740948886146\\
-15532.1694542791	-0.000294399737903101\\
-15526.309363936	-0.000344763410350772\\
-15520.449273593	-0.000398731341943496\\
-15514.5891832499	-0.000456195171261834\\
-15502.8690025638	-0.000581139735917738\\
-15491.1488218777	-0.000718584280846323\\
-15479.4286411916	-0.00086740798862672\\
-15467.7084605054	-0.00102639131353052\\
-15455.9882798193	-0.00119422610243507\\
-15444.2680991332	-0.00136952648722847\\
-15426.687828104	-0.00164328368143192\\
-15409.1075570748	-0.00192544384110782\\
-15391.5272860457	-0.00221071580829229\\
-15373.9470150165	-0.00249373244525841\\
-15356.3667439873	-0.00276915263604038\\
-15338.7864729581	-0.00303176311997486\\
-15315.3461115859	-0.00335341180118523\\
-15291.9057502136	-0.00363257930993113\\
-15268.4653888414	-0.00385977550368891\\
-15245.0250274692	-0.00402722956387104\\
-15221.5846660969	-0.00412916048897852\\
-15198.1443047247	-0.00416198145516016\\
-15174.7039433525	-0.0041244308448573\\
-15151.2635819802	-0.00401762519138437\\
-15127.823220608	-0.00384503190591027\\
-15104.3828592357	-0.00361236235425926\\
-15080.9424978635	-0.00332738854628988\\
-15057.5021364913	-0.00299968929659867\\
-15039.9218654621	-0.00273248936690219\\
-15022.3415944329	-0.00245256272851965\\
-15004.7613234037	-0.00216527240276336\\
-14987.1810523746	-0.00187614027613619\\
-14969.6007813454	-0.00159074244143019\\
-14957.8806006593	-0.00140528517846411\\
-14946.1604199731	-0.00122554832635953\\
-14934.440239287	-0.00105308939068388\\
-14922.7200586009	-0.000889407827800184\\
-14910.9998779148	-0.000735932211527269\\
-14899.2796972287	-0.000594007970482249\\
-14887.5595165425	-0.000464885801107435\\
-14881.6994261995	-0.000405487754519023\\
-14875.8393358564	-0.000349710856090986\\
-14869.9792455134	-0.00029768012719494\\
-14864.1191551703	-0.000249512801427311\\
-14858.2590648272	-0.000205318067391454\\
-14852.3989744842	-0.000165196828188351\\
-14846.5388841411	-0.000129241478160138\\
-14840.6787937981	-9.75356973902043e-05\\
-14834.818703455	-7.01542644141694e-05\\
-14828.958613112	-4.7162887539361e-05\\
-14823.0985227689	-2.86180551753425e-05\\
-14817.2384324258	-1.45669054970921e-05\\
-14811.3783420828	-5.04711574925861e-06\\
-14805.5182517397	-8.6811443922328e-08\\
-14799.6581613967	2.95504333086003e-07\\
-14793.7980710536	-3.90899866785218e-06\\
-14787.9379807105	-1.26994478629447e-05\\
-14782.0778903675	-2.60652583372505e-05\\
-14776.2178000244	-4.39861438923031e-05\\
-14770.3577096814	-6.64321486381091e-05\\
-14764.4976193383	-9.33636990873341e-05\\
-14758.6375289952	-0.000124731676664487\\
-14752.7774386522	-0.000160477510506018\\
-14746.9173483091	-0.000200533290366235\\
-14741.0572579661	-0.000244821899432677\\
-14735.197167623	-0.000293257166773564\\
-14729.3370772799	-0.00034574403913085\\
-14723.4769869369	-0.000402178771716434\\
-14717.6168965938	-0.000462449137631033\\
-14711.7568062508	-0.00052643465546198\\
-14700.0366255646	-0.00066502943812183\\
-14688.3164448785	-0.000816843928538207\\
-14676.5962641924	-0.000980644310672026\\
-14664.8760835063	-0.00115509232217937\\
-14653.1559028202	-0.00133875634153709\\
-14641.4357221341	-0.00153012329196357\\
-14623.8554511049	-0.00182814070969855\\
-14606.2751800757	-0.00213435837440913\\
-14588.6949090465	-0.00244304124377386\\
-14571.1146380173	-0.00274838688313236\\
-14553.5343669882	-0.00304463577926302\\
-14535.954095959	-0.00332618123636432\\
-14512.5137345867	-0.00366951270763007\\
-14489.0733732145	-0.00396561708263934\\
-14465.6330118423	-0.00420441347361459\\
-14442.19265047	-0.00437771359764474\\
-14418.7522890978	-0.00447950828613353\\
-14395.3119277256	-0.00450618167305567\\
-14371.8715663533	-0.00445664546272161\\
-14348.4312049811	-0.00433238833977242\\
-14324.9908436088	-0.00413743842996907\\
-14301.5504822366	-0.00387823964719625\\
-14278.1101208644	-0.00356344567879536\\
-14254.6697594921	-0.00320363816260212\\
-14237.0894884629	-0.00291150977669876\\
-14219.5092174338	-0.0026064623691151\\
-14201.9289464046	-0.00229434862998998\\
-14184.3486753754	-0.00198117875579589\\
-14166.7684043462	-0.001673006530483\\
-14155.0482236601	-0.00147329542822453\\
-14143.328042974	-0.00128019931801013\\
-14131.6078622879	-0.00109539790317468\\
-14119.8876816018	-0.000920504995117792\\
-14108.1675009156	-0.000757054672722113\\
-14096.4473202295	-0.000606488081124548\\
-14090.5872298865	-0.000536458049918619\\
-14084.7271395434	-0.000470140983008642\\
-14078.8670492003	-0.000407685483172795\\
-14073.0069588573	-0.00034923216998541\\
-14067.1468685142	-0.00029491337167422\\
-14061.2867781712	-0.000244852834125261\\
-14055.4266878281	-0.000199165447665733\\
-14049.566597485	-0.000157956992237147\\
-14043.706507142	-0.000121323901525574\\
-14037.8464167989	-8.93530465502477e-05\\
-14031.9863264559	-6.21215392153655e-05\\
-14026.1262361128	-3.96965562601339e-05\\
-14020.2661457697	-2.21351840003748e-05\\
-14014.4060554267	-9.48428422582442e-06\\
-14008.5459650836	-1.78038155813986e-06\\
-14002.6858747406	9.50427459075276e-07\\
-13996.8257843975	-1.30745667827178e-06\\
-13990.9656940544	-8.5590895997881e-06\\
-13985.1056037114	-2.07989584786331e-05\\
-13979.2455133683	-3.80109797838483e-05\\
-13973.3854230253	-6.01685194012163e-05\\
-13967.5253326822	-8.72344351123164e-05\\
-13961.6652423392	-0.000119161141364189\\
-13955.8051519961	-0.000155890696233735\\
-13949.945061653	-0.000197354910416347\\
-13944.08497131	-0.000243475478057306\\
-13938.2248809669	-0.000294164129152046\\
-13932.3647906239	-0.000349322803258595\\
-13926.5047002808	-0.000408843844137917\\
-13920.6446099377	-0.00047261021497196\\
-13914.7845195947	-0.000540495733723048\\
-13903.0643389086	-0.000688075310060081\\
-13891.3441582224	-0.000850400377821322\\
-13879.6239775363	-0.00102616068145417\\
-13867.9037968502	-0.00121392866245647\\
-13856.1836161641	-0.00141217123863335\\
-13844.463435478	-0.00161926250841621\\
-13826.8831644488	-0.00194274240345282\\
-13809.3028934196	-0.00227626601918929\\
-13791.7226223904	-0.00261360398102513\\
-13774.1423513612	-0.00294842926925367\\
-13756.5620803321	-0.00327443739756012\\
-13738.9818093029	-0.00358546665249678\\
-13715.5414479307	-0.00396672453832547\\
-13692.1010865584	-0.00429800319315724\\
-13668.6607251862	-0.00456802352855954\\
-13645.2203638139	-0.00476752102761162\\
-13621.7800024417	-0.00488957061478979\\
-13598.3396410695	-0.00492983415279872\\
-13574.8992796972	-0.00488672176510969\\
-13551.458918325	-0.00476146104900936\\
-13528.0185569528	-0.00455807132229245\\
-13504.5781955805	-0.00428324324341816\\
-13481.1378342083	-0.00394612734397704\\
-13457.697472836	-0.0035580381126893\\
-13440.1172018069	-0.00324134948453528\\
-13422.5369307777	-0.0029093868254581\\
-13404.9566597485	-0.00256852132462072\\
-13387.3763887193	-0.00222532187206752\\
-13369.7961176901	-0.00188643071881885\\
-13358.075937004	-0.00166615676518477\\
-13346.3557563179	-0.00145264326805179\\
-13334.6355756318	-0.0012477519367211\\
-13322.9153949457	-0.00105327689021096\\
-13311.1952142596	-0.000870929261528749\\
-13299.4750335734	-0.000702322459742393\\
-13287.7548528873	-0.000548958216667969\\
-13281.8947625443	-0.000478427577543718\\
-13276.0346722012	-0.000412213538406629\\
-13270.1745818581	-0.000350466762573949\\
-13264.3144915151	-0.000293328674839658\\
-13258.454401172	-0.000240931148450913\\
-13252.594310829	-0.000193396211866665\\
-13246.7342204859	-0.000150835775911191\\
-13240.8741301428	-0.000113351381929817\\
-13235.0140397998	-8.10339715372089e-05\\
-13229.1539494567	-5.39636784412149e-05\\
-13223.2938591137	-3.22096428508621e-05\\
-13217.4337687706	-1.58298488625258e-05\\
-13211.5736784275	-4.87098522725504e-06\\
-13205.7135880845	6.3167019038465e-07\\
-13199.8534977414	6.54341970378372e-07\\
-13193.9934073984	-4.8151755103121e-06\\
-13188.1333170553	-1.57774751126702e-05\\
-13182.2732267122	-3.22215184816521e-05\\
-13176.4131363692	-5.41246421823077e-05\\
-13170.5530460261	-8.14525884456714e-05\\
-13164.6929556831	-0.00011415956049917\\
-13158.83286534	-0.000152188302400557\\
-13152.9727749969	-0.000195470203237305\\
-13147.1126846539	-0.000243925425517103\\
-13141.2525943108	-0.000297463057510905\\
-13135.3925039678	-0.00035598128924929\\
-13129.5324136247	-0.000419367611854463\\
-13123.6723232817	-0.000487499039817409\\
-13117.8122329386	-0.000560242355756582\\
-13106.0920522525	-0.000718982245030837\\
-13094.3718715664	-0.000894327574196526\\
-13082.6516908802	-0.00108487418502644\\
-13070.9315101941	-0.00128908492892611\\
-13059.211329508	-0.00150530228981414\\
-13047.4911488219	-0.00173176206430656\\
-13035.7709681358	-0.0019666079865307\\
-13018.1906971066	-0.00233035573420049\\
-13000.6104260774	-0.00270184819879727\\
-12983.0301550482	-0.0030741463658034\\
-12965.449884019	-0.00344026461380296\\
-12947.8696129899	-0.00379330365670004\\
-12924.4292516176	-0.00423217700209131\\
-12900.9888902454	-0.00462113957348873\\
-12877.5485288732	-0.00494697168780536\\
-12854.1081675009	-0.00519851312115898\\
-12830.6678061287	-0.0053670539438863\\
-12807.2274447564	-0.00544664484648214\\
-12783.7870833842	-0.0054343159556759\\
-12760.346722012	-0.0053301961663468\\
-12736.9063606397	-0.00513752832805438\\
-12713.4659992675	-0.00486257911584817\\
-12690.0256378953	-0.00451444596939408\\
-12666.585276523	-0.00410476697843897\\
-12649.0050054938	-0.00376535139099728\\
-12631.4247344647	-0.00340552026183173\\
-12613.8444634355	-0.00303217118216235\\
-12596.2641924063	-0.00265249363169663\\
-12578.6839213771	-0.00227383334547552\\
-12561.1036503479	-0.00190355304595082\\
-12549.3834696618	-0.00166493937274596\\
-12537.6632889757	-0.00143535901862137\\
-12525.9431082896	-0.0012168282831699\\
-12514.2229276035	-0.00101127650213027\\
-12502.5027469173	-0.000820529415734817\\
-12490.7825662312	-0.000646293359927287\\
-12484.9224758882	-0.000565863872680999\\
-12479.0623855451	-0.00049014041669363\\
-12473.2022952021	-0.000419296266918004\\
-12467.342204859	-0.000353494652429067\\
-12461.4821145159	-0.000292888394196339\\
-12455.6220241729	-0.000237619564309568\\
-12449.7619338298	-0.000187819167386013\\
-12443.9018434868	-0.000143606844864621\\
-12438.0417531437	-0.000105090602865757\\
-12432.1816628006	-7.23665642196277e-05\\
-12426.3215724576	-4.55187452331639e-05\\
-12420.4614821145	-2.4618857708653e-05\\
-12414.6013917715	-9.72613667174347e-06\\
-12408.7413014284	-8.87194216363841e-07\\
-12402.8812110853	1.86410017836111e-06\\
-12397.0211207423	-1.49328743399849e-06\\
-12391.1610303992	-1.09674895595291e-05\\
-12385.3009400562	-2.65536985479246e-05\\
-12379.4408497131	-4.82341551947391e-05\\
-12373.58075937	-7.59781647374057e-05\\
-12367.720669027	-0.000109742140246979\\
-12361.8605786839	-0.00014946967337851\\
-12356.0004883409	-0.000195091632371791\\
-12350.1403979978	-0.000246526287144587\\
-12344.2803076547	-0.000303679461249289\\
-12338.4202173117	-0.000366444710434796\\
-12332.5601269686	-0.000434703527455033\\
-12326.7000366256	-0.000508325572768237\\
-12320.8399462825	-0.000587168930638313\\
-12314.9798559394	-0.000671080390199808\\
-12303.2596752533	-0.00085344015153012\\
-12291.5394945672	-0.00105396546207741\\
-12279.8193138811	-0.00127105816443956\\
-12268.099133195	-0.0015029743429901\\
-12256.3789525089	-0.00174783860663117\\
-12244.6587718227	-0.00200365953334869\\
-12227.0785007936	-0.00240333770227891\\
-12209.4982297644	-0.00281555983069255\\
-12191.9179587352	-0.00323266480331906\\
-12174.337687706	-0.00364685960373361\\
-12156.7574166768	-0.00405036721391871\\
-12139.1771456477	-0.00443557497323858\\
-12115.7367842754	-0.00490815999641268\\
-12092.2964229032	-0.00531928227074619\\
-12068.8560615309	-0.00565491675940761\\
-12045.4157001587	-0.00590350227901801\\
-12021.9753387865	-0.00605635053260512\\
-11998.5349774142	-0.00610796074592958\\
-11975.094616042	-0.00605622859673263\\
-11951.6542546698	-0.00590254162310263\\
-11928.2138932975	-0.00565175707812776\\
-11904.7735319253	-0.00531206214531751\\
-11881.333170553	-0.00489472040428386\\
-11857.8928091808	-0.00441371231139647\\
-11840.3125381516	-0.00402087609760015\\
-11822.7322671225	-0.00360884650759666\\
-11805.1519960933	-0.00318554936262997\\
-11787.5717250641	-0.00275916968489288\\
-11769.9914540349	-0.00233799705551539\\
-11758.2712733488	-0.00206418253330767\\
-11746.5510926627	-0.00179874230552348\\
-11734.8309119766	-0.0015440093137091\\
-11723.1107312904	-0.00130223446394367\\
-11711.3905506043	-0.0010755672505179\\
-11699.6703699182	-0.000866037166831683\\
-11693.8102795751	-0.000768295589181372\\
-11687.9501892321	-0.000675536063791807\\
-11682.090098889	-0.000587970851146087\\
-11676.230008546	-0.000505801608369831\\
-11670.3699182029	-0.00042921894251929\\
-11664.5098278598	-0.000358401986397821\\
-11658.6497375168	-0.000293517997887542\\
-11652.7896471737	-0.000234721983653205\\
-11646.9295568307	-0.000182156348062746\\
-11641.0694664876	-0.00013595056812842\\
-11635.2093761446	-9.62208952007329e-05\\
-11629.3492858015	-6.3070084082781e-05\\
-11623.4891954584	-3.65871502018753e-05\\
-11617.6291051154	-1.68471553991068e-05\\
-11611.7690147723	-3.91102284497095e-06\\
-11605.9089244293	2.17461847321223e-06\\
-11600.0488340862	1.37755931010916e-06\\
-11594.1887437431	-6.31989455759075e-06\\
-11588.3286534001	-2.09208575862028e-05\\
-11582.468563057	-4.24138304485396e-05\\
-11576.608472714	-7.07726968791103e-05\\
-11570.7483823709	-0.000105956751441114\\
-11564.8882920278	-0.000147910758218454\\
-11559.0282016848	-0.000196565040367412\\
-11553.1681113417	-0.000251835600389362\\
-11547.3080209987	-0.000313624270939049\\
-11541.4479306556	-0.000381818895900072\\
-11535.5878403125	-0.000456293541417036\\
-11529.7277499695	-0.000536908736441398\\
-11523.8676596264	-0.000623511742444622\\
-11518.0075692834	-0.000715936851649411\\
-11506.2873885972	-0.000917527687418584\\
-11494.5672079111	-0.00114010927397309\\
-11482.847027225	-0.00138192579859688\\
-11471.1268465389	-0.00164105178308009\\
-11459.4066658528	-0.00191540775218794\\
-11447.6864851666	-0.00220277726615378\\
-11435.9663044805	-0.00250082517838853\\
-11418.3860334513	-0.00296257132498431\\
-11400.8057624222	-0.00343431947684863\\
-11383.225491393	-0.00390730200535276\\
-11365.6452203638	-0.00437267762882282\\
-11348.0649493346	-0.0048216995929385\\
-11324.6245879624	-0.00538034952354886\\
-11301.1842265902	-0.00587602711381145\\
-11277.7438652179	-0.00629185935513033\\
-11254.3035038457	-0.0066135519109268\\
-11230.8631424734	-0.00682989408722613\\
-11207.4227811012	-0.00693316337742769\\
-11183.982419729	-0.00691941509582266\\
-11160.5420583567	-0.00678864638302746\\
-11137.1016969845	-0.00654482802896784\\
-11113.6613356123	-0.00619580197600453\\
-11090.22097424	-0.00575304688685145\\
-11066.7806128678	-0.00523131864439613\\
-11049.2003418386	-0.00479866978748966\\
-11031.6200708094	-0.00433969099928981\\
-11014.0397997802	-0.00386319812857683\\
-10996.4595287511	-0.00337839706934762\\
-10978.8792577219	-0.00289471060746283\\
-10961.2989866927	-0.00242160006130094\\
-10949.5788060066	-0.00211668208161304\\
-10937.8586253205	-0.00182330269912859\\
-10926.1384446344	-0.00154406501829236\\
-10914.4182639482	-0.00128146330477403\\
-10902.6980832621	-0.00103786139158772\\
-10890.977902576	-0.000815472099262405\\
-10885.1178122329	-0.000712877601540279\\
-10879.2577218899	-0.000616337848460132\\
-10873.3976315468	-0.000526078559795566\\
-10867.5375412038	-0.000442312634040603\\
-10861.6774508607	-0.000365239670544749\\
-10855.8173605176	-0.000295045518808255\\
-10849.9572701746	-0.000231901855888585\\
-10844.0971798315	-0.000175965792927513\\
-10838.2370894885	-0.000127379511635701\\
-10832.3769991454	-8.62699316219065e-05\\
-10826.5169088023	-5.27484092917896e-05\\
-10820.6568184593	-2.69104690456881e-05\\
-10814.7967281162	-8.83556739862741e-06\\
-10808.9366377732	1.41310940428786e-06\\
-10803.0765474301	3.78881377659286e-06\\
-10797.216457087	-1.73862924079821e-06\\
-10791.356366744	-1.51827227904997e-05\\
-10785.4962764009	-3.65402340709376e-05\\
-10779.6361860579	-6.57911654813149e-05\\
-10773.7760957148	-0.000102898761174201\\
-10767.9160053717	-0.000147809549082245\\
-10762.0559150287	-0.000200453418395766\\
-10756.1958246856	-0.000260743732411699\\
-10750.3357343426	-0.000328577476596296\\
-10744.4756439995	-0.000403835441610912\\
-10738.6155536565	-0.000486382440981392\\
-10732.7554633134	-0.000576067563078175\\
-10726.8953729703	-0.000672724456844635\\
-10721.0352826273	-0.000776171650899172\\
-10715.1751922842	-0.000886212905200961\\
-10703.4550115981	-0.00112522112495846\\
-10691.734830912	-0.00138789918018199\\
-10680.0146502259	-0.00167218857319558\\
-10668.2944695397	-0.00197583878746574\\
-10656.5742888536	-0.00229642552848497\\
-10644.8541081675	-0.00263137051536534\\
-10627.2738371383	-0.00315474838405956\\
-10609.6935661091	-0.0036947155029711\\
-10592.11329508	-0.00424130234704384\\
-10574.5330240508	-0.00478434851646065\\
-10556.9527530216	-0.00531369529496827\\
-10539.3724819924	-0.00581937945383468\\
-10515.9321206202	-0.00644034369319348\\
-10492.491759248	-0.00698124630269803\\
-10469.0513978757	-0.00742358494436153\\
-10445.6110365035	-0.00775204125828384\\
-10422.1706751312	-0.00795502908235959\\
-10398.730313759	-0.0080251212004464\\
-10375.2899523868	-0.00795933907754341\\
-10351.8495910145	-0.00775929453583125\\
-10328.4092296423	-0.00743117722572026\\
-10304.9688682701	-0.00698558690704567\\
-10281.5285068978	-0.00643721478891272\\
-10258.0881455256	-0.00580438332047388\\
-10240.5078744964	-0.00528709101054839\\
-10222.9276034672	-0.00474417623187421\\
-10205.347332438	-0.00418611565386137\\
-10187.7670614089	-0.00362375093245474\\
-10170.1867903797	-0.00306808447021394\\
-10158.4666096936	-0.00270677930417743\\
-10146.7464290074	-0.00235651582489034\\
-10135.0262483213	-0.0020204066127214\\
-10123.3060676352	-0.00170145931781868\\
-10111.5858869491	-0.00140255066733593\\
-10099.865706263	-0.00112640145748773\\
-10094.0056159199	-0.000997663367469788\\
-10088.1455255769	-0.000875552746390671\\
-10082.2854352338	-0.000760355466990062\\
-10076.4253448907	-0.000652343454739857\\
-10070.5652545477	-0.000551774078525254\\
-10064.7051642046	-0.000458889570673991\\
-10058.8450738616	-0.000373916477591258\\
-10052.9849835185	-0.000297065142248737\\
-10047.1248931754	-0.000228529219733049\\
-10041.2648028324	-0.00016848522688702\\
-10035.4047124893	-0.00011709212714911\\
-10029.5446221463	-7.44909514947628e-05\\
-10023.6845318032	-4.08044564115603e-05\\
-10017.8244414601	-1.61368196951621e-05\\
-10011.9643511171	-5.73374805836168e-07\\
-10006.104260774	5.81961555461356e-06\\
-10000.244170431	2.99514607758894e-06\\
-9994.3840800879	-9.07438479967228e-06\\
-9988.52398974484	-3.03970692294141e-05\\
-9982.66389940178	-6.09614252983426e-05\\
-9976.80380905872	-0.000100736373622287\\
-9970.94371871567	-0.000149671255333666\\
-9965.08362837261	-0.000207695891532479\\
-9959.22353802955	-0.000274720684174514\\
-9953.36344768649	-0.000350636758274933\\
-9947.50335734343	-0.0004353161452437\\
-9941.64326700037	-0.000528612007057587\\
-9935.78317665731	-0.000630358900904697\\
-9929.92308631425	-0.000740373083796323\\
-9924.06299597119	-0.000858452856706486\\
-9918.20290562813	-0.000984378947466937\\
-9906.48272494201	-0.00125880769181348\\
-9894.76254425589	-0.00156157060253946\\
-9883.04236356977	-0.0018903285430994\\
-9871.32218288365	-0.00224251054200134\\
-9859.60200219753	-0.00261533443557092\\
-9847.88182151142	-0.00300582940147332\\
-9836.1616408253	-0.00341086020214953\\
-9818.58136979612	-0.00403846816739487\\
-9801.00109876694	-0.00467989879912989\\
-9783.42082773776	-0.00532331275057724\\
-9765.84055670858	-0.00595674513375714\\
-9748.2602856794	-0.00656833288766489\\
-9724.81992430716	-0.00732994253437018\\
-9701.37956293493	-0.00800656018087187\\
-9677.93920156269	-0.00857510017185733\\
-9654.49884019045	-0.00901591725888633\\
-9631.05847881822	-0.0093135087368781\\
-9607.61811744598	-0.00945708349465288\\
-9584.17775607374	-0.00944097734400312\\
-9560.7373947015	-0.00926489897379619\\
-9537.29703332927	-0.00893399643703047\\
-9513.85667195703	-0.00845874003762402\\
-9490.41631058479	-0.00785462364526479\\
-9466.97594921255	-0.00714169262402612\\
-9449.39567818337	-0.00654987514561934\\
-9431.81540715419	-0.00592157794692808\\
-9414.23513612501	-0.00526890509852343\\
-9396.65486509584	-0.00460452743791083\\
-9379.07459406666	-0.00394144539697287\\
-9361.49432303748	-0.00329274361819302\\
-9349.77414235136	-0.00287465926948797\\
-9338.05396166524	-0.00247245309246311\\
-9326.33378097912	-0.00208974762412753\\
-9314.613600293	-0.00173002032881896\\
-9302.89341960689	-0.0013965733222765\\
-9297.03332926383	-0.00124068173402229\\
-9291.17323892077	-0.00109250441344431\\
-9285.31314857771	-0.000952392169648613\\
-9279.45305823465	-0.000820679716238142\\
-9273.59296789159	-0.000697684918003724\\
-9267.73287754853	-0.000583708071202371\\
-9261.87278720547	-0.000479031218977657\\
-9256.01269686241	-0.000383917503452488\\
-9250.15260651935	-0.000298610556014676\\
-9244.29251617629	-0.000223333927157371\\
-9238.43242583323	-0.000158290557183506\\
-9232.57233549017	-0.000103662289029873\\
-9226.71224514711	-5.96094243329286e-05\\
-9220.85215480405	-2.62703238093769e-05\\
-9214.992064461	-3.76105290718477e-06\\
-9209.13197411794	7.82492638399595e-06\\
-9203.27188377488	8.41701679615916e-06\\
-9197.41179343182	-2.03230056979739e-06\\
-9191.55170308876	-2.35473007952853e-05\\
-9185.6916127457	-5.61288983802229e-05\\
-9179.83152240264	-9.97545795444341e-05\\
-9173.97143205958	-0.000154378383879877\\
-9168.11134171652	-0.000219930935522611\\
-9162.25125137346	-0.00029631952393277\\
-9156.3911610304	-0.000383428234241947\\
-9150.53107068734	-0.000481118127020635\\
-9144.67098034428	-0.000589227467265248\\
-9138.81089000122	-0.00070757200220806\\
-9132.95079965816	-0.000835945287546409\\
-9127.0907093151	-0.000974119061537158\\
-9121.23061897204	-0.00112184366623683\\
-9115.37052862898	-0.00127884851523458\\
-9103.65034794286	-0.00161951508258513\\
-9091.93016725675	-0.00199355611571542\\
-9080.20998657063	-0.00239811150803011\\
-9068.48980588451	-0.00283004623514711\\
-9056.76962519839	-0.00328597535268875\\
-9045.04944451227	-0.00376229125416531\\
-9027.46917348309	-0.0045066304573585\\
-9009.88890245391	-0.00527477136284887\\
-8992.30863142473	-0.00605266598490444\\
-8974.72836039555	-0.00682595985488937\\
-8957.14808936638	-0.00758026334523564\\
-8939.5678183372	-0.00830142605598395\\
-8916.12745696496	-0.00918795747381797\\
-8892.68709559272	-0.00996134920026928\\
-8869.24673422049	-0.0105950345094054\\
-8845.80637284825	-0.01106688918645\\
-8822.36601147601	-0.011360036198386\\
-8798.92565010377	-0.0114634822592087\\
-8775.48528873154	-0.0113725627245209\\
-8752.0449273593	-0.0110891774419085\\
-8728.60456598706	-0.0106218070311142\\
-8705.16420461482	-0.00998530635664211\\
-8681.72384324258	-0.00920047943980881\\
-8658.28348187034	-0.00829344749103283\\
-8640.70321084117	-0.00755125312717801\\
-8623.12293981199	-0.00677173350075279\\
-8605.54266878281	-0.00597000655580826\\
-8587.96239775363	-0.00516176310263483\\
-8570.38212672445	-0.00436297278406238\\
-8558.66194603833	-0.00384358253648326\\
-8546.94176535222	-0.00334014404606398\\
-8535.2215846661	-0.00285720894460144\\
-8523.50140397998	-0.00239918488619662\\
-8511.78122329386	-0.00197029725704231\\
-8500.06104260774	-0.0015745521810667\\
-8494.20095226468	-0.00139029131535057\\
-8488.34086192162	-0.00121570114241605\\
-8482.48077157856	-0.00105120532135851\\
-8476.6206812355	-0.000897207534121875\\
-8470.76059089244	-0.000754090566093032\\
-8464.90050054939	-0.000622215427785242\\
-8459.04041020633	-0.0005019205197564\\
-8453.18031986327	-0.000393520842531137\\
-8447.32022952021	-0.000297307253461387\\
-8441.46013917715	-0.000213545772104161\\
-8435.60004883409	-0.000142476935878959\\
-8429.73995849103	-8.4315207436601e-05\\
-8423.87986814797	-3.92484351895344e-05\\
-8418.01977780491	-7.43736831905056e-06\\
-8412.15968746185	1.09847725478301e-05\\
-8406.29959711879	1.59126679020703e-05\\
-8400.43950677573	7.26921418331051e-06\\
-8394.57941643267	-1.49942199766626e-05\\
-8388.71932608961	-5.08975968296321e-05\\
-8382.85923574655	-0.000100432072515783\\
-8376.99914540349	-0.000163559920488515\\
-8371.13905506043	-0.000240214515715671\\
-8365.27896471737	-0.000330300379891579\\
-8359.41887437431	-0.000433693287773499\\
-8353.55878403125	-0.00055024043461665\\
-8347.69869368819	-0.000679760664552227\\
-8341.83860334514	-0.000822044759626161\\
-8335.97851300208	-0.000976855789136341\\
-8330.11842265902	-0.00114392951863701\\
-8324.25833231596	-0.00132297487807245\\
-8318.3982419729	-0.00151367448811043\\
-8306.67806128678	-0.00192863895519843\\
-8294.95788060066	-0.00238578128066021\\
-8283.23769991454	-0.00288168198627955\\
-8271.51751922842	-0.0034125701953642\\
-8259.7973385423	-0.00397435322599193\\
-8248.07715785618	-0.00456264910827734\\
-8230.496886827	-0.00548460284504284\\
-8212.91661579783	-0.00643921019716877\\
-8195.33634476865	-0.00740916301652995\\
-8177.75607373947	-0.00837668572914441\\
-8160.17580271029	-0.0093238699436654\\
-8142.59553168111	-0.0102330150762208\\
-8119.15517030888	-0.0113565110135151\\
-8095.71480893664	-0.0123439177454911\\
-8072.2744475644	-0.0131612757384649\\
-8048.83408619216	-0.0137799505293344\\
-8025.39372481993	-0.0141776817014191\\
-8001.95336344769	-0.0143394308680123\\
-7978.51300207545	-0.0142579971199111\\
-7955.07264070321	-0.0139343757865282\\
-7931.63227933097	-0.0133778446711558\\
-7908.19191795873	-0.0126057708848127\\
-7884.7515565865	-0.0116431407279807\\
-7861.31119521426	-0.0105218244416379\\
-7843.73092418508	-0.00959923387898797\\
-7826.1506531559	-0.00862627023100858\\
-7808.57038212672	-0.00762184287113324\\
-7790.99011109754	-0.00660568626821062\\
-7773.40984006837	-0.00559799263855317\\
-7761.68965938225	-0.00494091672700519\\
-7749.96947869613	-0.00430256280080076\\
-7738.24929801001	-0.0036887645673474\\
-7726.52911732389	-0.00310519358549772\\
-7714.80893663777	-0.00255730994362976\\
-7703.08875595166	-0.00205031434601259\\
-7697.2286656086	-0.00181369247194271\\
-7691.36857526554	-0.00158910202191079\\
-7685.50848492248	-0.00137709997138264\\
-7679.64839457942	-0.00117821886042376\\
-7673.78830423636	-0.000992965571739724\\
-7667.9282138933	-0.000821820157997939\\
-7662.06812355024	-0.000665234721270361\\
-7656.20803320718	-0.000523632347051589\\
-7650.34794286412	-0.000397406095374902\\
-7644.48785252106	-0.000286918051360987\\
-7638.627762178	-0.000192498437387347\\
-7632.76767183494	-0.000114444789048414\\
-7626.90758149188	-5.30211968507886e-05\\
-7621.04749114882	-8.45761550854606e-06\\
-7615.18740080576	1.90507574305546e-05\\
-7609.3273104627	2.93440320654991e-05\\
-7603.46722011964	2.22981042770951e-05\\
-7597.60712977658	-2.17493976741762e-06\\
-7591.74703943353	-4.41264938878022e-05\\
-7585.88694909047	-0.000103571177926833\\
-7580.02685874741	-0.0001804866603871\\
-7574.16676840435	-0.000274813558354109\\
-7568.30667806129	-0.000386455415196423\\
-7562.44658771823	-0.000515278756450276\\
-7556.58649737517	-0.000661113224018044\\
-7550.72640703211	-0.000823751788772461\\
-7544.86631668905	-0.00100295104139083\\
-7539.00622634599	-0.00119843156113038\\
-7533.14613600293	-0.001409878362134\\
-7527.28604565987	-0.00163694141651033\\
-7521.42595531681	-0.00187923625356278\\
-7509.70577463069	-0.00240781529858957\\
-7497.98559394457	-0.00299187833827533\\
-7486.26541325845	-0.00362718786104831\\
-7474.54523257233	-0.00430903784531584\\
-7462.82505188622	-0.00503228982371356\\
-7451.1048712001	-0.00579141289927255\\
-7439.38469051398	-0.00658052740899902\\
-7421.8044194848	-0.00780684088478494\\
-7404.22414845562	-0.00906479804055633\\
-7386.64387742644	-0.0103316585680163\\
-7369.06360639727	-0.0115842350325273\\
-7351.48333536809	-0.0127993302170293\\
-7328.04297399585	-0.014321925398716\\
-7304.60261262361	-0.0156861747399205\\
-7281.16225125137	-0.0168452844263202\\
-7257.72188987914	-0.0177586807862829\\
-7234.2815285069	-0.0183934984765587\\
-7210.84116713466	-0.0187258335602542\\
-7187.40080576242	-0.0187417150808529\\
-7163.96044439018	-0.0184377574007755\\
-7140.52008301795	-0.017821465654814\\
-7117.07972164571	-0.016911177844771\\
-7093.63936027347	-0.0157356390353644\\
-7070.19899890123	-0.0143332153963796\\
-7052.61872787205	-0.0131604138175955\\
-7035.03845684288	-0.0119086551565541\\
-7017.4581858137	-0.0106022338939304\\
-6999.87791478452	-0.00926684942112058\\
-6982.29764375534	-0.00792913526700807\\
-6970.57746306922	-0.00704939798005758\\
-6958.8572823831	-0.00618868403251014\\
-6947.13710169699	-0.00535493635989402\\
-6935.41692101087	-0.00455594769407947\\
-6923.69674032475	-0.00379929311741169\\
-6911.97655963863	-0.00309226380890344\\
-6900.25637895251	-0.00244180255162248\\
-6894.39628860945	-0.00213984394061383\\
-6888.53619826639	-0.00185444155975373\\
-6882.67610792333	-0.00158633805257091\\
-6876.81601758027	-0.00133624317001102\\
-6870.95592723721	-0.00110483211031015\\
-6865.09583689415	-0.000892743923840465\\
-6859.23574655109	-0.000700579986113411\\
-6853.37565620803	-0.000528902542717993\\
-6847.51556586497	-0.000378233329560688\\
-6841.65547552191	-0.000249052271612713\\
-6835.79538517886	-0.000141796263310171\\
-6829.9352948358	-5.68580335656158e-05\\
-6824.07520449274	5.41490181981804e-06\\
-6818.21511414968	4.47211977396864e-05\\
-6812.35502380662	6.08065448811398e-05\\
-6806.49493346356	5.34644400932771e-05\\
-6800.6348431205	2.25368599307568e-05\\
-6794.77475277744	-3.20851645051866e-05\\
-6788.91466243438	-0.000110461073133609\\
-6783.05457209132	-0.000212600404349719\\
-6777.19448174826	-0.000338462525666669\\
-6771.3343914052	-0.00048795646860541\\
-6765.47430106214	-0.000660940868717434\\
-6759.61421071908	-0.000857224011372839\\
-6753.75412037602	-0.00107656398377321\\
-6747.89403003296	-0.00131866893334073\\
-6742.0339396899	-0.00158319743259294\\
-6736.17384934684	-0.00186975895010211\\
-6730.31375900378	-0.00217791442732509\\
-6724.45366866072	-0.00250717696029412\\
-6718.59357831766	-0.0028570125857194\\
-6706.87339763155	-0.00361603739947661\\
-6695.15321694543	-0.00444981228665629\\
-6683.43303625931	-0.00535247504330896\\
-6671.71285557319	-0.00631752045353124\\
-6659.99267488707	-0.007337848763703\\
-6648.27249420095	-0.00840581967497942\\
-6630.69222317177	-0.0100792236035522\\
-6613.1119521426	-0.0118123505099848\\
-6595.53168111342	-0.013574512043281\\
-6577.95141008424	-0.0153339935881119\\
-6560.37113905506	-0.0170586438051054\\
-6542.79086802588	-0.0187164819759805\\
-6519.35050665365	-0.0207692729987549\\
-6495.91014528141	-0.0225781695175015\\
-6472.46978390917	-0.0240802241843144\\
-6449.02942253693	-0.0252216564092075\\
-6425.58906116469	-0.0259598893565532\\
-6402.14869979245	-0.0262652540818828\\
-6378.70833842022	-0.0261222960899504\\
-6355.26797704798	-0.0255306314425216\\
-6331.82761567574	-0.024505313412591\\
-6308.3872543035	-0.0230766861454268\\
-6284.94689293126	-0.021289718372482\\
-6267.36662190209	-0.0197490117059137\\
-6249.78635087291	-0.0180684446572701\\
-6232.20607984373	-0.0162803607417812\\
-6214.62580881455	-0.0144198008258561\\
-6197.04553778538	-0.0125238849122458\\
-6179.4652667562	-0.0106311431676606\\
-6167.74508607008	-0.00939044587901222\\
-6156.02490538396	-0.0081802436172531\\
-6144.30472469784	-0.00701207879060207\\
-6132.58454401172	-0.00589729026377876\\
-6120.8643633256	-0.00484691496030564\\
-6109.14418263948	-0.00387159075676286\\
-6103.28409229642	-0.00341526460706999\\
-6097.42400195336	-0.00298146149550286\\
-6091.5639116103	-0.0025713571050529\\
-6085.70382126725	-0.00218608493357806\\
-6079.84373092419	-0.00182673363421009\\
-6073.98364058113	-0.00149434443234112\\
-6068.12355023807	-0.00118990862541044\\
-6062.26345989501	-0.000914365170761012\\
-6056.40336955195	-0.000668598367619746\\
-6050.54327920889	-0.000453435638175552\\
-6044.68318886583	-0.000269645413337205\\
-6038.82309852277	-0.000117935128014562\\
-6032.96300817971	1.05066909580283e-06\\
-6027.10291783665	8.67320865952666e-05\\
-6021.24282749359	0.000138595537882602\\
-6015.38273715053	0.000156195260852918\\
-6009.52264680747	0.00013915470579904\\
-6003.66255646441	8.71677877848544e-05\\
nan	nan\\
-5991.94237577829	-0.000122510615126645\\
-5986.08228543523	-0.000280452640003633\\
-5980.22219509217	-0.000473840603818683\\
-5974.36210474911	-0.000702614454714258\\
-5968.50201440605	-0.000966639184324193\\
-5962.64192406299	-0.00126570460641393\\
-5956.78183371994	-0.00159952529117195\\
-5950.92174337688	-0.00196774065625192\\
-5945.06165303382	-0.00236991521549063\\
-5939.20156269076	-0.0028055389858606\\
-5933.3414723477	-0.00327402805282989\\
-5927.48138200464	-0.0037747252939739\\
-5921.62129166158	-0.00430690126063523\\
-5909.90111097546	-0.00546241633302454\\
-5898.18093028934	-0.00673333451729139\\
-5886.46074960322	-0.00811134241370701\\
-5874.7405689171	-0.00958710721534504\\
-5863.02038823099	-0.0111503412637656\\
-5851.30020754487	-0.012789875514294\\
-5833.71993651569	-0.0153658619123142\\
-5816.13966548651	-0.0180431423637368\\
-5798.55939445733	-0.0207755936994355\\
-5780.97912342815	-0.0235150732477431\\
-5763.39885239898	-0.0262122929199842\\
-5739.95849102674	-0.0296573449622676\\
-5716.5181296545	-0.0328243405983894\\
-5693.07776828226	-0.0356045307153503\\
-5669.63740691002	-0.037899288185729\\
-5646.19704553779	-0.0396237834005381\\
-5622.75668416555	-0.0407103125403666\\
-5599.31632279331	-0.0411111576362818\\
-5575.87596142107	-0.0408008691843478\\
-5552.43560004883	-0.0397778776368927\\
-5528.9952386766	-0.0380653591668352\\
-5505.55487730436	-0.0357113031012466\\
-5482.11451593212	-0.0327877528011607\\
-5464.53424490294	-0.0302774310715828\\
-5446.95397387376	-0.0275469352793039\\
-5429.37370284459	-0.0246493262921954\\
-5411.79343181541	-0.0216424565689595\\
-5394.21316078623	-0.0185880071358758\\
-5382.49298010011	-0.016557021200611\\
-5370.77279941399	-0.0145530461144421\\
-5359.05261872787	-0.0125957939561263\\
-5347.33243804175	-0.0107049802931535\\
-5335.61225735563	-0.00890015948248691\\
-5323.89207666952	-0.00720055824365951\\
-5318.03198632646	-0.00639609265514323\\
-5312.1718959834	-0.00562490882429211\\
-5306.31180564034	-0.0048892437984047\\
-5300.45171529728	-0.00419128309699356\\
-5294.59162495422	-0.00353315567387322\\
-5288.73153461116	-0.00291692894699858\\
-5282.8714442681	-0.00234460390672503\\
-5277.01135392504	-0.00181811031274662\\
-5271.15126358198	-0.00133930199060966\\
-5265.29117323892	-0.000909952237710022\\
-5259.43108289586	-0.000531749349427142\\
-5253.5709925528	-0.00020629227510786\\
-5247.71090220974	6.4913585937501e-05\\
-5241.85081186668	0.000280460438792809\\
-5235.99072152362	0.000439041988111794\\
-5230.13063118056	0.000539457164811521\\
-5224.2705408375	0.000580613675216832\\
-5218.41045049444	0.000561531363910543\\
-5212.55036015138	0.000481345381755638\\
-5206.69026980832	0.000339309150907825\\
-5200.83017946527	0.000134797118939733\\
-5194.97008912221	-0.000132692705418172\\
-5189.10999877915	-0.000463536442155097\\
-5183.24990843609	-0.000857982262230763\\
-5177.38981809303	-0.00131614854860384\\
-5171.52972774997	-0.00183802230208667\\
-5165.66963740691	-0.00242345779756414\\
-5159.80954706385	-0.00307217549572846\\
-5153.94945672079	-0.00378376121516422\\
-5148.08936637773	-0.00455766556858167\\
-5142.22927603467	-0.0053932036677192\\
-5136.36918569161	-0.00628955509950849\\
-5130.50909534855	-0.00724576417652113\\
-5124.64900500549	-0.00826074046401502\\
-5112.92882431937	-0.0104619643084524\\
-5101.20864363326	-0.012881845815389\\
-5089.48846294714	-0.0155069288851972\\
-5077.76828226102	-0.0183217482260002\\
-5066.0481015749	-0.0213089132012892\\
-5054.32792088878	-0.0244492091518728\\
-5036.7476498596	-0.0294005769326532\\
-5019.16737883043	-0.0345719762746803\\
-5001.58710780125	-0.0398795686526454\\
-4984.00683677207	-0.0452340866088032\\
-4966.42656574289	-0.0505423156890993\\
-4942.98620437065	-0.0573826692757273\\
-4919.54584299841	-0.0637433610438042\\
-4896.10548162617	-0.0694025438860658\\
-4872.66512025394	-0.0741516930284429\\
-4849.2247588817	-0.0778033315588932\\
-4825.78439750946	-0.0801984987493339\\
-4802.34403613722	-0.0812137232874689\\
-4778.90367476498	-0.080767268888778\\
-4755.46331339275	-0.078824432068685\\
-4732.02295202051	-0.0754016910161567\\
-4708.58259064827	-0.0705695301285212\\
-4691.0023196191	-0.0660940975749839\\
-4673.42204858992	-0.0609690344120997\\
-4655.84177756074	-0.0552849484671149\\
-4638.26150653156	-0.0491488528904369\\
-4620.68123550238	-0.0426831816400012\\
-4608.96105481626	-0.038256546367968\\
-4597.24087413014	-0.0337875244621813\\
-4585.52069344402	-0.0293220225675828\\
-4573.80051275791	-0.0249079484284059\\
-4562.08033207179	-0.020594933383191\\
-4556.22024172873	-0.0184922313870945\\
-4550.36015138567	-0.0164340353348914\\
-4544.50006104261	-0.0144268986882795\\
-4538.63997069955	-0.0124774238011881\\
-4532.77988035649	-0.0105922512223704\\
-4526.91979001343	-0.0087780487888698\\
-4521.05969967037	-0.00704150052832416\\
-4515.19960932731	-0.00538929538585851\\
-4509.33951898425	-0.00382811579186758\\
-4503.47942864119	-0.00236462608924002\\
-4497.61933829813	-0.0010054608368811\\
-4491.75924795507	0.000242786992088656\\
-4485.89915761201	0.00137357789895909\\
-4480.03906726895	0.00238043781621532\\
-4474.17897692589	0.00325697011347182\\
-4468.31888658283	0.00399686763332092\\
-4462.45879623977	0.00459392473771451\\
-4456.59870589671	0.00504204934483239\\
-4450.73861555366	0.0053352749374476\\
-4444.8785252106	0.00546777252190704\\
-4439.01843486754	0.00543386251869322\\
-4433.15834452448	0.00522802656357129\\
-4427.29825418142	0.00484491920031939\\
-4421.43816383836	0.00427937944383579\\
-4415.5780734953	0.00352644219481869\\
-4409.71798315224	0.00258134948502271\\
-4403.85789280918	0.00143956153411115\\
-4397.99780246612	9.67675980635635e-05\\
-4392.13771212306	-0.00145110341038938\\
-4386.27762178	-0.00320787254820711\\
-4380.41753143694	-0.00517710073465855\\
-4374.55744109388	-0.00736207892878549\\
-4368.69735075082	-0.00976581860991315\\
-4362.83726040776	-0.0123910425795004\\
-4356.97717006471	-0.0152401761018386\\
-4351.11707972165	-0.0183153384010472\\
-4345.25698937859	-0.0216183345308175\\
-4339.39689903553	-0.0251506476336813\\
-4333.53680869247	-0.0289134316054568\\
-4321.81662800635	-0.0371333404505316\\
-4310.09644732023	-0.0462804555158956\\
-4298.37626663411	-0.0563515072221691\\
-4286.65608594799	-0.067337407025305\\
-4274.93590526187	-0.0792231244319007\\
-4263.21572457576	-0.0919875972916452\\
-4245.63545354658	-0.112720914432465\\
-4228.0551825174	-0.13525153999808\\
-4210.47491148822	-0.159431415498824\\
-4192.89464045904	-0.185082409880712\\
-4169.4542790868	-0.221210061458708\\
-4146.01391771456	-0.259015019513751\\
-4122.57355634233	-0.297851811219159\\
-4099.13319497009	-0.337006267618582\\
-4075.69283359785	-0.375707748956619\\
-4046.39238188255	-0.42220602590329\\
-4017.09193016726	-0.465080906397207\\
-3987.79147845196	-0.50267429382598\\
-3958.49102673666	-0.533366552555596\\
-3929.19057502137	-0.555630751400453\\
-3899.89012330607	-0.568085780573532\\
-3870.58967159077	-0.569546432736429\\
-3841.28921987547	-0.559068617544459\\
-3817.84885850324	-0.541634026479852\\
-3794.408497131	-0.515929512003144\\
-3770.96813575876	-0.481891659052772\\
-3753.38786472958	-0.45093497498114\\
-3735.8075937004	-0.415415037318013\\
-3718.22732267122	-0.375460943051072\\
-3706.5071419851	-0.346445221558062\\
-3694.78696129899	-0.315596884907732\\
-3683.06678061287	-0.282986257974658\\
-3671.34659992675	-0.248691735036614\\
-3659.62641924063	-0.212799498214104\\
-3653.76632889757	-0.194283220345604\\
-3647.90623855451	-0.175403196746636\\
-3642.04614821145	-0.156172241528107\\
-3636.18605786839	-0.136603588761247\\
-3630.32596752533	-0.116710879352488\\
-3624.46587718227	-0.096508147431388\\
-3618.60578683921	-0.0760098062416358\\
-3612.74569649615	-0.0552306335631772\\
-3606.88560615309	-0.0341857566787052\\
-3601.02551581004	-0.012890636917385\\
-3595.16542546698	0.00863894622100714\\
-3589.30533512392	0.0303869113242634\\
-3583.44524478086	0.0523368916980917\\
-3577.5851544378	0.0744722523392822\\
-3571.72506409474	0.0967761072337489\\
-3565.86497375168	0.119231336952416\\
-3560.00488340862	0.141820606527005\\
-3554.14479306556	0.164526383572827\\
-3548.2847027225	0.187330956644857\\
-3542.42461237944	0.210216453800139\\
-3536.56452203638	0.233164861330508\\
-3524.84434135026	0.279177757358764\\
-3513.12416066415	0.325223420805006\\
-3501.40397997803	0.371154580007798\\
-3489.68379929191	0.416823532001606\\
-3472.10352826273	0.484512828700304\\
-3454.52325723355	0.550786572764059\\
-3436.94298620437	0.615159537146415\\
-3413.50262483213	0.697219797801138\\
-3390.06226345989	0.773994468866484\\
-3366.62190208766	0.844486334270942\\
-3337.32145037236	0.922389205523877\\
-3308.02099865706	0.987508530355332\\
-3278.72054694176	1.03857193114877\\
-3249.42009522647	1.07461235805407\\
-3220.11964351117	1.09498419145606\\
-3190.81919179587	1.09936990030888\\
-3155.65864973752	1.08356456399332\\
-3126.35819802222	1.05324975583795\\
-3097.05774630692	1.00805215848877\\
-3067.75729459163	0.948907070119819\\
-3038.45684287633	0.876979340563779\\
-3015.01648150409	0.811140388467105\\
-2991.57612013185	0.738753798811626\\
-2968.13575875961	0.660636423108755\\
-2950.55548773043	0.598799696405538\\
-2932.97521670126	0.534587711230066\\
-2915.39494567208	0.468373785426526\\
-2903.67476498596	0.423304094248923\\
-2891.95458429984	0.377622132326263\\
-2880.23440361372	0.331438200414991\\
-2868.5142229276	0.284861713767814\\
-2856.79404224148	0.238000904769999\\
-2850.93395189842	0.214497309746409\\
-2845.07386155537	0.190962544938027\\
-2839.21377121231	0.167409686196695\\
-2833.35368086925	0.143851686950385\\
-2827.49359052619	0.120301371266129\\
-2821.63350018313	0.096771427256034\\
-2815.77340984007	0.0732744008047981\\
-2809.91331949701	0.0498226896368607\\
-2804.05322915395	0.0264285377076926\\
-2798.19313881089	0.00310402993171366\\
-2792.33304846783	-0.0201389127631252\\
-2786.47295812477	-0.0432885380521901\\
-2780.61286778171	-0.0663332665046747\\
-2774.75277743865	-0.0892616955546188\\
-2768.89268709559	-0.112062603149251\\
-2763.03259675254	-0.134724951087915\\
-2757.17250640948	-0.15723788804906\\
-2751.31241606642	-0.179590752309485\\
-2745.45232572336	-0.201773074171616\\
-2739.5922353803	-0.223774578095359\\
-2727.87205469418	-0.267195011551427\\
-2716.15187400806	-0.309773794772448\\
-2704.43169332194	-0.351435866303824\\
-2692.71151263582	-0.392109382374847\\
-2675.13124160664	-0.451116888696373\\
-2657.55097057746	-0.507528458285575\\
-2639.97069954828	-0.561142139615882\\
-2616.53033817605	-0.627955264149058\\
-2593.08997680381	-0.689066371958402\\
-2569.64961543157	-0.74412334775921\\
-2540.34916371627	-0.803966438612835\\
-2511.04871200098	-0.853364159186347\\
-2481.74826028568	-0.891905295019938\\
-2446.58771822732	-0.923380677160696\\
-2411.42717616897	-0.93838352068375\\
-2376.26663411061	-0.936710277289482\\
-2341.10609205225	-0.918340945436338\\
-2305.9455499939	-0.883448279689694\\
-2276.6450982786	-0.842016394828914\\
-2247.3446465633	-0.789707646561303\\
-2218.044194848	-0.726975632300107\\
-2194.60383347576	-0.669668258444039\\
-2171.16347210353	-0.606410453249348\\
-2147.72311073129	-0.537610415757817\\
-2130.14283970211	-0.482649768977804\\
-2112.56256867293	-0.425046521524711\\
-2100.84238798681	-0.385290474439117\\
-2089.1222073007	-0.344536139459104\\
-2077.40202661458	-0.302859875357715\\
-2065.68184592846	-0.260341198076615\\
-2059.8217555854	-0.238791641223069\\
-2053.96166524234	-0.217062676501985\\
-2048.10157489928	-0.195165082091855\\
-2042.24148455622	-0.173109810811487\\
-2036.38139421316	-0.150907985369714\\
-2030.5213038701	-0.128570893332231\\
-2024.66121352704	-0.106109981819477\\
-2018.80112318398	-0.0835368519419716\\
-2012.94103284092	-0.0608632529603804\\
-2007.08094249787	-0.0381010761922624\\
-2001.22085215481	-0.0152623486548387\\
-1995.36076181175	0.00764077354533821\\
-1989.50067146869	0.0305960120719594\\
-1983.64058112563	0.0535909734650745\\
-1977.78049078257	0.0766131564837794\\
-1971.92040043951	0.0996499596837526\\
-1966.06031009645	0.122688689220735\\
-1960.20021975339	0.145716566871446\\
-1954.34012941033	0.168720738271243\\
-1948.48003906727	0.191688281346292\\
-1942.61994872421	0.214606214952346\\
-1936.75985838115	0.237461507690922\\
-1925.03967769503	0.282931847848651\\
-1913.31949700891	0.327994391559696\\
-1901.59931632279	0.372544016865965\\
-1889.87913563667	0.416475699368519\\
-1872.2988646075	0.480985937887762\\
-1854.71859357832	0.543520996748328\\
-1837.13832254914	0.603737485207334\\
-1813.6979611769	0.67984259998767\\
-1790.25759980466	0.750471012333545\\
-1766.81723843242	0.814907356671754\\
-1737.51678671713	0.885746226967516\\
-1708.21633500183	0.944740483360737\\
-1678.91588328653	0.9909303533315\\
-1649.61543157124	1.0235611680713\\
-1614.45488951288	1.04408629989649\\
-1579.29434745452	1.04385550840326\\
-1544.13380539617	1.02288928178193\\
-1514.83335368087	0.989933789663105\\
-1485.53290196557	0.943476070394558\\
-1456.23245025027	0.884279955332726\\
-1426.93199853498	0.813305426779727\\
-1403.49163716274	0.748802068986713\\
-1380.0512757905	0.678133481277436\\
-1356.61091441826	0.602001650428078\\
-1339.03064338909	0.541763999266042\\
-1321.45037235991	0.479197912944355\\
-1303.87010133073	0.414638066700273\\
-1292.14992064461	0.370658497079378\\
-1280.42973995849	0.326045931480171\\
-1268.70955927237	0.280902656332863\\
-1256.98937858626	0.235331102097145\\
-1251.1292882432	0.212416702411962\\
-1245.26919790014	0.189433572519851\\
-1239.40910755708	0.16639443766229\\
-1233.54901721402	0.143311982992198\\
-1227.68892687096	0.120198845958366\\
-1221.8288365279	0.0970676088605646\\
-1215.96874618484	0.0739307915839076\\
-1210.10865584178	0.0508008445204743\\
-1204.24856549872	0.0276901416734756\\
-1198.38847515566	0.00461097396102351\\
-1192.5283848126	-0.0184244572892841\\
-1186.66829446954	-0.0414040466446187\\
-1180.80820412648	-0.0643157906693696\\
-1174.94811378342	-0.0871477938224896\\
-1169.08802344036	-0.109888274150354\\
-1163.2279330973	-0.132525568771261\\
-1157.36784275424	-0.155048139153195\\
-1151.50775241118	-0.177444576173161\\
-1145.64766206812	-0.199703604986076\\
-1139.78757172506	-0.221814089656128\\
-1133.927481382	-0.243765037614798\\
-1122.20730069589	-0.287145095045285\\
-1110.48712000977	-0.329758859409109\\
-1098.76693932365	-0.371523955943757\\
-1087.04675863753	-0.412360659128184\\
-1069.46648760835	-0.471707347331137\\
-1051.88621657917	-0.528544775894659\\
-1034.30594554999	-0.58263991915285\\
-1010.86558417775	-0.650126723400459\\
-987.425222805519	-0.711879422613491\\
-963.98486143328	-0.767474173442195\\
-934.684409717982	-0.827743097312675\\
-905.383958002687	-0.877188857761063\\
-876.083506287388	-0.915311652346762\\
-840.92296422903	-0.945585980803626\\
-805.762422170676	-0.958582884712484\\
-770.601880112317	-0.954115110802406\\
-735.441338053963	-0.932239374418993\\
-700.280795995604	-0.893258152974448\\
-670.980344280306	-0.848097540576334\\
-641.679892565011	-0.79189033012327\\
-612.379440849712	-0.725214044692156\\
-588.939079477474	-0.664803222484261\\
-565.498718105238	-0.598556933139792\\
-547.918447076059	-0.545322583025903\\
-530.33817604688	-0.489284402511263\\
-512.757905017701	-0.4306740808352\\
-501.037724331585	-0.390292825658711\\
-489.317543645466	-0.348954505945881\\
-477.597362959346	-0.30673863218497\\
-465.877182273227	-0.263727393015123\\
-460.017091930167	-0.24194989970827\\
-454.157001587108	-0.220005540588357\\
-448.296911244048	-0.197905289303377\\
-442.436820900988	-0.175660263084189\\
-436.576730557928	-0.153281718031439\\
-430.716640214869	-0.130781044220187\\
-424.856549871809	-0.108169760600926\\
-418.996459528749	-0.0854595096805506\\
-413.136369185693	-0.0626620520201019\\
-407.276278842633	-0.0397892605221325\\
-401.416188499574	-0.0168531145238133\\
-395.556098156514	0.0061343063085423\\
-389.696007813454	0.0291608282749697\\
-383.835917470395	0.0522141901281179\\
-377.975827127335	0.075282049745088\\
-372.115736784275	0.0983519909824909\\
-366.255646441216	0.121411530709818\\
-360.395556098156	0.144448126009068\\
-354.535465755096	0.167449181552656\\
-348.675375412036	0.190402057137492\\
-342.815285068977	0.213294075374845\\
-336.955194725917	0.236112529531225\\
-325.235014039801	0.281477819998943\\
-313.514833353682	0.326395994598654\\
-301.794652667562	0.370765174506946\\
-290.074471981443	0.414483809178651\\
-272.494200952264	0.478621412341787\\
-254.913929923085	0.540731761689403\\
-237.333658893909	0.600484956266981\\
-213.89329752167	0.675935994324994\\
-190.452936149431	0.745898213831991\\
-167.012574777196	0.809687034745988\\
-137.712123061898	0.879785753679549\\
-108.411671346599	0.938161383224601\\
-79.1112196313043	0.983892347470839\\
-49.8107679160057	1.01625076655098\\
-14.6502258576475	1.03671504386981\\
20.5103162007072	1.03671504386981\\
55.6708582590654	1.01625076655098\\
84.9713099743603	0.983892347470839\\
114.271761689659	0.938161383224601\\
143.572213404957	0.879785753679549\\
172.872665120252	0.809687034745988\\
196.313026492491	0.745898213831991\\
219.75338786473	0.675935994324994\\
243.193749236969	0.600484956266981\\
260.774020266144	0.540731761689403\\
278.354291295323	0.478621412341787\\
295.934562324503	0.414483809178651\\
307.654743010622	0.370765174506946\\
319.374923696741	0.326395994598654\\
331.095104382861	0.281477819998943\\
342.815285068977	0.236112529531225\\
348.675375412036	0.213294075374845\\
354.535465755096	0.190402057137492\\
360.395556098156	0.167449181552656\\
366.255646441216	0.144448126009068\\
372.115736784275	0.121411530709818\\
377.975827127335	0.0983519909824909\\
383.835917470395	0.075282049745088\\
389.696007813454	0.0522141901281179\\
395.556098156514	0.0291608282749697\\
401.416188499574	0.0061343063085423\\
407.276278842633	-0.0168531145238133\\
413.136369185693	-0.0397892605221325\\
418.996459528749	-0.0626620520201019\\
424.856549871809	-0.0854595096805506\\
430.716640214869	-0.108169760600926\\
436.576730557928	-0.130781044220187\\
442.436820900988	-0.153281718031439\\
448.296911244048	-0.175660263084189\\
454.157001587108	-0.197905289303377\\
460.017091930167	-0.220005540588357\\
465.877182273227	-0.24194989970827\\
477.597362959346	-0.285327194918679\\
489.317543645466	-0.327951188018543\\
501.037724331585	-0.369738393487153\\
512.757905017701	-0.410607948229538\\
530.33817604688	-0.470021541141903\\
547.918447076059	-0.526943600588813\\
565.498718105238	-0.581137353262501\\
588.939079477474	-0.648771293214524\\
612.379440849712	-0.710682120815374\\
635.819802221951	-0.766439567974595\\
665.12025393725	-0.826910699068304\\
694.420705652545	-0.876551286936023\\
723.721157367843	-0.914858358479798\\
758.881699426198	-0.945340126926624\\
794.042241484556	-0.95853426742479\\
829.202783542914	-0.954261524511677\\
864.363325601269	-0.932587269953778\\
899.523867659627	-0.89382068239106\\
928.824319374922	-0.848850867708293\\
958.124771090221	-0.792843515319611\\
987.425222805519	-0.726370878342324\\
1010.86558417775	-0.666119878518875\\
1034.30594554999	-0.600024143061596\\
1051.88621657917	-0.546892261632748\\
1069.46648760835	-0.490943779711676\\
1087.04675863753	-0.432406696698731\\
1098.76693932365	-0.39206323341799\\
1110.48712000977	-0.350752537250848\\
1122.20730069589	-0.308552973193088\\
1133.927481382	-0.265545603869687\\
1139.78757172506	-0.243765037614798\\
1145.64766206812	-0.221814089656128\\
1151.50775241118	-0.199703604986076\\
1157.36784275424	-0.177444576173161\\
1163.2279330973	-0.155048139153195\\
1169.08802344036	-0.132525568771261\\
1174.94811378342	-0.109888274150354\\
1180.80820412648	-0.0871477938224896\\
1186.66829446954	-0.0643157906693696\\
1192.5283848126	-0.0414040466446187\\
1198.38847515566	-0.0184244572892841\\
1204.24856549872	0.00461097396102351\\
1210.10865584178	0.0276901416734756\\
1215.96874618484	0.0508008445204743\\
1221.8288365279	0.0739307915839076\\
1227.68892687096	0.0970676088605646\\
1233.54901721402	0.120198845958366\\
1239.40910755708	0.143311982992198\\
1245.26919790014	0.16639443766229\\
1251.1292882432	0.189433572519851\\
1256.98937858626	0.212416702411962\\
1262.84946892931	0.235331102097145\\
1274.56964961543	0.280902656332863\\
1286.28983030155	0.326045931480171\\
1298.01001098767	0.370658497079378\\
1309.73019167379	0.414638066700273\\
1327.31046270297	0.479197912944355\\
1344.89073373214	0.541763999266042\\
1362.47100476132	0.602001650428078\\
1385.91136613356	0.678133481277436\\
1409.3517275058	0.748802068986713\\
1432.79208887804	0.813305426779727\\
1462.09254059333	0.884279955332726\\
1491.39299230863	0.943476070394558\\
1520.69344402393	0.989933789663105\\
1549.99389573923	1.02288928178193\\
1585.15443779758	1.04385550840326\\
1620.31497985594	1.04408629989649\\
1655.4755219143	1.0235611680713\\
1684.77597362959	0.9909303533315\\
1714.07642534489	0.944740483360737\\
1743.37687706019	0.885746226967516\\
1772.67732877548	0.814907356671754\\
1796.11769014772	0.750471012333545\\
1819.55805151996	0.67984259998767\\
1842.9984128922	0.603737485207334\\
1860.57868392138	0.543520996748328\\
1878.15895495056	0.480985937887762\\
1895.73922597973	0.416475699368519\\
1907.45940666585	0.372544016865965\\
1919.17958735197	0.327994391559696\\
1930.89976803809	0.282931847848651\\
1942.61994872421	0.237461507690922\\
1948.48003906727	0.214606214952346\\
1954.34012941033	0.191688281346292\\
1960.20021975339	0.168720738271243\\
1966.06031009645	0.145716566871446\\
1971.92040043951	0.122688689220735\\
1977.78049078257	0.0996499596837526\\
1983.64058112563	0.0766131564837794\\
1989.50067146869	0.0535909734650745\\
1995.36076181175	0.0305960120719594\\
2001.22085215481	0.00764077354533821\\
2007.08094249787	-0.0152623486548387\\
2012.94103284092	-0.0381010761922624\\
2018.80112318398	-0.0608632529603804\\
2024.66121352704	-0.0835368519419716\\
2030.5213038701	-0.106109981819477\\
2036.38139421316	-0.128570893332231\\
2042.24148455622	-0.150907985369714\\
2048.10157489928	-0.173109810811487\\
2053.96166524234	-0.195165082091855\\
2059.8217555854	-0.217062676501985\\
2065.68184592846	-0.238791641223069\\
2077.40202661458	-0.281700748023349\\
2089.1222073007	-0.323808351663822\\
2100.84238798681	-0.365033395610847\\
2112.56256867293	-0.405297930566587\\
2130.14283970211	-0.46373116616562\\
2147.72311073129	-0.519597485589205\\
2165.30338176047	-0.572675563010941\\
2188.74374313271	-0.638758927273433\\
2212.18410450494	-0.699089135234759\\
2235.62446587718	-0.753285171235446\\
2264.92491759248	-0.811910015576784\\
2294.22536930778	-0.859915808117668\\
2329.38591136613	-0.902941341497078\\
2364.54645342449	-0.92960648866904\\
2399.70699548285	-0.939637041934716\\
2434.8675375412	-0.932958121884436\\
2470.02807959956	-0.909682291547193\\
2505.18862165792	-0.870103296113017\\
2534.48907337321	-0.825004112213072\\
2563.78952508851	-0.769283741039823\\
2593.08997680381	-0.703410949366457\\
2616.53033817605	-0.643782141835267\\
2639.97069954828	-0.578358991479302\\
2657.55097057746	-0.525720076236175\\
2675.13124160664	-0.470219484359756\\
2692.71151263582	-0.41205386256751\\
2704.43169332194	-0.371900558051928\\
2716.15187400806	-0.330723985293277\\
2727.87205469418	-0.2885943760248\\
2739.5922353803	-0.24558518454382\\
2745.45232572336	-0.223774578095359\\
2751.31241606642	-0.201773074171616\\
2757.17250640948	-0.179590752309485\\
2763.03259675254	-0.15723788804906\\
2768.89268709559	-0.134724951087915\\
2774.75277743865	-0.112062603149251\\
2780.61286778171	-0.0892616955546188\\
2786.47295812477	-0.0663332665046747\\
2792.33304846783	-0.0432885380521901\\
2798.19313881089	-0.0201389127631252\\
2804.05322915395	0.00310402993171366\\
2809.91331949701	0.0264285377076926\\
2815.77340984007	0.0498226896368607\\
2821.63350018313	0.0732744008047981\\
2827.49359052619	0.096771427256034\\
2833.35368086925	0.120301371266129\\
2839.21377121231	0.143851686950385\\
2845.07386155537	0.167409686196695\\
2850.93395189842	0.190962544938027\\
2856.79404224148	0.214497309746409\\
2862.65413258454	0.238000904769999\\
2874.37431327066	0.284861713767814\\
2886.09449395678	0.331438200414991\\
2897.8146746429	0.377622132326263\\
2909.53485532902	0.423304094248923\\
2927.1151263582	0.490644395715934\\
2944.69539738738	0.556232637070602\\
2962.27566841655	0.619694194736366\\
2985.71602978879	0.70035884027241\\
3009.15639116103	0.775715496874331\\
3032.59675253327	0.84492759849302\\
3061.89720424857	0.921597776224929\\
3091.19765596386	0.986007350486606\\
3120.49810767916	1.03690982273399\\
3149.79855939446	1.07327386384239\\
3179.09901110976	1.09431803521871\\
3214.25955316811	1.09866357258508\\
3243.56000488341	1.08466646787492\\
3272.8604565987	1.05483432904991\\
3302.160908314	1.00968181278724\\
3331.4613600293	0.950048800228693\\
3360.7618117446	0.877088071806481\\
3384.20217311684	0.810084836706592\\
3407.64253448907	0.736332031887184\\
3431.08289586131	0.656785040048408\\
3448.66316689049	0.593942541177733\\
3466.24343791967	0.528882787973249\\
3483.82370894885	0.462082759246009\\
3495.54388963497	0.416823532001606\\
3507.26407032109	0.371154580007798\\
3518.9842510072	0.325223420805006\\
3530.70443169332	0.279177757358764\\
3542.42461237944	0.233164861330508\\
3548.2847027225	0.210216453800139\\
3554.14479306556	0.187330956644857\\
3560.00488340862	0.164526383572827\\
3565.86497375168	0.141820606527005\\
3571.72506409474	0.119231336952416\\
3577.5851544378	0.0967761072337489\\
3583.44524478086	0.0744722523392822\\
3589.30533512392	0.0523368916980917\\
3595.16542546698	0.0303869113242634\\
3601.02551581004	0.00863894622100714\\
3606.88560615309	-0.012890636917385\\
3612.74569649615	-0.0341857566787052\\
3618.60578683921	-0.0552306335631772\\
3624.46587718227	-0.0760098062416358\\
3630.32596752533	-0.096508147431388\\
3636.18605786839	-0.116710879352488\\
3642.04614821145	-0.136603588761247\\
3647.90623855451	-0.156172241528107\\
3653.76632889757	-0.175403196746636\\
3659.62641924063	-0.194283220345604\\
3665.48650958369	-0.212799498214104\\
3677.20669026981	-0.248691735036614\\
3688.92687095593	-0.282986257974658\\
3700.64705164204	-0.315596884907732\\
3712.36723232816	-0.346445221558062\\
3729.94750335734	-0.389261863212917\\
3747.52777438652	-0.427754991910525\\
3765.1080454157	-0.461765987232056\\
3782.68831644488	-0.491180624234413\\
3806.12867781712	-0.523136771096618\\
3829.56903918935	-0.546762493786754\\
3853.00940056159	-0.56215544655377\\
3882.30985227689	-0.570182974569313\\
3911.61030399219	-0.566439601074606\\
3940.91075570748	-0.551920965364288\\
3970.21120742278	-0.527855253501815\\
3999.51165913808	-0.495656870274581\\
4028.81211085337	-0.456875600906461\\
4058.11256256867	-0.413143114833487\\
4087.41301428397	-0.366118727813539\\
4110.85337565621	-0.327229177393526\\
4134.29373702845	-0.28808386933809\\
4157.73409840068	-0.249440571688767\\
4181.17445977292	-0.211997351839675\\
4204.61482114516	-0.176380825722701\\
4222.19509217434	-0.151198593510241\\
4239.77536320352	-0.12755014350599\\
4257.3556342327	-0.10560367653161\\
4274.93590526187	-0.0854970701767792\\
4286.65608594799	-0.0731689627272707\\
4298.37626663411	-0.0617309000364902\\
4310.09644732023	-0.0512009193400539\\
4321.81662800635	-0.0415910668354822\\
4333.53680869247	-0.0329075041805321\\
4339.39689903553	-0.0289134316054568\\
4345.25698937859	-0.0251506476336813\\
4351.11707972165	-0.0216183345308175\\
4356.97717006471	-0.0183153384010472\\
4362.83726040776	-0.0152401761018386\\
4368.69735075082	-0.0123910425795004\\
4374.55744109388	-0.00976581860991315\\
4380.41753143694	-0.00736207892878549\\
4386.27762178	-0.00517710073465855\\
4392.13771212306	-0.00320787254820711\\
4397.99780246612	-0.00145110341038938\\
4403.85789280918	9.67675980635635e-05\\
4409.71798315224	0.00143956153411115\\
4415.5780734953	0.00258134948502271\\
4421.43816383836	0.00352644219481869\\
4427.29825418142	0.00427937944383579\\
4433.15834452448	0.00484491920031939\\
4439.01843486754	0.00522802656357129\\
4444.8785252106	0.00543386251869322\\
4450.73861555366	0.00546777252190704\\
4456.59870589671	0.0053352749374476\\
4462.45879623977	0.00504204934483239\\
4468.31888658283	0.00459392473771451\\
4474.17897692589	0.00399686763332092\\
4480.03906726895	0.00325697011347182\\
4485.89915761201	0.00238043781621532\\
4491.75924795507	0.00137357789895909\\
4497.61933829813	0.000242786992088656\\
4503.47942864119	-0.0010054608368811\\
4509.33951898425	-0.00236462608924002\\
4515.19960932731	-0.00382811579186758\\
4521.05969967037	-0.00538929538585851\\
4526.91979001343	-0.00704150052832416\\
4532.77988035649	-0.0087780487888698\\
4538.63997069955	-0.0105922512223704\\
4544.50006104261	-0.0124774238011881\\
4550.36015138567	-0.0144268986882795\\
4556.22024172873	-0.0164340353348914\\
4562.08033207179	-0.0184922313870945\\
4567.94042241485	-0.020594933383191\\
4579.66060310097	-0.0249079484284059\\
4591.38078378708	-0.0293220225675828\\
4603.1009644732	-0.0337875244621813\\
4614.82114515932	-0.038256546367968\\
4626.54132584544	-0.0426831816400012\\
4644.12159687462	-0.0491488528904369\\
4661.7018679038	-0.0552849484671149\\
4679.28213893298	-0.0609690344120997\\
4696.86240996215	-0.0660940975749839\\
4714.44268099133	-0.0705695301285212\\
4737.88304236357	-0.0754016910161567\\
4761.32340373581	-0.078824432068685\\
4784.76376510804	-0.080767268888778\\
4808.20412648028	-0.0812137232874689\\
4831.64448785252	-0.0801984987493339\\
4855.08484922476	-0.0778033315588932\\
4878.525210597	-0.0741516930284429\\
4901.96557196923	-0.0694025438860658\\
4925.40593334147	-0.0637433610438042\\
4948.84629471371	-0.0573826692757273\\
4972.28665608595	-0.0505423156890993\\
4989.86692711513	-0.0452340866088032\\
5007.44719814431	-0.0398795686526454\\
5025.02746917348	-0.0345719762746803\\
5042.60774020266	-0.0294005769326532\\
5060.18801123184	-0.0244492091518728\\
5071.90819191796	-0.0213089132012892\\
5083.62837260408	-0.0183217482260002\\
5095.3485532902	-0.0155069288851972\\
5107.06873397632	-0.012881845815389\\
5118.78891466243	-0.0104619643084524\\
5130.50909534855	-0.00826074046401502\\
5136.36918569161	-0.00724576417652113\\
5142.22927603467	-0.00628955509950849\\
5148.08936637773	-0.0053932036677192\\
5153.94945672079	-0.00455766556858167\\
5159.80954706385	-0.00378376121516422\\
5165.66963740691	-0.00307217549572846\\
5171.52972774997	-0.00242345779756414\\
5177.38981809303	-0.00183802230208667\\
5183.24990843609	-0.00131614854860384\\
5189.10999877915	-0.000857982262230763\\
5194.97008912221	-0.000463536442155097\\
5200.83017946527	-0.000132692705418172\\
5206.69026980832	0.000134797118939733\\
5212.55036015138	0.000339309150907825\\
5218.41045049444	0.000481345381755638\\
5224.2705408375	0.000561531363910543\\
5230.13063118056	0.000580613675216832\\
5235.99072152362	0.000539457164811521\\
5241.85081186668	0.000439041988111794\\
5247.71090220974	0.000280460438792809\\
5253.5709925528	6.4913585937501e-05\\
5259.43108289586	-0.00020629227510786\\
5265.29117323892	-0.000531749349427142\\
5271.15126358198	-0.000909952237710022\\
5277.01135392504	-0.00133930199060966\\
5282.8714442681	-0.00181811031274662\\
5288.73153461116	-0.00234460390672503\\
5294.59162495422	-0.00291692894699858\\
5300.45171529728	-0.00353315567387322\\
5306.31180564034	-0.00419128309699356\\
5312.1718959834	-0.0048892437984047\\
5318.03198632646	-0.00562490882429211\\
5323.89207666952	-0.00639609265514323\\
5335.61225735563	-0.00803602210988691\\
5347.33243804175	-0.00979060947598667\\
5359.05261872787	-0.011640854441922\\
5370.77279941399	-0.0135673456095694\\
5382.49298010011	-0.0155504272908981\\
5394.21316078623	-0.0175703652333819\\
5411.79343181541	-0.0206264570249625\\
5429.37370284459	-0.0236561994680502\\
5446.95397387376	-0.0265968666337916\\
5464.53424490294	-0.0293892177415632\\
5482.11451593212	-0.0319784943625235\\
5499.6947869613	-0.0343153102349622\\
5523.13514833354	-0.0369645697970126\\
5546.57550970577	-0.0390053866662389\\
5570.01587107801	-0.0403775975488117\\
5593.45623245025	-0.0410455768274131\\
5616.89659382249	-0.0409985804708773\\
5640.33695519473	-0.0402502743376171\\
5663.77731656696	-0.0388374875334381\\
5687.2176779392	-0.0368182547657073\\
5710.65803931144	-0.0342692322365158\\
5734.09840068368	-0.0312825889960758\\
5757.53876205592	-0.0279624893494367\\
5775.11903308509	-0.0253209774353723\\
5792.69930411427	-0.0226041612950089\\
5810.27957514345	-0.0198615860907476\\
5827.85984617263	-0.0171422580846296\\
5845.44011720181	-0.0144937414268623\\
5863.02038823099	-0.0119613011743077\\
5874.7405689171	-0.0103584689113255\\
5886.46074960322	-0.00883762661927349\\
5898.18093028934	-0.00740951237153606\\
5909.90111097546	-0.00608394503761024\\
5921.62129166158	-0.00486975521662219\\
5927.48138200464	-0.00430690126063523\\
5933.3414723477	-0.0037747252939739\\
5939.20156269076	-0.00327402805282989\\
5945.06165303382	-0.0028055389858606\\
5950.92174337688	-0.00236991521549063\\
5956.78183371994	-0.00196774065625192\\
5962.64192406299	-0.00159952529117195\\
5968.50201440605	-0.00126570460641393\\
5974.36210474911	-0.000966639184324193\\
5980.22219509217	-0.000702614454714258\\
5986.08228543523	-0.000473840603818683\\
5991.94237577829	-0.000280452640003633\\
5997.80246612135	-0.000122510615126645\\
nan	nan\\
6009.52264680747	8.71677877848544e-05\\
6015.38273715053	0.00013915470579904\\
6021.24282749359	0.000156195260852918\\
6027.10291783665	0.000138595537882602\\
6032.96300817971	8.67320865952666e-05\\
6038.82309852277	1.05066909580283e-06\\
6044.68318886583	-0.000117935128014562\\
6050.54327920889	-0.000269645413337205\\
6056.40336955195	-0.000453435638175552\\
6062.26345989501	-0.000668598367619746\\
6068.12355023807	-0.000914365170761012\\
6073.98364058113	-0.00118990862541044\\
6079.84373092419	-0.00149434443234112\\
6085.70382126725	-0.00182673363421009\\
6091.5639116103	-0.00218608493357806\\
6097.42400195336	-0.0025713571050529\\
6103.28409229642	-0.00298146149550286\\
6109.14418263948	-0.00341526460706999\\
6115.00427298254	-0.00387159075676286\\
6126.72445366866	-0.00484691496030564\\
6138.44463435478	-0.00589729026377876\\
6150.1648150409	-0.00701207879060207\\
6161.88499572702	-0.0081802436172531\\
6173.60517641314	-0.00939044587901222\\
6185.32535709926	-0.0106311431676606\\
6202.90562812843	-0.0125238849122458\\
6220.48589915761	-0.0144198008258561\\
6238.06617018679	-0.0162803607417812\\
6255.64644121597	-0.0180684446572701\\
6273.22671224515	-0.0197490117059137\\
6296.66707361738	-0.0217670759174429\\
6320.10743498962	-0.0234694346117178\\
6343.54779636186	-0.024800895361738\\
6366.9881577341	-0.0257200549059076\\
6390.42851910634	-0.0262003496211861\\
6413.86888047857	-0.026230592106393\\
6437.30924185081	-0.0258149879582533\\
6460.74960322305	-0.0249726439197466\\
6484.18996459529	-0.0237365949305218\\
6507.63032596753	-0.0221523927179412\\
6531.07068733976	-0.0202763119640315\\
6554.511048712	-0.0181732413749578\\
6572.09131974118	-0.0164896149110248\\
6589.67159077036	-0.0147497797944119\\
6607.25186179954	-0.0129858417428408\\
6624.83213282871	-0.0112298733508544\\
6642.41240385789	-0.00951331145706943\\
6659.99267488707	-0.00786637594471415\\
6671.71285557319	-0.00682123976455288\\
6683.43303625931	-0.00582763032354333\\
6695.15321694543	-0.00489292472691127\\
6706.87339763155	-0.00402393187227838\\
6718.59357831766	-0.00322684116990391\\
6724.45366866072	-0.0028570125857194\\
6730.31375900378	-0.00250717696029412\\
6736.17384934684	-0.00217791442732509\\
6742.0339396899	-0.00186975895010211\\
6747.89403003296	-0.00158319743259294\\
6753.75412037602	-0.00131866893334073\\
6759.61421071908	-0.00107656398377321\\
6765.47430106214	-0.000857224011372839\\
6771.3343914052	-0.000660940868717434\\
6777.19448174826	-0.00048795646860541\\
6783.05457209132	-0.000338462525666669\\
6788.91466243438	-0.000212600404349719\\
6794.77475277744	-0.000110461073133609\\
6800.6348431205	-3.20851645051866e-05\\
6806.49493346356	2.25368599307568e-05\\
6812.35502380662	5.34644400932771e-05\\
6818.21511414968	6.08065448811398e-05\\
6824.07520449274	4.47211977396864e-05\\
6829.9352948358	5.41490181981804e-06\\
6835.79538517886	-5.68580335656158e-05\\
6841.65547552191	-0.000141796263310171\\
6847.51556586497	-0.000249052271612713\\
6853.37565620803	-0.000378233329560688\\
6859.23574655109	-0.000528902542717993\\
6865.09583689415	-0.000700579986113411\\
6870.95592723721	-0.000892743923840465\\
6876.81601758027	-0.00110483211031015\\
6882.67610792333	-0.00133624317001102\\
6888.53619826639	-0.00158633805257091\\
6894.39628860945	-0.00185444155975373\\
6900.25637895251	-0.00213984394061383\\
6906.11646929557	-0.00244180255162248\\
6917.83664998169	-0.00309226380890344\\
6929.55683066781	-0.00379929311741169\\
6941.27701135393	-0.00455594769407947\\
6952.99719204004	-0.00535493635989402\\
6964.71737272616	-0.00618868403251014\\
6976.43755341228	-0.00704939798005758\\
6994.01782444146	-0.00837363117112833\\
7011.59809547064	-0.00971356553958538\\
7029.17836649982	-0.0110422187516779\\
7046.75863752899	-0.0123332146626785\\
7064.33890855817	-0.0135612713315236\\
7081.91917958735	-0.0147026617481923\\
7105.35954095959	-0.0160524478550469\\
7128.79990233183	-0.017164911640962\\
7152.24026370407	-0.0180039994718255\\
7175.6806250763	-0.0185435081848604\\
7199.12098644854	-0.0187677824095412\\
7222.56134782078	-0.0186720486407515\\
7246.00170919302	-0.0182623814027911\\
7269.44207056526	-0.0175553090788418\\
7292.88243193749	-0.0165770787554313\\
7316.32279330973	-0.0153626103708321\\
7339.76315468197	-0.0139541802068539\\
7363.20351605421	-0.0123998820348884\\
7380.78378708338	-0.011169733244554\\
7398.36405811256	-0.0099098093431418\\
7415.94432914174	-0.00864335340509519\\
7433.52460017092	-0.00739345189657561\\
7451.1048712001	-0.00618260106409093\\
7462.82505188622	-0.00540772731425485\\
7474.54523257233	-0.00466582556340511\\
7486.26541325845	-0.00396260676462862\\
7497.98559394457	-0.00330341083755303\\
7509.70577463069	-0.00269316478848812\\
7521.42595531681	-0.00213634463405881\\
7527.28604565987	-0.00187923625356278\\
7533.14613600293	-0.00163694141651033\\
7539.00622634599	-0.001409878362134\\
7544.86631668905	-0.00119843156113038\\
7550.72640703211	-0.00100295104139083\\
7556.58649737517	-0.000823751788772461\\
7562.44658771823	-0.000661113224018044\\
7568.30667806129	-0.000515278756450276\\
7574.16676840435	-0.000386455415196423\\
7580.02685874741	-0.000274813558354109\\
7585.88694909047	-0.0001804866603871\\
7591.74703943353	-0.000103571177926833\\
7597.60712977658	-4.41264938878022e-05\\
7603.46722011964	-2.17493976741762e-06\\
7609.3273104627	2.22981042770951e-05\\
7615.18740080576	2.93440320654991e-05\\
7621.04749114882	1.90507574305546e-05\\
7626.90758149188	-8.45761550854606e-06\\
7632.76767183494	-5.30211968507886e-05\\
7638.627762178	-0.000114444789048414\\
7644.48785252106	-0.000192498437387347\\
7650.34794286412	-0.000286918051360987\\
7656.20803320718	-0.000397406095374902\\
7662.06812355024	-0.000523632347051589\\
7667.9282138933	-0.000665234721270361\\
7673.78830423636	-0.000821820157997939\\
7679.64839457942	-0.000992965571739724\\
7685.50848492248	-0.00117821886042376\\
7691.36857526554	-0.00137709997138264\\
7697.2286656086	-0.00158910202191079\\
7703.08875595166	-0.00181369247194271\\
7714.80893663777	-0.00229838750201581\\
7726.52911732389	-0.00282645917370073\\
7738.24929801001	-0.00339285388872418\\
7749.96947869613	-0.0039922353658527\\
7761.68965938225	-0.0046190316944184\\
7773.40984006837	-0.0052674840045619\\
7790.99011109754	-0.00626784498545595\\
7808.57038212672	-0.00728343591455004\\
7826.1506531559	-0.00829399605778873\\
7843.73092418508	-0.0092795967675229\\
7861.31119521426	-0.0102210139008855\\
7878.89146624344	-0.0111000827878135\\
7902.33182761568	-0.0121464344563453\\
7925.77218898791	-0.0130173601795693\\
7949.21255036015	-0.0136843239245271\\
7972.65291173239	-0.014126073263905\\
7996.09327310463	-0.0143292441174659\\
8019.53363447687	-0.014288693288608\\
8042.9739958491	-0.0140075516483069\\
8066.41435722134	-0.013497000204998\\
8089.85471859358	-0.0127757804825216\\
8113.29507996582	-0.0118694592992409\\
8136.73544133805	-0.0108094759260288\\
8160.17580271029	-0.0096320064451336\\
8177.75607373947	-0.00869555529431648\\
8195.33634476865	-0.00773282548940009\\
8212.91661579783	-0.00676169411315181\\
8230.496886827	-0.00580001861840479\\
8248.07715785618	-0.00486529950990427\\
8259.7973385423	-0.00426547128903829\\
8271.51751922842	-0.00368986735086344\\
8283.23769991454	-0.00314300127055771\\
8294.95788060066	-0.00262911461177559\\
8306.67806128678	-0.00215214304136545\\
8318.3982419729	-0.0017156852439257\\
8324.25833231596	-0.00151367448811043\\
8330.11842265902	-0.00132297487807245\\
8335.97851300208	-0.00114392951863701\\
8341.83860334514	-0.000976855789136341\\
8347.69869368819	-0.000822044759626161\\
8353.55878403125	-0.000679760664552227\\
8359.41887437431	-0.00055024043461665\\
8365.27896471737	-0.000433693287773499\\
8371.13905506043	-0.000330300379891579\\
8376.99914540349	-0.000240214515715671\\
8382.85923574655	-0.000163559920488515\\
8388.71932608961	-0.000100432072515783\\
8394.57941643267	-5.08975968296321e-05\\
8400.43950677573	-1.49942199766626e-05\\
8406.29959711879	7.26921418331051e-06\\
8412.15968746185	1.59126679020703e-05\\
8418.01977780491	1.09847725478301e-05\\
8423.87986814797	-7.43736831905056e-06\\
8429.73995849103	-3.92484351895344e-05\\
8435.60004883409	-8.4315207436601e-05\\
8441.46013917715	-0.000142476935878959\\
8447.32022952021	-0.000213545772104161\\
8453.18031986327	-0.000297307253461387\\
8459.04041020633	-0.000393520842531137\\
8464.90050054939	-0.0005019205197564\\
8470.76059089244	-0.000622215427785242\\
8476.6206812355	-0.000754090566093032\\
8482.48077157856	-0.000897207534121875\\
8488.34086192162	-0.00105120532135851\\
8494.20095226468	-0.00121570114241605\\
8500.06104260774	-0.00139029131535057\\
8511.78122329386	-0.00176804106188055\\
8523.50140397998	-0.00218084307301458\\
8535.2215846661	-0.00262481423580925\\
8546.94176535222	-0.00309583426461801\\
8558.66194603833	-0.00358958192631633\\
8570.38212672445	-0.00410157272344232\\
8582.10230741057	-0.00462719770940002\\
8599.68257843975	-0.00543091830670217\\
8617.26284946893	-0.0062387563887071\\
8634.84312049811	-0.0070348019703062\\
8652.42339152729	-0.00780352308687485\\
8670.00366255646	-0.00853005589558444\\
8693.4440239287	-0.00940924703232076\\
8716.88438530094	-0.0101592461965354\\
8740.32474667318	-0.0107551915574221\\
8763.76510804542	-0.0111777149487521\\
8787.20546941765	-0.0114135236097825\\
8810.64583078989	-0.0114557714405134\\
8834.08619216213	-0.0113042099318977\\
8857.52655353437	-0.0109651164000528\\
8880.96691490661	-0.0104510046110101\\
8904.40727627884	-0.00978013008476346\\
8927.84763765108	-0.00897580907050897\\
8951.28799902332	-0.00806557617746692\\
8968.8682700525	-0.00733182030297483\\
8986.44854108167	-0.00656959696098284\\
9004.02881211085	-0.00579316675727577\\
9021.60908314003	-0.00501693804219783\\
9039.18935416921	-0.00425519297103749\\
9050.90953485533	-0.00376229125416531\\
9062.62971554145	-0.00328597535268875\\
9074.34989622757	-0.00283004623514711\\
9086.07007691369	-0.00239811150803011\\
9097.79025759981	-0.00199355611571542\\
9109.51043828592	-0.00161951508258513\\
9121.23061897204	-0.00127884851523458\\
9127.0907093151	-0.00112184366623683\\
9132.95079965816	-0.000974119061537158\\
9138.81089000122	-0.000835945287546409\\
9144.67098034428	-0.00070757200220806\\
9150.53107068734	-0.000589227467265248\\
9156.3911610304	-0.000481118127020635\\
9162.25125137346	-0.000383428234241947\\
9168.11134171652	-0.00029631952393277\\
9173.97143205958	-0.000219930935522611\\
9179.83152240264	-0.000154378383879877\\
9185.6916127457	-9.97545795444341e-05\\
9191.55170308876	-5.61288983802229e-05\\
9197.41179343182	-2.35473007952853e-05\\
9203.27188377488	-2.03230056979739e-06\\
9209.13197411794	8.41701679615916e-06\\
9214.992064461	7.82492638399595e-06\\
9220.85215480405	-3.76105290718477e-06\\
9226.71224514711	-2.62703238093769e-05\\
9232.57233549017	-5.96094243329286e-05\\
9238.43242583323	-0.000103662289029873\\
9244.29251617629	-0.000158290557183506\\
9250.15260651935	-0.000223333927157371\\
9256.01269686241	-0.000298610556014676\\
9261.87278720547	-0.000383917503452488\\
9267.73287754853	-0.000479031218977657\\
9273.59296789159	-0.000583708071202371\\
9279.45305823465	-0.000697684918003724\\
9285.31314857771	-0.000820679716238142\\
9291.17323892077	-0.000952392169648613\\
9297.03332926383	-0.00109250441344431\\
9302.89341960689	-0.00124068173402229\\
9314.613600293	-0.00155981305872976\\
9326.33378097912	-0.0019068008664674\\
9338.05396166524	-0.00227844166743545\\
9349.77414235136	-0.00267134196348359\\
9361.49432303748	-0.00308194789715026\\
9373.2145037236	-0.00350657608903786\\
9390.79477475278	-0.0041615200974647\\
9408.37504578196	-0.00482648881086774\\
9425.95531681113	-0.00548839807309779\\
9443.53558784031	-0.00613433036449451\\
9461.11585886949	-0.00675177795219038\\
9478.69612989867	-0.00732887633375124\\
9502.13649127091	-0.00801670951534958\\
9525.57685264315	-0.00859032199472302\\
9549.01721401538	-0.00903075302024208\\
9572.45757538762	-0.00932372693291143\\
9595.89793675986	-0.00946008298952435\\
9619.3382981321	-0.00943602918368367\\
9642.77865950433	-0.00925321340688587\\
9666.21902087657	-0.00891861146941937\\
9689.65938224881	-0.00844423763092448\\
9713.09974362105	-0.00784668916752573\\
9736.54010499328	-0.00714654193616621\\
9759.98046636552	-0.00636761870536307\\
9777.5607373947	-0.00574744695550532\\
9795.14100842388	-0.0051093595419891\\
9812.72127945306	-0.00446527880737508\\
9830.30155048224	-0.00382715294154637\\
9847.88182151142	-0.00320672950243876\\
9859.60200219753	-0.00280856542828907\\
9871.32218288365	-0.00242652378287964\\
9883.04236356977	-0.00206366072592535\\
9894.76254425589	-0.00172285567843108\\
9906.48272494201	-0.00140678791055024\\
9918.20290562813	-0.00111791493181928\\
9924.06299597119	-0.000984378947466937\\
9929.92308631425	-0.000858452856706486\\
9935.78317665731	-0.000740373083796323\\
9941.64326700037	-0.000630358900904697\\
9947.50335734343	-0.000528612007057587\\
9953.36344768649	-0.0004353161452437\\
9959.22353802955	-0.000350636758274933\\
9965.08362837261	-0.000274720684174514\\
9970.94371871567	-0.000207695891532479\\
9976.80380905872	-0.000149671255333666\\
9982.66389940178	-0.000100736373622287\\
9988.52398974484	-6.09614252983426e-05\\
9994.3840800879	-3.03970692294141e-05\\
10000.244170431	-9.07438479967228e-06\\
10006.104260774	2.99514607758894e-06\\
10011.9643511171	5.81961555461356e-06\\
10017.8244414601	-5.73374805836168e-07\\
10023.6845318032	-1.61368196951621e-05\\
10029.5446221463	-4.08044564115603e-05\\
10035.4047124893	-7.44909514947628e-05\\
10041.2648028324	-0.00011709212714911\\
10047.1248931754	-0.00016848522688702\\
10052.9849835185	-0.000228529219733049\\
10058.8450738616	-0.000297065142248737\\
10064.7051642046	-0.000373916477591258\\
10070.5652545477	-0.000458889570673991\\
10076.4253448907	-0.000551774078525254\\
10082.2854352338	-0.000652343454739857\\
10088.1455255769	-0.000760355466990062\\
10094.0056159199	-0.000875552746390671\\
10099.865706263	-0.000997663367469788\\
10111.5858869491	-0.00126146783372883\\
10123.3060676352	-0.00154932617225742\\
10135.0262483213	-0.00185860456324919\\
10146.7464290074	-0.00218650118910178\\
10158.4666096936	-0.00253007066847834\\
10170.1867903797	-0.00288624958414371\\
10187.7670614089	-0.00343724056602401\\
10205.347332438	-0.00399859697078768\\
10222.9276034672	-0.00455930909219037\\
10240.5078744964	-0.00510845869602223\\
10258.0881455256	-0.00563542519493778\\
10275.6684165548	-0.00613008474583198\\
10299.108777927	-0.00672318222266476\\
10322.5491392992	-0.00722223915348606\\
10345.9895006715	-0.00761068283890666\\
10369.4298620437	-0.00787581883625711\\
10392.8702234159	-0.00800922188084799\\
10416.3105847882	-0.00800697993267605\\
10439.7509461604	-0.00786978450421604\\
10463.1913075327	-0.00760286563735366\\
10486.6316689049	-0.0072157751154065\\
10510.0720302771	-0.00672202654818299\\
10533.5123916494	-0.00613860568526369\\
10556.9527530216	-0.00548536854095951\\
10574.5330240508	-0.00496281748628308\\
10592.11329508	-0.00442321470788328\\
10609.6935661091	-0.00387667524878795\\
10627.2738371383	-0.00333338047999822\\
10644.8541081675	-0.00280338471955571\\
10656.5742888536	-0.00246226982840681\\
10668.2944695397	-0.00213417286747633\\
10680.0146502259	-0.00182174124790012\\
10691.734830912	-0.00152747882211605\\
10703.4550115981	-0.00125372537631666\\
10715.1751922842	-0.00100263759482658\\
10721.0352826273	-0.000886212905200961\\
10726.8953729703	-0.000776171650899172\\
10732.7554633134	-0.000672724456844635\\
10738.6155536565	-0.000576067563078175\\
10744.4756439995	-0.000486382440981392\\
10750.3357343426	-0.000403835441610912\\
10756.1958246856	-0.000328577476596296\\
10762.0559150287	-0.000260743732411699\\
10767.9160053717	-0.000200453418395766\\
10773.7760957148	-0.000147809549082245\\
10779.6361860579	-0.000102898761174201\\
10785.4962764009	-6.57911654813149e-05\\
10791.356366744	-3.65402340709376e-05\\
10797.216457087	-1.51827227904997e-05\\
10803.0765474301	-1.73862924079821e-06\\
10808.9366377732	3.78881377659286e-06\\
10814.7967281162	1.41310940428786e-06\\
10820.6568184593	-8.83556739862741e-06\\
10826.5169088023	-2.69104690456881e-05\\
10832.3769991454	-5.27484092917896e-05\\
10838.2370894885	-8.62699316219065e-05\\
10844.0971798315	-0.000127379511635701\\
10849.9572701746	-0.000175965792927513\\
10855.8173605176	-0.000231901855888585\\
10861.6774508607	-0.000295045518808255\\
10867.5375412038	-0.000365239670544749\\
10873.3976315468	-0.000442312634040603\\
10879.2577218899	-0.000526078559795566\\
10885.1178122329	-0.000616337848460132\\
10890.977902576	-0.000712877601540279\\
10896.8379929191	-0.000815472099262405\\
10908.5581736052	-0.00103786139158772\\
10920.2783542913	-0.00128146330477403\\
10931.9985349774	-0.00154406501829236\\
10943.7187156635	-0.00182330269912859\\
10955.4388963497	-0.00211668208161304\\
10967.1590770358	-0.00242160006130094\\
10978.8792577219	-0.00273536711684524\\
10996.4595287511	-0.00321657624242851\\
11014.0397997802	-0.00370205835039019\\
11031.6200708094	-0.00418235783207747\\
11049.2003418386	-0.00464817693161243\\
11066.7806128678	-0.00509055100250311\\
11090.22097424	-0.00562937647035128\\
11113.6613356123	-0.00609337664524348\\
11137.1016969845	-0.00646704616429113\\
11160.5420583567	-0.00673804868416983\\
11183.982419729	-0.00689760606791184\\
11207.4227811012	-0.00694076489891446\\
11230.8631424734	-0.00686653236261184\\
11254.3035038457	-0.00667787797909878\\
11277.7438652179	-0.00638160219779822\\
11301.1842265902	-0.00598807732340355\\
11324.6245879624	-0.00551087048014656\\
11348.0649493346	-0.00496626220167837\\
11365.6452203638	-0.00452459467654208\\
11383.225491393	-0.0040637079253725\\
11400.8057624222	-0.00359227986894734\\
11418.3860334513	-0.00311913899205397\\
11435.9663044805	-0.00265309607988796\\
11453.5465755097	-0.00220277726615378\\
11465.2667561958	-0.00191540775218794\\
11476.9869368819	-0.00164105178308009\\
11488.7071175681	-0.00138192579859688\\
11500.4272982542	-0.00114010927397309\\
11512.1474789403	-0.000917527687418584\\
11523.8676596264	-0.000715936851649411\\
11529.7277499695	-0.000623511742444622\\
11535.5878403125	-0.000536908736441398\\
11541.4479306556	-0.000456293541417036\\
11547.3080209987	-0.000381818895900072\\
11553.1681113417	-0.000313624270939049\\
11559.0282016848	-0.000251835600389362\\
11564.8882920278	-0.000196565040367412\\
11570.7483823709	-0.000147910758218454\\
11576.608472714	-0.000105956751441114\\
11582.468563057	-7.07726968791103e-05\\
11588.3286534001	-4.24138304485396e-05\\
11594.1887437431	-2.09208575862028e-05\\
11600.0488340862	-6.31989455759075e-06\\
11605.9089244293	1.37755931010916e-06\\
11611.7690147723	2.17461847321223e-06\\
11617.6291051154	-3.91102284497095e-06\\
11623.4891954584	-1.68471553991068e-05\\
11629.3492858015	-3.65871502018753e-05\\
11635.2093761446	-6.3070084082781e-05\\
11641.0694664876	-9.62208952007329e-05\\
11646.9295568307	-0.00013595056812842\\
11652.7896471737	-0.000182156348062746\\
11658.6497375168	-0.000234721983653205\\
11664.5098278598	-0.000293517997887542\\
11670.3699182029	-0.000358401986397821\\
11676.230008546	-0.00042921894251929\\
11682.090098889	-0.000505801608369831\\
11687.9501892321	-0.000587970851146087\\
11693.8102795751	-0.000675536063791807\\
11699.6703699182	-0.000768295589181372\\
11711.3905506043	-0.000968538401198433\\
11723.1107312904	-0.00118688253509249\\
11734.8309119766	-0.00142136517875597\\
11746.5510926627	-0.00166989457027569\\
11758.2712733488	-0.0019302681327759\\
11769.9914540349	-0.00220019147766297\\
11787.5717250641	-0.00261779089249372\\
11805.1519960933	-0.00304335659034545\\
11822.7322671225	-0.00346859907837759\\
11840.3125381516	-0.00388528109790063\\
11857.8928091808	-0.00428537354115635\\
11875.47308021	-0.0046612065759582\\
11898.9134415822	-0.00511228213182316\\
11922.3538029545	-0.00549238781977268\\
11945.7941643267	-0.00578885046268634\\
11969.2345256989	-0.0059919061371329\\
11992.6748870712	-0.00609500790493707\\
12016.1152484434	-0.00609502320517862\\
12039.5556098157	-0.00599231505753933\\
12062.9959711879	-0.00579070517139075\\
12086.4363325601	-0.00549732103098363\\
12109.8766939324	-0.0051223329093238\\
12133.3170553046	-0.00467859041602923\\
12156.7574166768	-0.00418117149073315\\
12174.337687706	-0.00378293033943343\\
12191.9179587352	-0.00337143889706234\\
12209.4982297644	-0.00295443582151013\\
12227.0785007936	-0.00253972531043569\\
12244.6587718227	-0.00213502873525683\\
12256.3789525089	-0.00187450803930013\\
12268.099133195	-0.0016239098969604\\
12279.8193138811	-0.00138527762529527\\
12291.5394945672	-0.00116054662920597\\
12303.2596752533	-0.000951528422873768\\
12314.9798559394	-0.000759895750868374\\
12320.8399462825	-0.000671080390199808\\
12326.7000366256	-0.000587168930638313\\
12332.5601269686	-0.000508325572768237\\
12338.4202173117	-0.000434703527455033\\
12344.2803076547	-0.000366444710434796\\
12350.1403979978	-0.000303679461249289\\
12356.0004883409	-0.000246526287144587\\
12361.8605786839	-0.000195091632371791\\
12367.720669027	-0.00014946967337851\\
12373.58075937	-0.000109742140246979\\
12379.4408497131	-7.59781647374057e-05\\
12385.3009400562	-4.82341551947391e-05\\
12391.1610303992	-2.65536985479246e-05\\
12397.0211207423	-1.09674895595291e-05\\
12402.8812110853	-1.49328743399849e-06\\
12408.7413014284	1.86410017836111e-06\\
12414.6013917715	-8.87194216363841e-07\\
12420.4614821145	-9.72613667174347e-06\\
12426.3215724576	-2.4618857708653e-05\\
12432.1816628006	-4.55187452331639e-05\\
12438.0417531437	-7.23665642196277e-05\\
12443.9018434868	-0.000105090602865757\\
12449.7619338298	-0.000143606844864621\\
12455.6220241729	-0.000187819167386013\\
12461.4821145159	-0.000237619564309568\\
12467.342204859	-0.000292888394196339\\
12473.2022952021	-0.000353494652429067\\
12479.0623855451	-0.000419296266918004\\
12484.9224758882	-0.00049014041669363\\
12490.7825662312	-0.000565863872680999\\
12496.6426565743	-0.000646293359927287\\
12508.3628372604	-0.000820529415734817\\
12520.0830179465	-0.00101127650213027\\
12531.8031986326	-0.0012168282831699\\
12543.5233793188	-0.00143535901862137\\
12555.2435600049	-0.00166493937274596\\
12566.963740691	-0.00190355304595082\\
12584.5440117202	-0.00227383334547552\\
12602.1242827494	-0.00265249363169663\\
12619.7045537785	-0.00303217118216235\\
12637.2848248077	-0.00340552026183173\\
12654.8650958369	-0.00376535139099728\\
12672.4453668661	-0.00410476697843897\\
12695.8857282383	-0.00451444596939408\\
12719.3260896105	-0.00486257911584817\\
12742.7664509828	-0.00513752832805438\\
12766.206812355	-0.0053301961663468\\
12789.6471737273	-0.0054343159556759\\
12813.0875350995	-0.00544664484648214\\
12836.5278964717	-0.0053670539438863\\
12859.968257844	-0.00519851312115898\\
12883.4086192162	-0.00494697168780536\\
12906.8489805884	-0.00462113957348873\\
12930.2893419607	-0.00423217700209131\\
12953.7297033329	-0.00379330365670004\\
12971.3099743621	-0.00344026461380296\\
12988.8902453913	-0.0030741463658034\\
13006.4705164205	-0.00270184819879727\\
13024.0507874496	-0.00233035573420049\\
13041.6310584788	-0.0019666079865307\\
13053.3512391649	-0.00173176206430656\\
13065.0714198511	-0.00150530228981414\\
13076.7916005372	-0.00128908492892611\\
13088.5117812233	-0.00108487418502644\\
13100.2319619094	-0.000894327574196526\\
13111.9521425955	-0.000718982245030837\\
13123.6723232817	-0.000560242355756582\\
13129.5324136247	-0.000487499039817409\\
13135.3925039678	-0.000419367611854463\\
13141.2525943108	-0.00035598128924929\\
13147.1126846539	-0.000297463057510905\\
13152.9727749969	-0.000243925425517103\\
13158.83286534	-0.000195470203237305\\
13164.6929556831	-0.000152188302400557\\
13170.5530460261	-0.00011415956049917\\
13176.4131363692	-8.14525884456714e-05\\
13182.2732267122	-5.41246421823077e-05\\
13188.1333170553	-3.22215184816521e-05\\
13193.9934073984	-1.57774751126702e-05\\
13199.8534977414	-4.8151755103121e-06\\
13205.7135880845	6.54341970378372e-07\\
13211.5736784275	6.3167019038465e-07\\
13217.4337687706	-4.87098522725504e-06\\
13223.2938591137	-1.58298488625258e-05\\
13229.1539494567	-3.22096428508621e-05\\
13235.0140397998	-5.39636784412149e-05\\
13240.8741301428	-8.10339715372089e-05\\
13246.7342204859	-0.000113351381929817\\
13252.594310829	-0.000150835775911191\\
13258.454401172	-0.000193396211866665\\
13264.3144915151	-0.000240931148450913\\
13270.1745818581	-0.000293328674839658\\
13276.0346722012	-0.000350466762573949\\
13281.8947625443	-0.000412213538406629\\
13287.7548528873	-0.000478427577543718\\
13293.6149432304	-0.000548958216667969\\
13305.3351239165	-0.000702322459742393\\
13317.0553046026	-0.000870929261528749\\
13328.7754852887	-0.00105327689021096\\
13340.4956659748	-0.0012477519367211\\
13352.215846661	-0.00145264326805179\\
13363.9360273471	-0.00166615676518477\\
13375.6562080332	-0.00188643071881885\\
13393.2364790624	-0.00222532187206752\\
13410.8167500916	-0.00256852132462072\\
13428.3970211207	-0.0029093868254581\\
13445.9772921499	-0.00324134948453528\\
13463.5575631791	-0.0035580381126893\\
13486.9979245513	-0.00394612734397704\\
13510.4382859236	-0.00428324324341816\\
13533.8786472958	-0.00455807132229245\\
13557.3190086681	-0.00476146104900936\\
13580.7593700403	-0.00488672176510969\\
13604.1997314125	-0.00492983415279872\\
13627.6400927848	-0.00488957061478979\\
13651.080454157	-0.00476752102761162\\
13674.5208155292	-0.00456802352855954\\
13697.9611769015	-0.00429800319315724\\
13721.4015382737	-0.00396672453832547\\
13744.841899646	-0.00358546665249678\\
13762.4221706751	-0.00327443739756012\\
13780.0024417043	-0.00294842926925367\\
13797.5827127335	-0.00261360398102513\\
13815.1629837627	-0.00227626601918929\\
13832.7432547918	-0.00194274240345282\\
13850.323525821	-0.00161926250841621\\
13862.0437065071	-0.00141217123863335\\
13873.7638871933	-0.00121392866245647\\
13885.4840678794	-0.00102616068145417\\
13897.2042485655	-0.000850400377821322\\
13908.9244292516	-0.000688075310060081\\
13920.6446099377	-0.000540495733723048\\
13926.5047002808	-0.00047261021497196\\
13932.3647906239	-0.000408843844137917\\
13938.2248809669	-0.000349322803258595\\
13944.08497131	-0.000294164129152046\\
13949.945061653	-0.000243475478057306\\
13955.8051519961	-0.000197354910416347\\
13961.6652423392	-0.000155890696233735\\
13967.5253326822	-0.000119161141364189\\
13973.3854230253	-8.72344351123164e-05\\
13979.2455133683	-6.01685194012163e-05\\
13985.1056037114	-3.80109797838483e-05\\
13990.9656940544	-2.07989584786331e-05\\
13996.8257843975	-8.5590895997881e-06\\
14002.6858747406	-1.30745667827178e-06\\
14008.5459650836	9.50427459075276e-07\\
14014.4060554267	-1.78038155813986e-06\\
14020.2661457697	-9.48428422582442e-06\\
14026.1262361128	-2.21351840003748e-05\\
14031.9863264559	-3.96965562601339e-05\\
14037.8464167989	-6.21215392153655e-05\\
14043.706507142	-8.93530465502477e-05\\
14049.566597485	-0.000121323901525574\\
14055.4266878281	-0.000157956992237147\\
14061.2867781712	-0.000199165447665733\\
14067.1468685142	-0.000244852834125261\\
14073.0069588573	-0.00029491337167422\\
14078.8670492003	-0.00034923216998541\\
14084.7271395434	-0.000407685483172795\\
14090.5872298865	-0.000470140983008642\\
14096.4473202295	-0.000536458049918619\\
14108.1675009156	-0.000680074815249637\\
14119.8876816018	-0.000837257111180151\\
14131.6078622879	-0.0010066149790278\\
14143.328042974	-0.00118665920108699\\
14155.0482236601	-0.00137581413935497\\
14166.7684043462	-0.00157243127103134\\
14184.3486753754	-0.0018776013910078\\
14201.9289464046	-0.00218977685976661\\
14219.5092174338	-0.00250291496146803\\
14237.0894884629	-0.00281097820540363\\
14254.6697594921	-0.00310804890926719\\
14272.2500305213	-0.0033884410718132\\
14295.6903918935	-0.00372717211121178\\
14319.1307532658	-0.00401537813165297\\
14342.571114638	-0.00424340183347805\\
14366.0114760103	-0.00440366374096463\\
14389.4518373825	-0.00449090710846186\\
14412.8921987547	-0.00450236340228994\\
14436.332560127	-0.00443783318470261\\
14459.7729214992	-0.00429968009947604\\
14483.2132828714	-0.00409273859038365\\
14506.6536442437	-0.00382413888456508\\
14530.0940056159	-0.0035030555364934\\
14553.5343669882	-0.00314038836962424\\
14571.1146380173	-0.00284842832876437\\
14588.6949090465	-0.00254548029175913\\
14606.2751800757	-0.00223726496221597\\
14623.8554511049	-0.00192958274735686\\
14641.4357221341	-0.00162820335005547\\
14653.1559028202	-0.00143357435883956\\
14664.8760835063	-0.0012458651337258\\
14676.5962641924	-0.00106662478674308\\
14688.3164448785	-0.000897327202285143\\
14700.0366255646	-0.000739358761476168\\
14711.7568062508	-0.000594006834663746\\
14717.6168965938	-0.00052643465546198\\
14723.4769869369	-0.000462449137631033\\
14729.3370772799	-0.000402178771716434\\
14735.197167623	-0.00034574403913085\\
14741.0572579661	-0.000293257166773564\\
14746.9173483091	-0.000244821899432677\\
14752.7774386522	-0.000200533290366235\\
14758.6375289952	-0.000160477510506018\\
14764.4976193383	-0.000124731676664487\\
14770.3577096814	-9.33636990873341e-05\\
14776.2178000244	-6.64321486381091e-05\\
14782.0778903675	-4.39861438923031e-05\\
14787.9379807105	-2.60652583372505e-05\\
14793.7980710536	-1.26994478629447e-05\\
14799.6581613967	-3.90899866785218e-06\\
14805.5182517397	2.95504333086003e-07\\
14811.3783420828	-8.6811443922328e-08\\
14817.2384324258	-5.04711574925861e-06\\
14823.0985227689	-1.45669054970921e-05\\
14828.958613112	-2.86180551753425e-05\\
14834.818703455	-4.7162887539361e-05\\
14840.6787937981	-7.01542644141694e-05\\
14846.5388841411	-9.75356973902043e-05\\
14852.3989744842	-0.000129241478160138\\
14858.2590648272	-0.000165196828188351\\
14864.1191551703	-0.000205318067391454\\
14869.9792455134	-0.000249512801427311\\
14875.8393358564	-0.00029768012719494\\
14881.6994261995	-0.000349710856090986\\
14887.5595165425	-0.000405487754519023\\
14893.4196068856	-0.000464885801107435\\
14905.1397875717	-0.000594007970482249\\
14916.8599682578	-0.000735932211527269\\
14928.580148944	-0.000889407827800184\\
14940.3003296301	-0.00105308939068388\\
14952.0205103162	-0.00122554832635953\\
14963.7406910023	-0.00140528517846411\\
14975.4608716884	-0.00159074244143019\\
14993.0411427176	-0.00187614027613619\\
15010.6214137468	-0.00216527240276336\\
15028.201684776	-0.00245256272851965\\
15045.7819558052	-0.00273248936690219\\
15063.3622268343	-0.00299968929659867\\
15086.8025882066	-0.00332738854628988\\
15110.2429495788	-0.00361236235425926\\
15133.683310951	-0.00384503190591027\\
15157.1236723233	-0.00401762519138437\\
15180.5640336955	-0.0041244308448573\\
15204.0043950678	-0.00416198145516016\\
15227.44475644	-0.00412916048897852\\
15250.8851178122	-0.00402722956387104\\
15274.3254791845	-0.00385977550368891\\
15297.7658405567	-0.00363257930993113\\
15321.2062019289	-0.00335341180118523\\
15344.6465633012	-0.00303176311997486\\
15362.2268343304	-0.00276915263604038\\
15379.8071053595	-0.00249373244525841\\
15397.3873763887	-0.00221071580829229\\
15414.9676474179	-0.00192544384110782\\
15432.5479184471	-0.00164328368143192\\
15450.1281894763	-0.00136952648722847\\
15461.8483701624	-0.00119422610243507\\
15473.5685508485	-0.00102639131353052\\
15485.2887315346	-0.00086740798862672\\
15497.0089122207	-0.000718584280846323\\
15508.7290929068	-0.000581139735917738\\
15520.449273593	-0.000456195171261834\\
15526.309363936	-0.000398731341943496\\
15532.1694542791	-0.000344763410350772\\
15538.0295446221	-0.000294399737903101\\
15543.8896349652	-0.000247740948886146\\
15549.7497253083	-0.000204879725700828\\
15555.6098156513	-0.000165900621115956\\
15561.4699059944	-0.000130879887902337\\
15567.3299963374	-9.98853261971183e-05\\
15573.1900866805	-7.2976148895099e-05\\
15579.0501770236	-5.0202865334378e-05\\
15584.9102673666	-3.16071835082891e-05\\
15590.7703577097	-1.72219309810085e-05\\
15596.6304480527	-7.07099466777598e-06\\
15602.4905383958	-1.16927957373461e-06\\
15608.3506287389	4.773134301401e-07\\
15614.2107190819	-2.1281092255186e-06\\
15620.070809425	-8.97344969084332e-06\\
15625.930899768	-2.00376535697323e-05\\
15631.7909901111	-3.52907636955655e-05\\
15637.6510804542	-5.46939926726159e-05\\
15643.5111707972	-7.81998140127464e-05\\
15649.3712611403	-0.000105752071654079\\
15655.2313514833	-0.000137286107612723\\
15661.0914418264	-0.00017272890747811\\
15666.9515321695	-0.000211999263428327\\
15672.8116225125	-0.000255007954396796\\
15678.6717128556	-0.000301657943000936\\
15684.5318031986	-0.000351844588788032\\
15690.3918935417	-0.000405455877350871\\
15696.2519838848	-0.000462372664789287\\
15707.9721645709	-0.000585612083186852\\
15719.692345257	-0.000720477307532184\\
15731.4125259431	-0.000865786753700883\\
15743.1327066292	-0.00102027297176538\\
15754.8528873153	-0.00118259353263005\\
15766.5730680015	-0.00135134252604895\\
15784.1533390306	-0.00161332028411956\\
15801.7336100598	-0.00188140378297334\\
15819.313881089	-0.00215042256932537\\
15836.8941521182	-0.00241520420116699\\
15854.4744231474	-0.00267067248342958\\
15872.0546941765	-0.0029119435820022\\
15895.4950555488	-0.00320365634488596\\
15918.935416921	-0.00345215671861917\\
15942.3757782932	-0.00364910212262316\\
15965.8161396655	-0.0037879219159577\\
15989.2565010377	-0.00386403189965901\\
16012.69686241	-0.00387498105405611\\
16036.1372237822	-0.00382052583439182\\
16059.5775851544	-0.00370262980444198\\
16083.0179465267	-0.00352538891069844\\
16106.4583078989	-0.00329488520091443\\
16129.8986692712	-0.00301897418900374\\
16153.3390306434	-0.002707013282268\\
16170.9193016726	-0.00245568979215539\\
16188.4995727017	-0.00219476553655616\\
16206.0798437309	-0.00192917612599855\\
16223.6601147601	-0.00166393207529827\\
16241.2403857893	-0.00140402342593567\\
16252.9605664754	-0.00123612998662608\\
16264.6807471615	-0.00107417557523395\\
16276.4009278476	-0.000919504264411494\\
16288.1211085338	-0.000773396076263431\\
16299.8412892199	-0.000637056281542482\\
16311.561469906	-0.000511605350409784\\
16317.4215602491	-0.00045328785944793\\
16323.2816505921	-0.000398069638266874\\
16329.1417409352	-0.000346062880858526\\
16335.0018312782	-0.000297372884124726\\
16340.8619216213	-0.000252097831477033\\
16346.7220119643	-0.000210328591719765\\
16352.5821023074	-0.00017214853354908\\
16358.4421926505	-0.000137633356111667\\
16364.3022829935	-0.000106850935909349\\
16370.1623733366	-7.98611904064394e-05\\
16376.0224636796	-5.6715958578948e-05\\
16381.8825540227	-3.74588986587459e-05\\
16387.7426443658	-2.21254032706356e-05\\
16393.6027347088	-1.07425321233418e-05\\
16399.4628250519	-3.32896238068003e-06\\
16405.3229153949	1.05043195003138e-07\\
16411.183005738	-4.42349719875634e-07\\
16417.0430960811	-4.96455080859517e-06\\
16422.9031864241	-1.344656672264e-05\\
16428.7632767672	-2.58650404379933e-05\\
16434.6233671102	-4.21883082618168e-05\\
16440.4834574533	-6.2376474350329e-05\\
16446.3435477964	-8.63815025671135e-05\\
16452.2036381394	-0.000114147325468672\\
16458.0637284825	-0.000145609970177582\\
16463.9238188255	-0.000180697700851392\\
16469.7839091686	-0.000219331177446408\\
16475.6439995117	-0.000261423630404602\\
16481.5040898547	-0.000306881050907357\\
16487.3641801978	-0.000355602396256526\\
16493.2242705408	-0.000407479809935973\\
16504.944451227	-0.000520238766584936\\
16516.6646319131	-0.00064416803343282\\
16528.3848125992	-0.00077818507271979\\
16540.1049932853	-0.000921124013042951\\
16551.8251739714	-0.00107174565352221\\
16563.5453546575	-0.00122874806979584\\
16575.2655353437	-0.00139077773104547\\
16592.8458063728	-0.00164019175297522\\
16610.426077402	-0.00189296510718047\\
16628.0063484312	-0.00214423714665986\\
16645.5866194604	-0.00238918878177739\\
16663.1668904896	-0.00262313390868933\\
16686.6072518618	-0.00291026568369465\\
16710.047613234	-0.00316022850187203\\
16733.4879746063	-0.00336460994172886\\
16756.9283359785	-0.00351656461909876\\
16780.3686973508	-0.00361103955292185\\
16803.809058723	-0.00364493834982327\\
16827.2494200952	-0.00361721888623165\\
16850.6897814675	-0.00352892142605913\\
16874.1301428397	-0.00338312646340058\\
16897.5705042119	-0.0031848439498138\\
16921.0108655842	-0.00294083786981587\\
16944.4512269564	-0.00265939229299421\\
16962.0314979856	-0.00242942768344873\\
16979.6117690148	-0.00218810788894791\\
16997.1920400439	-0.00194000642106733\\
17014.7723110731	-0.00168981450016808\\
17032.3525821023	-0.00144225164530541\\
17049.9328531315	-0.00120197594788073\\
17061.6530338176	-0.00104807566845909\\
17073.3732145037	-0.000900703820765691\\
17085.0933951898	-0.000761084283144257\\
17096.813575876	-0.000630373411576671\\
17108.5337565621	-0.000509650373644626\\
17120.2539372482	-0.00039990815044365\\
17126.1140275913	-0.000349439536447095\\
17131.9741179343	-0.000302045281462762\\
17137.8342082774	-0.000257821725022549\\
17143.6942986204	-0.000216858420919931\\
17149.5543889635	-0.000179237953221132\\
17155.4144793066	-0.000145035767163624\\
17161.2745696496	-0.0001143200153026\\
17167.1346599927	-8.71514191847529e-05\\
17172.9947503357	-6.35831468621092e-05\\
17178.8548406788	-4.36607064672648e-05\\
17184.7149310219	-2.74218560787596e-05\\
17190.5750213649	-1.48965300430669e-05\\
17196.435111708	-6.10678190334502e-06\\
17202.295202051	-1.06674403792564e-06\\
17208.1552923941	2.17395917606153e-07\\
17214.0153827372	-2.25259817053585e-06\\
17219.8754730802	-8.467020996523e-06\\
17225.7355634233	-1.84082526626683e-05\\
17231.5956537663	-3.20508022419686e-05\\
17237.4557441094	-4.93613679520115e-05\\
17243.3158344524	-7.02989137991503e-05\\
17249.1759247955	-9.48147625138472e-05\\
17255.0360151386	-0.000122852704579042\\
17260.8961054816	-0.000154349123092558\\
17266.7561958247	-0.0001892331341979\\
17272.6162861677	-0.000227426742778135\\
17278.4763765108	-0.000268845013054363\\
17284.3364668539	-0.000313396253729843\\
17290.1965571969	-0.000360982217294307\\
17296.05664754	-0.000411498312991594\\
17307.7768282261	-0.000520872191859627\\
17319.4970089122	-0.000640563211729651\\
17331.2171895983	-0.000769530962286718\\
17342.9373702845	-0.000906658231333596\\
17354.6575509706	-0.00105076056980354\\
17366.3777316567	-0.00120059641004529\\
17383.9580026859	-0.0014332727330667\\
17401.5382737151	-0.00167145732269408\\
17419.1185447442	-0.00191057048899077\\
17436.6988157734	-0.00214602562042023\\
17454.2790868026	-0.0023733163380707\\
17471.8593578318	-0.00258810192568318\\
17495.299719204	-0.00284800142720994\\
17518.7400805762	-0.00306966032369933\\
17542.1804419485	-0.00324562526874034\\
17565.6208033207	-0.00337000730766418\\
17589.061164693	-0.00343867576068483\\
17612.5015260652	-0.00344939225448289\\
17635.9418874374	-0.00340188056494002\\
17659.3822488097	-0.00329783010838866\\
17682.8226101819	-0.00314083316079808\\
17706.2629715541	-0.00293625811604845\\
17729.7033329264	-0.00269106324199608\\
17753.1436942986	-0.00241355738528607\\
17770.7239653278	-0.00218983247968374\\
17788.304236357	-0.00195743755798805\\
17805.8845073862	-0.00172077508727904\\
17823.4647784153	-0.00148431933057469\\
17841.0450494445	-0.00125253118054514\\
17852.7652301306	-0.00110275988059628\\
17864.4854108167	-0.000958257168175217\\
17876.2055915029	-0.00082022808488528\\
17887.925772189	-0.000689821242698202\\
17899.6459528751	-0.000568119191981236\\
17911.3661335612	-0.00045612936152668\\
17917.2262239043	-0.000404068724687542\\
17923.0863142473	-0.000354775646231425\\
17928.9464045904	-0.00030835129829279\\
17934.8064949335	-0.000264890712639448\\
17940.6665852765	-0.000224482583897122\\
17946.5266756196	-0.000187209086336214\\
17952.3867659626	-0.00015314570456439\\
17958.2468563057	-0.000122361078494211\\
17964.1069466488	-9.49168628787064e-05\\
17969.9670369918	-7.08676017363276e-05\\
17975.8271273349	-5.0260617890078e-05\\
17981.6872176779	-3.31359178562001e-05\\
17987.547308021	-1.95261122754273e-05\\
17993.4073983641	-9.45635203853158e-06\\
17999.2674887071	-2.94428023271915e-06\\
nan	nan\\
18010.9876693932	-6.26058363293904e-07\\
18016.8477597363	-4.81744604422046e-06\\
18022.7078500794	-1.256161326216e-05\\
18028.5679404224	-2.38385014690801e-05\\
18034.4280307655	-3.86205909431193e-05\\
18040.2881211085	-5.68729641296078e-05\\
18046.1482114516	-7.85533845819969e-05\\
18052.0083017947	-0.000103612391332356\\
18057.8683921377	-0.000131993408469772\\
18063.7284824808	-0.000163632869696458\\
18069.5885728238	-0.000198460357582423\\
18075.4486631669	-0.000236398757207158\\
18081.3087535099	-0.000277364423876758\\
18087.168843853	-0.000321267364520118\\
18093.0289341961	-0.000368011432396361\\
18104.7491148822	-0.000469608852506454\\
18116.4692955683	-0.000581272632318132\\
18128.1894762544	-0.000702034797422572\\
18139.9096569405	-0.00083085174866964\\
18151.6298376267	-0.000966613188023733\\
18163.3500183128	-0.00110815159609184\\
18175.0701989989	-0.00125425218065833\\
18192.6504700281	-0.00147921149032229\\
18210.2307410573	-0.0017072890074432\\
18227.8110120864	-0.00193411008510074\\
18245.3912831156	-0.0021553327840013\\
18262.9715541448	-0.00236673031190815\\
18286.411915517	-0.0026263796659168\\
18309.8522768893	-0.00285265343988584\\
18333.2926382615	-0.00303792833347716\\
18356.7329996337	-0.00317598460315991\\
18380.173361006	-0.00326221221261659\\
18403.6137223782	-0.00329376217527827\\
18427.0540837505	-0.00326963812941574\\
18450.4944451227	-0.00319072521218933\\
18473.9348064949	-0.00305975541951707\\
18497.3751678672	-0.00288121078107255\\
18520.8155292394	-0.00266116777054495\\
18544.2558906116	-0.00240708834446833\\
18561.8361616408	-0.00219932674509548\\
18579.41643267	-0.00198118421492547\\
18596.9967036992	-0.00175679955082239\\
18614.5769747284	-0.00153042270630186\\
18632.1572457575	-0.00130633382351794\\
18649.7375167867	-0.00108876184037811\\
18661.4576974728	-0.000949365567354957\\
18673.177878159	-0.000815856617485192\\
18684.8980588451	-0.000689349321815146\\
18696.6182395312	-0.000570897465252642\\
18708.3384202173	-0.000461485448615252\\
18720.0586009034	-0.000362020047838695\\
18725.9186912465	-0.000316277591548625\\
18731.7787815895	-0.000273322839292687\\
18737.6388719326	-0.000233244052445928\\
18743.4989622757	-0.00019612335432281\\
18749.3590526187	-0.000162036560084215\\
18755.2191429618	-0.000131053020082639\\
18761.0792333048	-0.000103235476965918\\
18766.9393236479	-7.86399368576844e-05\\
18772.799413991	-5.73155548495061e-05\\
18778.659504334	-3.93045350541097e-05\\
18784.5195946771	-2.46420454223196e-05\\
18790.3796850201	-1.33561474888933e-05\\
18796.2397753632	-5.46774119205874e-06\\
18802.0998657063	-9.90524872404033e-07\\
18807.9599560493	6.90294731979261e-08\\
18813.8200463924	-2.28831436094514e-06\\
18819.6801367354	-8.054562652008e-06\\
18825.5402270785	-1.72145128995256e-05\\
18831.4003174216	-2.97457902099721e-05\\
18837.2604077646	-4.56188988320399e-05\\
18843.1204981077	-6.47972887199579e-05\\
18848.9805884507	-8.72374369815938e-05\\
18854.8406787938	-0.000112888944024474\\
18860.7007691369	-0.000141694644190633\\
18866.5608594799	-0.000173590730640519\\
18872.420949823	-0.000208506894207686\\
18878.281040166	-0.000246366475928381\\
18884.1411305091	-0.000287086632907327\\
18890.0012208522	-0.000330578517176891\\
18895.8613111952	-0.000376747467148975\\
18907.5814918813	-0.000476710083292226\\
18919.3016725674	-0.000586108944346025\\
18931.0218532536	-0.000703999708305402\\
18942.7420339397	-0.000829367304165914\\
18954.4622146258	-0.000961134596530641\\
18966.1823953119	-0.0010981715650256\\
18983.7626663411	-0.00131103170657061\\
19001.3429373703	-0.00152901138428191\\
19018.9232083995	-0.0017479307907947\\
19036.5034794286	-0.00196359942621619\\
19054.0837504578	-0.00217189575563247\\
19071.664021487	-0.00236884542259296\\
19095.1043828592	-0.00260735083508344\\
19118.5447442315	-0.00281099647918245\\
19141.9851056037	-0.00297292514471675\\
19165.4254669759	-0.00308770287047421\\
19188.8658283482	-0.0031514991531111\\
19212.3061897204	-0.00316221268620306\\
19235.7465510927	-0.00311953850888947\\
19259.1869124649	-0.00302497442691119\\
19282.6272738371	-0.00288176662031038\\
19306.0676352094	-0.00269479640136526\\
19329.5079965816	-0.00247041206654119\\
19352.9483579539	-0.00221621162719072\\
19370.528628983	-0.00201113198262937\\
19388.1089000122	-0.00179799330299313\\
19405.6891710414	-0.00158083860544934\\
19423.2694420706	-0.00136378112227274\\
19440.8497130997	-0.00115092601489146\\
19452.5698937859	-0.00101334785545283\\
19464.290074472	-0.000880580832073882\\
19476.0102551581	-0.000753736902092438\\
19487.7304358442	-0.00063387681071593\\
19499.4506165303	-0.000522001172095845\\
19511.1707972165	-0.000419042067812197\\
19517.0308875595	-0.000371176816142205\\
19522.8909779026	-0.000325855233192107\\
19528.7510682456	-0.000283171181858921\\
19534.6111585887	-0.000243212896917138\\
19540.4712489318	-0.000206062801140336\\
19546.3313392748	-0.000171797333797657\\
19552.1914296179	-0.000140486791887828\\
19558.0515199609	-0.000112195184460174\\
19563.911610304	-8.69801002897976e-05\\
19569.7717006471	-6.48925892133559e-05\\
19575.6317909901	-4.59770573428382e-05\\
19581.4918813332	-3.02711763954224e-05\\
19587.3519716762	-1.78058073018194e-05\\
19593.2120620193	-8.60493826722153e-06\\
19599.0721523623	-2.68563739496502e-06\\
19604.9322427054	-5.8019972881849e-08\\
19610.7923330485	-7.25230481056351e-07\\
19616.6524233915	-4.6834393550504e-06\\
19622.5125137346	-1.19218545007031e-05\\
19628.3726040776	-2.24227475303034e-05\\
19634.2326944207	-3.61614946600077e-05\\
19640.0927847638	-5.31066321693841e-05\\
19645.9528751068	-7.32199263042771e-05\\
19651.8129654499	-9.64564574642837e-05\\
19657.6730557929	-0.000122764718492462\\
19663.533146136	-0.000152086726849342\\
19669.3932364791	-0.000184358150436336\\
19675.2533268221	-0.000219508446775727\\
19681.1134171652	-0.000257461015274727\\
19686.9735075082	-0.000298133362220343\\
19692.8335978513	-0.00034143727816274\\
19704.5537785374	-0.000435559548498835\\
19716.2739592235	-0.000539015319244037\\
19727.9941399097	-0.000650913919879023\\
19739.7143205958	-0.000770294123527442\\
19751.4345012819	-0.000896132343363573\\
19763.154681968	-0.00102735134846444\\
19774.8748626541	-0.00116282942486231\\
19792.4551336833	-0.00137149493905905\\
19810.0354047125	-0.00158313571932502\\
19827.6156757417	-0.00179370166096843\\
19845.1959467708	-0.00199916885448461\\
19862.7762178	-0.00219561603434716\\
19886.2165791723	-0.00243707738301771\\
19909.6569405445	-0.00264771649768994\\
19933.0973019167	-0.0028204300975864\\
19956.537663289	-0.00294940824394311\\
19979.9780246612	-0.00303032804499619\\
20003.4183860334	-0.00306049685036157\\
20026.8587474057	-0.00303894019690302\\
20050.2991087779	-0.00296643163928709\\
20073.7394701502	-0.00284546356189473\\
20097.1798315224	-0.00268016006111007\\
20120.6201928946	-0.00247613493676936\\
20144.0605542669	-0.00224029967761424\\
20161.6408252961	-0.00204731205105907\\
20179.2210963252	-0.00184456966670262\\
20196.8013673544	-0.00163592296076743\\
20214.3816383836	-0.00142532986583139\\
20231.9619094128	-0.00121678043811988\\
20249.5421804419	-0.00101422096927117\\
20261.2623611281	-0.000884406333322104\\
20272.9825418142	-0.000760048725137236\\
20284.7027225003	-0.000642190833221445\\
20296.4229031864	-0.000531819517057388\\
20308.1430838725	-0.000429857506229384\\
20319.8632645587	-0.000337155648650876\\
20325.7233549017	-0.000294521742025921\\
20331.5834452448	-0.000254485772812664\\
20337.4435355878	-0.000217130812177925\\
20343.3036259309	-0.000182534224030026\\
20349.163716274	-0.000150767503627476\\
20355.023806617	-0.000121896128662621\\
20360.8838969601	-9.59794231308124e-05\\
20366.7439873031	-7.30704342765375e-05\\
20372.6040776462	-5.32158228612711e-05\\
20378.4641679893	-3.64557669931793e-05\\
20384.3242583323	-2.28238796987978e-05\\
20390.1843486754	-1.23471404209158e-05\\
20396.0444390184	-5.04584057474679e-06\\
20401.9045293615	-9.33543273517639e-07\\
20407.7646197046	-1.70573014601601e-08\\
20413.6247100476	-2.29642538632653e-06\\
20419.4848003907	-7.76492678911273e-06\\
20425.3448907337	-1.64090941985982e-05\\
20431.2049810768	-2.82087448914954e-05\\
20437.0650714198	-4.31370260855577e-05\\
20442.9251617629	-6.11604743852655e-05\\
20448.785252106	-8.22390891847683e-05\\
20454.645342449	-0.000106326419876124\\
20460.5054327921	-0.000133369666660804\\
20466.3655231351	-0.000163309794756426\\
20472.2256134782	-0.00019608166175015\\
20478.0857038213	-0.000231614157814808\\
20483.9457941643	-0.000269830358505824\\
20489.8058845074	-0.000310647689780881\\
20495.6659748504	-0.000353978104927964\\
20507.3861555366	-0.000447799778269396\\
20519.1063362227	-0.000550488985691926\\
20530.8265169088	-0.000661164947323027\\
20542.5466975949	-0.000778880002983106\\
20554.266878281	-0.000902627678661026\\
20565.9870589672	-0.00103135124413232\\
20583.5673299963	-0.0012313554878744\\
20601.1476010255	-0.0014362469475744\\
20618.7278720547	-0.0016421065777656\\
20636.3081430839	-0.00184500140783182\\
20653.8884141131	-0.0020410593251268\\
20671.4686851422	-0.00222654261847336\\
20694.9090465145	-0.00245133686907901\\
20718.3494078867	-0.00264349171910218\\
20741.7897692589	-0.00279652872353766\\
20765.2301306312	-0.00290529968183118\\
20788.6704920034	-0.00296615847732917\\
20812.1108533757	-0.00297708199229692\\
20835.5512147479	-0.00293773609248631\\
20858.9915761201	-0.00284948453350451\\
20882.4319374924	-0.00271534056849344\\
20905.8722988646	-0.00253986296944102\\
20929.3126602368	-0.00232900004545486\\
20952.7530216091	-0.00208988699101465\\
20970.3332926383	-0.00189684553715358\\
20987.9135636674	-0.00169611399452162\\
21005.4938346966	-0.0014915046565848\\
21023.0741057258	-0.0012868998843836\\
21040.654376755	-0.00108617822644426\\
21052.3745574411	-0.000956403497428226\\
21064.0947381272	-0.000831139163906947\\
21075.8149188133	-0.000711438520904894\\
21087.5350994994	-0.000598307129074585\\
21099.2552801856	-0.000492694339589269\\
21110.9754608717	-0.000395485299783425\\
21116.8355512147	-0.000350289394060698\\
21122.6956415578	-0.000307493506831298\\
21128.5557319009	-0.000267187008242536\\
21134.4158222439	-0.000229453971963535\\
21140.275912587	-0.000194372998928431\\
21146.13600293	-0.000162017052727458\\
21151.9960932731	-0.000132453306964239\\
21157.8561836162	-0.000105743004929365\\
21163.7162739592	-8.19413318335964e-05\\
21169.5763643023	-6.1097299924636e-05\\
21175.4364546453	-4.32536466786409e-05\\
21181.2965449884	-2.84467463083629e-05\\
21187.1566353315	-1.6706534756168e-05\\
21193.0167256745	-8.05644832862293e-06\\
21198.8768160176	-2.51337610274945e-06\\
21204.7369063606	-8.76261961876233e-08\\
21210.5969967037	-7.82905968610599e-07\\
21216.4570870468	-4.5963161948821e-06\\
21222.3171773898	-1.1518359214173e-05\\
21228.1772677329	-2.15329610357304e-05\\
21234.0373580759	-3.46175073520181e-05\\
21239.897448419	-5.07428933736981e-05\\
21245.7575387621	-6.98735873809977e-05\\
21251.6176291051	-9.19677078558858e-05\\
21257.4777194482	-0.000116977114014739\\
21263.3378097912	-0.000144847509567032\\
21269.1979001343	-0.000175518559451495\\
21275.0579904774	-0.000208924019328482\\
21280.9180808204	-0.000244991877513645\\
21286.7781711635	-0.000283644509092519\\
21292.6382615065	-0.000324798841846169\\
21304.3584421926	-0.000414254161034963\\
21316.0786228788	-0.000512591326913035\\
21327.7988035649	-0.00061896914622417\\
21339.518984251	-0.000732478899955249\\
21351.2391649371	-0.000852152069549209\\
21362.9593456232	-0.000976968564233929\\
21374.6795263094	-0.00110586537972965\\
21392.2597973385	-0.00130445474806615\\
21409.8400683677	-0.00150595524508733\\
21427.4203393969	-0.00170651903154536\\
21445.0006104261	-0.00190231942998053\\
21462.5808814553	-0.0020896236641936\\
21486.0212428275	-0.00232001175123109\\
21509.4616041997	-0.00252119409978788\\
21532.901965572	-0.00268638046715559\\
21556.3423269442	-0.00281000368688873\\
21579.7826883164	-0.00288790627088031\\
21603.2230496887	-0.00291747924063\\
21626.6634110609	-0.00289774855537648\\
21650.1037724332	-0.00282940627672859\\
21673.5441338054	-0.00271478547671585\\
21696.9844951776	-0.00255777979373768\\
21720.4248565499	-0.00236371040728484\\
21743.8652179221	-0.00213914497301313\\
21761.4454889513	-0.00195524304155938\\
21779.0257599805	-0.00176193957099966\\
21796.6060310096	-0.00156290913142642\\
21814.1863020388	-0.00136193264305908\\
21831.766573068	-0.00116282540728234\\
21849.3468440972	-0.000969364536003562\\
21861.0670247833	-0.000845345067356821\\
21872.7872054694	-0.000726513609956042\\
21884.5073861555	-0.000613870769674534\\
21896.2275668417	-0.000508364338302483\\
21907.9477475278	-0.000410881300080587\\
21919.6679282139	-0.000322240356207789\\
21925.5280185569	-0.000281471133043174\\
21931.3881089	-0.000243185030209084\\
21937.2481992431	-0.000207462215150681\\
21943.1082895861	-0.000174377412194723\\
21948.9683799292	-0.000143999745620105\\
21954.8284702722	-0.000116392594641658\\
21960.6885606153	-9.16134605594766e-05\\
21966.5486509584	-6.97138463841963e-05\\
21972.4087413014	-5.07391491756901e-05\\
21978.2688316445	-3.47285653200621e-05\\
21984.1289219875	-2.17150089512198e-05\\
21989.9890123306	-1.1725043677168e-05\\
21995.8491026737	-4.77882775783598e-06\\
22001.7091930167	-8.90072845246499e-07\\
22007.5692833598	-6.60163722716593e-08\\
22013.4293737028	-2.30740764265196e-06\\
22019.2894640459	-7.60850765080387e-06\\
22025.149554389	-1.59571026288878e-05\\
22031.009644732	-2.73345312889737e-05\\
22036.8697350751	-4.17157256969613e-05\\
22042.7298254181	-5.90692656961301e-05\\
22048.5899157612	-7.93574467543084e-05\\
22054.4500061043	-0.000102536361090831\\
22060.3100964473	-0.000128555991917456\\
22066.1701867904	-0.000157360320569037\\
22072.0302771334	-0.000188887446328079\\
22077.8903674765	-0.000223069718647253\\
22083.7504578196	-0.000259833881528768\\
22089.6105481626	-0.000299101229714281\\
22095.4706385057	-0.00034078777638871\\
22107.1908191918	-0.00043105719387192\\
22118.9109998779	-0.000529871669546716\\
22130.631180564	-0.00063638876178975\\
22142.3513612502	-0.000749701208001817\\
22154.0715419363	-0.000868844626112565\\
22165.7917226224	-0.000992805695996385\\
22183.3719936516	-0.00118546666293637\\
22200.9522646807	-0.00138291026549748\\
22218.5325357099	-0.00158136882234624\\
22236.1128067391	-0.00177705763348737\\
22253.6930777683	-0.00196624695904826\\
22271.2733487975	-0.002145332911633\\
22294.7137101697	-0.00236254086203783\\
22318.1540715419	-0.00254841778113761\\
22341.5944329142	-0.00269668955074039\\
22365.0347942864	-0.00280235687916964\\
22388.4751556586	-0.00286186316254732\\
22411.9155170309	-0.00287321353214946\\
22435.3558784031	-0.00283604110599474\\
22458.7962397754	-0.00275161824591257\\
22482.2366011476	-0.00262281247894199\\
22505.6769625198	-0.00245398860755634\\
22529.1173238921	-0.00225086034869867\\
22552.5576852643	-0.00202029654036547\\
22570.1379562935	-0.00183402866461428\\
22587.7182273227	-0.00164024049670971\\
22605.2984983518	-0.00144261668501902\\
22622.878769381	-0.00124491323039203\\
22640.4590404102	-0.00105088602549815\\
22652.1792210963	-0.000925401425474034\\
22663.8994017824	-0.000804250644338998\\
22675.6195824686	-0.000688456230458653\\
22687.3397631547	-0.000578995122488002\\
22699.0599438408	-0.000476790397472854\\
22710.7801245269	-0.000382703476811251\\
22716.64021487	-0.00033895472741485\\
22722.500305213	-0.000297526855936778\\
22728.3603955561	-0.000258507057649185\\
22734.2204858992	-0.000221977418955338\\
22740.0805762422	-0.000188014744367536\\
22745.9406665853	-0.000156690394694511\\
22751.8007569283	-0.000128070136766963\\
22757.6608472714	-0.000102214005008151\\
22763.5209376145	-7.91761751625779e-05\\
22769.3810279575	-5.90048504323481e-05\\
22775.2411183006	-4.17421602660924e-05\\
22781.1012086436	-2.74240720093062e-05\\
22786.9612989867	-1.60803156061958e-05\\
22792.8213893298	-7.73432150630204e-06\\
22798.6814796728	-2.40317190797747e-06\\
22804.5415700159	-9.75654389382338e-08\\
22810.4016603589	-8.21795347178264e-07\\
22816.261750702	-4.57374124570659e-06\\
22822.121841045	-1.13448744295206e-05\\
22827.9819313881	-2.11202767477588e-05\\
22833.8420217312	-3.38786729934728e-05\\
22839.7021120742	-4.95924767370388e-05\\
22845.5622024173	-6.82278495082941e-05\\
22851.4222927603	-8.97447731970047e-05\\
22857.2823831034	-0.000114097135518203\\
22863.1424734465	-0.000141232828365827\\
22869.0025637895	-0.000171093858826122\\
22874.8626541326	-0.000203616472647206\\
22880.7227444756	-0.000238731289862733\\
22886.5828348187	-0.000276363452316818\\
22892.4429251618	-0.000316432782750981\\
22904.1631058479	-0.000403536675847783\\
22915.883286534	-0.00049930191015217\\
22927.6034672201	-0.000602914317891804\\
22939.3236479062	-0.000713493527961728\\
22951.0438285924	-0.000830100429676494\\
22962.7640092785	-0.000951745131603513\\
22974.4841899646	-0.00107739534806444\\
22992.0644609938	-0.00127104229099939\\
23009.6447320229	-0.00146760582281903\\
23027.2250030521	-0.00166334005929498\\
23044.8052740813	-0.00185451617663497\\
23062.3855451105	-0.00203749332194539\\
23085.8259064827	-0.00226271870021161\\
23109.266267855	-0.00245958960357174\\
23132.7066292272	-0.0026214559182546\\
23156.1469905994	-0.00274285271512935\\
23179.5873519717	-0.00281968433563411\\
23203.0277133439	-0.00284936218507815\\
23226.4680747162	-0.00283089160029641\\
23249.9084360884	-0.00276490487818962\\
23273.3487974606	-0.002653639369865\\
23296.7891588329	-0.00250086139780925\\
23320.2295202051	-0.00231173858207379\\
23343.6698815773	-0.00209266490066311\\
23361.2501526065	-0.00191312729735072\\
23378.8304236357	-0.00172430801620584\\
23396.4106946649	-0.00152979955553598\\
23413.9909656941	-0.00133330200090315\\
23431.5712367232	-0.00113855262994018\\
23449.1515077524	-0.000949254823742937\\
23460.8716884385	-0.000827868621857204\\
23472.5918691246	-0.000711534937569924\\
23484.3120498108	-0.000601237361657966\\
23496.0322304969	-0.000497908301027104\\
23507.752411183	-0.000402421095557273\\
23519.4725918691	-0.000315582635360732\\
23525.3326822122	-0.000275638977228373\\
23531.1927725552	-0.000238126540927679\\
23537.0528628983	-0.000203124563407944\\
23542.9129532414	-0.000170706963778393\\
23548.7730435844	-0.000140942187121089\\
23554.6331339275	-0.000113893059913721\\
23560.4932242705	-8.96166573390238e-05\\
23566.3533146136	-6.81641827723735e-05\\
23572.2134049567	-4.95808596941853e-05\\
23578.0734952997	-3.39058362646826e-05\\
23583.9335856428	-2.1172102755075e-05\\
23589.7936759858	-1.14064220109323e-05\\
23595.6537663289	-4.62927309620972e-06\\
23601.513856672	-8.54808233626689e-07\\
23607.373947015	-9.0823133571775e-08\\
23613.2340373581	-2.33874077388739e-06\\
23619.0941277011	-7.59360866299011e-06\\
23624.9542180442	-1.58441095928798e-05\\
23630.8143083873	-2.7072585857694e-05\\
23636.6743987303	-4.12550768870749e-05\\
23642.5344890734	-5.83613702124949e-05\\
23648.3945794164	-7.83550656602712e-05\\
23654.2546697595	-0.000101193652633783\\
23660.1147601026	-0.000126828600321508\\
23665.9748504456	-0.00015520546063389\\
23671.8349407887	-0.000186263983670552\\
23677.6950311317	-0.000219938245436648\\
23683.5551214748	-0.000256156787581085\\
23689.4152118178	-0.000294842768824362\\
23695.2753021609	-0.000335914127787074\\
23706.995482847	-0.000424859686585813\\
23718.7156635331	-0.000522239440905372\\
23730.4358442193	-0.000627228047288775\\
23742.1560249054	-0.000738935829724721\\
23753.8762055915	-0.000856416311212569\\
23765.5963862776	-0.000978674225344446\\
23783.1766573068	-0.00116874445146766\\
23800.756928336	-0.00136360727140985\\
23818.3371993652	-0.0015595527887299\\
23835.9174703943	-0.00175285082931565\\
23853.4977414235	-0.00193982189351401\\
23871.0780124527	-0.00211690714030613\\
23894.5183738249	-0.0023318525077528\\
23917.9587351972	-0.00251599650596467\\
23941.3990965694	-0.00266311657780403\\
23964.8394579416	-0.00276824144011184\\
23988.2798193139	-0.00282781894062616\\
24000	-0.0028398359951788\\
};
\end{axis}
\end{tikzpicture}%

Les sorties des filtres sont les suivantes

%%%% This file was created by matlab2tikz.
%
%The latest updates can be retrieved from
%  http://www.mathworks.com/matlabcentral/fileexchange/22022-matlab2tikz-matlab2tikz
%where you can also make suggestions and rate matlab2tikz.
%
\definecolor{mycolor1}{rgb}{0.00000,0.44700,0.74100}%
\definecolor{mycolor2}{rgb}{0.85000,0.32500,0.09800}%
%
\begin{tikzpicture}

\begin{axis}[%
width=4.521in,
height=3.559in,
at={(0.758in,0.488in)},
scale only axis,
xmin=0,
xmax=800,
xlabel style={font=\color{white!15!black}},
xlabel={Temps [s]},
ymin=-1.5,
ymax=1.5,
ylabel style={font=\color{white!15!black}},
ylabel={Amplitude},
axis background/.style={fill=white},
title style={font=\bfseries},
title={Signal démodulé},
legend style={legend cell align=left, align=left, draw=white!15!black}
]
\addplot [color=mycolor1]
  table[row sep=crcr]{%
1	-0\\
2	-0.00516601315871412\\
3	-0.0142301633827628\\
4	-0.0218526290866521\\
5	-0.0216211509514324\\
6	-0.010190337881113\\
8	0.0307439906353011\\
9	0.0423913515019194\\
10	0.0381249306856262\\
11	0.0180284843521576\\
12	-0.0106283293206388\\
13	-0.0366354367853319\\
14	-0.0499946129270938\\
15	-0.0462885480906152\\
16	-0.0280869139857032\\
17	-0.00290729959090186\\
18	0.0208576070853042\\
19	0.0376312007393835\\
20	0.0459804366053049\\
21	0.0471664525310871\\
22	0.0422460987181239\\
23	0.0299147018047279\\
24	0.00675049988376486\\
25	-0.0300749496432218\\
26	-0.0792457778771904\\
27	-0.133277379493506\\
28	-0.17969068832258\\
29	-0.205287839758171\\
30	-0.201599608778338\\
31	-0.16875404040411\\
32	-0.115713759025425\\
33	-0.056684712993615\\
34	-0.00548185359355102\\
35	0.0293595206604778\\
36	0.0464620754146381\\
37	0.0499374313844783\\
38	0.0455621718339216\\
39	0.0372260787070218\\
40	0.0256643054286769\\
41	0.00991772565441806\\
43	-0.0299371926961385\\
44	-0.0440737467312147\\
45	-0.0457004302486439\\
46	-0.0321252536018619\\
47	-0.00688965172469125\\
48	0.0209248204699861\\
49	0.0404787739346375\\
50	0.0442057780693403\\
51	0.0314357180155866\\
53	-0.0139010129585131\\
54	-0.0271294554013366\\
55	-0.0274201036049817\\
56	-0.0176410022626214\\
57	-0.00479034548618529\\
58	0.00429560231327741\\
59	0.00663108858771011\\
63	-0.00428804464672794\\
64	-0.00411936934756341\\
67	0.00428804464672794\\
68	0.00411936934756341\\
71	-0.00428804464672794\\
72	-0.00411936934756341\\
75	0.00428804464672794\\
76	0.00411936934756341\\
79	-0.00428804464672794\\
80	-0.00411936934756341\\
83	0.00428804464672794\\
84	0.00411936934756341\\
87	-0.00428804464672794\\
88	-0.00411936934756341\\
91	0.00428804464672794\\
92	0.00411936934756341\\
95	-0.00428804464672794\\
96	-0.00411936934756341\\
99	0.00428804464672794\\
100	0.00411936934756341\\
103	-0.00428804464672794\\
104	-0.00411936934756341\\
107	0.00428804464672794\\
108	0.00411936934756341\\
111	-0.00428804464672794\\
112	-0.00411936934756341\\
115	0.00428804464672794\\
116	0.00411936934756341\\
119	-0.00428804464672794\\
120	-0.00411936934756341\\
123	0.00428804464672794\\
124	0.00411936934756341\\
127	-0.00428804464672794\\
128	-0.00411936934756341\\
131	0.00428804464672794\\
132	0.00411936934756341\\
135	-0.00428804464672794\\
136	-0.00411936934756341\\
139	0.00428804464672794\\
140	0.00411936934756341\\
143	-0.00428804464672794\\
144	-0.00411936934756341\\
147	0.00428804464672794\\
148	0.00411936934756341\\
151	-0.00428804464672794\\
152	-0.00411936934756341\\
155	0.00428804464672794\\
156	0.00411936934756341\\
159	-0.00428804464672794\\
160	-0.00411936934756341\\
161	-0.00153762335298779\\
162	0.0120655889402315\\
163	0.0315816323479794\\
164	0.0472161126191395\\
165	0.0491149094439152\\
166	0.0324663120760533\\
167	0.00035886459181711\\
168	-0.0371753352352471\\
169	-0.0674779089146114\\
170	-0.0807240709107191\\
171	-0.0737822915559718\\
172	-0.0507090057221831\\
174	0.00957576620226064\\
175	0.0321965179547306\\
176	0.0457178178854747\\
177	0.0501012864422137\\
178	0.045169566380423\\
179	0.0296839258747923\\
180	0.00275309211895092\\
182	-0.0694775647629058\\
183	-0.0918842550184991\\
184	-0.0836577343654881\\
185	-0.0320386274925113\\
186	0.0656921866926723\\
187	0.199118074632452\\
188	0.346241840237553\\
189	0.479166921693604\\
190	0.571496037754173\\
191	0.604430424884754\\
192	0.569494489710905\\
193	0.467709304841492\\
194	0.306789584340322\\
195	0.0985725796227825\\
196	-0.141874931753932\\
197	-0.395626411622175\\
198	-0.639291831771516\\
199	-0.846279315378752\\
200	-0.990731947914924\\
201	-1.05320808755175\\
202	-1.02578558830783\\
203	-0.914126005328967\\
204	-0.73530948037137\\
205	-0.512283663873518\\
206	-0.267434145566085\\
207	-0.0181366135494727\\
208	0.224017203640756\\
209	0.450665295745466\\
210	0.652903123807846\\
211	0.818917907066748\\
212	0.93466316448314\\
213	0.987319182359556\\
214	0.96960959447199\\
215	0.882545049253622\\
216	0.735102549556586\\
217	0.541138872522765\\
218	0.315402109943534\\
219	0.0709175774895812\\
220	-0.180950736503746\\
221	-0.428227334104008\\
222	-0.646320946564288\\
223	-0.82036885461207\\
224	-0.938509980941603\\
225	-0.992693203031081\\
226	-0.9792260238371\\
227	-0.89902620936607\\
228	-0.757559244437857\\
229	-0.564465868927073\\
230	-0.332905077272812\\
231	-0.0786573547540002\\
232	0.180950736503746\\
233	0.428227334104008\\
234	0.646320946564288\\
235	0.82036885461207\\
236	0.938509980941603\\
237	0.992693203031081\\
238	0.9792260238371\\
239	0.89902620936607\\
240	0.757559244437857\\
241	0.564465868927073\\
242	0.332905077272812\\
243	0.0786573547540002\\
244	-0.180950736503746\\
245	-0.428227334104008\\
246	-0.646320946564288\\
247	-0.82036885461207\\
248	-0.938509980941603\\
249	-0.992693203031081\\
250	-0.9792260238371\\
251	-0.89902620936607\\
252	-0.757559244437857\\
253	-0.564465868927073\\
254	-0.332905077272812\\
255	-0.0786573547540002\\
256	0.180950736503746\\
257	0.428227334104008\\
258	0.646320946564288\\
259	0.82036885461207\\
260	0.938509980941603\\
261	0.992693203031081\\
262	0.9792260238371\\
263	0.89902620936607\\
264	0.757559244437857\\
265	0.564465868927073\\
266	0.332905077272812\\
267	0.0786573547540002\\
268	-0.180950736503746\\
269	-0.428227334104008\\
270	-0.646320946564288\\
271	-0.82036885461207\\
272	-0.938509980941603\\
273	-0.992693203031081\\
274	-0.9792260238371\\
275	-0.89902620936607\\
276	-0.757559244437857\\
277	-0.564465868927073\\
278	-0.332905077272812\\
279	-0.0786573547540002\\
280	0.180950736503746\\
281	0.428227334104008\\
282	0.646320946564288\\
283	0.82036885461207\\
284	0.938509980941603\\
285	0.992693203031081\\
286	0.9792260238371\\
287	0.89902620936607\\
288	0.757559244437857\\
289	0.564465868927073\\
290	0.332905077272812\\
291	0.0786573547540002\\
292	-0.180950736503746\\
293	-0.428227334104008\\
294	-0.646320946564288\\
295	-0.82036885461207\\
296	-0.938509980941603\\
297	-0.992693203031081\\
298	-0.9792260238371\\
299	-0.89902620936607\\
300	-0.757559244437857\\
301	-0.564465868927073\\
302	-0.332905077272812\\
303	-0.0786573547540002\\
304	0.180950736503746\\
305	0.428227334104008\\
306	0.646320946564288\\
307	0.82036885461207\\
308	0.938509980941603\\
309	0.992693203031081\\
310	0.9792260238371\\
311	0.89902620936607\\
312	0.757559244437857\\
313	0.564465868927073\\
314	0.332905077272812\\
315	0.0786573547540002\\
316	-0.180950736503746\\
317	-0.428227334104008\\
318	-0.646320946564288\\
319	-0.82036885461207\\
320	-0.938509980941603\\
321	-0.992693203031081\\
322	-0.983958538745924\\
323	-0.913475975928804\\
324	-0.782629705857858\\
325	-0.595187552900256\\
326	-0.3605000942174\\
327	-0.0953805803686691\\
328	0.177067055615112\\
329	0.431697745040424\\
330	0.64687572680748\\
331	0.808716908848282\\
332	0.91205979536312\\
333	0.957973070507364\\
334	0.949492847885381\\
335	0.888195540185393\\
336	0.773626653063275\\
337	0.60586137664734\\
338	0.389579224262093\\
339	0.137085165925441\\
340	-0.131753732593666\\
341	-0.392984026134286\\
342	-0.623903144079463\\
343	-0.807734419742928\\
344	-0.935590427151851\\
345	-1.00528774631255\\
346	-1.01835701385346\\
347	-0.977485921351331\\
348	-0.886123335351499\\
349	-0.750386321943324\\
350	-0.581723281051495\\
351	-0.398040024600277\\
352	-0.221683084320716\\
353	-0.0744352633640801\\
354	0.0284986428778211\\
355	0.0827278046449464\\
356	0.0950562939427755\\
357	0.0798046262131038\\
359	0.0263576578840912\\
360	0.00608521470451251\\
361	-0.00739147722447342\\
362	-0.0155749278997064\\
363	-0.0189704627463243\\
364	-0.0161337092962412\\
365	-0.00522866327969496\\
366	0.0132137285373801\\
367	0.0344934398087844\\
368	0.0505115334977972\\
369	0.053123016224049\\
370	0.0383842101805385\\
371	0.00925670404296852\\
372	-0.0249798574286615\\
373	-0.0525825519198406\\
374	-0.0645996515569323\\
375	-0.0588120896710507\\
376	-0.0403126082965173\\
377	-0.0184873310217881\\
378	-0.0021066069415383\\
379	0.00479637435387303\\
380	0.00411936934756341\\
383	-0.00428804464672794\\
384	-0.00411936934756341\\
387	0.00428804464672794\\
388	0.00411936934756341\\
391	-0.00428804464672794\\
392	-0.00411936934756341\\
395	0.00428804464672794\\
396	0.00411936934756341\\
399	-0.00428804464672794\\
400	-0.00411936934756341\\
403	0.00428804464672794\\
404	0.00411936934756341\\
407	-0.00428804464672794\\
408	-0.00411936934756341\\
411	0.00428804464672794\\
412	0.00411936934756341\\
415	-0.00428804464672794\\
416	-0.00411936934756341\\
419	0.00428804464672794\\
420	0.00411936934756341\\
423	-0.00428804464672794\\
424	-0.00411936934756341\\
427	0.00428804464672794\\
428	0.00411936934756341\\
431	-0.00428804464672794\\
432	-0.00411936934756341\\
435	0.00428804464672794\\
436	0.00411936934756341\\
439	-0.00428804464672794\\
440	-0.00411936934756341\\
443	0.00428804464672794\\
444	0.00411936934756341\\
447	-0.00428804464672794\\
448	-0.00411936934756341\\
451	0.00428804464672794\\
452	0.00411936934756341\\
455	-0.00428804464672794\\
456	-0.00411936934756341\\
459	0.00428804464672794\\
460	0.00411936934756341\\
463	-0.00428804464672794\\
464	-0.00411936934756341\\
467	0.00428804464672794\\
468	0.00411936934756341\\
471	-0.00428804464672794\\
472	-0.00411936934756341\\
475	0.00428804464672794\\
476	0.00411936934756341\\
479	-0.00428804464672794\\
480	-0.00411936934756341\\
481	-0.00153762335298779\\
482	0.00258961872282271\\
483	0.00523518053114458\\
484	0.00151005191594322\\
485	-0.0118978938569398\\
487	-0.0554414831589156\\
488	-0.0671790562545311\\
489	-0.0593009813604795\\
490	-0.0292062578068908\\
491	0.0166209820290533\\
492	0.0643825468007435\\
493	0.0982620358153099\\
494	0.106785213531566\\
495	0.0872623678428681\\
496	0.0463715940020393\\
497	-0.00305912665521646\\
498	-0.0471785575510921\\
499	-0.0755736276280459\\
500	-0.0832586593296583\\
501	-0.0700532116325121\\
502	-0.0387326120024909\\
503	0.006198495179774\\
504	0.0580870498087052\\
505	0.106593686434621\\
506	0.136803840018388\\
507	0.130830441156149\\
508	0.0725048725116721\\
509	-0.0461486087293679\\
510	-0.219405098293805\\
512	-0.633225772154333\\
513	-0.803393010034597\\
514	-0.904276644150059\\
515	-0.915716102277884\\
516	-0.832726242793683\\
517	-0.663999212643375\\
518	-0.427639368990754\\
519	-0.14664356151809\\
520	0.154026263200194\\
521	0.448522317730863\\
522	0.710559327428655\\
523	0.914885855210514\\
524	1.04058648646765\\
525	1.07545887214962\\
526	1.0194548983776\\
527	0.885091773850718\\
528	0.693991110693446\\
529	0.470639354999776\\
530	0.235979459097621\\
531	0.003633941588987\\
532	-0.219748790378617\\
533	-0.430657579770696\\
534	-0.623805485647949\\
535	-0.78901556264907\\
536	-0.91168162626434\\
537	-0.976557747895981\\
538	-0.972826787293343\\
539	-0.897807234217453\\
540	-0.757559244437857\\
541	-0.564465868927073\\
542	-0.332905077272812\\
543	-0.0786573547541138\\
544	0.180950736503746\\
545	0.428227334104008\\
546	0.646320946564288\\
547	0.820368854611957\\
548	0.938509980941603\\
549	0.992693203031081\\
550	0.9792260238371\\
551	0.89902620936607\\
552	0.757559244437857\\
553	0.564465868927073\\
554	0.332905077272812\\
555	0.0786573547541138\\
556	-0.180950736503746\\
557	-0.428227334104008\\
558	-0.646320946564288\\
559	-0.820368854611957\\
560	-0.938509980941603\\
561	-0.992693203031081\\
562	-0.9792260238371\\
563	-0.89902620936607\\
564	-0.757559244437857\\
565	-0.564465868927073\\
566	-0.332905077272812\\
567	-0.0786573547541138\\
568	0.180950736503746\\
569	0.428227334104008\\
570	0.646320946564288\\
571	0.820368854611957\\
572	0.938509980941603\\
573	0.992693203031081\\
574	0.9792260238371\\
575	0.89902620936607\\
576	0.757559244437857\\
577	0.564465868927073\\
578	0.332905077272812\\
579	0.0786573547541138\\
580	-0.180950736503746\\
581	-0.428227334104008\\
582	-0.646320946564288\\
583	-0.820368854611957\\
584	-0.938509980941603\\
585	-0.992693203031081\\
586	-0.9792260238371\\
587	-0.89902620936607\\
588	-0.757559244437857\\
589	-0.564465868927073\\
590	-0.332905077272812\\
591	-0.0786573547541138\\
592	0.180950736503746\\
593	0.428227334104008\\
594	0.646320946564288\\
595	0.820368854611957\\
596	0.938509980941603\\
597	0.992693203031081\\
598	0.9792260238371\\
599	0.89902620936607\\
600	0.757559244437857\\
601	0.564465868927073\\
602	0.332905077272812\\
603	0.0786573547541138\\
604	-0.180950736503746\\
605	-0.428227334104008\\
606	-0.646320946564288\\
607	-0.820368854611957\\
608	-0.938509980941603\\
609	-0.992693203031081\\
610	-0.9792260238371\\
611	-0.89902620936607\\
612	-0.757559244437857\\
613	-0.564465868927073\\
614	-0.332905077272812\\
615	-0.0786573547541138\\
616	0.180950736503746\\
617	0.428227334104008\\
618	0.646320946564288\\
619	0.820368854611957\\
620	0.938509980941603\\
621	0.992693203031081\\
622	0.9792260238371\\
623	0.89902620936607\\
624	0.757559244437857\\
625	0.564465868927073\\
626	0.332905077272812\\
627	0.0786573547541138\\
628	-0.180950736503746\\
629	-0.428227334104008\\
630	-0.646320946564288\\
631	-0.820368854611957\\
632	-0.938509980941603\\
633	-0.992693203031081\\
634	-0.9792260238371\\
635	-0.89902620936607\\
636	-0.757559244437857\\
637	-0.564465868927073\\
638	-0.332905077272812\\
639	-0.0786573547541138\\
640	0.180950736503746\\
641	0.428227334104008\\
642	0.646320946564288\\
643	0.820368854611957\\
644	0.938509980941603\\
645	0.992693203031081\\
646	0.9792260238371\\
647	0.89902620936607\\
648	0.757559244437857\\
649	0.564465868927073\\
650	0.332905077272812\\
651	0.0786573547541138\\
652	-0.180950736503746\\
653	-0.428227334104008\\
654	-0.646320946564288\\
655	-0.820368854611957\\
656	-0.938509980941603\\
657	-0.992693203031081\\
658	-0.9792260238371\\
659	-0.89902620936607\\
660	-0.757559244437857\\
661	-0.564465868927073\\
662	-0.332905077272812\\
663	-0.0786573547541138\\
664	0.180950736503746\\
665	0.428227334104008\\
666	0.646320946564288\\
667	0.820368854611957\\
668	0.938509980941603\\
669	0.992693203031081\\
670	0.9792260238371\\
671	0.89902620936607\\
672	0.757559244437857\\
673	0.564465868927073\\
674	0.332905077272812\\
675	0.0786573547541138\\
676	-0.180950736503746\\
677	-0.428227334104008\\
678	-0.646320946564288\\
679	-0.820368854611957\\
680	-0.938509980941603\\
681	-0.992693203031081\\
682	-0.9792260238371\\
683	-0.89902620936607\\
684	-0.757559244437857\\
685	-0.564465868927073\\
686	-0.332905077272812\\
687	-0.0786573547541138\\
688	0.180950736503746\\
689	0.428227334104008\\
690	0.646320946564288\\
691	0.820368854611957\\
692	0.938509980941603\\
693	0.992693203031081\\
694	0.9792260238371\\
695	0.89902620936607\\
696	0.757559244437857\\
697	0.564465868927073\\
698	0.332905077272812\\
699	0.0786573547541138\\
700	-0.180950736503746\\
701	-0.428227334104008\\
702	-0.646320946564288\\
703	-0.820368854611957\\
704	-0.938509980941603\\
705	-0.992693203031081\\
706	-0.9792260238371\\
707	-0.89902620936607\\
708	-0.757559244437857\\
709	-0.564465868927073\\
710	-0.332905077272812\\
711	-0.0786573547541138\\
712	0.180950736503746\\
713	0.428227334104008\\
714	0.646320946564288\\
715	0.820368854611957\\
716	0.938509980941603\\
717	0.992693203031081\\
718	0.9792260238371\\
719	0.89902620936607\\
720	0.757559244437857\\
721	0.564465868927073\\
722	0.332905077272812\\
723	0.0786573547541138\\
724	-0.180950736503746\\
725	-0.428227334104008\\
726	-0.646320946564288\\
727	-0.820368854611957\\
728	-0.938509980941603\\
729	-0.992693203031081\\
730	-0.9792260238371\\
731	-0.89902620936607\\
732	-0.757559244437857\\
733	-0.564465868927073\\
734	-0.332905077272812\\
735	-0.0786573547541138\\
736	0.180950736503746\\
737	0.428227334104008\\
738	0.646320946564288\\
739	0.820368854611957\\
740	0.938509980941603\\
741	0.992693203031081\\
742	0.9792260238371\\
743	0.89902620936607\\
744	0.757559244437857\\
745	0.564465868927073\\
746	0.332905077272812\\
747	0.0786573547541138\\
748	-0.180950736503746\\
749	-0.428227334104008\\
750	-0.646320946564288\\
751	-0.820368854611957\\
752	-0.938509980941603\\
753	-0.992693203031081\\
754	-0.9792260238371\\
755	-0.89902620936607\\
756	-0.757559244437857\\
757	-0.564465868927073\\
758	-0.332905077272812\\
759	-0.0786573547541138\\
760	0.180950736503746\\
761	0.428227334104008\\
762	0.646320946564288\\
763	0.820368854611957\\
764	0.938509980941603\\
765	0.992693203031081\\
766	0.9792260238371\\
767	0.89902620936607\\
768	0.757559244437857\\
769	0.564465868927073\\
770	0.332905077272812\\
771	0.0786573547541138\\
772	-0.180950736503746\\
773	-0.428227334104008\\
774	-0.646320946564288\\
775	-0.820368854611957\\
776	-0.938509980941603\\
777	-0.992693203031081\\
778	-0.9792260238371\\
779	-0.89902620936607\\
780	-0.757559244437857\\
781	-0.564465868927073\\
782	-0.332905077272812\\
783	-0.0786573547541138\\
784	0.180950736503746\\
785	0.428227334104008\\
786	0.646320946564288\\
787	0.820368854611957\\
788	0.938509980941603\\
789	0.992693203031081\\
790	0.9792260238371\\
791	0.89902620936607\\
792	0.757559244437857\\
793	0.564465868927073\\
794	0.332905077272812\\
795	0.0786573547541138\\
796	-0.180950736503746\\
797	-0.428227334104008\\
798	-0.646320946564288\\
799	-0.820368854611957\\
800	-0.938509980941603\\
};
\addlegendentry{1}

\addplot [color=mycolor2]
  table[row sep=crcr]{%
1	0\\
2	0.00516601315871412\\
3	0.0142301633827628\\
4	0.0218526290866521\\
5	0.0216211509514324\\
6	0.010190337881113\\
8	-0.0307439906353011\\
9	-0.0423913515019194\\
10	-0.0381249306856262\\
11	-0.0180284843521576\\
12	0.0106283293206388\\
13	0.0366354367853319\\
14	0.0499946129270938\\
15	0.0462885480906152\\
16	0.0280869139857032\\
17	0.00290729959090186\\
18	-0.0208576070853042\\
19	-0.0376312007393835\\
20	-0.0459804366053049\\
21	-0.0471664525310871\\
22	-0.0422460987181239\\
23	-0.0299147018047279\\
24	-0.00675049988376486\\
25	0.0300749496432218\\
26	0.0792457778771904\\
27	0.133277379493506\\
28	0.17969068832258\\
29	0.205287839758171\\
30	0.201599608778338\\
31	-0.772557480878618\\
32	-0.788570162134988\\
33	-0.280854351264452\\
34	0.432413452249421\\
35	0.91195200062225\\
36	0.857821845745775\\
37	0.287601632873702\\
38	-0.472493770489677\\
39	-0.978537599989863\\
40	-0.92994822658909\\
41	-0.347456789912599\\
42	0.436719619054998\\
43	0.97124871397898\\
44	0.948357667891628\\
45	0.383239494506711\\
46	-0.394806345053894\\
47	-0.93442186955815\\
48	-0.925208741630399\\
49	-0.378017838192818\\
50	0.382725820586529\\
51	0.909875803267255\\
52	0.895649337293776\\
53	0.351440077216694\\
54	-0.399802143254419\\
55	-0.913891417677746\\
56	-0.886642918897792\\
57	-0.332748718771995\\
58	0.422635996342478\\
59	0.934680432695018\\
60	0.90016455181285\\
61	0.336001440905193\\
62	-0.424986757107831\\
63	-0.937023476636\\
64	-0.90016455181285\\
65	-0.336001440905193\\
66	0.424986757107831\\
67	0.937023476636\\
68	0.90016455181285\\
69	0.336001440905193\\
70	-0.424986757107831\\
71	-0.937023476636\\
72	-0.90016455181285\\
73	-0.336001440905193\\
74	0.424986757107831\\
75	0.937023476636\\
76	0.90016455181285\\
77	0.336001440905193\\
78	-0.424986757107831\\
79	-0.937023476636\\
80	-0.90016455181285\\
81	-0.336001440905193\\
82	0.424986757107831\\
83	0.937023476636\\
84	0.90016455181285\\
85	0.336001440905193\\
86	-0.424986757107831\\
87	-0.937023476636\\
88	-0.90016455181285\\
89	-0.336001440905193\\
90	0.424986757107831\\
91	0.937023476636\\
92	0.90016455181285\\
93	0.336001440905193\\
94	-0.424986757107831\\
95	-0.937023476636\\
96	-0.90016455181285\\
97	-0.336001440905193\\
98	0.424986757107831\\
99	0.937023476636\\
100	0.90016455181285\\
101	0.336001440905193\\
102	-0.424986757107831\\
103	-0.937023476636\\
104	-0.90016455181285\\
105	-0.336001440905193\\
106	0.424986757107831\\
107	0.937023476636\\
108	0.90016455181285\\
109	0.336001440905193\\
110	-0.424986757107831\\
111	-0.937023476636\\
112	-0.90016455181285\\
113	-0.336001440905193\\
114	0.424986757107831\\
115	0.937023476636\\
116	0.90016455181285\\
117	0.336001440905193\\
118	-0.424986757107831\\
119	-0.937023476636\\
120	-0.90016455181285\\
121	-0.336001440905193\\
122	0.424986757107831\\
123	0.937023476636\\
124	0.90016455181285\\
125	0.336001440905193\\
126	-0.424986757107831\\
127	-0.937023476636\\
128	-0.90016455181285\\
129	-0.336001440905193\\
130	0.424986757107831\\
131	0.937023476636\\
132	0.90016455181285\\
133	0.336001440905193\\
134	-0.424986757107831\\
135	-0.937023476636\\
136	-0.90016455181285\\
137	-0.336001440905193\\
138	0.424986757107831\\
139	0.937023476636\\
140	0.90016455181285\\
141	0.336001440905193\\
142	-0.424986757107831\\
143	-0.937023476636\\
144	-0.90016455181285\\
145	-0.336001440905193\\
146	0.424986757107831\\
147	0.937023476636\\
148	0.90016455181285\\
149	0.336001440905193\\
150	-0.424986757107831\\
151	-0.937023476636\\
152	-0.90016455181285\\
153	-0.336001440905193\\
154	0.424986757107831\\
155	0.937023476636\\
156	0.90016455181285\\
157	0.336001440905193\\
158	-0.424986757107831\\
159	-0.937023476636\\
160	-0.90016455181285\\
161	-0.336001440905193\\
162	0.414866009715524\\
163	0.909729888934748\\
164	0.857067808541274\\
165	0.288424154814265\\
166	-0.459397910731809\\
167	-0.941670385874545\\
168	-0.867108585925166\\
169	-0.270061155343456\\
170	0.507655669566475\\
171	1.01509381283881\\
172	0.954992926882596\\
173	0.357539675535008\\
174	-0.436507364858016\\
175	-0.973508039237458\\
176	-0.950001739045888\\
177	-0.387640350700394\\
178	0.381762032275333\\
179	0.911627595408049\\
180	0.901530829041462\\
181	0.370666655543118\\
182	-0.35745403389285\\
183	-0.849427266264343\\
184	-0.820626186794925\\
185	-0.305500436765556\\
186	0.361239411963084\\
187	0.742193446650276\\
188	0.558042080922746\\
189	-0.141627857435424\\
190	-0.998427636410042\\
191	0.298383461911044\\
192	0.191256419243359\\
193	0.0991347090264298\\
194	0.0275180516501905\\
195	-0.0195838344329786\\
196	-0.0398381663065948\\
197	-0.034405082350986\\
198	-0.00975212246055435\\
199	0.0224541737727577\\
200	0.0482679409001321\\
201	0.05633257971067\\
202	0.0424339980852437\\
203	0.0113121185331693\\
204	-0.0254414285828943\\
205	-0.0545603499945173\\
206	-0.0668734904244275\\
207	-0.0608521316404449\\
208	-0.0423041055802287\\
209	-0.0206338017724192\\
210	-0.00385916957588961\\
211	0.00490723453913233\\
214	0.013741995750479\\
215	0.0202688375421758\\
216	0.0256483593976782\\
217	0.0257051413451563\\
218	0.0189055260469786\\
219	0.00807116770022276\\
220	-0.000762361556780888\\
224	-0.00395402607318829\\
227	-0.00378767742972741\\
231	-0.000331390435803769\\
235	0.00345628699392364\\
238	0.00412556638548267\\
241	0.00237814494096256\\
248	-0.00395402607318829\\
251	-0.00378767742972741\\
255	-0.000331390435803769\\
259	0.00345628699392364\\
262	0.00412556638548267\\
265	0.00237814494096256\\
272	-0.00395402607318829\\
275	-0.00378767742972741\\
279	-0.000331390435803769\\
283	0.00345628699392364\\
286	0.00412556638548267\\
289	0.00237814494096256\\
296	-0.00395402607318829\\
299	-0.00378767742972741\\
303	-0.000331390435803769\\
307	0.00345628699392364\\
310	0.00412556638548267\\
313	0.00237814494096256\\
320	-0.00395402607318829\\
321	-0.00418230481000137\\
322	0.000606948523341089\\
323	0.0106620891331204\\
324	0.021878796903593\\
325	0.0283435390323348\\
326	0.0261924582268875\\
327	0.0163918351788652\\
328	0.00464604244530165\\
329	-0.00166625106737683\\
330	0.00216822742459044\\
331	0.0151082327577114\\
332	0.0304042116516712\\
333	0.0389024373336042\\
334	0.0338587423372019\\
335	0.0146183466104048\\
336	-0.0128757441090102\\
337	-0.0390173627793047\\
338	-0.0552715882715802\\
339	-0.0580964207356374\\
340	-0.0499593654667478\\
342	-0.0251408101526067\\
344	-0.00687357986305415\\
345	0.00841223847146466\\
346	0.0350054236308779\\
347	0.0746720345556469\\
348	0.125372426397234\\
349	0.183542308075403\\
350	0.247415645060983\\
351	-0.543271496682451\\
352	-0.682600836839697\\
353	-0.2631038008941\\
354	0.398432955778048\\
355	0.858583716637781\\
356	0.809227627217638\\
357	0.257734438045077\\
358	-0.479802620328201\\
359	-0.967669179166819\\
360	-0.910369135864926\\
361	-0.330147587033707\\
362	0.442506526555576\\
363	0.960281984029052\\
364	0.920417630456654\\
365	0.342767727537876\\
366	-0.44014532719325\\
367	-0.975804961091626\\
368	-0.95479545465821\\
369	-0.39066208048223\\
370	0.388547388475217\\
371	0.932054817239873\\
372	0.929263778589075\\
373	0.390121616178021\\
374	-0.362331947098824\\
375	-0.882499431611677\\
376	-0.863971312863896\\
377	-0.319051733236392\\
378	0.429038205597294\\
379	0.936515146928855\\
380	0.90016455181285\\
381	0.336001440905193\\
382	-0.424986757107831\\
383	-0.937023476636\\
384	-0.90016455181285\\
385	-0.336001440905193\\
386	0.424986757107831\\
387	0.937023476636\\
388	0.90016455181285\\
389	0.336001440905193\\
390	-0.424986757107831\\
391	-0.937023476636\\
392	-0.90016455181285\\
393	-0.336001440905193\\
394	0.424986757107831\\
395	0.937023476636\\
396	0.90016455181285\\
397	0.336001440905193\\
398	-0.424986757107831\\
399	-0.937023476636\\
400	-0.90016455181285\\
401	-0.336001440905193\\
402	0.424986757107831\\
403	0.937023476636\\
404	0.90016455181285\\
405	0.336001440905193\\
406	-0.424986757107831\\
407	-0.937023476636\\
408	-0.90016455181285\\
409	-0.336001440905193\\
410	0.424986757107831\\
411	0.937023476636\\
412	0.90016455181285\\
413	0.336001440905193\\
414	-0.424986757107831\\
415	-0.937023476636\\
416	-0.90016455181285\\
417	-0.336001440905193\\
418	0.424986757107831\\
419	0.937023476636\\
420	0.90016455181285\\
421	0.336001440905193\\
422	-0.424986757107831\\
423	-0.937023476636\\
424	-0.90016455181285\\
425	-0.336001440905193\\
426	0.424986757107831\\
427	0.937023476636\\
428	0.90016455181285\\
429	0.336001440905193\\
430	-0.424986757107831\\
431	-0.937023476636\\
432	-0.90016455181285\\
433	-0.336001440905193\\
434	0.424986757107717\\
435	0.937023476636\\
436	0.90016455181285\\
437	0.336001440905193\\
438	-0.424986757107831\\
439	-0.937023476636\\
440	-0.90016455181285\\
441	-0.336001440905193\\
442	0.424986757107831\\
443	0.937023476636\\
444	0.90016455181285\\
445	0.336001440905193\\
446	-0.424986757107831\\
447	-0.937023476636\\
448	-0.90016455181285\\
449	-0.336001440905193\\
450	0.424986757107831\\
451	0.937023476636\\
452	0.90016455181285\\
453	0.336001440905193\\
454	-0.424986757107831\\
455	-0.937023476636\\
456	-0.90016455181285\\
457	-0.336001440905193\\
458	0.424986757107831\\
459	0.937023476636\\
460	0.90016455181285\\
461	0.336001440905193\\
462	-0.424986757107831\\
463	-0.937023476636\\
464	-0.90016455181285\\
465	-0.336001440905193\\
466	0.424986757107717\\
467	0.937023476636\\
468	0.90016455181285\\
469	0.336001440905306\\
470	-0.424986757107717\\
471	-0.937023476636\\
472	-0.90016455181285\\
473	-0.336001440905193\\
474	0.424986757107831\\
475	0.937023476636\\
476	0.90016455181285\\
477	0.336001440905193\\
478	-0.424986757107717\\
479	-0.937023476636\\
480	-0.90016455181285\\
481	-0.336001440905193\\
482	0.424341979932933\\
483	0.936076340751697\\
484	0.90277386924447\\
485	0.34943695811512\\
486	-0.393551600182491\\
487	-0.885870038123812\\
488	-0.837104864905882\\
489	-0.278238082897701\\
490	0.456137856462647\\
491	0.924690539253675\\
492	0.83990137435967\\
493	0.239277028442984\\
494	-0.533716812187322\\
495	-1.0285738891256\\
496	-0.950655515162566\\
497	-0.334479937602964\\
498	0.474110156206734\\
499	1.01688514891077\\
500	0.987542580490071\\
501	0.407592275890693\\
502	-0.388198986653151\\
503	-0.947510016462502\\
504	-0.962370970969118\\
505	-0.444132750692916\\
506	0.290127758637368\\
507	0.810481080126578\\
508	0.831779048648741\\
509	0.383687672987548\\
510	-0.207526500361837\\
511	-0.397698078793951\\
512	-0.309238234860459\\
513	-0.193482497806485\\
514	-0.0790749460725237\\
515	0.0129022154820859\\
516	0.0719753338394185\\
517	0.0971551987753401\\
518	0.093331733000241\\
519	0.0676548163281723\\
520	0.0276868348602193\\
521	-0.0184908237577019\\
522	-0.0615153731966984\\
523	-0.0910607136046337\\
524	-0.0981224794527407\\
525	-0.0785833643085425\\
526	-0.0361033081550204\\
527	0.01772211294508\\
528	0.0667597982608186\\
529	0.0962046588682597\\
530	0.0983281768928919\\
531	0.0753548036009306\\
533	0.000626085797648557\\
534	-0.025238468584007\\
535	-0.0348095789568106\\
536	-0.0307823807505656\\
538	-0.0105248029292397\\
539	-0.00500665257834498\\
541	-0.00237814494096256\\
548	0.00395402607318829\\
551	0.00378767742972741\\
555	0.000331390435803769\\
559	-0.00345628699392364\\
562	-0.00412556638548267\\
565	-0.00237814494096256\\
572	0.00395402607318829\\
575	0.00378767742972741\\
579	0.000331390435803769\\
583	-0.00345628699392364\\
586	-0.00412556638548267\\
589	-0.00237814494096256\\
596	0.00395402607318829\\
599	0.00378767742972741\\
603	0.000331390435803769\\
607	-0.00345628699392364\\
610	-0.00412556638548267\\
613	-0.00237814494096256\\
620	0.00395402607318829\\
623	0.00378767742972741\\
627	0.000331390435803769\\
631	-0.00345628699392364\\
634	-0.00412556638548267\\
637	-0.00237814494096256\\
644	0.00395402607318829\\
647	0.00378767742972741\\
651	0.000331390435803769\\
655	-0.00345628699392364\\
658	-0.00412556638548267\\
661	-0.00237814494096256\\
668	0.00395402607318829\\
671	0.00378767742972741\\
675	0.000331390435803769\\
679	-0.00345628699392364\\
682	-0.00412556638548267\\
685	-0.00237814494096256\\
692	0.00395402607318829\\
695	0.00378767742972741\\
699	0.000331390435803769\\
703	-0.00345628699392364\\
706	-0.00412556638548267\\
709	-0.00237814494096256\\
716	0.00395402607318829\\
719	0.00378767742972741\\
723	0.000331390435803769\\
727	-0.00345628699392364\\
730	-0.00412556638548267\\
733	-0.00237814494096256\\
740	0.00395402607318829\\
743	0.00378767742972741\\
747	0.000331390435803769\\
751	-0.00345628699392364\\
754	-0.00412556638548267\\
757	-0.00237814494096256\\
764	0.00395402607318829\\
767	0.00378767742972741\\
771	0.000331390435803769\\
775	-0.00345628699392364\\
778	-0.00412556638548267\\
781	-0.00237814494096256\\
788	0.00395402607318829\\
791	0.00378767742972741\\
795	0.000331390435803769\\
799	-0.00345628699392364\\
800	-0.00395402607318829\\
};
\addlegendentry{0}

\end{axis}
\end{tikzpicture}%

\subsection{Détection d'énergie}


Avec ces signaux ainsi filtrés, nous désirons reconstituer le signal de base. Pour cela nous allons utiliser un détecteur d'énergie.
Nous divisons nos signaux en périodes $T_s$ et sur chaque période nous calculons l'énergie suivant la formule suivante:
\[
E=\sum_{i=1}^{N_s} x_n^2
\]
Enfin, on compare cette énergie à un seuil $K$ qu'on fixera à la moyenne des énergies du signal.
Pour le signal en sortie du passe bas par exemple, si $E>K$ alors le signal reconstitué sera égal à 1 sur cette période $T_s$, sinon il sera égal à 0.

Voici les figures obtenues grâce à cette méthode.
%%%% This file was created by matlab2tikz.
%
%The latest updates can be retrieved from
%  http://www.mathworks.com/matlabcentral/fileexchange/22022-matlab2tikz-matlab2tikz
%where you can also make suggestions and rate matlab2tikz.
%
\definecolor{mycolor1}{rgb}{0.00000,0.44700,0.74100}%
%
\begin{tikzpicture}

\begin{axis}[%
width=4.521in,
height=3.559in,
at={(0.758in,0.488in)},
scale only axis,
xmin=0,
xmax=4000,
xlabel style={font=\color{white!15!black}},
xlabel={temps [s]},
ymin=-0.1,
ymax=1.1,
ylabel style={font=\color{white!15!black}},
ylabel={$\text{\$m}_\text{i}\text{(t)\$}$},
axis background/.style={fill=white},
title style={font=\bfseries},
title={Signal NRZ aléatoire}
]
\addplot [color=mycolor1, forget plot]
  table[row sep=crcr]{%
1	0\\
160	0\\
161	1\\
480	1\\
481	0\\
800	0\\
801	1\\
1280	1\\
1281	0\\
1440	0\\
1441	1\\
1760	1\\
1761	0\\
1920	0\\
1921	1\\
2080	1\\
2081	0\\
2400	0\\
2401	1\\
3040	1\\
3041	0\\
3200	0\\
3201	1\\
4000	1\\
};
\end{axis}
\end{tikzpicture}%

%%%% This file was created by matlab2tikz.
%
%The latest updates can be retrieved from
%  http://www.mathworks.com/matlabcentral/fileexchange/22022-matlab2tikz-matlab2tikz
%where you can also make suggestions and rate matlab2tikz.
%
\definecolor{mycolor1}{rgb}{0.00000,0.44700,0.74100}%
%
\begin{tikzpicture}

\begin{axis}[%
width=4.521in,
height=3.566in,
at={(0.758in,0.481in)},
scale only axis,
xmin=0,
xmax=4000,
ymin=-0.1,
ymax=1.1,
axis background/.style={fill=white}
]
\addplot [color=mycolor1, forget plot]
  table[row sep=crcr]{%
1	0\\
250	0\\
251	1\\
575	1\\
576	0\\
900	0\\
901	1\\
1375	1\\
1376	0\\
1550	0\\
1551	1\\
1850	1\\
1851	0\\
2025	0\\
2026	1\\
2175	1\\
2176	0\\
2500	0\\
2501	1\\
3125	1\\
3126	0\\
3300	0\\
3301	1\\
4000	1\\
};
\end{axis}
\end{tikzpicture}%

On peut désormais calculer le taux d'erreur binaire $\eta$ correspondant au nombre de bits mal transmits sur le nombre de bits total.
Avec $F_1=980Hz$ et $F_0=1180Hz$ on obtient un taux d'erreur binaire $\eta=9\%$
\subsection{Modification du démodulateur} 

\subsubsection{Modification du nombre de coefficients}

On remarque qu'avec $F_1=980Hz$ et $F_0=1180Hz$, le nombre de coefficient n'a que très peu d'influence sur le taux d'erreur binaire.

%%%%%%%%%%%%%%%%%%%%%%%%%% A CHECK %%%%%%%%%%%%%%%%%%%%%%%%%%%%%%%%%%%%%%%%%%%%%%%%%%%%%%%%%%%%%%%%%%%%%%%%%%%%%%%%%%%%%%%%%%%%%%%%%%
En passant à 201 coefficients, le taux d'erreur reste de 0.2662.
En passant à 61 coefficients, le taux d'erreur reste de 0.0725. 
% V21 : 0.0912 pour les deux quantités de coefficients

%%%% This file was created by matlab2tikz.
%
%The latest updates can be retrieved from
%  http://www.mathworks.com/matlabcentral/fileexchange/22022-matlab2tikz-matlab2tikz
%where you can also make suggestions and rate matlab2tikz.
%
\definecolor{mycolor1}{rgb}{0.00000,0.44700,0.74100}%
%
\begin{tikzpicture}

\begin{axis}[%
width=4.521in,
height=3.566in,
at={(0.758in,0.481in)},
scale only axis,
xmin=0,
xmax=4000,
ymin=-0.1,
ymax=1.1,
axis background/.style={fill=white}
]
\addplot [color=mycolor1, forget plot]
  table[row sep=crcr]{%
1	0\\
250	0\\
251	1\\
575	1\\
576	0\\
900	0\\
901	1\\
1375	1\\
1376	0\\
1550	0\\
1551	1\\
1850	1\\
1851	0\\
2025	0\\
2026	1\\
2175	1\\
2176	0\\
2500	0\\
2501	1\\
3125	1\\
3126	0\\
3300	0\\
3301	1\\
4000	1\\
};
\end{axis}
\end{tikzpicture}%

%%%% This file was created by matlab2tikz.
%
%The latest updates can be retrieved from
%  http://www.mathworks.com/matlabcentral/fileexchange/22022-matlab2tikz-matlab2tikz
%where you can also make suggestions and rate matlab2tikz.
%
\definecolor{mycolor1}{rgb}{0.00000,0.44700,0.74100}%
\definecolor{mycolor2}{rgb}{0.85000,0.32500,0.09800}%
%
\begin{tikzpicture}

\begin{axis}[%
width=4.521in,
height=3.548in,
at={(0.758in,0.499in)},
scale only axis,
unbounded coords=jump,
xmin=-25000,
xmax=25000,
ymode=log,
ymin=1e-08,
ymax=100,
yminorticks=true,
axis background/.style={fill=white}
]
\addplot [color=mycolor1, forget plot]
  table[row sep=crcr]{%
-24000	1.00284134051483\\
-9707.23965327799	1.00784668917116\\
-9443.53558784031	1.00634391250713\\
-9086.07007691369	1.00219220929594\\
-8664.1435722134	1.00853005589454\\
-8294.95788060066	1.0023857812845\\
-7878.89146624344	1.011376305327\\
-7509.70577463069	1.00240781529685\\
-7081.91917958735	1.01506004269969\\
-6730.31375900378	1.00217791442368\\
-6267.36662190209	1.01974901170664\\
-5968.50201440605	1.00096663918285\\
-5821.99975582957	1.01714225808576\\
-5611.03650347943	1.04099858046966\\
-5487.97460627518	1.03356722618652\\
-5177.38981809303	1.0013161485488\\
-5077.76828226102	1.01832174822303\\
-4802.34403613722	1.08121372328696\\
-4720.30277133439	1.07315522551496\\
-4626.54132584544	1.04486686922427\\
-4468.31888658283	0.996003132367466\\
-4409.71798315224	0.997418650516003\\
-4356.97717006471	1.0152401761057\\
-4304.23635697717	1.05120091934154\\
-4245.63545354658	1.11272091443234\\
-4169.4542790868	1.22121006145484\\
-4040.53229153949	1.43112405022432\\
-3981.9313881089	1.50941537509066\\
-3929.19057502137	1.55563075140529\\
-3882.30985227689	1.57034780457899\\
-3835.42912953241	1.55547668029589\\
-3794.408497131	1.515929511999\\
-3753.38786472958	1.45093497498028\\
-3718.22732267122	1.37546094305573\\
-3683.06678061287	1.28298625797261\\
-3647.90623855451	1.17540319674613\\
-3618.60578683921	1.07600980623875\\
-3589.30533512392	0.969613088671864\\
-3560.00488340862	0.858179393470407\\
-3536.56452203638	0.766835138670634\\
-3513.12416066415	0.674776579192027\\
-3489.68379929191	0.583176467998959\\
-3472.10352826273	0.51548717129855\\
-3454.52325723355	0.449213427234115\\
-3436.94298620437	0.384840462855972\\
-3419.36271517519	0.322840402176546\\
-3407.64253448907	0.283050703579549\\
-3395.92235380295	0.244647760125015\\
-3384.20217311684	0.20775640496713\\
-3372.48199243072	0.172496231612077\\
-3360.7618117446	0.138981138678295\\
-3354.90172140154	0.122911928192072\\
-3349.04163105848	0.107318904642378\\
-3343.18154071542	0.0922145195238283\\
-3337.32145037236	0.0776107944725708\\
-3331.4613600293	0.0635193108314036\\
-3325.60126968624	0.049951199770272\\
-3319.74117934318	0.0369171329732436\\
-3313.88108900012	0.0244273139028829\\
-3308.02099865706	0.0124914696441364\\
-3302.160908314	0.00111884333889737\\
-3296.30081797094	-0.00968181278675526\\
-3290.44072762788	-0.019902243794301\\
-3284.58063728482	-0.0295346978027968\\
-3278.72054694176	-0.0385719311511417\\
-3272.8604565987	-0.0470072129547623\\
-3267.00036625564	-0.0548343290500605\\
-3261.14027591259	-0.0620475853296203\\
-3255.28018556953	-0.068641810471287\\
-3249.42009522647	-0.0746123580559206\\
-3243.56000488341	-0.0799551080805467\\
-3237.69991454035	-0.0846664678712921\\
-3231.83982419729	-0.0887433723951299\\
-3220.11964351117	-0.0949841914586418\\
-3208.39946282505	-0.0986635725869213\\
-3196.67928213893	-0.0997758869663738\\
-3184.95910145281	-0.0983237780917271\\
-3173.2389207667	-0.0943180352165872\\
-3161.51874008058	-0.0877774331808311\\
-3155.65864973752	-0.083564563990025\\
-3149.79855939446	-0.0787285401467724\\
-3143.9384690514	-0.0732738638425516\\
-3138.07837870834	-0.0672054949415363\\
-3132.21828836528	-0.0605288424918328\\
-3126.35819802222	-0.0532497558352813\\
-3120.49810767916	-0.0453745153341191\\
-3114.6380173361	-0.036909822733815\\
-3108.77792699304	-0.0278627911743993\\
-3102.91783664998	-0.0182409348719961\\
-3097.05774630692	-0.00805215848653547\\
-3091.19765596386	0.00269525380486109\\
-3085.3375656208	0.0139926495122191\\
-3079.47747527774	0.0258310192938315\\
-3073.61738493469	0.0382010086571355\\
-3067.75729459163	0.051092929879601\\
-3061.89720424857	0.0644967741307062\\
-3056.03711390551	0.0784022237753743\\
-3050.17702356245	0.0927986648406104\\
-3044.31693321939	0.107675199625698\\
-3038.45684287633	0.123020659436337\\
-3032.59675253327	0.138823617423142\\
-3020.87657184715	0.171755107382671\\
-3009.15639116103	0.206373584339024\\
-2997.43621047491	0.242579665254818\\
-2985.71602978879	0.280271087478488\\
-2973.99584910267	0.319343125452807\\
-2956.41557807349	0.380305805263977\\
-2938.83530704432	0.443767362927457\\
-2921.25503601514	0.509355604286703\\
-2897.8146746429	0.599467263808462\\
-2874.37431327066	0.691807769167712\\
-2850.93395189842	0.785502690256061\\
-2821.63350018313	0.903228572744357\\
-2792.33304846783	1.02013891276253\\
-2757.17250640948	1.15723788805243\\
-2722.01196435112	1.2885943760245\\
-2686.85142229276	1.41205386257115\\
-2645.83078989134	1.54359360095767\\
-2604.81015748993	1.65924664745704\\
-2557.92943474545	1.76928374103485\\
-2511.04871200098	1.85336415918948\\
-2464.1679892565	1.90968229154065\\
-2411.42717616897	1.93838352068749\\
-2358.68636308143	1.92960648867852\\
-2311.80564033696	1.89039774211996\\
-2264.92491759248	1.8223753875837\\
-2218.044194848	1.72697563229736\\
-2171.16347210353	1.6064104532494\\
-2130.14283970211	1.48264976897469\\
-2089.1222073007	1.34453613945666\\
-2053.96166524234	1.21706267650456\\
-2018.80112318398	1.08353685194144\\
-1983.64058112563	0.946409026532242\\
-1954.34012941033	0.831279261727773\\
-1925.03967769503	0.71706815215248\\
-1901.59931632279	0.62745598313109\\
-1878.15895495056	0.54031517120234\\
-1854.71859357832	0.456479003249757\\
-1837.13832254914	0.396262514792276\\
-1819.55805151996	0.338699265508007\\
-1801.97778049078	0.284112704628335\\
-1784.3975094616	0.232812600360447\\
-1772.67732877548	0.200584104372002\\
-1760.95714808936	0.17002973578995\\
-1749.23696740325	0.141228143238201\\
-1737.51678671713	0.114253773031927\\
-1731.65669637407	0.101473846214038\\
-1725.79660603101	0.089176604251686\\
-1719.93651568795	0.0773700555255055\\
-1714.07642534489	0.066061898860856\\
-1708.21633500183	0.0552595166417749\\
-1702.35624465877	0.0449699681951868\\
-1696.49615431571	0.0351999834814129\\
-1690.63606397265	0.0259559570753332\\
-1684.77597362959	0.0172439424659039\\
-1678.91588328653	0.00906964666595916\\
-1673.05579294348	0.00143842514990056\\
-1667.19570260042	-0.0056447228804401\\
-1661.33561225736	-0.0121751588931026\\
-1655.4755219143	-0.0181486090815703\\
-1649.61543157124	-0.0235611680734111\\
-1643.75534122818	-0.0284093022881394\\
-1637.89525088512	-0.0326898529396692\\
-1632.03516054206	-0.03640003867962\\
-1626.175070199	-0.0395374578802563\\
-1620.31497985594	-0.0421000905512338\\
-1614.45488951288	-0.0440862998926634\\
-1608.59479916982	-0.0454948334797297\\
-1602.73470882676	-0.0463248240782844\\
-1596.8746184837	-0.0465757900940885\\
-1591.01452814064	-0.0462476356531751\\
-1585.15443779758	-0.0453406503121169\\
-1579.29434745452	-0.0438555084071351\\
-1573.43425711146	-0.0417932680339995\\
-1567.57416676841	-0.0391553696685281\\
-1561.71407642535	-0.0359436344278396\\
-1555.85398608229	-0.0321602619750341\\
-1549.99389573923	-0.0278078280714245\\
-1544.13380539617	-0.022889281782841\\
-1538.27371505311	-0.0174079423414557\\
-1532.41362471005	-0.0113674956705275\\
-1526.55353436699	-0.00477199057791364\\
-1520.69344402393	0.00237416537717297\\
-1514.83335368087	0.0100662103354489\\
-1508.97326333781	0.018299032901186\\
-1503.11317299475	0.0270671772101309\\
-1497.25308265169	0.036364848295125\\
-1491.39299230863	0.0461859177382354\\
-1485.53290196557	0.0565239296041517\\
-1479.67281162251	0.0673721066447664\\
-1473.81272127945	0.0787233567635679\\
-1467.95263093639	0.0905702797341611\\
-1462.09254059333	0.10290517415985\\
-1450.37235990721	0.129006609313262\\
-1438.6521792211	0.156960306392116\\
-1426.93199853498	0.186694573222286\\
-1415.21181784886	0.218133640395567\\
-1403.49163716274	0.251197931015164\\
-1385.91136613356	0.30365897324134\\
-1368.33109510438	0.359295171602081\\
-1350.7508240752	0.417798500253555\\
-1333.17055304603	0.478849030566786\\
-1309.73019167379	0.563641378644391\\
-1286.28983030155	0.651575047520894\\
-1262.84946892931	0.741835985968858\\
-1233.54901721402	0.856688017009298\\
-1204.24856549872	0.97230985832773\\
-1169.08802344036	1.1098882741476\\
-1133.927481382	1.24376503761306\\
-1098.76693932365	1.37152395594126\\
-1057.74630692223	1.50989246662556\\
-1016.72567452081	1.63377494418881\\
-969.84495177634	1.75417409627901\\
-922.964229031862	1.84885086771245\\
-876.083506287388	1.91531165235299\\
-829.202783542914	1.95185225640328\\
-782.322060798437	1.95754641903256\\
-735.441338053963	1.93223937442644\\
-688.560615309485	1.87655128693148\\
-641.679892565011	1.79189033012841\\
-594.799169820533	1.68047018598328\\
-553.778537419119	1.56338990081338\\
-512.757905017701	1.43067408083479\\
-471.737272616287	1.28532719492052\\
-436.576730557928	1.15328171802749\\
-401.416188499574	1.01685311452385\\
-372.115736784275	0.90164800901762\\
-342.815285068977	0.786705924625891\\
-313.514833353682	0.673604005400197\\
-290.074471981443	0.58551619082351\\
-266.634110609204	0.500433436681232\\
-243.193749236969	0.419150848011362\\
-225.61347820779	0.361148050051534\\
-208.033207178611	0.306031303990695\\
-190.452936149431	0.254101786166131\\
-178.732755463312	0.221394895211982\\
-167.012574777196	0.190312965254502\\
-155.292394091077	0.160933171753497\\
-143.572213404957	0.133328829423154\\
-131.852032718838	0.107569144067225\\
-125.991942375778	0.0954014925732736\\
-120.131852032719	0.0837189760648634\\
-114.271761689659	0.0725289864276536\\
-108.411671346599	0.0618386167720763\\
-102.551581003539	0.0516546552306847\\
-96.6914906604798	0.0419835790077327\\
-90.83140031742	0.0328315486793118\\
-84.9713099743603	0.0242044027596892\\
-79.1112196313043	0.0161076525315924\\
-73.2511292882446	0.00854647715867775\\
-67.3910389451848	0.00152571907700035\\
-61.5309486021251	-0.00495012032126534\\
-55.6708582590654	-0.0108768847103271\\
-49.8107679160057	-0.0162507665513866\\
-43.950677572946	-0.021068310507595\\
-38.0905872298863	-0.0253264165483282\\
-32.2304968868266	-0.0290223427374395\\
-26.3704065437669	-0.032153707704045\\
-20.5103162007072	-0.0347184927898288\\
-14.6502258576475	-0.0367150438710176\\
-8.79013551458775	-0.0381420728534124\\
-2.93004517152804	-0.0389986588368639\\
2.93004517152804	-0.039284248946786\\
8.79013551458775	-0.0389986588368639\\
14.6502258576475	-0.0381420728534124\\
20.5103162007072	-0.0367150438710176\\
26.3704065437669	-0.0347184927898288\\
32.2304968868266	-0.032153707704045\\
38.0905872298863	-0.0290223427374395\\
43.950677572946	-0.0253264165483282\\
49.8107679160057	-0.021068310507595\\
55.6708582590654	-0.0162507665513866\\
61.5309486021251	-0.0108768847103271\\
67.3910389451848	-0.00495012032126534\\
73.2511292882446	0.00152571907700035\\
79.1112196313043	0.00854647715867775\\
84.9713099743603	0.0161076525315924\\
90.83140031742	0.0242044027596892\\
96.6914906604798	0.0328315486793118\\
102.551581003539	0.0419835790077327\\
108.411671346599	0.0516546552306847\\
114.271761689659	0.0618386167720763\\
120.131852032719	0.0725289864276536\\
125.991942375778	0.0837189760648634\\
131.852032718838	0.0954014925732736\\
137.712123061898	0.107569144067225\\
149.432303748017	0.133328829423154\\
161.152484434137	0.160933171753497\\
172.872665120252	0.190312965254502\\
184.592845806372	0.221394895211982\\
196.313026492491	0.254101786166131\\
213.89329752167	0.306031303990695\\
231.473568550849	0.361148050051534\\
249.053839580029	0.419150848011362\\
272.494200952264	0.500433436681232\\
295.934562324503	0.58551619082351\\
319.374923696741	0.673604005400197\\
348.675375412036	0.786705924625891\\
377.975827127335	0.90164800901762\\
407.276278842633	1.01685311452385\\
442.436820900988	1.15328171802749\\
477.597362959346	1.28532719492052\\
518.617995360761	1.43067408083479\\
559.638627762179	1.56338990081338\\
600.659260163593	1.68047018598328\\
647.539982908071	1.79189033012841\\
694.420705652545	1.87655128693148\\
741.301428397022	1.93223937442644\\
788.182151141496	1.95754641903256\\
835.062873885974	1.95185225640328\\
881.943596630448	1.91531165235299\\
928.824319374922	1.84885086771245\\
975.7050421194	1.75417409627901\\
1022.58576486387	1.63377494418881\\
1063.60639726529	1.50989246662556\\
1104.62702966671	1.37152395594126\\
1145.64766206812	1.22181408965479\\
1180.80820412648	1.08714779382094\\
1215.96874618484	0.949199155478187\\
1245.26919790014	0.833605562337403\\
1274.56964961543	0.719097343665334\\
1298.01001098767	0.62934150291742\\
1321.45037235991	0.54211722323692\\
1344.89073373214	0.458236000732144\\
1362.47100476132	0.397998349569571\\
1380.0512757905	0.340415724088584\\
1397.63154681968	0.285804339612201\\
1415.21181784886	0.234467747490499\\
1426.93199853498	0.202205894311817\\
1438.6521792211	0.171609511729068\\
1450.37235990721	0.142756307220438\\
1462.09254059333	0.115720044667061\\
1467.95263093639	0.10290517415985\\
1473.81272127945	0.0905702797341611\\
1479.67281162251	0.0787233567635679\\
1485.53290196557	0.0673721066447664\\
1491.39299230863	0.0565239296041517\\
1497.25308265169	0.0461859177382354\\
1503.11317299475	0.036364848295125\\
1508.97326333781	0.0270671772101309\\
1514.83335368087	0.018299032901186\\
1520.69344402393	0.0100662103354489\\
1526.55353436699	0.00237416537717297\\
1532.41362471005	-0.00477199057791364\\
1538.27371505311	-0.0113674956705275\\
1544.13380539617	-0.0174079423414557\\
1549.99389573923	-0.022889281782841\\
1555.85398608229	-0.0278078280714245\\
1561.71407642535	-0.0321602619750341\\
1567.57416676841	-0.0359436344278396\\
1573.43425711146	-0.0391553696685281\\
1579.29434745452	-0.0417932680339995\\
1585.15443779758	-0.0438555084071351\\
1591.01452814064	-0.0453406503121169\\
1596.8746184837	-0.0462476356531751\\
1602.73470882676	-0.0465757900940885\\
1608.59479916982	-0.0463248240782844\\
1614.45488951288	-0.0454948334797297\\
1620.31497985594	-0.0440862998926634\\
1626.175070199	-0.0421000905512338\\
1632.03516054206	-0.0395374578802563\\
1637.89525088512	-0.03640003867962\\
1643.75534122818	-0.0326898529396692\\
1649.61543157124	-0.0284093022881394\\
1655.4755219143	-0.0235611680734111\\
1661.33561225736	-0.0181486090815703\\
1667.19570260042	-0.0121751588931026\\
1673.05579294348	-0.0056447228804401\\
1678.91588328653	0.00143842514990056\\
1684.77597362959	0.00906964666595916\\
1690.63606397265	0.0172439424659039\\
1696.49615431571	0.0259559570753332\\
1702.35624465877	0.0351999834814129\\
1708.21633500183	0.0449699681951868\\
1714.07642534489	0.0552595166417749\\
1719.93651568795	0.066061898860856\\
1725.79660603101	0.0773700555255055\\
1731.65669637407	0.089176604251686\\
1737.51678671713	0.101473846214038\\
1743.37687706019	0.114253773031927\\
1755.09705774631	0.141228143238201\\
1766.81723843242	0.17002973578995\\
1778.53741911854	0.200584104372002\\
1790.25759980466	0.232812600360447\\
1801.97778049078	0.26663265139941\\
1819.55805151996	0.320157400015862\\
1837.13832254914	0.376763733173268\\
1854.71859357832	0.43613249110897\\
1878.15895495056	0.519014062109301\\
1901.59931632279	0.605406347429255\\
1925.03967769503	0.694479336996092\\
1954.34012941033	0.808311718652108\\
1983.64058112563	0.923386843517411\\
2012.94103284092	1.03810107619541\\
2048.10157489928	1.17310981081574\\
2083.26211695764	1.30285987535807\\
2124.28274935905	1.44452721005051\\
2165.30338176047	1.57267556300461\\
2206.32401416188	1.68456770292539\\
2253.20473690636	1.78970764655892\\
2300.08545965084	1.86820488191576\\
2346.96618239531	1.91834094544438\\
2399.70699548285	1.93963704193348\\
2452.44780857038	1.92338067716478\\
2499.32853131486	1.87781407428582\\
2546.20925405933	1.8039664386101\\
2593.08997680381	1.70341094936495\\
2639.97069954828	1.57835899148366\\
2680.9913319497	1.4511168886958\\
2722.01196435112	1.30977379476686\\
2757.17250640948	1.17959075230625\\
2792.33304846783	1.04328853805076\\
2821.63350018313	0.926725599194983\\
2850.93395189842	0.809037455065404\\
2874.37431327066	0.715138286231451\\
2897.8146746429	0.622377867674489\\
2921.25503601514	0.531626214573054\\
2938.83530704432	0.465412288772645\\
2956.41557807349	0.401200303594941\\
2973.99584910267	0.339363576894535\\
2985.71602978879	0.299641159725105\\
2997.43621047491	0.261246201189658\\
3009.15639116103	0.224284503124075\\
3020.87657184715	0.188859611535734\\
3032.59675253327	0.155072401510128\\
3038.45684287633	0.138823617423142\\
3044.31693321939	0.123020659436337\\
3050.17702356245	0.107675199625698\\
3056.03711390551	0.0927986648406104\\
3061.89720424857	0.0784022237753743\\
3067.75729459163	0.0644967741307062\\
3073.61738493469	0.051092929879601\\
3079.47747527774	0.0382010086571355\\
3085.3375656208	0.0258310192938315\\
3091.19765596386	0.0139926495122191\\
3097.05774630692	0.00269525380486109\\
3102.91783664998	-0.00805215848653547\\
3108.77792699304	-0.0182409348719961\\
3114.6380173361	-0.0278627911743993\\
3120.49810767916	-0.036909822733815\\
3126.35819802222	-0.0453745153341191\\
3132.21828836528	-0.0532497558352813\\
3138.07837870834	-0.0605288424918328\\
3143.9384690514	-0.0672054949415363\\
3149.79855939446	-0.0732738638425516\\
3155.65864973752	-0.0787285401467724\\
3161.51874008058	-0.083564563990025\\
3167.37883042364	-0.0877774331808311\\
3179.09901110976	-0.0943180352165872\\
3190.81919179587	-0.0983237780917271\\
3202.53937248199	-0.0997758869663738\\
3214.25955316811	-0.0986635725869213\\
3225.97973385423	-0.0949841914586418\\
3237.69991454035	-0.0887433723951299\\
3243.56000488341	-0.0846664678712921\\
3249.42009522647	-0.0799551080805467\\
3255.28018556953	-0.0746123580559206\\
3261.14027591259	-0.068641810471287\\
3267.00036625564	-0.0620475853296203\\
3272.8604565987	-0.0548343290500605\\
3278.72054694176	-0.0470072129547623\\
3284.58063728482	-0.0385719311511417\\
3290.44072762788	-0.0295346978027968\\
3296.30081797094	-0.019902243794301\\
3302.160908314	-0.00968181278675526\\
3308.02099865706	0.00111884333889737\\
3313.88108900012	0.0124914696441364\\
3319.74117934318	0.0244273139028829\\
3325.60126968624	0.0369171329732436\\
3331.4613600293	0.049951199770272\\
3337.32145037236	0.0635193108314036\\
3343.18154071542	0.0776107944725708\\
3349.04163105848	0.0922145195238283\\
3354.90172140154	0.107318904642378\\
3360.7618117446	0.122911928192072\\
3366.62190208766	0.138981138678295\\
3378.34208277378	0.172496231612077\\
3390.06226345989	0.20775640496713\\
3401.78244414601	0.244647760125015\\
3413.50262483213	0.283050703579549\\
3425.22280551825	0.322840402176546\\
3442.80307654743	0.384840462855972\\
3460.38334757661	0.449213427234115\\
3477.96361860579	0.51548717129855\\
3501.40397997803	0.605968936689599\\
3524.84434135026	0.697794319707401\\
3548.2847027225	0.789783546201827\\
3577.5851544378	0.903223892767596\\
3606.88560615309	1.01289063691337\\
3636.18605786839	1.11671087935326\\
3671.34659992675	1.23093964877382\\
3706.5071419851	1.33124592697357\\
3741.66768404346	1.41541503731414\\
3782.68831644488	1.4911806242377\\
3823.7089488463	1.5416340264825\\
3864.72958124771	1.56683176098259\\
3911.61030399219	1.56643960107551\\
3958.49102673666	1.53854026129323\\
4011.2318398242	1.4808441855122\\
4069.83274325479	1.3946357643897\\
4157.73409840068	1.2494405716888\\
4257.3556342327	1.10560367653108\\
4315.95653766329	1.04628045551816\\
4368.69735075082	1.01239104258024\\
4421.43816383836	0.996473557808924\\
4480.03906726895	0.996743029888078\\
4556.22024172873	1.01643403533274\\
4749.60322304969	1.07729493343014\\
4831.64448785252	1.08019849875199\\
4925.40593334147	1.06374336103988\\
5212.55036015138	0.999660690850909\\
5318.03198632646	1.00562490882278\\
5640.33695519473	1.04025027433413\\
5780.97912342815	1.02442129157638\\
5997.80246612135	1.00012251061262\\
6138.44463435478	1.0058972902668\\
6437.30924185081	1.02581498795457\\
6618.97204248565	1.01181235050572\\
6824.07520449274	0.999955278799365\\
7011.59809547064	1.00971356554002\\
7240.14161884996	1.01839349847716\\
7480.40532291539	1.00430903784712\\
7679.64839457942	1.00099296556962\\
8183.61616408253	1.00837668572661\\
8476.6206812355	1.00075409056754\\
9015.74899279697	1.00527477136276\\
9302.89341960689	1.00124068173093\\
9783.42082773776	1.00553615605415\\
10129.1661579783	1.00170145932106\\
10574.5330240508	1.00496281748533\\
10955.4388963497	1.00211668208415\\
11377.3654010499	1.00421894135086\\
11775.851544378	1.00233799705817\\
12186.0578683921	1.00350960004946\\
12596.2641924063	1.00252579591822\\
13006.4705164205	1.00270184819896\\
13416.6768404346	1.00268272959236\\
13826.8831644488	1.00205310528284\\
14242.949578806	1.0029115097727\\
14659.0159931632	1.00133875633935\\
15086.8025882066	1.00332738854949\\
15532.1694542791	1.00034476341141\\
16516.6646319131	1.00064416803259\\
17249.1759247955	1.00009481476089\\
19399.8290806983	1.00165341641159\\
19991.6982053473	1.00305185813636\\
24000	1.00283983599301\\
};
\addplot [color=mycolor2, forget plot]
  table[row sep=crcr]{%
-24000	-0.00284134051541254\\
-23976.5596386278	-0.00281733185141145\\
-23953.1192772555	-0.00274611709168084\\
-23929.6789158833	-0.00263010254264651\\
-23906.238554511	-0.00247320833745444\\
-23882.7981931388	-0.00228073607346066\\
-23859.3578317666	-0.00205918980113674\\
-23841.7775607374	-0.00187842412257593\\
-23824.1972897082	-0.00168894126064206\\
-23806.617018679	-0.00149434793701714\\
-23789.0367476499	-0.0012983483943787\\
-23771.4564766207	-0.00110467394012934\\
-23759.7362959346	-0.000978674225344446\\
-23748.0161152484	-0.000856416311212569\\
-23736.2959345623	-0.000738935829724721\\
-23724.5757538762	-0.000627228047288775\\
-23712.8555731901	-0.000522239440905372\\
-23701.135392504	-0.000424859686585813\\
-23689.4152118178	-0.000335914127787074\\
-23683.5551214748	-0.000294842768824362\\
-23677.6950311317	-0.000256156787581085\\
-23671.8349407887	-0.000219938245436648\\
-23665.9748504456	-0.000186263983670552\\
-23660.1147601026	-0.00015520546063389\\
-23654.2546697595	-0.000126828600321508\\
-23648.3945794164	-0.000101193652633783\\
-23642.5344890734	-7.83550656602712e-05\\
-23636.6743987303	-5.83613702124949e-05\\
-23630.8143083873	-4.12550768870749e-05\\
-23624.9542180442	-2.7072585857694e-05\\
-23619.0941277011	-1.58441095928798e-05\\
-23613.2340373581	-7.59360866299011e-06\\
-23607.373947015	-2.33874077388739e-06\\
-23601.513856672	-9.0823133571775e-08\\
-23595.6537663289	-8.54808233626689e-07\\
-23589.7936759858	-4.62927309620972e-06\\
-23583.9335856428	-1.14064220109323e-05\\
-23578.0734952997	-2.1172102755075e-05\\
-23572.2134049567	-3.39058362646826e-05\\
-23566.3533146136	-4.95808596941853e-05\\
-23560.4932242705	-6.81641827723735e-05\\
-23554.6331339275	-8.96166573390238e-05\\
-23548.7730435844	-0.000113893059913721\\
-23542.9129532414	-0.000140942187121089\\
-23537.0528628983	-0.000170706963778393\\
-23531.1927725552	-0.000203124563407944\\
-23525.3326822122	-0.000238126540927679\\
-23519.4725918691	-0.000275638977228373\\
-23513.6125015261	-0.000315582635360732\\
-23501.8923208399	-0.000402421095557273\\
-23490.1721401538	-0.000497908301027104\\
-23478.4519594677	-0.000601237361657966\\
-23466.7317787816	-0.000711534937569924\\
-23455.0115980955	-0.000827868621857204\\
-23443.2914174093	-0.000949254823742937\\
-23425.7111463802	-0.00113855262994018\\
-23408.130875351	-0.00133330200090315\\
-23390.5506043218	-0.00152979955553598\\
-23372.9703332926	-0.00172430801620584\\
-23355.3900622635	-0.00191312729735072\\
-23337.8097912343	-0.00209266490066311\\
-23314.369429862	-0.00231173858207379\\
-23290.9290684898	-0.00250086139780925\\
-23267.4887071176	-0.002653639369865\\
-23244.0483457453	-0.00276490487818962\\
-23220.6079843731	-0.00283089160029641\\
-23197.1676230009	-0.00284936218507815\\
-23173.7272616286	-0.00281968433563411\\
-23150.2869002564	-0.00274285271512935\\
-23126.8465388841	-0.0026214559182546\\
-23103.4061775119	-0.00245958960357174\\
-23079.9658161397	-0.00226271870021161\\
-23056.5254547674	-0.00203749332194539\\
-23038.9451837382	-0.00185451617663497\\
-23021.3649127091	-0.00166334005929498\\
-23003.7846416799	-0.00146760582281903\\
-22986.2043706507	-0.00127104229099939\\
-22968.6240996215	-0.00107739534806444\\
-22956.9039189354	-0.000951745131603513\\
-22945.1837382493	-0.000830100429676494\\
-22933.4635575632	-0.000713493527961728\\
-22921.7433768771	-0.000602914317891804\\
-22910.0231961909	-0.00049930191015217\\
-22898.3030155048	-0.000403536675847783\\
-22886.5828348187	-0.000316432782750981\\
-22880.7227444756	-0.000276363452316818\\
-22874.8626541326	-0.000238731289862733\\
-22869.0025637895	-0.000203616472647206\\
-22863.1424734465	-0.000171093858826122\\
-22857.2823831034	-0.000141232828365827\\
-22851.4222927603	-0.000114097135518203\\
-22845.5622024173	-8.97447731970047e-05\\
-22839.7021120742	-6.82278495082941e-05\\
-22833.8420217312	-4.95924767370388e-05\\
-22827.9819313881	-3.38786729934728e-05\\
-22822.121841045	-2.11202767477588e-05\\
-22816.261750702	-1.13448744295206e-05\\
-22810.4016603589	-4.57374124570659e-06\\
-22804.5415700159	-8.21795347178264e-07\\
-22798.6814796728	-9.75654389382338e-08\\
-22792.8213893298	-2.40317190797747e-06\\
-22786.9612989867	-7.73432150630204e-06\\
-22781.1012086436	-1.60803156061958e-05\\
-22775.2411183006	-2.74240720093062e-05\\
-22769.3810279575	-4.17421602660924e-05\\
-22763.5209376145	-5.90048504323481e-05\\
-22757.6608472714	-7.91761751625779e-05\\
-22751.8007569283	-0.000102214005008151\\
-22745.9406665853	-0.000128070136766963\\
-22740.0805762422	-0.000156690394694511\\
-22734.2204858992	-0.000188014744367536\\
-22728.3603955561	-0.000221977418955338\\
-22722.500305213	-0.000258507057649185\\
-22716.64021487	-0.000297526855936778\\
-22710.7801245269	-0.00033895472741485\\
-22699.0599438408	-0.000428680983837285\\
-22687.3397631547	-0.000526930340305221\\
-22675.6195824686	-0.000632874964801485\\
-22663.8994017824	-0.000745621665077606\\
-22652.1792210963	-0.000864219428830506\\
-22640.4590404102	-0.000987667455107538\\
-22622.878769381	-0.00117964492654475\\
-22605.2984983518	-0.00137653821552113\\
-22587.7182273227	-0.00157460737881394\\
-22570.1379562935	-0.00177008848521255\\
-22552.5576852643	-0.00195926521608474\\
-22534.9774142351	-0.00213853959297041\\
-22511.5370528629	-0.00235630774269005\\
-22488.0966914907	-0.00254307525048042\\
-22464.6563301184	-0.00269252401529647\\
-22441.2159687462	-0.00279959325577728\\
-22417.7756073739	-0.00286065131759589\\
-22394.3352460017	-0.0028736193259035\\
-22370.8948846295	-0.00283804247544993\\
-22347.4545232572	-0.00275510651585241\\
-22324.014161885	-0.00262759883751439\\
-22300.5738005128	-0.00245981543344755\\
-22277.1334391405	-0.00225741684033257\\
-22253.6930777683	-0.00202723788587094\\
-22236.1128067391	-0.0018410241368157\\
-22218.5325357099	-0.0016470921792747\\
-22200.9522646807	-0.00144913758640776\\
-22183.3719936516	-0.00125093471025438\\
-22165.7917226224	-0.00105626491351087\\
-22154.0715419363	-0.000930287966284631\\
-22142.3513612502	-0.000808606329948592\\
-22130.631180564	-0.000692254594489036\\
-22118.9109998779	-0.000582222661009546\\
-22107.1908191918	-0.000479447346117513\\
-22095.4706385057	-0.000384804432001888\\
-22089.6105481626	-0.00034078777638871\\
-22083.7504578196	-0.000299101229714281\\
-22077.8903674765	-0.000259833881528768\\
-22072.0302771334	-0.000223069718647253\\
-22066.1701867904	-0.000188887446328079\\
-22060.3100964473	-0.000157360320569037\\
-22054.4500061043	-0.000128555991917456\\
-22048.5899157612	-0.000102536361090831\\
-22042.7298254181	-7.93574467543084e-05\\
-22036.8697350751	-5.90692656961301e-05\\
-22031.009644732	-4.17157256969613e-05\\
-22025.149554389	-2.73345312889737e-05\\
-22019.2894640459	-1.59571026288878e-05\\
-22013.4293737028	-7.60850765080387e-06\\
-22007.5692833598	-2.30740764265196e-06\\
-22001.7091930167	-6.60163722716593e-08\\
-21995.8491026737	-8.90072845246499e-07\\
-21989.9890123306	-4.77882775783598e-06\\
-21984.1289219875	-1.1725043677168e-05\\
-21978.2688316445	-2.17150089512198e-05\\
-21972.4087413014	-3.47285653200621e-05\\
-21966.5486509584	-5.07391491756901e-05\\
-21960.6885606153	-6.97138463841963e-05\\
-21954.8284702722	-9.16134605594766e-05\\
-21948.9683799292	-0.000116392594641658\\
-21943.1082895861	-0.000143999745620105\\
-21937.2481992431	-0.000174377412194723\\
-21931.3881089	-0.000207462215150681\\
-21925.5280185569	-0.000243185030209084\\
-21919.6679282139	-0.000281471133043174\\
-21913.8078378708	-0.000322240356207789\\
-21902.0876571847	-0.000410881300080587\\
-21890.3674764986	-0.000508364338302483\\
-21878.6472958125	-0.000613870769674534\\
-21866.9271151264	-0.000726513609956042\\
-21855.2069344402	-0.000845345067356821\\
-21843.4867537541	-0.000969364536003562\\
-21825.9064827249	-0.00116282540728234\\
-21808.3262116958	-0.00136193264305908\\
-21790.7459406666	-0.00156290913142642\\
-21773.1656696374	-0.00176193957099966\\
-21755.5853986082	-0.00195524304155938\\
-21738.0051275791	-0.00213914497301313\\
-21714.5647662068	-0.00236371040728484\\
-21691.1244048346	-0.00255777979373768\\
-21667.6840434623	-0.00271478547671585\\
-21644.2436820901	-0.00282940627672859\\
-21620.8033207179	-0.00289774855537648\\
-21597.3629593456	-0.00291747924063\\
-21573.9225979734	-0.00288790627088031\\
-21550.4822366011	-0.00281000368688873\\
-21527.0418752289	-0.00268638046715559\\
-21503.6015138567	-0.00252119409978788\\
-21480.1611524844	-0.00232001175123109\\
-21456.7207911122	-0.0020896236641936\\
-21439.140520083	-0.00190231942998053\\
-21421.5602490538	-0.00170651903154536\\
-21403.9799780247	-0.00150595524508733\\
-21386.3997069955	-0.00130445474806615\\
-21368.8194359663	-0.00110586537972965\\
-21357.0992552802	-0.000976968564233929\\
-21345.3790745941	-0.000852152069549209\\
-21333.6588939079	-0.000732478899955249\\
-21321.9387132218	-0.00061896914622417\\
-21310.2185325357	-0.000512591326913035\\
-21298.4983518496	-0.000414254161034963\\
-21286.7781711635	-0.000324798841846169\\
-21280.9180808204	-0.000283644509092519\\
-21275.0579904774	-0.000244991877513645\\
-21269.1979001343	-0.000208924019328482\\
-21263.3378097912	-0.000175518559451495\\
-21257.4777194482	-0.000144847509567032\\
-21251.6176291051	-0.000116977114014739\\
-21245.7575387621	-9.19677078558858e-05\\
-21239.897448419	-6.98735873809977e-05\\
-21234.0373580759	-5.07428933736981e-05\\
-21228.1772677329	-3.46175073520181e-05\\
-21222.3171773898	-2.15329610357304e-05\\
-21216.4570870468	-1.1518359214173e-05\\
-21210.5969967037	-4.5963161948821e-06\\
-21204.7369063606	-7.82905968610599e-07\\
-21198.8768160176	-8.76261961876233e-08\\
-21193.0167256745	-2.51337610274945e-06\\
-21187.1566353315	-8.05644832862293e-06\\
-21181.2965449884	-1.6706534756168e-05\\
-21175.4364546453	-2.84467463083629e-05\\
-21169.5763643023	-4.32536466786409e-05\\
-21163.7162739592	-6.1097299924636e-05\\
-21157.8561836162	-8.19413318335964e-05\\
-21151.9960932731	-0.000105743004929365\\
-21146.13600293	-0.000132453306964239\\
-21140.275912587	-0.000162017052727458\\
-21134.4158222439	-0.000194372998928431\\
-21128.5557319009	-0.000229453971963535\\
-21122.6956415578	-0.000267187008242536\\
-21116.8355512147	-0.000307493506831298\\
-21110.9754608717	-0.000350289394060698\\
-21099.2552801856	-0.00044298674486833\\
-21087.5350994994	-0.000544503992459431\\
-21075.8149188133	-0.000653990920535655\\
-21064.0947381272	-0.000770529313278586\\
-21052.3745574411	-0.000893140683809489\\
-21040.654376755	-0.00102079451760563\\
-21023.0741057258	-0.001219371205951\\
-21005.4938346966	-0.0014231102353403\\
-20987.9135636674	-0.00162815076132841\\
-20970.3332926383	-0.001830603462985\\
-20952.7530216091	-0.00202662453111203\\
-20935.1727505799	-0.0022124888636091\\
-20911.7323892077	-0.00243843387551629\\
-20888.2920278354	-0.00263242710895824\\
-20864.8516664632	-0.00278789881637128\\
-20841.411305091	-0.00289957289827388\\
-20817.9709437187	-0.00296364698937434\\
-20794.5305823465	-0.00297792299951192\\
-20771.0902209742	-0.00294188363257645\\
-20747.649859602	-0.00285671222316417\\
-20724.2094982298	-0.00272525514233901\\
-20700.7691368575	-0.00255192796067947\\
-20677.3287754853	-0.00234256845891991\\
-20653.8884141131	-0.00210424137417143\\
-20636.3081430839	-0.00191130174259514\\
-20618.7278720547	-0.00171026079202968\\
-20601.1476010255	-0.00150495390425756\\
-20583.5673299963	-0.00129930204105443\\
-20565.9870589672	-0.00109723721700748\\
-20554.266878281	-0.000966435241031535\\
-20542.5466975949	-0.000840064734401158\\
-20530.8265169088	-0.000719204212635943\\
-20519.1063362227	-0.000604886375442926\\
-20507.3861555366	-0.000498089330591291\\
-20495.6659748504	-0.000399728272984745\\
-20489.8058845074	-0.000353978104927964\\
-20483.9457941643	-0.000310647689780881\\
-20478.0857038213	-0.000269830358505824\\
-20472.2256134782	-0.000231614157814808\\
-20466.3655231351	-0.00019608166175015\\
-20460.5054327921	-0.000163309794756426\\
-20454.645342449	-0.000133369666660804\\
-20448.785252106	-0.000106326419876124\\
-20442.9251617629	-8.22390891847683e-05\\
-20437.0650714198	-6.11604743852655e-05\\
-20431.2049810768	-4.31370260855577e-05\\
-20425.3448907337	-2.82087448914954e-05\\
-20419.4848003907	-1.64090941985982e-05\\
-20413.6247100476	-7.76492678911273e-06\\
-20407.7646197046	-2.29642538632653e-06\\
-20401.9045293615	-1.70573014601601e-08\\
-20396.0444390184	-9.33543273517639e-07\\
-20390.1843486754	-5.04584057474679e-06\\
-20384.3242583323	-1.23471404209158e-05\\
-20378.4641679893	-2.28238796987978e-05\\
-20372.6040776462	-3.64557669931793e-05\\
-20366.7439873031	-5.32158228612711e-05\\
-20360.8838969601	-7.30704342765375e-05\\
-20355.023806617	-9.59794231308124e-05\\
-20349.163716274	-0.000121896128662621\\
-20343.3036259309	-0.000150767503627476\\
-20337.4435355878	-0.000182534224030026\\
-20331.5834452448	-0.000217130812177925\\
-20325.7233549017	-0.000254485772812664\\
-20319.8632645587	-0.000294521742025921\\
-20314.0031742156	-0.000337155648650876\\
-20302.2829935295	-0.000429857506229384\\
-20290.5628128434	-0.000531819517057388\\
-20278.8426321572	-0.000642190833221445\\
-20267.1224514711	-0.000760048725137236\\
-20255.402270785	-0.000884406333322104\\
-20243.6820900989	-0.00101422096927117\\
-20226.1018190697	-0.00121678043811988\\
-20208.5215480405	-0.00142532986583139\\
-20190.9412770114	-0.00163592296076743\\
-20173.3610059822	-0.00184456966670262\\
-20155.780734953	-0.00204731205105907\\
-20138.2004639238	-0.00224029967761424\\
-20114.7601025516	-0.00247613493676936\\
-20091.3197411793	-0.00268016006111007\\
-20067.8793798071	-0.00284546356189473\\
-20044.4390184349	-0.00296643163928709\\
-20020.9986570626	-0.00303894019690302\\
-19997.5582956904	-0.00306049685036157\\
-19974.1179343182	-0.00303032804499619\\
-19950.6775729459	-0.00294940824394311\\
-19927.2372115737	-0.0028204300975864\\
-19903.7968502014	-0.00264771649768994\\
-19880.3564888292	-0.00243707738301771\\
-19856.916127457	-0.00219561603434716\\
-19839.3358564278	-0.00199916885448461\\
-19821.7555853986	-0.00179370166096843\\
-19804.1753143694	-0.00158313571932502\\
-19786.5950433403	-0.00137149493905905\\
-19769.0147723111	-0.00116282942486231\\
-19757.294591625	-0.00102735134846444\\
-19745.5744109388	-0.000896132343363573\\
-19733.8542302527	-0.000770294123527442\\
-19722.1340495666	-0.000650913919879023\\
-19710.4138688805	-0.000539015319244037\\
-19698.6936881944	-0.000435559548498835\\
-19686.9735075082	-0.00034143727816274\\
-19681.1134171652	-0.000298133362220343\\
-19675.2533268221	-0.000257461015274727\\
-19669.3932364791	-0.000219508446775727\\
-19663.533146136	-0.000184358150436336\\
-19657.6730557929	-0.000152086726849342\\
-19651.8129654499	-0.000122764718492462\\
-19645.9528751068	-9.64564574642837e-05\\
-19640.0927847638	-7.32199263042771e-05\\
-19634.2326944207	-5.31066321693841e-05\\
-19628.3726040776	-3.61614946600077e-05\\
-19622.5125137346	-2.24227475303034e-05\\
-19616.6524233915	-1.19218545007031e-05\\
-19610.7923330485	-4.6834393550504e-06\\
-19604.9322427054	-7.25230481056351e-07\\
-19599.0721523623	-5.8019972881849e-08\\
-19593.2120620193	-2.68563739496502e-06\\
-19587.3519716762	-8.60493826722153e-06\\
-19581.4918813332	-1.78058073018194e-05\\
-19575.6317909901	-3.02711763954224e-05\\
-19569.7717006471	-4.59770573428382e-05\\
-19563.911610304	-6.48925892133559e-05\\
-19558.0515199609	-8.69801002897976e-05\\
-19552.1914296179	-0.000112195184460174\\
-19546.3313392748	-0.000140486791887828\\
-19540.4712489318	-0.000171797333797657\\
-19534.6111585887	-0.000206062801140336\\
-19528.7510682456	-0.000243212896917138\\
-19522.8909779026	-0.000283171181858921\\
-19517.0308875595	-0.000325855233192107\\
-19511.1707972165	-0.000371176816142205\\
-19499.4506165303	-0.000469351693087314\\
-19487.7304358442	-0.000576880978879248\\
-19476.0102551581	-0.000692869827713654\\
-19464.290074472	-0.000816350876680697\\
-19452.5698937859	-0.000946292361867014\\
-19440.8497130997	-0.00108160678768538\\
-19423.2694420706	-0.00129216150219575\\
-19405.6891710414	-0.00150827166609827\\
-19388.1089000122	-0.00172585220064794\\
-19370.528628983	-0.00194078384828432\\
-19352.9483579539	-0.00214899153281789\\
-19335.3680869247	-0.00234652200028608\\
-19311.9277255524	-0.00258683256160068\\
-19288.4873641802	-0.0027933871035165\\
-19265.047002808	-0.00295918251593026\\
-19241.6066414357	-0.00307857974996451\\
-19218.1662800635	-0.00314749735773166\\
-19194.7259186912	-0.0031635527994453\\
-19171.285557319	-0.00312614663227855\\
-19147.8451959468	-0.00303648661415781\\
-19124.4048345745	-0.00289755077932135\\
-19100.9644732023	-0.0027139906008825\\
-19077.5241118301	-0.00249197738222417\\
-19054.0837504578	-0.00223899694348069\\
-19036.5034794286	-0.00203405071237181\\
-19018.9232083995	-0.00182038724457605\\
-19001.3429373703	-0.00160208759142087\\
-18983.7626663411	-0.00138332807492845\\
-18966.1823953119	-0.00116830094451714\\
-18954.4622146258	-0.00102906680304969\\
-18942.7420339397	-0.000894520485818125\\
-18931.0218532536	-0.000765814873823112\\
-18919.3016725674	-0.000644054721529025\\
-18907.5814918813	-0.000530287249178608\\
-18895.8613111952	-0.00042549321124583\\
-18890.0012208522	-0.000376747467148975\\
-18884.1411305091	-0.000330578517176891\\
-18878.281040166	-0.000287086632907327\\
-18872.420949823	-0.000246366475928381\\
-18866.5608594799	-0.000208506894207686\\
-18860.7007691369	-0.000173590730640519\\
-18854.8406787938	-0.000141694644190633\\
-18848.9805884507	-0.000112888944024474\\
-18843.1204981077	-8.72374369815938e-05\\
-18837.2604077646	-6.47972887199579e-05\\
-18831.4003174216	-4.56188988320399e-05\\
-18825.5402270785	-2.97457902099721e-05\\
-18819.6801367354	-1.72145128995256e-05\\
-18813.8200463924	-8.054562652008e-06\\
-18807.9599560493	-2.28831436094514e-06\\
-18802.0998657063	6.90294731979261e-08\\
-18796.2397753632	-9.90524872404033e-07\\
-18790.3796850201	-5.46774119205874e-06\\
-18784.5195946771	-1.33561474888933e-05\\
-18778.659504334	-2.46420454223196e-05\\
-18772.799413991	-3.93045350541097e-05\\
-18766.9393236479	-5.73155548495061e-05\\
-18761.0792333048	-7.86399368576844e-05\\
-18755.2191429618	-0.000103235476965918\\
-18749.3590526187	-0.000131053020082639\\
-18743.4989622757	-0.000162036560084215\\
-18737.6388719326	-0.00019612335432281\\
-18731.7787815895	-0.000233244052445928\\
-18725.9186912465	-0.000273322839292687\\
-18720.0586009034	-0.000316277591548625\\
-18714.1985105604	-0.000362020047838695\\
-18702.4783298743	-0.000461485448615252\\
-18690.7581491881	-0.000570897465252642\\
-18679.037968502	-0.000689349321815146\\
-18667.3177878159	-0.000815856617485192\\
-18655.5976071298	-0.000949365567354957\\
-18643.8774264437	-0.00108876184037811\\
-18626.2971554145	-0.00130633382351794\\
-18608.7168843853	-0.00153042270630186\\
-18591.1366133561	-0.00175679955082239\\
-18573.5563423269	-0.00198118421492547\\
-18555.9760712978	-0.00219932674509548\\
-18538.3958002686	-0.00240708834446833\\
-18514.9554388963	-0.00266116777054495\\
-18491.5150775241	-0.00288121078107255\\
-18468.0747161519	-0.00305975541951707\\
-18444.6343547796	-0.00319072521218933\\
-18421.1939934074	-0.00326963812941574\\
-18397.7536320352	-0.00329376217527827\\
-18374.3132706629	-0.00326221221261659\\
-18350.8729092907	-0.00317598460315991\\
-18327.4325479184	-0.00303792833347716\\
-18303.9921865462	-0.00285265343988584\\
-18280.551825174	-0.0026263796659168\\
-18257.1114638017	-0.00236673031190815\\
-18239.5311927726	-0.0021553327840013\\
-18221.9509217434	-0.00193411008510074\\
-18204.3706507142	-0.0017072890074432\\
-18186.790379685	-0.00147921149032229\\
-18169.2101086558	-0.00125425218065833\\
-18157.4899279697	-0.00110815159609184\\
-18145.7697472836	-0.000966613188023733\\
-18134.0495665975	-0.00083085174866964\\
-18122.3293859114	-0.000702034797422572\\
-18110.6092052252	-0.000581272632318132\\
-18098.8890245391	-0.000469608852506454\\
-18087.168843853	-0.000368011432396361\\
-18081.3087535099	-0.000321267364520118\\
-18075.4486631669	-0.000277364423876758\\
-18069.5885728238	-0.000236398757207158\\
-18063.7284824808	-0.000198460357582423\\
-18057.8683921377	-0.000163632869696458\\
-18052.0083017947	-0.000131993408469772\\
-18046.1482114516	-0.000103612391332356\\
-18040.2881211085	-7.85533845819969e-05\\
-18034.4280307655	-5.68729641296078e-05\\
-18028.5679404224	-3.86205909431193e-05\\
-18022.7078500794	-2.38385014690801e-05\\
-18016.8477597363	-1.256161326216e-05\\
-18010.9876693932	-4.81744604422046e-06\\
-18005.1275790502	-6.26058363293904e-07\\
nan	nan\\
-17993.4073983641	-2.94428023271915e-06\\
-17987.547308021	-9.45635203853158e-06\\
-17981.6872176779	-1.95261122754273e-05\\
-17975.8271273349	-3.31359178562001e-05\\
-17969.9670369918	-5.0260617890078e-05\\
-17964.1069466488	-7.08676017363276e-05\\
-17958.2468563057	-9.49168628787064e-05\\
-17952.3867659626	-0.000122361078494211\\
-17946.5266756196	-0.00015314570456439\\
-17940.6665852765	-0.000187209086336214\\
-17934.8064949335	-0.000224482583897122\\
-17928.9464045904	-0.000264890712639448\\
-17923.0863142473	-0.00030835129829279\\
-17917.2262239043	-0.000354775646231425\\
-17911.3661335612	-0.000404068724687542\\
-17899.6459528751	-0.000510850454125667\\
-17887.925772189	-0.000627817291522315\\
-17876.2055915029	-0.000754002566785259\\
-17864.4854108167	-0.000888360196592567\\
-17852.7652301306	-0.00102977342971059\\
-17841.0450494445	-0.00117706420270445\\
-17823.4647784153	-0.00140631971993997\\
-17805.8845073862	-0.00164171214863288\\
-17788.304236357	-0.00187880394986281\\
-17770.7239653278	-0.00211311584269286\\
-17753.1436942986	-0.00234021200491103\\
-17735.5634232694	-0.00255578463012177\\
-17712.1230618972	-0.00281824737597132\\
-17688.682700525	-0.00304409270797898\\
-17665.2423391527	-0.00322565414337754\\
-17641.8019777805	-0.00335674158194467\\
-17618.3616164083	-0.00343285499842081\\
-17594.921255036	-0.0034513416214765\\
-17571.4808936638	-0.00341149112122121\\
-17548.0405322915	-0.00331456540429649\\
-17524.6001709193	-0.0031637618160903\\
-17501.1598095471	-0.00296411079482232\\
-17477.7194481748	-0.00272231124081693\\
-17454.2790868026	-0.00244650897752015\\
-17436.6988157734	-0.00222291492083514\\
-17419.1185447442	-0.00198968700024413\\
-17401.5382737151	-0.00175128573374146\\
-17383.9580026859	-0.00151228062114851\\
-17366.3777316567	-0.00127726338193065\\
-17354.6575509706	-0.00112504134473395\\
-17342.9373702845	-0.000977913898139755\\
-17331.2171895983	-0.00083714722528202\\
-17319.4970089122	-0.000703955630570854\\
-17307.7768282261	-0.000579491176827423\\
-17296.05664754	-0.00046483383308878\\
-17290.1965571969	-0.000411498312991594\\
-17284.3364668539	-0.000360982217294307\\
-17278.4763765108	-0.000313396253729843\\
-17272.6162861677	-0.000268845013054363\\
-17266.7561958247	-0.000227426742778135\\
-17260.8961054816	-0.0001892331341979\\
-17255.0360151386	-0.000154349123092558\\
-17249.1759247955	-0.000122852704579042\\
-17243.3158344524	-9.48147625138472e-05\\
-17237.4557441094	-7.02989137991503e-05\\
-17231.5956537663	-4.93613679520115e-05\\
-17225.7355634233	-3.20508022419686e-05\\
-17219.8754730802	-1.84082526626683e-05\\
-17214.0153827372	-8.467020996523e-06\\
-17208.1552923941	-2.25259817053585e-06\\
-17202.295202051	2.17395917606153e-07\\
-17196.435111708	-1.06674403792564e-06\\
-17190.5750213649	-6.10678190334502e-06\\
-17184.7149310219	-1.48965300430669e-05\\
-17178.8548406788	-2.74218560787596e-05\\
-17172.9947503357	-4.36607064672648e-05\\
-17167.1346599927	-6.35831468621092e-05\\
-17161.2745696496	-8.71514191847529e-05\\
-17155.4144793066	-0.0001143200153026\\
-17149.5543889635	-0.000145035767163624\\
-17143.6942986204	-0.000179237953221132\\
-17137.8342082774	-0.000216858420919931\\
-17131.9741179343	-0.000257821725022549\\
-17126.1140275913	-0.000302045281462762\\
-17120.2539372482	-0.000349439536447095\\
-17114.3938469051	-0.00039990815044365\\
-17102.673666219	-0.000509650373644626\\
-17090.9534855329	-0.000630373411576671\\
-17079.2333048468	-0.000761084283144257\\
-17067.5131241607	-0.000900703820765691\\
-17055.7929434745	-0.00104807566845909\\
-17044.0727627884	-0.00120197594788073\\
-17026.4924917592	-0.00144225164530541\\
-17008.9122207301	-0.00168981450016808\\
-16991.3319497009	-0.00194000642106733\\
-16973.7516786717	-0.00218810788894791\\
-16956.1714076425	-0.00242942768344873\\
-16938.5911366134	-0.00265939229299421\\
-16915.1507752411	-0.00294083786981587\\
-16891.7104138689	-0.0031848439498138\\
-16868.2700524966	-0.00338312646340058\\
-16844.8296911244	-0.00352892142605913\\
-16821.3893297522	-0.00361721888623165\\
-16797.9489683799	-0.00364493834982327\\
-16774.5086070077	-0.00361103955292185\\
-16751.0682456355	-0.00351656461909876\\
-16727.6278842632	-0.00336460994172886\\
-16704.187522891	-0.00316022850187203\\
-16680.7471615187	-0.00291026568369465\\
-16657.3068001465	-0.00262313390868933\\
-16639.7265291173	-0.00238918878177739\\
-16622.1462580881	-0.00214423714665986\\
-16604.565987059	-0.00189296510718047\\
-16586.9857160298	-0.00164019175297522\\
-16569.4054450006	-0.00139077773104547\\
-16557.6852643145	-0.00122874806979584\\
-16545.9650836284	-0.00107174565352221\\
-16534.2449029423	-0.000921124013042951\\
-16522.5247222561	-0.00077818507271979\\
-16510.80454157	-0.00064416803343282\\
-16499.0843608839	-0.000520238766584936\\
-16487.3641801978	-0.000407479809935973\\
-16481.5040898547	-0.000355602396256526\\
-16475.6439995117	-0.000306881050907357\\
-16469.7839091686	-0.000261423630404602\\
-16463.9238188255	-0.000219331177446408\\
-16458.0637284825	-0.000180697700851392\\
-16452.2036381394	-0.000145609970177582\\
-16446.3435477964	-0.000114147325468672\\
-16440.4834574533	-8.63815025671135e-05\\
-16434.6233671102	-6.2376474350329e-05\\
-16428.7632767672	-4.21883082618168e-05\\
-16422.9031864241	-2.58650404379933e-05\\
-16417.0430960811	-1.344656672264e-05\\
-16411.183005738	-4.96455080859517e-06\\
-16405.3229153949	-4.42349719875634e-07\\
-16399.4628250519	1.05043195003138e-07\\
-16393.6027347088	-3.32896238068003e-06\\
-16387.7426443658	-1.07425321233418e-05\\
-16381.8825540227	-2.21254032706356e-05\\
-16376.0224636796	-3.74588986587459e-05\\
-16370.1623733366	-5.6715958578948e-05\\
-16364.3022829935	-7.98611904064394e-05\\
-16358.4421926505	-0.000106850935909349\\
-16352.5821023074	-0.000137633356111667\\
-16346.7220119643	-0.00017214853354908\\
-16340.8619216213	-0.000210328591719765\\
-16335.0018312782	-0.000252097831477033\\
-16329.1417409352	-0.000297372884124726\\
-16323.2816505921	-0.000346062880858526\\
-16317.4215602491	-0.000398069638266874\\
-16311.561469906	-0.00045328785944793\\
-16299.8412892199	-0.00057290325122892\\
-16288.1211085338	-0.000703932999889502\\
-16276.4009278476	-0.000845302577438621\\
-16264.6807471615	-0.000995847901314616\\
-16252.9605664754	-0.00115432502617893\\
-16241.2403857893	-0.00131942052976436\\
-16223.6601147601	-0.0015764590329577\\
-16206.0798437309	-0.00184047531004613\\
-16188.4995727017	-0.00210650727740089\\
-16170.9193016726	-0.00236954064236121\\
-16153.3390306434	-0.00262460426284351\\
-16135.7587596142	-0.00286686495555605\\
-16112.318398242	-0.00316205242376054\\
-16088.8780368697	-0.00341634258813389\\
-16065.4376754975	-0.00362109170310465\\
-16041.9973141253	-0.00376930052668983\\
-16018.556952753	-0.00385585748052977\\
-15995.1165913808	-0.00387771890197957\\
-15971.6762300085	-0.00383402005160874\\
-15948.2358686363	-0.00372611285472477\\
-15924.7955072641	-0.00355752881660454\\
-15901.3551458918	-0.00333386807227135\\
-15877.9147845196	-0.00306261803341563\\
-15854.4744231474	-0.00275290748649402\\
-15836.8941521182	-0.002501641152423\\
-15819.313881089	-0.0022394075006136\\
-15801.7336100598	-0.00197122912742261\\
-15784.1533390306	-0.00170225724090561\\
-15766.5730680015	-0.00143767392686483\\
-15754.8528873153	-0.00126625449902387\\
-15743.1327066292	-0.00110054047939344\\
-15731.4125259431	-0.000941965042732009\\
-15719.692345257	-0.000791903830388467\\
-15707.9721645709	-0.000651663183166302\\
-15696.2519838848	-0.000522468936991754\\
-15690.3918935417	-0.000462372664789287\\
-15684.5318031986	-0.000405455877350871\\
-15678.6717128556	-0.000351844588788032\\
-15672.8116225125	-0.000301657943000936\\
-15666.9515321695	-0.000255007954396796\\
-15661.0914418264	-0.000211999263428327\\
-15655.2313514833	-0.00017272890747811\\
-15649.3712611403	-0.000137286107612723\\
-15643.5111707972	-0.000105752071654079\\
-15637.6510804542	-7.81998140127464e-05\\
-15631.7909901111	-5.46939926726159e-05\\
-15625.930899768	-3.52907636955655e-05\\
-15620.070809425	-2.00376535697323e-05\\
-15614.2107190819	-8.97344969084332e-06\\
-15608.3506287389	-2.1281092255186e-06\\
-15602.4905383958	4.773134301401e-07\\
-15596.6304480527	-1.16927957373461e-06\\
-15590.7703577097	-7.07099466777598e-06\\
-15584.9102673666	-1.72219309810085e-05\\
-15579.0501770236	-3.16071835082891e-05\\
-15573.1900866805	-5.0202865334378e-05\\
-15567.3299963374	-7.2976148895099e-05\\
-15561.4699059944	-9.98853261971183e-05\\
-15555.6098156513	-0.000130879887902337\\
-15549.7497253083	-0.000165900621115956\\
-15543.8896349652	-0.000204879725700828\\
-15538.0295446221	-0.000247740948886146\\
-15532.1694542791	-0.000294399737903101\\
-15526.309363936	-0.000344763410350772\\
-15520.449273593	-0.000398731341943496\\
-15514.5891832499	-0.000456195171261834\\
-15502.8690025638	-0.000581139735917738\\
-15491.1488218777	-0.000718584280846323\\
-15479.4286411916	-0.00086740798862672\\
-15467.7084605054	-0.00102639131353052\\
-15455.9882798193	-0.00119422610243507\\
-15444.2680991332	-0.00136952648722847\\
-15426.687828104	-0.00164328368143192\\
-15409.1075570748	-0.00192544384110782\\
-15391.5272860457	-0.00221071580829229\\
-15373.9470150165	-0.00249373244525841\\
-15356.3667439873	-0.00276915263604038\\
-15338.7864729581	-0.00303176311997486\\
-15315.3461115859	-0.00335341180118523\\
-15291.9057502136	-0.00363257930993113\\
-15268.4653888414	-0.00385977550368891\\
-15245.0250274692	-0.00402722956387104\\
-15221.5846660969	-0.00412916048897852\\
-15198.1443047247	-0.00416198145516016\\
-15174.7039433525	-0.0041244308448573\\
-15151.2635819802	-0.00401762519138437\\
-15127.823220608	-0.00384503190591027\\
-15104.3828592357	-0.00361236235425926\\
-15080.9424978635	-0.00332738854628988\\
-15057.5021364913	-0.00299968929659867\\
-15039.9218654621	-0.00273248936690219\\
-15022.3415944329	-0.00245256272851965\\
-15004.7613234037	-0.00216527240276336\\
-14987.1810523746	-0.00187614027613619\\
-14969.6007813454	-0.00159074244143019\\
-14957.8806006593	-0.00140528517846411\\
-14946.1604199731	-0.00122554832635953\\
-14934.440239287	-0.00105308939068388\\
-14922.7200586009	-0.000889407827800184\\
-14910.9998779148	-0.000735932211527269\\
-14899.2796972287	-0.000594007970482249\\
-14887.5595165425	-0.000464885801107435\\
-14881.6994261995	-0.000405487754519023\\
-14875.8393358564	-0.000349710856090986\\
-14869.9792455134	-0.00029768012719494\\
-14864.1191551703	-0.000249512801427311\\
-14858.2590648272	-0.000205318067391454\\
-14852.3989744842	-0.000165196828188351\\
-14846.5388841411	-0.000129241478160138\\
-14840.6787937981	-9.75356973902043e-05\\
-14834.818703455	-7.01542644141694e-05\\
-14828.958613112	-4.7162887539361e-05\\
-14823.0985227689	-2.86180551753425e-05\\
-14817.2384324258	-1.45669054970921e-05\\
-14811.3783420828	-5.04711574925861e-06\\
-14805.5182517397	-8.6811443922328e-08\\
-14799.6581613967	2.95504333086003e-07\\
-14793.7980710536	-3.90899866785218e-06\\
-14787.9379807105	-1.26994478629447e-05\\
-14782.0778903675	-2.60652583372505e-05\\
-14776.2178000244	-4.39861438923031e-05\\
-14770.3577096814	-6.64321486381091e-05\\
-14764.4976193383	-9.33636990873341e-05\\
-14758.6375289952	-0.000124731676664487\\
-14752.7774386522	-0.000160477510506018\\
-14746.9173483091	-0.000200533290366235\\
-14741.0572579661	-0.000244821899432677\\
-14735.197167623	-0.000293257166773564\\
-14729.3370772799	-0.00034574403913085\\
-14723.4769869369	-0.000402178771716434\\
-14717.6168965938	-0.000462449137631033\\
-14711.7568062508	-0.00052643465546198\\
-14700.0366255646	-0.00066502943812183\\
-14688.3164448785	-0.000816843928538207\\
-14676.5962641924	-0.000980644310672026\\
-14664.8760835063	-0.00115509232217937\\
-14653.1559028202	-0.00133875634153709\\
-14641.4357221341	-0.00153012329196357\\
-14623.8554511049	-0.00182814070969855\\
-14606.2751800757	-0.00213435837440913\\
-14588.6949090465	-0.00244304124377386\\
-14571.1146380173	-0.00274838688313236\\
-14553.5343669882	-0.00304463577926302\\
-14535.954095959	-0.00332618123636432\\
-14512.5137345867	-0.00366951270763007\\
-14489.0733732145	-0.00396561708263934\\
-14465.6330118423	-0.00420441347361459\\
-14442.19265047	-0.00437771359764474\\
-14418.7522890978	-0.00447950828613353\\
-14395.3119277256	-0.00450618167305567\\
-14371.8715663533	-0.00445664546272161\\
-14348.4312049811	-0.00433238833977242\\
-14324.9908436088	-0.00413743842996907\\
-14301.5504822366	-0.00387823964719625\\
-14278.1101208644	-0.00356344567879536\\
-14254.6697594921	-0.00320363816260212\\
-14237.0894884629	-0.00291150977669876\\
-14219.5092174338	-0.0026064623691151\\
-14201.9289464046	-0.00229434862998998\\
-14184.3486753754	-0.00198117875579589\\
-14166.7684043462	-0.001673006530483\\
-14155.0482236601	-0.00147329542822453\\
-14143.328042974	-0.00128019931801013\\
-14131.6078622879	-0.00109539790317468\\
-14119.8876816018	-0.000920504995117792\\
-14108.1675009156	-0.000757054672722113\\
-14096.4473202295	-0.000606488081124548\\
-14090.5872298865	-0.000536458049918619\\
-14084.7271395434	-0.000470140983008642\\
-14078.8670492003	-0.000407685483172795\\
-14073.0069588573	-0.00034923216998541\\
-14067.1468685142	-0.00029491337167422\\
-14061.2867781712	-0.000244852834125261\\
-14055.4266878281	-0.000199165447665733\\
-14049.566597485	-0.000157956992237147\\
-14043.706507142	-0.000121323901525574\\
-14037.8464167989	-8.93530465502477e-05\\
-14031.9863264559	-6.21215392153655e-05\\
-14026.1262361128	-3.96965562601339e-05\\
-14020.2661457697	-2.21351840003748e-05\\
-14014.4060554267	-9.48428422582442e-06\\
-14008.5459650836	-1.78038155813986e-06\\
-14002.6858747406	9.50427459075276e-07\\
-13996.8257843975	-1.30745667827178e-06\\
-13990.9656940544	-8.5590895997881e-06\\
-13985.1056037114	-2.07989584786331e-05\\
-13979.2455133683	-3.80109797838483e-05\\
-13973.3854230253	-6.01685194012163e-05\\
-13967.5253326822	-8.72344351123164e-05\\
-13961.6652423392	-0.000119161141364189\\
-13955.8051519961	-0.000155890696233735\\
-13949.945061653	-0.000197354910416347\\
-13944.08497131	-0.000243475478057306\\
-13938.2248809669	-0.000294164129152046\\
-13932.3647906239	-0.000349322803258595\\
-13926.5047002808	-0.000408843844137917\\
-13920.6446099377	-0.00047261021497196\\
-13914.7845195947	-0.000540495733723048\\
-13903.0643389086	-0.000688075310060081\\
-13891.3441582224	-0.000850400377821322\\
-13879.6239775363	-0.00102616068145417\\
-13867.9037968502	-0.00121392866245647\\
-13856.1836161641	-0.00141217123863335\\
-13844.463435478	-0.00161926250841621\\
-13826.8831644488	-0.00194274240345282\\
-13809.3028934196	-0.00227626601918929\\
-13791.7226223904	-0.00261360398102513\\
-13774.1423513612	-0.00294842926925367\\
-13756.5620803321	-0.00327443739756012\\
-13738.9818093029	-0.00358546665249678\\
-13715.5414479307	-0.00396672453832547\\
-13692.1010865584	-0.00429800319315724\\
-13668.6607251862	-0.00456802352855954\\
-13645.2203638139	-0.00476752102761162\\
-13621.7800024417	-0.00488957061478979\\
-13598.3396410695	-0.00492983415279872\\
-13574.8992796972	-0.00488672176510969\\
-13551.458918325	-0.00476146104900936\\
-13528.0185569528	-0.00455807132229245\\
-13504.5781955805	-0.00428324324341816\\
-13481.1378342083	-0.00394612734397704\\
-13457.697472836	-0.0035580381126893\\
-13440.1172018069	-0.00324134948453528\\
-13422.5369307777	-0.0029093868254581\\
-13404.9566597485	-0.00256852132462072\\
-13387.3763887193	-0.00222532187206752\\
-13369.7961176901	-0.00188643071881885\\
-13358.075937004	-0.00166615676518477\\
-13346.3557563179	-0.00145264326805179\\
-13334.6355756318	-0.0012477519367211\\
-13322.9153949457	-0.00105327689021096\\
-13311.1952142596	-0.000870929261528749\\
-13299.4750335734	-0.000702322459742393\\
-13287.7548528873	-0.000548958216667969\\
-13281.8947625443	-0.000478427577543718\\
-13276.0346722012	-0.000412213538406629\\
-13270.1745818581	-0.000350466762573949\\
-13264.3144915151	-0.000293328674839658\\
-13258.454401172	-0.000240931148450913\\
-13252.594310829	-0.000193396211866665\\
-13246.7342204859	-0.000150835775911191\\
-13240.8741301428	-0.000113351381929817\\
-13235.0140397998	-8.10339715372089e-05\\
-13229.1539494567	-5.39636784412149e-05\\
-13223.2938591137	-3.22096428508621e-05\\
-13217.4337687706	-1.58298488625258e-05\\
-13211.5736784275	-4.87098522725504e-06\\
-13205.7135880845	6.3167019038465e-07\\
-13199.8534977414	6.54341970378372e-07\\
-13193.9934073984	-4.8151755103121e-06\\
-13188.1333170553	-1.57774751126702e-05\\
-13182.2732267122	-3.22215184816521e-05\\
-13176.4131363692	-5.41246421823077e-05\\
-13170.5530460261	-8.14525884456714e-05\\
-13164.6929556831	-0.00011415956049917\\
-13158.83286534	-0.000152188302400557\\
-13152.9727749969	-0.000195470203237305\\
-13147.1126846539	-0.000243925425517103\\
-13141.2525943108	-0.000297463057510905\\
-13135.3925039678	-0.00035598128924929\\
-13129.5324136247	-0.000419367611854463\\
-13123.6723232817	-0.000487499039817409\\
-13117.8122329386	-0.000560242355756582\\
-13106.0920522525	-0.000718982245030837\\
-13094.3718715664	-0.000894327574196526\\
-13082.6516908802	-0.00108487418502644\\
-13070.9315101941	-0.00128908492892611\\
-13059.211329508	-0.00150530228981414\\
-13047.4911488219	-0.00173176206430656\\
-13035.7709681358	-0.0019666079865307\\
-13018.1906971066	-0.00233035573420049\\
-13000.6104260774	-0.00270184819879727\\
-12983.0301550482	-0.0030741463658034\\
-12965.449884019	-0.00344026461380296\\
-12947.8696129899	-0.00379330365670004\\
-12924.4292516176	-0.00423217700209131\\
-12900.9888902454	-0.00462113957348873\\
-12877.5485288732	-0.00494697168780536\\
-12854.1081675009	-0.00519851312115898\\
-12830.6678061287	-0.0053670539438863\\
-12807.2274447564	-0.00544664484648214\\
-12783.7870833842	-0.0054343159556759\\
-12760.346722012	-0.0053301961663468\\
-12736.9063606397	-0.00513752832805438\\
-12713.4659992675	-0.00486257911584817\\
-12690.0256378953	-0.00451444596939408\\
-12666.585276523	-0.00410476697843897\\
-12649.0050054938	-0.00376535139099728\\
-12631.4247344647	-0.00340552026183173\\
-12613.8444634355	-0.00303217118216235\\
-12596.2641924063	-0.00265249363169663\\
-12578.6839213771	-0.00227383334547552\\
-12561.1036503479	-0.00190355304595082\\
-12549.3834696618	-0.00166493937274596\\
-12537.6632889757	-0.00143535901862137\\
-12525.9431082896	-0.0012168282831699\\
-12514.2229276035	-0.00101127650213027\\
-12502.5027469173	-0.000820529415734817\\
-12490.7825662312	-0.000646293359927287\\
-12484.9224758882	-0.000565863872680999\\
-12479.0623855451	-0.00049014041669363\\
-12473.2022952021	-0.000419296266918004\\
-12467.342204859	-0.000353494652429067\\
-12461.4821145159	-0.000292888394196339\\
-12455.6220241729	-0.000237619564309568\\
-12449.7619338298	-0.000187819167386013\\
-12443.9018434868	-0.000143606844864621\\
-12438.0417531437	-0.000105090602865757\\
-12432.1816628006	-7.23665642196277e-05\\
-12426.3215724576	-4.55187452331639e-05\\
-12420.4614821145	-2.4618857708653e-05\\
-12414.6013917715	-9.72613667174347e-06\\
-12408.7413014284	-8.87194216363841e-07\\
-12402.8812110853	1.86410017836111e-06\\
-12397.0211207423	-1.49328743399849e-06\\
-12391.1610303992	-1.09674895595291e-05\\
-12385.3009400562	-2.65536985479246e-05\\
-12379.4408497131	-4.82341551947391e-05\\
-12373.58075937	-7.59781647374057e-05\\
-12367.720669027	-0.000109742140246979\\
-12361.8605786839	-0.00014946967337851\\
-12356.0004883409	-0.000195091632371791\\
-12350.1403979978	-0.000246526287144587\\
-12344.2803076547	-0.000303679461249289\\
-12338.4202173117	-0.000366444710434796\\
-12332.5601269686	-0.000434703527455033\\
-12326.7000366256	-0.000508325572768237\\
-12320.8399462825	-0.000587168930638313\\
-12314.9798559394	-0.000671080390199808\\
-12303.2596752533	-0.00085344015153012\\
-12291.5394945672	-0.00105396546207741\\
-12279.8193138811	-0.00127105816443956\\
-12268.099133195	-0.0015029743429901\\
-12256.3789525089	-0.00174783860663117\\
-12244.6587718227	-0.00200365953334869\\
-12227.0785007936	-0.00240333770227891\\
-12209.4982297644	-0.00281555983069255\\
-12191.9179587352	-0.00323266480331906\\
-12174.337687706	-0.00364685960373361\\
-12156.7574166768	-0.00405036721391871\\
-12139.1771456477	-0.00443557497323858\\
-12115.7367842754	-0.00490815999641268\\
-12092.2964229032	-0.00531928227074619\\
-12068.8560615309	-0.00565491675940761\\
-12045.4157001587	-0.00590350227901801\\
-12021.9753387865	-0.00605635053260512\\
-11998.5349774142	-0.00610796074592958\\
-11975.094616042	-0.00605622859673263\\
-11951.6542546698	-0.00590254162310263\\
-11928.2138932975	-0.00565175707812776\\
-11904.7735319253	-0.00531206214531751\\
-11881.333170553	-0.00489472040428386\\
-11857.8928091808	-0.00441371231139647\\
-11840.3125381516	-0.00402087609760015\\
-11822.7322671225	-0.00360884650759666\\
-11805.1519960933	-0.00318554936262997\\
-11787.5717250641	-0.00275916968489288\\
-11769.9914540349	-0.00233799705551539\\
-11758.2712733488	-0.00206418253330767\\
-11746.5510926627	-0.00179874230552348\\
-11734.8309119766	-0.0015440093137091\\
-11723.1107312904	-0.00130223446394367\\
-11711.3905506043	-0.0010755672505179\\
-11699.6703699182	-0.000866037166831683\\
-11693.8102795751	-0.000768295589181372\\
-11687.9501892321	-0.000675536063791807\\
-11682.090098889	-0.000587970851146087\\
-11676.230008546	-0.000505801608369831\\
-11670.3699182029	-0.00042921894251929\\
-11664.5098278598	-0.000358401986397821\\
-11658.6497375168	-0.000293517997887542\\
-11652.7896471737	-0.000234721983653205\\
-11646.9295568307	-0.000182156348062746\\
-11641.0694664876	-0.00013595056812842\\
-11635.2093761446	-9.62208952007329e-05\\
-11629.3492858015	-6.3070084082781e-05\\
-11623.4891954584	-3.65871502018753e-05\\
-11617.6291051154	-1.68471553991068e-05\\
-11611.7690147723	-3.91102284497095e-06\\
-11605.9089244293	2.17461847321223e-06\\
-11600.0488340862	1.37755931010916e-06\\
-11594.1887437431	-6.31989455759075e-06\\
-11588.3286534001	-2.09208575862028e-05\\
-11582.468563057	-4.24138304485396e-05\\
-11576.608472714	-7.07726968791103e-05\\
-11570.7483823709	-0.000105956751441114\\
-11564.8882920278	-0.000147910758218454\\
-11559.0282016848	-0.000196565040367412\\
-11553.1681113417	-0.000251835600389362\\
-11547.3080209987	-0.000313624270939049\\
-11541.4479306556	-0.000381818895900072\\
-11535.5878403125	-0.000456293541417036\\
-11529.7277499695	-0.000536908736441398\\
-11523.8676596264	-0.000623511742444622\\
-11518.0075692834	-0.000715936851649411\\
-11506.2873885972	-0.000917527687418584\\
-11494.5672079111	-0.00114010927397309\\
-11482.847027225	-0.00138192579859688\\
-11471.1268465389	-0.00164105178308009\\
-11459.4066658528	-0.00191540775218794\\
-11447.6864851666	-0.00220277726615378\\
-11435.9663044805	-0.00250082517838853\\
-11418.3860334513	-0.00296257132498431\\
-11400.8057624222	-0.00343431947684863\\
-11383.225491393	-0.00390730200535276\\
-11365.6452203638	-0.00437267762882282\\
-11348.0649493346	-0.0048216995929385\\
-11324.6245879624	-0.00538034952354886\\
-11301.1842265902	-0.00587602711381145\\
-11277.7438652179	-0.00629185935513033\\
-11254.3035038457	-0.0066135519109268\\
-11230.8631424734	-0.00682989408722613\\
-11207.4227811012	-0.00693316337742769\\
-11183.982419729	-0.00691941509582266\\
-11160.5420583567	-0.00678864638302746\\
-11137.1016969845	-0.00654482802896784\\
-11113.6613356123	-0.00619580197600453\\
-11090.22097424	-0.00575304688685145\\
-11066.7806128678	-0.00523131864439613\\
-11049.2003418386	-0.00479866978748966\\
-11031.6200708094	-0.00433969099928981\\
-11014.0397997802	-0.00386319812857683\\
-10996.4595287511	-0.00337839706934762\\
-10978.8792577219	-0.00289471060746283\\
-10961.2989866927	-0.00242160006130094\\
-10949.5788060066	-0.00211668208161304\\
-10937.8586253205	-0.00182330269912859\\
-10926.1384446344	-0.00154406501829236\\
-10914.4182639482	-0.00128146330477403\\
-10902.6980832621	-0.00103786139158772\\
-10890.977902576	-0.000815472099262405\\
-10885.1178122329	-0.000712877601540279\\
-10879.2577218899	-0.000616337848460132\\
-10873.3976315468	-0.000526078559795566\\
-10867.5375412038	-0.000442312634040603\\
-10861.6774508607	-0.000365239670544749\\
-10855.8173605176	-0.000295045518808255\\
-10849.9572701746	-0.000231901855888585\\
-10844.0971798315	-0.000175965792927513\\
-10838.2370894885	-0.000127379511635701\\
-10832.3769991454	-8.62699316219065e-05\\
-10826.5169088023	-5.27484092917896e-05\\
-10820.6568184593	-2.69104690456881e-05\\
-10814.7967281162	-8.83556739862741e-06\\
-10808.9366377732	1.41310940428786e-06\\
-10803.0765474301	3.78881377659286e-06\\
-10797.216457087	-1.73862924079821e-06\\
-10791.356366744	-1.51827227904997e-05\\
-10785.4962764009	-3.65402340709376e-05\\
-10779.6361860579	-6.57911654813149e-05\\
-10773.7760957148	-0.000102898761174201\\
-10767.9160053717	-0.000147809549082245\\
-10762.0559150287	-0.000200453418395766\\
-10756.1958246856	-0.000260743732411699\\
-10750.3357343426	-0.000328577476596296\\
-10744.4756439995	-0.000403835441610912\\
-10738.6155536565	-0.000486382440981392\\
-10732.7554633134	-0.000576067563078175\\
-10726.8953729703	-0.000672724456844635\\
-10721.0352826273	-0.000776171650899172\\
-10715.1751922842	-0.000886212905200961\\
-10703.4550115981	-0.00112522112495846\\
-10691.734830912	-0.00138789918018199\\
-10680.0146502259	-0.00167218857319558\\
-10668.2944695397	-0.00197583878746574\\
-10656.5742888536	-0.00229642552848497\\
-10644.8541081675	-0.00263137051536534\\
-10627.2738371383	-0.00315474838405956\\
-10609.6935661091	-0.0036947155029711\\
-10592.11329508	-0.00424130234704384\\
-10574.5330240508	-0.00478434851646065\\
-10556.9527530216	-0.00531369529496827\\
-10539.3724819924	-0.00581937945383468\\
-10515.9321206202	-0.00644034369319348\\
-10492.491759248	-0.00698124630269803\\
-10469.0513978757	-0.00742358494436153\\
-10445.6110365035	-0.00775204125828384\\
-10422.1706751312	-0.00795502908235959\\
-10398.730313759	-0.0080251212004464\\
-10375.2899523868	-0.00795933907754341\\
-10351.8495910145	-0.00775929453583125\\
-10328.4092296423	-0.00743117722572026\\
-10304.9688682701	-0.00698558690704567\\
-10281.5285068978	-0.00643721478891272\\
-10258.0881455256	-0.00580438332047388\\
-10240.5078744964	-0.00528709101054839\\
-10222.9276034672	-0.00474417623187421\\
-10205.347332438	-0.00418611565386137\\
-10187.7670614089	-0.00362375093245474\\
-10170.1867903797	-0.00306808447021394\\
-10158.4666096936	-0.00270677930417743\\
-10146.7464290074	-0.00235651582489034\\
-10135.0262483213	-0.0020204066127214\\
-10123.3060676352	-0.00170145931781868\\
-10111.5858869491	-0.00140255066733593\\
-10099.865706263	-0.00112640145748773\\
-10094.0056159199	-0.000997663367469788\\
-10088.1455255769	-0.000875552746390671\\
-10082.2854352338	-0.000760355466990062\\
-10076.4253448907	-0.000652343454739857\\
-10070.5652545477	-0.000551774078525254\\
-10064.7051642046	-0.000458889570673991\\
-10058.8450738616	-0.000373916477591258\\
-10052.9849835185	-0.000297065142248737\\
-10047.1248931754	-0.000228529219733049\\
-10041.2648028324	-0.00016848522688702\\
-10035.4047124893	-0.00011709212714911\\
-10029.5446221463	-7.44909514947628e-05\\
-10023.6845318032	-4.08044564115603e-05\\
-10017.8244414601	-1.61368196951621e-05\\
-10011.9643511171	-5.73374805836168e-07\\
-10006.104260774	5.81961555461356e-06\\
-10000.244170431	2.99514607758894e-06\\
-9994.3840800879	-9.07438479967228e-06\\
-9988.52398974484	-3.03970692294141e-05\\
-9982.66389940178	-6.09614252983426e-05\\
-9976.80380905872	-0.000100736373622287\\
-9970.94371871567	-0.000149671255333666\\
-9965.08362837261	-0.000207695891532479\\
-9959.22353802955	-0.000274720684174514\\
-9953.36344768649	-0.000350636758274933\\
-9947.50335734343	-0.0004353161452437\\
-9941.64326700037	-0.000528612007057587\\
-9935.78317665731	-0.000630358900904697\\
-9929.92308631425	-0.000740373083796323\\
-9924.06299597119	-0.000858452856706486\\
-9918.20290562813	-0.000984378947466937\\
-9906.48272494201	-0.00125880769181348\\
-9894.76254425589	-0.00156157060253946\\
-9883.04236356977	-0.0018903285430994\\
-9871.32218288365	-0.00224251054200134\\
-9859.60200219753	-0.00261533443557092\\
-9847.88182151142	-0.00300582940147332\\
-9836.1616408253	-0.00341086020214953\\
-9818.58136979612	-0.00403846816739487\\
-9801.00109876694	-0.00467989879912989\\
-9783.42082773776	-0.00532331275057724\\
-9765.84055670858	-0.00595674513375714\\
-9748.2602856794	-0.00656833288766489\\
-9724.81992430716	-0.00732994253437018\\
-9701.37956293493	-0.00800656018087187\\
-9677.93920156269	-0.00857510017185733\\
-9654.49884019045	-0.00901591725888633\\
-9631.05847881822	-0.0093135087368781\\
-9607.61811744598	-0.00945708349465288\\
-9584.17775607374	-0.00944097734400312\\
-9560.7373947015	-0.00926489897379619\\
-9537.29703332927	-0.00893399643703047\\
-9513.85667195703	-0.00845874003762402\\
-9490.41631058479	-0.00785462364526479\\
-9466.97594921255	-0.00714169262402612\\
-9449.39567818337	-0.00654987514561934\\
-9431.81540715419	-0.00592157794692808\\
-9414.23513612501	-0.00526890509852343\\
-9396.65486509584	-0.00460452743791083\\
-9379.07459406666	-0.00394144539697287\\
-9361.49432303748	-0.00329274361819302\\
-9349.77414235136	-0.00287465926948797\\
-9338.05396166524	-0.00247245309246311\\
-9326.33378097912	-0.00208974762412753\\
-9314.613600293	-0.00173002032881896\\
-9302.89341960689	-0.0013965733222765\\
-9297.03332926383	-0.00124068173402229\\
-9291.17323892077	-0.00109250441344431\\
-9285.31314857771	-0.000952392169648613\\
-9279.45305823465	-0.000820679716238142\\
-9273.59296789159	-0.000697684918003724\\
-9267.73287754853	-0.000583708071202371\\
-9261.87278720547	-0.000479031218977657\\
-9256.01269686241	-0.000383917503452488\\
-9250.15260651935	-0.000298610556014676\\
-9244.29251617629	-0.000223333927157371\\
-9238.43242583323	-0.000158290557183506\\
-9232.57233549017	-0.000103662289029873\\
-9226.71224514711	-5.96094243329286e-05\\
-9220.85215480405	-2.62703238093769e-05\\
-9214.992064461	-3.76105290718477e-06\\
-9209.13197411794	7.82492638399595e-06\\
-9203.27188377488	8.41701679615916e-06\\
-9197.41179343182	-2.03230056979739e-06\\
-9191.55170308876	-2.35473007952853e-05\\
-9185.6916127457	-5.61288983802229e-05\\
-9179.83152240264	-9.97545795444341e-05\\
-9173.97143205958	-0.000154378383879877\\
-9168.11134171652	-0.000219930935522611\\
-9162.25125137346	-0.00029631952393277\\
-9156.3911610304	-0.000383428234241947\\
-9150.53107068734	-0.000481118127020635\\
-9144.67098034428	-0.000589227467265248\\
-9138.81089000122	-0.00070757200220806\\
-9132.95079965816	-0.000835945287546409\\
-9127.0907093151	-0.000974119061537158\\
-9121.23061897204	-0.00112184366623683\\
-9115.37052862898	-0.00127884851523458\\
-9103.65034794286	-0.00161951508258513\\
-9091.93016725675	-0.00199355611571542\\
-9080.20998657063	-0.00239811150803011\\
-9068.48980588451	-0.00283004623514711\\
-9056.76962519839	-0.00328597535268875\\
-9045.04944451227	-0.00376229125416531\\
-9027.46917348309	-0.0045066304573585\\
-9009.88890245391	-0.00527477136284887\\
-8992.30863142473	-0.00605266598490444\\
-8974.72836039555	-0.00682595985488937\\
-8957.14808936638	-0.00758026334523564\\
-8939.5678183372	-0.00830142605598395\\
-8916.12745696496	-0.00918795747381797\\
-8892.68709559272	-0.00996134920026928\\
-8869.24673422049	-0.0105950345094054\\
-8845.80637284825	-0.01106688918645\\
-8822.36601147601	-0.011360036198386\\
-8798.92565010377	-0.0114634822592087\\
-8775.48528873154	-0.0113725627245209\\
-8752.0449273593	-0.0110891774419085\\
-8728.60456598706	-0.0106218070311142\\
-8705.16420461482	-0.00998530635664211\\
-8681.72384324258	-0.00920047943980881\\
-8658.28348187034	-0.00829344749103283\\
-8640.70321084117	-0.00755125312717801\\
-8623.12293981199	-0.00677173350075279\\
-8605.54266878281	-0.00597000655580826\\
-8587.96239775363	-0.00516176310263483\\
-8570.38212672445	-0.00436297278406238\\
-8558.66194603833	-0.00384358253648326\\
-8546.94176535222	-0.00334014404606398\\
-8535.2215846661	-0.00285720894460144\\
-8523.50140397998	-0.00239918488619662\\
-8511.78122329386	-0.00197029725704231\\
-8500.06104260774	-0.0015745521810667\\
-8494.20095226468	-0.00139029131535057\\
-8488.34086192162	-0.00121570114241605\\
-8482.48077157856	-0.00105120532135851\\
-8476.6206812355	-0.000897207534121875\\
-8470.76059089244	-0.000754090566093032\\
-8464.90050054939	-0.000622215427785242\\
-8459.04041020633	-0.0005019205197564\\
-8453.18031986327	-0.000393520842531137\\
-8447.32022952021	-0.000297307253461387\\
-8441.46013917715	-0.000213545772104161\\
-8435.60004883409	-0.000142476935878959\\
-8429.73995849103	-8.4315207436601e-05\\
-8423.87986814797	-3.92484351895344e-05\\
-8418.01977780491	-7.43736831905056e-06\\
-8412.15968746185	1.09847725478301e-05\\
-8406.29959711879	1.59126679020703e-05\\
-8400.43950677573	7.26921418331051e-06\\
-8394.57941643267	-1.49942199766626e-05\\
-8388.71932608961	-5.08975968296321e-05\\
-8382.85923574655	-0.000100432072515783\\
-8376.99914540349	-0.000163559920488515\\
-8371.13905506043	-0.000240214515715671\\
-8365.27896471737	-0.000330300379891579\\
-8359.41887437431	-0.000433693287773499\\
-8353.55878403125	-0.00055024043461665\\
-8347.69869368819	-0.000679760664552227\\
-8341.83860334514	-0.000822044759626161\\
-8335.97851300208	-0.000976855789136341\\
-8330.11842265902	-0.00114392951863701\\
-8324.25833231596	-0.00132297487807245\\
-8318.3982419729	-0.00151367448811043\\
-8306.67806128678	-0.00192863895519843\\
-8294.95788060066	-0.00238578128066021\\
-8283.23769991454	-0.00288168198627955\\
-8271.51751922842	-0.0034125701953642\\
-8259.7973385423	-0.00397435322599193\\
-8248.07715785618	-0.00456264910827734\\
-8230.496886827	-0.00548460284504284\\
-8212.91661579783	-0.00643921019716877\\
-8195.33634476865	-0.00740916301652995\\
-8177.75607373947	-0.00837668572914441\\
-8160.17580271029	-0.0093238699436654\\
-8142.59553168111	-0.0102330150762208\\
-8119.15517030888	-0.0113565110135151\\
-8095.71480893664	-0.0123439177454911\\
-8072.2744475644	-0.0131612757384649\\
-8048.83408619216	-0.0137799505293344\\
-8025.39372481993	-0.0141776817014191\\
-8001.95336344769	-0.0143394308680123\\
-7978.51300207545	-0.0142579971199111\\
-7955.07264070321	-0.0139343757865282\\
-7931.63227933097	-0.0133778446711558\\
-7908.19191795873	-0.0126057708848127\\
-7884.7515565865	-0.0116431407279807\\
-7861.31119521426	-0.0105218244416379\\
-7843.73092418508	-0.00959923387898797\\
-7826.1506531559	-0.00862627023100858\\
-7808.57038212672	-0.00762184287113324\\
-7790.99011109754	-0.00660568626821062\\
-7773.40984006837	-0.00559799263855317\\
-7761.68965938225	-0.00494091672700519\\
-7749.96947869613	-0.00430256280080076\\
-7738.24929801001	-0.0036887645673474\\
-7726.52911732389	-0.00310519358549772\\
-7714.80893663777	-0.00255730994362976\\
-7703.08875595166	-0.00205031434601259\\
-7697.2286656086	-0.00181369247194271\\
-7691.36857526554	-0.00158910202191079\\
-7685.50848492248	-0.00137709997138264\\
-7679.64839457942	-0.00117821886042376\\
-7673.78830423636	-0.000992965571739724\\
-7667.9282138933	-0.000821820157997939\\
-7662.06812355024	-0.000665234721270361\\
-7656.20803320718	-0.000523632347051589\\
-7650.34794286412	-0.000397406095374902\\
-7644.48785252106	-0.000286918051360987\\
-7638.627762178	-0.000192498437387347\\
-7632.76767183494	-0.000114444789048414\\
-7626.90758149188	-5.30211968507886e-05\\
-7621.04749114882	-8.45761550854606e-06\\
-7615.18740080576	1.90507574305546e-05\\
-7609.3273104627	2.93440320654991e-05\\
-7603.46722011964	2.22981042770951e-05\\
-7597.60712977658	-2.17493976741762e-06\\
-7591.74703943353	-4.41264938878022e-05\\
-7585.88694909047	-0.000103571177926833\\
-7580.02685874741	-0.0001804866603871\\
-7574.16676840435	-0.000274813558354109\\
-7568.30667806129	-0.000386455415196423\\
-7562.44658771823	-0.000515278756450276\\
-7556.58649737517	-0.000661113224018044\\
-7550.72640703211	-0.000823751788772461\\
-7544.86631668905	-0.00100295104139083\\
-7539.00622634599	-0.00119843156113038\\
-7533.14613600293	-0.001409878362134\\
-7527.28604565987	-0.00163694141651033\\
-7521.42595531681	-0.00187923625356278\\
-7509.70577463069	-0.00240781529858957\\
-7497.98559394457	-0.00299187833827533\\
-7486.26541325845	-0.00362718786104831\\
-7474.54523257233	-0.00430903784531584\\
-7462.82505188622	-0.00503228982371356\\
-7451.1048712001	-0.00579141289927255\\
-7439.38469051398	-0.00658052740899902\\
-7421.8044194848	-0.00780684088478494\\
-7404.22414845562	-0.00906479804055633\\
-7386.64387742644	-0.0103316585680163\\
-7369.06360639727	-0.0115842350325273\\
-7351.48333536809	-0.0127993302170293\\
-7328.04297399585	-0.014321925398716\\
-7304.60261262361	-0.0156861747399205\\
-7281.16225125137	-0.0168452844263202\\
-7257.72188987914	-0.0177586807862829\\
-7234.2815285069	-0.0183934984765587\\
-7210.84116713466	-0.0187258335602542\\
-7187.40080576242	-0.0187417150808529\\
-7163.96044439018	-0.0184377574007755\\
-7140.52008301795	-0.017821465654814\\
-7117.07972164571	-0.016911177844771\\
-7093.63936027347	-0.0157356390353644\\
-7070.19899890123	-0.0143332153963796\\
-7052.61872787205	-0.0131604138175955\\
-7035.03845684288	-0.0119086551565541\\
-7017.4581858137	-0.0106022338939304\\
-6999.87791478452	-0.00926684942112058\\
-6982.29764375534	-0.00792913526700807\\
-6970.57746306922	-0.00704939798005758\\
-6958.8572823831	-0.00618868403251014\\
-6947.13710169699	-0.00535493635989402\\
-6935.41692101087	-0.00455594769407947\\
-6923.69674032475	-0.00379929311741169\\
-6911.97655963863	-0.00309226380890344\\
-6900.25637895251	-0.00244180255162248\\
-6894.39628860945	-0.00213984394061383\\
-6888.53619826639	-0.00185444155975373\\
-6882.67610792333	-0.00158633805257091\\
-6876.81601758027	-0.00133624317001102\\
-6870.95592723721	-0.00110483211031015\\
-6865.09583689415	-0.000892743923840465\\
-6859.23574655109	-0.000700579986113411\\
-6853.37565620803	-0.000528902542717993\\
-6847.51556586497	-0.000378233329560688\\
-6841.65547552191	-0.000249052271612713\\
-6835.79538517886	-0.000141796263310171\\
-6829.9352948358	-5.68580335656158e-05\\
-6824.07520449274	5.41490181981804e-06\\
-6818.21511414968	4.47211977396864e-05\\
-6812.35502380662	6.08065448811398e-05\\
-6806.49493346356	5.34644400932771e-05\\
-6800.6348431205	2.25368599307568e-05\\
-6794.77475277744	-3.20851645051866e-05\\
-6788.91466243438	-0.000110461073133609\\
-6783.05457209132	-0.000212600404349719\\
-6777.19448174826	-0.000338462525666669\\
-6771.3343914052	-0.00048795646860541\\
-6765.47430106214	-0.000660940868717434\\
-6759.61421071908	-0.000857224011372839\\
-6753.75412037602	-0.00107656398377321\\
-6747.89403003296	-0.00131866893334073\\
-6742.0339396899	-0.00158319743259294\\
-6736.17384934684	-0.00186975895010211\\
-6730.31375900378	-0.00217791442732509\\
-6724.45366866072	-0.00250717696029412\\
-6718.59357831766	-0.0028570125857194\\
-6706.87339763155	-0.00361603739947661\\
-6695.15321694543	-0.00444981228665629\\
-6683.43303625931	-0.00535247504330896\\
-6671.71285557319	-0.00631752045353124\\
-6659.99267488707	-0.007337848763703\\
-6648.27249420095	-0.00840581967497942\\
-6630.69222317177	-0.0100792236035522\\
-6613.1119521426	-0.0118123505099848\\
-6595.53168111342	-0.013574512043281\\
-6577.95141008424	-0.0153339935881119\\
-6560.37113905506	-0.0170586438051054\\
-6542.79086802588	-0.0187164819759805\\
-6519.35050665365	-0.0207692729987549\\
-6495.91014528141	-0.0225781695175015\\
-6472.46978390917	-0.0240802241843144\\
-6449.02942253693	-0.0252216564092075\\
-6425.58906116469	-0.0259598893565532\\
-6402.14869979245	-0.0262652540818828\\
-6378.70833842022	-0.0261222960899504\\
-6355.26797704798	-0.0255306314425216\\
-6331.82761567574	-0.024505313412591\\
-6308.3872543035	-0.0230766861454268\\
-6284.94689293126	-0.021289718372482\\
-6267.36662190209	-0.0197490117059137\\
-6249.78635087291	-0.0180684446572701\\
-6232.20607984373	-0.0162803607417812\\
-6214.62580881455	-0.0144198008258561\\
-6197.04553778538	-0.0125238849122458\\
-6179.4652667562	-0.0106311431676606\\
-6167.74508607008	-0.00939044587901222\\
-6156.02490538396	-0.0081802436172531\\
-6144.30472469784	-0.00701207879060207\\
-6132.58454401172	-0.00589729026377876\\
-6120.8643633256	-0.00484691496030564\\
-6109.14418263948	-0.00387159075676286\\
-6103.28409229642	-0.00341526460706999\\
-6097.42400195336	-0.00298146149550286\\
-6091.5639116103	-0.0025713571050529\\
-6085.70382126725	-0.00218608493357806\\
-6079.84373092419	-0.00182673363421009\\
-6073.98364058113	-0.00149434443234112\\
-6068.12355023807	-0.00118990862541044\\
-6062.26345989501	-0.000914365170761012\\
-6056.40336955195	-0.000668598367619746\\
-6050.54327920889	-0.000453435638175552\\
-6044.68318886583	-0.000269645413337205\\
-6038.82309852277	-0.000117935128014562\\
-6032.96300817971	1.05066909580283e-06\\
-6027.10291783665	8.67320865952666e-05\\
-6021.24282749359	0.000138595537882602\\
-6015.38273715053	0.000156195260852918\\
-6009.52264680747	0.00013915470579904\\
-6003.66255646441	8.71677877848544e-05\\
nan	nan\\
-5991.94237577829	-0.000122510615126645\\
-5986.08228543523	-0.000280452640003633\\
-5980.22219509217	-0.000473840603818683\\
-5974.36210474911	-0.000702614454714258\\
-5968.50201440605	-0.000966639184324193\\
-5962.64192406299	-0.00126570460641393\\
-5956.78183371994	-0.00159952529117195\\
-5950.92174337688	-0.00196774065625192\\
-5945.06165303382	-0.00236991521549063\\
-5939.20156269076	-0.0028055389858606\\
-5933.3414723477	-0.00327402805282989\\
-5927.48138200464	-0.0037747252939739\\
-5921.62129166158	-0.00430690126063523\\
-5909.90111097546	-0.00546241633302454\\
-5898.18093028934	-0.00673333451729139\\
-5886.46074960322	-0.00811134241370701\\
-5874.7405689171	-0.00958710721534504\\
-5863.02038823099	-0.0111503412637656\\
-5851.30020754487	-0.012789875514294\\
-5833.71993651569	-0.0153658619123142\\
-5816.13966548651	-0.0180431423637368\\
-5798.55939445733	-0.0207755936994355\\
-5780.97912342815	-0.0235150732477431\\
-5763.39885239898	-0.0262122929199842\\
-5739.95849102674	-0.0296573449622676\\
-5716.5181296545	-0.0328243405983894\\
-5693.07776828226	-0.0356045307153503\\
-5669.63740691002	-0.037899288185729\\
-5646.19704553779	-0.0396237834005381\\
-5622.75668416555	-0.0407103125403666\\
-5599.31632279331	-0.0411111576362818\\
-5575.87596142107	-0.0408008691843478\\
-5552.43560004883	-0.0397778776368927\\
-5528.9952386766	-0.0380653591668352\\
-5505.55487730436	-0.0357113031012466\\
-5482.11451593212	-0.0327877528011607\\
-5464.53424490294	-0.0302774310715828\\
-5446.95397387376	-0.0275469352793039\\
-5429.37370284459	-0.0246493262921954\\
-5411.79343181541	-0.0216424565689595\\
-5394.21316078623	-0.0185880071358758\\
-5382.49298010011	-0.016557021200611\\
-5370.77279941399	-0.0145530461144421\\
-5359.05261872787	-0.0125957939561263\\
-5347.33243804175	-0.0107049802931535\\
-5335.61225735563	-0.00890015948248691\\
-5323.89207666952	-0.00720055824365951\\
-5318.03198632646	-0.00639609265514323\\
-5312.1718959834	-0.00562490882429211\\
-5306.31180564034	-0.0048892437984047\\
-5300.45171529728	-0.00419128309699356\\
-5294.59162495422	-0.00353315567387322\\
-5288.73153461116	-0.00291692894699858\\
-5282.8714442681	-0.00234460390672503\\
-5277.01135392504	-0.00181811031274662\\
-5271.15126358198	-0.00133930199060966\\
-5265.29117323892	-0.000909952237710022\\
-5259.43108289586	-0.000531749349427142\\
-5253.5709925528	-0.00020629227510786\\
-5247.71090220974	6.4913585937501e-05\\
-5241.85081186668	0.000280460438792809\\
-5235.99072152362	0.000439041988111794\\
-5230.13063118056	0.000539457164811521\\
-5224.2705408375	0.000580613675216832\\
-5218.41045049444	0.000561531363910543\\
-5212.55036015138	0.000481345381755638\\
-5206.69026980832	0.000339309150907825\\
-5200.83017946527	0.000134797118939733\\
-5194.97008912221	-0.000132692705418172\\
-5189.10999877915	-0.000463536442155097\\
-5183.24990843609	-0.000857982262230763\\
-5177.38981809303	-0.00131614854860384\\
-5171.52972774997	-0.00183802230208667\\
-5165.66963740691	-0.00242345779756414\\
-5159.80954706385	-0.00307217549572846\\
-5153.94945672079	-0.00378376121516422\\
-5148.08936637773	-0.00455766556858167\\
-5142.22927603467	-0.0053932036677192\\
-5136.36918569161	-0.00628955509950849\\
-5130.50909534855	-0.00724576417652113\\
-5124.64900500549	-0.00826074046401502\\
-5112.92882431937	-0.0104619643084524\\
-5101.20864363326	-0.012881845815389\\
-5089.48846294714	-0.0155069288851972\\
-5077.76828226102	-0.0183217482260002\\
-5066.0481015749	-0.0213089132012892\\
-5054.32792088878	-0.0244492091518728\\
-5036.7476498596	-0.0294005769326532\\
-5019.16737883043	-0.0345719762746803\\
-5001.58710780125	-0.0398795686526454\\
-4984.00683677207	-0.0452340866088032\\
-4966.42656574289	-0.0505423156890993\\
-4942.98620437065	-0.0573826692757273\\
-4919.54584299841	-0.0637433610438042\\
-4896.10548162617	-0.0694025438860658\\
-4872.66512025394	-0.0741516930284429\\
-4849.2247588817	-0.0778033315588932\\
-4825.78439750946	-0.0801984987493339\\
-4802.34403613722	-0.0812137232874689\\
-4778.90367476498	-0.080767268888778\\
-4755.46331339275	-0.078824432068685\\
-4732.02295202051	-0.0754016910161567\\
-4708.58259064827	-0.0705695301285212\\
-4691.0023196191	-0.0660940975749839\\
-4673.42204858992	-0.0609690344120997\\
-4655.84177756074	-0.0552849484671149\\
-4638.26150653156	-0.0491488528904369\\
-4620.68123550238	-0.0426831816400012\\
-4608.96105481626	-0.038256546367968\\
-4597.24087413014	-0.0337875244621813\\
-4585.52069344402	-0.0293220225675828\\
-4573.80051275791	-0.0249079484284059\\
-4562.08033207179	-0.020594933383191\\
-4556.22024172873	-0.0184922313870945\\
-4550.36015138567	-0.0164340353348914\\
-4544.50006104261	-0.0144268986882795\\
-4538.63997069955	-0.0124774238011881\\
-4532.77988035649	-0.0105922512223704\\
-4526.91979001343	-0.0087780487888698\\
-4521.05969967037	-0.00704150052832416\\
-4515.19960932731	-0.00538929538585851\\
-4509.33951898425	-0.00382811579186758\\
-4503.47942864119	-0.00236462608924002\\
-4497.61933829813	-0.0010054608368811\\
-4491.75924795507	0.000242786992088656\\
-4485.89915761201	0.00137357789895909\\
-4480.03906726895	0.00238043781621532\\
-4474.17897692589	0.00325697011347182\\
-4468.31888658283	0.00399686763332092\\
-4462.45879623977	0.00459392473771451\\
-4456.59870589671	0.00504204934483239\\
-4450.73861555366	0.0053352749374476\\
-4444.8785252106	0.00546777252190704\\
-4439.01843486754	0.00543386251869322\\
-4433.15834452448	0.00522802656357129\\
-4427.29825418142	0.00484491920031939\\
-4421.43816383836	0.00427937944383579\\
-4415.5780734953	0.00352644219481869\\
-4409.71798315224	0.00258134948502271\\
-4403.85789280918	0.00143956153411115\\
-4397.99780246612	9.67675980635635e-05\\
-4392.13771212306	-0.00145110341038938\\
-4386.27762178	-0.00320787254820711\\
-4380.41753143694	-0.00517710073465855\\
-4374.55744109388	-0.00736207892878549\\
-4368.69735075082	-0.00976581860991315\\
-4362.83726040776	-0.0123910425795004\\
-4356.97717006471	-0.0152401761018386\\
-4351.11707972165	-0.0183153384010472\\
-4345.25698937859	-0.0216183345308175\\
-4339.39689903553	-0.0251506476336813\\
-4333.53680869247	-0.0289134316054568\\
-4321.81662800635	-0.0371333404505316\\
-4310.09644732023	-0.0462804555158956\\
-4298.37626663411	-0.0563515072221691\\
-4286.65608594799	-0.067337407025305\\
-4274.93590526187	-0.0792231244319007\\
-4263.21572457576	-0.0919875972916452\\
-4245.63545354658	-0.112720914432465\\
-4228.0551825174	-0.13525153999808\\
-4210.47491148822	-0.159431415498824\\
-4192.89464045904	-0.185082409880712\\
-4169.4542790868	-0.221210061458708\\
-4146.01391771456	-0.259015019513751\\
-4122.57355634233	-0.297851811219159\\
-4099.13319497009	-0.337006267618582\\
-4075.69283359785	-0.375707748956619\\
-4046.39238188255	-0.42220602590329\\
-4017.09193016726	-0.465080906397207\\
-3987.79147845196	-0.50267429382598\\
-3958.49102673666	-0.533366552555596\\
-3929.19057502137	-0.555630751400453\\
-3899.89012330607	-0.568085780573532\\
-3870.58967159077	-0.569546432736429\\
-3841.28921987547	-0.559068617544459\\
-3817.84885850324	-0.541634026479852\\
-3794.408497131	-0.515929512003144\\
-3770.96813575876	-0.481891659052772\\
-3753.38786472958	-0.45093497498114\\
-3735.8075937004	-0.415415037318013\\
-3718.22732267122	-0.375460943051072\\
-3706.5071419851	-0.346445221558062\\
-3694.78696129899	-0.315596884907732\\
-3683.06678061287	-0.282986257974658\\
-3671.34659992675	-0.248691735036614\\
-3659.62641924063	-0.212799498214104\\
-3653.76632889757	-0.194283220345604\\
-3647.90623855451	-0.175403196746636\\
-3642.04614821145	-0.156172241528107\\
-3636.18605786839	-0.136603588761247\\
-3630.32596752533	-0.116710879352488\\
-3624.46587718227	-0.096508147431388\\
-3618.60578683921	-0.0760098062416358\\
-3612.74569649615	-0.0552306335631772\\
-3606.88560615309	-0.0341857566787052\\
-3601.02551581004	-0.012890636917385\\
-3595.16542546698	0.00863894622100714\\
-3589.30533512392	0.0303869113242634\\
-3583.44524478086	0.0523368916980917\\
-3577.5851544378	0.0744722523392822\\
-3571.72506409474	0.0967761072337489\\
-3565.86497375168	0.119231336952416\\
-3560.00488340862	0.141820606527005\\
-3554.14479306556	0.164526383572827\\
-3548.2847027225	0.187330956644857\\
-3542.42461237944	0.210216453800139\\
-3536.56452203638	0.233164861330508\\
-3524.84434135026	0.279177757358764\\
-3513.12416066415	0.325223420805006\\
-3501.40397997803	0.371154580007798\\
-3489.68379929191	0.416823532001606\\
-3472.10352826273	0.484512828700304\\
-3454.52325723355	0.550786572764059\\
-3436.94298620437	0.615159537146415\\
-3413.50262483213	0.697219797801138\\
-3390.06226345989	0.773994468866484\\
-3366.62190208766	0.844486334270942\\
-3337.32145037236	0.922389205523877\\
-3308.02099865706	0.987508530355332\\
-3278.72054694176	1.03857193114877\\
-3249.42009522647	1.07461235805407\\
-3220.11964351117	1.09498419145606\\
-3190.81919179587	1.09936990030888\\
-3155.65864973752	1.08356456399332\\
-3126.35819802222	1.05324975583795\\
-3097.05774630692	1.00805215848877\\
-3067.75729459163	0.948907070119819\\
-3038.45684287633	0.876979340563779\\
-3015.01648150409	0.811140388467105\\
-2991.57612013185	0.738753798811626\\
-2968.13575875961	0.660636423108755\\
-2950.55548773043	0.598799696405538\\
-2932.97521670126	0.534587711230066\\
-2915.39494567208	0.468373785426526\\
-2903.67476498596	0.423304094248923\\
-2891.95458429984	0.377622132326263\\
-2880.23440361372	0.331438200414991\\
-2868.5142229276	0.284861713767814\\
-2856.79404224148	0.238000904769999\\
-2850.93395189842	0.214497309746409\\
-2845.07386155537	0.190962544938027\\
-2839.21377121231	0.167409686196695\\
-2833.35368086925	0.143851686950385\\
-2827.49359052619	0.120301371266129\\
-2821.63350018313	0.096771427256034\\
-2815.77340984007	0.0732744008047981\\
-2809.91331949701	0.0498226896368607\\
-2804.05322915395	0.0264285377076926\\
-2798.19313881089	0.00310402993171366\\
-2792.33304846783	-0.0201389127631252\\
-2786.47295812477	-0.0432885380521901\\
-2780.61286778171	-0.0663332665046747\\
-2774.75277743865	-0.0892616955546188\\
-2768.89268709559	-0.112062603149251\\
-2763.03259675254	-0.134724951087915\\
-2757.17250640948	-0.15723788804906\\
-2751.31241606642	-0.179590752309485\\
-2745.45232572336	-0.201773074171616\\
-2739.5922353803	-0.223774578095359\\
-2727.87205469418	-0.267195011551427\\
-2716.15187400806	-0.309773794772448\\
-2704.43169332194	-0.351435866303824\\
-2692.71151263582	-0.392109382374847\\
-2675.13124160664	-0.451116888696373\\
-2657.55097057746	-0.507528458285575\\
-2639.97069954828	-0.561142139615882\\
-2616.53033817605	-0.627955264149058\\
-2593.08997680381	-0.689066371958402\\
-2569.64961543157	-0.74412334775921\\
-2540.34916371627	-0.803966438612835\\
-2511.04871200098	-0.853364159186347\\
-2481.74826028568	-0.891905295019938\\
-2446.58771822732	-0.923380677160696\\
-2411.42717616897	-0.93838352068375\\
-2376.26663411061	-0.936710277289482\\
-2341.10609205225	-0.918340945436338\\
-2305.9455499939	-0.883448279689694\\
-2276.6450982786	-0.842016394828914\\
-2247.3446465633	-0.789707646561303\\
-2218.044194848	-0.726975632300107\\
-2194.60383347576	-0.669668258444039\\
-2171.16347210353	-0.606410453249348\\
-2147.72311073129	-0.537610415757817\\
-2130.14283970211	-0.482649768977804\\
-2112.56256867293	-0.425046521524711\\
-2100.84238798681	-0.385290474439117\\
-2089.1222073007	-0.344536139459104\\
-2077.40202661458	-0.302859875357715\\
-2065.68184592846	-0.260341198076615\\
-2059.8217555854	-0.238791641223069\\
-2053.96166524234	-0.217062676501985\\
-2048.10157489928	-0.195165082091855\\
-2042.24148455622	-0.173109810811487\\
-2036.38139421316	-0.150907985369714\\
-2030.5213038701	-0.128570893332231\\
-2024.66121352704	-0.106109981819477\\
-2018.80112318398	-0.0835368519419716\\
-2012.94103284092	-0.0608632529603804\\
-2007.08094249787	-0.0381010761922624\\
-2001.22085215481	-0.0152623486548387\\
-1995.36076181175	0.00764077354533821\\
-1989.50067146869	0.0305960120719594\\
-1983.64058112563	0.0535909734650745\\
-1977.78049078257	0.0766131564837794\\
-1971.92040043951	0.0996499596837526\\
-1966.06031009645	0.122688689220735\\
-1960.20021975339	0.145716566871446\\
-1954.34012941033	0.168720738271243\\
-1948.48003906727	0.191688281346292\\
-1942.61994872421	0.214606214952346\\
-1936.75985838115	0.237461507690922\\
-1925.03967769503	0.282931847848651\\
-1913.31949700891	0.327994391559696\\
-1901.59931632279	0.372544016865965\\
-1889.87913563667	0.416475699368519\\
-1872.2988646075	0.480985937887762\\
-1854.71859357832	0.543520996748328\\
-1837.13832254914	0.603737485207334\\
-1813.6979611769	0.67984259998767\\
-1790.25759980466	0.750471012333545\\
-1766.81723843242	0.814907356671754\\
-1737.51678671713	0.885746226967516\\
-1708.21633500183	0.944740483360737\\
-1678.91588328653	0.9909303533315\\
-1649.61543157124	1.0235611680713\\
-1614.45488951288	1.04408629989649\\
-1579.29434745452	1.04385550840326\\
-1544.13380539617	1.02288928178193\\
-1514.83335368087	0.989933789663105\\
-1485.53290196557	0.943476070394558\\
-1456.23245025027	0.884279955332726\\
-1426.93199853498	0.813305426779727\\
-1403.49163716274	0.748802068986713\\
-1380.0512757905	0.678133481277436\\
-1356.61091441826	0.602001650428078\\
-1339.03064338909	0.541763999266042\\
-1321.45037235991	0.479197912944355\\
-1303.87010133073	0.414638066700273\\
-1292.14992064461	0.370658497079378\\
-1280.42973995849	0.326045931480171\\
-1268.70955927237	0.280902656332863\\
-1256.98937858626	0.235331102097145\\
-1251.1292882432	0.212416702411962\\
-1245.26919790014	0.189433572519851\\
-1239.40910755708	0.16639443766229\\
-1233.54901721402	0.143311982992198\\
-1227.68892687096	0.120198845958366\\
-1221.8288365279	0.0970676088605646\\
-1215.96874618484	0.0739307915839076\\
-1210.10865584178	0.0508008445204743\\
-1204.24856549872	0.0276901416734756\\
-1198.38847515566	0.00461097396102351\\
-1192.5283848126	-0.0184244572892841\\
-1186.66829446954	-0.0414040466446187\\
-1180.80820412648	-0.0643157906693696\\
-1174.94811378342	-0.0871477938224896\\
-1169.08802344036	-0.109888274150354\\
-1163.2279330973	-0.132525568771261\\
-1157.36784275424	-0.155048139153195\\
-1151.50775241118	-0.177444576173161\\
-1145.64766206812	-0.199703604986076\\
-1139.78757172506	-0.221814089656128\\
-1133.927481382	-0.243765037614798\\
-1122.20730069589	-0.287145095045285\\
-1110.48712000977	-0.329758859409109\\
-1098.76693932365	-0.371523955943757\\
-1087.04675863753	-0.412360659128184\\
-1069.46648760835	-0.471707347331137\\
-1051.88621657917	-0.528544775894659\\
-1034.30594554999	-0.58263991915285\\
-1010.86558417775	-0.650126723400459\\
-987.425222805519	-0.711879422613491\\
-963.98486143328	-0.767474173442195\\
-934.684409717982	-0.827743097312675\\
-905.383958002687	-0.877188857761063\\
-876.083506287388	-0.915311652346762\\
-840.92296422903	-0.945585980803626\\
-805.762422170676	-0.958582884712484\\
-770.601880112317	-0.954115110802406\\
-735.441338053963	-0.932239374418993\\
-700.280795995604	-0.893258152974448\\
-670.980344280306	-0.848097540576334\\
-641.679892565011	-0.79189033012327\\
-612.379440849712	-0.725214044692156\\
-588.939079477474	-0.664803222484261\\
-565.498718105238	-0.598556933139792\\
-547.918447076059	-0.545322583025903\\
-530.33817604688	-0.489284402511263\\
-512.757905017701	-0.4306740808352\\
-501.037724331585	-0.390292825658711\\
-489.317543645466	-0.348954505945881\\
-477.597362959346	-0.30673863218497\\
-465.877182273227	-0.263727393015123\\
-460.017091930167	-0.24194989970827\\
-454.157001587108	-0.220005540588357\\
-448.296911244048	-0.197905289303377\\
-442.436820900988	-0.175660263084189\\
-436.576730557928	-0.153281718031439\\
-430.716640214869	-0.130781044220187\\
-424.856549871809	-0.108169760600926\\
-418.996459528749	-0.0854595096805506\\
-413.136369185693	-0.0626620520201019\\
-407.276278842633	-0.0397892605221325\\
-401.416188499574	-0.0168531145238133\\
-395.556098156514	0.0061343063085423\\
-389.696007813454	0.0291608282749697\\
-383.835917470395	0.0522141901281179\\
-377.975827127335	0.075282049745088\\
-372.115736784275	0.0983519909824909\\
-366.255646441216	0.121411530709818\\
-360.395556098156	0.144448126009068\\
-354.535465755096	0.167449181552656\\
-348.675375412036	0.190402057137492\\
-342.815285068977	0.213294075374845\\
-336.955194725917	0.236112529531225\\
-325.235014039801	0.281477819998943\\
-313.514833353682	0.326395994598654\\
-301.794652667562	0.370765174506946\\
-290.074471981443	0.414483809178651\\
-272.494200952264	0.478621412341787\\
-254.913929923085	0.540731761689403\\
-237.333658893909	0.600484956266981\\
-213.89329752167	0.675935994324994\\
-190.452936149431	0.745898213831991\\
-167.012574777196	0.809687034745988\\
-137.712123061898	0.879785753679549\\
-108.411671346599	0.938161383224601\\
-79.1112196313043	0.983892347470839\\
-49.8107679160057	1.01625076655098\\
-14.6502258576475	1.03671504386981\\
20.5103162007072	1.03671504386981\\
55.6708582590654	1.01625076655098\\
84.9713099743603	0.983892347470839\\
114.271761689659	0.938161383224601\\
143.572213404957	0.879785753679549\\
172.872665120252	0.809687034745988\\
196.313026492491	0.745898213831991\\
219.75338786473	0.675935994324994\\
243.193749236969	0.600484956266981\\
260.774020266144	0.540731761689403\\
278.354291295323	0.478621412341787\\
295.934562324503	0.414483809178651\\
307.654743010622	0.370765174506946\\
319.374923696741	0.326395994598654\\
331.095104382861	0.281477819998943\\
342.815285068977	0.236112529531225\\
348.675375412036	0.213294075374845\\
354.535465755096	0.190402057137492\\
360.395556098156	0.167449181552656\\
366.255646441216	0.144448126009068\\
372.115736784275	0.121411530709818\\
377.975827127335	0.0983519909824909\\
383.835917470395	0.075282049745088\\
389.696007813454	0.0522141901281179\\
395.556098156514	0.0291608282749697\\
401.416188499574	0.0061343063085423\\
407.276278842633	-0.0168531145238133\\
413.136369185693	-0.0397892605221325\\
418.996459528749	-0.0626620520201019\\
424.856549871809	-0.0854595096805506\\
430.716640214869	-0.108169760600926\\
436.576730557928	-0.130781044220187\\
442.436820900988	-0.153281718031439\\
448.296911244048	-0.175660263084189\\
454.157001587108	-0.197905289303377\\
460.017091930167	-0.220005540588357\\
465.877182273227	-0.24194989970827\\
477.597362959346	-0.285327194918679\\
489.317543645466	-0.327951188018543\\
501.037724331585	-0.369738393487153\\
512.757905017701	-0.410607948229538\\
530.33817604688	-0.470021541141903\\
547.918447076059	-0.526943600588813\\
565.498718105238	-0.581137353262501\\
588.939079477474	-0.648771293214524\\
612.379440849712	-0.710682120815374\\
635.819802221951	-0.766439567974595\\
665.12025393725	-0.826910699068304\\
694.420705652545	-0.876551286936023\\
723.721157367843	-0.914858358479798\\
758.881699426198	-0.945340126926624\\
794.042241484556	-0.95853426742479\\
829.202783542914	-0.954261524511677\\
864.363325601269	-0.932587269953778\\
899.523867659627	-0.89382068239106\\
928.824319374922	-0.848850867708293\\
958.124771090221	-0.792843515319611\\
987.425222805519	-0.726370878342324\\
1010.86558417775	-0.666119878518875\\
1034.30594554999	-0.600024143061596\\
1051.88621657917	-0.546892261632748\\
1069.46648760835	-0.490943779711676\\
1087.04675863753	-0.432406696698731\\
1098.76693932365	-0.39206323341799\\
1110.48712000977	-0.350752537250848\\
1122.20730069589	-0.308552973193088\\
1133.927481382	-0.265545603869687\\
1139.78757172506	-0.243765037614798\\
1145.64766206812	-0.221814089656128\\
1151.50775241118	-0.199703604986076\\
1157.36784275424	-0.177444576173161\\
1163.2279330973	-0.155048139153195\\
1169.08802344036	-0.132525568771261\\
1174.94811378342	-0.109888274150354\\
1180.80820412648	-0.0871477938224896\\
1186.66829446954	-0.0643157906693696\\
1192.5283848126	-0.0414040466446187\\
1198.38847515566	-0.0184244572892841\\
1204.24856549872	0.00461097396102351\\
1210.10865584178	0.0276901416734756\\
1215.96874618484	0.0508008445204743\\
1221.8288365279	0.0739307915839076\\
1227.68892687096	0.0970676088605646\\
1233.54901721402	0.120198845958366\\
1239.40910755708	0.143311982992198\\
1245.26919790014	0.16639443766229\\
1251.1292882432	0.189433572519851\\
1256.98937858626	0.212416702411962\\
1262.84946892931	0.235331102097145\\
1274.56964961543	0.280902656332863\\
1286.28983030155	0.326045931480171\\
1298.01001098767	0.370658497079378\\
1309.73019167379	0.414638066700273\\
1327.31046270297	0.479197912944355\\
1344.89073373214	0.541763999266042\\
1362.47100476132	0.602001650428078\\
1385.91136613356	0.678133481277436\\
1409.3517275058	0.748802068986713\\
1432.79208887804	0.813305426779727\\
1462.09254059333	0.884279955332726\\
1491.39299230863	0.943476070394558\\
1520.69344402393	0.989933789663105\\
1549.99389573923	1.02288928178193\\
1585.15443779758	1.04385550840326\\
1620.31497985594	1.04408629989649\\
1655.4755219143	1.0235611680713\\
1684.77597362959	0.9909303533315\\
1714.07642534489	0.944740483360737\\
1743.37687706019	0.885746226967516\\
1772.67732877548	0.814907356671754\\
1796.11769014772	0.750471012333545\\
1819.55805151996	0.67984259998767\\
1842.9984128922	0.603737485207334\\
1860.57868392138	0.543520996748328\\
1878.15895495056	0.480985937887762\\
1895.73922597973	0.416475699368519\\
1907.45940666585	0.372544016865965\\
1919.17958735197	0.327994391559696\\
1930.89976803809	0.282931847848651\\
1942.61994872421	0.237461507690922\\
1948.48003906727	0.214606214952346\\
1954.34012941033	0.191688281346292\\
1960.20021975339	0.168720738271243\\
1966.06031009645	0.145716566871446\\
1971.92040043951	0.122688689220735\\
1977.78049078257	0.0996499596837526\\
1983.64058112563	0.0766131564837794\\
1989.50067146869	0.0535909734650745\\
1995.36076181175	0.0305960120719594\\
2001.22085215481	0.00764077354533821\\
2007.08094249787	-0.0152623486548387\\
2012.94103284092	-0.0381010761922624\\
2018.80112318398	-0.0608632529603804\\
2024.66121352704	-0.0835368519419716\\
2030.5213038701	-0.106109981819477\\
2036.38139421316	-0.128570893332231\\
2042.24148455622	-0.150907985369714\\
2048.10157489928	-0.173109810811487\\
2053.96166524234	-0.195165082091855\\
2059.8217555854	-0.217062676501985\\
2065.68184592846	-0.238791641223069\\
2077.40202661458	-0.281700748023349\\
2089.1222073007	-0.323808351663822\\
2100.84238798681	-0.365033395610847\\
2112.56256867293	-0.405297930566587\\
2130.14283970211	-0.46373116616562\\
2147.72311073129	-0.519597485589205\\
2165.30338176047	-0.572675563010941\\
2188.74374313271	-0.638758927273433\\
2212.18410450494	-0.699089135234759\\
2235.62446587718	-0.753285171235446\\
2264.92491759248	-0.811910015576784\\
2294.22536930778	-0.859915808117668\\
2329.38591136613	-0.902941341497078\\
2364.54645342449	-0.92960648866904\\
2399.70699548285	-0.939637041934716\\
2434.8675375412	-0.932958121884436\\
2470.02807959956	-0.909682291547193\\
2505.18862165792	-0.870103296113017\\
2534.48907337321	-0.825004112213072\\
2563.78952508851	-0.769283741039823\\
2593.08997680381	-0.703410949366457\\
2616.53033817605	-0.643782141835267\\
2639.97069954828	-0.578358991479302\\
2657.55097057746	-0.525720076236175\\
2675.13124160664	-0.470219484359756\\
2692.71151263582	-0.41205386256751\\
2704.43169332194	-0.371900558051928\\
2716.15187400806	-0.330723985293277\\
2727.87205469418	-0.2885943760248\\
2739.5922353803	-0.24558518454382\\
2745.45232572336	-0.223774578095359\\
2751.31241606642	-0.201773074171616\\
2757.17250640948	-0.179590752309485\\
2763.03259675254	-0.15723788804906\\
2768.89268709559	-0.134724951087915\\
2774.75277743865	-0.112062603149251\\
2780.61286778171	-0.0892616955546188\\
2786.47295812477	-0.0663332665046747\\
2792.33304846783	-0.0432885380521901\\
2798.19313881089	-0.0201389127631252\\
2804.05322915395	0.00310402993171366\\
2809.91331949701	0.0264285377076926\\
2815.77340984007	0.0498226896368607\\
2821.63350018313	0.0732744008047981\\
2827.49359052619	0.096771427256034\\
2833.35368086925	0.120301371266129\\
2839.21377121231	0.143851686950385\\
2845.07386155537	0.167409686196695\\
2850.93395189842	0.190962544938027\\
2856.79404224148	0.214497309746409\\
2862.65413258454	0.238000904769999\\
2874.37431327066	0.284861713767814\\
2886.09449395678	0.331438200414991\\
2897.8146746429	0.377622132326263\\
2909.53485532902	0.423304094248923\\
2927.1151263582	0.490644395715934\\
2944.69539738738	0.556232637070602\\
2962.27566841655	0.619694194736366\\
2985.71602978879	0.70035884027241\\
3009.15639116103	0.775715496874331\\
3032.59675253327	0.84492759849302\\
3061.89720424857	0.921597776224929\\
3091.19765596386	0.986007350486606\\
3120.49810767916	1.03690982273399\\
3149.79855939446	1.07327386384239\\
3179.09901110976	1.09431803521871\\
3214.25955316811	1.09866357258508\\
3243.56000488341	1.08466646787492\\
3272.8604565987	1.05483432904991\\
3302.160908314	1.00968181278724\\
3331.4613600293	0.950048800228693\\
3360.7618117446	0.877088071806481\\
3384.20217311684	0.810084836706592\\
3407.64253448907	0.736332031887184\\
3431.08289586131	0.656785040048408\\
3448.66316689049	0.593942541177733\\
3466.24343791967	0.528882787973249\\
3483.82370894885	0.462082759246009\\
3495.54388963497	0.416823532001606\\
3507.26407032109	0.371154580007798\\
3518.9842510072	0.325223420805006\\
3530.70443169332	0.279177757358764\\
3542.42461237944	0.233164861330508\\
3548.2847027225	0.210216453800139\\
3554.14479306556	0.187330956644857\\
3560.00488340862	0.164526383572827\\
3565.86497375168	0.141820606527005\\
3571.72506409474	0.119231336952416\\
3577.5851544378	0.0967761072337489\\
3583.44524478086	0.0744722523392822\\
3589.30533512392	0.0523368916980917\\
3595.16542546698	0.0303869113242634\\
3601.02551581004	0.00863894622100714\\
3606.88560615309	-0.012890636917385\\
3612.74569649615	-0.0341857566787052\\
3618.60578683921	-0.0552306335631772\\
3624.46587718227	-0.0760098062416358\\
3630.32596752533	-0.096508147431388\\
3636.18605786839	-0.116710879352488\\
3642.04614821145	-0.136603588761247\\
3647.90623855451	-0.156172241528107\\
3653.76632889757	-0.175403196746636\\
3659.62641924063	-0.194283220345604\\
3665.48650958369	-0.212799498214104\\
3677.20669026981	-0.248691735036614\\
3688.92687095593	-0.282986257974658\\
3700.64705164204	-0.315596884907732\\
3712.36723232816	-0.346445221558062\\
3729.94750335734	-0.389261863212917\\
3747.52777438652	-0.427754991910525\\
3765.1080454157	-0.461765987232056\\
3782.68831644488	-0.491180624234413\\
3806.12867781712	-0.523136771096618\\
3829.56903918935	-0.546762493786754\\
3853.00940056159	-0.56215544655377\\
3882.30985227689	-0.570182974569313\\
3911.61030399219	-0.566439601074606\\
3940.91075570748	-0.551920965364288\\
3970.21120742278	-0.527855253501815\\
3999.51165913808	-0.495656870274581\\
4028.81211085337	-0.456875600906461\\
4058.11256256867	-0.413143114833487\\
4087.41301428397	-0.366118727813539\\
4110.85337565621	-0.327229177393526\\
4134.29373702845	-0.28808386933809\\
4157.73409840068	-0.249440571688767\\
4181.17445977292	-0.211997351839675\\
4204.61482114516	-0.176380825722701\\
4222.19509217434	-0.151198593510241\\
4239.77536320352	-0.12755014350599\\
4257.3556342327	-0.10560367653161\\
4274.93590526187	-0.0854970701767792\\
4286.65608594799	-0.0731689627272707\\
4298.37626663411	-0.0617309000364902\\
4310.09644732023	-0.0512009193400539\\
4321.81662800635	-0.0415910668354822\\
4333.53680869247	-0.0329075041805321\\
4339.39689903553	-0.0289134316054568\\
4345.25698937859	-0.0251506476336813\\
4351.11707972165	-0.0216183345308175\\
4356.97717006471	-0.0183153384010472\\
4362.83726040776	-0.0152401761018386\\
4368.69735075082	-0.0123910425795004\\
4374.55744109388	-0.00976581860991315\\
4380.41753143694	-0.00736207892878549\\
4386.27762178	-0.00517710073465855\\
4392.13771212306	-0.00320787254820711\\
4397.99780246612	-0.00145110341038938\\
4403.85789280918	9.67675980635635e-05\\
4409.71798315224	0.00143956153411115\\
4415.5780734953	0.00258134948502271\\
4421.43816383836	0.00352644219481869\\
4427.29825418142	0.00427937944383579\\
4433.15834452448	0.00484491920031939\\
4439.01843486754	0.00522802656357129\\
4444.8785252106	0.00543386251869322\\
4450.73861555366	0.00546777252190704\\
4456.59870589671	0.0053352749374476\\
4462.45879623977	0.00504204934483239\\
4468.31888658283	0.00459392473771451\\
4474.17897692589	0.00399686763332092\\
4480.03906726895	0.00325697011347182\\
4485.89915761201	0.00238043781621532\\
4491.75924795507	0.00137357789895909\\
4497.61933829813	0.000242786992088656\\
4503.47942864119	-0.0010054608368811\\
4509.33951898425	-0.00236462608924002\\
4515.19960932731	-0.00382811579186758\\
4521.05969967037	-0.00538929538585851\\
4526.91979001343	-0.00704150052832416\\
4532.77988035649	-0.0087780487888698\\
4538.63997069955	-0.0105922512223704\\
4544.50006104261	-0.0124774238011881\\
4550.36015138567	-0.0144268986882795\\
4556.22024172873	-0.0164340353348914\\
4562.08033207179	-0.0184922313870945\\
4567.94042241485	-0.020594933383191\\
4579.66060310097	-0.0249079484284059\\
4591.38078378708	-0.0293220225675828\\
4603.1009644732	-0.0337875244621813\\
4614.82114515932	-0.038256546367968\\
4626.54132584544	-0.0426831816400012\\
4644.12159687462	-0.0491488528904369\\
4661.7018679038	-0.0552849484671149\\
4679.28213893298	-0.0609690344120997\\
4696.86240996215	-0.0660940975749839\\
4714.44268099133	-0.0705695301285212\\
4737.88304236357	-0.0754016910161567\\
4761.32340373581	-0.078824432068685\\
4784.76376510804	-0.080767268888778\\
4808.20412648028	-0.0812137232874689\\
4831.64448785252	-0.0801984987493339\\
4855.08484922476	-0.0778033315588932\\
4878.525210597	-0.0741516930284429\\
4901.96557196923	-0.0694025438860658\\
4925.40593334147	-0.0637433610438042\\
4948.84629471371	-0.0573826692757273\\
4972.28665608595	-0.0505423156890993\\
4989.86692711513	-0.0452340866088032\\
5007.44719814431	-0.0398795686526454\\
5025.02746917348	-0.0345719762746803\\
5042.60774020266	-0.0294005769326532\\
5060.18801123184	-0.0244492091518728\\
5071.90819191796	-0.0213089132012892\\
5083.62837260408	-0.0183217482260002\\
5095.3485532902	-0.0155069288851972\\
5107.06873397632	-0.012881845815389\\
5118.78891466243	-0.0104619643084524\\
5130.50909534855	-0.00826074046401502\\
5136.36918569161	-0.00724576417652113\\
5142.22927603467	-0.00628955509950849\\
5148.08936637773	-0.0053932036677192\\
5153.94945672079	-0.00455766556858167\\
5159.80954706385	-0.00378376121516422\\
5165.66963740691	-0.00307217549572846\\
5171.52972774997	-0.00242345779756414\\
5177.38981809303	-0.00183802230208667\\
5183.24990843609	-0.00131614854860384\\
5189.10999877915	-0.000857982262230763\\
5194.97008912221	-0.000463536442155097\\
5200.83017946527	-0.000132692705418172\\
5206.69026980832	0.000134797118939733\\
5212.55036015138	0.000339309150907825\\
5218.41045049444	0.000481345381755638\\
5224.2705408375	0.000561531363910543\\
5230.13063118056	0.000580613675216832\\
5235.99072152362	0.000539457164811521\\
5241.85081186668	0.000439041988111794\\
5247.71090220974	0.000280460438792809\\
5253.5709925528	6.4913585937501e-05\\
5259.43108289586	-0.00020629227510786\\
5265.29117323892	-0.000531749349427142\\
5271.15126358198	-0.000909952237710022\\
5277.01135392504	-0.00133930199060966\\
5282.8714442681	-0.00181811031274662\\
5288.73153461116	-0.00234460390672503\\
5294.59162495422	-0.00291692894699858\\
5300.45171529728	-0.00353315567387322\\
5306.31180564034	-0.00419128309699356\\
5312.1718959834	-0.0048892437984047\\
5318.03198632646	-0.00562490882429211\\
5323.89207666952	-0.00639609265514323\\
5335.61225735563	-0.00803602210988691\\
5347.33243804175	-0.00979060947598667\\
5359.05261872787	-0.011640854441922\\
5370.77279941399	-0.0135673456095694\\
5382.49298010011	-0.0155504272908981\\
5394.21316078623	-0.0175703652333819\\
5411.79343181541	-0.0206264570249625\\
5429.37370284459	-0.0236561994680502\\
5446.95397387376	-0.0265968666337916\\
5464.53424490294	-0.0293892177415632\\
5482.11451593212	-0.0319784943625235\\
5499.6947869613	-0.0343153102349622\\
5523.13514833354	-0.0369645697970126\\
5546.57550970577	-0.0390053866662389\\
5570.01587107801	-0.0403775975488117\\
5593.45623245025	-0.0410455768274131\\
5616.89659382249	-0.0409985804708773\\
5640.33695519473	-0.0402502743376171\\
5663.77731656696	-0.0388374875334381\\
5687.2176779392	-0.0368182547657073\\
5710.65803931144	-0.0342692322365158\\
5734.09840068368	-0.0312825889960758\\
5757.53876205592	-0.0279624893494367\\
5775.11903308509	-0.0253209774353723\\
5792.69930411427	-0.0226041612950089\\
5810.27957514345	-0.0198615860907476\\
5827.85984617263	-0.0171422580846296\\
5845.44011720181	-0.0144937414268623\\
5863.02038823099	-0.0119613011743077\\
5874.7405689171	-0.0103584689113255\\
5886.46074960322	-0.00883762661927349\\
5898.18093028934	-0.00740951237153606\\
5909.90111097546	-0.00608394503761024\\
5921.62129166158	-0.00486975521662219\\
5927.48138200464	-0.00430690126063523\\
5933.3414723477	-0.0037747252939739\\
5939.20156269076	-0.00327402805282989\\
5945.06165303382	-0.0028055389858606\\
5950.92174337688	-0.00236991521549063\\
5956.78183371994	-0.00196774065625192\\
5962.64192406299	-0.00159952529117195\\
5968.50201440605	-0.00126570460641393\\
5974.36210474911	-0.000966639184324193\\
5980.22219509217	-0.000702614454714258\\
5986.08228543523	-0.000473840603818683\\
5991.94237577829	-0.000280452640003633\\
5997.80246612135	-0.000122510615126645\\
nan	nan\\
6009.52264680747	8.71677877848544e-05\\
6015.38273715053	0.00013915470579904\\
6021.24282749359	0.000156195260852918\\
6027.10291783665	0.000138595537882602\\
6032.96300817971	8.67320865952666e-05\\
6038.82309852277	1.05066909580283e-06\\
6044.68318886583	-0.000117935128014562\\
6050.54327920889	-0.000269645413337205\\
6056.40336955195	-0.000453435638175552\\
6062.26345989501	-0.000668598367619746\\
6068.12355023807	-0.000914365170761012\\
6073.98364058113	-0.00118990862541044\\
6079.84373092419	-0.00149434443234112\\
6085.70382126725	-0.00182673363421009\\
6091.5639116103	-0.00218608493357806\\
6097.42400195336	-0.0025713571050529\\
6103.28409229642	-0.00298146149550286\\
6109.14418263948	-0.00341526460706999\\
6115.00427298254	-0.00387159075676286\\
6126.72445366866	-0.00484691496030564\\
6138.44463435478	-0.00589729026377876\\
6150.1648150409	-0.00701207879060207\\
6161.88499572702	-0.0081802436172531\\
6173.60517641314	-0.00939044587901222\\
6185.32535709926	-0.0106311431676606\\
6202.90562812843	-0.0125238849122458\\
6220.48589915761	-0.0144198008258561\\
6238.06617018679	-0.0162803607417812\\
6255.64644121597	-0.0180684446572701\\
6273.22671224515	-0.0197490117059137\\
6296.66707361738	-0.0217670759174429\\
6320.10743498962	-0.0234694346117178\\
6343.54779636186	-0.024800895361738\\
6366.9881577341	-0.0257200549059076\\
6390.42851910634	-0.0262003496211861\\
6413.86888047857	-0.026230592106393\\
6437.30924185081	-0.0258149879582533\\
6460.74960322305	-0.0249726439197466\\
6484.18996459529	-0.0237365949305218\\
6507.63032596753	-0.0221523927179412\\
6531.07068733976	-0.0202763119640315\\
6554.511048712	-0.0181732413749578\\
6572.09131974118	-0.0164896149110248\\
6589.67159077036	-0.0147497797944119\\
6607.25186179954	-0.0129858417428408\\
6624.83213282871	-0.0112298733508544\\
6642.41240385789	-0.00951331145706943\\
6659.99267488707	-0.00786637594471415\\
6671.71285557319	-0.00682123976455288\\
6683.43303625931	-0.00582763032354333\\
6695.15321694543	-0.00489292472691127\\
6706.87339763155	-0.00402393187227838\\
6718.59357831766	-0.00322684116990391\\
6724.45366866072	-0.0028570125857194\\
6730.31375900378	-0.00250717696029412\\
6736.17384934684	-0.00217791442732509\\
6742.0339396899	-0.00186975895010211\\
6747.89403003296	-0.00158319743259294\\
6753.75412037602	-0.00131866893334073\\
6759.61421071908	-0.00107656398377321\\
6765.47430106214	-0.000857224011372839\\
6771.3343914052	-0.000660940868717434\\
6777.19448174826	-0.00048795646860541\\
6783.05457209132	-0.000338462525666669\\
6788.91466243438	-0.000212600404349719\\
6794.77475277744	-0.000110461073133609\\
6800.6348431205	-3.20851645051866e-05\\
6806.49493346356	2.25368599307568e-05\\
6812.35502380662	5.34644400932771e-05\\
6818.21511414968	6.08065448811398e-05\\
6824.07520449274	4.47211977396864e-05\\
6829.9352948358	5.41490181981804e-06\\
6835.79538517886	-5.68580335656158e-05\\
6841.65547552191	-0.000141796263310171\\
6847.51556586497	-0.000249052271612713\\
6853.37565620803	-0.000378233329560688\\
6859.23574655109	-0.000528902542717993\\
6865.09583689415	-0.000700579986113411\\
6870.95592723721	-0.000892743923840465\\
6876.81601758027	-0.00110483211031015\\
6882.67610792333	-0.00133624317001102\\
6888.53619826639	-0.00158633805257091\\
6894.39628860945	-0.00185444155975373\\
6900.25637895251	-0.00213984394061383\\
6906.11646929557	-0.00244180255162248\\
6917.83664998169	-0.00309226380890344\\
6929.55683066781	-0.00379929311741169\\
6941.27701135393	-0.00455594769407947\\
6952.99719204004	-0.00535493635989402\\
6964.71737272616	-0.00618868403251014\\
6976.43755341228	-0.00704939798005758\\
6994.01782444146	-0.00837363117112833\\
7011.59809547064	-0.00971356553958538\\
7029.17836649982	-0.0110422187516779\\
7046.75863752899	-0.0123332146626785\\
7064.33890855817	-0.0135612713315236\\
7081.91917958735	-0.0147026617481923\\
7105.35954095959	-0.0160524478550469\\
7128.79990233183	-0.017164911640962\\
7152.24026370407	-0.0180039994718255\\
7175.6806250763	-0.0185435081848604\\
7199.12098644854	-0.0187677824095412\\
7222.56134782078	-0.0186720486407515\\
7246.00170919302	-0.0182623814027911\\
7269.44207056526	-0.0175553090788418\\
7292.88243193749	-0.0165770787554313\\
7316.32279330973	-0.0153626103708321\\
7339.76315468197	-0.0139541802068539\\
7363.20351605421	-0.0123998820348884\\
7380.78378708338	-0.011169733244554\\
7398.36405811256	-0.0099098093431418\\
7415.94432914174	-0.00864335340509519\\
7433.52460017092	-0.00739345189657561\\
7451.1048712001	-0.00618260106409093\\
7462.82505188622	-0.00540772731425485\\
7474.54523257233	-0.00466582556340511\\
7486.26541325845	-0.00396260676462862\\
7497.98559394457	-0.00330341083755303\\
7509.70577463069	-0.00269316478848812\\
7521.42595531681	-0.00213634463405881\\
7527.28604565987	-0.00187923625356278\\
7533.14613600293	-0.00163694141651033\\
7539.00622634599	-0.001409878362134\\
7544.86631668905	-0.00119843156113038\\
7550.72640703211	-0.00100295104139083\\
7556.58649737517	-0.000823751788772461\\
7562.44658771823	-0.000661113224018044\\
7568.30667806129	-0.000515278756450276\\
7574.16676840435	-0.000386455415196423\\
7580.02685874741	-0.000274813558354109\\
7585.88694909047	-0.0001804866603871\\
7591.74703943353	-0.000103571177926833\\
7597.60712977658	-4.41264938878022e-05\\
7603.46722011964	-2.17493976741762e-06\\
7609.3273104627	2.22981042770951e-05\\
7615.18740080576	2.93440320654991e-05\\
7621.04749114882	1.90507574305546e-05\\
7626.90758149188	-8.45761550854606e-06\\
7632.76767183494	-5.30211968507886e-05\\
7638.627762178	-0.000114444789048414\\
7644.48785252106	-0.000192498437387347\\
7650.34794286412	-0.000286918051360987\\
7656.20803320718	-0.000397406095374902\\
7662.06812355024	-0.000523632347051589\\
7667.9282138933	-0.000665234721270361\\
7673.78830423636	-0.000821820157997939\\
7679.64839457942	-0.000992965571739724\\
7685.50848492248	-0.00117821886042376\\
7691.36857526554	-0.00137709997138264\\
7697.2286656086	-0.00158910202191079\\
7703.08875595166	-0.00181369247194271\\
7714.80893663777	-0.00229838750201581\\
7726.52911732389	-0.00282645917370073\\
7738.24929801001	-0.00339285388872418\\
7749.96947869613	-0.0039922353658527\\
7761.68965938225	-0.0046190316944184\\
7773.40984006837	-0.0052674840045619\\
7790.99011109754	-0.00626784498545595\\
7808.57038212672	-0.00728343591455004\\
7826.1506531559	-0.00829399605778873\\
7843.73092418508	-0.0092795967675229\\
7861.31119521426	-0.0102210139008855\\
7878.89146624344	-0.0111000827878135\\
7902.33182761568	-0.0121464344563453\\
7925.77218898791	-0.0130173601795693\\
7949.21255036015	-0.0136843239245271\\
7972.65291173239	-0.014126073263905\\
7996.09327310463	-0.0143292441174659\\
8019.53363447687	-0.014288693288608\\
8042.9739958491	-0.0140075516483069\\
8066.41435722134	-0.013497000204998\\
8089.85471859358	-0.0127757804825216\\
8113.29507996582	-0.0118694592992409\\
8136.73544133805	-0.0108094759260288\\
8160.17580271029	-0.0096320064451336\\
8177.75607373947	-0.00869555529431648\\
8195.33634476865	-0.00773282548940009\\
8212.91661579783	-0.00676169411315181\\
8230.496886827	-0.00580001861840479\\
8248.07715785618	-0.00486529950990427\\
8259.7973385423	-0.00426547128903829\\
8271.51751922842	-0.00368986735086344\\
8283.23769991454	-0.00314300127055771\\
8294.95788060066	-0.00262911461177559\\
8306.67806128678	-0.00215214304136545\\
8318.3982419729	-0.0017156852439257\\
8324.25833231596	-0.00151367448811043\\
8330.11842265902	-0.00132297487807245\\
8335.97851300208	-0.00114392951863701\\
8341.83860334514	-0.000976855789136341\\
8347.69869368819	-0.000822044759626161\\
8353.55878403125	-0.000679760664552227\\
8359.41887437431	-0.00055024043461665\\
8365.27896471737	-0.000433693287773499\\
8371.13905506043	-0.000330300379891579\\
8376.99914540349	-0.000240214515715671\\
8382.85923574655	-0.000163559920488515\\
8388.71932608961	-0.000100432072515783\\
8394.57941643267	-5.08975968296321e-05\\
8400.43950677573	-1.49942199766626e-05\\
8406.29959711879	7.26921418331051e-06\\
8412.15968746185	1.59126679020703e-05\\
8418.01977780491	1.09847725478301e-05\\
8423.87986814797	-7.43736831905056e-06\\
8429.73995849103	-3.92484351895344e-05\\
8435.60004883409	-8.4315207436601e-05\\
8441.46013917715	-0.000142476935878959\\
8447.32022952021	-0.000213545772104161\\
8453.18031986327	-0.000297307253461387\\
8459.04041020633	-0.000393520842531137\\
8464.90050054939	-0.0005019205197564\\
8470.76059089244	-0.000622215427785242\\
8476.6206812355	-0.000754090566093032\\
8482.48077157856	-0.000897207534121875\\
8488.34086192162	-0.00105120532135851\\
8494.20095226468	-0.00121570114241605\\
8500.06104260774	-0.00139029131535057\\
8511.78122329386	-0.00176804106188055\\
8523.50140397998	-0.00218084307301458\\
8535.2215846661	-0.00262481423580925\\
8546.94176535222	-0.00309583426461801\\
8558.66194603833	-0.00358958192631633\\
8570.38212672445	-0.00410157272344232\\
8582.10230741057	-0.00462719770940002\\
8599.68257843975	-0.00543091830670217\\
8617.26284946893	-0.0062387563887071\\
8634.84312049811	-0.0070348019703062\\
8652.42339152729	-0.00780352308687485\\
8670.00366255646	-0.00853005589558444\\
8693.4440239287	-0.00940924703232076\\
8716.88438530094	-0.0101592461965354\\
8740.32474667318	-0.0107551915574221\\
8763.76510804542	-0.0111777149487521\\
8787.20546941765	-0.0114135236097825\\
8810.64583078989	-0.0114557714405134\\
8834.08619216213	-0.0113042099318977\\
8857.52655353437	-0.0109651164000528\\
8880.96691490661	-0.0104510046110101\\
8904.40727627884	-0.00978013008476346\\
8927.84763765108	-0.00897580907050897\\
8951.28799902332	-0.00806557617746692\\
8968.8682700525	-0.00733182030297483\\
8986.44854108167	-0.00656959696098284\\
9004.02881211085	-0.00579316675727577\\
9021.60908314003	-0.00501693804219783\\
9039.18935416921	-0.00425519297103749\\
9050.90953485533	-0.00376229125416531\\
9062.62971554145	-0.00328597535268875\\
9074.34989622757	-0.00283004623514711\\
9086.07007691369	-0.00239811150803011\\
9097.79025759981	-0.00199355611571542\\
9109.51043828592	-0.00161951508258513\\
9121.23061897204	-0.00127884851523458\\
9127.0907093151	-0.00112184366623683\\
9132.95079965816	-0.000974119061537158\\
9138.81089000122	-0.000835945287546409\\
9144.67098034428	-0.00070757200220806\\
9150.53107068734	-0.000589227467265248\\
9156.3911610304	-0.000481118127020635\\
9162.25125137346	-0.000383428234241947\\
9168.11134171652	-0.00029631952393277\\
9173.97143205958	-0.000219930935522611\\
9179.83152240264	-0.000154378383879877\\
9185.6916127457	-9.97545795444341e-05\\
9191.55170308876	-5.61288983802229e-05\\
9197.41179343182	-2.35473007952853e-05\\
9203.27188377488	-2.03230056979739e-06\\
9209.13197411794	8.41701679615916e-06\\
9214.992064461	7.82492638399595e-06\\
9220.85215480405	-3.76105290718477e-06\\
9226.71224514711	-2.62703238093769e-05\\
9232.57233549017	-5.96094243329286e-05\\
9238.43242583323	-0.000103662289029873\\
9244.29251617629	-0.000158290557183506\\
9250.15260651935	-0.000223333927157371\\
9256.01269686241	-0.000298610556014676\\
9261.87278720547	-0.000383917503452488\\
9267.73287754853	-0.000479031218977657\\
9273.59296789159	-0.000583708071202371\\
9279.45305823465	-0.000697684918003724\\
9285.31314857771	-0.000820679716238142\\
9291.17323892077	-0.000952392169648613\\
9297.03332926383	-0.00109250441344431\\
9302.89341960689	-0.00124068173402229\\
9314.613600293	-0.00155981305872976\\
9326.33378097912	-0.0019068008664674\\
9338.05396166524	-0.00227844166743545\\
9349.77414235136	-0.00267134196348359\\
9361.49432303748	-0.00308194789715026\\
9373.2145037236	-0.00350657608903786\\
9390.79477475278	-0.0041615200974647\\
9408.37504578196	-0.00482648881086774\\
9425.95531681113	-0.00548839807309779\\
9443.53558784031	-0.00613433036449451\\
9461.11585886949	-0.00675177795219038\\
9478.69612989867	-0.00732887633375124\\
9502.13649127091	-0.00801670951534958\\
9525.57685264315	-0.00859032199472302\\
9549.01721401538	-0.00903075302024208\\
9572.45757538762	-0.00932372693291143\\
9595.89793675986	-0.00946008298952435\\
9619.3382981321	-0.00943602918368367\\
9642.77865950433	-0.00925321340688587\\
9666.21902087657	-0.00891861146941937\\
9689.65938224881	-0.00844423763092448\\
9713.09974362105	-0.00784668916752573\\
9736.54010499328	-0.00714654193616621\\
9759.98046636552	-0.00636761870536307\\
9777.5607373947	-0.00574744695550532\\
9795.14100842388	-0.0051093595419891\\
9812.72127945306	-0.00446527880737508\\
9830.30155048224	-0.00382715294154637\\
9847.88182151142	-0.00320672950243876\\
9859.60200219753	-0.00280856542828907\\
9871.32218288365	-0.00242652378287964\\
9883.04236356977	-0.00206366072592535\\
9894.76254425589	-0.00172285567843108\\
9906.48272494201	-0.00140678791055024\\
9918.20290562813	-0.00111791493181928\\
9924.06299597119	-0.000984378947466937\\
9929.92308631425	-0.000858452856706486\\
9935.78317665731	-0.000740373083796323\\
9941.64326700037	-0.000630358900904697\\
9947.50335734343	-0.000528612007057587\\
9953.36344768649	-0.0004353161452437\\
9959.22353802955	-0.000350636758274933\\
9965.08362837261	-0.000274720684174514\\
9970.94371871567	-0.000207695891532479\\
9976.80380905872	-0.000149671255333666\\
9982.66389940178	-0.000100736373622287\\
9988.52398974484	-6.09614252983426e-05\\
9994.3840800879	-3.03970692294141e-05\\
10000.244170431	-9.07438479967228e-06\\
10006.104260774	2.99514607758894e-06\\
10011.9643511171	5.81961555461356e-06\\
10017.8244414601	-5.73374805836168e-07\\
10023.6845318032	-1.61368196951621e-05\\
10029.5446221463	-4.08044564115603e-05\\
10035.4047124893	-7.44909514947628e-05\\
10041.2648028324	-0.00011709212714911\\
10047.1248931754	-0.00016848522688702\\
10052.9849835185	-0.000228529219733049\\
10058.8450738616	-0.000297065142248737\\
10064.7051642046	-0.000373916477591258\\
10070.5652545477	-0.000458889570673991\\
10076.4253448907	-0.000551774078525254\\
10082.2854352338	-0.000652343454739857\\
10088.1455255769	-0.000760355466990062\\
10094.0056159199	-0.000875552746390671\\
10099.865706263	-0.000997663367469788\\
10111.5858869491	-0.00126146783372883\\
10123.3060676352	-0.00154932617225742\\
10135.0262483213	-0.00185860456324919\\
10146.7464290074	-0.00218650118910178\\
10158.4666096936	-0.00253007066847834\\
10170.1867903797	-0.00288624958414371\\
10187.7670614089	-0.00343724056602401\\
10205.347332438	-0.00399859697078768\\
10222.9276034672	-0.00455930909219037\\
10240.5078744964	-0.00510845869602223\\
10258.0881455256	-0.00563542519493778\\
10275.6684165548	-0.00613008474583198\\
10299.108777927	-0.00672318222266476\\
10322.5491392992	-0.00722223915348606\\
10345.9895006715	-0.00761068283890666\\
10369.4298620437	-0.00787581883625711\\
10392.8702234159	-0.00800922188084799\\
10416.3105847882	-0.00800697993267605\\
10439.7509461604	-0.00786978450421604\\
10463.1913075327	-0.00760286563735366\\
10486.6316689049	-0.0072157751154065\\
10510.0720302771	-0.00672202654818299\\
10533.5123916494	-0.00613860568526369\\
10556.9527530216	-0.00548536854095951\\
10574.5330240508	-0.00496281748628308\\
10592.11329508	-0.00442321470788328\\
10609.6935661091	-0.00387667524878795\\
10627.2738371383	-0.00333338047999822\\
10644.8541081675	-0.00280338471955571\\
10656.5742888536	-0.00246226982840681\\
10668.2944695397	-0.00213417286747633\\
10680.0146502259	-0.00182174124790012\\
10691.734830912	-0.00152747882211605\\
10703.4550115981	-0.00125372537631666\\
10715.1751922842	-0.00100263759482658\\
10721.0352826273	-0.000886212905200961\\
10726.8953729703	-0.000776171650899172\\
10732.7554633134	-0.000672724456844635\\
10738.6155536565	-0.000576067563078175\\
10744.4756439995	-0.000486382440981392\\
10750.3357343426	-0.000403835441610912\\
10756.1958246856	-0.000328577476596296\\
10762.0559150287	-0.000260743732411699\\
10767.9160053717	-0.000200453418395766\\
10773.7760957148	-0.000147809549082245\\
10779.6361860579	-0.000102898761174201\\
10785.4962764009	-6.57911654813149e-05\\
10791.356366744	-3.65402340709376e-05\\
10797.216457087	-1.51827227904997e-05\\
10803.0765474301	-1.73862924079821e-06\\
10808.9366377732	3.78881377659286e-06\\
10814.7967281162	1.41310940428786e-06\\
10820.6568184593	-8.83556739862741e-06\\
10826.5169088023	-2.69104690456881e-05\\
10832.3769991454	-5.27484092917896e-05\\
10838.2370894885	-8.62699316219065e-05\\
10844.0971798315	-0.000127379511635701\\
10849.9572701746	-0.000175965792927513\\
10855.8173605176	-0.000231901855888585\\
10861.6774508607	-0.000295045518808255\\
10867.5375412038	-0.000365239670544749\\
10873.3976315468	-0.000442312634040603\\
10879.2577218899	-0.000526078559795566\\
10885.1178122329	-0.000616337848460132\\
10890.977902576	-0.000712877601540279\\
10896.8379929191	-0.000815472099262405\\
10908.5581736052	-0.00103786139158772\\
10920.2783542913	-0.00128146330477403\\
10931.9985349774	-0.00154406501829236\\
10943.7187156635	-0.00182330269912859\\
10955.4388963497	-0.00211668208161304\\
10967.1590770358	-0.00242160006130094\\
10978.8792577219	-0.00273536711684524\\
10996.4595287511	-0.00321657624242851\\
11014.0397997802	-0.00370205835039019\\
11031.6200708094	-0.00418235783207747\\
11049.2003418386	-0.00464817693161243\\
11066.7806128678	-0.00509055100250311\\
11090.22097424	-0.00562937647035128\\
11113.6613356123	-0.00609337664524348\\
11137.1016969845	-0.00646704616429113\\
11160.5420583567	-0.00673804868416983\\
11183.982419729	-0.00689760606791184\\
11207.4227811012	-0.00694076489891446\\
11230.8631424734	-0.00686653236261184\\
11254.3035038457	-0.00667787797909878\\
11277.7438652179	-0.00638160219779822\\
11301.1842265902	-0.00598807732340355\\
11324.6245879624	-0.00551087048014656\\
11348.0649493346	-0.00496626220167837\\
11365.6452203638	-0.00452459467654208\\
11383.225491393	-0.0040637079253725\\
11400.8057624222	-0.00359227986894734\\
11418.3860334513	-0.00311913899205397\\
11435.9663044805	-0.00265309607988796\\
11453.5465755097	-0.00220277726615378\\
11465.2667561958	-0.00191540775218794\\
11476.9869368819	-0.00164105178308009\\
11488.7071175681	-0.00138192579859688\\
11500.4272982542	-0.00114010927397309\\
11512.1474789403	-0.000917527687418584\\
11523.8676596264	-0.000715936851649411\\
11529.7277499695	-0.000623511742444622\\
11535.5878403125	-0.000536908736441398\\
11541.4479306556	-0.000456293541417036\\
11547.3080209987	-0.000381818895900072\\
11553.1681113417	-0.000313624270939049\\
11559.0282016848	-0.000251835600389362\\
11564.8882920278	-0.000196565040367412\\
11570.7483823709	-0.000147910758218454\\
11576.608472714	-0.000105956751441114\\
11582.468563057	-7.07726968791103e-05\\
11588.3286534001	-4.24138304485396e-05\\
11594.1887437431	-2.09208575862028e-05\\
11600.0488340862	-6.31989455759075e-06\\
11605.9089244293	1.37755931010916e-06\\
11611.7690147723	2.17461847321223e-06\\
11617.6291051154	-3.91102284497095e-06\\
11623.4891954584	-1.68471553991068e-05\\
11629.3492858015	-3.65871502018753e-05\\
11635.2093761446	-6.3070084082781e-05\\
11641.0694664876	-9.62208952007329e-05\\
11646.9295568307	-0.00013595056812842\\
11652.7896471737	-0.000182156348062746\\
11658.6497375168	-0.000234721983653205\\
11664.5098278598	-0.000293517997887542\\
11670.3699182029	-0.000358401986397821\\
11676.230008546	-0.00042921894251929\\
11682.090098889	-0.000505801608369831\\
11687.9501892321	-0.000587970851146087\\
11693.8102795751	-0.000675536063791807\\
11699.6703699182	-0.000768295589181372\\
11711.3905506043	-0.000968538401198433\\
11723.1107312904	-0.00118688253509249\\
11734.8309119766	-0.00142136517875597\\
11746.5510926627	-0.00166989457027569\\
11758.2712733488	-0.0019302681327759\\
11769.9914540349	-0.00220019147766297\\
11787.5717250641	-0.00261779089249372\\
11805.1519960933	-0.00304335659034545\\
11822.7322671225	-0.00346859907837759\\
11840.3125381516	-0.00388528109790063\\
11857.8928091808	-0.00428537354115635\\
11875.47308021	-0.0046612065759582\\
11898.9134415822	-0.00511228213182316\\
11922.3538029545	-0.00549238781977268\\
11945.7941643267	-0.00578885046268634\\
11969.2345256989	-0.0059919061371329\\
11992.6748870712	-0.00609500790493707\\
12016.1152484434	-0.00609502320517862\\
12039.5556098157	-0.00599231505753933\\
12062.9959711879	-0.00579070517139075\\
12086.4363325601	-0.00549732103098363\\
12109.8766939324	-0.0051223329093238\\
12133.3170553046	-0.00467859041602923\\
12156.7574166768	-0.00418117149073315\\
12174.337687706	-0.00378293033943343\\
12191.9179587352	-0.00337143889706234\\
12209.4982297644	-0.00295443582151013\\
12227.0785007936	-0.00253972531043569\\
12244.6587718227	-0.00213502873525683\\
12256.3789525089	-0.00187450803930013\\
12268.099133195	-0.0016239098969604\\
12279.8193138811	-0.00138527762529527\\
12291.5394945672	-0.00116054662920597\\
12303.2596752533	-0.000951528422873768\\
12314.9798559394	-0.000759895750868374\\
12320.8399462825	-0.000671080390199808\\
12326.7000366256	-0.000587168930638313\\
12332.5601269686	-0.000508325572768237\\
12338.4202173117	-0.000434703527455033\\
12344.2803076547	-0.000366444710434796\\
12350.1403979978	-0.000303679461249289\\
12356.0004883409	-0.000246526287144587\\
12361.8605786839	-0.000195091632371791\\
12367.720669027	-0.00014946967337851\\
12373.58075937	-0.000109742140246979\\
12379.4408497131	-7.59781647374057e-05\\
12385.3009400562	-4.82341551947391e-05\\
12391.1610303992	-2.65536985479246e-05\\
12397.0211207423	-1.09674895595291e-05\\
12402.8812110853	-1.49328743399849e-06\\
12408.7413014284	1.86410017836111e-06\\
12414.6013917715	-8.87194216363841e-07\\
12420.4614821145	-9.72613667174347e-06\\
12426.3215724576	-2.4618857708653e-05\\
12432.1816628006	-4.55187452331639e-05\\
12438.0417531437	-7.23665642196277e-05\\
12443.9018434868	-0.000105090602865757\\
12449.7619338298	-0.000143606844864621\\
12455.6220241729	-0.000187819167386013\\
12461.4821145159	-0.000237619564309568\\
12467.342204859	-0.000292888394196339\\
12473.2022952021	-0.000353494652429067\\
12479.0623855451	-0.000419296266918004\\
12484.9224758882	-0.00049014041669363\\
12490.7825662312	-0.000565863872680999\\
12496.6426565743	-0.000646293359927287\\
12508.3628372604	-0.000820529415734817\\
12520.0830179465	-0.00101127650213027\\
12531.8031986326	-0.0012168282831699\\
12543.5233793188	-0.00143535901862137\\
12555.2435600049	-0.00166493937274596\\
12566.963740691	-0.00190355304595082\\
12584.5440117202	-0.00227383334547552\\
12602.1242827494	-0.00265249363169663\\
12619.7045537785	-0.00303217118216235\\
12637.2848248077	-0.00340552026183173\\
12654.8650958369	-0.00376535139099728\\
12672.4453668661	-0.00410476697843897\\
12695.8857282383	-0.00451444596939408\\
12719.3260896105	-0.00486257911584817\\
12742.7664509828	-0.00513752832805438\\
12766.206812355	-0.0053301961663468\\
12789.6471737273	-0.0054343159556759\\
12813.0875350995	-0.00544664484648214\\
12836.5278964717	-0.0053670539438863\\
12859.968257844	-0.00519851312115898\\
12883.4086192162	-0.00494697168780536\\
12906.8489805884	-0.00462113957348873\\
12930.2893419607	-0.00423217700209131\\
12953.7297033329	-0.00379330365670004\\
12971.3099743621	-0.00344026461380296\\
12988.8902453913	-0.0030741463658034\\
13006.4705164205	-0.00270184819879727\\
13024.0507874496	-0.00233035573420049\\
13041.6310584788	-0.0019666079865307\\
13053.3512391649	-0.00173176206430656\\
13065.0714198511	-0.00150530228981414\\
13076.7916005372	-0.00128908492892611\\
13088.5117812233	-0.00108487418502644\\
13100.2319619094	-0.000894327574196526\\
13111.9521425955	-0.000718982245030837\\
13123.6723232817	-0.000560242355756582\\
13129.5324136247	-0.000487499039817409\\
13135.3925039678	-0.000419367611854463\\
13141.2525943108	-0.00035598128924929\\
13147.1126846539	-0.000297463057510905\\
13152.9727749969	-0.000243925425517103\\
13158.83286534	-0.000195470203237305\\
13164.6929556831	-0.000152188302400557\\
13170.5530460261	-0.00011415956049917\\
13176.4131363692	-8.14525884456714e-05\\
13182.2732267122	-5.41246421823077e-05\\
13188.1333170553	-3.22215184816521e-05\\
13193.9934073984	-1.57774751126702e-05\\
13199.8534977414	-4.8151755103121e-06\\
13205.7135880845	6.54341970378372e-07\\
13211.5736784275	6.3167019038465e-07\\
13217.4337687706	-4.87098522725504e-06\\
13223.2938591137	-1.58298488625258e-05\\
13229.1539494567	-3.22096428508621e-05\\
13235.0140397998	-5.39636784412149e-05\\
13240.8741301428	-8.10339715372089e-05\\
13246.7342204859	-0.000113351381929817\\
13252.594310829	-0.000150835775911191\\
13258.454401172	-0.000193396211866665\\
13264.3144915151	-0.000240931148450913\\
13270.1745818581	-0.000293328674839658\\
13276.0346722012	-0.000350466762573949\\
13281.8947625443	-0.000412213538406629\\
13287.7548528873	-0.000478427577543718\\
13293.6149432304	-0.000548958216667969\\
13305.3351239165	-0.000702322459742393\\
13317.0553046026	-0.000870929261528749\\
13328.7754852887	-0.00105327689021096\\
13340.4956659748	-0.0012477519367211\\
13352.215846661	-0.00145264326805179\\
13363.9360273471	-0.00166615676518477\\
13375.6562080332	-0.00188643071881885\\
13393.2364790624	-0.00222532187206752\\
13410.8167500916	-0.00256852132462072\\
13428.3970211207	-0.0029093868254581\\
13445.9772921499	-0.00324134948453528\\
13463.5575631791	-0.0035580381126893\\
13486.9979245513	-0.00394612734397704\\
13510.4382859236	-0.00428324324341816\\
13533.8786472958	-0.00455807132229245\\
13557.3190086681	-0.00476146104900936\\
13580.7593700403	-0.00488672176510969\\
13604.1997314125	-0.00492983415279872\\
13627.6400927848	-0.00488957061478979\\
13651.080454157	-0.00476752102761162\\
13674.5208155292	-0.00456802352855954\\
13697.9611769015	-0.00429800319315724\\
13721.4015382737	-0.00396672453832547\\
13744.841899646	-0.00358546665249678\\
13762.4221706751	-0.00327443739756012\\
13780.0024417043	-0.00294842926925367\\
13797.5827127335	-0.00261360398102513\\
13815.1629837627	-0.00227626601918929\\
13832.7432547918	-0.00194274240345282\\
13850.323525821	-0.00161926250841621\\
13862.0437065071	-0.00141217123863335\\
13873.7638871933	-0.00121392866245647\\
13885.4840678794	-0.00102616068145417\\
13897.2042485655	-0.000850400377821322\\
13908.9244292516	-0.000688075310060081\\
13920.6446099377	-0.000540495733723048\\
13926.5047002808	-0.00047261021497196\\
13932.3647906239	-0.000408843844137917\\
13938.2248809669	-0.000349322803258595\\
13944.08497131	-0.000294164129152046\\
13949.945061653	-0.000243475478057306\\
13955.8051519961	-0.000197354910416347\\
13961.6652423392	-0.000155890696233735\\
13967.5253326822	-0.000119161141364189\\
13973.3854230253	-8.72344351123164e-05\\
13979.2455133683	-6.01685194012163e-05\\
13985.1056037114	-3.80109797838483e-05\\
13990.9656940544	-2.07989584786331e-05\\
13996.8257843975	-8.5590895997881e-06\\
14002.6858747406	-1.30745667827178e-06\\
14008.5459650836	9.50427459075276e-07\\
14014.4060554267	-1.78038155813986e-06\\
14020.2661457697	-9.48428422582442e-06\\
14026.1262361128	-2.21351840003748e-05\\
14031.9863264559	-3.96965562601339e-05\\
14037.8464167989	-6.21215392153655e-05\\
14043.706507142	-8.93530465502477e-05\\
14049.566597485	-0.000121323901525574\\
14055.4266878281	-0.000157956992237147\\
14061.2867781712	-0.000199165447665733\\
14067.1468685142	-0.000244852834125261\\
14073.0069588573	-0.00029491337167422\\
14078.8670492003	-0.00034923216998541\\
14084.7271395434	-0.000407685483172795\\
14090.5872298865	-0.000470140983008642\\
14096.4473202295	-0.000536458049918619\\
14108.1675009156	-0.000680074815249637\\
14119.8876816018	-0.000837257111180151\\
14131.6078622879	-0.0010066149790278\\
14143.328042974	-0.00118665920108699\\
14155.0482236601	-0.00137581413935497\\
14166.7684043462	-0.00157243127103134\\
14184.3486753754	-0.0018776013910078\\
14201.9289464046	-0.00218977685976661\\
14219.5092174338	-0.00250291496146803\\
14237.0894884629	-0.00281097820540363\\
14254.6697594921	-0.00310804890926719\\
14272.2500305213	-0.0033884410718132\\
14295.6903918935	-0.00372717211121178\\
14319.1307532658	-0.00401537813165297\\
14342.571114638	-0.00424340183347805\\
14366.0114760103	-0.00440366374096463\\
14389.4518373825	-0.00449090710846186\\
14412.8921987547	-0.00450236340228994\\
14436.332560127	-0.00443783318470261\\
14459.7729214992	-0.00429968009947604\\
14483.2132828714	-0.00409273859038365\\
14506.6536442437	-0.00382413888456508\\
14530.0940056159	-0.0035030555364934\\
14553.5343669882	-0.00314038836962424\\
14571.1146380173	-0.00284842832876437\\
14588.6949090465	-0.00254548029175913\\
14606.2751800757	-0.00223726496221597\\
14623.8554511049	-0.00192958274735686\\
14641.4357221341	-0.00162820335005547\\
14653.1559028202	-0.00143357435883956\\
14664.8760835063	-0.0012458651337258\\
14676.5962641924	-0.00106662478674308\\
14688.3164448785	-0.000897327202285143\\
14700.0366255646	-0.000739358761476168\\
14711.7568062508	-0.000594006834663746\\
14717.6168965938	-0.00052643465546198\\
14723.4769869369	-0.000462449137631033\\
14729.3370772799	-0.000402178771716434\\
14735.197167623	-0.00034574403913085\\
14741.0572579661	-0.000293257166773564\\
14746.9173483091	-0.000244821899432677\\
14752.7774386522	-0.000200533290366235\\
14758.6375289952	-0.000160477510506018\\
14764.4976193383	-0.000124731676664487\\
14770.3577096814	-9.33636990873341e-05\\
14776.2178000244	-6.64321486381091e-05\\
14782.0778903675	-4.39861438923031e-05\\
14787.9379807105	-2.60652583372505e-05\\
14793.7980710536	-1.26994478629447e-05\\
14799.6581613967	-3.90899866785218e-06\\
14805.5182517397	2.95504333086003e-07\\
14811.3783420828	-8.6811443922328e-08\\
14817.2384324258	-5.04711574925861e-06\\
14823.0985227689	-1.45669054970921e-05\\
14828.958613112	-2.86180551753425e-05\\
14834.818703455	-4.7162887539361e-05\\
14840.6787937981	-7.01542644141694e-05\\
14846.5388841411	-9.75356973902043e-05\\
14852.3989744842	-0.000129241478160138\\
14858.2590648272	-0.000165196828188351\\
14864.1191551703	-0.000205318067391454\\
14869.9792455134	-0.000249512801427311\\
14875.8393358564	-0.00029768012719494\\
14881.6994261995	-0.000349710856090986\\
14887.5595165425	-0.000405487754519023\\
14893.4196068856	-0.000464885801107435\\
14905.1397875717	-0.000594007970482249\\
14916.8599682578	-0.000735932211527269\\
14928.580148944	-0.000889407827800184\\
14940.3003296301	-0.00105308939068388\\
14952.0205103162	-0.00122554832635953\\
14963.7406910023	-0.00140528517846411\\
14975.4608716884	-0.00159074244143019\\
14993.0411427176	-0.00187614027613619\\
15010.6214137468	-0.00216527240276336\\
15028.201684776	-0.00245256272851965\\
15045.7819558052	-0.00273248936690219\\
15063.3622268343	-0.00299968929659867\\
15086.8025882066	-0.00332738854628988\\
15110.2429495788	-0.00361236235425926\\
15133.683310951	-0.00384503190591027\\
15157.1236723233	-0.00401762519138437\\
15180.5640336955	-0.0041244308448573\\
15204.0043950678	-0.00416198145516016\\
15227.44475644	-0.00412916048897852\\
15250.8851178122	-0.00402722956387104\\
15274.3254791845	-0.00385977550368891\\
15297.7658405567	-0.00363257930993113\\
15321.2062019289	-0.00335341180118523\\
15344.6465633012	-0.00303176311997486\\
15362.2268343304	-0.00276915263604038\\
15379.8071053595	-0.00249373244525841\\
15397.3873763887	-0.00221071580829229\\
15414.9676474179	-0.00192544384110782\\
15432.5479184471	-0.00164328368143192\\
15450.1281894763	-0.00136952648722847\\
15461.8483701624	-0.00119422610243507\\
15473.5685508485	-0.00102639131353052\\
15485.2887315346	-0.00086740798862672\\
15497.0089122207	-0.000718584280846323\\
15508.7290929068	-0.000581139735917738\\
15520.449273593	-0.000456195171261834\\
15526.309363936	-0.000398731341943496\\
15532.1694542791	-0.000344763410350772\\
15538.0295446221	-0.000294399737903101\\
15543.8896349652	-0.000247740948886146\\
15549.7497253083	-0.000204879725700828\\
15555.6098156513	-0.000165900621115956\\
15561.4699059944	-0.000130879887902337\\
15567.3299963374	-9.98853261971183e-05\\
15573.1900866805	-7.2976148895099e-05\\
15579.0501770236	-5.0202865334378e-05\\
15584.9102673666	-3.16071835082891e-05\\
15590.7703577097	-1.72219309810085e-05\\
15596.6304480527	-7.07099466777598e-06\\
15602.4905383958	-1.16927957373461e-06\\
15608.3506287389	4.773134301401e-07\\
15614.2107190819	-2.1281092255186e-06\\
15620.070809425	-8.97344969084332e-06\\
15625.930899768	-2.00376535697323e-05\\
15631.7909901111	-3.52907636955655e-05\\
15637.6510804542	-5.46939926726159e-05\\
15643.5111707972	-7.81998140127464e-05\\
15649.3712611403	-0.000105752071654079\\
15655.2313514833	-0.000137286107612723\\
15661.0914418264	-0.00017272890747811\\
15666.9515321695	-0.000211999263428327\\
15672.8116225125	-0.000255007954396796\\
15678.6717128556	-0.000301657943000936\\
15684.5318031986	-0.000351844588788032\\
15690.3918935417	-0.000405455877350871\\
15696.2519838848	-0.000462372664789287\\
15707.9721645709	-0.000585612083186852\\
15719.692345257	-0.000720477307532184\\
15731.4125259431	-0.000865786753700883\\
15743.1327066292	-0.00102027297176538\\
15754.8528873153	-0.00118259353263005\\
15766.5730680015	-0.00135134252604895\\
15784.1533390306	-0.00161332028411956\\
15801.7336100598	-0.00188140378297334\\
15819.313881089	-0.00215042256932537\\
15836.8941521182	-0.00241520420116699\\
15854.4744231474	-0.00267067248342958\\
15872.0546941765	-0.0029119435820022\\
15895.4950555488	-0.00320365634488596\\
15918.935416921	-0.00345215671861917\\
15942.3757782932	-0.00364910212262316\\
15965.8161396655	-0.0037879219159577\\
15989.2565010377	-0.00386403189965901\\
16012.69686241	-0.00387498105405611\\
16036.1372237822	-0.00382052583439182\\
16059.5775851544	-0.00370262980444198\\
16083.0179465267	-0.00352538891069844\\
16106.4583078989	-0.00329488520091443\\
16129.8986692712	-0.00301897418900374\\
16153.3390306434	-0.002707013282268\\
16170.9193016726	-0.00245568979215539\\
16188.4995727017	-0.00219476553655616\\
16206.0798437309	-0.00192917612599855\\
16223.6601147601	-0.00166393207529827\\
16241.2403857893	-0.00140402342593567\\
16252.9605664754	-0.00123612998662608\\
16264.6807471615	-0.00107417557523395\\
16276.4009278476	-0.000919504264411494\\
16288.1211085338	-0.000773396076263431\\
16299.8412892199	-0.000637056281542482\\
16311.561469906	-0.000511605350409784\\
16317.4215602491	-0.00045328785944793\\
16323.2816505921	-0.000398069638266874\\
16329.1417409352	-0.000346062880858526\\
16335.0018312782	-0.000297372884124726\\
16340.8619216213	-0.000252097831477033\\
16346.7220119643	-0.000210328591719765\\
16352.5821023074	-0.00017214853354908\\
16358.4421926505	-0.000137633356111667\\
16364.3022829935	-0.000106850935909349\\
16370.1623733366	-7.98611904064394e-05\\
16376.0224636796	-5.6715958578948e-05\\
16381.8825540227	-3.74588986587459e-05\\
16387.7426443658	-2.21254032706356e-05\\
16393.6027347088	-1.07425321233418e-05\\
16399.4628250519	-3.32896238068003e-06\\
16405.3229153949	1.05043195003138e-07\\
16411.183005738	-4.42349719875634e-07\\
16417.0430960811	-4.96455080859517e-06\\
16422.9031864241	-1.344656672264e-05\\
16428.7632767672	-2.58650404379933e-05\\
16434.6233671102	-4.21883082618168e-05\\
16440.4834574533	-6.2376474350329e-05\\
16446.3435477964	-8.63815025671135e-05\\
16452.2036381394	-0.000114147325468672\\
16458.0637284825	-0.000145609970177582\\
16463.9238188255	-0.000180697700851392\\
16469.7839091686	-0.000219331177446408\\
16475.6439995117	-0.000261423630404602\\
16481.5040898547	-0.000306881050907357\\
16487.3641801978	-0.000355602396256526\\
16493.2242705408	-0.000407479809935973\\
16504.944451227	-0.000520238766584936\\
16516.6646319131	-0.00064416803343282\\
16528.3848125992	-0.00077818507271979\\
16540.1049932853	-0.000921124013042951\\
16551.8251739714	-0.00107174565352221\\
16563.5453546575	-0.00122874806979584\\
16575.2655353437	-0.00139077773104547\\
16592.8458063728	-0.00164019175297522\\
16610.426077402	-0.00189296510718047\\
16628.0063484312	-0.00214423714665986\\
16645.5866194604	-0.00238918878177739\\
16663.1668904896	-0.00262313390868933\\
16686.6072518618	-0.00291026568369465\\
16710.047613234	-0.00316022850187203\\
16733.4879746063	-0.00336460994172886\\
16756.9283359785	-0.00351656461909876\\
16780.3686973508	-0.00361103955292185\\
16803.809058723	-0.00364493834982327\\
16827.2494200952	-0.00361721888623165\\
16850.6897814675	-0.00352892142605913\\
16874.1301428397	-0.00338312646340058\\
16897.5705042119	-0.0031848439498138\\
16921.0108655842	-0.00294083786981587\\
16944.4512269564	-0.00265939229299421\\
16962.0314979856	-0.00242942768344873\\
16979.6117690148	-0.00218810788894791\\
16997.1920400439	-0.00194000642106733\\
17014.7723110731	-0.00168981450016808\\
17032.3525821023	-0.00144225164530541\\
17049.9328531315	-0.00120197594788073\\
17061.6530338176	-0.00104807566845909\\
17073.3732145037	-0.000900703820765691\\
17085.0933951898	-0.000761084283144257\\
17096.813575876	-0.000630373411576671\\
17108.5337565621	-0.000509650373644626\\
17120.2539372482	-0.00039990815044365\\
17126.1140275913	-0.000349439536447095\\
17131.9741179343	-0.000302045281462762\\
17137.8342082774	-0.000257821725022549\\
17143.6942986204	-0.000216858420919931\\
17149.5543889635	-0.000179237953221132\\
17155.4144793066	-0.000145035767163624\\
17161.2745696496	-0.0001143200153026\\
17167.1346599927	-8.71514191847529e-05\\
17172.9947503357	-6.35831468621092e-05\\
17178.8548406788	-4.36607064672648e-05\\
17184.7149310219	-2.74218560787596e-05\\
17190.5750213649	-1.48965300430669e-05\\
17196.435111708	-6.10678190334502e-06\\
17202.295202051	-1.06674403792564e-06\\
17208.1552923941	2.17395917606153e-07\\
17214.0153827372	-2.25259817053585e-06\\
17219.8754730802	-8.467020996523e-06\\
17225.7355634233	-1.84082526626683e-05\\
17231.5956537663	-3.20508022419686e-05\\
17237.4557441094	-4.93613679520115e-05\\
17243.3158344524	-7.02989137991503e-05\\
17249.1759247955	-9.48147625138472e-05\\
17255.0360151386	-0.000122852704579042\\
17260.8961054816	-0.000154349123092558\\
17266.7561958247	-0.0001892331341979\\
17272.6162861677	-0.000227426742778135\\
17278.4763765108	-0.000268845013054363\\
17284.3364668539	-0.000313396253729843\\
17290.1965571969	-0.000360982217294307\\
17296.05664754	-0.000411498312991594\\
17307.7768282261	-0.000520872191859627\\
17319.4970089122	-0.000640563211729651\\
17331.2171895983	-0.000769530962286718\\
17342.9373702845	-0.000906658231333596\\
17354.6575509706	-0.00105076056980354\\
17366.3777316567	-0.00120059641004529\\
17383.9580026859	-0.0014332727330667\\
17401.5382737151	-0.00167145732269408\\
17419.1185447442	-0.00191057048899077\\
17436.6988157734	-0.00214602562042023\\
17454.2790868026	-0.0023733163380707\\
17471.8593578318	-0.00258810192568318\\
17495.299719204	-0.00284800142720994\\
17518.7400805762	-0.00306966032369933\\
17542.1804419485	-0.00324562526874034\\
17565.6208033207	-0.00337000730766418\\
17589.061164693	-0.00343867576068483\\
17612.5015260652	-0.00344939225448289\\
17635.9418874374	-0.00340188056494002\\
17659.3822488097	-0.00329783010838866\\
17682.8226101819	-0.00314083316079808\\
17706.2629715541	-0.00293625811604845\\
17729.7033329264	-0.00269106324199608\\
17753.1436942986	-0.00241355738528607\\
17770.7239653278	-0.00218983247968374\\
17788.304236357	-0.00195743755798805\\
17805.8845073862	-0.00172077508727904\\
17823.4647784153	-0.00148431933057469\\
17841.0450494445	-0.00125253118054514\\
17852.7652301306	-0.00110275988059628\\
17864.4854108167	-0.000958257168175217\\
17876.2055915029	-0.00082022808488528\\
17887.925772189	-0.000689821242698202\\
17899.6459528751	-0.000568119191981236\\
17911.3661335612	-0.00045612936152668\\
17917.2262239043	-0.000404068724687542\\
17923.0863142473	-0.000354775646231425\\
17928.9464045904	-0.00030835129829279\\
17934.8064949335	-0.000264890712639448\\
17940.6665852765	-0.000224482583897122\\
17946.5266756196	-0.000187209086336214\\
17952.3867659626	-0.00015314570456439\\
17958.2468563057	-0.000122361078494211\\
17964.1069466488	-9.49168628787064e-05\\
17969.9670369918	-7.08676017363276e-05\\
17975.8271273349	-5.0260617890078e-05\\
17981.6872176779	-3.31359178562001e-05\\
17987.547308021	-1.95261122754273e-05\\
17993.4073983641	-9.45635203853158e-06\\
17999.2674887071	-2.94428023271915e-06\\
nan	nan\\
18010.9876693932	-6.26058363293904e-07\\
18016.8477597363	-4.81744604422046e-06\\
18022.7078500794	-1.256161326216e-05\\
18028.5679404224	-2.38385014690801e-05\\
18034.4280307655	-3.86205909431193e-05\\
18040.2881211085	-5.68729641296078e-05\\
18046.1482114516	-7.85533845819969e-05\\
18052.0083017947	-0.000103612391332356\\
18057.8683921377	-0.000131993408469772\\
18063.7284824808	-0.000163632869696458\\
18069.5885728238	-0.000198460357582423\\
18075.4486631669	-0.000236398757207158\\
18081.3087535099	-0.000277364423876758\\
18087.168843853	-0.000321267364520118\\
18093.0289341961	-0.000368011432396361\\
18104.7491148822	-0.000469608852506454\\
18116.4692955683	-0.000581272632318132\\
18128.1894762544	-0.000702034797422572\\
18139.9096569405	-0.00083085174866964\\
18151.6298376267	-0.000966613188023733\\
18163.3500183128	-0.00110815159609184\\
18175.0701989989	-0.00125425218065833\\
18192.6504700281	-0.00147921149032229\\
18210.2307410573	-0.0017072890074432\\
18227.8110120864	-0.00193411008510074\\
18245.3912831156	-0.0021553327840013\\
18262.9715541448	-0.00236673031190815\\
18286.411915517	-0.0026263796659168\\
18309.8522768893	-0.00285265343988584\\
18333.2926382615	-0.00303792833347716\\
18356.7329996337	-0.00317598460315991\\
18380.173361006	-0.00326221221261659\\
18403.6137223782	-0.00329376217527827\\
18427.0540837505	-0.00326963812941574\\
18450.4944451227	-0.00319072521218933\\
18473.9348064949	-0.00305975541951707\\
18497.3751678672	-0.00288121078107255\\
18520.8155292394	-0.00266116777054495\\
18544.2558906116	-0.00240708834446833\\
18561.8361616408	-0.00219932674509548\\
18579.41643267	-0.00198118421492547\\
18596.9967036992	-0.00175679955082239\\
18614.5769747284	-0.00153042270630186\\
18632.1572457575	-0.00130633382351794\\
18649.7375167867	-0.00108876184037811\\
18661.4576974728	-0.000949365567354957\\
18673.177878159	-0.000815856617485192\\
18684.8980588451	-0.000689349321815146\\
18696.6182395312	-0.000570897465252642\\
18708.3384202173	-0.000461485448615252\\
18720.0586009034	-0.000362020047838695\\
18725.9186912465	-0.000316277591548625\\
18731.7787815895	-0.000273322839292687\\
18737.6388719326	-0.000233244052445928\\
18743.4989622757	-0.00019612335432281\\
18749.3590526187	-0.000162036560084215\\
18755.2191429618	-0.000131053020082639\\
18761.0792333048	-0.000103235476965918\\
18766.9393236479	-7.86399368576844e-05\\
18772.799413991	-5.73155548495061e-05\\
18778.659504334	-3.93045350541097e-05\\
18784.5195946771	-2.46420454223196e-05\\
18790.3796850201	-1.33561474888933e-05\\
18796.2397753632	-5.46774119205874e-06\\
18802.0998657063	-9.90524872404033e-07\\
18807.9599560493	6.90294731979261e-08\\
18813.8200463924	-2.28831436094514e-06\\
18819.6801367354	-8.054562652008e-06\\
18825.5402270785	-1.72145128995256e-05\\
18831.4003174216	-2.97457902099721e-05\\
18837.2604077646	-4.56188988320399e-05\\
18843.1204981077	-6.47972887199579e-05\\
18848.9805884507	-8.72374369815938e-05\\
18854.8406787938	-0.000112888944024474\\
18860.7007691369	-0.000141694644190633\\
18866.5608594799	-0.000173590730640519\\
18872.420949823	-0.000208506894207686\\
18878.281040166	-0.000246366475928381\\
18884.1411305091	-0.000287086632907327\\
18890.0012208522	-0.000330578517176891\\
18895.8613111952	-0.000376747467148975\\
18907.5814918813	-0.000476710083292226\\
18919.3016725674	-0.000586108944346025\\
18931.0218532536	-0.000703999708305402\\
18942.7420339397	-0.000829367304165914\\
18954.4622146258	-0.000961134596530641\\
18966.1823953119	-0.0010981715650256\\
18983.7626663411	-0.00131103170657061\\
19001.3429373703	-0.00152901138428191\\
19018.9232083995	-0.0017479307907947\\
19036.5034794286	-0.00196359942621619\\
19054.0837504578	-0.00217189575563247\\
19071.664021487	-0.00236884542259296\\
19095.1043828592	-0.00260735083508344\\
19118.5447442315	-0.00281099647918245\\
19141.9851056037	-0.00297292514471675\\
19165.4254669759	-0.00308770287047421\\
19188.8658283482	-0.0031514991531111\\
19212.3061897204	-0.00316221268620306\\
19235.7465510927	-0.00311953850888947\\
19259.1869124649	-0.00302497442691119\\
19282.6272738371	-0.00288176662031038\\
19306.0676352094	-0.00269479640136526\\
19329.5079965816	-0.00247041206654119\\
19352.9483579539	-0.00221621162719072\\
19370.528628983	-0.00201113198262937\\
19388.1089000122	-0.00179799330299313\\
19405.6891710414	-0.00158083860544934\\
19423.2694420706	-0.00136378112227274\\
19440.8497130997	-0.00115092601489146\\
19452.5698937859	-0.00101334785545283\\
19464.290074472	-0.000880580832073882\\
19476.0102551581	-0.000753736902092438\\
19487.7304358442	-0.00063387681071593\\
19499.4506165303	-0.000522001172095845\\
19511.1707972165	-0.000419042067812197\\
19517.0308875595	-0.000371176816142205\\
19522.8909779026	-0.000325855233192107\\
19528.7510682456	-0.000283171181858921\\
19534.6111585887	-0.000243212896917138\\
19540.4712489318	-0.000206062801140336\\
19546.3313392748	-0.000171797333797657\\
19552.1914296179	-0.000140486791887828\\
19558.0515199609	-0.000112195184460174\\
19563.911610304	-8.69801002897976e-05\\
19569.7717006471	-6.48925892133559e-05\\
19575.6317909901	-4.59770573428382e-05\\
19581.4918813332	-3.02711763954224e-05\\
19587.3519716762	-1.78058073018194e-05\\
19593.2120620193	-8.60493826722153e-06\\
19599.0721523623	-2.68563739496502e-06\\
19604.9322427054	-5.8019972881849e-08\\
19610.7923330485	-7.25230481056351e-07\\
19616.6524233915	-4.6834393550504e-06\\
19622.5125137346	-1.19218545007031e-05\\
19628.3726040776	-2.24227475303034e-05\\
19634.2326944207	-3.61614946600077e-05\\
19640.0927847638	-5.31066321693841e-05\\
19645.9528751068	-7.32199263042771e-05\\
19651.8129654499	-9.64564574642837e-05\\
19657.6730557929	-0.000122764718492462\\
19663.533146136	-0.000152086726849342\\
19669.3932364791	-0.000184358150436336\\
19675.2533268221	-0.000219508446775727\\
19681.1134171652	-0.000257461015274727\\
19686.9735075082	-0.000298133362220343\\
19692.8335978513	-0.00034143727816274\\
19704.5537785374	-0.000435559548498835\\
19716.2739592235	-0.000539015319244037\\
19727.9941399097	-0.000650913919879023\\
19739.7143205958	-0.000770294123527442\\
19751.4345012819	-0.000896132343363573\\
19763.154681968	-0.00102735134846444\\
19774.8748626541	-0.00116282942486231\\
19792.4551336833	-0.00137149493905905\\
19810.0354047125	-0.00158313571932502\\
19827.6156757417	-0.00179370166096843\\
19845.1959467708	-0.00199916885448461\\
19862.7762178	-0.00219561603434716\\
19886.2165791723	-0.00243707738301771\\
19909.6569405445	-0.00264771649768994\\
19933.0973019167	-0.0028204300975864\\
19956.537663289	-0.00294940824394311\\
19979.9780246612	-0.00303032804499619\\
20003.4183860334	-0.00306049685036157\\
20026.8587474057	-0.00303894019690302\\
20050.2991087779	-0.00296643163928709\\
20073.7394701502	-0.00284546356189473\\
20097.1798315224	-0.00268016006111007\\
20120.6201928946	-0.00247613493676936\\
20144.0605542669	-0.00224029967761424\\
20161.6408252961	-0.00204731205105907\\
20179.2210963252	-0.00184456966670262\\
20196.8013673544	-0.00163592296076743\\
20214.3816383836	-0.00142532986583139\\
20231.9619094128	-0.00121678043811988\\
20249.5421804419	-0.00101422096927117\\
20261.2623611281	-0.000884406333322104\\
20272.9825418142	-0.000760048725137236\\
20284.7027225003	-0.000642190833221445\\
20296.4229031864	-0.000531819517057388\\
20308.1430838725	-0.000429857506229384\\
20319.8632645587	-0.000337155648650876\\
20325.7233549017	-0.000294521742025921\\
20331.5834452448	-0.000254485772812664\\
20337.4435355878	-0.000217130812177925\\
20343.3036259309	-0.000182534224030026\\
20349.163716274	-0.000150767503627476\\
20355.023806617	-0.000121896128662621\\
20360.8838969601	-9.59794231308124e-05\\
20366.7439873031	-7.30704342765375e-05\\
20372.6040776462	-5.32158228612711e-05\\
20378.4641679893	-3.64557669931793e-05\\
20384.3242583323	-2.28238796987978e-05\\
20390.1843486754	-1.23471404209158e-05\\
20396.0444390184	-5.04584057474679e-06\\
20401.9045293615	-9.33543273517639e-07\\
20407.7646197046	-1.70573014601601e-08\\
20413.6247100476	-2.29642538632653e-06\\
20419.4848003907	-7.76492678911273e-06\\
20425.3448907337	-1.64090941985982e-05\\
20431.2049810768	-2.82087448914954e-05\\
20437.0650714198	-4.31370260855577e-05\\
20442.9251617629	-6.11604743852655e-05\\
20448.785252106	-8.22390891847683e-05\\
20454.645342449	-0.000106326419876124\\
20460.5054327921	-0.000133369666660804\\
20466.3655231351	-0.000163309794756426\\
20472.2256134782	-0.00019608166175015\\
20478.0857038213	-0.000231614157814808\\
20483.9457941643	-0.000269830358505824\\
20489.8058845074	-0.000310647689780881\\
20495.6659748504	-0.000353978104927964\\
20507.3861555366	-0.000447799778269396\\
20519.1063362227	-0.000550488985691926\\
20530.8265169088	-0.000661164947323027\\
20542.5466975949	-0.000778880002983106\\
20554.266878281	-0.000902627678661026\\
20565.9870589672	-0.00103135124413232\\
20583.5673299963	-0.0012313554878744\\
20601.1476010255	-0.0014362469475744\\
20618.7278720547	-0.0016421065777656\\
20636.3081430839	-0.00184500140783182\\
20653.8884141131	-0.0020410593251268\\
20671.4686851422	-0.00222654261847336\\
20694.9090465145	-0.00245133686907901\\
20718.3494078867	-0.00264349171910218\\
20741.7897692589	-0.00279652872353766\\
20765.2301306312	-0.00290529968183118\\
20788.6704920034	-0.00296615847732917\\
20812.1108533757	-0.00297708199229692\\
20835.5512147479	-0.00293773609248631\\
20858.9915761201	-0.00284948453350451\\
20882.4319374924	-0.00271534056849344\\
20905.8722988646	-0.00253986296944102\\
20929.3126602368	-0.00232900004545486\\
20952.7530216091	-0.00208988699101465\\
20970.3332926383	-0.00189684553715358\\
20987.9135636674	-0.00169611399452162\\
21005.4938346966	-0.0014915046565848\\
21023.0741057258	-0.0012868998843836\\
21040.654376755	-0.00108617822644426\\
21052.3745574411	-0.000956403497428226\\
21064.0947381272	-0.000831139163906947\\
21075.8149188133	-0.000711438520904894\\
21087.5350994994	-0.000598307129074585\\
21099.2552801856	-0.000492694339589269\\
21110.9754608717	-0.000395485299783425\\
21116.8355512147	-0.000350289394060698\\
21122.6956415578	-0.000307493506831298\\
21128.5557319009	-0.000267187008242536\\
21134.4158222439	-0.000229453971963535\\
21140.275912587	-0.000194372998928431\\
21146.13600293	-0.000162017052727458\\
21151.9960932731	-0.000132453306964239\\
21157.8561836162	-0.000105743004929365\\
21163.7162739592	-8.19413318335964e-05\\
21169.5763643023	-6.1097299924636e-05\\
21175.4364546453	-4.32536466786409e-05\\
21181.2965449884	-2.84467463083629e-05\\
21187.1566353315	-1.6706534756168e-05\\
21193.0167256745	-8.05644832862293e-06\\
21198.8768160176	-2.51337610274945e-06\\
21204.7369063606	-8.76261961876233e-08\\
21210.5969967037	-7.82905968610599e-07\\
21216.4570870468	-4.5963161948821e-06\\
21222.3171773898	-1.1518359214173e-05\\
21228.1772677329	-2.15329610357304e-05\\
21234.0373580759	-3.46175073520181e-05\\
21239.897448419	-5.07428933736981e-05\\
21245.7575387621	-6.98735873809977e-05\\
21251.6176291051	-9.19677078558858e-05\\
21257.4777194482	-0.000116977114014739\\
21263.3378097912	-0.000144847509567032\\
21269.1979001343	-0.000175518559451495\\
21275.0579904774	-0.000208924019328482\\
21280.9180808204	-0.000244991877513645\\
21286.7781711635	-0.000283644509092519\\
21292.6382615065	-0.000324798841846169\\
21304.3584421926	-0.000414254161034963\\
21316.0786228788	-0.000512591326913035\\
21327.7988035649	-0.00061896914622417\\
21339.518984251	-0.000732478899955249\\
21351.2391649371	-0.000852152069549209\\
21362.9593456232	-0.000976968564233929\\
21374.6795263094	-0.00110586537972965\\
21392.2597973385	-0.00130445474806615\\
21409.8400683677	-0.00150595524508733\\
21427.4203393969	-0.00170651903154536\\
21445.0006104261	-0.00190231942998053\\
21462.5808814553	-0.0020896236641936\\
21486.0212428275	-0.00232001175123109\\
21509.4616041997	-0.00252119409978788\\
21532.901965572	-0.00268638046715559\\
21556.3423269442	-0.00281000368688873\\
21579.7826883164	-0.00288790627088031\\
21603.2230496887	-0.00291747924063\\
21626.6634110609	-0.00289774855537648\\
21650.1037724332	-0.00282940627672859\\
21673.5441338054	-0.00271478547671585\\
21696.9844951776	-0.00255777979373768\\
21720.4248565499	-0.00236371040728484\\
21743.8652179221	-0.00213914497301313\\
21761.4454889513	-0.00195524304155938\\
21779.0257599805	-0.00176193957099966\\
21796.6060310096	-0.00156290913142642\\
21814.1863020388	-0.00136193264305908\\
21831.766573068	-0.00116282540728234\\
21849.3468440972	-0.000969364536003562\\
21861.0670247833	-0.000845345067356821\\
21872.7872054694	-0.000726513609956042\\
21884.5073861555	-0.000613870769674534\\
21896.2275668417	-0.000508364338302483\\
21907.9477475278	-0.000410881300080587\\
21919.6679282139	-0.000322240356207789\\
21925.5280185569	-0.000281471133043174\\
21931.3881089	-0.000243185030209084\\
21937.2481992431	-0.000207462215150681\\
21943.1082895861	-0.000174377412194723\\
21948.9683799292	-0.000143999745620105\\
21954.8284702722	-0.000116392594641658\\
21960.6885606153	-9.16134605594766e-05\\
21966.5486509584	-6.97138463841963e-05\\
21972.4087413014	-5.07391491756901e-05\\
21978.2688316445	-3.47285653200621e-05\\
21984.1289219875	-2.17150089512198e-05\\
21989.9890123306	-1.1725043677168e-05\\
21995.8491026737	-4.77882775783598e-06\\
22001.7091930167	-8.90072845246499e-07\\
22007.5692833598	-6.60163722716593e-08\\
22013.4293737028	-2.30740764265196e-06\\
22019.2894640459	-7.60850765080387e-06\\
22025.149554389	-1.59571026288878e-05\\
22031.009644732	-2.73345312889737e-05\\
22036.8697350751	-4.17157256969613e-05\\
22042.7298254181	-5.90692656961301e-05\\
22048.5899157612	-7.93574467543084e-05\\
22054.4500061043	-0.000102536361090831\\
22060.3100964473	-0.000128555991917456\\
22066.1701867904	-0.000157360320569037\\
22072.0302771334	-0.000188887446328079\\
22077.8903674765	-0.000223069718647253\\
22083.7504578196	-0.000259833881528768\\
22089.6105481626	-0.000299101229714281\\
22095.4706385057	-0.00034078777638871\\
22107.1908191918	-0.00043105719387192\\
22118.9109998779	-0.000529871669546716\\
22130.631180564	-0.00063638876178975\\
22142.3513612502	-0.000749701208001817\\
22154.0715419363	-0.000868844626112565\\
22165.7917226224	-0.000992805695996385\\
22183.3719936516	-0.00118546666293637\\
22200.9522646807	-0.00138291026549748\\
22218.5325357099	-0.00158136882234624\\
22236.1128067391	-0.00177705763348737\\
22253.6930777683	-0.00196624695904826\\
22271.2733487975	-0.002145332911633\\
22294.7137101697	-0.00236254086203783\\
22318.1540715419	-0.00254841778113761\\
22341.5944329142	-0.00269668955074039\\
22365.0347942864	-0.00280235687916964\\
22388.4751556586	-0.00286186316254732\\
22411.9155170309	-0.00287321353214946\\
22435.3558784031	-0.00283604110599474\\
22458.7962397754	-0.00275161824591257\\
22482.2366011476	-0.00262281247894199\\
22505.6769625198	-0.00245398860755634\\
22529.1173238921	-0.00225086034869867\\
22552.5576852643	-0.00202029654036547\\
22570.1379562935	-0.00183402866461428\\
22587.7182273227	-0.00164024049670971\\
22605.2984983518	-0.00144261668501902\\
22622.878769381	-0.00124491323039203\\
22640.4590404102	-0.00105088602549815\\
22652.1792210963	-0.000925401425474034\\
22663.8994017824	-0.000804250644338998\\
22675.6195824686	-0.000688456230458653\\
22687.3397631547	-0.000578995122488002\\
22699.0599438408	-0.000476790397472854\\
22710.7801245269	-0.000382703476811251\\
22716.64021487	-0.00033895472741485\\
22722.500305213	-0.000297526855936778\\
22728.3603955561	-0.000258507057649185\\
22734.2204858992	-0.000221977418955338\\
22740.0805762422	-0.000188014744367536\\
22745.9406665853	-0.000156690394694511\\
22751.8007569283	-0.000128070136766963\\
22757.6608472714	-0.000102214005008151\\
22763.5209376145	-7.91761751625779e-05\\
22769.3810279575	-5.90048504323481e-05\\
22775.2411183006	-4.17421602660924e-05\\
22781.1012086436	-2.74240720093062e-05\\
22786.9612989867	-1.60803156061958e-05\\
22792.8213893298	-7.73432150630204e-06\\
22798.6814796728	-2.40317190797747e-06\\
22804.5415700159	-9.75654389382338e-08\\
22810.4016603589	-8.21795347178264e-07\\
22816.261750702	-4.57374124570659e-06\\
22822.121841045	-1.13448744295206e-05\\
22827.9819313881	-2.11202767477588e-05\\
22833.8420217312	-3.38786729934728e-05\\
22839.7021120742	-4.95924767370388e-05\\
22845.5622024173	-6.82278495082941e-05\\
22851.4222927603	-8.97447731970047e-05\\
22857.2823831034	-0.000114097135518203\\
22863.1424734465	-0.000141232828365827\\
22869.0025637895	-0.000171093858826122\\
22874.8626541326	-0.000203616472647206\\
22880.7227444756	-0.000238731289862733\\
22886.5828348187	-0.000276363452316818\\
22892.4429251618	-0.000316432782750981\\
22904.1631058479	-0.000403536675847783\\
22915.883286534	-0.00049930191015217\\
22927.6034672201	-0.000602914317891804\\
22939.3236479062	-0.000713493527961728\\
22951.0438285924	-0.000830100429676494\\
22962.7640092785	-0.000951745131603513\\
22974.4841899646	-0.00107739534806444\\
22992.0644609938	-0.00127104229099939\\
23009.6447320229	-0.00146760582281903\\
23027.2250030521	-0.00166334005929498\\
23044.8052740813	-0.00185451617663497\\
23062.3855451105	-0.00203749332194539\\
23085.8259064827	-0.00226271870021161\\
23109.266267855	-0.00245958960357174\\
23132.7066292272	-0.0026214559182546\\
23156.1469905994	-0.00274285271512935\\
23179.5873519717	-0.00281968433563411\\
23203.0277133439	-0.00284936218507815\\
23226.4680747162	-0.00283089160029641\\
23249.9084360884	-0.00276490487818962\\
23273.3487974606	-0.002653639369865\\
23296.7891588329	-0.00250086139780925\\
23320.2295202051	-0.00231173858207379\\
23343.6698815773	-0.00209266490066311\\
23361.2501526065	-0.00191312729735072\\
23378.8304236357	-0.00172430801620584\\
23396.4106946649	-0.00152979955553598\\
23413.9909656941	-0.00133330200090315\\
23431.5712367232	-0.00113855262994018\\
23449.1515077524	-0.000949254823742937\\
23460.8716884385	-0.000827868621857204\\
23472.5918691246	-0.000711534937569924\\
23484.3120498108	-0.000601237361657966\\
23496.0322304969	-0.000497908301027104\\
23507.752411183	-0.000402421095557273\\
23519.4725918691	-0.000315582635360732\\
23525.3326822122	-0.000275638977228373\\
23531.1927725552	-0.000238126540927679\\
23537.0528628983	-0.000203124563407944\\
23542.9129532414	-0.000170706963778393\\
23548.7730435844	-0.000140942187121089\\
23554.6331339275	-0.000113893059913721\\
23560.4932242705	-8.96166573390238e-05\\
23566.3533146136	-6.81641827723735e-05\\
23572.2134049567	-4.95808596941853e-05\\
23578.0734952997	-3.39058362646826e-05\\
23583.9335856428	-2.1172102755075e-05\\
23589.7936759858	-1.14064220109323e-05\\
23595.6537663289	-4.62927309620972e-06\\
23601.513856672	-8.54808233626689e-07\\
23607.373947015	-9.0823133571775e-08\\
23613.2340373581	-2.33874077388739e-06\\
23619.0941277011	-7.59360866299011e-06\\
23624.9542180442	-1.58441095928798e-05\\
23630.8143083873	-2.7072585857694e-05\\
23636.6743987303	-4.12550768870749e-05\\
23642.5344890734	-5.83613702124949e-05\\
23648.3945794164	-7.83550656602712e-05\\
23654.2546697595	-0.000101193652633783\\
23660.1147601026	-0.000126828600321508\\
23665.9748504456	-0.00015520546063389\\
23671.8349407887	-0.000186263983670552\\
23677.6950311317	-0.000219938245436648\\
23683.5551214748	-0.000256156787581085\\
23689.4152118178	-0.000294842768824362\\
23695.2753021609	-0.000335914127787074\\
23706.995482847	-0.000424859686585813\\
23718.7156635331	-0.000522239440905372\\
23730.4358442193	-0.000627228047288775\\
23742.1560249054	-0.000738935829724721\\
23753.8762055915	-0.000856416311212569\\
23765.5963862776	-0.000978674225344446\\
23783.1766573068	-0.00116874445146766\\
23800.756928336	-0.00136360727140985\\
23818.3371993652	-0.0015595527887299\\
23835.9174703943	-0.00175285082931565\\
23853.4977414235	-0.00193982189351401\\
23871.0780124527	-0.00211690714030613\\
23894.5183738249	-0.0023318525077528\\
23917.9587351972	-0.00251599650596467\\
23941.3990965694	-0.00266311657780403\\
23964.8394579416	-0.00276824144011184\\
23988.2798193139	-0.00282781894062616\\
24000	-0.0028398359951788\\
};
\end{axis}
\end{tikzpicture}%

%%%% This file was created by matlab2tikz.
%
%The latest updates can be retrieved from
%  http://www.mathworks.com/matlabcentral/fileexchange/22022-matlab2tikz-matlab2tikz
%where you can also make suggestions and rate matlab2tikz.
%
\definecolor{mycolor1}{rgb}{0.00000,0.44700,0.74100}%
\definecolor{mycolor2}{rgb}{0.85000,0.32500,0.09800}%
%
\begin{tikzpicture}

\begin{axis}[%
width=4.521in,
height=3.548in,
at={(0.758in,0.499in)},
scale only axis,
unbounded coords=jump,
xmin=-25000,
xmax=25000,
ymode=log,
ymin=1e-08,
ymax=100,
yminorticks=true,
axis background/.style={fill=white}
]
\addplot [color=mycolor1, forget plot]
  table[row sep=crcr]{%
-24000	1.00284134051483\\
-9707.23965327799	1.00784668917116\\
-9443.53558784031	1.00634391250713\\
-9086.07007691369	1.00219220929594\\
-8664.1435722134	1.00853005589454\\
-8294.95788060066	1.0023857812845\\
-7878.89146624344	1.011376305327\\
-7509.70577463069	1.00240781529685\\
-7081.91917958735	1.01506004269969\\
-6730.31375900378	1.00217791442368\\
-6267.36662190209	1.01974901170664\\
-5968.50201440605	1.00096663918285\\
-5821.99975582957	1.01714225808576\\
-5611.03650347943	1.04099858046966\\
-5487.97460627518	1.03356722618652\\
-5177.38981809303	1.0013161485488\\
-5077.76828226102	1.01832174822303\\
-4802.34403613722	1.08121372328696\\
-4720.30277133439	1.07315522551496\\
-4626.54132584544	1.04486686922427\\
-4468.31888658283	0.996003132367466\\
-4409.71798315224	0.997418650516003\\
-4356.97717006471	1.0152401761057\\
-4304.23635697717	1.05120091934154\\
-4245.63545354658	1.11272091443234\\
-4169.4542790868	1.22121006145484\\
-4040.53229153949	1.43112405022432\\
-3981.9313881089	1.50941537509066\\
-3929.19057502137	1.55563075140529\\
-3882.30985227689	1.57034780457899\\
-3835.42912953241	1.55547668029589\\
-3794.408497131	1.515929511999\\
-3753.38786472958	1.45093497498028\\
-3718.22732267122	1.37546094305573\\
-3683.06678061287	1.28298625797261\\
-3647.90623855451	1.17540319674613\\
-3618.60578683921	1.07600980623875\\
-3589.30533512392	0.969613088671864\\
-3560.00488340862	0.858179393470407\\
-3536.56452203638	0.766835138670634\\
-3513.12416066415	0.674776579192027\\
-3489.68379929191	0.583176467998959\\
-3472.10352826273	0.51548717129855\\
-3454.52325723355	0.449213427234115\\
-3436.94298620437	0.384840462855972\\
-3419.36271517519	0.322840402176546\\
-3407.64253448907	0.283050703579549\\
-3395.92235380295	0.244647760125015\\
-3384.20217311684	0.20775640496713\\
-3372.48199243072	0.172496231612077\\
-3360.7618117446	0.138981138678295\\
-3354.90172140154	0.122911928192072\\
-3349.04163105848	0.107318904642378\\
-3343.18154071542	0.0922145195238283\\
-3337.32145037236	0.0776107944725708\\
-3331.4613600293	0.0635193108314036\\
-3325.60126968624	0.049951199770272\\
-3319.74117934318	0.0369171329732436\\
-3313.88108900012	0.0244273139028829\\
-3308.02099865706	0.0124914696441364\\
-3302.160908314	0.00111884333889737\\
-3296.30081797094	-0.00968181278675526\\
-3290.44072762788	-0.019902243794301\\
-3284.58063728482	-0.0295346978027968\\
-3278.72054694176	-0.0385719311511417\\
-3272.8604565987	-0.0470072129547623\\
-3267.00036625564	-0.0548343290500605\\
-3261.14027591259	-0.0620475853296203\\
-3255.28018556953	-0.068641810471287\\
-3249.42009522647	-0.0746123580559206\\
-3243.56000488341	-0.0799551080805467\\
-3237.69991454035	-0.0846664678712921\\
-3231.83982419729	-0.0887433723951299\\
-3220.11964351117	-0.0949841914586418\\
-3208.39946282505	-0.0986635725869213\\
-3196.67928213893	-0.0997758869663738\\
-3184.95910145281	-0.0983237780917271\\
-3173.2389207667	-0.0943180352165872\\
-3161.51874008058	-0.0877774331808311\\
-3155.65864973752	-0.083564563990025\\
-3149.79855939446	-0.0787285401467724\\
-3143.9384690514	-0.0732738638425516\\
-3138.07837870834	-0.0672054949415363\\
-3132.21828836528	-0.0605288424918328\\
-3126.35819802222	-0.0532497558352813\\
-3120.49810767916	-0.0453745153341191\\
-3114.6380173361	-0.036909822733815\\
-3108.77792699304	-0.0278627911743993\\
-3102.91783664998	-0.0182409348719961\\
-3097.05774630692	-0.00805215848653547\\
-3091.19765596386	0.00269525380486109\\
-3085.3375656208	0.0139926495122191\\
-3079.47747527774	0.0258310192938315\\
-3073.61738493469	0.0382010086571355\\
-3067.75729459163	0.051092929879601\\
-3061.89720424857	0.0644967741307062\\
-3056.03711390551	0.0784022237753743\\
-3050.17702356245	0.0927986648406104\\
-3044.31693321939	0.107675199625698\\
-3038.45684287633	0.123020659436337\\
-3032.59675253327	0.138823617423142\\
-3020.87657184715	0.171755107382671\\
-3009.15639116103	0.206373584339024\\
-2997.43621047491	0.242579665254818\\
-2985.71602978879	0.280271087478488\\
-2973.99584910267	0.319343125452807\\
-2956.41557807349	0.380305805263977\\
-2938.83530704432	0.443767362927457\\
-2921.25503601514	0.509355604286703\\
-2897.8146746429	0.599467263808462\\
-2874.37431327066	0.691807769167712\\
-2850.93395189842	0.785502690256061\\
-2821.63350018313	0.903228572744357\\
-2792.33304846783	1.02013891276253\\
-2757.17250640948	1.15723788805243\\
-2722.01196435112	1.2885943760245\\
-2686.85142229276	1.41205386257115\\
-2645.83078989134	1.54359360095767\\
-2604.81015748993	1.65924664745704\\
-2557.92943474545	1.76928374103485\\
-2511.04871200098	1.85336415918948\\
-2464.1679892565	1.90968229154065\\
-2411.42717616897	1.93838352068749\\
-2358.68636308143	1.92960648867852\\
-2311.80564033696	1.89039774211996\\
-2264.92491759248	1.8223753875837\\
-2218.044194848	1.72697563229736\\
-2171.16347210353	1.6064104532494\\
-2130.14283970211	1.48264976897469\\
-2089.1222073007	1.34453613945666\\
-2053.96166524234	1.21706267650456\\
-2018.80112318398	1.08353685194144\\
-1983.64058112563	0.946409026532242\\
-1954.34012941033	0.831279261727773\\
-1925.03967769503	0.71706815215248\\
-1901.59931632279	0.62745598313109\\
-1878.15895495056	0.54031517120234\\
-1854.71859357832	0.456479003249757\\
-1837.13832254914	0.396262514792276\\
-1819.55805151996	0.338699265508007\\
-1801.97778049078	0.284112704628335\\
-1784.3975094616	0.232812600360447\\
-1772.67732877548	0.200584104372002\\
-1760.95714808936	0.17002973578995\\
-1749.23696740325	0.141228143238201\\
-1737.51678671713	0.114253773031927\\
-1731.65669637407	0.101473846214038\\
-1725.79660603101	0.089176604251686\\
-1719.93651568795	0.0773700555255055\\
-1714.07642534489	0.066061898860856\\
-1708.21633500183	0.0552595166417749\\
-1702.35624465877	0.0449699681951868\\
-1696.49615431571	0.0351999834814129\\
-1690.63606397265	0.0259559570753332\\
-1684.77597362959	0.0172439424659039\\
-1678.91588328653	0.00906964666595916\\
-1673.05579294348	0.00143842514990056\\
-1667.19570260042	-0.0056447228804401\\
-1661.33561225736	-0.0121751588931026\\
-1655.4755219143	-0.0181486090815703\\
-1649.61543157124	-0.0235611680734111\\
-1643.75534122818	-0.0284093022881394\\
-1637.89525088512	-0.0326898529396692\\
-1632.03516054206	-0.03640003867962\\
-1626.175070199	-0.0395374578802563\\
-1620.31497985594	-0.0421000905512338\\
-1614.45488951288	-0.0440862998926634\\
-1608.59479916982	-0.0454948334797297\\
-1602.73470882676	-0.0463248240782844\\
-1596.8746184837	-0.0465757900940885\\
-1591.01452814064	-0.0462476356531751\\
-1585.15443779758	-0.0453406503121169\\
-1579.29434745452	-0.0438555084071351\\
-1573.43425711146	-0.0417932680339995\\
-1567.57416676841	-0.0391553696685281\\
-1561.71407642535	-0.0359436344278396\\
-1555.85398608229	-0.0321602619750341\\
-1549.99389573923	-0.0278078280714245\\
-1544.13380539617	-0.022889281782841\\
-1538.27371505311	-0.0174079423414557\\
-1532.41362471005	-0.0113674956705275\\
-1526.55353436699	-0.00477199057791364\\
-1520.69344402393	0.00237416537717297\\
-1514.83335368087	0.0100662103354489\\
-1508.97326333781	0.018299032901186\\
-1503.11317299475	0.0270671772101309\\
-1497.25308265169	0.036364848295125\\
-1491.39299230863	0.0461859177382354\\
-1485.53290196557	0.0565239296041517\\
-1479.67281162251	0.0673721066447664\\
-1473.81272127945	0.0787233567635679\\
-1467.95263093639	0.0905702797341611\\
-1462.09254059333	0.10290517415985\\
-1450.37235990721	0.129006609313262\\
-1438.6521792211	0.156960306392116\\
-1426.93199853498	0.186694573222286\\
-1415.21181784886	0.218133640395567\\
-1403.49163716274	0.251197931015164\\
-1385.91136613356	0.30365897324134\\
-1368.33109510438	0.359295171602081\\
-1350.7508240752	0.417798500253555\\
-1333.17055304603	0.478849030566786\\
-1309.73019167379	0.563641378644391\\
-1286.28983030155	0.651575047520894\\
-1262.84946892931	0.741835985968858\\
-1233.54901721402	0.856688017009298\\
-1204.24856549872	0.97230985832773\\
-1169.08802344036	1.1098882741476\\
-1133.927481382	1.24376503761306\\
-1098.76693932365	1.37152395594126\\
-1057.74630692223	1.50989246662556\\
-1016.72567452081	1.63377494418881\\
-969.84495177634	1.75417409627901\\
-922.964229031862	1.84885086771245\\
-876.083506287388	1.91531165235299\\
-829.202783542914	1.95185225640328\\
-782.322060798437	1.95754641903256\\
-735.441338053963	1.93223937442644\\
-688.560615309485	1.87655128693148\\
-641.679892565011	1.79189033012841\\
-594.799169820533	1.68047018598328\\
-553.778537419119	1.56338990081338\\
-512.757905017701	1.43067408083479\\
-471.737272616287	1.28532719492052\\
-436.576730557928	1.15328171802749\\
-401.416188499574	1.01685311452385\\
-372.115736784275	0.90164800901762\\
-342.815285068977	0.786705924625891\\
-313.514833353682	0.673604005400197\\
-290.074471981443	0.58551619082351\\
-266.634110609204	0.500433436681232\\
-243.193749236969	0.419150848011362\\
-225.61347820779	0.361148050051534\\
-208.033207178611	0.306031303990695\\
-190.452936149431	0.254101786166131\\
-178.732755463312	0.221394895211982\\
-167.012574777196	0.190312965254502\\
-155.292394091077	0.160933171753497\\
-143.572213404957	0.133328829423154\\
-131.852032718838	0.107569144067225\\
-125.991942375778	0.0954014925732736\\
-120.131852032719	0.0837189760648634\\
-114.271761689659	0.0725289864276536\\
-108.411671346599	0.0618386167720763\\
-102.551581003539	0.0516546552306847\\
-96.6914906604798	0.0419835790077327\\
-90.83140031742	0.0328315486793118\\
-84.9713099743603	0.0242044027596892\\
-79.1112196313043	0.0161076525315924\\
-73.2511292882446	0.00854647715867775\\
-67.3910389451848	0.00152571907700035\\
-61.5309486021251	-0.00495012032126534\\
-55.6708582590654	-0.0108768847103271\\
-49.8107679160057	-0.0162507665513866\\
-43.950677572946	-0.021068310507595\\
-38.0905872298863	-0.0253264165483282\\
-32.2304968868266	-0.0290223427374395\\
-26.3704065437669	-0.032153707704045\\
-20.5103162007072	-0.0347184927898288\\
-14.6502258576475	-0.0367150438710176\\
-8.79013551458775	-0.0381420728534124\\
-2.93004517152804	-0.0389986588368639\\
2.93004517152804	-0.039284248946786\\
8.79013551458775	-0.0389986588368639\\
14.6502258576475	-0.0381420728534124\\
20.5103162007072	-0.0367150438710176\\
26.3704065437669	-0.0347184927898288\\
32.2304968868266	-0.032153707704045\\
38.0905872298863	-0.0290223427374395\\
43.950677572946	-0.0253264165483282\\
49.8107679160057	-0.021068310507595\\
55.6708582590654	-0.0162507665513866\\
61.5309486021251	-0.0108768847103271\\
67.3910389451848	-0.00495012032126534\\
73.2511292882446	0.00152571907700035\\
79.1112196313043	0.00854647715867775\\
84.9713099743603	0.0161076525315924\\
90.83140031742	0.0242044027596892\\
96.6914906604798	0.0328315486793118\\
102.551581003539	0.0419835790077327\\
108.411671346599	0.0516546552306847\\
114.271761689659	0.0618386167720763\\
120.131852032719	0.0725289864276536\\
125.991942375778	0.0837189760648634\\
131.852032718838	0.0954014925732736\\
137.712123061898	0.107569144067225\\
149.432303748017	0.133328829423154\\
161.152484434137	0.160933171753497\\
172.872665120252	0.190312965254502\\
184.592845806372	0.221394895211982\\
196.313026492491	0.254101786166131\\
213.89329752167	0.306031303990695\\
231.473568550849	0.361148050051534\\
249.053839580029	0.419150848011362\\
272.494200952264	0.500433436681232\\
295.934562324503	0.58551619082351\\
319.374923696741	0.673604005400197\\
348.675375412036	0.786705924625891\\
377.975827127335	0.90164800901762\\
407.276278842633	1.01685311452385\\
442.436820900988	1.15328171802749\\
477.597362959346	1.28532719492052\\
518.617995360761	1.43067408083479\\
559.638627762179	1.56338990081338\\
600.659260163593	1.68047018598328\\
647.539982908071	1.79189033012841\\
694.420705652545	1.87655128693148\\
741.301428397022	1.93223937442644\\
788.182151141496	1.95754641903256\\
835.062873885974	1.95185225640328\\
881.943596630448	1.91531165235299\\
928.824319374922	1.84885086771245\\
975.7050421194	1.75417409627901\\
1022.58576486387	1.63377494418881\\
1063.60639726529	1.50989246662556\\
1104.62702966671	1.37152395594126\\
1145.64766206812	1.22181408965479\\
1180.80820412648	1.08714779382094\\
1215.96874618484	0.949199155478187\\
1245.26919790014	0.833605562337403\\
1274.56964961543	0.719097343665334\\
1298.01001098767	0.62934150291742\\
1321.45037235991	0.54211722323692\\
1344.89073373214	0.458236000732144\\
1362.47100476132	0.397998349569571\\
1380.0512757905	0.340415724088584\\
1397.63154681968	0.285804339612201\\
1415.21181784886	0.234467747490499\\
1426.93199853498	0.202205894311817\\
1438.6521792211	0.171609511729068\\
1450.37235990721	0.142756307220438\\
1462.09254059333	0.115720044667061\\
1467.95263093639	0.10290517415985\\
1473.81272127945	0.0905702797341611\\
1479.67281162251	0.0787233567635679\\
1485.53290196557	0.0673721066447664\\
1491.39299230863	0.0565239296041517\\
1497.25308265169	0.0461859177382354\\
1503.11317299475	0.036364848295125\\
1508.97326333781	0.0270671772101309\\
1514.83335368087	0.018299032901186\\
1520.69344402393	0.0100662103354489\\
1526.55353436699	0.00237416537717297\\
1532.41362471005	-0.00477199057791364\\
1538.27371505311	-0.0113674956705275\\
1544.13380539617	-0.0174079423414557\\
1549.99389573923	-0.022889281782841\\
1555.85398608229	-0.0278078280714245\\
1561.71407642535	-0.0321602619750341\\
1567.57416676841	-0.0359436344278396\\
1573.43425711146	-0.0391553696685281\\
1579.29434745452	-0.0417932680339995\\
1585.15443779758	-0.0438555084071351\\
1591.01452814064	-0.0453406503121169\\
1596.8746184837	-0.0462476356531751\\
1602.73470882676	-0.0465757900940885\\
1608.59479916982	-0.0463248240782844\\
1614.45488951288	-0.0454948334797297\\
1620.31497985594	-0.0440862998926634\\
1626.175070199	-0.0421000905512338\\
1632.03516054206	-0.0395374578802563\\
1637.89525088512	-0.03640003867962\\
1643.75534122818	-0.0326898529396692\\
1649.61543157124	-0.0284093022881394\\
1655.4755219143	-0.0235611680734111\\
1661.33561225736	-0.0181486090815703\\
1667.19570260042	-0.0121751588931026\\
1673.05579294348	-0.0056447228804401\\
1678.91588328653	0.00143842514990056\\
1684.77597362959	0.00906964666595916\\
1690.63606397265	0.0172439424659039\\
1696.49615431571	0.0259559570753332\\
1702.35624465877	0.0351999834814129\\
1708.21633500183	0.0449699681951868\\
1714.07642534489	0.0552595166417749\\
1719.93651568795	0.066061898860856\\
1725.79660603101	0.0773700555255055\\
1731.65669637407	0.089176604251686\\
1737.51678671713	0.101473846214038\\
1743.37687706019	0.114253773031927\\
1755.09705774631	0.141228143238201\\
1766.81723843242	0.17002973578995\\
1778.53741911854	0.200584104372002\\
1790.25759980466	0.232812600360447\\
1801.97778049078	0.26663265139941\\
1819.55805151996	0.320157400015862\\
1837.13832254914	0.376763733173268\\
1854.71859357832	0.43613249110897\\
1878.15895495056	0.519014062109301\\
1901.59931632279	0.605406347429255\\
1925.03967769503	0.694479336996092\\
1954.34012941033	0.808311718652108\\
1983.64058112563	0.923386843517411\\
2012.94103284092	1.03810107619541\\
2048.10157489928	1.17310981081574\\
2083.26211695764	1.30285987535807\\
2124.28274935905	1.44452721005051\\
2165.30338176047	1.57267556300461\\
2206.32401416188	1.68456770292539\\
2253.20473690636	1.78970764655892\\
2300.08545965084	1.86820488191576\\
2346.96618239531	1.91834094544438\\
2399.70699548285	1.93963704193348\\
2452.44780857038	1.92338067716478\\
2499.32853131486	1.87781407428582\\
2546.20925405933	1.8039664386101\\
2593.08997680381	1.70341094936495\\
2639.97069954828	1.57835899148366\\
2680.9913319497	1.4511168886958\\
2722.01196435112	1.30977379476686\\
2757.17250640948	1.17959075230625\\
2792.33304846783	1.04328853805076\\
2821.63350018313	0.926725599194983\\
2850.93395189842	0.809037455065404\\
2874.37431327066	0.715138286231451\\
2897.8146746429	0.622377867674489\\
2921.25503601514	0.531626214573054\\
2938.83530704432	0.465412288772645\\
2956.41557807349	0.401200303594941\\
2973.99584910267	0.339363576894535\\
2985.71602978879	0.299641159725105\\
2997.43621047491	0.261246201189658\\
3009.15639116103	0.224284503124075\\
3020.87657184715	0.188859611535734\\
3032.59675253327	0.155072401510128\\
3038.45684287633	0.138823617423142\\
3044.31693321939	0.123020659436337\\
3050.17702356245	0.107675199625698\\
3056.03711390551	0.0927986648406104\\
3061.89720424857	0.0784022237753743\\
3067.75729459163	0.0644967741307062\\
3073.61738493469	0.051092929879601\\
3079.47747527774	0.0382010086571355\\
3085.3375656208	0.0258310192938315\\
3091.19765596386	0.0139926495122191\\
3097.05774630692	0.00269525380486109\\
3102.91783664998	-0.00805215848653547\\
3108.77792699304	-0.0182409348719961\\
3114.6380173361	-0.0278627911743993\\
3120.49810767916	-0.036909822733815\\
3126.35819802222	-0.0453745153341191\\
3132.21828836528	-0.0532497558352813\\
3138.07837870834	-0.0605288424918328\\
3143.9384690514	-0.0672054949415363\\
3149.79855939446	-0.0732738638425516\\
3155.65864973752	-0.0787285401467724\\
3161.51874008058	-0.083564563990025\\
3167.37883042364	-0.0877774331808311\\
3179.09901110976	-0.0943180352165872\\
3190.81919179587	-0.0983237780917271\\
3202.53937248199	-0.0997758869663738\\
3214.25955316811	-0.0986635725869213\\
3225.97973385423	-0.0949841914586418\\
3237.69991454035	-0.0887433723951299\\
3243.56000488341	-0.0846664678712921\\
3249.42009522647	-0.0799551080805467\\
3255.28018556953	-0.0746123580559206\\
3261.14027591259	-0.068641810471287\\
3267.00036625564	-0.0620475853296203\\
3272.8604565987	-0.0548343290500605\\
3278.72054694176	-0.0470072129547623\\
3284.58063728482	-0.0385719311511417\\
3290.44072762788	-0.0295346978027968\\
3296.30081797094	-0.019902243794301\\
3302.160908314	-0.00968181278675526\\
3308.02099865706	0.00111884333889737\\
3313.88108900012	0.0124914696441364\\
3319.74117934318	0.0244273139028829\\
3325.60126968624	0.0369171329732436\\
3331.4613600293	0.049951199770272\\
3337.32145037236	0.0635193108314036\\
3343.18154071542	0.0776107944725708\\
3349.04163105848	0.0922145195238283\\
3354.90172140154	0.107318904642378\\
3360.7618117446	0.122911928192072\\
3366.62190208766	0.138981138678295\\
3378.34208277378	0.172496231612077\\
3390.06226345989	0.20775640496713\\
3401.78244414601	0.244647760125015\\
3413.50262483213	0.283050703579549\\
3425.22280551825	0.322840402176546\\
3442.80307654743	0.384840462855972\\
3460.38334757661	0.449213427234115\\
3477.96361860579	0.51548717129855\\
3501.40397997803	0.605968936689599\\
3524.84434135026	0.697794319707401\\
3548.2847027225	0.789783546201827\\
3577.5851544378	0.903223892767596\\
3606.88560615309	1.01289063691337\\
3636.18605786839	1.11671087935326\\
3671.34659992675	1.23093964877382\\
3706.5071419851	1.33124592697357\\
3741.66768404346	1.41541503731414\\
3782.68831644488	1.4911806242377\\
3823.7089488463	1.5416340264825\\
3864.72958124771	1.56683176098259\\
3911.61030399219	1.56643960107551\\
3958.49102673666	1.53854026129323\\
4011.2318398242	1.4808441855122\\
4069.83274325479	1.3946357643897\\
4157.73409840068	1.2494405716888\\
4257.3556342327	1.10560367653108\\
4315.95653766329	1.04628045551816\\
4368.69735075082	1.01239104258024\\
4421.43816383836	0.996473557808924\\
4480.03906726895	0.996743029888078\\
4556.22024172873	1.01643403533274\\
4749.60322304969	1.07729493343014\\
4831.64448785252	1.08019849875199\\
4925.40593334147	1.06374336103988\\
5212.55036015138	0.999660690850909\\
5318.03198632646	1.00562490882278\\
5640.33695519473	1.04025027433413\\
5780.97912342815	1.02442129157638\\
5997.80246612135	1.00012251061262\\
6138.44463435478	1.0058972902668\\
6437.30924185081	1.02581498795457\\
6618.97204248565	1.01181235050572\\
6824.07520449274	0.999955278799365\\
7011.59809547064	1.00971356554002\\
7240.14161884996	1.01839349847716\\
7480.40532291539	1.00430903784712\\
7679.64839457942	1.00099296556962\\
8183.61616408253	1.00837668572661\\
8476.6206812355	1.00075409056754\\
9015.74899279697	1.00527477136276\\
9302.89341960689	1.00124068173093\\
9783.42082773776	1.00553615605415\\
10129.1661579783	1.00170145932106\\
10574.5330240508	1.00496281748533\\
10955.4388963497	1.00211668208415\\
11377.3654010499	1.00421894135086\\
11775.851544378	1.00233799705817\\
12186.0578683921	1.00350960004946\\
12596.2641924063	1.00252579591822\\
13006.4705164205	1.00270184819896\\
13416.6768404346	1.00268272959236\\
13826.8831644488	1.00205310528284\\
14242.949578806	1.0029115097727\\
14659.0159931632	1.00133875633935\\
15086.8025882066	1.00332738854949\\
15532.1694542791	1.00034476341141\\
16516.6646319131	1.00064416803259\\
17249.1759247955	1.00009481476089\\
19399.8290806983	1.00165341641159\\
19991.6982053473	1.00305185813636\\
24000	1.00283983599301\\
};
\addplot [color=mycolor2, forget plot]
  table[row sep=crcr]{%
-24000	-0.00284134051541254\\
-23976.5596386278	-0.00281733185141145\\
-23953.1192772555	-0.00274611709168084\\
-23929.6789158833	-0.00263010254264651\\
-23906.238554511	-0.00247320833745444\\
-23882.7981931388	-0.00228073607346066\\
-23859.3578317666	-0.00205918980113674\\
-23841.7775607374	-0.00187842412257593\\
-23824.1972897082	-0.00168894126064206\\
-23806.617018679	-0.00149434793701714\\
-23789.0367476499	-0.0012983483943787\\
-23771.4564766207	-0.00110467394012934\\
-23759.7362959346	-0.000978674225344446\\
-23748.0161152484	-0.000856416311212569\\
-23736.2959345623	-0.000738935829724721\\
-23724.5757538762	-0.000627228047288775\\
-23712.8555731901	-0.000522239440905372\\
-23701.135392504	-0.000424859686585813\\
-23689.4152118178	-0.000335914127787074\\
-23683.5551214748	-0.000294842768824362\\
-23677.6950311317	-0.000256156787581085\\
-23671.8349407887	-0.000219938245436648\\
-23665.9748504456	-0.000186263983670552\\
-23660.1147601026	-0.00015520546063389\\
-23654.2546697595	-0.000126828600321508\\
-23648.3945794164	-0.000101193652633783\\
-23642.5344890734	-7.83550656602712e-05\\
-23636.6743987303	-5.83613702124949e-05\\
-23630.8143083873	-4.12550768870749e-05\\
-23624.9542180442	-2.7072585857694e-05\\
-23619.0941277011	-1.58441095928798e-05\\
-23613.2340373581	-7.59360866299011e-06\\
-23607.373947015	-2.33874077388739e-06\\
-23601.513856672	-9.0823133571775e-08\\
-23595.6537663289	-8.54808233626689e-07\\
-23589.7936759858	-4.62927309620972e-06\\
-23583.9335856428	-1.14064220109323e-05\\
-23578.0734952997	-2.1172102755075e-05\\
-23572.2134049567	-3.39058362646826e-05\\
-23566.3533146136	-4.95808596941853e-05\\
-23560.4932242705	-6.81641827723735e-05\\
-23554.6331339275	-8.96166573390238e-05\\
-23548.7730435844	-0.000113893059913721\\
-23542.9129532414	-0.000140942187121089\\
-23537.0528628983	-0.000170706963778393\\
-23531.1927725552	-0.000203124563407944\\
-23525.3326822122	-0.000238126540927679\\
-23519.4725918691	-0.000275638977228373\\
-23513.6125015261	-0.000315582635360732\\
-23501.8923208399	-0.000402421095557273\\
-23490.1721401538	-0.000497908301027104\\
-23478.4519594677	-0.000601237361657966\\
-23466.7317787816	-0.000711534937569924\\
-23455.0115980955	-0.000827868621857204\\
-23443.2914174093	-0.000949254823742937\\
-23425.7111463802	-0.00113855262994018\\
-23408.130875351	-0.00133330200090315\\
-23390.5506043218	-0.00152979955553598\\
-23372.9703332926	-0.00172430801620584\\
-23355.3900622635	-0.00191312729735072\\
-23337.8097912343	-0.00209266490066311\\
-23314.369429862	-0.00231173858207379\\
-23290.9290684898	-0.00250086139780925\\
-23267.4887071176	-0.002653639369865\\
-23244.0483457453	-0.00276490487818962\\
-23220.6079843731	-0.00283089160029641\\
-23197.1676230009	-0.00284936218507815\\
-23173.7272616286	-0.00281968433563411\\
-23150.2869002564	-0.00274285271512935\\
-23126.8465388841	-0.0026214559182546\\
-23103.4061775119	-0.00245958960357174\\
-23079.9658161397	-0.00226271870021161\\
-23056.5254547674	-0.00203749332194539\\
-23038.9451837382	-0.00185451617663497\\
-23021.3649127091	-0.00166334005929498\\
-23003.7846416799	-0.00146760582281903\\
-22986.2043706507	-0.00127104229099939\\
-22968.6240996215	-0.00107739534806444\\
-22956.9039189354	-0.000951745131603513\\
-22945.1837382493	-0.000830100429676494\\
-22933.4635575632	-0.000713493527961728\\
-22921.7433768771	-0.000602914317891804\\
-22910.0231961909	-0.00049930191015217\\
-22898.3030155048	-0.000403536675847783\\
-22886.5828348187	-0.000316432782750981\\
-22880.7227444756	-0.000276363452316818\\
-22874.8626541326	-0.000238731289862733\\
-22869.0025637895	-0.000203616472647206\\
-22863.1424734465	-0.000171093858826122\\
-22857.2823831034	-0.000141232828365827\\
-22851.4222927603	-0.000114097135518203\\
-22845.5622024173	-8.97447731970047e-05\\
-22839.7021120742	-6.82278495082941e-05\\
-22833.8420217312	-4.95924767370388e-05\\
-22827.9819313881	-3.38786729934728e-05\\
-22822.121841045	-2.11202767477588e-05\\
-22816.261750702	-1.13448744295206e-05\\
-22810.4016603589	-4.57374124570659e-06\\
-22804.5415700159	-8.21795347178264e-07\\
-22798.6814796728	-9.75654389382338e-08\\
-22792.8213893298	-2.40317190797747e-06\\
-22786.9612989867	-7.73432150630204e-06\\
-22781.1012086436	-1.60803156061958e-05\\
-22775.2411183006	-2.74240720093062e-05\\
-22769.3810279575	-4.17421602660924e-05\\
-22763.5209376145	-5.90048504323481e-05\\
-22757.6608472714	-7.91761751625779e-05\\
-22751.8007569283	-0.000102214005008151\\
-22745.9406665853	-0.000128070136766963\\
-22740.0805762422	-0.000156690394694511\\
-22734.2204858992	-0.000188014744367536\\
-22728.3603955561	-0.000221977418955338\\
-22722.500305213	-0.000258507057649185\\
-22716.64021487	-0.000297526855936778\\
-22710.7801245269	-0.00033895472741485\\
-22699.0599438408	-0.000428680983837285\\
-22687.3397631547	-0.000526930340305221\\
-22675.6195824686	-0.000632874964801485\\
-22663.8994017824	-0.000745621665077606\\
-22652.1792210963	-0.000864219428830506\\
-22640.4590404102	-0.000987667455107538\\
-22622.878769381	-0.00117964492654475\\
-22605.2984983518	-0.00137653821552113\\
-22587.7182273227	-0.00157460737881394\\
-22570.1379562935	-0.00177008848521255\\
-22552.5576852643	-0.00195926521608474\\
-22534.9774142351	-0.00213853959297041\\
-22511.5370528629	-0.00235630774269005\\
-22488.0966914907	-0.00254307525048042\\
-22464.6563301184	-0.00269252401529647\\
-22441.2159687462	-0.00279959325577728\\
-22417.7756073739	-0.00286065131759589\\
-22394.3352460017	-0.0028736193259035\\
-22370.8948846295	-0.00283804247544993\\
-22347.4545232572	-0.00275510651585241\\
-22324.014161885	-0.00262759883751439\\
-22300.5738005128	-0.00245981543344755\\
-22277.1334391405	-0.00225741684033257\\
-22253.6930777683	-0.00202723788587094\\
-22236.1128067391	-0.0018410241368157\\
-22218.5325357099	-0.0016470921792747\\
-22200.9522646807	-0.00144913758640776\\
-22183.3719936516	-0.00125093471025438\\
-22165.7917226224	-0.00105626491351087\\
-22154.0715419363	-0.000930287966284631\\
-22142.3513612502	-0.000808606329948592\\
-22130.631180564	-0.000692254594489036\\
-22118.9109998779	-0.000582222661009546\\
-22107.1908191918	-0.000479447346117513\\
-22095.4706385057	-0.000384804432001888\\
-22089.6105481626	-0.00034078777638871\\
-22083.7504578196	-0.000299101229714281\\
-22077.8903674765	-0.000259833881528768\\
-22072.0302771334	-0.000223069718647253\\
-22066.1701867904	-0.000188887446328079\\
-22060.3100964473	-0.000157360320569037\\
-22054.4500061043	-0.000128555991917456\\
-22048.5899157612	-0.000102536361090831\\
-22042.7298254181	-7.93574467543084e-05\\
-22036.8697350751	-5.90692656961301e-05\\
-22031.009644732	-4.17157256969613e-05\\
-22025.149554389	-2.73345312889737e-05\\
-22019.2894640459	-1.59571026288878e-05\\
-22013.4293737028	-7.60850765080387e-06\\
-22007.5692833598	-2.30740764265196e-06\\
-22001.7091930167	-6.60163722716593e-08\\
-21995.8491026737	-8.90072845246499e-07\\
-21989.9890123306	-4.77882775783598e-06\\
-21984.1289219875	-1.1725043677168e-05\\
-21978.2688316445	-2.17150089512198e-05\\
-21972.4087413014	-3.47285653200621e-05\\
-21966.5486509584	-5.07391491756901e-05\\
-21960.6885606153	-6.97138463841963e-05\\
-21954.8284702722	-9.16134605594766e-05\\
-21948.9683799292	-0.000116392594641658\\
-21943.1082895861	-0.000143999745620105\\
-21937.2481992431	-0.000174377412194723\\
-21931.3881089	-0.000207462215150681\\
-21925.5280185569	-0.000243185030209084\\
-21919.6679282139	-0.000281471133043174\\
-21913.8078378708	-0.000322240356207789\\
-21902.0876571847	-0.000410881300080587\\
-21890.3674764986	-0.000508364338302483\\
-21878.6472958125	-0.000613870769674534\\
-21866.9271151264	-0.000726513609956042\\
-21855.2069344402	-0.000845345067356821\\
-21843.4867537541	-0.000969364536003562\\
-21825.9064827249	-0.00116282540728234\\
-21808.3262116958	-0.00136193264305908\\
-21790.7459406666	-0.00156290913142642\\
-21773.1656696374	-0.00176193957099966\\
-21755.5853986082	-0.00195524304155938\\
-21738.0051275791	-0.00213914497301313\\
-21714.5647662068	-0.00236371040728484\\
-21691.1244048346	-0.00255777979373768\\
-21667.6840434623	-0.00271478547671585\\
-21644.2436820901	-0.00282940627672859\\
-21620.8033207179	-0.00289774855537648\\
-21597.3629593456	-0.00291747924063\\
-21573.9225979734	-0.00288790627088031\\
-21550.4822366011	-0.00281000368688873\\
-21527.0418752289	-0.00268638046715559\\
-21503.6015138567	-0.00252119409978788\\
-21480.1611524844	-0.00232001175123109\\
-21456.7207911122	-0.0020896236641936\\
-21439.140520083	-0.00190231942998053\\
-21421.5602490538	-0.00170651903154536\\
-21403.9799780247	-0.00150595524508733\\
-21386.3997069955	-0.00130445474806615\\
-21368.8194359663	-0.00110586537972965\\
-21357.0992552802	-0.000976968564233929\\
-21345.3790745941	-0.000852152069549209\\
-21333.6588939079	-0.000732478899955249\\
-21321.9387132218	-0.00061896914622417\\
-21310.2185325357	-0.000512591326913035\\
-21298.4983518496	-0.000414254161034963\\
-21286.7781711635	-0.000324798841846169\\
-21280.9180808204	-0.000283644509092519\\
-21275.0579904774	-0.000244991877513645\\
-21269.1979001343	-0.000208924019328482\\
-21263.3378097912	-0.000175518559451495\\
-21257.4777194482	-0.000144847509567032\\
-21251.6176291051	-0.000116977114014739\\
-21245.7575387621	-9.19677078558858e-05\\
-21239.897448419	-6.98735873809977e-05\\
-21234.0373580759	-5.07428933736981e-05\\
-21228.1772677329	-3.46175073520181e-05\\
-21222.3171773898	-2.15329610357304e-05\\
-21216.4570870468	-1.1518359214173e-05\\
-21210.5969967037	-4.5963161948821e-06\\
-21204.7369063606	-7.82905968610599e-07\\
-21198.8768160176	-8.76261961876233e-08\\
-21193.0167256745	-2.51337610274945e-06\\
-21187.1566353315	-8.05644832862293e-06\\
-21181.2965449884	-1.6706534756168e-05\\
-21175.4364546453	-2.84467463083629e-05\\
-21169.5763643023	-4.32536466786409e-05\\
-21163.7162739592	-6.1097299924636e-05\\
-21157.8561836162	-8.19413318335964e-05\\
-21151.9960932731	-0.000105743004929365\\
-21146.13600293	-0.000132453306964239\\
-21140.275912587	-0.000162017052727458\\
-21134.4158222439	-0.000194372998928431\\
-21128.5557319009	-0.000229453971963535\\
-21122.6956415578	-0.000267187008242536\\
-21116.8355512147	-0.000307493506831298\\
-21110.9754608717	-0.000350289394060698\\
-21099.2552801856	-0.00044298674486833\\
-21087.5350994994	-0.000544503992459431\\
-21075.8149188133	-0.000653990920535655\\
-21064.0947381272	-0.000770529313278586\\
-21052.3745574411	-0.000893140683809489\\
-21040.654376755	-0.00102079451760563\\
-21023.0741057258	-0.001219371205951\\
-21005.4938346966	-0.0014231102353403\\
-20987.9135636674	-0.00162815076132841\\
-20970.3332926383	-0.001830603462985\\
-20952.7530216091	-0.00202662453111203\\
-20935.1727505799	-0.0022124888636091\\
-20911.7323892077	-0.00243843387551629\\
-20888.2920278354	-0.00263242710895824\\
-20864.8516664632	-0.00278789881637128\\
-20841.411305091	-0.00289957289827388\\
-20817.9709437187	-0.00296364698937434\\
-20794.5305823465	-0.00297792299951192\\
-20771.0902209742	-0.00294188363257645\\
-20747.649859602	-0.00285671222316417\\
-20724.2094982298	-0.00272525514233901\\
-20700.7691368575	-0.00255192796067947\\
-20677.3287754853	-0.00234256845891991\\
-20653.8884141131	-0.00210424137417143\\
-20636.3081430839	-0.00191130174259514\\
-20618.7278720547	-0.00171026079202968\\
-20601.1476010255	-0.00150495390425756\\
-20583.5673299963	-0.00129930204105443\\
-20565.9870589672	-0.00109723721700748\\
-20554.266878281	-0.000966435241031535\\
-20542.5466975949	-0.000840064734401158\\
-20530.8265169088	-0.000719204212635943\\
-20519.1063362227	-0.000604886375442926\\
-20507.3861555366	-0.000498089330591291\\
-20495.6659748504	-0.000399728272984745\\
-20489.8058845074	-0.000353978104927964\\
-20483.9457941643	-0.000310647689780881\\
-20478.0857038213	-0.000269830358505824\\
-20472.2256134782	-0.000231614157814808\\
-20466.3655231351	-0.00019608166175015\\
-20460.5054327921	-0.000163309794756426\\
-20454.645342449	-0.000133369666660804\\
-20448.785252106	-0.000106326419876124\\
-20442.9251617629	-8.22390891847683e-05\\
-20437.0650714198	-6.11604743852655e-05\\
-20431.2049810768	-4.31370260855577e-05\\
-20425.3448907337	-2.82087448914954e-05\\
-20419.4848003907	-1.64090941985982e-05\\
-20413.6247100476	-7.76492678911273e-06\\
-20407.7646197046	-2.29642538632653e-06\\
-20401.9045293615	-1.70573014601601e-08\\
-20396.0444390184	-9.33543273517639e-07\\
-20390.1843486754	-5.04584057474679e-06\\
-20384.3242583323	-1.23471404209158e-05\\
-20378.4641679893	-2.28238796987978e-05\\
-20372.6040776462	-3.64557669931793e-05\\
-20366.7439873031	-5.32158228612711e-05\\
-20360.8838969601	-7.30704342765375e-05\\
-20355.023806617	-9.59794231308124e-05\\
-20349.163716274	-0.000121896128662621\\
-20343.3036259309	-0.000150767503627476\\
-20337.4435355878	-0.000182534224030026\\
-20331.5834452448	-0.000217130812177925\\
-20325.7233549017	-0.000254485772812664\\
-20319.8632645587	-0.000294521742025921\\
-20314.0031742156	-0.000337155648650876\\
-20302.2829935295	-0.000429857506229384\\
-20290.5628128434	-0.000531819517057388\\
-20278.8426321572	-0.000642190833221445\\
-20267.1224514711	-0.000760048725137236\\
-20255.402270785	-0.000884406333322104\\
-20243.6820900989	-0.00101422096927117\\
-20226.1018190697	-0.00121678043811988\\
-20208.5215480405	-0.00142532986583139\\
-20190.9412770114	-0.00163592296076743\\
-20173.3610059822	-0.00184456966670262\\
-20155.780734953	-0.00204731205105907\\
-20138.2004639238	-0.00224029967761424\\
-20114.7601025516	-0.00247613493676936\\
-20091.3197411793	-0.00268016006111007\\
-20067.8793798071	-0.00284546356189473\\
-20044.4390184349	-0.00296643163928709\\
-20020.9986570626	-0.00303894019690302\\
-19997.5582956904	-0.00306049685036157\\
-19974.1179343182	-0.00303032804499619\\
-19950.6775729459	-0.00294940824394311\\
-19927.2372115737	-0.0028204300975864\\
-19903.7968502014	-0.00264771649768994\\
-19880.3564888292	-0.00243707738301771\\
-19856.916127457	-0.00219561603434716\\
-19839.3358564278	-0.00199916885448461\\
-19821.7555853986	-0.00179370166096843\\
-19804.1753143694	-0.00158313571932502\\
-19786.5950433403	-0.00137149493905905\\
-19769.0147723111	-0.00116282942486231\\
-19757.294591625	-0.00102735134846444\\
-19745.5744109388	-0.000896132343363573\\
-19733.8542302527	-0.000770294123527442\\
-19722.1340495666	-0.000650913919879023\\
-19710.4138688805	-0.000539015319244037\\
-19698.6936881944	-0.000435559548498835\\
-19686.9735075082	-0.00034143727816274\\
-19681.1134171652	-0.000298133362220343\\
-19675.2533268221	-0.000257461015274727\\
-19669.3932364791	-0.000219508446775727\\
-19663.533146136	-0.000184358150436336\\
-19657.6730557929	-0.000152086726849342\\
-19651.8129654499	-0.000122764718492462\\
-19645.9528751068	-9.64564574642837e-05\\
-19640.0927847638	-7.32199263042771e-05\\
-19634.2326944207	-5.31066321693841e-05\\
-19628.3726040776	-3.61614946600077e-05\\
-19622.5125137346	-2.24227475303034e-05\\
-19616.6524233915	-1.19218545007031e-05\\
-19610.7923330485	-4.6834393550504e-06\\
-19604.9322427054	-7.25230481056351e-07\\
-19599.0721523623	-5.8019972881849e-08\\
-19593.2120620193	-2.68563739496502e-06\\
-19587.3519716762	-8.60493826722153e-06\\
-19581.4918813332	-1.78058073018194e-05\\
-19575.6317909901	-3.02711763954224e-05\\
-19569.7717006471	-4.59770573428382e-05\\
-19563.911610304	-6.48925892133559e-05\\
-19558.0515199609	-8.69801002897976e-05\\
-19552.1914296179	-0.000112195184460174\\
-19546.3313392748	-0.000140486791887828\\
-19540.4712489318	-0.000171797333797657\\
-19534.6111585887	-0.000206062801140336\\
-19528.7510682456	-0.000243212896917138\\
-19522.8909779026	-0.000283171181858921\\
-19517.0308875595	-0.000325855233192107\\
-19511.1707972165	-0.000371176816142205\\
-19499.4506165303	-0.000469351693087314\\
-19487.7304358442	-0.000576880978879248\\
-19476.0102551581	-0.000692869827713654\\
-19464.290074472	-0.000816350876680697\\
-19452.5698937859	-0.000946292361867014\\
-19440.8497130997	-0.00108160678768538\\
-19423.2694420706	-0.00129216150219575\\
-19405.6891710414	-0.00150827166609827\\
-19388.1089000122	-0.00172585220064794\\
-19370.528628983	-0.00194078384828432\\
-19352.9483579539	-0.00214899153281789\\
-19335.3680869247	-0.00234652200028608\\
-19311.9277255524	-0.00258683256160068\\
-19288.4873641802	-0.0027933871035165\\
-19265.047002808	-0.00295918251593026\\
-19241.6066414357	-0.00307857974996451\\
-19218.1662800635	-0.00314749735773166\\
-19194.7259186912	-0.0031635527994453\\
-19171.285557319	-0.00312614663227855\\
-19147.8451959468	-0.00303648661415781\\
-19124.4048345745	-0.00289755077932135\\
-19100.9644732023	-0.0027139906008825\\
-19077.5241118301	-0.00249197738222417\\
-19054.0837504578	-0.00223899694348069\\
-19036.5034794286	-0.00203405071237181\\
-19018.9232083995	-0.00182038724457605\\
-19001.3429373703	-0.00160208759142087\\
-18983.7626663411	-0.00138332807492845\\
-18966.1823953119	-0.00116830094451714\\
-18954.4622146258	-0.00102906680304969\\
-18942.7420339397	-0.000894520485818125\\
-18931.0218532536	-0.000765814873823112\\
-18919.3016725674	-0.000644054721529025\\
-18907.5814918813	-0.000530287249178608\\
-18895.8613111952	-0.00042549321124583\\
-18890.0012208522	-0.000376747467148975\\
-18884.1411305091	-0.000330578517176891\\
-18878.281040166	-0.000287086632907327\\
-18872.420949823	-0.000246366475928381\\
-18866.5608594799	-0.000208506894207686\\
-18860.7007691369	-0.000173590730640519\\
-18854.8406787938	-0.000141694644190633\\
-18848.9805884507	-0.000112888944024474\\
-18843.1204981077	-8.72374369815938e-05\\
-18837.2604077646	-6.47972887199579e-05\\
-18831.4003174216	-4.56188988320399e-05\\
-18825.5402270785	-2.97457902099721e-05\\
-18819.6801367354	-1.72145128995256e-05\\
-18813.8200463924	-8.054562652008e-06\\
-18807.9599560493	-2.28831436094514e-06\\
-18802.0998657063	6.90294731979261e-08\\
-18796.2397753632	-9.90524872404033e-07\\
-18790.3796850201	-5.46774119205874e-06\\
-18784.5195946771	-1.33561474888933e-05\\
-18778.659504334	-2.46420454223196e-05\\
-18772.799413991	-3.93045350541097e-05\\
-18766.9393236479	-5.73155548495061e-05\\
-18761.0792333048	-7.86399368576844e-05\\
-18755.2191429618	-0.000103235476965918\\
-18749.3590526187	-0.000131053020082639\\
-18743.4989622757	-0.000162036560084215\\
-18737.6388719326	-0.00019612335432281\\
-18731.7787815895	-0.000233244052445928\\
-18725.9186912465	-0.000273322839292687\\
-18720.0586009034	-0.000316277591548625\\
-18714.1985105604	-0.000362020047838695\\
-18702.4783298743	-0.000461485448615252\\
-18690.7581491881	-0.000570897465252642\\
-18679.037968502	-0.000689349321815146\\
-18667.3177878159	-0.000815856617485192\\
-18655.5976071298	-0.000949365567354957\\
-18643.8774264437	-0.00108876184037811\\
-18626.2971554145	-0.00130633382351794\\
-18608.7168843853	-0.00153042270630186\\
-18591.1366133561	-0.00175679955082239\\
-18573.5563423269	-0.00198118421492547\\
-18555.9760712978	-0.00219932674509548\\
-18538.3958002686	-0.00240708834446833\\
-18514.9554388963	-0.00266116777054495\\
-18491.5150775241	-0.00288121078107255\\
-18468.0747161519	-0.00305975541951707\\
-18444.6343547796	-0.00319072521218933\\
-18421.1939934074	-0.00326963812941574\\
-18397.7536320352	-0.00329376217527827\\
-18374.3132706629	-0.00326221221261659\\
-18350.8729092907	-0.00317598460315991\\
-18327.4325479184	-0.00303792833347716\\
-18303.9921865462	-0.00285265343988584\\
-18280.551825174	-0.0026263796659168\\
-18257.1114638017	-0.00236673031190815\\
-18239.5311927726	-0.0021553327840013\\
-18221.9509217434	-0.00193411008510074\\
-18204.3706507142	-0.0017072890074432\\
-18186.790379685	-0.00147921149032229\\
-18169.2101086558	-0.00125425218065833\\
-18157.4899279697	-0.00110815159609184\\
-18145.7697472836	-0.000966613188023733\\
-18134.0495665975	-0.00083085174866964\\
-18122.3293859114	-0.000702034797422572\\
-18110.6092052252	-0.000581272632318132\\
-18098.8890245391	-0.000469608852506454\\
-18087.168843853	-0.000368011432396361\\
-18081.3087535099	-0.000321267364520118\\
-18075.4486631669	-0.000277364423876758\\
-18069.5885728238	-0.000236398757207158\\
-18063.7284824808	-0.000198460357582423\\
-18057.8683921377	-0.000163632869696458\\
-18052.0083017947	-0.000131993408469772\\
-18046.1482114516	-0.000103612391332356\\
-18040.2881211085	-7.85533845819969e-05\\
-18034.4280307655	-5.68729641296078e-05\\
-18028.5679404224	-3.86205909431193e-05\\
-18022.7078500794	-2.38385014690801e-05\\
-18016.8477597363	-1.256161326216e-05\\
-18010.9876693932	-4.81744604422046e-06\\
-18005.1275790502	-6.26058363293904e-07\\
nan	nan\\
-17993.4073983641	-2.94428023271915e-06\\
-17987.547308021	-9.45635203853158e-06\\
-17981.6872176779	-1.95261122754273e-05\\
-17975.8271273349	-3.31359178562001e-05\\
-17969.9670369918	-5.0260617890078e-05\\
-17964.1069466488	-7.08676017363276e-05\\
-17958.2468563057	-9.49168628787064e-05\\
-17952.3867659626	-0.000122361078494211\\
-17946.5266756196	-0.00015314570456439\\
-17940.6665852765	-0.000187209086336214\\
-17934.8064949335	-0.000224482583897122\\
-17928.9464045904	-0.000264890712639448\\
-17923.0863142473	-0.00030835129829279\\
-17917.2262239043	-0.000354775646231425\\
-17911.3661335612	-0.000404068724687542\\
-17899.6459528751	-0.000510850454125667\\
-17887.925772189	-0.000627817291522315\\
-17876.2055915029	-0.000754002566785259\\
-17864.4854108167	-0.000888360196592567\\
-17852.7652301306	-0.00102977342971059\\
-17841.0450494445	-0.00117706420270445\\
-17823.4647784153	-0.00140631971993997\\
-17805.8845073862	-0.00164171214863288\\
-17788.304236357	-0.00187880394986281\\
-17770.7239653278	-0.00211311584269286\\
-17753.1436942986	-0.00234021200491103\\
-17735.5634232694	-0.00255578463012177\\
-17712.1230618972	-0.00281824737597132\\
-17688.682700525	-0.00304409270797898\\
-17665.2423391527	-0.00322565414337754\\
-17641.8019777805	-0.00335674158194467\\
-17618.3616164083	-0.00343285499842081\\
-17594.921255036	-0.0034513416214765\\
-17571.4808936638	-0.00341149112122121\\
-17548.0405322915	-0.00331456540429649\\
-17524.6001709193	-0.0031637618160903\\
-17501.1598095471	-0.00296411079482232\\
-17477.7194481748	-0.00272231124081693\\
-17454.2790868026	-0.00244650897752015\\
-17436.6988157734	-0.00222291492083514\\
-17419.1185447442	-0.00198968700024413\\
-17401.5382737151	-0.00175128573374146\\
-17383.9580026859	-0.00151228062114851\\
-17366.3777316567	-0.00127726338193065\\
-17354.6575509706	-0.00112504134473395\\
-17342.9373702845	-0.000977913898139755\\
-17331.2171895983	-0.00083714722528202\\
-17319.4970089122	-0.000703955630570854\\
-17307.7768282261	-0.000579491176827423\\
-17296.05664754	-0.00046483383308878\\
-17290.1965571969	-0.000411498312991594\\
-17284.3364668539	-0.000360982217294307\\
-17278.4763765108	-0.000313396253729843\\
-17272.6162861677	-0.000268845013054363\\
-17266.7561958247	-0.000227426742778135\\
-17260.8961054816	-0.0001892331341979\\
-17255.0360151386	-0.000154349123092558\\
-17249.1759247955	-0.000122852704579042\\
-17243.3158344524	-9.48147625138472e-05\\
-17237.4557441094	-7.02989137991503e-05\\
-17231.5956537663	-4.93613679520115e-05\\
-17225.7355634233	-3.20508022419686e-05\\
-17219.8754730802	-1.84082526626683e-05\\
-17214.0153827372	-8.467020996523e-06\\
-17208.1552923941	-2.25259817053585e-06\\
-17202.295202051	2.17395917606153e-07\\
-17196.435111708	-1.06674403792564e-06\\
-17190.5750213649	-6.10678190334502e-06\\
-17184.7149310219	-1.48965300430669e-05\\
-17178.8548406788	-2.74218560787596e-05\\
-17172.9947503357	-4.36607064672648e-05\\
-17167.1346599927	-6.35831468621092e-05\\
-17161.2745696496	-8.71514191847529e-05\\
-17155.4144793066	-0.0001143200153026\\
-17149.5543889635	-0.000145035767163624\\
-17143.6942986204	-0.000179237953221132\\
-17137.8342082774	-0.000216858420919931\\
-17131.9741179343	-0.000257821725022549\\
-17126.1140275913	-0.000302045281462762\\
-17120.2539372482	-0.000349439536447095\\
-17114.3938469051	-0.00039990815044365\\
-17102.673666219	-0.000509650373644626\\
-17090.9534855329	-0.000630373411576671\\
-17079.2333048468	-0.000761084283144257\\
-17067.5131241607	-0.000900703820765691\\
-17055.7929434745	-0.00104807566845909\\
-17044.0727627884	-0.00120197594788073\\
-17026.4924917592	-0.00144225164530541\\
-17008.9122207301	-0.00168981450016808\\
-16991.3319497009	-0.00194000642106733\\
-16973.7516786717	-0.00218810788894791\\
-16956.1714076425	-0.00242942768344873\\
-16938.5911366134	-0.00265939229299421\\
-16915.1507752411	-0.00294083786981587\\
-16891.7104138689	-0.0031848439498138\\
-16868.2700524966	-0.00338312646340058\\
-16844.8296911244	-0.00352892142605913\\
-16821.3893297522	-0.00361721888623165\\
-16797.9489683799	-0.00364493834982327\\
-16774.5086070077	-0.00361103955292185\\
-16751.0682456355	-0.00351656461909876\\
-16727.6278842632	-0.00336460994172886\\
-16704.187522891	-0.00316022850187203\\
-16680.7471615187	-0.00291026568369465\\
-16657.3068001465	-0.00262313390868933\\
-16639.7265291173	-0.00238918878177739\\
-16622.1462580881	-0.00214423714665986\\
-16604.565987059	-0.00189296510718047\\
-16586.9857160298	-0.00164019175297522\\
-16569.4054450006	-0.00139077773104547\\
-16557.6852643145	-0.00122874806979584\\
-16545.9650836284	-0.00107174565352221\\
-16534.2449029423	-0.000921124013042951\\
-16522.5247222561	-0.00077818507271979\\
-16510.80454157	-0.00064416803343282\\
-16499.0843608839	-0.000520238766584936\\
-16487.3641801978	-0.000407479809935973\\
-16481.5040898547	-0.000355602396256526\\
-16475.6439995117	-0.000306881050907357\\
-16469.7839091686	-0.000261423630404602\\
-16463.9238188255	-0.000219331177446408\\
-16458.0637284825	-0.000180697700851392\\
-16452.2036381394	-0.000145609970177582\\
-16446.3435477964	-0.000114147325468672\\
-16440.4834574533	-8.63815025671135e-05\\
-16434.6233671102	-6.2376474350329e-05\\
-16428.7632767672	-4.21883082618168e-05\\
-16422.9031864241	-2.58650404379933e-05\\
-16417.0430960811	-1.344656672264e-05\\
-16411.183005738	-4.96455080859517e-06\\
-16405.3229153949	-4.42349719875634e-07\\
-16399.4628250519	1.05043195003138e-07\\
-16393.6027347088	-3.32896238068003e-06\\
-16387.7426443658	-1.07425321233418e-05\\
-16381.8825540227	-2.21254032706356e-05\\
-16376.0224636796	-3.74588986587459e-05\\
-16370.1623733366	-5.6715958578948e-05\\
-16364.3022829935	-7.98611904064394e-05\\
-16358.4421926505	-0.000106850935909349\\
-16352.5821023074	-0.000137633356111667\\
-16346.7220119643	-0.00017214853354908\\
-16340.8619216213	-0.000210328591719765\\
-16335.0018312782	-0.000252097831477033\\
-16329.1417409352	-0.000297372884124726\\
-16323.2816505921	-0.000346062880858526\\
-16317.4215602491	-0.000398069638266874\\
-16311.561469906	-0.00045328785944793\\
-16299.8412892199	-0.00057290325122892\\
-16288.1211085338	-0.000703932999889502\\
-16276.4009278476	-0.000845302577438621\\
-16264.6807471615	-0.000995847901314616\\
-16252.9605664754	-0.00115432502617893\\
-16241.2403857893	-0.00131942052976436\\
-16223.6601147601	-0.0015764590329577\\
-16206.0798437309	-0.00184047531004613\\
-16188.4995727017	-0.00210650727740089\\
-16170.9193016726	-0.00236954064236121\\
-16153.3390306434	-0.00262460426284351\\
-16135.7587596142	-0.00286686495555605\\
-16112.318398242	-0.00316205242376054\\
-16088.8780368697	-0.00341634258813389\\
-16065.4376754975	-0.00362109170310465\\
-16041.9973141253	-0.00376930052668983\\
-16018.556952753	-0.00385585748052977\\
-15995.1165913808	-0.00387771890197957\\
-15971.6762300085	-0.00383402005160874\\
-15948.2358686363	-0.00372611285472477\\
-15924.7955072641	-0.00355752881660454\\
-15901.3551458918	-0.00333386807227135\\
-15877.9147845196	-0.00306261803341563\\
-15854.4744231474	-0.00275290748649402\\
-15836.8941521182	-0.002501641152423\\
-15819.313881089	-0.0022394075006136\\
-15801.7336100598	-0.00197122912742261\\
-15784.1533390306	-0.00170225724090561\\
-15766.5730680015	-0.00143767392686483\\
-15754.8528873153	-0.00126625449902387\\
-15743.1327066292	-0.00110054047939344\\
-15731.4125259431	-0.000941965042732009\\
-15719.692345257	-0.000791903830388467\\
-15707.9721645709	-0.000651663183166302\\
-15696.2519838848	-0.000522468936991754\\
-15690.3918935417	-0.000462372664789287\\
-15684.5318031986	-0.000405455877350871\\
-15678.6717128556	-0.000351844588788032\\
-15672.8116225125	-0.000301657943000936\\
-15666.9515321695	-0.000255007954396796\\
-15661.0914418264	-0.000211999263428327\\
-15655.2313514833	-0.00017272890747811\\
-15649.3712611403	-0.000137286107612723\\
-15643.5111707972	-0.000105752071654079\\
-15637.6510804542	-7.81998140127464e-05\\
-15631.7909901111	-5.46939926726159e-05\\
-15625.930899768	-3.52907636955655e-05\\
-15620.070809425	-2.00376535697323e-05\\
-15614.2107190819	-8.97344969084332e-06\\
-15608.3506287389	-2.1281092255186e-06\\
-15602.4905383958	4.773134301401e-07\\
-15596.6304480527	-1.16927957373461e-06\\
-15590.7703577097	-7.07099466777598e-06\\
-15584.9102673666	-1.72219309810085e-05\\
-15579.0501770236	-3.16071835082891e-05\\
-15573.1900866805	-5.0202865334378e-05\\
-15567.3299963374	-7.2976148895099e-05\\
-15561.4699059944	-9.98853261971183e-05\\
-15555.6098156513	-0.000130879887902337\\
-15549.7497253083	-0.000165900621115956\\
-15543.8896349652	-0.000204879725700828\\
-15538.0295446221	-0.000247740948886146\\
-15532.1694542791	-0.000294399737903101\\
-15526.309363936	-0.000344763410350772\\
-15520.449273593	-0.000398731341943496\\
-15514.5891832499	-0.000456195171261834\\
-15502.8690025638	-0.000581139735917738\\
-15491.1488218777	-0.000718584280846323\\
-15479.4286411916	-0.00086740798862672\\
-15467.7084605054	-0.00102639131353052\\
-15455.9882798193	-0.00119422610243507\\
-15444.2680991332	-0.00136952648722847\\
-15426.687828104	-0.00164328368143192\\
-15409.1075570748	-0.00192544384110782\\
-15391.5272860457	-0.00221071580829229\\
-15373.9470150165	-0.00249373244525841\\
-15356.3667439873	-0.00276915263604038\\
-15338.7864729581	-0.00303176311997486\\
-15315.3461115859	-0.00335341180118523\\
-15291.9057502136	-0.00363257930993113\\
-15268.4653888414	-0.00385977550368891\\
-15245.0250274692	-0.00402722956387104\\
-15221.5846660969	-0.00412916048897852\\
-15198.1443047247	-0.00416198145516016\\
-15174.7039433525	-0.0041244308448573\\
-15151.2635819802	-0.00401762519138437\\
-15127.823220608	-0.00384503190591027\\
-15104.3828592357	-0.00361236235425926\\
-15080.9424978635	-0.00332738854628988\\
-15057.5021364913	-0.00299968929659867\\
-15039.9218654621	-0.00273248936690219\\
-15022.3415944329	-0.00245256272851965\\
-15004.7613234037	-0.00216527240276336\\
-14987.1810523746	-0.00187614027613619\\
-14969.6007813454	-0.00159074244143019\\
-14957.8806006593	-0.00140528517846411\\
-14946.1604199731	-0.00122554832635953\\
-14934.440239287	-0.00105308939068388\\
-14922.7200586009	-0.000889407827800184\\
-14910.9998779148	-0.000735932211527269\\
-14899.2796972287	-0.000594007970482249\\
-14887.5595165425	-0.000464885801107435\\
-14881.6994261995	-0.000405487754519023\\
-14875.8393358564	-0.000349710856090986\\
-14869.9792455134	-0.00029768012719494\\
-14864.1191551703	-0.000249512801427311\\
-14858.2590648272	-0.000205318067391454\\
-14852.3989744842	-0.000165196828188351\\
-14846.5388841411	-0.000129241478160138\\
-14840.6787937981	-9.75356973902043e-05\\
-14834.818703455	-7.01542644141694e-05\\
-14828.958613112	-4.7162887539361e-05\\
-14823.0985227689	-2.86180551753425e-05\\
-14817.2384324258	-1.45669054970921e-05\\
-14811.3783420828	-5.04711574925861e-06\\
-14805.5182517397	-8.6811443922328e-08\\
-14799.6581613967	2.95504333086003e-07\\
-14793.7980710536	-3.90899866785218e-06\\
-14787.9379807105	-1.26994478629447e-05\\
-14782.0778903675	-2.60652583372505e-05\\
-14776.2178000244	-4.39861438923031e-05\\
-14770.3577096814	-6.64321486381091e-05\\
-14764.4976193383	-9.33636990873341e-05\\
-14758.6375289952	-0.000124731676664487\\
-14752.7774386522	-0.000160477510506018\\
-14746.9173483091	-0.000200533290366235\\
-14741.0572579661	-0.000244821899432677\\
-14735.197167623	-0.000293257166773564\\
-14729.3370772799	-0.00034574403913085\\
-14723.4769869369	-0.000402178771716434\\
-14717.6168965938	-0.000462449137631033\\
-14711.7568062508	-0.00052643465546198\\
-14700.0366255646	-0.00066502943812183\\
-14688.3164448785	-0.000816843928538207\\
-14676.5962641924	-0.000980644310672026\\
-14664.8760835063	-0.00115509232217937\\
-14653.1559028202	-0.00133875634153709\\
-14641.4357221341	-0.00153012329196357\\
-14623.8554511049	-0.00182814070969855\\
-14606.2751800757	-0.00213435837440913\\
-14588.6949090465	-0.00244304124377386\\
-14571.1146380173	-0.00274838688313236\\
-14553.5343669882	-0.00304463577926302\\
-14535.954095959	-0.00332618123636432\\
-14512.5137345867	-0.00366951270763007\\
-14489.0733732145	-0.00396561708263934\\
-14465.6330118423	-0.00420441347361459\\
-14442.19265047	-0.00437771359764474\\
-14418.7522890978	-0.00447950828613353\\
-14395.3119277256	-0.00450618167305567\\
-14371.8715663533	-0.00445664546272161\\
-14348.4312049811	-0.00433238833977242\\
-14324.9908436088	-0.00413743842996907\\
-14301.5504822366	-0.00387823964719625\\
-14278.1101208644	-0.00356344567879536\\
-14254.6697594921	-0.00320363816260212\\
-14237.0894884629	-0.00291150977669876\\
-14219.5092174338	-0.0026064623691151\\
-14201.9289464046	-0.00229434862998998\\
-14184.3486753754	-0.00198117875579589\\
-14166.7684043462	-0.001673006530483\\
-14155.0482236601	-0.00147329542822453\\
-14143.328042974	-0.00128019931801013\\
-14131.6078622879	-0.00109539790317468\\
-14119.8876816018	-0.000920504995117792\\
-14108.1675009156	-0.000757054672722113\\
-14096.4473202295	-0.000606488081124548\\
-14090.5872298865	-0.000536458049918619\\
-14084.7271395434	-0.000470140983008642\\
-14078.8670492003	-0.000407685483172795\\
-14073.0069588573	-0.00034923216998541\\
-14067.1468685142	-0.00029491337167422\\
-14061.2867781712	-0.000244852834125261\\
-14055.4266878281	-0.000199165447665733\\
-14049.566597485	-0.000157956992237147\\
-14043.706507142	-0.000121323901525574\\
-14037.8464167989	-8.93530465502477e-05\\
-14031.9863264559	-6.21215392153655e-05\\
-14026.1262361128	-3.96965562601339e-05\\
-14020.2661457697	-2.21351840003748e-05\\
-14014.4060554267	-9.48428422582442e-06\\
-14008.5459650836	-1.78038155813986e-06\\
-14002.6858747406	9.50427459075276e-07\\
-13996.8257843975	-1.30745667827178e-06\\
-13990.9656940544	-8.5590895997881e-06\\
-13985.1056037114	-2.07989584786331e-05\\
-13979.2455133683	-3.80109797838483e-05\\
-13973.3854230253	-6.01685194012163e-05\\
-13967.5253326822	-8.72344351123164e-05\\
-13961.6652423392	-0.000119161141364189\\
-13955.8051519961	-0.000155890696233735\\
-13949.945061653	-0.000197354910416347\\
-13944.08497131	-0.000243475478057306\\
-13938.2248809669	-0.000294164129152046\\
-13932.3647906239	-0.000349322803258595\\
-13926.5047002808	-0.000408843844137917\\
-13920.6446099377	-0.00047261021497196\\
-13914.7845195947	-0.000540495733723048\\
-13903.0643389086	-0.000688075310060081\\
-13891.3441582224	-0.000850400377821322\\
-13879.6239775363	-0.00102616068145417\\
-13867.9037968502	-0.00121392866245647\\
-13856.1836161641	-0.00141217123863335\\
-13844.463435478	-0.00161926250841621\\
-13826.8831644488	-0.00194274240345282\\
-13809.3028934196	-0.00227626601918929\\
-13791.7226223904	-0.00261360398102513\\
-13774.1423513612	-0.00294842926925367\\
-13756.5620803321	-0.00327443739756012\\
-13738.9818093029	-0.00358546665249678\\
-13715.5414479307	-0.00396672453832547\\
-13692.1010865584	-0.00429800319315724\\
-13668.6607251862	-0.00456802352855954\\
-13645.2203638139	-0.00476752102761162\\
-13621.7800024417	-0.00488957061478979\\
-13598.3396410695	-0.00492983415279872\\
-13574.8992796972	-0.00488672176510969\\
-13551.458918325	-0.00476146104900936\\
-13528.0185569528	-0.00455807132229245\\
-13504.5781955805	-0.00428324324341816\\
-13481.1378342083	-0.00394612734397704\\
-13457.697472836	-0.0035580381126893\\
-13440.1172018069	-0.00324134948453528\\
-13422.5369307777	-0.0029093868254581\\
-13404.9566597485	-0.00256852132462072\\
-13387.3763887193	-0.00222532187206752\\
-13369.7961176901	-0.00188643071881885\\
-13358.075937004	-0.00166615676518477\\
-13346.3557563179	-0.00145264326805179\\
-13334.6355756318	-0.0012477519367211\\
-13322.9153949457	-0.00105327689021096\\
-13311.1952142596	-0.000870929261528749\\
-13299.4750335734	-0.000702322459742393\\
-13287.7548528873	-0.000548958216667969\\
-13281.8947625443	-0.000478427577543718\\
-13276.0346722012	-0.000412213538406629\\
-13270.1745818581	-0.000350466762573949\\
-13264.3144915151	-0.000293328674839658\\
-13258.454401172	-0.000240931148450913\\
-13252.594310829	-0.000193396211866665\\
-13246.7342204859	-0.000150835775911191\\
-13240.8741301428	-0.000113351381929817\\
-13235.0140397998	-8.10339715372089e-05\\
-13229.1539494567	-5.39636784412149e-05\\
-13223.2938591137	-3.22096428508621e-05\\
-13217.4337687706	-1.58298488625258e-05\\
-13211.5736784275	-4.87098522725504e-06\\
-13205.7135880845	6.3167019038465e-07\\
-13199.8534977414	6.54341970378372e-07\\
-13193.9934073984	-4.8151755103121e-06\\
-13188.1333170553	-1.57774751126702e-05\\
-13182.2732267122	-3.22215184816521e-05\\
-13176.4131363692	-5.41246421823077e-05\\
-13170.5530460261	-8.14525884456714e-05\\
-13164.6929556831	-0.00011415956049917\\
-13158.83286534	-0.000152188302400557\\
-13152.9727749969	-0.000195470203237305\\
-13147.1126846539	-0.000243925425517103\\
-13141.2525943108	-0.000297463057510905\\
-13135.3925039678	-0.00035598128924929\\
-13129.5324136247	-0.000419367611854463\\
-13123.6723232817	-0.000487499039817409\\
-13117.8122329386	-0.000560242355756582\\
-13106.0920522525	-0.000718982245030837\\
-13094.3718715664	-0.000894327574196526\\
-13082.6516908802	-0.00108487418502644\\
-13070.9315101941	-0.00128908492892611\\
-13059.211329508	-0.00150530228981414\\
-13047.4911488219	-0.00173176206430656\\
-13035.7709681358	-0.0019666079865307\\
-13018.1906971066	-0.00233035573420049\\
-13000.6104260774	-0.00270184819879727\\
-12983.0301550482	-0.0030741463658034\\
-12965.449884019	-0.00344026461380296\\
-12947.8696129899	-0.00379330365670004\\
-12924.4292516176	-0.00423217700209131\\
-12900.9888902454	-0.00462113957348873\\
-12877.5485288732	-0.00494697168780536\\
-12854.1081675009	-0.00519851312115898\\
-12830.6678061287	-0.0053670539438863\\
-12807.2274447564	-0.00544664484648214\\
-12783.7870833842	-0.0054343159556759\\
-12760.346722012	-0.0053301961663468\\
-12736.9063606397	-0.00513752832805438\\
-12713.4659992675	-0.00486257911584817\\
-12690.0256378953	-0.00451444596939408\\
-12666.585276523	-0.00410476697843897\\
-12649.0050054938	-0.00376535139099728\\
-12631.4247344647	-0.00340552026183173\\
-12613.8444634355	-0.00303217118216235\\
-12596.2641924063	-0.00265249363169663\\
-12578.6839213771	-0.00227383334547552\\
-12561.1036503479	-0.00190355304595082\\
-12549.3834696618	-0.00166493937274596\\
-12537.6632889757	-0.00143535901862137\\
-12525.9431082896	-0.0012168282831699\\
-12514.2229276035	-0.00101127650213027\\
-12502.5027469173	-0.000820529415734817\\
-12490.7825662312	-0.000646293359927287\\
-12484.9224758882	-0.000565863872680999\\
-12479.0623855451	-0.00049014041669363\\
-12473.2022952021	-0.000419296266918004\\
-12467.342204859	-0.000353494652429067\\
-12461.4821145159	-0.000292888394196339\\
-12455.6220241729	-0.000237619564309568\\
-12449.7619338298	-0.000187819167386013\\
-12443.9018434868	-0.000143606844864621\\
-12438.0417531437	-0.000105090602865757\\
-12432.1816628006	-7.23665642196277e-05\\
-12426.3215724576	-4.55187452331639e-05\\
-12420.4614821145	-2.4618857708653e-05\\
-12414.6013917715	-9.72613667174347e-06\\
-12408.7413014284	-8.87194216363841e-07\\
-12402.8812110853	1.86410017836111e-06\\
-12397.0211207423	-1.49328743399849e-06\\
-12391.1610303992	-1.09674895595291e-05\\
-12385.3009400562	-2.65536985479246e-05\\
-12379.4408497131	-4.82341551947391e-05\\
-12373.58075937	-7.59781647374057e-05\\
-12367.720669027	-0.000109742140246979\\
-12361.8605786839	-0.00014946967337851\\
-12356.0004883409	-0.000195091632371791\\
-12350.1403979978	-0.000246526287144587\\
-12344.2803076547	-0.000303679461249289\\
-12338.4202173117	-0.000366444710434796\\
-12332.5601269686	-0.000434703527455033\\
-12326.7000366256	-0.000508325572768237\\
-12320.8399462825	-0.000587168930638313\\
-12314.9798559394	-0.000671080390199808\\
-12303.2596752533	-0.00085344015153012\\
-12291.5394945672	-0.00105396546207741\\
-12279.8193138811	-0.00127105816443956\\
-12268.099133195	-0.0015029743429901\\
-12256.3789525089	-0.00174783860663117\\
-12244.6587718227	-0.00200365953334869\\
-12227.0785007936	-0.00240333770227891\\
-12209.4982297644	-0.00281555983069255\\
-12191.9179587352	-0.00323266480331906\\
-12174.337687706	-0.00364685960373361\\
-12156.7574166768	-0.00405036721391871\\
-12139.1771456477	-0.00443557497323858\\
-12115.7367842754	-0.00490815999641268\\
-12092.2964229032	-0.00531928227074619\\
-12068.8560615309	-0.00565491675940761\\
-12045.4157001587	-0.00590350227901801\\
-12021.9753387865	-0.00605635053260512\\
-11998.5349774142	-0.00610796074592958\\
-11975.094616042	-0.00605622859673263\\
-11951.6542546698	-0.00590254162310263\\
-11928.2138932975	-0.00565175707812776\\
-11904.7735319253	-0.00531206214531751\\
-11881.333170553	-0.00489472040428386\\
-11857.8928091808	-0.00441371231139647\\
-11840.3125381516	-0.00402087609760015\\
-11822.7322671225	-0.00360884650759666\\
-11805.1519960933	-0.00318554936262997\\
-11787.5717250641	-0.00275916968489288\\
-11769.9914540349	-0.00233799705551539\\
-11758.2712733488	-0.00206418253330767\\
-11746.5510926627	-0.00179874230552348\\
-11734.8309119766	-0.0015440093137091\\
-11723.1107312904	-0.00130223446394367\\
-11711.3905506043	-0.0010755672505179\\
-11699.6703699182	-0.000866037166831683\\
-11693.8102795751	-0.000768295589181372\\
-11687.9501892321	-0.000675536063791807\\
-11682.090098889	-0.000587970851146087\\
-11676.230008546	-0.000505801608369831\\
-11670.3699182029	-0.00042921894251929\\
-11664.5098278598	-0.000358401986397821\\
-11658.6497375168	-0.000293517997887542\\
-11652.7896471737	-0.000234721983653205\\
-11646.9295568307	-0.000182156348062746\\
-11641.0694664876	-0.00013595056812842\\
-11635.2093761446	-9.62208952007329e-05\\
-11629.3492858015	-6.3070084082781e-05\\
-11623.4891954584	-3.65871502018753e-05\\
-11617.6291051154	-1.68471553991068e-05\\
-11611.7690147723	-3.91102284497095e-06\\
-11605.9089244293	2.17461847321223e-06\\
-11600.0488340862	1.37755931010916e-06\\
-11594.1887437431	-6.31989455759075e-06\\
-11588.3286534001	-2.09208575862028e-05\\
-11582.468563057	-4.24138304485396e-05\\
-11576.608472714	-7.07726968791103e-05\\
-11570.7483823709	-0.000105956751441114\\
-11564.8882920278	-0.000147910758218454\\
-11559.0282016848	-0.000196565040367412\\
-11553.1681113417	-0.000251835600389362\\
-11547.3080209987	-0.000313624270939049\\
-11541.4479306556	-0.000381818895900072\\
-11535.5878403125	-0.000456293541417036\\
-11529.7277499695	-0.000536908736441398\\
-11523.8676596264	-0.000623511742444622\\
-11518.0075692834	-0.000715936851649411\\
-11506.2873885972	-0.000917527687418584\\
-11494.5672079111	-0.00114010927397309\\
-11482.847027225	-0.00138192579859688\\
-11471.1268465389	-0.00164105178308009\\
-11459.4066658528	-0.00191540775218794\\
-11447.6864851666	-0.00220277726615378\\
-11435.9663044805	-0.00250082517838853\\
-11418.3860334513	-0.00296257132498431\\
-11400.8057624222	-0.00343431947684863\\
-11383.225491393	-0.00390730200535276\\
-11365.6452203638	-0.00437267762882282\\
-11348.0649493346	-0.0048216995929385\\
-11324.6245879624	-0.00538034952354886\\
-11301.1842265902	-0.00587602711381145\\
-11277.7438652179	-0.00629185935513033\\
-11254.3035038457	-0.0066135519109268\\
-11230.8631424734	-0.00682989408722613\\
-11207.4227811012	-0.00693316337742769\\
-11183.982419729	-0.00691941509582266\\
-11160.5420583567	-0.00678864638302746\\
-11137.1016969845	-0.00654482802896784\\
-11113.6613356123	-0.00619580197600453\\
-11090.22097424	-0.00575304688685145\\
-11066.7806128678	-0.00523131864439613\\
-11049.2003418386	-0.00479866978748966\\
-11031.6200708094	-0.00433969099928981\\
-11014.0397997802	-0.00386319812857683\\
-10996.4595287511	-0.00337839706934762\\
-10978.8792577219	-0.00289471060746283\\
-10961.2989866927	-0.00242160006130094\\
-10949.5788060066	-0.00211668208161304\\
-10937.8586253205	-0.00182330269912859\\
-10926.1384446344	-0.00154406501829236\\
-10914.4182639482	-0.00128146330477403\\
-10902.6980832621	-0.00103786139158772\\
-10890.977902576	-0.000815472099262405\\
-10885.1178122329	-0.000712877601540279\\
-10879.2577218899	-0.000616337848460132\\
-10873.3976315468	-0.000526078559795566\\
-10867.5375412038	-0.000442312634040603\\
-10861.6774508607	-0.000365239670544749\\
-10855.8173605176	-0.000295045518808255\\
-10849.9572701746	-0.000231901855888585\\
-10844.0971798315	-0.000175965792927513\\
-10838.2370894885	-0.000127379511635701\\
-10832.3769991454	-8.62699316219065e-05\\
-10826.5169088023	-5.27484092917896e-05\\
-10820.6568184593	-2.69104690456881e-05\\
-10814.7967281162	-8.83556739862741e-06\\
-10808.9366377732	1.41310940428786e-06\\
-10803.0765474301	3.78881377659286e-06\\
-10797.216457087	-1.73862924079821e-06\\
-10791.356366744	-1.51827227904997e-05\\
-10785.4962764009	-3.65402340709376e-05\\
-10779.6361860579	-6.57911654813149e-05\\
-10773.7760957148	-0.000102898761174201\\
-10767.9160053717	-0.000147809549082245\\
-10762.0559150287	-0.000200453418395766\\
-10756.1958246856	-0.000260743732411699\\
-10750.3357343426	-0.000328577476596296\\
-10744.4756439995	-0.000403835441610912\\
-10738.6155536565	-0.000486382440981392\\
-10732.7554633134	-0.000576067563078175\\
-10726.8953729703	-0.000672724456844635\\
-10721.0352826273	-0.000776171650899172\\
-10715.1751922842	-0.000886212905200961\\
-10703.4550115981	-0.00112522112495846\\
-10691.734830912	-0.00138789918018199\\
-10680.0146502259	-0.00167218857319558\\
-10668.2944695397	-0.00197583878746574\\
-10656.5742888536	-0.00229642552848497\\
-10644.8541081675	-0.00263137051536534\\
-10627.2738371383	-0.00315474838405956\\
-10609.6935661091	-0.0036947155029711\\
-10592.11329508	-0.00424130234704384\\
-10574.5330240508	-0.00478434851646065\\
-10556.9527530216	-0.00531369529496827\\
-10539.3724819924	-0.00581937945383468\\
-10515.9321206202	-0.00644034369319348\\
-10492.491759248	-0.00698124630269803\\
-10469.0513978757	-0.00742358494436153\\
-10445.6110365035	-0.00775204125828384\\
-10422.1706751312	-0.00795502908235959\\
-10398.730313759	-0.0080251212004464\\
-10375.2899523868	-0.00795933907754341\\
-10351.8495910145	-0.00775929453583125\\
-10328.4092296423	-0.00743117722572026\\
-10304.9688682701	-0.00698558690704567\\
-10281.5285068978	-0.00643721478891272\\
-10258.0881455256	-0.00580438332047388\\
-10240.5078744964	-0.00528709101054839\\
-10222.9276034672	-0.00474417623187421\\
-10205.347332438	-0.00418611565386137\\
-10187.7670614089	-0.00362375093245474\\
-10170.1867903797	-0.00306808447021394\\
-10158.4666096936	-0.00270677930417743\\
-10146.7464290074	-0.00235651582489034\\
-10135.0262483213	-0.0020204066127214\\
-10123.3060676352	-0.00170145931781868\\
-10111.5858869491	-0.00140255066733593\\
-10099.865706263	-0.00112640145748773\\
-10094.0056159199	-0.000997663367469788\\
-10088.1455255769	-0.000875552746390671\\
-10082.2854352338	-0.000760355466990062\\
-10076.4253448907	-0.000652343454739857\\
-10070.5652545477	-0.000551774078525254\\
-10064.7051642046	-0.000458889570673991\\
-10058.8450738616	-0.000373916477591258\\
-10052.9849835185	-0.000297065142248737\\
-10047.1248931754	-0.000228529219733049\\
-10041.2648028324	-0.00016848522688702\\
-10035.4047124893	-0.00011709212714911\\
-10029.5446221463	-7.44909514947628e-05\\
-10023.6845318032	-4.08044564115603e-05\\
-10017.8244414601	-1.61368196951621e-05\\
-10011.9643511171	-5.73374805836168e-07\\
-10006.104260774	5.81961555461356e-06\\
-10000.244170431	2.99514607758894e-06\\
-9994.3840800879	-9.07438479967228e-06\\
-9988.52398974484	-3.03970692294141e-05\\
-9982.66389940178	-6.09614252983426e-05\\
-9976.80380905872	-0.000100736373622287\\
-9970.94371871567	-0.000149671255333666\\
-9965.08362837261	-0.000207695891532479\\
-9959.22353802955	-0.000274720684174514\\
-9953.36344768649	-0.000350636758274933\\
-9947.50335734343	-0.0004353161452437\\
-9941.64326700037	-0.000528612007057587\\
-9935.78317665731	-0.000630358900904697\\
-9929.92308631425	-0.000740373083796323\\
-9924.06299597119	-0.000858452856706486\\
-9918.20290562813	-0.000984378947466937\\
-9906.48272494201	-0.00125880769181348\\
-9894.76254425589	-0.00156157060253946\\
-9883.04236356977	-0.0018903285430994\\
-9871.32218288365	-0.00224251054200134\\
-9859.60200219753	-0.00261533443557092\\
-9847.88182151142	-0.00300582940147332\\
-9836.1616408253	-0.00341086020214953\\
-9818.58136979612	-0.00403846816739487\\
-9801.00109876694	-0.00467989879912989\\
-9783.42082773776	-0.00532331275057724\\
-9765.84055670858	-0.00595674513375714\\
-9748.2602856794	-0.00656833288766489\\
-9724.81992430716	-0.00732994253437018\\
-9701.37956293493	-0.00800656018087187\\
-9677.93920156269	-0.00857510017185733\\
-9654.49884019045	-0.00901591725888633\\
-9631.05847881822	-0.0093135087368781\\
-9607.61811744598	-0.00945708349465288\\
-9584.17775607374	-0.00944097734400312\\
-9560.7373947015	-0.00926489897379619\\
-9537.29703332927	-0.00893399643703047\\
-9513.85667195703	-0.00845874003762402\\
-9490.41631058479	-0.00785462364526479\\
-9466.97594921255	-0.00714169262402612\\
-9449.39567818337	-0.00654987514561934\\
-9431.81540715419	-0.00592157794692808\\
-9414.23513612501	-0.00526890509852343\\
-9396.65486509584	-0.00460452743791083\\
-9379.07459406666	-0.00394144539697287\\
-9361.49432303748	-0.00329274361819302\\
-9349.77414235136	-0.00287465926948797\\
-9338.05396166524	-0.00247245309246311\\
-9326.33378097912	-0.00208974762412753\\
-9314.613600293	-0.00173002032881896\\
-9302.89341960689	-0.0013965733222765\\
-9297.03332926383	-0.00124068173402229\\
-9291.17323892077	-0.00109250441344431\\
-9285.31314857771	-0.000952392169648613\\
-9279.45305823465	-0.000820679716238142\\
-9273.59296789159	-0.000697684918003724\\
-9267.73287754853	-0.000583708071202371\\
-9261.87278720547	-0.000479031218977657\\
-9256.01269686241	-0.000383917503452488\\
-9250.15260651935	-0.000298610556014676\\
-9244.29251617629	-0.000223333927157371\\
-9238.43242583323	-0.000158290557183506\\
-9232.57233549017	-0.000103662289029873\\
-9226.71224514711	-5.96094243329286e-05\\
-9220.85215480405	-2.62703238093769e-05\\
-9214.992064461	-3.76105290718477e-06\\
-9209.13197411794	7.82492638399595e-06\\
-9203.27188377488	8.41701679615916e-06\\
-9197.41179343182	-2.03230056979739e-06\\
-9191.55170308876	-2.35473007952853e-05\\
-9185.6916127457	-5.61288983802229e-05\\
-9179.83152240264	-9.97545795444341e-05\\
-9173.97143205958	-0.000154378383879877\\
-9168.11134171652	-0.000219930935522611\\
-9162.25125137346	-0.00029631952393277\\
-9156.3911610304	-0.000383428234241947\\
-9150.53107068734	-0.000481118127020635\\
-9144.67098034428	-0.000589227467265248\\
-9138.81089000122	-0.00070757200220806\\
-9132.95079965816	-0.000835945287546409\\
-9127.0907093151	-0.000974119061537158\\
-9121.23061897204	-0.00112184366623683\\
-9115.37052862898	-0.00127884851523458\\
-9103.65034794286	-0.00161951508258513\\
-9091.93016725675	-0.00199355611571542\\
-9080.20998657063	-0.00239811150803011\\
-9068.48980588451	-0.00283004623514711\\
-9056.76962519839	-0.00328597535268875\\
-9045.04944451227	-0.00376229125416531\\
-9027.46917348309	-0.0045066304573585\\
-9009.88890245391	-0.00527477136284887\\
-8992.30863142473	-0.00605266598490444\\
-8974.72836039555	-0.00682595985488937\\
-8957.14808936638	-0.00758026334523564\\
-8939.5678183372	-0.00830142605598395\\
-8916.12745696496	-0.00918795747381797\\
-8892.68709559272	-0.00996134920026928\\
-8869.24673422049	-0.0105950345094054\\
-8845.80637284825	-0.01106688918645\\
-8822.36601147601	-0.011360036198386\\
-8798.92565010377	-0.0114634822592087\\
-8775.48528873154	-0.0113725627245209\\
-8752.0449273593	-0.0110891774419085\\
-8728.60456598706	-0.0106218070311142\\
-8705.16420461482	-0.00998530635664211\\
-8681.72384324258	-0.00920047943980881\\
-8658.28348187034	-0.00829344749103283\\
-8640.70321084117	-0.00755125312717801\\
-8623.12293981199	-0.00677173350075279\\
-8605.54266878281	-0.00597000655580826\\
-8587.96239775363	-0.00516176310263483\\
-8570.38212672445	-0.00436297278406238\\
-8558.66194603833	-0.00384358253648326\\
-8546.94176535222	-0.00334014404606398\\
-8535.2215846661	-0.00285720894460144\\
-8523.50140397998	-0.00239918488619662\\
-8511.78122329386	-0.00197029725704231\\
-8500.06104260774	-0.0015745521810667\\
-8494.20095226468	-0.00139029131535057\\
-8488.34086192162	-0.00121570114241605\\
-8482.48077157856	-0.00105120532135851\\
-8476.6206812355	-0.000897207534121875\\
-8470.76059089244	-0.000754090566093032\\
-8464.90050054939	-0.000622215427785242\\
-8459.04041020633	-0.0005019205197564\\
-8453.18031986327	-0.000393520842531137\\
-8447.32022952021	-0.000297307253461387\\
-8441.46013917715	-0.000213545772104161\\
-8435.60004883409	-0.000142476935878959\\
-8429.73995849103	-8.4315207436601e-05\\
-8423.87986814797	-3.92484351895344e-05\\
-8418.01977780491	-7.43736831905056e-06\\
-8412.15968746185	1.09847725478301e-05\\
-8406.29959711879	1.59126679020703e-05\\
-8400.43950677573	7.26921418331051e-06\\
-8394.57941643267	-1.49942199766626e-05\\
-8388.71932608961	-5.08975968296321e-05\\
-8382.85923574655	-0.000100432072515783\\
-8376.99914540349	-0.000163559920488515\\
-8371.13905506043	-0.000240214515715671\\
-8365.27896471737	-0.000330300379891579\\
-8359.41887437431	-0.000433693287773499\\
-8353.55878403125	-0.00055024043461665\\
-8347.69869368819	-0.000679760664552227\\
-8341.83860334514	-0.000822044759626161\\
-8335.97851300208	-0.000976855789136341\\
-8330.11842265902	-0.00114392951863701\\
-8324.25833231596	-0.00132297487807245\\
-8318.3982419729	-0.00151367448811043\\
-8306.67806128678	-0.00192863895519843\\
-8294.95788060066	-0.00238578128066021\\
-8283.23769991454	-0.00288168198627955\\
-8271.51751922842	-0.0034125701953642\\
-8259.7973385423	-0.00397435322599193\\
-8248.07715785618	-0.00456264910827734\\
-8230.496886827	-0.00548460284504284\\
-8212.91661579783	-0.00643921019716877\\
-8195.33634476865	-0.00740916301652995\\
-8177.75607373947	-0.00837668572914441\\
-8160.17580271029	-0.0093238699436654\\
-8142.59553168111	-0.0102330150762208\\
-8119.15517030888	-0.0113565110135151\\
-8095.71480893664	-0.0123439177454911\\
-8072.2744475644	-0.0131612757384649\\
-8048.83408619216	-0.0137799505293344\\
-8025.39372481993	-0.0141776817014191\\
-8001.95336344769	-0.0143394308680123\\
-7978.51300207545	-0.0142579971199111\\
-7955.07264070321	-0.0139343757865282\\
-7931.63227933097	-0.0133778446711558\\
-7908.19191795873	-0.0126057708848127\\
-7884.7515565865	-0.0116431407279807\\
-7861.31119521426	-0.0105218244416379\\
-7843.73092418508	-0.00959923387898797\\
-7826.1506531559	-0.00862627023100858\\
-7808.57038212672	-0.00762184287113324\\
-7790.99011109754	-0.00660568626821062\\
-7773.40984006837	-0.00559799263855317\\
-7761.68965938225	-0.00494091672700519\\
-7749.96947869613	-0.00430256280080076\\
-7738.24929801001	-0.0036887645673474\\
-7726.52911732389	-0.00310519358549772\\
-7714.80893663777	-0.00255730994362976\\
-7703.08875595166	-0.00205031434601259\\
-7697.2286656086	-0.00181369247194271\\
-7691.36857526554	-0.00158910202191079\\
-7685.50848492248	-0.00137709997138264\\
-7679.64839457942	-0.00117821886042376\\
-7673.78830423636	-0.000992965571739724\\
-7667.9282138933	-0.000821820157997939\\
-7662.06812355024	-0.000665234721270361\\
-7656.20803320718	-0.000523632347051589\\
-7650.34794286412	-0.000397406095374902\\
-7644.48785252106	-0.000286918051360987\\
-7638.627762178	-0.000192498437387347\\
-7632.76767183494	-0.000114444789048414\\
-7626.90758149188	-5.30211968507886e-05\\
-7621.04749114882	-8.45761550854606e-06\\
-7615.18740080576	1.90507574305546e-05\\
-7609.3273104627	2.93440320654991e-05\\
-7603.46722011964	2.22981042770951e-05\\
-7597.60712977658	-2.17493976741762e-06\\
-7591.74703943353	-4.41264938878022e-05\\
-7585.88694909047	-0.000103571177926833\\
-7580.02685874741	-0.0001804866603871\\
-7574.16676840435	-0.000274813558354109\\
-7568.30667806129	-0.000386455415196423\\
-7562.44658771823	-0.000515278756450276\\
-7556.58649737517	-0.000661113224018044\\
-7550.72640703211	-0.000823751788772461\\
-7544.86631668905	-0.00100295104139083\\
-7539.00622634599	-0.00119843156113038\\
-7533.14613600293	-0.001409878362134\\
-7527.28604565987	-0.00163694141651033\\
-7521.42595531681	-0.00187923625356278\\
-7509.70577463069	-0.00240781529858957\\
-7497.98559394457	-0.00299187833827533\\
-7486.26541325845	-0.00362718786104831\\
-7474.54523257233	-0.00430903784531584\\
-7462.82505188622	-0.00503228982371356\\
-7451.1048712001	-0.00579141289927255\\
-7439.38469051398	-0.00658052740899902\\
-7421.8044194848	-0.00780684088478494\\
-7404.22414845562	-0.00906479804055633\\
-7386.64387742644	-0.0103316585680163\\
-7369.06360639727	-0.0115842350325273\\
-7351.48333536809	-0.0127993302170293\\
-7328.04297399585	-0.014321925398716\\
-7304.60261262361	-0.0156861747399205\\
-7281.16225125137	-0.0168452844263202\\
-7257.72188987914	-0.0177586807862829\\
-7234.2815285069	-0.0183934984765587\\
-7210.84116713466	-0.0187258335602542\\
-7187.40080576242	-0.0187417150808529\\
-7163.96044439018	-0.0184377574007755\\
-7140.52008301795	-0.017821465654814\\
-7117.07972164571	-0.016911177844771\\
-7093.63936027347	-0.0157356390353644\\
-7070.19899890123	-0.0143332153963796\\
-7052.61872787205	-0.0131604138175955\\
-7035.03845684288	-0.0119086551565541\\
-7017.4581858137	-0.0106022338939304\\
-6999.87791478452	-0.00926684942112058\\
-6982.29764375534	-0.00792913526700807\\
-6970.57746306922	-0.00704939798005758\\
-6958.8572823831	-0.00618868403251014\\
-6947.13710169699	-0.00535493635989402\\
-6935.41692101087	-0.00455594769407947\\
-6923.69674032475	-0.00379929311741169\\
-6911.97655963863	-0.00309226380890344\\
-6900.25637895251	-0.00244180255162248\\
-6894.39628860945	-0.00213984394061383\\
-6888.53619826639	-0.00185444155975373\\
-6882.67610792333	-0.00158633805257091\\
-6876.81601758027	-0.00133624317001102\\
-6870.95592723721	-0.00110483211031015\\
-6865.09583689415	-0.000892743923840465\\
-6859.23574655109	-0.000700579986113411\\
-6853.37565620803	-0.000528902542717993\\
-6847.51556586497	-0.000378233329560688\\
-6841.65547552191	-0.000249052271612713\\
-6835.79538517886	-0.000141796263310171\\
-6829.9352948358	-5.68580335656158e-05\\
-6824.07520449274	5.41490181981804e-06\\
-6818.21511414968	4.47211977396864e-05\\
-6812.35502380662	6.08065448811398e-05\\
-6806.49493346356	5.34644400932771e-05\\
-6800.6348431205	2.25368599307568e-05\\
-6794.77475277744	-3.20851645051866e-05\\
-6788.91466243438	-0.000110461073133609\\
-6783.05457209132	-0.000212600404349719\\
-6777.19448174826	-0.000338462525666669\\
-6771.3343914052	-0.00048795646860541\\
-6765.47430106214	-0.000660940868717434\\
-6759.61421071908	-0.000857224011372839\\
-6753.75412037602	-0.00107656398377321\\
-6747.89403003296	-0.00131866893334073\\
-6742.0339396899	-0.00158319743259294\\
-6736.17384934684	-0.00186975895010211\\
-6730.31375900378	-0.00217791442732509\\
-6724.45366866072	-0.00250717696029412\\
-6718.59357831766	-0.0028570125857194\\
-6706.87339763155	-0.00361603739947661\\
-6695.15321694543	-0.00444981228665629\\
-6683.43303625931	-0.00535247504330896\\
-6671.71285557319	-0.00631752045353124\\
-6659.99267488707	-0.007337848763703\\
-6648.27249420095	-0.00840581967497942\\
-6630.69222317177	-0.0100792236035522\\
-6613.1119521426	-0.0118123505099848\\
-6595.53168111342	-0.013574512043281\\
-6577.95141008424	-0.0153339935881119\\
-6560.37113905506	-0.0170586438051054\\
-6542.79086802588	-0.0187164819759805\\
-6519.35050665365	-0.0207692729987549\\
-6495.91014528141	-0.0225781695175015\\
-6472.46978390917	-0.0240802241843144\\
-6449.02942253693	-0.0252216564092075\\
-6425.58906116469	-0.0259598893565532\\
-6402.14869979245	-0.0262652540818828\\
-6378.70833842022	-0.0261222960899504\\
-6355.26797704798	-0.0255306314425216\\
-6331.82761567574	-0.024505313412591\\
-6308.3872543035	-0.0230766861454268\\
-6284.94689293126	-0.021289718372482\\
-6267.36662190209	-0.0197490117059137\\
-6249.78635087291	-0.0180684446572701\\
-6232.20607984373	-0.0162803607417812\\
-6214.62580881455	-0.0144198008258561\\
-6197.04553778538	-0.0125238849122458\\
-6179.4652667562	-0.0106311431676606\\
-6167.74508607008	-0.00939044587901222\\
-6156.02490538396	-0.0081802436172531\\
-6144.30472469784	-0.00701207879060207\\
-6132.58454401172	-0.00589729026377876\\
-6120.8643633256	-0.00484691496030564\\
-6109.14418263948	-0.00387159075676286\\
-6103.28409229642	-0.00341526460706999\\
-6097.42400195336	-0.00298146149550286\\
-6091.5639116103	-0.0025713571050529\\
-6085.70382126725	-0.00218608493357806\\
-6079.84373092419	-0.00182673363421009\\
-6073.98364058113	-0.00149434443234112\\
-6068.12355023807	-0.00118990862541044\\
-6062.26345989501	-0.000914365170761012\\
-6056.40336955195	-0.000668598367619746\\
-6050.54327920889	-0.000453435638175552\\
-6044.68318886583	-0.000269645413337205\\
-6038.82309852277	-0.000117935128014562\\
-6032.96300817971	1.05066909580283e-06\\
-6027.10291783665	8.67320865952666e-05\\
-6021.24282749359	0.000138595537882602\\
-6015.38273715053	0.000156195260852918\\
-6009.52264680747	0.00013915470579904\\
-6003.66255646441	8.71677877848544e-05\\
nan	nan\\
-5991.94237577829	-0.000122510615126645\\
-5986.08228543523	-0.000280452640003633\\
-5980.22219509217	-0.000473840603818683\\
-5974.36210474911	-0.000702614454714258\\
-5968.50201440605	-0.000966639184324193\\
-5962.64192406299	-0.00126570460641393\\
-5956.78183371994	-0.00159952529117195\\
-5950.92174337688	-0.00196774065625192\\
-5945.06165303382	-0.00236991521549063\\
-5939.20156269076	-0.0028055389858606\\
-5933.3414723477	-0.00327402805282989\\
-5927.48138200464	-0.0037747252939739\\
-5921.62129166158	-0.00430690126063523\\
-5909.90111097546	-0.00546241633302454\\
-5898.18093028934	-0.00673333451729139\\
-5886.46074960322	-0.00811134241370701\\
-5874.7405689171	-0.00958710721534504\\
-5863.02038823099	-0.0111503412637656\\
-5851.30020754487	-0.012789875514294\\
-5833.71993651569	-0.0153658619123142\\
-5816.13966548651	-0.0180431423637368\\
-5798.55939445733	-0.0207755936994355\\
-5780.97912342815	-0.0235150732477431\\
-5763.39885239898	-0.0262122929199842\\
-5739.95849102674	-0.0296573449622676\\
-5716.5181296545	-0.0328243405983894\\
-5693.07776828226	-0.0356045307153503\\
-5669.63740691002	-0.037899288185729\\
-5646.19704553779	-0.0396237834005381\\
-5622.75668416555	-0.0407103125403666\\
-5599.31632279331	-0.0411111576362818\\
-5575.87596142107	-0.0408008691843478\\
-5552.43560004883	-0.0397778776368927\\
-5528.9952386766	-0.0380653591668352\\
-5505.55487730436	-0.0357113031012466\\
-5482.11451593212	-0.0327877528011607\\
-5464.53424490294	-0.0302774310715828\\
-5446.95397387376	-0.0275469352793039\\
-5429.37370284459	-0.0246493262921954\\
-5411.79343181541	-0.0216424565689595\\
-5394.21316078623	-0.0185880071358758\\
-5382.49298010011	-0.016557021200611\\
-5370.77279941399	-0.0145530461144421\\
-5359.05261872787	-0.0125957939561263\\
-5347.33243804175	-0.0107049802931535\\
-5335.61225735563	-0.00890015948248691\\
-5323.89207666952	-0.00720055824365951\\
-5318.03198632646	-0.00639609265514323\\
-5312.1718959834	-0.00562490882429211\\
-5306.31180564034	-0.0048892437984047\\
-5300.45171529728	-0.00419128309699356\\
-5294.59162495422	-0.00353315567387322\\
-5288.73153461116	-0.00291692894699858\\
-5282.8714442681	-0.00234460390672503\\
-5277.01135392504	-0.00181811031274662\\
-5271.15126358198	-0.00133930199060966\\
-5265.29117323892	-0.000909952237710022\\
-5259.43108289586	-0.000531749349427142\\
-5253.5709925528	-0.00020629227510786\\
-5247.71090220974	6.4913585937501e-05\\
-5241.85081186668	0.000280460438792809\\
-5235.99072152362	0.000439041988111794\\
-5230.13063118056	0.000539457164811521\\
-5224.2705408375	0.000580613675216832\\
-5218.41045049444	0.000561531363910543\\
-5212.55036015138	0.000481345381755638\\
-5206.69026980832	0.000339309150907825\\
-5200.83017946527	0.000134797118939733\\
-5194.97008912221	-0.000132692705418172\\
-5189.10999877915	-0.000463536442155097\\
-5183.24990843609	-0.000857982262230763\\
-5177.38981809303	-0.00131614854860384\\
-5171.52972774997	-0.00183802230208667\\
-5165.66963740691	-0.00242345779756414\\
-5159.80954706385	-0.00307217549572846\\
-5153.94945672079	-0.00378376121516422\\
-5148.08936637773	-0.00455766556858167\\
-5142.22927603467	-0.0053932036677192\\
-5136.36918569161	-0.00628955509950849\\
-5130.50909534855	-0.00724576417652113\\
-5124.64900500549	-0.00826074046401502\\
-5112.92882431937	-0.0104619643084524\\
-5101.20864363326	-0.012881845815389\\
-5089.48846294714	-0.0155069288851972\\
-5077.76828226102	-0.0183217482260002\\
-5066.0481015749	-0.0213089132012892\\
-5054.32792088878	-0.0244492091518728\\
-5036.7476498596	-0.0294005769326532\\
-5019.16737883043	-0.0345719762746803\\
-5001.58710780125	-0.0398795686526454\\
-4984.00683677207	-0.0452340866088032\\
-4966.42656574289	-0.0505423156890993\\
-4942.98620437065	-0.0573826692757273\\
-4919.54584299841	-0.0637433610438042\\
-4896.10548162617	-0.0694025438860658\\
-4872.66512025394	-0.0741516930284429\\
-4849.2247588817	-0.0778033315588932\\
-4825.78439750946	-0.0801984987493339\\
-4802.34403613722	-0.0812137232874689\\
-4778.90367476498	-0.080767268888778\\
-4755.46331339275	-0.078824432068685\\
-4732.02295202051	-0.0754016910161567\\
-4708.58259064827	-0.0705695301285212\\
-4691.0023196191	-0.0660940975749839\\
-4673.42204858992	-0.0609690344120997\\
-4655.84177756074	-0.0552849484671149\\
-4638.26150653156	-0.0491488528904369\\
-4620.68123550238	-0.0426831816400012\\
-4608.96105481626	-0.038256546367968\\
-4597.24087413014	-0.0337875244621813\\
-4585.52069344402	-0.0293220225675828\\
-4573.80051275791	-0.0249079484284059\\
-4562.08033207179	-0.020594933383191\\
-4556.22024172873	-0.0184922313870945\\
-4550.36015138567	-0.0164340353348914\\
-4544.50006104261	-0.0144268986882795\\
-4538.63997069955	-0.0124774238011881\\
-4532.77988035649	-0.0105922512223704\\
-4526.91979001343	-0.0087780487888698\\
-4521.05969967037	-0.00704150052832416\\
-4515.19960932731	-0.00538929538585851\\
-4509.33951898425	-0.00382811579186758\\
-4503.47942864119	-0.00236462608924002\\
-4497.61933829813	-0.0010054608368811\\
-4491.75924795507	0.000242786992088656\\
-4485.89915761201	0.00137357789895909\\
-4480.03906726895	0.00238043781621532\\
-4474.17897692589	0.00325697011347182\\
-4468.31888658283	0.00399686763332092\\
-4462.45879623977	0.00459392473771451\\
-4456.59870589671	0.00504204934483239\\
-4450.73861555366	0.0053352749374476\\
-4444.8785252106	0.00546777252190704\\
-4439.01843486754	0.00543386251869322\\
-4433.15834452448	0.00522802656357129\\
-4427.29825418142	0.00484491920031939\\
-4421.43816383836	0.00427937944383579\\
-4415.5780734953	0.00352644219481869\\
-4409.71798315224	0.00258134948502271\\
-4403.85789280918	0.00143956153411115\\
-4397.99780246612	9.67675980635635e-05\\
-4392.13771212306	-0.00145110341038938\\
-4386.27762178	-0.00320787254820711\\
-4380.41753143694	-0.00517710073465855\\
-4374.55744109388	-0.00736207892878549\\
-4368.69735075082	-0.00976581860991315\\
-4362.83726040776	-0.0123910425795004\\
-4356.97717006471	-0.0152401761018386\\
-4351.11707972165	-0.0183153384010472\\
-4345.25698937859	-0.0216183345308175\\
-4339.39689903553	-0.0251506476336813\\
-4333.53680869247	-0.0289134316054568\\
-4321.81662800635	-0.0371333404505316\\
-4310.09644732023	-0.0462804555158956\\
-4298.37626663411	-0.0563515072221691\\
-4286.65608594799	-0.067337407025305\\
-4274.93590526187	-0.0792231244319007\\
-4263.21572457576	-0.0919875972916452\\
-4245.63545354658	-0.112720914432465\\
-4228.0551825174	-0.13525153999808\\
-4210.47491148822	-0.159431415498824\\
-4192.89464045904	-0.185082409880712\\
-4169.4542790868	-0.221210061458708\\
-4146.01391771456	-0.259015019513751\\
-4122.57355634233	-0.297851811219159\\
-4099.13319497009	-0.337006267618582\\
-4075.69283359785	-0.375707748956619\\
-4046.39238188255	-0.42220602590329\\
-4017.09193016726	-0.465080906397207\\
-3987.79147845196	-0.50267429382598\\
-3958.49102673666	-0.533366552555596\\
-3929.19057502137	-0.555630751400453\\
-3899.89012330607	-0.568085780573532\\
-3870.58967159077	-0.569546432736429\\
-3841.28921987547	-0.559068617544459\\
-3817.84885850324	-0.541634026479852\\
-3794.408497131	-0.515929512003144\\
-3770.96813575876	-0.481891659052772\\
-3753.38786472958	-0.45093497498114\\
-3735.8075937004	-0.415415037318013\\
-3718.22732267122	-0.375460943051072\\
-3706.5071419851	-0.346445221558062\\
-3694.78696129899	-0.315596884907732\\
-3683.06678061287	-0.282986257974658\\
-3671.34659992675	-0.248691735036614\\
-3659.62641924063	-0.212799498214104\\
-3653.76632889757	-0.194283220345604\\
-3647.90623855451	-0.175403196746636\\
-3642.04614821145	-0.156172241528107\\
-3636.18605786839	-0.136603588761247\\
-3630.32596752533	-0.116710879352488\\
-3624.46587718227	-0.096508147431388\\
-3618.60578683921	-0.0760098062416358\\
-3612.74569649615	-0.0552306335631772\\
-3606.88560615309	-0.0341857566787052\\
-3601.02551581004	-0.012890636917385\\
-3595.16542546698	0.00863894622100714\\
-3589.30533512392	0.0303869113242634\\
-3583.44524478086	0.0523368916980917\\
-3577.5851544378	0.0744722523392822\\
-3571.72506409474	0.0967761072337489\\
-3565.86497375168	0.119231336952416\\
-3560.00488340862	0.141820606527005\\
-3554.14479306556	0.164526383572827\\
-3548.2847027225	0.187330956644857\\
-3542.42461237944	0.210216453800139\\
-3536.56452203638	0.233164861330508\\
-3524.84434135026	0.279177757358764\\
-3513.12416066415	0.325223420805006\\
-3501.40397997803	0.371154580007798\\
-3489.68379929191	0.416823532001606\\
-3472.10352826273	0.484512828700304\\
-3454.52325723355	0.550786572764059\\
-3436.94298620437	0.615159537146415\\
-3413.50262483213	0.697219797801138\\
-3390.06226345989	0.773994468866484\\
-3366.62190208766	0.844486334270942\\
-3337.32145037236	0.922389205523877\\
-3308.02099865706	0.987508530355332\\
-3278.72054694176	1.03857193114877\\
-3249.42009522647	1.07461235805407\\
-3220.11964351117	1.09498419145606\\
-3190.81919179587	1.09936990030888\\
-3155.65864973752	1.08356456399332\\
-3126.35819802222	1.05324975583795\\
-3097.05774630692	1.00805215848877\\
-3067.75729459163	0.948907070119819\\
-3038.45684287633	0.876979340563779\\
-3015.01648150409	0.811140388467105\\
-2991.57612013185	0.738753798811626\\
-2968.13575875961	0.660636423108755\\
-2950.55548773043	0.598799696405538\\
-2932.97521670126	0.534587711230066\\
-2915.39494567208	0.468373785426526\\
-2903.67476498596	0.423304094248923\\
-2891.95458429984	0.377622132326263\\
-2880.23440361372	0.331438200414991\\
-2868.5142229276	0.284861713767814\\
-2856.79404224148	0.238000904769999\\
-2850.93395189842	0.214497309746409\\
-2845.07386155537	0.190962544938027\\
-2839.21377121231	0.167409686196695\\
-2833.35368086925	0.143851686950385\\
-2827.49359052619	0.120301371266129\\
-2821.63350018313	0.096771427256034\\
-2815.77340984007	0.0732744008047981\\
-2809.91331949701	0.0498226896368607\\
-2804.05322915395	0.0264285377076926\\
-2798.19313881089	0.00310402993171366\\
-2792.33304846783	-0.0201389127631252\\
-2786.47295812477	-0.0432885380521901\\
-2780.61286778171	-0.0663332665046747\\
-2774.75277743865	-0.0892616955546188\\
-2768.89268709559	-0.112062603149251\\
-2763.03259675254	-0.134724951087915\\
-2757.17250640948	-0.15723788804906\\
-2751.31241606642	-0.179590752309485\\
-2745.45232572336	-0.201773074171616\\
-2739.5922353803	-0.223774578095359\\
-2727.87205469418	-0.267195011551427\\
-2716.15187400806	-0.309773794772448\\
-2704.43169332194	-0.351435866303824\\
-2692.71151263582	-0.392109382374847\\
-2675.13124160664	-0.451116888696373\\
-2657.55097057746	-0.507528458285575\\
-2639.97069954828	-0.561142139615882\\
-2616.53033817605	-0.627955264149058\\
-2593.08997680381	-0.689066371958402\\
-2569.64961543157	-0.74412334775921\\
-2540.34916371627	-0.803966438612835\\
-2511.04871200098	-0.853364159186347\\
-2481.74826028568	-0.891905295019938\\
-2446.58771822732	-0.923380677160696\\
-2411.42717616897	-0.93838352068375\\
-2376.26663411061	-0.936710277289482\\
-2341.10609205225	-0.918340945436338\\
-2305.9455499939	-0.883448279689694\\
-2276.6450982786	-0.842016394828914\\
-2247.3446465633	-0.789707646561303\\
-2218.044194848	-0.726975632300107\\
-2194.60383347576	-0.669668258444039\\
-2171.16347210353	-0.606410453249348\\
-2147.72311073129	-0.537610415757817\\
-2130.14283970211	-0.482649768977804\\
-2112.56256867293	-0.425046521524711\\
-2100.84238798681	-0.385290474439117\\
-2089.1222073007	-0.344536139459104\\
-2077.40202661458	-0.302859875357715\\
-2065.68184592846	-0.260341198076615\\
-2059.8217555854	-0.238791641223069\\
-2053.96166524234	-0.217062676501985\\
-2048.10157489928	-0.195165082091855\\
-2042.24148455622	-0.173109810811487\\
-2036.38139421316	-0.150907985369714\\
-2030.5213038701	-0.128570893332231\\
-2024.66121352704	-0.106109981819477\\
-2018.80112318398	-0.0835368519419716\\
-2012.94103284092	-0.0608632529603804\\
-2007.08094249787	-0.0381010761922624\\
-2001.22085215481	-0.0152623486548387\\
-1995.36076181175	0.00764077354533821\\
-1989.50067146869	0.0305960120719594\\
-1983.64058112563	0.0535909734650745\\
-1977.78049078257	0.0766131564837794\\
-1971.92040043951	0.0996499596837526\\
-1966.06031009645	0.122688689220735\\
-1960.20021975339	0.145716566871446\\
-1954.34012941033	0.168720738271243\\
-1948.48003906727	0.191688281346292\\
-1942.61994872421	0.214606214952346\\
-1936.75985838115	0.237461507690922\\
-1925.03967769503	0.282931847848651\\
-1913.31949700891	0.327994391559696\\
-1901.59931632279	0.372544016865965\\
-1889.87913563667	0.416475699368519\\
-1872.2988646075	0.480985937887762\\
-1854.71859357832	0.543520996748328\\
-1837.13832254914	0.603737485207334\\
-1813.6979611769	0.67984259998767\\
-1790.25759980466	0.750471012333545\\
-1766.81723843242	0.814907356671754\\
-1737.51678671713	0.885746226967516\\
-1708.21633500183	0.944740483360737\\
-1678.91588328653	0.9909303533315\\
-1649.61543157124	1.0235611680713\\
-1614.45488951288	1.04408629989649\\
-1579.29434745452	1.04385550840326\\
-1544.13380539617	1.02288928178193\\
-1514.83335368087	0.989933789663105\\
-1485.53290196557	0.943476070394558\\
-1456.23245025027	0.884279955332726\\
-1426.93199853498	0.813305426779727\\
-1403.49163716274	0.748802068986713\\
-1380.0512757905	0.678133481277436\\
-1356.61091441826	0.602001650428078\\
-1339.03064338909	0.541763999266042\\
-1321.45037235991	0.479197912944355\\
-1303.87010133073	0.414638066700273\\
-1292.14992064461	0.370658497079378\\
-1280.42973995849	0.326045931480171\\
-1268.70955927237	0.280902656332863\\
-1256.98937858626	0.235331102097145\\
-1251.1292882432	0.212416702411962\\
-1245.26919790014	0.189433572519851\\
-1239.40910755708	0.16639443766229\\
-1233.54901721402	0.143311982992198\\
-1227.68892687096	0.120198845958366\\
-1221.8288365279	0.0970676088605646\\
-1215.96874618484	0.0739307915839076\\
-1210.10865584178	0.0508008445204743\\
-1204.24856549872	0.0276901416734756\\
-1198.38847515566	0.00461097396102351\\
-1192.5283848126	-0.0184244572892841\\
-1186.66829446954	-0.0414040466446187\\
-1180.80820412648	-0.0643157906693696\\
-1174.94811378342	-0.0871477938224896\\
-1169.08802344036	-0.109888274150354\\
-1163.2279330973	-0.132525568771261\\
-1157.36784275424	-0.155048139153195\\
-1151.50775241118	-0.177444576173161\\
-1145.64766206812	-0.199703604986076\\
-1139.78757172506	-0.221814089656128\\
-1133.927481382	-0.243765037614798\\
-1122.20730069589	-0.287145095045285\\
-1110.48712000977	-0.329758859409109\\
-1098.76693932365	-0.371523955943757\\
-1087.04675863753	-0.412360659128184\\
-1069.46648760835	-0.471707347331137\\
-1051.88621657917	-0.528544775894659\\
-1034.30594554999	-0.58263991915285\\
-1010.86558417775	-0.650126723400459\\
-987.425222805519	-0.711879422613491\\
-963.98486143328	-0.767474173442195\\
-934.684409717982	-0.827743097312675\\
-905.383958002687	-0.877188857761063\\
-876.083506287388	-0.915311652346762\\
-840.92296422903	-0.945585980803626\\
-805.762422170676	-0.958582884712484\\
-770.601880112317	-0.954115110802406\\
-735.441338053963	-0.932239374418993\\
-700.280795995604	-0.893258152974448\\
-670.980344280306	-0.848097540576334\\
-641.679892565011	-0.79189033012327\\
-612.379440849712	-0.725214044692156\\
-588.939079477474	-0.664803222484261\\
-565.498718105238	-0.598556933139792\\
-547.918447076059	-0.545322583025903\\
-530.33817604688	-0.489284402511263\\
-512.757905017701	-0.4306740808352\\
-501.037724331585	-0.390292825658711\\
-489.317543645466	-0.348954505945881\\
-477.597362959346	-0.30673863218497\\
-465.877182273227	-0.263727393015123\\
-460.017091930167	-0.24194989970827\\
-454.157001587108	-0.220005540588357\\
-448.296911244048	-0.197905289303377\\
-442.436820900988	-0.175660263084189\\
-436.576730557928	-0.153281718031439\\
-430.716640214869	-0.130781044220187\\
-424.856549871809	-0.108169760600926\\
-418.996459528749	-0.0854595096805506\\
-413.136369185693	-0.0626620520201019\\
-407.276278842633	-0.0397892605221325\\
-401.416188499574	-0.0168531145238133\\
-395.556098156514	0.0061343063085423\\
-389.696007813454	0.0291608282749697\\
-383.835917470395	0.0522141901281179\\
-377.975827127335	0.075282049745088\\
-372.115736784275	0.0983519909824909\\
-366.255646441216	0.121411530709818\\
-360.395556098156	0.144448126009068\\
-354.535465755096	0.167449181552656\\
-348.675375412036	0.190402057137492\\
-342.815285068977	0.213294075374845\\
-336.955194725917	0.236112529531225\\
-325.235014039801	0.281477819998943\\
-313.514833353682	0.326395994598654\\
-301.794652667562	0.370765174506946\\
-290.074471981443	0.414483809178651\\
-272.494200952264	0.478621412341787\\
-254.913929923085	0.540731761689403\\
-237.333658893909	0.600484956266981\\
-213.89329752167	0.675935994324994\\
-190.452936149431	0.745898213831991\\
-167.012574777196	0.809687034745988\\
-137.712123061898	0.879785753679549\\
-108.411671346599	0.938161383224601\\
-79.1112196313043	0.983892347470839\\
-49.8107679160057	1.01625076655098\\
-14.6502258576475	1.03671504386981\\
20.5103162007072	1.03671504386981\\
55.6708582590654	1.01625076655098\\
84.9713099743603	0.983892347470839\\
114.271761689659	0.938161383224601\\
143.572213404957	0.879785753679549\\
172.872665120252	0.809687034745988\\
196.313026492491	0.745898213831991\\
219.75338786473	0.675935994324994\\
243.193749236969	0.600484956266981\\
260.774020266144	0.540731761689403\\
278.354291295323	0.478621412341787\\
295.934562324503	0.414483809178651\\
307.654743010622	0.370765174506946\\
319.374923696741	0.326395994598654\\
331.095104382861	0.281477819998943\\
342.815285068977	0.236112529531225\\
348.675375412036	0.213294075374845\\
354.535465755096	0.190402057137492\\
360.395556098156	0.167449181552656\\
366.255646441216	0.144448126009068\\
372.115736784275	0.121411530709818\\
377.975827127335	0.0983519909824909\\
383.835917470395	0.075282049745088\\
389.696007813454	0.0522141901281179\\
395.556098156514	0.0291608282749697\\
401.416188499574	0.0061343063085423\\
407.276278842633	-0.0168531145238133\\
413.136369185693	-0.0397892605221325\\
418.996459528749	-0.0626620520201019\\
424.856549871809	-0.0854595096805506\\
430.716640214869	-0.108169760600926\\
436.576730557928	-0.130781044220187\\
442.436820900988	-0.153281718031439\\
448.296911244048	-0.175660263084189\\
454.157001587108	-0.197905289303377\\
460.017091930167	-0.220005540588357\\
465.877182273227	-0.24194989970827\\
477.597362959346	-0.285327194918679\\
489.317543645466	-0.327951188018543\\
501.037724331585	-0.369738393487153\\
512.757905017701	-0.410607948229538\\
530.33817604688	-0.470021541141903\\
547.918447076059	-0.526943600588813\\
565.498718105238	-0.581137353262501\\
588.939079477474	-0.648771293214524\\
612.379440849712	-0.710682120815374\\
635.819802221951	-0.766439567974595\\
665.12025393725	-0.826910699068304\\
694.420705652545	-0.876551286936023\\
723.721157367843	-0.914858358479798\\
758.881699426198	-0.945340126926624\\
794.042241484556	-0.95853426742479\\
829.202783542914	-0.954261524511677\\
864.363325601269	-0.932587269953778\\
899.523867659627	-0.89382068239106\\
928.824319374922	-0.848850867708293\\
958.124771090221	-0.792843515319611\\
987.425222805519	-0.726370878342324\\
1010.86558417775	-0.666119878518875\\
1034.30594554999	-0.600024143061596\\
1051.88621657917	-0.546892261632748\\
1069.46648760835	-0.490943779711676\\
1087.04675863753	-0.432406696698731\\
1098.76693932365	-0.39206323341799\\
1110.48712000977	-0.350752537250848\\
1122.20730069589	-0.308552973193088\\
1133.927481382	-0.265545603869687\\
1139.78757172506	-0.243765037614798\\
1145.64766206812	-0.221814089656128\\
1151.50775241118	-0.199703604986076\\
1157.36784275424	-0.177444576173161\\
1163.2279330973	-0.155048139153195\\
1169.08802344036	-0.132525568771261\\
1174.94811378342	-0.109888274150354\\
1180.80820412648	-0.0871477938224896\\
1186.66829446954	-0.0643157906693696\\
1192.5283848126	-0.0414040466446187\\
1198.38847515566	-0.0184244572892841\\
1204.24856549872	0.00461097396102351\\
1210.10865584178	0.0276901416734756\\
1215.96874618484	0.0508008445204743\\
1221.8288365279	0.0739307915839076\\
1227.68892687096	0.0970676088605646\\
1233.54901721402	0.120198845958366\\
1239.40910755708	0.143311982992198\\
1245.26919790014	0.16639443766229\\
1251.1292882432	0.189433572519851\\
1256.98937858626	0.212416702411962\\
1262.84946892931	0.235331102097145\\
1274.56964961543	0.280902656332863\\
1286.28983030155	0.326045931480171\\
1298.01001098767	0.370658497079378\\
1309.73019167379	0.414638066700273\\
1327.31046270297	0.479197912944355\\
1344.89073373214	0.541763999266042\\
1362.47100476132	0.602001650428078\\
1385.91136613356	0.678133481277436\\
1409.3517275058	0.748802068986713\\
1432.79208887804	0.813305426779727\\
1462.09254059333	0.884279955332726\\
1491.39299230863	0.943476070394558\\
1520.69344402393	0.989933789663105\\
1549.99389573923	1.02288928178193\\
1585.15443779758	1.04385550840326\\
1620.31497985594	1.04408629989649\\
1655.4755219143	1.0235611680713\\
1684.77597362959	0.9909303533315\\
1714.07642534489	0.944740483360737\\
1743.37687706019	0.885746226967516\\
1772.67732877548	0.814907356671754\\
1796.11769014772	0.750471012333545\\
1819.55805151996	0.67984259998767\\
1842.9984128922	0.603737485207334\\
1860.57868392138	0.543520996748328\\
1878.15895495056	0.480985937887762\\
1895.73922597973	0.416475699368519\\
1907.45940666585	0.372544016865965\\
1919.17958735197	0.327994391559696\\
1930.89976803809	0.282931847848651\\
1942.61994872421	0.237461507690922\\
1948.48003906727	0.214606214952346\\
1954.34012941033	0.191688281346292\\
1960.20021975339	0.168720738271243\\
1966.06031009645	0.145716566871446\\
1971.92040043951	0.122688689220735\\
1977.78049078257	0.0996499596837526\\
1983.64058112563	0.0766131564837794\\
1989.50067146869	0.0535909734650745\\
1995.36076181175	0.0305960120719594\\
2001.22085215481	0.00764077354533821\\
2007.08094249787	-0.0152623486548387\\
2012.94103284092	-0.0381010761922624\\
2018.80112318398	-0.0608632529603804\\
2024.66121352704	-0.0835368519419716\\
2030.5213038701	-0.106109981819477\\
2036.38139421316	-0.128570893332231\\
2042.24148455622	-0.150907985369714\\
2048.10157489928	-0.173109810811487\\
2053.96166524234	-0.195165082091855\\
2059.8217555854	-0.217062676501985\\
2065.68184592846	-0.238791641223069\\
2077.40202661458	-0.281700748023349\\
2089.1222073007	-0.323808351663822\\
2100.84238798681	-0.365033395610847\\
2112.56256867293	-0.405297930566587\\
2130.14283970211	-0.46373116616562\\
2147.72311073129	-0.519597485589205\\
2165.30338176047	-0.572675563010941\\
2188.74374313271	-0.638758927273433\\
2212.18410450494	-0.699089135234759\\
2235.62446587718	-0.753285171235446\\
2264.92491759248	-0.811910015576784\\
2294.22536930778	-0.859915808117668\\
2329.38591136613	-0.902941341497078\\
2364.54645342449	-0.92960648866904\\
2399.70699548285	-0.939637041934716\\
2434.8675375412	-0.932958121884436\\
2470.02807959956	-0.909682291547193\\
2505.18862165792	-0.870103296113017\\
2534.48907337321	-0.825004112213072\\
2563.78952508851	-0.769283741039823\\
2593.08997680381	-0.703410949366457\\
2616.53033817605	-0.643782141835267\\
2639.97069954828	-0.578358991479302\\
2657.55097057746	-0.525720076236175\\
2675.13124160664	-0.470219484359756\\
2692.71151263582	-0.41205386256751\\
2704.43169332194	-0.371900558051928\\
2716.15187400806	-0.330723985293277\\
2727.87205469418	-0.2885943760248\\
2739.5922353803	-0.24558518454382\\
2745.45232572336	-0.223774578095359\\
2751.31241606642	-0.201773074171616\\
2757.17250640948	-0.179590752309485\\
2763.03259675254	-0.15723788804906\\
2768.89268709559	-0.134724951087915\\
2774.75277743865	-0.112062603149251\\
2780.61286778171	-0.0892616955546188\\
2786.47295812477	-0.0663332665046747\\
2792.33304846783	-0.0432885380521901\\
2798.19313881089	-0.0201389127631252\\
2804.05322915395	0.00310402993171366\\
2809.91331949701	0.0264285377076926\\
2815.77340984007	0.0498226896368607\\
2821.63350018313	0.0732744008047981\\
2827.49359052619	0.096771427256034\\
2833.35368086925	0.120301371266129\\
2839.21377121231	0.143851686950385\\
2845.07386155537	0.167409686196695\\
2850.93395189842	0.190962544938027\\
2856.79404224148	0.214497309746409\\
2862.65413258454	0.238000904769999\\
2874.37431327066	0.284861713767814\\
2886.09449395678	0.331438200414991\\
2897.8146746429	0.377622132326263\\
2909.53485532902	0.423304094248923\\
2927.1151263582	0.490644395715934\\
2944.69539738738	0.556232637070602\\
2962.27566841655	0.619694194736366\\
2985.71602978879	0.70035884027241\\
3009.15639116103	0.775715496874331\\
3032.59675253327	0.84492759849302\\
3061.89720424857	0.921597776224929\\
3091.19765596386	0.986007350486606\\
3120.49810767916	1.03690982273399\\
3149.79855939446	1.07327386384239\\
3179.09901110976	1.09431803521871\\
3214.25955316811	1.09866357258508\\
3243.56000488341	1.08466646787492\\
3272.8604565987	1.05483432904991\\
3302.160908314	1.00968181278724\\
3331.4613600293	0.950048800228693\\
3360.7618117446	0.877088071806481\\
3384.20217311684	0.810084836706592\\
3407.64253448907	0.736332031887184\\
3431.08289586131	0.656785040048408\\
3448.66316689049	0.593942541177733\\
3466.24343791967	0.528882787973249\\
3483.82370894885	0.462082759246009\\
3495.54388963497	0.416823532001606\\
3507.26407032109	0.371154580007798\\
3518.9842510072	0.325223420805006\\
3530.70443169332	0.279177757358764\\
3542.42461237944	0.233164861330508\\
3548.2847027225	0.210216453800139\\
3554.14479306556	0.187330956644857\\
3560.00488340862	0.164526383572827\\
3565.86497375168	0.141820606527005\\
3571.72506409474	0.119231336952416\\
3577.5851544378	0.0967761072337489\\
3583.44524478086	0.0744722523392822\\
3589.30533512392	0.0523368916980917\\
3595.16542546698	0.0303869113242634\\
3601.02551581004	0.00863894622100714\\
3606.88560615309	-0.012890636917385\\
3612.74569649615	-0.0341857566787052\\
3618.60578683921	-0.0552306335631772\\
3624.46587718227	-0.0760098062416358\\
3630.32596752533	-0.096508147431388\\
3636.18605786839	-0.116710879352488\\
3642.04614821145	-0.136603588761247\\
3647.90623855451	-0.156172241528107\\
3653.76632889757	-0.175403196746636\\
3659.62641924063	-0.194283220345604\\
3665.48650958369	-0.212799498214104\\
3677.20669026981	-0.248691735036614\\
3688.92687095593	-0.282986257974658\\
3700.64705164204	-0.315596884907732\\
3712.36723232816	-0.346445221558062\\
3729.94750335734	-0.389261863212917\\
3747.52777438652	-0.427754991910525\\
3765.1080454157	-0.461765987232056\\
3782.68831644488	-0.491180624234413\\
3806.12867781712	-0.523136771096618\\
3829.56903918935	-0.546762493786754\\
3853.00940056159	-0.56215544655377\\
3882.30985227689	-0.570182974569313\\
3911.61030399219	-0.566439601074606\\
3940.91075570748	-0.551920965364288\\
3970.21120742278	-0.527855253501815\\
3999.51165913808	-0.495656870274581\\
4028.81211085337	-0.456875600906461\\
4058.11256256867	-0.413143114833487\\
4087.41301428397	-0.366118727813539\\
4110.85337565621	-0.327229177393526\\
4134.29373702845	-0.28808386933809\\
4157.73409840068	-0.249440571688767\\
4181.17445977292	-0.211997351839675\\
4204.61482114516	-0.176380825722701\\
4222.19509217434	-0.151198593510241\\
4239.77536320352	-0.12755014350599\\
4257.3556342327	-0.10560367653161\\
4274.93590526187	-0.0854970701767792\\
4286.65608594799	-0.0731689627272707\\
4298.37626663411	-0.0617309000364902\\
4310.09644732023	-0.0512009193400539\\
4321.81662800635	-0.0415910668354822\\
4333.53680869247	-0.0329075041805321\\
4339.39689903553	-0.0289134316054568\\
4345.25698937859	-0.0251506476336813\\
4351.11707972165	-0.0216183345308175\\
4356.97717006471	-0.0183153384010472\\
4362.83726040776	-0.0152401761018386\\
4368.69735075082	-0.0123910425795004\\
4374.55744109388	-0.00976581860991315\\
4380.41753143694	-0.00736207892878549\\
4386.27762178	-0.00517710073465855\\
4392.13771212306	-0.00320787254820711\\
4397.99780246612	-0.00145110341038938\\
4403.85789280918	9.67675980635635e-05\\
4409.71798315224	0.00143956153411115\\
4415.5780734953	0.00258134948502271\\
4421.43816383836	0.00352644219481869\\
4427.29825418142	0.00427937944383579\\
4433.15834452448	0.00484491920031939\\
4439.01843486754	0.00522802656357129\\
4444.8785252106	0.00543386251869322\\
4450.73861555366	0.00546777252190704\\
4456.59870589671	0.0053352749374476\\
4462.45879623977	0.00504204934483239\\
4468.31888658283	0.00459392473771451\\
4474.17897692589	0.00399686763332092\\
4480.03906726895	0.00325697011347182\\
4485.89915761201	0.00238043781621532\\
4491.75924795507	0.00137357789895909\\
4497.61933829813	0.000242786992088656\\
4503.47942864119	-0.0010054608368811\\
4509.33951898425	-0.00236462608924002\\
4515.19960932731	-0.00382811579186758\\
4521.05969967037	-0.00538929538585851\\
4526.91979001343	-0.00704150052832416\\
4532.77988035649	-0.0087780487888698\\
4538.63997069955	-0.0105922512223704\\
4544.50006104261	-0.0124774238011881\\
4550.36015138567	-0.0144268986882795\\
4556.22024172873	-0.0164340353348914\\
4562.08033207179	-0.0184922313870945\\
4567.94042241485	-0.020594933383191\\
4579.66060310097	-0.0249079484284059\\
4591.38078378708	-0.0293220225675828\\
4603.1009644732	-0.0337875244621813\\
4614.82114515932	-0.038256546367968\\
4626.54132584544	-0.0426831816400012\\
4644.12159687462	-0.0491488528904369\\
4661.7018679038	-0.0552849484671149\\
4679.28213893298	-0.0609690344120997\\
4696.86240996215	-0.0660940975749839\\
4714.44268099133	-0.0705695301285212\\
4737.88304236357	-0.0754016910161567\\
4761.32340373581	-0.078824432068685\\
4784.76376510804	-0.080767268888778\\
4808.20412648028	-0.0812137232874689\\
4831.64448785252	-0.0801984987493339\\
4855.08484922476	-0.0778033315588932\\
4878.525210597	-0.0741516930284429\\
4901.96557196923	-0.0694025438860658\\
4925.40593334147	-0.0637433610438042\\
4948.84629471371	-0.0573826692757273\\
4972.28665608595	-0.0505423156890993\\
4989.86692711513	-0.0452340866088032\\
5007.44719814431	-0.0398795686526454\\
5025.02746917348	-0.0345719762746803\\
5042.60774020266	-0.0294005769326532\\
5060.18801123184	-0.0244492091518728\\
5071.90819191796	-0.0213089132012892\\
5083.62837260408	-0.0183217482260002\\
5095.3485532902	-0.0155069288851972\\
5107.06873397632	-0.012881845815389\\
5118.78891466243	-0.0104619643084524\\
5130.50909534855	-0.00826074046401502\\
5136.36918569161	-0.00724576417652113\\
5142.22927603467	-0.00628955509950849\\
5148.08936637773	-0.0053932036677192\\
5153.94945672079	-0.00455766556858167\\
5159.80954706385	-0.00378376121516422\\
5165.66963740691	-0.00307217549572846\\
5171.52972774997	-0.00242345779756414\\
5177.38981809303	-0.00183802230208667\\
5183.24990843609	-0.00131614854860384\\
5189.10999877915	-0.000857982262230763\\
5194.97008912221	-0.000463536442155097\\
5200.83017946527	-0.000132692705418172\\
5206.69026980832	0.000134797118939733\\
5212.55036015138	0.000339309150907825\\
5218.41045049444	0.000481345381755638\\
5224.2705408375	0.000561531363910543\\
5230.13063118056	0.000580613675216832\\
5235.99072152362	0.000539457164811521\\
5241.85081186668	0.000439041988111794\\
5247.71090220974	0.000280460438792809\\
5253.5709925528	6.4913585937501e-05\\
5259.43108289586	-0.00020629227510786\\
5265.29117323892	-0.000531749349427142\\
5271.15126358198	-0.000909952237710022\\
5277.01135392504	-0.00133930199060966\\
5282.8714442681	-0.00181811031274662\\
5288.73153461116	-0.00234460390672503\\
5294.59162495422	-0.00291692894699858\\
5300.45171529728	-0.00353315567387322\\
5306.31180564034	-0.00419128309699356\\
5312.1718959834	-0.0048892437984047\\
5318.03198632646	-0.00562490882429211\\
5323.89207666952	-0.00639609265514323\\
5335.61225735563	-0.00803602210988691\\
5347.33243804175	-0.00979060947598667\\
5359.05261872787	-0.011640854441922\\
5370.77279941399	-0.0135673456095694\\
5382.49298010011	-0.0155504272908981\\
5394.21316078623	-0.0175703652333819\\
5411.79343181541	-0.0206264570249625\\
5429.37370284459	-0.0236561994680502\\
5446.95397387376	-0.0265968666337916\\
5464.53424490294	-0.0293892177415632\\
5482.11451593212	-0.0319784943625235\\
5499.6947869613	-0.0343153102349622\\
5523.13514833354	-0.0369645697970126\\
5546.57550970577	-0.0390053866662389\\
5570.01587107801	-0.0403775975488117\\
5593.45623245025	-0.0410455768274131\\
5616.89659382249	-0.0409985804708773\\
5640.33695519473	-0.0402502743376171\\
5663.77731656696	-0.0388374875334381\\
5687.2176779392	-0.0368182547657073\\
5710.65803931144	-0.0342692322365158\\
5734.09840068368	-0.0312825889960758\\
5757.53876205592	-0.0279624893494367\\
5775.11903308509	-0.0253209774353723\\
5792.69930411427	-0.0226041612950089\\
5810.27957514345	-0.0198615860907476\\
5827.85984617263	-0.0171422580846296\\
5845.44011720181	-0.0144937414268623\\
5863.02038823099	-0.0119613011743077\\
5874.7405689171	-0.0103584689113255\\
5886.46074960322	-0.00883762661927349\\
5898.18093028934	-0.00740951237153606\\
5909.90111097546	-0.00608394503761024\\
5921.62129166158	-0.00486975521662219\\
5927.48138200464	-0.00430690126063523\\
5933.3414723477	-0.0037747252939739\\
5939.20156269076	-0.00327402805282989\\
5945.06165303382	-0.0028055389858606\\
5950.92174337688	-0.00236991521549063\\
5956.78183371994	-0.00196774065625192\\
5962.64192406299	-0.00159952529117195\\
5968.50201440605	-0.00126570460641393\\
5974.36210474911	-0.000966639184324193\\
5980.22219509217	-0.000702614454714258\\
5986.08228543523	-0.000473840603818683\\
5991.94237577829	-0.000280452640003633\\
5997.80246612135	-0.000122510615126645\\
nan	nan\\
6009.52264680747	8.71677877848544e-05\\
6015.38273715053	0.00013915470579904\\
6021.24282749359	0.000156195260852918\\
6027.10291783665	0.000138595537882602\\
6032.96300817971	8.67320865952666e-05\\
6038.82309852277	1.05066909580283e-06\\
6044.68318886583	-0.000117935128014562\\
6050.54327920889	-0.000269645413337205\\
6056.40336955195	-0.000453435638175552\\
6062.26345989501	-0.000668598367619746\\
6068.12355023807	-0.000914365170761012\\
6073.98364058113	-0.00118990862541044\\
6079.84373092419	-0.00149434443234112\\
6085.70382126725	-0.00182673363421009\\
6091.5639116103	-0.00218608493357806\\
6097.42400195336	-0.0025713571050529\\
6103.28409229642	-0.00298146149550286\\
6109.14418263948	-0.00341526460706999\\
6115.00427298254	-0.00387159075676286\\
6126.72445366866	-0.00484691496030564\\
6138.44463435478	-0.00589729026377876\\
6150.1648150409	-0.00701207879060207\\
6161.88499572702	-0.0081802436172531\\
6173.60517641314	-0.00939044587901222\\
6185.32535709926	-0.0106311431676606\\
6202.90562812843	-0.0125238849122458\\
6220.48589915761	-0.0144198008258561\\
6238.06617018679	-0.0162803607417812\\
6255.64644121597	-0.0180684446572701\\
6273.22671224515	-0.0197490117059137\\
6296.66707361738	-0.0217670759174429\\
6320.10743498962	-0.0234694346117178\\
6343.54779636186	-0.024800895361738\\
6366.9881577341	-0.0257200549059076\\
6390.42851910634	-0.0262003496211861\\
6413.86888047857	-0.026230592106393\\
6437.30924185081	-0.0258149879582533\\
6460.74960322305	-0.0249726439197466\\
6484.18996459529	-0.0237365949305218\\
6507.63032596753	-0.0221523927179412\\
6531.07068733976	-0.0202763119640315\\
6554.511048712	-0.0181732413749578\\
6572.09131974118	-0.0164896149110248\\
6589.67159077036	-0.0147497797944119\\
6607.25186179954	-0.0129858417428408\\
6624.83213282871	-0.0112298733508544\\
6642.41240385789	-0.00951331145706943\\
6659.99267488707	-0.00786637594471415\\
6671.71285557319	-0.00682123976455288\\
6683.43303625931	-0.00582763032354333\\
6695.15321694543	-0.00489292472691127\\
6706.87339763155	-0.00402393187227838\\
6718.59357831766	-0.00322684116990391\\
6724.45366866072	-0.0028570125857194\\
6730.31375900378	-0.00250717696029412\\
6736.17384934684	-0.00217791442732509\\
6742.0339396899	-0.00186975895010211\\
6747.89403003296	-0.00158319743259294\\
6753.75412037602	-0.00131866893334073\\
6759.61421071908	-0.00107656398377321\\
6765.47430106214	-0.000857224011372839\\
6771.3343914052	-0.000660940868717434\\
6777.19448174826	-0.00048795646860541\\
6783.05457209132	-0.000338462525666669\\
6788.91466243438	-0.000212600404349719\\
6794.77475277744	-0.000110461073133609\\
6800.6348431205	-3.20851645051866e-05\\
6806.49493346356	2.25368599307568e-05\\
6812.35502380662	5.34644400932771e-05\\
6818.21511414968	6.08065448811398e-05\\
6824.07520449274	4.47211977396864e-05\\
6829.9352948358	5.41490181981804e-06\\
6835.79538517886	-5.68580335656158e-05\\
6841.65547552191	-0.000141796263310171\\
6847.51556586497	-0.000249052271612713\\
6853.37565620803	-0.000378233329560688\\
6859.23574655109	-0.000528902542717993\\
6865.09583689415	-0.000700579986113411\\
6870.95592723721	-0.000892743923840465\\
6876.81601758027	-0.00110483211031015\\
6882.67610792333	-0.00133624317001102\\
6888.53619826639	-0.00158633805257091\\
6894.39628860945	-0.00185444155975373\\
6900.25637895251	-0.00213984394061383\\
6906.11646929557	-0.00244180255162248\\
6917.83664998169	-0.00309226380890344\\
6929.55683066781	-0.00379929311741169\\
6941.27701135393	-0.00455594769407947\\
6952.99719204004	-0.00535493635989402\\
6964.71737272616	-0.00618868403251014\\
6976.43755341228	-0.00704939798005758\\
6994.01782444146	-0.00837363117112833\\
7011.59809547064	-0.00971356553958538\\
7029.17836649982	-0.0110422187516779\\
7046.75863752899	-0.0123332146626785\\
7064.33890855817	-0.0135612713315236\\
7081.91917958735	-0.0147026617481923\\
7105.35954095959	-0.0160524478550469\\
7128.79990233183	-0.017164911640962\\
7152.24026370407	-0.0180039994718255\\
7175.6806250763	-0.0185435081848604\\
7199.12098644854	-0.0187677824095412\\
7222.56134782078	-0.0186720486407515\\
7246.00170919302	-0.0182623814027911\\
7269.44207056526	-0.0175553090788418\\
7292.88243193749	-0.0165770787554313\\
7316.32279330973	-0.0153626103708321\\
7339.76315468197	-0.0139541802068539\\
7363.20351605421	-0.0123998820348884\\
7380.78378708338	-0.011169733244554\\
7398.36405811256	-0.0099098093431418\\
7415.94432914174	-0.00864335340509519\\
7433.52460017092	-0.00739345189657561\\
7451.1048712001	-0.00618260106409093\\
7462.82505188622	-0.00540772731425485\\
7474.54523257233	-0.00466582556340511\\
7486.26541325845	-0.00396260676462862\\
7497.98559394457	-0.00330341083755303\\
7509.70577463069	-0.00269316478848812\\
7521.42595531681	-0.00213634463405881\\
7527.28604565987	-0.00187923625356278\\
7533.14613600293	-0.00163694141651033\\
7539.00622634599	-0.001409878362134\\
7544.86631668905	-0.00119843156113038\\
7550.72640703211	-0.00100295104139083\\
7556.58649737517	-0.000823751788772461\\
7562.44658771823	-0.000661113224018044\\
7568.30667806129	-0.000515278756450276\\
7574.16676840435	-0.000386455415196423\\
7580.02685874741	-0.000274813558354109\\
7585.88694909047	-0.0001804866603871\\
7591.74703943353	-0.000103571177926833\\
7597.60712977658	-4.41264938878022e-05\\
7603.46722011964	-2.17493976741762e-06\\
7609.3273104627	2.22981042770951e-05\\
7615.18740080576	2.93440320654991e-05\\
7621.04749114882	1.90507574305546e-05\\
7626.90758149188	-8.45761550854606e-06\\
7632.76767183494	-5.30211968507886e-05\\
7638.627762178	-0.000114444789048414\\
7644.48785252106	-0.000192498437387347\\
7650.34794286412	-0.000286918051360987\\
7656.20803320718	-0.000397406095374902\\
7662.06812355024	-0.000523632347051589\\
7667.9282138933	-0.000665234721270361\\
7673.78830423636	-0.000821820157997939\\
7679.64839457942	-0.000992965571739724\\
7685.50848492248	-0.00117821886042376\\
7691.36857526554	-0.00137709997138264\\
7697.2286656086	-0.00158910202191079\\
7703.08875595166	-0.00181369247194271\\
7714.80893663777	-0.00229838750201581\\
7726.52911732389	-0.00282645917370073\\
7738.24929801001	-0.00339285388872418\\
7749.96947869613	-0.0039922353658527\\
7761.68965938225	-0.0046190316944184\\
7773.40984006837	-0.0052674840045619\\
7790.99011109754	-0.00626784498545595\\
7808.57038212672	-0.00728343591455004\\
7826.1506531559	-0.00829399605778873\\
7843.73092418508	-0.0092795967675229\\
7861.31119521426	-0.0102210139008855\\
7878.89146624344	-0.0111000827878135\\
7902.33182761568	-0.0121464344563453\\
7925.77218898791	-0.0130173601795693\\
7949.21255036015	-0.0136843239245271\\
7972.65291173239	-0.014126073263905\\
7996.09327310463	-0.0143292441174659\\
8019.53363447687	-0.014288693288608\\
8042.9739958491	-0.0140075516483069\\
8066.41435722134	-0.013497000204998\\
8089.85471859358	-0.0127757804825216\\
8113.29507996582	-0.0118694592992409\\
8136.73544133805	-0.0108094759260288\\
8160.17580271029	-0.0096320064451336\\
8177.75607373947	-0.00869555529431648\\
8195.33634476865	-0.00773282548940009\\
8212.91661579783	-0.00676169411315181\\
8230.496886827	-0.00580001861840479\\
8248.07715785618	-0.00486529950990427\\
8259.7973385423	-0.00426547128903829\\
8271.51751922842	-0.00368986735086344\\
8283.23769991454	-0.00314300127055771\\
8294.95788060066	-0.00262911461177559\\
8306.67806128678	-0.00215214304136545\\
8318.3982419729	-0.0017156852439257\\
8324.25833231596	-0.00151367448811043\\
8330.11842265902	-0.00132297487807245\\
8335.97851300208	-0.00114392951863701\\
8341.83860334514	-0.000976855789136341\\
8347.69869368819	-0.000822044759626161\\
8353.55878403125	-0.000679760664552227\\
8359.41887437431	-0.00055024043461665\\
8365.27896471737	-0.000433693287773499\\
8371.13905506043	-0.000330300379891579\\
8376.99914540349	-0.000240214515715671\\
8382.85923574655	-0.000163559920488515\\
8388.71932608961	-0.000100432072515783\\
8394.57941643267	-5.08975968296321e-05\\
8400.43950677573	-1.49942199766626e-05\\
8406.29959711879	7.26921418331051e-06\\
8412.15968746185	1.59126679020703e-05\\
8418.01977780491	1.09847725478301e-05\\
8423.87986814797	-7.43736831905056e-06\\
8429.73995849103	-3.92484351895344e-05\\
8435.60004883409	-8.4315207436601e-05\\
8441.46013917715	-0.000142476935878959\\
8447.32022952021	-0.000213545772104161\\
8453.18031986327	-0.000297307253461387\\
8459.04041020633	-0.000393520842531137\\
8464.90050054939	-0.0005019205197564\\
8470.76059089244	-0.000622215427785242\\
8476.6206812355	-0.000754090566093032\\
8482.48077157856	-0.000897207534121875\\
8488.34086192162	-0.00105120532135851\\
8494.20095226468	-0.00121570114241605\\
8500.06104260774	-0.00139029131535057\\
8511.78122329386	-0.00176804106188055\\
8523.50140397998	-0.00218084307301458\\
8535.2215846661	-0.00262481423580925\\
8546.94176535222	-0.00309583426461801\\
8558.66194603833	-0.00358958192631633\\
8570.38212672445	-0.00410157272344232\\
8582.10230741057	-0.00462719770940002\\
8599.68257843975	-0.00543091830670217\\
8617.26284946893	-0.0062387563887071\\
8634.84312049811	-0.0070348019703062\\
8652.42339152729	-0.00780352308687485\\
8670.00366255646	-0.00853005589558444\\
8693.4440239287	-0.00940924703232076\\
8716.88438530094	-0.0101592461965354\\
8740.32474667318	-0.0107551915574221\\
8763.76510804542	-0.0111777149487521\\
8787.20546941765	-0.0114135236097825\\
8810.64583078989	-0.0114557714405134\\
8834.08619216213	-0.0113042099318977\\
8857.52655353437	-0.0109651164000528\\
8880.96691490661	-0.0104510046110101\\
8904.40727627884	-0.00978013008476346\\
8927.84763765108	-0.00897580907050897\\
8951.28799902332	-0.00806557617746692\\
8968.8682700525	-0.00733182030297483\\
8986.44854108167	-0.00656959696098284\\
9004.02881211085	-0.00579316675727577\\
9021.60908314003	-0.00501693804219783\\
9039.18935416921	-0.00425519297103749\\
9050.90953485533	-0.00376229125416531\\
9062.62971554145	-0.00328597535268875\\
9074.34989622757	-0.00283004623514711\\
9086.07007691369	-0.00239811150803011\\
9097.79025759981	-0.00199355611571542\\
9109.51043828592	-0.00161951508258513\\
9121.23061897204	-0.00127884851523458\\
9127.0907093151	-0.00112184366623683\\
9132.95079965816	-0.000974119061537158\\
9138.81089000122	-0.000835945287546409\\
9144.67098034428	-0.00070757200220806\\
9150.53107068734	-0.000589227467265248\\
9156.3911610304	-0.000481118127020635\\
9162.25125137346	-0.000383428234241947\\
9168.11134171652	-0.00029631952393277\\
9173.97143205958	-0.000219930935522611\\
9179.83152240264	-0.000154378383879877\\
9185.6916127457	-9.97545795444341e-05\\
9191.55170308876	-5.61288983802229e-05\\
9197.41179343182	-2.35473007952853e-05\\
9203.27188377488	-2.03230056979739e-06\\
9209.13197411794	8.41701679615916e-06\\
9214.992064461	7.82492638399595e-06\\
9220.85215480405	-3.76105290718477e-06\\
9226.71224514711	-2.62703238093769e-05\\
9232.57233549017	-5.96094243329286e-05\\
9238.43242583323	-0.000103662289029873\\
9244.29251617629	-0.000158290557183506\\
9250.15260651935	-0.000223333927157371\\
9256.01269686241	-0.000298610556014676\\
9261.87278720547	-0.000383917503452488\\
9267.73287754853	-0.000479031218977657\\
9273.59296789159	-0.000583708071202371\\
9279.45305823465	-0.000697684918003724\\
9285.31314857771	-0.000820679716238142\\
9291.17323892077	-0.000952392169648613\\
9297.03332926383	-0.00109250441344431\\
9302.89341960689	-0.00124068173402229\\
9314.613600293	-0.00155981305872976\\
9326.33378097912	-0.0019068008664674\\
9338.05396166524	-0.00227844166743545\\
9349.77414235136	-0.00267134196348359\\
9361.49432303748	-0.00308194789715026\\
9373.2145037236	-0.00350657608903786\\
9390.79477475278	-0.0041615200974647\\
9408.37504578196	-0.00482648881086774\\
9425.95531681113	-0.00548839807309779\\
9443.53558784031	-0.00613433036449451\\
9461.11585886949	-0.00675177795219038\\
9478.69612989867	-0.00732887633375124\\
9502.13649127091	-0.00801670951534958\\
9525.57685264315	-0.00859032199472302\\
9549.01721401538	-0.00903075302024208\\
9572.45757538762	-0.00932372693291143\\
9595.89793675986	-0.00946008298952435\\
9619.3382981321	-0.00943602918368367\\
9642.77865950433	-0.00925321340688587\\
9666.21902087657	-0.00891861146941937\\
9689.65938224881	-0.00844423763092448\\
9713.09974362105	-0.00784668916752573\\
9736.54010499328	-0.00714654193616621\\
9759.98046636552	-0.00636761870536307\\
9777.5607373947	-0.00574744695550532\\
9795.14100842388	-0.0051093595419891\\
9812.72127945306	-0.00446527880737508\\
9830.30155048224	-0.00382715294154637\\
9847.88182151142	-0.00320672950243876\\
9859.60200219753	-0.00280856542828907\\
9871.32218288365	-0.00242652378287964\\
9883.04236356977	-0.00206366072592535\\
9894.76254425589	-0.00172285567843108\\
9906.48272494201	-0.00140678791055024\\
9918.20290562813	-0.00111791493181928\\
9924.06299597119	-0.000984378947466937\\
9929.92308631425	-0.000858452856706486\\
9935.78317665731	-0.000740373083796323\\
9941.64326700037	-0.000630358900904697\\
9947.50335734343	-0.000528612007057587\\
9953.36344768649	-0.0004353161452437\\
9959.22353802955	-0.000350636758274933\\
9965.08362837261	-0.000274720684174514\\
9970.94371871567	-0.000207695891532479\\
9976.80380905872	-0.000149671255333666\\
9982.66389940178	-0.000100736373622287\\
9988.52398974484	-6.09614252983426e-05\\
9994.3840800879	-3.03970692294141e-05\\
10000.244170431	-9.07438479967228e-06\\
10006.104260774	2.99514607758894e-06\\
10011.9643511171	5.81961555461356e-06\\
10017.8244414601	-5.73374805836168e-07\\
10023.6845318032	-1.61368196951621e-05\\
10029.5446221463	-4.08044564115603e-05\\
10035.4047124893	-7.44909514947628e-05\\
10041.2648028324	-0.00011709212714911\\
10047.1248931754	-0.00016848522688702\\
10052.9849835185	-0.000228529219733049\\
10058.8450738616	-0.000297065142248737\\
10064.7051642046	-0.000373916477591258\\
10070.5652545477	-0.000458889570673991\\
10076.4253448907	-0.000551774078525254\\
10082.2854352338	-0.000652343454739857\\
10088.1455255769	-0.000760355466990062\\
10094.0056159199	-0.000875552746390671\\
10099.865706263	-0.000997663367469788\\
10111.5858869491	-0.00126146783372883\\
10123.3060676352	-0.00154932617225742\\
10135.0262483213	-0.00185860456324919\\
10146.7464290074	-0.00218650118910178\\
10158.4666096936	-0.00253007066847834\\
10170.1867903797	-0.00288624958414371\\
10187.7670614089	-0.00343724056602401\\
10205.347332438	-0.00399859697078768\\
10222.9276034672	-0.00455930909219037\\
10240.5078744964	-0.00510845869602223\\
10258.0881455256	-0.00563542519493778\\
10275.6684165548	-0.00613008474583198\\
10299.108777927	-0.00672318222266476\\
10322.5491392992	-0.00722223915348606\\
10345.9895006715	-0.00761068283890666\\
10369.4298620437	-0.00787581883625711\\
10392.8702234159	-0.00800922188084799\\
10416.3105847882	-0.00800697993267605\\
10439.7509461604	-0.00786978450421604\\
10463.1913075327	-0.00760286563735366\\
10486.6316689049	-0.0072157751154065\\
10510.0720302771	-0.00672202654818299\\
10533.5123916494	-0.00613860568526369\\
10556.9527530216	-0.00548536854095951\\
10574.5330240508	-0.00496281748628308\\
10592.11329508	-0.00442321470788328\\
10609.6935661091	-0.00387667524878795\\
10627.2738371383	-0.00333338047999822\\
10644.8541081675	-0.00280338471955571\\
10656.5742888536	-0.00246226982840681\\
10668.2944695397	-0.00213417286747633\\
10680.0146502259	-0.00182174124790012\\
10691.734830912	-0.00152747882211605\\
10703.4550115981	-0.00125372537631666\\
10715.1751922842	-0.00100263759482658\\
10721.0352826273	-0.000886212905200961\\
10726.8953729703	-0.000776171650899172\\
10732.7554633134	-0.000672724456844635\\
10738.6155536565	-0.000576067563078175\\
10744.4756439995	-0.000486382440981392\\
10750.3357343426	-0.000403835441610912\\
10756.1958246856	-0.000328577476596296\\
10762.0559150287	-0.000260743732411699\\
10767.9160053717	-0.000200453418395766\\
10773.7760957148	-0.000147809549082245\\
10779.6361860579	-0.000102898761174201\\
10785.4962764009	-6.57911654813149e-05\\
10791.356366744	-3.65402340709376e-05\\
10797.216457087	-1.51827227904997e-05\\
10803.0765474301	-1.73862924079821e-06\\
10808.9366377732	3.78881377659286e-06\\
10814.7967281162	1.41310940428786e-06\\
10820.6568184593	-8.83556739862741e-06\\
10826.5169088023	-2.69104690456881e-05\\
10832.3769991454	-5.27484092917896e-05\\
10838.2370894885	-8.62699316219065e-05\\
10844.0971798315	-0.000127379511635701\\
10849.9572701746	-0.000175965792927513\\
10855.8173605176	-0.000231901855888585\\
10861.6774508607	-0.000295045518808255\\
10867.5375412038	-0.000365239670544749\\
10873.3976315468	-0.000442312634040603\\
10879.2577218899	-0.000526078559795566\\
10885.1178122329	-0.000616337848460132\\
10890.977902576	-0.000712877601540279\\
10896.8379929191	-0.000815472099262405\\
10908.5581736052	-0.00103786139158772\\
10920.2783542913	-0.00128146330477403\\
10931.9985349774	-0.00154406501829236\\
10943.7187156635	-0.00182330269912859\\
10955.4388963497	-0.00211668208161304\\
10967.1590770358	-0.00242160006130094\\
10978.8792577219	-0.00273536711684524\\
10996.4595287511	-0.00321657624242851\\
11014.0397997802	-0.00370205835039019\\
11031.6200708094	-0.00418235783207747\\
11049.2003418386	-0.00464817693161243\\
11066.7806128678	-0.00509055100250311\\
11090.22097424	-0.00562937647035128\\
11113.6613356123	-0.00609337664524348\\
11137.1016969845	-0.00646704616429113\\
11160.5420583567	-0.00673804868416983\\
11183.982419729	-0.00689760606791184\\
11207.4227811012	-0.00694076489891446\\
11230.8631424734	-0.00686653236261184\\
11254.3035038457	-0.00667787797909878\\
11277.7438652179	-0.00638160219779822\\
11301.1842265902	-0.00598807732340355\\
11324.6245879624	-0.00551087048014656\\
11348.0649493346	-0.00496626220167837\\
11365.6452203638	-0.00452459467654208\\
11383.225491393	-0.0040637079253725\\
11400.8057624222	-0.00359227986894734\\
11418.3860334513	-0.00311913899205397\\
11435.9663044805	-0.00265309607988796\\
11453.5465755097	-0.00220277726615378\\
11465.2667561958	-0.00191540775218794\\
11476.9869368819	-0.00164105178308009\\
11488.7071175681	-0.00138192579859688\\
11500.4272982542	-0.00114010927397309\\
11512.1474789403	-0.000917527687418584\\
11523.8676596264	-0.000715936851649411\\
11529.7277499695	-0.000623511742444622\\
11535.5878403125	-0.000536908736441398\\
11541.4479306556	-0.000456293541417036\\
11547.3080209987	-0.000381818895900072\\
11553.1681113417	-0.000313624270939049\\
11559.0282016848	-0.000251835600389362\\
11564.8882920278	-0.000196565040367412\\
11570.7483823709	-0.000147910758218454\\
11576.608472714	-0.000105956751441114\\
11582.468563057	-7.07726968791103e-05\\
11588.3286534001	-4.24138304485396e-05\\
11594.1887437431	-2.09208575862028e-05\\
11600.0488340862	-6.31989455759075e-06\\
11605.9089244293	1.37755931010916e-06\\
11611.7690147723	2.17461847321223e-06\\
11617.6291051154	-3.91102284497095e-06\\
11623.4891954584	-1.68471553991068e-05\\
11629.3492858015	-3.65871502018753e-05\\
11635.2093761446	-6.3070084082781e-05\\
11641.0694664876	-9.62208952007329e-05\\
11646.9295568307	-0.00013595056812842\\
11652.7896471737	-0.000182156348062746\\
11658.6497375168	-0.000234721983653205\\
11664.5098278598	-0.000293517997887542\\
11670.3699182029	-0.000358401986397821\\
11676.230008546	-0.00042921894251929\\
11682.090098889	-0.000505801608369831\\
11687.9501892321	-0.000587970851146087\\
11693.8102795751	-0.000675536063791807\\
11699.6703699182	-0.000768295589181372\\
11711.3905506043	-0.000968538401198433\\
11723.1107312904	-0.00118688253509249\\
11734.8309119766	-0.00142136517875597\\
11746.5510926627	-0.00166989457027569\\
11758.2712733488	-0.0019302681327759\\
11769.9914540349	-0.00220019147766297\\
11787.5717250641	-0.00261779089249372\\
11805.1519960933	-0.00304335659034545\\
11822.7322671225	-0.00346859907837759\\
11840.3125381516	-0.00388528109790063\\
11857.8928091808	-0.00428537354115635\\
11875.47308021	-0.0046612065759582\\
11898.9134415822	-0.00511228213182316\\
11922.3538029545	-0.00549238781977268\\
11945.7941643267	-0.00578885046268634\\
11969.2345256989	-0.0059919061371329\\
11992.6748870712	-0.00609500790493707\\
12016.1152484434	-0.00609502320517862\\
12039.5556098157	-0.00599231505753933\\
12062.9959711879	-0.00579070517139075\\
12086.4363325601	-0.00549732103098363\\
12109.8766939324	-0.0051223329093238\\
12133.3170553046	-0.00467859041602923\\
12156.7574166768	-0.00418117149073315\\
12174.337687706	-0.00378293033943343\\
12191.9179587352	-0.00337143889706234\\
12209.4982297644	-0.00295443582151013\\
12227.0785007936	-0.00253972531043569\\
12244.6587718227	-0.00213502873525683\\
12256.3789525089	-0.00187450803930013\\
12268.099133195	-0.0016239098969604\\
12279.8193138811	-0.00138527762529527\\
12291.5394945672	-0.00116054662920597\\
12303.2596752533	-0.000951528422873768\\
12314.9798559394	-0.000759895750868374\\
12320.8399462825	-0.000671080390199808\\
12326.7000366256	-0.000587168930638313\\
12332.5601269686	-0.000508325572768237\\
12338.4202173117	-0.000434703527455033\\
12344.2803076547	-0.000366444710434796\\
12350.1403979978	-0.000303679461249289\\
12356.0004883409	-0.000246526287144587\\
12361.8605786839	-0.000195091632371791\\
12367.720669027	-0.00014946967337851\\
12373.58075937	-0.000109742140246979\\
12379.4408497131	-7.59781647374057e-05\\
12385.3009400562	-4.82341551947391e-05\\
12391.1610303992	-2.65536985479246e-05\\
12397.0211207423	-1.09674895595291e-05\\
12402.8812110853	-1.49328743399849e-06\\
12408.7413014284	1.86410017836111e-06\\
12414.6013917715	-8.87194216363841e-07\\
12420.4614821145	-9.72613667174347e-06\\
12426.3215724576	-2.4618857708653e-05\\
12432.1816628006	-4.55187452331639e-05\\
12438.0417531437	-7.23665642196277e-05\\
12443.9018434868	-0.000105090602865757\\
12449.7619338298	-0.000143606844864621\\
12455.6220241729	-0.000187819167386013\\
12461.4821145159	-0.000237619564309568\\
12467.342204859	-0.000292888394196339\\
12473.2022952021	-0.000353494652429067\\
12479.0623855451	-0.000419296266918004\\
12484.9224758882	-0.00049014041669363\\
12490.7825662312	-0.000565863872680999\\
12496.6426565743	-0.000646293359927287\\
12508.3628372604	-0.000820529415734817\\
12520.0830179465	-0.00101127650213027\\
12531.8031986326	-0.0012168282831699\\
12543.5233793188	-0.00143535901862137\\
12555.2435600049	-0.00166493937274596\\
12566.963740691	-0.00190355304595082\\
12584.5440117202	-0.00227383334547552\\
12602.1242827494	-0.00265249363169663\\
12619.7045537785	-0.00303217118216235\\
12637.2848248077	-0.00340552026183173\\
12654.8650958369	-0.00376535139099728\\
12672.4453668661	-0.00410476697843897\\
12695.8857282383	-0.00451444596939408\\
12719.3260896105	-0.00486257911584817\\
12742.7664509828	-0.00513752832805438\\
12766.206812355	-0.0053301961663468\\
12789.6471737273	-0.0054343159556759\\
12813.0875350995	-0.00544664484648214\\
12836.5278964717	-0.0053670539438863\\
12859.968257844	-0.00519851312115898\\
12883.4086192162	-0.00494697168780536\\
12906.8489805884	-0.00462113957348873\\
12930.2893419607	-0.00423217700209131\\
12953.7297033329	-0.00379330365670004\\
12971.3099743621	-0.00344026461380296\\
12988.8902453913	-0.0030741463658034\\
13006.4705164205	-0.00270184819879727\\
13024.0507874496	-0.00233035573420049\\
13041.6310584788	-0.0019666079865307\\
13053.3512391649	-0.00173176206430656\\
13065.0714198511	-0.00150530228981414\\
13076.7916005372	-0.00128908492892611\\
13088.5117812233	-0.00108487418502644\\
13100.2319619094	-0.000894327574196526\\
13111.9521425955	-0.000718982245030837\\
13123.6723232817	-0.000560242355756582\\
13129.5324136247	-0.000487499039817409\\
13135.3925039678	-0.000419367611854463\\
13141.2525943108	-0.00035598128924929\\
13147.1126846539	-0.000297463057510905\\
13152.9727749969	-0.000243925425517103\\
13158.83286534	-0.000195470203237305\\
13164.6929556831	-0.000152188302400557\\
13170.5530460261	-0.00011415956049917\\
13176.4131363692	-8.14525884456714e-05\\
13182.2732267122	-5.41246421823077e-05\\
13188.1333170553	-3.22215184816521e-05\\
13193.9934073984	-1.57774751126702e-05\\
13199.8534977414	-4.8151755103121e-06\\
13205.7135880845	6.54341970378372e-07\\
13211.5736784275	6.3167019038465e-07\\
13217.4337687706	-4.87098522725504e-06\\
13223.2938591137	-1.58298488625258e-05\\
13229.1539494567	-3.22096428508621e-05\\
13235.0140397998	-5.39636784412149e-05\\
13240.8741301428	-8.10339715372089e-05\\
13246.7342204859	-0.000113351381929817\\
13252.594310829	-0.000150835775911191\\
13258.454401172	-0.000193396211866665\\
13264.3144915151	-0.000240931148450913\\
13270.1745818581	-0.000293328674839658\\
13276.0346722012	-0.000350466762573949\\
13281.8947625443	-0.000412213538406629\\
13287.7548528873	-0.000478427577543718\\
13293.6149432304	-0.000548958216667969\\
13305.3351239165	-0.000702322459742393\\
13317.0553046026	-0.000870929261528749\\
13328.7754852887	-0.00105327689021096\\
13340.4956659748	-0.0012477519367211\\
13352.215846661	-0.00145264326805179\\
13363.9360273471	-0.00166615676518477\\
13375.6562080332	-0.00188643071881885\\
13393.2364790624	-0.00222532187206752\\
13410.8167500916	-0.00256852132462072\\
13428.3970211207	-0.0029093868254581\\
13445.9772921499	-0.00324134948453528\\
13463.5575631791	-0.0035580381126893\\
13486.9979245513	-0.00394612734397704\\
13510.4382859236	-0.00428324324341816\\
13533.8786472958	-0.00455807132229245\\
13557.3190086681	-0.00476146104900936\\
13580.7593700403	-0.00488672176510969\\
13604.1997314125	-0.00492983415279872\\
13627.6400927848	-0.00488957061478979\\
13651.080454157	-0.00476752102761162\\
13674.5208155292	-0.00456802352855954\\
13697.9611769015	-0.00429800319315724\\
13721.4015382737	-0.00396672453832547\\
13744.841899646	-0.00358546665249678\\
13762.4221706751	-0.00327443739756012\\
13780.0024417043	-0.00294842926925367\\
13797.5827127335	-0.00261360398102513\\
13815.1629837627	-0.00227626601918929\\
13832.7432547918	-0.00194274240345282\\
13850.323525821	-0.00161926250841621\\
13862.0437065071	-0.00141217123863335\\
13873.7638871933	-0.00121392866245647\\
13885.4840678794	-0.00102616068145417\\
13897.2042485655	-0.000850400377821322\\
13908.9244292516	-0.000688075310060081\\
13920.6446099377	-0.000540495733723048\\
13926.5047002808	-0.00047261021497196\\
13932.3647906239	-0.000408843844137917\\
13938.2248809669	-0.000349322803258595\\
13944.08497131	-0.000294164129152046\\
13949.945061653	-0.000243475478057306\\
13955.8051519961	-0.000197354910416347\\
13961.6652423392	-0.000155890696233735\\
13967.5253326822	-0.000119161141364189\\
13973.3854230253	-8.72344351123164e-05\\
13979.2455133683	-6.01685194012163e-05\\
13985.1056037114	-3.80109797838483e-05\\
13990.9656940544	-2.07989584786331e-05\\
13996.8257843975	-8.5590895997881e-06\\
14002.6858747406	-1.30745667827178e-06\\
14008.5459650836	9.50427459075276e-07\\
14014.4060554267	-1.78038155813986e-06\\
14020.2661457697	-9.48428422582442e-06\\
14026.1262361128	-2.21351840003748e-05\\
14031.9863264559	-3.96965562601339e-05\\
14037.8464167989	-6.21215392153655e-05\\
14043.706507142	-8.93530465502477e-05\\
14049.566597485	-0.000121323901525574\\
14055.4266878281	-0.000157956992237147\\
14061.2867781712	-0.000199165447665733\\
14067.1468685142	-0.000244852834125261\\
14073.0069588573	-0.00029491337167422\\
14078.8670492003	-0.00034923216998541\\
14084.7271395434	-0.000407685483172795\\
14090.5872298865	-0.000470140983008642\\
14096.4473202295	-0.000536458049918619\\
14108.1675009156	-0.000680074815249637\\
14119.8876816018	-0.000837257111180151\\
14131.6078622879	-0.0010066149790278\\
14143.328042974	-0.00118665920108699\\
14155.0482236601	-0.00137581413935497\\
14166.7684043462	-0.00157243127103134\\
14184.3486753754	-0.0018776013910078\\
14201.9289464046	-0.00218977685976661\\
14219.5092174338	-0.00250291496146803\\
14237.0894884629	-0.00281097820540363\\
14254.6697594921	-0.00310804890926719\\
14272.2500305213	-0.0033884410718132\\
14295.6903918935	-0.00372717211121178\\
14319.1307532658	-0.00401537813165297\\
14342.571114638	-0.00424340183347805\\
14366.0114760103	-0.00440366374096463\\
14389.4518373825	-0.00449090710846186\\
14412.8921987547	-0.00450236340228994\\
14436.332560127	-0.00443783318470261\\
14459.7729214992	-0.00429968009947604\\
14483.2132828714	-0.00409273859038365\\
14506.6536442437	-0.00382413888456508\\
14530.0940056159	-0.0035030555364934\\
14553.5343669882	-0.00314038836962424\\
14571.1146380173	-0.00284842832876437\\
14588.6949090465	-0.00254548029175913\\
14606.2751800757	-0.00223726496221597\\
14623.8554511049	-0.00192958274735686\\
14641.4357221341	-0.00162820335005547\\
14653.1559028202	-0.00143357435883956\\
14664.8760835063	-0.0012458651337258\\
14676.5962641924	-0.00106662478674308\\
14688.3164448785	-0.000897327202285143\\
14700.0366255646	-0.000739358761476168\\
14711.7568062508	-0.000594006834663746\\
14717.6168965938	-0.00052643465546198\\
14723.4769869369	-0.000462449137631033\\
14729.3370772799	-0.000402178771716434\\
14735.197167623	-0.00034574403913085\\
14741.0572579661	-0.000293257166773564\\
14746.9173483091	-0.000244821899432677\\
14752.7774386522	-0.000200533290366235\\
14758.6375289952	-0.000160477510506018\\
14764.4976193383	-0.000124731676664487\\
14770.3577096814	-9.33636990873341e-05\\
14776.2178000244	-6.64321486381091e-05\\
14782.0778903675	-4.39861438923031e-05\\
14787.9379807105	-2.60652583372505e-05\\
14793.7980710536	-1.26994478629447e-05\\
14799.6581613967	-3.90899866785218e-06\\
14805.5182517397	2.95504333086003e-07\\
14811.3783420828	-8.6811443922328e-08\\
14817.2384324258	-5.04711574925861e-06\\
14823.0985227689	-1.45669054970921e-05\\
14828.958613112	-2.86180551753425e-05\\
14834.818703455	-4.7162887539361e-05\\
14840.6787937981	-7.01542644141694e-05\\
14846.5388841411	-9.75356973902043e-05\\
14852.3989744842	-0.000129241478160138\\
14858.2590648272	-0.000165196828188351\\
14864.1191551703	-0.000205318067391454\\
14869.9792455134	-0.000249512801427311\\
14875.8393358564	-0.00029768012719494\\
14881.6994261995	-0.000349710856090986\\
14887.5595165425	-0.000405487754519023\\
14893.4196068856	-0.000464885801107435\\
14905.1397875717	-0.000594007970482249\\
14916.8599682578	-0.000735932211527269\\
14928.580148944	-0.000889407827800184\\
14940.3003296301	-0.00105308939068388\\
14952.0205103162	-0.00122554832635953\\
14963.7406910023	-0.00140528517846411\\
14975.4608716884	-0.00159074244143019\\
14993.0411427176	-0.00187614027613619\\
15010.6214137468	-0.00216527240276336\\
15028.201684776	-0.00245256272851965\\
15045.7819558052	-0.00273248936690219\\
15063.3622268343	-0.00299968929659867\\
15086.8025882066	-0.00332738854628988\\
15110.2429495788	-0.00361236235425926\\
15133.683310951	-0.00384503190591027\\
15157.1236723233	-0.00401762519138437\\
15180.5640336955	-0.0041244308448573\\
15204.0043950678	-0.00416198145516016\\
15227.44475644	-0.00412916048897852\\
15250.8851178122	-0.00402722956387104\\
15274.3254791845	-0.00385977550368891\\
15297.7658405567	-0.00363257930993113\\
15321.2062019289	-0.00335341180118523\\
15344.6465633012	-0.00303176311997486\\
15362.2268343304	-0.00276915263604038\\
15379.8071053595	-0.00249373244525841\\
15397.3873763887	-0.00221071580829229\\
15414.9676474179	-0.00192544384110782\\
15432.5479184471	-0.00164328368143192\\
15450.1281894763	-0.00136952648722847\\
15461.8483701624	-0.00119422610243507\\
15473.5685508485	-0.00102639131353052\\
15485.2887315346	-0.00086740798862672\\
15497.0089122207	-0.000718584280846323\\
15508.7290929068	-0.000581139735917738\\
15520.449273593	-0.000456195171261834\\
15526.309363936	-0.000398731341943496\\
15532.1694542791	-0.000344763410350772\\
15538.0295446221	-0.000294399737903101\\
15543.8896349652	-0.000247740948886146\\
15549.7497253083	-0.000204879725700828\\
15555.6098156513	-0.000165900621115956\\
15561.4699059944	-0.000130879887902337\\
15567.3299963374	-9.98853261971183e-05\\
15573.1900866805	-7.2976148895099e-05\\
15579.0501770236	-5.0202865334378e-05\\
15584.9102673666	-3.16071835082891e-05\\
15590.7703577097	-1.72219309810085e-05\\
15596.6304480527	-7.07099466777598e-06\\
15602.4905383958	-1.16927957373461e-06\\
15608.3506287389	4.773134301401e-07\\
15614.2107190819	-2.1281092255186e-06\\
15620.070809425	-8.97344969084332e-06\\
15625.930899768	-2.00376535697323e-05\\
15631.7909901111	-3.52907636955655e-05\\
15637.6510804542	-5.46939926726159e-05\\
15643.5111707972	-7.81998140127464e-05\\
15649.3712611403	-0.000105752071654079\\
15655.2313514833	-0.000137286107612723\\
15661.0914418264	-0.00017272890747811\\
15666.9515321695	-0.000211999263428327\\
15672.8116225125	-0.000255007954396796\\
15678.6717128556	-0.000301657943000936\\
15684.5318031986	-0.000351844588788032\\
15690.3918935417	-0.000405455877350871\\
15696.2519838848	-0.000462372664789287\\
15707.9721645709	-0.000585612083186852\\
15719.692345257	-0.000720477307532184\\
15731.4125259431	-0.000865786753700883\\
15743.1327066292	-0.00102027297176538\\
15754.8528873153	-0.00118259353263005\\
15766.5730680015	-0.00135134252604895\\
15784.1533390306	-0.00161332028411956\\
15801.7336100598	-0.00188140378297334\\
15819.313881089	-0.00215042256932537\\
15836.8941521182	-0.00241520420116699\\
15854.4744231474	-0.00267067248342958\\
15872.0546941765	-0.0029119435820022\\
15895.4950555488	-0.00320365634488596\\
15918.935416921	-0.00345215671861917\\
15942.3757782932	-0.00364910212262316\\
15965.8161396655	-0.0037879219159577\\
15989.2565010377	-0.00386403189965901\\
16012.69686241	-0.00387498105405611\\
16036.1372237822	-0.00382052583439182\\
16059.5775851544	-0.00370262980444198\\
16083.0179465267	-0.00352538891069844\\
16106.4583078989	-0.00329488520091443\\
16129.8986692712	-0.00301897418900374\\
16153.3390306434	-0.002707013282268\\
16170.9193016726	-0.00245568979215539\\
16188.4995727017	-0.00219476553655616\\
16206.0798437309	-0.00192917612599855\\
16223.6601147601	-0.00166393207529827\\
16241.2403857893	-0.00140402342593567\\
16252.9605664754	-0.00123612998662608\\
16264.6807471615	-0.00107417557523395\\
16276.4009278476	-0.000919504264411494\\
16288.1211085338	-0.000773396076263431\\
16299.8412892199	-0.000637056281542482\\
16311.561469906	-0.000511605350409784\\
16317.4215602491	-0.00045328785944793\\
16323.2816505921	-0.000398069638266874\\
16329.1417409352	-0.000346062880858526\\
16335.0018312782	-0.000297372884124726\\
16340.8619216213	-0.000252097831477033\\
16346.7220119643	-0.000210328591719765\\
16352.5821023074	-0.00017214853354908\\
16358.4421926505	-0.000137633356111667\\
16364.3022829935	-0.000106850935909349\\
16370.1623733366	-7.98611904064394e-05\\
16376.0224636796	-5.6715958578948e-05\\
16381.8825540227	-3.74588986587459e-05\\
16387.7426443658	-2.21254032706356e-05\\
16393.6027347088	-1.07425321233418e-05\\
16399.4628250519	-3.32896238068003e-06\\
16405.3229153949	1.05043195003138e-07\\
16411.183005738	-4.42349719875634e-07\\
16417.0430960811	-4.96455080859517e-06\\
16422.9031864241	-1.344656672264e-05\\
16428.7632767672	-2.58650404379933e-05\\
16434.6233671102	-4.21883082618168e-05\\
16440.4834574533	-6.2376474350329e-05\\
16446.3435477964	-8.63815025671135e-05\\
16452.2036381394	-0.000114147325468672\\
16458.0637284825	-0.000145609970177582\\
16463.9238188255	-0.000180697700851392\\
16469.7839091686	-0.000219331177446408\\
16475.6439995117	-0.000261423630404602\\
16481.5040898547	-0.000306881050907357\\
16487.3641801978	-0.000355602396256526\\
16493.2242705408	-0.000407479809935973\\
16504.944451227	-0.000520238766584936\\
16516.6646319131	-0.00064416803343282\\
16528.3848125992	-0.00077818507271979\\
16540.1049932853	-0.000921124013042951\\
16551.8251739714	-0.00107174565352221\\
16563.5453546575	-0.00122874806979584\\
16575.2655353437	-0.00139077773104547\\
16592.8458063728	-0.00164019175297522\\
16610.426077402	-0.00189296510718047\\
16628.0063484312	-0.00214423714665986\\
16645.5866194604	-0.00238918878177739\\
16663.1668904896	-0.00262313390868933\\
16686.6072518618	-0.00291026568369465\\
16710.047613234	-0.00316022850187203\\
16733.4879746063	-0.00336460994172886\\
16756.9283359785	-0.00351656461909876\\
16780.3686973508	-0.00361103955292185\\
16803.809058723	-0.00364493834982327\\
16827.2494200952	-0.00361721888623165\\
16850.6897814675	-0.00352892142605913\\
16874.1301428397	-0.00338312646340058\\
16897.5705042119	-0.0031848439498138\\
16921.0108655842	-0.00294083786981587\\
16944.4512269564	-0.00265939229299421\\
16962.0314979856	-0.00242942768344873\\
16979.6117690148	-0.00218810788894791\\
16997.1920400439	-0.00194000642106733\\
17014.7723110731	-0.00168981450016808\\
17032.3525821023	-0.00144225164530541\\
17049.9328531315	-0.00120197594788073\\
17061.6530338176	-0.00104807566845909\\
17073.3732145037	-0.000900703820765691\\
17085.0933951898	-0.000761084283144257\\
17096.813575876	-0.000630373411576671\\
17108.5337565621	-0.000509650373644626\\
17120.2539372482	-0.00039990815044365\\
17126.1140275913	-0.000349439536447095\\
17131.9741179343	-0.000302045281462762\\
17137.8342082774	-0.000257821725022549\\
17143.6942986204	-0.000216858420919931\\
17149.5543889635	-0.000179237953221132\\
17155.4144793066	-0.000145035767163624\\
17161.2745696496	-0.0001143200153026\\
17167.1346599927	-8.71514191847529e-05\\
17172.9947503357	-6.35831468621092e-05\\
17178.8548406788	-4.36607064672648e-05\\
17184.7149310219	-2.74218560787596e-05\\
17190.5750213649	-1.48965300430669e-05\\
17196.435111708	-6.10678190334502e-06\\
17202.295202051	-1.06674403792564e-06\\
17208.1552923941	2.17395917606153e-07\\
17214.0153827372	-2.25259817053585e-06\\
17219.8754730802	-8.467020996523e-06\\
17225.7355634233	-1.84082526626683e-05\\
17231.5956537663	-3.20508022419686e-05\\
17237.4557441094	-4.93613679520115e-05\\
17243.3158344524	-7.02989137991503e-05\\
17249.1759247955	-9.48147625138472e-05\\
17255.0360151386	-0.000122852704579042\\
17260.8961054816	-0.000154349123092558\\
17266.7561958247	-0.0001892331341979\\
17272.6162861677	-0.000227426742778135\\
17278.4763765108	-0.000268845013054363\\
17284.3364668539	-0.000313396253729843\\
17290.1965571969	-0.000360982217294307\\
17296.05664754	-0.000411498312991594\\
17307.7768282261	-0.000520872191859627\\
17319.4970089122	-0.000640563211729651\\
17331.2171895983	-0.000769530962286718\\
17342.9373702845	-0.000906658231333596\\
17354.6575509706	-0.00105076056980354\\
17366.3777316567	-0.00120059641004529\\
17383.9580026859	-0.0014332727330667\\
17401.5382737151	-0.00167145732269408\\
17419.1185447442	-0.00191057048899077\\
17436.6988157734	-0.00214602562042023\\
17454.2790868026	-0.0023733163380707\\
17471.8593578318	-0.00258810192568318\\
17495.299719204	-0.00284800142720994\\
17518.7400805762	-0.00306966032369933\\
17542.1804419485	-0.00324562526874034\\
17565.6208033207	-0.00337000730766418\\
17589.061164693	-0.00343867576068483\\
17612.5015260652	-0.00344939225448289\\
17635.9418874374	-0.00340188056494002\\
17659.3822488097	-0.00329783010838866\\
17682.8226101819	-0.00314083316079808\\
17706.2629715541	-0.00293625811604845\\
17729.7033329264	-0.00269106324199608\\
17753.1436942986	-0.00241355738528607\\
17770.7239653278	-0.00218983247968374\\
17788.304236357	-0.00195743755798805\\
17805.8845073862	-0.00172077508727904\\
17823.4647784153	-0.00148431933057469\\
17841.0450494445	-0.00125253118054514\\
17852.7652301306	-0.00110275988059628\\
17864.4854108167	-0.000958257168175217\\
17876.2055915029	-0.00082022808488528\\
17887.925772189	-0.000689821242698202\\
17899.6459528751	-0.000568119191981236\\
17911.3661335612	-0.00045612936152668\\
17917.2262239043	-0.000404068724687542\\
17923.0863142473	-0.000354775646231425\\
17928.9464045904	-0.00030835129829279\\
17934.8064949335	-0.000264890712639448\\
17940.6665852765	-0.000224482583897122\\
17946.5266756196	-0.000187209086336214\\
17952.3867659626	-0.00015314570456439\\
17958.2468563057	-0.000122361078494211\\
17964.1069466488	-9.49168628787064e-05\\
17969.9670369918	-7.08676017363276e-05\\
17975.8271273349	-5.0260617890078e-05\\
17981.6872176779	-3.31359178562001e-05\\
17987.547308021	-1.95261122754273e-05\\
17993.4073983641	-9.45635203853158e-06\\
17999.2674887071	-2.94428023271915e-06\\
nan	nan\\
18010.9876693932	-6.26058363293904e-07\\
18016.8477597363	-4.81744604422046e-06\\
18022.7078500794	-1.256161326216e-05\\
18028.5679404224	-2.38385014690801e-05\\
18034.4280307655	-3.86205909431193e-05\\
18040.2881211085	-5.68729641296078e-05\\
18046.1482114516	-7.85533845819969e-05\\
18052.0083017947	-0.000103612391332356\\
18057.8683921377	-0.000131993408469772\\
18063.7284824808	-0.000163632869696458\\
18069.5885728238	-0.000198460357582423\\
18075.4486631669	-0.000236398757207158\\
18081.3087535099	-0.000277364423876758\\
18087.168843853	-0.000321267364520118\\
18093.0289341961	-0.000368011432396361\\
18104.7491148822	-0.000469608852506454\\
18116.4692955683	-0.000581272632318132\\
18128.1894762544	-0.000702034797422572\\
18139.9096569405	-0.00083085174866964\\
18151.6298376267	-0.000966613188023733\\
18163.3500183128	-0.00110815159609184\\
18175.0701989989	-0.00125425218065833\\
18192.6504700281	-0.00147921149032229\\
18210.2307410573	-0.0017072890074432\\
18227.8110120864	-0.00193411008510074\\
18245.3912831156	-0.0021553327840013\\
18262.9715541448	-0.00236673031190815\\
18286.411915517	-0.0026263796659168\\
18309.8522768893	-0.00285265343988584\\
18333.2926382615	-0.00303792833347716\\
18356.7329996337	-0.00317598460315991\\
18380.173361006	-0.00326221221261659\\
18403.6137223782	-0.00329376217527827\\
18427.0540837505	-0.00326963812941574\\
18450.4944451227	-0.00319072521218933\\
18473.9348064949	-0.00305975541951707\\
18497.3751678672	-0.00288121078107255\\
18520.8155292394	-0.00266116777054495\\
18544.2558906116	-0.00240708834446833\\
18561.8361616408	-0.00219932674509548\\
18579.41643267	-0.00198118421492547\\
18596.9967036992	-0.00175679955082239\\
18614.5769747284	-0.00153042270630186\\
18632.1572457575	-0.00130633382351794\\
18649.7375167867	-0.00108876184037811\\
18661.4576974728	-0.000949365567354957\\
18673.177878159	-0.000815856617485192\\
18684.8980588451	-0.000689349321815146\\
18696.6182395312	-0.000570897465252642\\
18708.3384202173	-0.000461485448615252\\
18720.0586009034	-0.000362020047838695\\
18725.9186912465	-0.000316277591548625\\
18731.7787815895	-0.000273322839292687\\
18737.6388719326	-0.000233244052445928\\
18743.4989622757	-0.00019612335432281\\
18749.3590526187	-0.000162036560084215\\
18755.2191429618	-0.000131053020082639\\
18761.0792333048	-0.000103235476965918\\
18766.9393236479	-7.86399368576844e-05\\
18772.799413991	-5.73155548495061e-05\\
18778.659504334	-3.93045350541097e-05\\
18784.5195946771	-2.46420454223196e-05\\
18790.3796850201	-1.33561474888933e-05\\
18796.2397753632	-5.46774119205874e-06\\
18802.0998657063	-9.90524872404033e-07\\
18807.9599560493	6.90294731979261e-08\\
18813.8200463924	-2.28831436094514e-06\\
18819.6801367354	-8.054562652008e-06\\
18825.5402270785	-1.72145128995256e-05\\
18831.4003174216	-2.97457902099721e-05\\
18837.2604077646	-4.56188988320399e-05\\
18843.1204981077	-6.47972887199579e-05\\
18848.9805884507	-8.72374369815938e-05\\
18854.8406787938	-0.000112888944024474\\
18860.7007691369	-0.000141694644190633\\
18866.5608594799	-0.000173590730640519\\
18872.420949823	-0.000208506894207686\\
18878.281040166	-0.000246366475928381\\
18884.1411305091	-0.000287086632907327\\
18890.0012208522	-0.000330578517176891\\
18895.8613111952	-0.000376747467148975\\
18907.5814918813	-0.000476710083292226\\
18919.3016725674	-0.000586108944346025\\
18931.0218532536	-0.000703999708305402\\
18942.7420339397	-0.000829367304165914\\
18954.4622146258	-0.000961134596530641\\
18966.1823953119	-0.0010981715650256\\
18983.7626663411	-0.00131103170657061\\
19001.3429373703	-0.00152901138428191\\
19018.9232083995	-0.0017479307907947\\
19036.5034794286	-0.00196359942621619\\
19054.0837504578	-0.00217189575563247\\
19071.664021487	-0.00236884542259296\\
19095.1043828592	-0.00260735083508344\\
19118.5447442315	-0.00281099647918245\\
19141.9851056037	-0.00297292514471675\\
19165.4254669759	-0.00308770287047421\\
19188.8658283482	-0.0031514991531111\\
19212.3061897204	-0.00316221268620306\\
19235.7465510927	-0.00311953850888947\\
19259.1869124649	-0.00302497442691119\\
19282.6272738371	-0.00288176662031038\\
19306.0676352094	-0.00269479640136526\\
19329.5079965816	-0.00247041206654119\\
19352.9483579539	-0.00221621162719072\\
19370.528628983	-0.00201113198262937\\
19388.1089000122	-0.00179799330299313\\
19405.6891710414	-0.00158083860544934\\
19423.2694420706	-0.00136378112227274\\
19440.8497130997	-0.00115092601489146\\
19452.5698937859	-0.00101334785545283\\
19464.290074472	-0.000880580832073882\\
19476.0102551581	-0.000753736902092438\\
19487.7304358442	-0.00063387681071593\\
19499.4506165303	-0.000522001172095845\\
19511.1707972165	-0.000419042067812197\\
19517.0308875595	-0.000371176816142205\\
19522.8909779026	-0.000325855233192107\\
19528.7510682456	-0.000283171181858921\\
19534.6111585887	-0.000243212896917138\\
19540.4712489318	-0.000206062801140336\\
19546.3313392748	-0.000171797333797657\\
19552.1914296179	-0.000140486791887828\\
19558.0515199609	-0.000112195184460174\\
19563.911610304	-8.69801002897976e-05\\
19569.7717006471	-6.48925892133559e-05\\
19575.6317909901	-4.59770573428382e-05\\
19581.4918813332	-3.02711763954224e-05\\
19587.3519716762	-1.78058073018194e-05\\
19593.2120620193	-8.60493826722153e-06\\
19599.0721523623	-2.68563739496502e-06\\
19604.9322427054	-5.8019972881849e-08\\
19610.7923330485	-7.25230481056351e-07\\
19616.6524233915	-4.6834393550504e-06\\
19622.5125137346	-1.19218545007031e-05\\
19628.3726040776	-2.24227475303034e-05\\
19634.2326944207	-3.61614946600077e-05\\
19640.0927847638	-5.31066321693841e-05\\
19645.9528751068	-7.32199263042771e-05\\
19651.8129654499	-9.64564574642837e-05\\
19657.6730557929	-0.000122764718492462\\
19663.533146136	-0.000152086726849342\\
19669.3932364791	-0.000184358150436336\\
19675.2533268221	-0.000219508446775727\\
19681.1134171652	-0.000257461015274727\\
19686.9735075082	-0.000298133362220343\\
19692.8335978513	-0.00034143727816274\\
19704.5537785374	-0.000435559548498835\\
19716.2739592235	-0.000539015319244037\\
19727.9941399097	-0.000650913919879023\\
19739.7143205958	-0.000770294123527442\\
19751.4345012819	-0.000896132343363573\\
19763.154681968	-0.00102735134846444\\
19774.8748626541	-0.00116282942486231\\
19792.4551336833	-0.00137149493905905\\
19810.0354047125	-0.00158313571932502\\
19827.6156757417	-0.00179370166096843\\
19845.1959467708	-0.00199916885448461\\
19862.7762178	-0.00219561603434716\\
19886.2165791723	-0.00243707738301771\\
19909.6569405445	-0.00264771649768994\\
19933.0973019167	-0.0028204300975864\\
19956.537663289	-0.00294940824394311\\
19979.9780246612	-0.00303032804499619\\
20003.4183860334	-0.00306049685036157\\
20026.8587474057	-0.00303894019690302\\
20050.2991087779	-0.00296643163928709\\
20073.7394701502	-0.00284546356189473\\
20097.1798315224	-0.00268016006111007\\
20120.6201928946	-0.00247613493676936\\
20144.0605542669	-0.00224029967761424\\
20161.6408252961	-0.00204731205105907\\
20179.2210963252	-0.00184456966670262\\
20196.8013673544	-0.00163592296076743\\
20214.3816383836	-0.00142532986583139\\
20231.9619094128	-0.00121678043811988\\
20249.5421804419	-0.00101422096927117\\
20261.2623611281	-0.000884406333322104\\
20272.9825418142	-0.000760048725137236\\
20284.7027225003	-0.000642190833221445\\
20296.4229031864	-0.000531819517057388\\
20308.1430838725	-0.000429857506229384\\
20319.8632645587	-0.000337155648650876\\
20325.7233549017	-0.000294521742025921\\
20331.5834452448	-0.000254485772812664\\
20337.4435355878	-0.000217130812177925\\
20343.3036259309	-0.000182534224030026\\
20349.163716274	-0.000150767503627476\\
20355.023806617	-0.000121896128662621\\
20360.8838969601	-9.59794231308124e-05\\
20366.7439873031	-7.30704342765375e-05\\
20372.6040776462	-5.32158228612711e-05\\
20378.4641679893	-3.64557669931793e-05\\
20384.3242583323	-2.28238796987978e-05\\
20390.1843486754	-1.23471404209158e-05\\
20396.0444390184	-5.04584057474679e-06\\
20401.9045293615	-9.33543273517639e-07\\
20407.7646197046	-1.70573014601601e-08\\
20413.6247100476	-2.29642538632653e-06\\
20419.4848003907	-7.76492678911273e-06\\
20425.3448907337	-1.64090941985982e-05\\
20431.2049810768	-2.82087448914954e-05\\
20437.0650714198	-4.31370260855577e-05\\
20442.9251617629	-6.11604743852655e-05\\
20448.785252106	-8.22390891847683e-05\\
20454.645342449	-0.000106326419876124\\
20460.5054327921	-0.000133369666660804\\
20466.3655231351	-0.000163309794756426\\
20472.2256134782	-0.00019608166175015\\
20478.0857038213	-0.000231614157814808\\
20483.9457941643	-0.000269830358505824\\
20489.8058845074	-0.000310647689780881\\
20495.6659748504	-0.000353978104927964\\
20507.3861555366	-0.000447799778269396\\
20519.1063362227	-0.000550488985691926\\
20530.8265169088	-0.000661164947323027\\
20542.5466975949	-0.000778880002983106\\
20554.266878281	-0.000902627678661026\\
20565.9870589672	-0.00103135124413232\\
20583.5673299963	-0.0012313554878744\\
20601.1476010255	-0.0014362469475744\\
20618.7278720547	-0.0016421065777656\\
20636.3081430839	-0.00184500140783182\\
20653.8884141131	-0.0020410593251268\\
20671.4686851422	-0.00222654261847336\\
20694.9090465145	-0.00245133686907901\\
20718.3494078867	-0.00264349171910218\\
20741.7897692589	-0.00279652872353766\\
20765.2301306312	-0.00290529968183118\\
20788.6704920034	-0.00296615847732917\\
20812.1108533757	-0.00297708199229692\\
20835.5512147479	-0.00293773609248631\\
20858.9915761201	-0.00284948453350451\\
20882.4319374924	-0.00271534056849344\\
20905.8722988646	-0.00253986296944102\\
20929.3126602368	-0.00232900004545486\\
20952.7530216091	-0.00208988699101465\\
20970.3332926383	-0.00189684553715358\\
20987.9135636674	-0.00169611399452162\\
21005.4938346966	-0.0014915046565848\\
21023.0741057258	-0.0012868998843836\\
21040.654376755	-0.00108617822644426\\
21052.3745574411	-0.000956403497428226\\
21064.0947381272	-0.000831139163906947\\
21075.8149188133	-0.000711438520904894\\
21087.5350994994	-0.000598307129074585\\
21099.2552801856	-0.000492694339589269\\
21110.9754608717	-0.000395485299783425\\
21116.8355512147	-0.000350289394060698\\
21122.6956415578	-0.000307493506831298\\
21128.5557319009	-0.000267187008242536\\
21134.4158222439	-0.000229453971963535\\
21140.275912587	-0.000194372998928431\\
21146.13600293	-0.000162017052727458\\
21151.9960932731	-0.000132453306964239\\
21157.8561836162	-0.000105743004929365\\
21163.7162739592	-8.19413318335964e-05\\
21169.5763643023	-6.1097299924636e-05\\
21175.4364546453	-4.32536466786409e-05\\
21181.2965449884	-2.84467463083629e-05\\
21187.1566353315	-1.6706534756168e-05\\
21193.0167256745	-8.05644832862293e-06\\
21198.8768160176	-2.51337610274945e-06\\
21204.7369063606	-8.76261961876233e-08\\
21210.5969967037	-7.82905968610599e-07\\
21216.4570870468	-4.5963161948821e-06\\
21222.3171773898	-1.1518359214173e-05\\
21228.1772677329	-2.15329610357304e-05\\
21234.0373580759	-3.46175073520181e-05\\
21239.897448419	-5.07428933736981e-05\\
21245.7575387621	-6.98735873809977e-05\\
21251.6176291051	-9.19677078558858e-05\\
21257.4777194482	-0.000116977114014739\\
21263.3378097912	-0.000144847509567032\\
21269.1979001343	-0.000175518559451495\\
21275.0579904774	-0.000208924019328482\\
21280.9180808204	-0.000244991877513645\\
21286.7781711635	-0.000283644509092519\\
21292.6382615065	-0.000324798841846169\\
21304.3584421926	-0.000414254161034963\\
21316.0786228788	-0.000512591326913035\\
21327.7988035649	-0.00061896914622417\\
21339.518984251	-0.000732478899955249\\
21351.2391649371	-0.000852152069549209\\
21362.9593456232	-0.000976968564233929\\
21374.6795263094	-0.00110586537972965\\
21392.2597973385	-0.00130445474806615\\
21409.8400683677	-0.00150595524508733\\
21427.4203393969	-0.00170651903154536\\
21445.0006104261	-0.00190231942998053\\
21462.5808814553	-0.0020896236641936\\
21486.0212428275	-0.00232001175123109\\
21509.4616041997	-0.00252119409978788\\
21532.901965572	-0.00268638046715559\\
21556.3423269442	-0.00281000368688873\\
21579.7826883164	-0.00288790627088031\\
21603.2230496887	-0.00291747924063\\
21626.6634110609	-0.00289774855537648\\
21650.1037724332	-0.00282940627672859\\
21673.5441338054	-0.00271478547671585\\
21696.9844951776	-0.00255777979373768\\
21720.4248565499	-0.00236371040728484\\
21743.8652179221	-0.00213914497301313\\
21761.4454889513	-0.00195524304155938\\
21779.0257599805	-0.00176193957099966\\
21796.6060310096	-0.00156290913142642\\
21814.1863020388	-0.00136193264305908\\
21831.766573068	-0.00116282540728234\\
21849.3468440972	-0.000969364536003562\\
21861.0670247833	-0.000845345067356821\\
21872.7872054694	-0.000726513609956042\\
21884.5073861555	-0.000613870769674534\\
21896.2275668417	-0.000508364338302483\\
21907.9477475278	-0.000410881300080587\\
21919.6679282139	-0.000322240356207789\\
21925.5280185569	-0.000281471133043174\\
21931.3881089	-0.000243185030209084\\
21937.2481992431	-0.000207462215150681\\
21943.1082895861	-0.000174377412194723\\
21948.9683799292	-0.000143999745620105\\
21954.8284702722	-0.000116392594641658\\
21960.6885606153	-9.16134605594766e-05\\
21966.5486509584	-6.97138463841963e-05\\
21972.4087413014	-5.07391491756901e-05\\
21978.2688316445	-3.47285653200621e-05\\
21984.1289219875	-2.17150089512198e-05\\
21989.9890123306	-1.1725043677168e-05\\
21995.8491026737	-4.77882775783598e-06\\
22001.7091930167	-8.90072845246499e-07\\
22007.5692833598	-6.60163722716593e-08\\
22013.4293737028	-2.30740764265196e-06\\
22019.2894640459	-7.60850765080387e-06\\
22025.149554389	-1.59571026288878e-05\\
22031.009644732	-2.73345312889737e-05\\
22036.8697350751	-4.17157256969613e-05\\
22042.7298254181	-5.90692656961301e-05\\
22048.5899157612	-7.93574467543084e-05\\
22054.4500061043	-0.000102536361090831\\
22060.3100964473	-0.000128555991917456\\
22066.1701867904	-0.000157360320569037\\
22072.0302771334	-0.000188887446328079\\
22077.8903674765	-0.000223069718647253\\
22083.7504578196	-0.000259833881528768\\
22089.6105481626	-0.000299101229714281\\
22095.4706385057	-0.00034078777638871\\
22107.1908191918	-0.00043105719387192\\
22118.9109998779	-0.000529871669546716\\
22130.631180564	-0.00063638876178975\\
22142.3513612502	-0.000749701208001817\\
22154.0715419363	-0.000868844626112565\\
22165.7917226224	-0.000992805695996385\\
22183.3719936516	-0.00118546666293637\\
22200.9522646807	-0.00138291026549748\\
22218.5325357099	-0.00158136882234624\\
22236.1128067391	-0.00177705763348737\\
22253.6930777683	-0.00196624695904826\\
22271.2733487975	-0.002145332911633\\
22294.7137101697	-0.00236254086203783\\
22318.1540715419	-0.00254841778113761\\
22341.5944329142	-0.00269668955074039\\
22365.0347942864	-0.00280235687916964\\
22388.4751556586	-0.00286186316254732\\
22411.9155170309	-0.00287321353214946\\
22435.3558784031	-0.00283604110599474\\
22458.7962397754	-0.00275161824591257\\
22482.2366011476	-0.00262281247894199\\
22505.6769625198	-0.00245398860755634\\
22529.1173238921	-0.00225086034869867\\
22552.5576852643	-0.00202029654036547\\
22570.1379562935	-0.00183402866461428\\
22587.7182273227	-0.00164024049670971\\
22605.2984983518	-0.00144261668501902\\
22622.878769381	-0.00124491323039203\\
22640.4590404102	-0.00105088602549815\\
22652.1792210963	-0.000925401425474034\\
22663.8994017824	-0.000804250644338998\\
22675.6195824686	-0.000688456230458653\\
22687.3397631547	-0.000578995122488002\\
22699.0599438408	-0.000476790397472854\\
22710.7801245269	-0.000382703476811251\\
22716.64021487	-0.00033895472741485\\
22722.500305213	-0.000297526855936778\\
22728.3603955561	-0.000258507057649185\\
22734.2204858992	-0.000221977418955338\\
22740.0805762422	-0.000188014744367536\\
22745.9406665853	-0.000156690394694511\\
22751.8007569283	-0.000128070136766963\\
22757.6608472714	-0.000102214005008151\\
22763.5209376145	-7.91761751625779e-05\\
22769.3810279575	-5.90048504323481e-05\\
22775.2411183006	-4.17421602660924e-05\\
22781.1012086436	-2.74240720093062e-05\\
22786.9612989867	-1.60803156061958e-05\\
22792.8213893298	-7.73432150630204e-06\\
22798.6814796728	-2.40317190797747e-06\\
22804.5415700159	-9.75654389382338e-08\\
22810.4016603589	-8.21795347178264e-07\\
22816.261750702	-4.57374124570659e-06\\
22822.121841045	-1.13448744295206e-05\\
22827.9819313881	-2.11202767477588e-05\\
22833.8420217312	-3.38786729934728e-05\\
22839.7021120742	-4.95924767370388e-05\\
22845.5622024173	-6.82278495082941e-05\\
22851.4222927603	-8.97447731970047e-05\\
22857.2823831034	-0.000114097135518203\\
22863.1424734465	-0.000141232828365827\\
22869.0025637895	-0.000171093858826122\\
22874.8626541326	-0.000203616472647206\\
22880.7227444756	-0.000238731289862733\\
22886.5828348187	-0.000276363452316818\\
22892.4429251618	-0.000316432782750981\\
22904.1631058479	-0.000403536675847783\\
22915.883286534	-0.00049930191015217\\
22927.6034672201	-0.000602914317891804\\
22939.3236479062	-0.000713493527961728\\
22951.0438285924	-0.000830100429676494\\
22962.7640092785	-0.000951745131603513\\
22974.4841899646	-0.00107739534806444\\
22992.0644609938	-0.00127104229099939\\
23009.6447320229	-0.00146760582281903\\
23027.2250030521	-0.00166334005929498\\
23044.8052740813	-0.00185451617663497\\
23062.3855451105	-0.00203749332194539\\
23085.8259064827	-0.00226271870021161\\
23109.266267855	-0.00245958960357174\\
23132.7066292272	-0.0026214559182546\\
23156.1469905994	-0.00274285271512935\\
23179.5873519717	-0.00281968433563411\\
23203.0277133439	-0.00284936218507815\\
23226.4680747162	-0.00283089160029641\\
23249.9084360884	-0.00276490487818962\\
23273.3487974606	-0.002653639369865\\
23296.7891588329	-0.00250086139780925\\
23320.2295202051	-0.00231173858207379\\
23343.6698815773	-0.00209266490066311\\
23361.2501526065	-0.00191312729735072\\
23378.8304236357	-0.00172430801620584\\
23396.4106946649	-0.00152979955553598\\
23413.9909656941	-0.00133330200090315\\
23431.5712367232	-0.00113855262994018\\
23449.1515077524	-0.000949254823742937\\
23460.8716884385	-0.000827868621857204\\
23472.5918691246	-0.000711534937569924\\
23484.3120498108	-0.000601237361657966\\
23496.0322304969	-0.000497908301027104\\
23507.752411183	-0.000402421095557273\\
23519.4725918691	-0.000315582635360732\\
23525.3326822122	-0.000275638977228373\\
23531.1927725552	-0.000238126540927679\\
23537.0528628983	-0.000203124563407944\\
23542.9129532414	-0.000170706963778393\\
23548.7730435844	-0.000140942187121089\\
23554.6331339275	-0.000113893059913721\\
23560.4932242705	-8.96166573390238e-05\\
23566.3533146136	-6.81641827723735e-05\\
23572.2134049567	-4.95808596941853e-05\\
23578.0734952997	-3.39058362646826e-05\\
23583.9335856428	-2.1172102755075e-05\\
23589.7936759858	-1.14064220109323e-05\\
23595.6537663289	-4.62927309620972e-06\\
23601.513856672	-8.54808233626689e-07\\
23607.373947015	-9.0823133571775e-08\\
23613.2340373581	-2.33874077388739e-06\\
23619.0941277011	-7.59360866299011e-06\\
23624.9542180442	-1.58441095928798e-05\\
23630.8143083873	-2.7072585857694e-05\\
23636.6743987303	-4.12550768870749e-05\\
23642.5344890734	-5.83613702124949e-05\\
23648.3945794164	-7.83550656602712e-05\\
23654.2546697595	-0.000101193652633783\\
23660.1147601026	-0.000126828600321508\\
23665.9748504456	-0.00015520546063389\\
23671.8349407887	-0.000186263983670552\\
23677.6950311317	-0.000219938245436648\\
23683.5551214748	-0.000256156787581085\\
23689.4152118178	-0.000294842768824362\\
23695.2753021609	-0.000335914127787074\\
23706.995482847	-0.000424859686585813\\
23718.7156635331	-0.000522239440905372\\
23730.4358442193	-0.000627228047288775\\
23742.1560249054	-0.000738935829724721\\
23753.8762055915	-0.000856416311212569\\
23765.5963862776	-0.000978674225344446\\
23783.1766573068	-0.00116874445146766\\
23800.756928336	-0.00136360727140985\\
23818.3371993652	-0.0015595527887299\\
23835.9174703943	-0.00175285082931565\\
23853.4977414235	-0.00193982189351401\\
23871.0780124527	-0.00211690714030613\\
23894.5183738249	-0.0023318525077528\\
23917.9587351972	-0.00251599650596467\\
23941.3990965694	-0.00266311657780403\\
23964.8394579416	-0.00276824144011184\\
23988.2798193139	-0.00282781894062616\\
24000	-0.0028398359951788\\
};
\end{axis}
\end{tikzpicture}%

%%%% This file was created by matlab2tikz.
%
%The latest updates can be retrieved from
%  http://www.mathworks.com/matlabcentral/fileexchange/22022-matlab2tikz-matlab2tikz
%where you can also make suggestions and rate matlab2tikz.
%
\definecolor{mycolor1}{rgb}{0.00000,0.44700,0.74100}%
\definecolor{mycolor2}{rgb}{0.85000,0.32500,0.09800}%
%
\begin{tikzpicture}

\begin{axis}[%
width=4.521in,
height=3.548in,
at={(0.758in,0.499in)},
scale only axis,
xmin=-0.0025,
xmax=0.0025,
ymin=-0.2,
ymax=1,
axis background/.style={fill=white}
]
\addplot [color=mycolor1, forget plot]
  table[row sep=crcr]{%
-0.00208333333333333	-0.00318309886183787\\
-0.00206249999999997	-0.00318317492734366\\
-0.00204166666666672	-0.00311909156929424\\
-0.00202083333333336	-0.00299081106446764\\
-0.002	-0.00279956172779527\\
-0.00197916666666664	-0.00254783933014335\\
-0.00195833333333328	-0.00223938409131019\\
-0.00193750000000004	-0.00187913323343325\\
-0.00191666666666668	-0.00147314954171185\\
-0.00189583333333332	-0.00102852684019938\\
-0.00185416666666671	-5.61774650087443e-05\\
-0.00177083333333339	0.00195666421959917\\
-0.00175000000000003	0.00241545663842524\\
-0.00172916666666667	0.00283652879109197\\
-0.00170833333333331	0.00321058443246536\\
-0.00168749999999995	0.00352902537568234\\
-0.00166666666666671	0.00378413364320329\\
-0.00164583333333335	0.00396924098084706\\
-0.00162499999999999	0.00407888230089704\\
-0.00160416666666663	0.00410892985197497\\
-0.00158333333333338	0.00405670520761037\\
-0.00156250000000002	0.00392106651788271\\
-0.00154166666666666	0.00370246887363823\\
-0.0015208333333333	0.00340299608403061\\
-0.00149999999999995	0.0030263626579945\\
-0.0014791666666667	0.00257788530046443\\
-0.00145833333333334	0.00206442377572313\\
-0.00143749999999998	0.00149429154376235\\
-0.00139583333333337	0.000223797745730003\\
-0.00129166666666669	-0.00314668345950841\\
-0.00127083333333333	-0.00374732295256697\\
-0.00124999999999997	-0.00429196845667068\\
-0.00122916666666661	-0.00476798876485895\\
-0.00120833333333337	-0.00516364914617817\\
-0.00118750000000001	-0.00546835616486996\\
-0.00116666666666665	-0.00567288882739969\\
-0.00114583333333329	-0.00576961169154977\\
-0.00112500000000004	-0.00575266580402922\\
-0.00110416666666668	-0.00561813364122188\\
-0.00108333333333333	-0.00536417460743954\\
-0.00106249999999997	-0.0049911280907049\\
-0.00104166666666672	-0.00450158158078551\\
-0.00102083333333336	-0.00390040190991692\\
-0.001	-0.00319472827441547\\
-0.000979166666666642	-0.00239392632536617\\
-0.000958333333333283	-0.00150950326827604\\
-0.000916666666666677	0.000454246434654526\\
-0.000812499999999994	0.00567989546351644\\
-0.000791666666666635	0.00661879347518046\\
-0.000770833333333387	0.00747281077017503\\
-0.000750000000000028	0.00822102909442923\\
-0.00072916666666667	0.00884328453400929\\
-0.000708333333333311	0.00932050634240345\\
-0.000687499999999952	0.00963504613361477\\
-0.000666666666666704	0.00977099196810116\\
-0.000645833333333345	0.00971446202866277\\
-0.000624999999999987	0.00945387284344712\\
-0.000604166666666628	0.00898017735982271\\
-0.00058333333333338	0.00828706860093087\\
-0.000562500000000021	0.00737114513983217\\
-0.000541666666666663	0.00623203519655136\\
-0.000520833333333304	0.00487247679202218\\
-0.000499999999999945	0.00329835206993767\\
-0.000479166666666697	0.00151867461190636\\
-0.000458333333333338	-0.000454470688504793\\
-0.00043749999999998	-0.00260603487341049\\
-0.000395833333333373	-0.00737036774129618\\
-0.000354166666666655	-0.0126015609750331\\
-0.000229166666666614	-0.0289336924620972\\
-0.000187500000000007	-0.0338042628617465\\
-0.000166666666666648	-0.0360019245108092\\
-0.00014583333333329	-0.0380065351526726\\
-0.000125000000000042	-0.0397946281998789\\
-0.000104166666666683	-0.0413451468692532\\
-8.33333333333242e-05	-0.0426397415309975\\
-6.24999999999654e-05	-0.043663029550196\\
-4.16666666667176e-05	-0.0444028136048185\\
-2.08333333333588e-05	-0.0448502551012079\\
0	0.955\\
2.08333333333588e-05	-0.0448502551012079\\
4.16666666667176e-05	-0.0444028136048185\\
6.24999999999654e-05	-0.043663029550196\\
8.33333333333242e-05	-0.0426397415309975\\
0.000104166666666683	-0.0413451468692532\\
0.000125000000000042	-0.0397946281998789\\
0.00014583333333329	-0.0380065351526726\\
0.000166666666666648	-0.0360019245108092\\
0.000187500000000007	-0.0338042628617465\\
0.000229166666666614	-0.0289336924620972\\
0.000270833333333331	-0.0236175509297717\\
0.000375000000000014	-0.0099398164461062\\
0.000416666666666621	-0.0049181582154173\\
0.00043749999999998	-0.00260603487341049\\
0.000458333333333338	-0.000454470688504793\\
0.000479166666666697	0.00151867461190636\\
0.000499999999999945	0.00329835206993767\\
0.000520833333333304	0.00487247679202218\\
0.000541666666666663	0.00623203519655136\\
0.000562500000000021	0.00737114513983217\\
0.00058333333333338	0.00828706860093087\\
0.000604166666666628	0.00898017735982271\\
0.000624999999999987	0.00945387284344712\\
0.000645833333333345	0.00971446202866277\\
0.000666666666666704	0.00977099196810116\\
0.000687499999999952	0.00963504613361477\\
0.000708333333333311	0.00932050634240345\\
0.00072916666666667	0.00884328453400929\\
0.000750000000000028	0.00822102909442923\\
0.000770833333333387	0.00747281077017503\\
0.000791666666666635	0.00661879347518046\\
0.000833333333333353	0.00467744641894319\\
0.000937500000000036	-0.000554984574026918\\
0.000958333333333283	-0.00150950326827604\\
0.000979166666666642	-0.00239392632536617\\
0.001	-0.00319472827441547\\
0.00102083333333336	-0.00390040190991692\\
0.00104166666666672	-0.00450158158078551\\
0.00106249999999997	-0.0049911280907049\\
0.00108333333333333	-0.00536417460743954\\
0.00110416666666668	-0.00561813364122188\\
0.00112500000000004	-0.00575266580402922\\
0.00114583333333329	-0.00576961169154977\\
0.00116666666666665	-0.00567288882739969\\
0.00118750000000001	-0.00546835616486996\\
0.00120833333333337	-0.00516364914617817\\
0.00122916666666661	-0.00476798876485895\\
0.00124999999999997	-0.00429196845667068\\
0.00127083333333333	-0.00374732295256697\\
0.00129166666666669	-0.00314668345950841\\
0.0013333333333333	-0.00183090131978936\\
0.00141666666666662	0.000877137131158112\\
0.00143749999999998	0.00149429154376235\\
0.00145833333333334	0.00206442377572313\\
0.0014791666666667	0.00257788530046443\\
0.00149999999999995	0.0030263626579945\\
0.0015208333333333	0.00340299608403061\\
0.00154166666666666	0.00370246887363823\\
0.00156250000000002	0.00392106651788271\\
0.00158333333333338	0.00405670520761037\\
0.00160416666666663	0.00410892985197497\\
0.00162499999999999	0.00407888230089704\\
0.00164583333333335	0.00396924098084706\\
0.00166666666666671	0.00378413364320329\\
0.00168749999999995	0.00352902537568234\\
0.00170833333333331	0.00321058443246536\\
0.00172916666666667	0.00283652879109197\\
0.00175000000000003	0.00241545663842524\\
0.00179166666666664	0.00146995464960564\\
0.00189583333333332	-0.00102852684019938\\
0.00191666666666668	-0.00147314954171185\\
0.00193750000000004	-0.00187913323343325\\
0.00195833333333328	-0.00223938409131019\\
0.00197916666666664	-0.00254783933014335\\
0.002	-0.00279956172779527\\
0.00202083333333336	-0.00299081106446764\\
0.00204166666666672	-0.00311909156929424\\
0.00206249999999997	-0.00318317492734366\\
0.00208333333333333	-0.00318309886183787\\
};
\addplot [color=mycolor2, forget plot]
  table[row sep=crcr]{%
-0.00208333333333333	0.00318309886183791\\
-0.0020625	0.00318317492734361\\
-0.00204166666666666	0.00311909156929429\\
-0.00202083333333333	0.00299081106446766\\
-0.002	0.00279956172779523\\
-0.00197916666666666	0.00254783933014338\\
-0.00195833333333333	0.00223938409131023\\
-0.0019375	0.00187913323343322\\
-0.00191666666666666	0.00147314954171188\\
-0.00189583333333333	0.00102852684019933\\
-0.00185416666666666	5.61774650087374e-05\\
-0.00177083333333333	-0.00195666421959916\\
-0.00175	-0.00241545663842527\\
-0.00172916666666666	-0.00283652879109197\\
-0.00170833333333333	-0.00321058443246536\\
-0.0016875	-0.00352902537568231\\
-0.00166666666666666	-0.00378413364320329\\
-0.00164583333333333	-0.00396924098084705\\
-0.001625	-0.00407888230089703\\
-0.00160416666666666	-0.00410892985197494\\
-0.00158333333333333	-0.0040567052076104\\
-0.0015625	-0.00392106651788268\\
-0.00154166666666666	-0.00370246887363825\\
-0.00152083333333333	-0.00340299608403059\\
-0.0015	-0.0030263626579945\\
-0.00147916666666666	-0.00257788530046447\\
-0.00145833333333333	-0.0020644237757231\\
-0.0014375	-0.00149429154376232\\
-0.00139583333333333	-0.000223797745729983\\
-0.00129166666666667	0.00314668345950841\\
-0.00127083333333333	0.00374732295256701\\
-0.00125	0.0042919684566707\\
-0.00122916666666667	0.00476798876485895\\
-0.00120833333333333	0.00516364914617822\\
-0.0011875	0.00546835616486993\\
-0.00116666666666667	0.00567288882739972\\
-0.00114583333333333	0.00576961169154982\\
-0.001125	0.00575266580402919\\
-0.00110416666666667	0.00561813364122191\\
-0.00108333333333333	0.00536417460743957\\
-0.0010625	0.00499112809070489\\
-0.00104166666666667	0.00450158158078553\\
-0.00102083333333333	0.0039004019099169\\
-0.001	0.00319472827441542\\
-0.00097916666666667	0.00239392632536613\\
-0.000958333333333332	0.00150950326827601\\
-0.00091666666666667	-0.000454246434654471\\
-0.000812500000000001	-0.00567989546351649\\
-0.00079166666666667	-0.00661879347518047\\
-0.000770833333333332	-0.00747281077017498\\
-0.000750000000000001	-0.00822102909442928\\
-0.00072916666666667	-0.00884328453400932\\
-0.000708333333333332	-0.00932050634240345\\
-0.000687500000000001	-0.00963504613361472\\
-0.00066666666666667	-0.00977099196810118\\
-0.000645833333333332	-0.00971446202866279\\
-0.000625000000000001	-0.00945387284344715\\
-0.00060416666666667	-0.00898017735982275\\
-0.000583333333333332	-0.00828706860093084\\
-0.0005625	-0.0073711451398322\\
-0.000541666666666669	-0.0062320351965514\\
-0.000520833333333331	-0.00487247679202217\\
-0.0005	-0.0032983520699377\\
-0.000479166666666669	-0.0015186746119064\\
-0.000458333333333331	0.000454470688504793\\
-0.0004375	0.0026060348734105\\
-0.000395833333333331	0.0073703677412962\\
-0.000354166666666669	0.0126015609750331\\
-0.000229166666666669	0.0289336924620972\\
-0.0001875	0.0338042628617465\\
-0.000166666666666669	0.0360019245108092\\
-0.000145833333333331	0.0380065351526726\\
-0.000125	0.0397946281998788\\
-0.000104166666666669	0.0413451468692531\\
-8.33333333333311e-05	0.0426397415309974\\
-6.25000000000001e-05	0.0436630295501959\\
-4.1666666666669e-05	0.0444028136048185\\
-2.0833333333331e-05	0.0448502551012078\\
0	0.045\\
2.0833333333331e-05	0.0448502551012078\\
4.1666666666669e-05	0.0444028136048185\\
6.25000000000001e-05	0.0436630295501959\\
8.33333333333311e-05	0.0426397415309974\\
0.000104166666666669	0.0413451468692531\\
0.000125	0.0397946281998788\\
0.000145833333333331	0.0380065351526726\\
0.000166666666666669	0.0360019245108092\\
0.0001875	0.0338042628617465\\
0.000229166666666669	0.0289336924620972\\
0.000270833333333331	0.0236175509297717\\
0.000375	0.00993981644610623\\
0.000416666666666669	0.00491815821541732\\
0.0004375	0.0026060348734105\\
0.000458333333333331	0.000454470688504793\\
0.000479166666666669	-0.0015186746119064\\
0.0005	-0.0032983520699377\\
0.000520833333333331	-0.00487247679202217\\
0.000541666666666669	-0.0062320351965514\\
0.0005625	-0.0073711451398322\\
0.000583333333333332	-0.00828706860093084\\
0.00060416666666667	-0.00898017735982275\\
0.000625000000000001	-0.00945387284344715\\
0.000645833333333332	-0.00971446202866279\\
0.00066666666666667	-0.00977099196810118\\
0.000687500000000001	-0.00963504613361472\\
0.000708333333333332	-0.00932050634240345\\
0.00072916666666667	-0.00884328453400932\\
0.000750000000000001	-0.00822102909442928\\
0.000770833333333332	-0.00747281077017498\\
0.00079166666666667	-0.00661879347518047\\
0.000833333333333332	-0.0046774464189432\\
0.000937500000000001	0.000554984574026959\\
0.000958333333333332	0.00150950326827601\\
0.00097916666666667	0.00239392632536613\\
0.001	0.00319472827441542\\
0.00102083333333333	0.0039004019099169\\
0.00104166666666667	0.00450158158078553\\
0.0010625	0.00499112809070489\\
0.00108333333333333	0.00536417460743957\\
0.00110416666666667	0.00561813364122191\\
0.001125	0.00575266580402919\\
0.00114583333333333	0.00576961169154982\\
0.00116666666666667	0.00567288882739972\\
0.0011875	0.00546835616486993\\
0.00120833333333333	0.00516364914617822\\
0.00122916666666667	0.00476798876485895\\
0.00125	0.0042919684566707\\
0.00127083333333333	0.00374732295256701\\
0.00129166666666667	0.00314668345950841\\
0.00133333333333333	0.00183090131978936\\
0.00141666666666667	-0.000877137131158105\\
0.0014375	-0.00149429154376232\\
0.00145833333333333	-0.0020644237757231\\
0.00147916666666666	-0.00257788530046447\\
0.0015	-0.0030263626579945\\
0.00152083333333333	-0.00340299608403059\\
0.00154166666666666	-0.00370246887363825\\
0.0015625	-0.00392106651788268\\
0.00158333333333333	-0.0040567052076104\\
0.00160416666666666	-0.00410892985197494\\
0.001625	-0.00407888230089703\\
0.00164583333333333	-0.00396924098084705\\
0.00166666666666666	-0.00378413364320329\\
0.0016875	-0.00352902537568231\\
0.00170833333333333	-0.00321058443246536\\
0.00172916666666666	-0.00283652879109197\\
0.00175	-0.00241545663842527\\
0.00179166666666666	-0.00146995464960569\\
0.00189583333333333	0.00102852684019933\\
0.00191666666666666	0.00147314954171188\\
0.0019375	0.00187913323343322\\
0.00195833333333333	0.00223938409131023\\
0.00197916666666666	0.00254783933014338\\
0.002	0.00279956172779523\\
0.00202083333333333	0.00299081106446766\\
0.00204166666666666	0.00311909156929429\\
0.0020625	0.00318317492734361\\
0.00208333333333333	0.00318309886183791\\
};
\end{axis}
\end{tikzpicture}%

%\subsubsection{Modification de la fréquence de coupure}

% TODO


\section{Démodulateur de fréquence adapté à la norme V21}

\subsection{Contexte de synchronisation idéale}
\label{synchro-ideale}

Nous allons ici introduire une nouvelle méthode de démodulation plus adaptée à la norme V21.
La multiplication du signal avec le cosinus correspondant (par exemple) au bit 1 donne, en fonction du temps, une mesure de la synchronisation entre le signal et ce cosinus: plus le résultat est proche de 1, plus les signaux sont synchronisés à cet instant, et donc plus le signal est susceptible d'être un bit 1.
On a:
\[ 
\int_{0}^{T_s} cos^2(2 \pi F_0 t + \phi_0) \, \mathrm{d}t  = \frac{1} {8 \pi F_0} sin(2(2 \pi F_0 T + phi_0))+4 \pi F_0 T - sin(2 \phi_0) 

\int_{0}^{T_s} cos^2(2 \pi F_0 t + \phi_1) \, \mathrm{d}t  = \frac{1} {8 \pi F_0} sin(2(2 \pi F_0 T + phi_1))+4 \pi F_0 T - sin(2 \phi_1) 

\int_{0}^{T_s} cos(2 \pi F_0 t + \phi_0) cos(2 \pi F_0 t + \phi_1)\, \mathrm{d}t  = \frac{1} {8 \pi F_0} sin(4 \pi F_0 T + phi_0 + phi_1)+4 \pi F_0 T cos(phi_0 - phi_1) - sin (\phi_0 + \phi_1) 

\]


\begin{center}
\begin{figure}[H]
	\centering
	\begin{tikzpicture}
        \begin{axis}[
            ymax=1.1, ymin=-1.1,
            ylabel={Amplitude},
            xlabel={t [s]},
            samples=2000,
        ]
        \addplot[domain=0:100, red] {cos(2*pi*x)};
        \addplot[domain=0:100, blue] {cos(4*pi*x)};
        \addplot[domain=0:100, black, ultra thick] {cos(2*pi*x)*cos(4*pi*x)};
        \legend{cos(2t), cos(4t), produit}
        \end{axis}
        \node[] at (axis cs:55,1) {synchronises};
        \node[] at (axis cs:29,1) {desynchronises};
    \end{tikzpicture}
	\caption{Synchronisation entre deux cosinus de fréquences différentes}
	\label{fig:synchronisation}
\end{figure}
\end{center}

On fait ensuite une accumulation (une moyenne en quelque sorte) de ces mesures sur une période $T_s$ en intégrant sur $T_s$, pour avoir une idée de la synchronisation avec le signal d'un bit 1 et d'un bit 0 à chaque instant échantilloné.

% TODO montrer l'intégrale pas juste le produit

\begin{center}
\begin{figure}[H]
	\centering
	\begin{tikzpicture}
        \begin{axis}[
            ymax=1.1, ymin=-1.1,
            ylabel={Amplitude},
            xlabel={t [s]},
            samples=1000,
        ]
        \addplot[domain=0:200, red, thick] {cos(2*pi*x) * (x<30*pi) + cos(4*pi*x) * (x>=30*pi)};
        \addplot[domain=0:200, blue, thick] {cos(4*pi*x)};
        \addplot[domain=0:200, black, ultra thick] {(cos(2*pi*x) * (x<30*pi) + cos(4*pi*x) * (x>=30*pi)) * cos(4*pi*x)};
        \legend{signal, cos(4t), produit}
        \end{axis}
    \end{tikzpicture}
	\caption{Synchronisation entre le signal obtenu et le signal théorique d'un bit}
\end{figure}
\end{center}

Finalement, pour chaque échantillon temporel (de durée $T_s$), on effectue une comparaison:

$$
\operatorname{reconstitué}(kT_s) = \begin{cases}
    1 & \text{si } \operatorname{sync}_1(kT_s) - \operatorname{sync}_0(kT_s) > 0 \\
    0 & \text{sinon}
\end{cases}
$$

En notant $\operatorname{sync}_b$ le produit pour le bit $b$.

Le signal démodulé avec cette méthode est le suivant:


\begin{figure}
%%%% This file was created by matlab2tikz.
%
%The latest updates can be retrieved from
%  http://www.mathworks.com/matlabcentral/fileexchange/22022-matlab2tikz-matlab2tikz
%where you can also make suggestions and rate matlab2tikz.
%
\definecolor{mycolor1}{rgb}{0.00000,0.44700,0.74100}%
%
\begin{tikzpicture}

\begin{axis}[%
width=4.521in,
height=3.566in,
at={(0.758in,0.481in)},
scale only axis,
xmin=0,
xmax=160,
ymin=-0.1,
ymax=1.1,
axis background/.style={fill=white},
title style={font=\bfseries},
title={Signal démodulé avec synchronisation parfaite}
]
\addplot [color=mycolor1, forget plot]
  table[row sep=crcr]{%
1	1\\
4	1\\
5	0\\
6	1\\
9	1\\
10	0\\
11	1\\
28	1\\
29	0\\
30	1\\
47	1\\
48	0\\
49	1\\
51	1\\
52	0\\
53	0\\
54	1\\
71	1\\
72	0\\
73	1\\
76	1\\
77	0\\
78	1\\
100	1\\
101	0\\
102	1\\
105	1\\
106	0\\
107	1\\
110	1\\
111	0\\
112	1\\
114	1\\
115	0\\
116	1\\
139	1\\
140	0\\
141	1\\
148	1\\
149	0\\
150	1\\
153	1\\
154	0\\
155	1\\
158	1\\
159	0\\
160	1\\
};
\end{axis}
\end{tikzpicture}%
\end{figure}



\subsection{Gestion d'une erreur de synchronisation de phase porteuse}

Un sinus est déphasé d'un quart de phase, comparé à un cosinus de même fréquence et même déphasage.

En rajoutant ces mesures de désynchronisation, on prend en compte les signaux d'entrée qui seraient déphasés: si le signal n'est pas synchronisé avec le cosinus à cause d'un léger déphasage, l'ajout d'une mesure de synchronisation avec ce même cosinus, mais déphasé de $\frac{\pi}{2}$ compensera la faible valeur de synchronisation.

\begin{figure}[H]
	\centering
	\begin{tikzpicture}
		\begin{axis}[ymin=-1.1, ymax=1.1, xlabel={Temps [s]}, ylabel={Amplitude}, samples=1000]
		\addplot[domain=0:100, blue, thick] {cos(2*pi*x)};
		\addplot[domain=0:100, blue, thick, dashed] {sin(2*pi*x)};
		\addplot[domain=0:100, red, thick] {cos(2*pi*x - 20*pi)};
		\legend{cosinus, sinus, signal};
	\end{axis}	
	\end{tikzpicture}
	\caption{Compensation d'un déphasage du signal d'entrée}
	\label{fig:compensation-dephasage}
\end{figure}

Le reste du processus reste le même qu'en \ref{synchro-ideale}



% ne pas écrire en dessous

\begin{center}
    \includegraphics[width=0.125\textwidth]{frog.jpg}
\end{center}
\end{document}
